% Generated mechanically from the frozen chow.tex target.
% Do not edit this generated file; edit the bound locale source instead.

% OK, start here.
%

\setcounter{chapter}{41}
\chapter{Chow 同调与陈类}


\phantomsection
\label{chow-section-phantom}




\section{引言}
\label{chow-section-introduction}

\noindent
本章讨论 Chow 同调群，以及把向量丛的陈类构造为运算 Chow 上同调群中元素的过程
（系数一律取自 $\mathbf{Z}$）。

\medskip\noindent
本章首先给出关键引理 \ref{chow-lemma-milnor-gersten-low-degree} 的一个尽可能简短的代数证明。
我们先定义 Herbrand 商（第 \ref{chow-section-periodic-complexes} 节），并在若干情形下计算它
（第 \ref{chow-section-calculation} 节）。随后证明一些简单的代数引理，说明经修改后存在适当的分解
（第 \ref{chow-section-preparation-tame-symbol} 节）。利用这些结果，我们在第
\ref{chow-section-tame-symbol} 节构造并定义驯顺符号。证明关键引理只需要驯顺符号最基本的性质；
该证明见第 \ref{chow-section-key-lemma} 节。

\medskip\noindent
接着，我们在第 \ref{chow-section-setup} 节引入本章其余部分所采用的基本设置。为了使材料更具挑战性，
我们决定处理一个比通常情形稍为一般的情况。具体而言，假设概形 $X$ 在一个固定的局部诺特基概形上
局部有限型，并且该基概形泛链状且配有一个维数函数。这些假设足以定义 Chow 同调群 $\CH_*(X)$，
以及以陈类同它们作帽积的作用。这表明，对这种基概形上局部有限型的代数栈，我们也应当能够作出同样的定义。

\medskip\noindent
随后我们依照 \cite{F} 的开头几章定义循环、平坦拉回、固有推前和有理等价；但对于所涉及循环的支撑，
我们的处理没有那么精细。

\medskip\noindent
我们与 \cite{F} 的论述方式有所不同：在第 \ref{chow-section-key} 节利用上述关键引理证明一个基本的交换关系。
由此可证，与可逆层取交这一运算可下降到有理等价，并且满足交换律；见第
\ref{chow-section-commutativity} 节。再次应用关键引理可证，有效 Cartier 除子的 Gysin 映射可下降到有理等价；
见第 \ref{chow-section-gysin} 节。证明这些结果以后，就可以直接定义向量丛的陈类，证明其可加性和分裂原理，
引入 Chern 特征与 Todd 类，并陈述 Grothendieck--Riemann--Roch 定理。

\medskip\noindent
本章有两个附录。附录 A（第 \ref{chow-section-appendix-A} 节）讨论驯顺符号的另一种较长构造，
以及与之相应的关键引理证明。最后，在附录 B（第 \ref{chow-section-appendix-chow} 节）中，
我们简要讨论它同凝聚层 $K$-理论的关系，并讨论若干爆破引理。建议读者参阅这两个附录的引言以了解详情。

\medskip\noindent
下一章将重新研究代数闭域上光滑射影簇的 Chow 群 $\CH_*(X)$。利用 \cite{Samuel}、
\cite{ChevalleyI} 和 \cite{ChevalleyII} 中的移动引理，以及 Serre 的 Tor 公式
（见 \cite{Serre_local_algebra} 或 \cite{Serre_algebre_locale}），我们将在 $\CH_*(X)$ 上定义环结构。
参见 Intersection Theory，第 \ref{intersection-section-introduction} 节及以下各节。








\section{周期复形与 Herbrand 商}
\label{chow-section-periodic-complexes}

\noindent
周期复形当然有一个非常一般的概念。可以要求映射具有周期性，也可以要求对象具有周期性。
我们会在需要时补充这些情形；目前只需要以下几种情况。

\begin{definition}
\label{chow-definition-periodic-complex}
设 $R$ 为环。
\begin{enumerate}
\item $R$ 上的一个 {\it $2$-周期复形}由四元组 $(M, N, \varphi, \psi)$ 给出，
其中 $M$、$N$ 是 $R$-模，$\varphi : M \to N$、$\psi : N \to M$ 是 $R$-模同态，并且
$$
\xymatrix{
\ldots \ar[r] &
M \ar[r]^\varphi &
N \ar[r]^\psi &
M \ar[r]^\varphi &
N \ar[r] & \ldots
}
$$
构成一个复形。在这种情形下，定义该复形的{\it 上同调模}为 $R$-模
$$
H^0(M, N, \varphi, \psi) = \Ker(\varphi)/\Im(\psi)
\quad\text{且}\quad
H^1(M, N, \varphi, \psi) = \Ker(\psi)/\Im(\varphi).
$$
如果上同调群均为零，就称这个 $2$-周期复形是{\it 正合的}。
\item $R$ 上的一个 {\it $(2, 1)$-周期复形}由三元组 $(M, \varphi, \psi)$ 给出，
其中 $M$ 是 $R$-模，$\varphi : M \to M$、$\psi : M \to M$ 是 $R$-模同态，并且
$$
\xymatrix{
\ldots \ar[r] &
M \ar[r]^\varphi &
M \ar[r]^\psi &
M \ar[r]^\varphi &
M \ar[r] & \ldots
}
$$
构成一个复形。由于这是 $2$-周期复形的一种特殊情形，我们有它的{\it 上同调模}
$H^0(M, \varphi, \psi)$、$H^1(M, \varphi, \psi)$，以及相应的正合性概念。
\end{enumerate}
\end{definition}

\noindent
下文凡是对 $2$-周期复形证明的结果，都将不再另行说明而直接用于 $(2, 1)$-周期复形。
显然，所有 $2$-周期复形组成一个范畴，其中态射
$(f, g) : (M, N, \varphi, \psi) \to (M', N', \varphi', \psi')$
是满足 $\varphi' \circ f = g \circ \varphi$ 和 $\psi' \circ g = f \circ \psi$ 的一对态射
$f : M \to M'$ 与 $g : N \to N'$。这样得到一个阿贝尔范畴；其中的核与余核如《同调代数》引理 \ref{homology-lemma-cat-chain-abelian} 所述。

\begin{definition}
\label{chow-definition-periodic-length}
设 $(M, N, \varphi, \psi)$ 是环 $R$ 上的一个 $2$-周期复形，并且它的上同调模均具有有限长度。
在这种情形下，定义 $(M, N, \varphi, \psi)$ 的{\it 重数}为整数
$$
e_R(M, N, \varphi, \psi) =
\text{length}_R(H^0(M, N, \varphi, \psi))
-
\text{length}_R(H^1(M, N, \varphi, \psi))
$$
对于 $(2, 1)$-周期复形 $(M, \varphi, \psi)$，将这个整数记为
$e_R(M, \varphi, \psi)$，有时也称它为{\it（加法）Herbrand 商}。
\end{definition}

\noindent
如果 $(M, \varphi, \psi)$ 的上同调群是有限阿贝尔群，则通常把
$$
q(M, \varphi, \psi) =
\frac{\# H^0(M, \varphi, \psi)}{\# H^1(M, \varphi, \psi)}
$$
称为{\it（乘法）Herbrand 商}。换言之，乘法 Herbrand 商就是 $H^0$ 的元素个数除以
$H^1$ 的元素个数。如果 $R$ 是局部环，而且 $R$ 的剩余域含有 $q$ 个元素，那么
$$
q(M, \varphi, \psi) = q^{e_R(M, \varphi, \psi)}
$$

\medskip\noindent
环 $R$ 上 $(2, 1)$-周期复形的一个例子是任意形如 $(M, 0, \psi)$ 的三元组，
其中 $M$ 是 $R$-模，$\psi$ 是 $R$-线性映射。如果 $\psi$ 的核和余核都具有有限长度，则有
\begin{equation}
\label{chow-equation-multiplicity-coker-ker}
e_R(M, 0, \psi) = \text{length}_R(\Coker(\psi)) - \text{length}_R(\Ker(\psi))
\end{equation}
下面陈述并证明关于这些记号的几个必要引理。

\begin{lemma}
\label{chow-lemma-additivity-periodic-length}
设 $R$ 为环，并设有一个 $2$-周期复形的短正合列
$$
0 \to (M_1, N_1, \varphi_1, \psi_1)
\to (M_2, N_2, \varphi_2, \psi_2)
\to (M_3, N_3, \varphi_3, \psi_3)
\to 0
$$
如果其中两个复形的上同调模具有有限长度，那么第三个复形也具有这个性质，并且
$$
e_R(M_2, N_2, \varphi_2, \psi_2) =
e_R(M_1, N_1, \varphi_1, \psi_1) +
e_R(M_3, N_3, \varphi_3, \psi_3).
$$
\end{lemma}

\begin{proof}
简记 $A = (M_1, N_1, \varphi_1, \psi_1)$、
$B = (M_2, N_2, \varphi_2, \psi_2)$ 和 $C = (M_3, N_3, \varphi_3, \psi_3)$。
我们有上同调长正合列
$$
\ldots
\to H^1(C)
\to H^0(A)
\to H^0(B)
\to H^0(C)
\to H^1(A)
\to H^1(B)
\to H^1(C)
\to \ldots
$$
由此得到有限正合列
$$
0 \to I
\to H^0(A)
\to H^0(B)
\to H^0(C)
\to H^1(A)
\to H^1(B)
\to K \to 0
$$
以及滤过 $0 \to K \to H^1(C) \to I \to 0$。由长度函数的可加性
（《交换代数》引理 \ref{algebra-lemma-length-additive}）即得结论。
\end{proof}

\begin{lemma}
\label{chow-lemma-finite-periodic-length}
设 $R$ 为环。如果 $(M, N, \varphi, \psi)$ 是一个 $2$-周期复形，且 $M$、$N$ 具有有限长度，那么
$e_R(M, N, \varphi, \psi) = \text{length}_R(M) - \text{length}_R(N)$。
特别地，如果 $(M, \varphi, \psi)$ 是一个 $(2, 1)$-周期复形，且 $M$ 具有有限长度，那么
$e_R(M, \varphi, \psi) = 0$。
\end{lemma}

\begin{proof}
在由 $M$、$N$ 均具有有限长度的 $2$-周期复形所成的范畴中，量
``$\text{length}_R(M) - \text{length}_R(N)$'' 对短正合列具有可加性
（精确陈述留给读者）。考虑短正合列
$$
0 \to (M, \Im(\varphi), \varphi, 0) \to
(M, N, \varphi, \psi) \to (0, N/\Im(\varphi), 0, 0) \to 0
$$
结合开头的观察与引理 \ref{chow-lemma-additivity-periodic-length} 的可加性，
可将问题约化为以下两种情形：(a) $M = 0$；(b) $\varphi$ 是满射。这两种情形留给读者。
\end{proof}

\begin{lemma}
\label{chow-lemma-compare-periodic-lengths}
设 $R$ 为环。设
$f : (M, \varphi, \psi) \to (M', \varphi', \psi')$
是两个 $(2, 1)$-周期复形之间的映射，并且这两个复形的上同调模均具有有限长度。
如果 $\Ker(f)$ 和 $\Coker(f)$ 具有有限长度，那么
$e_R(M, \varphi, \psi) = e_R(M', \varphi', \psi')$。
\end{lemma}

\begin{proof}
应用引理 \ref{chow-lemma-additivity-periodic-length} 的可加性，并注意由
引理 \ref{chow-lemma-finite-periodic-length}，
$(\Ker(f), \varphi, \psi)$ 和 $(\Coker(f), \varphi', \psi')$ 的重数均为零。
\end{proof}




\section{若干重数的计算}
\label{chow-section-calculation}

\noindent
为了在后文证明某些循环相等，我们需要计算若干重数。除上一节关于重数的初等引理外，
主要工具是《交换代数》引理 \ref{algebra-lemma-order-vanishing-determinant}。

\begin{lemma}
\label{chow-lemma-length-multiplication}
设 $R$ 为诺特局部环，$M$ 为有限 $R$-模，并设 $x \in R$。假设
\begin{enumerate}
\item $\dim(\text{Supp}(M)) \leq 1$，并且
\item $\dim(\text{Supp}(M/xM)) \leq 0$。
\end{enumerate}
写成
$\text{Supp}(M) = \{\mathfrak m, \mathfrak q_1, \ldots, \mathfrak q_t\}$。
那么
$$
e_R(M, 0, x) =
\sum\nolimits_{i = 1, \ldots, t}
\text{ord}_{R/\mathfrak q_i}(x)
\text{length}_{R_{\mathfrak q_i}}(M_{\mathfrak q_i}).
$$
\end{lemma}

\begin{proof}
先作一些准备性说明。如果 $M$ 具有有限长度，即 $t = 0$，则引理成立；
因为此时等式两边都为零，见引理 \ref{chow-lemma-finite-periodic-length}。
此外，若满足 (1)、(2) 的模之间有短正合列 $0 \to M \to M' \to M'' \to 0$，
则只要引理对其中两个模成立，由长度的可加性（《交换代数》引理 \ref{algebra-lemma-length-additive}）和重数的可加性（引理 \ref{chow-lemma-additivity-periodic-length}），它对第三个模也成立。

\medskip\noindent
以 $M_i$ 表示 $M$ 在 $M_{\mathfrak q_i}$ 中的像，于是
$\text{Supp}(M_i) = \{\mathfrak m, \mathfrak q_i\}$。
映射 $M \to \bigoplus M_i$ 的核与余核的支撑均为 $\{\mathfrak m\}$，因而具有有限长度。
由上述准备性说明，只需对每个 $M_i$ 证明引理。因此可以假设
$\text{Supp}(M) = \{\mathfrak m, \mathfrak q\}$。
此时，由《交换代数》引理 \ref{algebra-lemma-Noetherian-power-ideal-kills-module}，有有限滤过
$M \supset \mathfrak qM \supset \mathfrak q^2M \supset \ldots \supset
\mathfrak q^nM = 0$。再次利用可加性可知，只需证明 $M$ 被 $\mathfrak q$ 零化的情形。
这时可把 $M$ 看成 $R/\mathfrak q$-模；换言之，可以假设 $R$ 是维数为 $1$、分式域为 $K$ 的诺特局部整环。
再商去挠子模，即商去 $M \to M \otimes_R K = V$ 的核（挠子模具有有限长度，
故可由准备性说明处理），便可假设 $M \subset V$ 是一个格
（《交换代数》定义 \ref{algebra-definition-lattice}）。于是 $x : M \to M$ 是单射，并且
$\text{length}_R(M/xM) = d(M, xM)$
（《交换代数》定义 \ref{algebra-definition-distance}）。由于
$\text{length}_K(V) = \dim_K(V)$，可知
$\det(x : V \to V) = x^{\dim_K(V)}$，且
$\text{ord}_R(\det(x : V \to V)) = \dim_K(V) \text{ord}_R(x)$。
因而在此情形下，所需等式由《交换代数》引理 \ref{algebra-lemma-order-vanishing-determinant} 得出。
\end{proof}

\begin{lemma}
\label{chow-lemma-additivity-divisors-restricted}
设 $R$ 为诺特局部环，$x \in R$。如果 $M$ 是有限 Cohen--Macaulay $R$-模，且
$\dim(\text{Supp}(M)) = 1$、$\dim(\text{Supp}(M/xM)) = 0$，那么
$$
\text{length}_R(M/xM)
=
\sum\nolimits_i \text{length}_R(R/(x, \mathfrak q_i))
\text{length}_{R_{\mathfrak q_i}}(M_{\mathfrak q_i}).
$$
这里 $\mathfrak q_1, \ldots, \mathfrak q_t$ 是 $M$ 的支撑的极小素理想。如果 $I \subset R$ 是一个理想，
$x$ 在 $R/I$ 上是非零因子，并且 $\dim(R/I) = 1$，那么
$$
\text{length}_R(R/(x, I))
=
\sum\nolimits_i \text{length}_R(R/(x, \mathfrak q_i))
\text{length}_{R_{\mathfrak q_i}}((R/I)_{\mathfrak q_i})
$$
这里 $\mathfrak q_1, \ldots, \mathfrak q_n$ 是包含 $I$ 的极小素理想。
\end{lemma}

\begin{proof}
这些都是引理 \ref{chow-lemma-length-multiplication} 的特殊情形。
\end{proof}

\noindent
下面是在另一种可以确定重数取值的情形。

\begin{lemma}
\label{chow-lemma-powers-period-length-zero}
设 $R$ 为环，$M$ 为 $R$-模。设 $\varphi : M \to M$ 是自同态，$n > 0$，满足
$\varphi^n = 0$，并且 $\Ker(\varphi)/\Im(\varphi^{n - 1})$ 作为 $R$-模具有有限长度。
那么对 $i = 0, \ldots, n$ 均有
$$
e_R(M, \varphi^i, \varphi^{n - i}) = 0
$$
\end{lemma}

\begin{proof}
按约定 $\varphi^0 = \text{id}_M$，所以 $i = 0, n$ 的情形是显然的。
把 $M$ 看成 $R[t]$-模，其中乘以 $t$ 的作用由 $\varphi$ 给出。记
$K_i = \Ker(t^i : M \to M)$，并令
$$
a_i = \text{length}_R(K_i/t^{n - i}M),\quad
b_i = \text{length}_R(K_i/tK_{i + 1}),\quad
c_i = \text{length}_R(K_1/t^iK_{i + 1})
$$
边界值为 $a_0 = a_n = b_0 = c_0 = 0$。当 $i < n$ 时，$c_i$ 是整数，
因为 $K_1/t^iK_{i + 1}$ 是按假设具有有限长度的 $K_1/t^{n - 1}M$ 的商。
下面将反复使用 $K_i \cap t^jM = t^jK_{i + j}$。对 $0 < i < n - 1$，有正合列
$$
0 \to
K_1/t^{n - i - 1}K_{n - i} \to
K_{i + 1}/t^{n - i - 1}M \xrightarrow{t} K_i/t^{n - i}M
\to K_i/tK_{i + 1} \to 0
$$
对 $i$ 作归纳可知，当 $i < n$ 时 $a_i$、$b_i$ 都是整数，并且
$$
c_{n - i - 1} - a_{i + 1} + a_i - b_i = 0
$$
对 $0 < i < n - 1$，还有短正合列
$$
0 \to
K_i/tK_{i + 1} \to
K_{i + 1}/tK_{i + 2} \xrightarrow{t^i}
K_1/t^{i + 1}K_{i + 2} \to
K_1/t^iK_{i + 1} \to 0
$$
因而
$$
b_i - b_{i + 1} + c_{i + 1} - c_i = 0
$$
由于 $b_0 = c_0$，可得对 $i < n$ 均有 $b_i = c_i$。于是
$$
a_2 = a_1 + b_{n - 2} - b_1,\quad
a_3 = a_2 + b_{n - 3} - b_2,\quad \ldots
$$
直接可见这蕴含所需的 $a_i = a_{n - i}$。
\end{proof}

\begin{lemma}
\label{chow-lemma-multiply-period-length}
设 $(R, \mathfrak m)$ 为诺特局部环。设 $(M, \varphi, \psi)$ 是 $R$ 上的一个
$(2, 1)$-周期复形，其中 $M$ 有限，上同调群作为 $R$-模具有有限长度。
设 $x \in R$ 满足 $\dim(\text{Supp}(M/xM)) \leq 0$。那么
$$
e_R(M, x\varphi, \psi) = e_R(M, \varphi, \psi) - e_R(\Im(\varphi), 0, x)
$$
并且
$$
e_R(M, \varphi, x\psi) = e_R(M, \varphi, \psi) + e_R(\Im(\psi), 0, x)
$$
\end{lemma}

\begin{proof}
只证明第一个公式；第二个公式的证明完全相同。
设 $M' = M[x^\infty]$ 为 $M$ 的 $x$-幂挠子模，并考虑短正合列
$0 \to M' \to M \to M'' \to 0$。那么 $M''$ 无 $x$-幂挠
（《代数进阶》引理 \ref{more-algebra-lemma-divide-by-torsion}）。
由于 $\varphi$、$\psi$ 都把 $M'$ 映入 $M'$，可得 $(2, 1)$-周期复形的短正合列
$$
0 \to (M', \varphi', \psi') \to (M, \varphi, \psi) \to
(M'', \varphi'', \psi'') \to 0
$$
此外还有短正合列
$0 \to M' \cap \Im(\varphi) \to \Im(\varphi) \to \Im(\varphi'') \to 0$。
由引理 \ref{chow-lemma-compare-periodic-lengths}，有
$e_R(M', \varphi, \psi) = e_R(M', x\varphi, \psi) =
e_R(M' \cap \Im(\varphi), 0, x) = 0$。
再由可加性（引理 \ref{chow-lemma-additivity-periodic-length}）可知，
只需对 $(M'', \varphi'', \psi'')$ 证明引理。这就约化到下一段所讨论的情形。

\medskip\noindent
假设 $x : M \to M$ 是单射。此时 $\Ker(x\varphi) = \Ker(\varphi)$。
另一方面，有短正合列
$$
0 \to \Im(\varphi)/x\Im(\varphi) \to
\Ker(\psi)/\Im(x\varphi) \to \Ker(\psi)/\Im(\varphi) \to 0
$$
这与 (\ref{chow-equation-multiplicity-coker-ker}) 一起证明了该公式。
\end{proof}







\section{驯顺符号的准备工作}
\label{chow-section-preparation-tame-symbol}

\noindent
本节给出一些引理，以帮助我们在下一节定义驯顺符号。

\begin{lemma}
\label{chow-lemma-glue-at-max}
设 $A$ 为诺特环，$\mathfrak m_1, \ldots, \mathfrak m_r$ 为 $A$ 中两两不同的极大理想。
对 $i = 1, \ldots, r$，设 $\varphi_i : A_{\mathfrak m_i} \to B_i$ 是环同态，
其核与余核均被 $\mathfrak m_i$ 的某个幂零化。那么存在环同态 $\varphi : A \to B$，使得
\begin{enumerate}
\item $\varphi$ 在 $\mathfrak m_i$ 处的局部化同构于 $\varphi_i$；并且
\item $\Ker(\varphi)$ 与 $\Coker(\varphi)$ 均被
$\mathfrak m_1 \cap \ldots \cap \mathfrak m_r$ 的某个幂零化。
\end{enumerate}
此外，如果每个 $\varphi_i$ 都是有限的、单射的或满射的，那么 $\varphi$ 也分别具有相应性质。
\end{lemma}

\begin{proof}
令 $I = \mathfrak m_1 \cap \ldots \cap \mathfrak m_r$，并令
$A_i = A_{\mathfrak m_i}$、$A' = \prod A_i$。那么
$IA' = \prod \mathfrak m_i A_i$，且 $A \to A'$ 是满足 $A/I \cong A'/IA'$ 的平坦环同态。
因此，由《代数进阶》引理 \ref{more-algebra-lemma-application-formal-glueing}，存在 $A$-模同态
$\varphi : A \to B$，使 $\varphi_i$ 同构于 $\varphi$ 在 $\mathfrak m_i$ 处的局部化。
再利用《代数进阶》注 \ref{more-algebra-remark-formal-glueing-algebras}
中的讨论，可在 $B$ 上赋予一个 $A$-代数结构，使其与 $B_i$ 上给定的 $A$-代数结构相符。
由 $\Ker(\varphi)_{\mathfrak m_i} \cong \Ker(\varphi_i)$ 和
$\Coker(\varphi)_{\mathfrak m_i} \cong \Coker(\varphi_i)$，
很容易得到引理的最后一个断言。
\end{proof}

\noindent
下面的引理与《交换代数》引理 \ref{algebra-lemma-nonregular-dimension-one} 非常相似。

\begin{lemma}
\label{chow-lemma-Noetherian-domain-dim-1-two-elements}
设 $(R, \mathfrak m)$ 是维数为 $1$ 的诺特局部环，并设 $a, b \in R$ 为非零因子。
那么存在有限环扩张 $R \subset R'$，其中 $R'/R$ 被 $\mathfrak m$ 的某个幂零化；
并且存在非零因子 $t, a', b' \in R'$，使得
$a = ta'$、$b = tb'$ 且 $R' = a'R' + b'R'$。
\end{lemma}

\begin{proof}
如果 $a$ 或 $b$ 是单位，则取 $R = R'$ 即得引理。因此可以假设 $a, b \in \mathfrak m$。
令 $I = (a, b)$。基本思路是沿 $I$ 爆破 $R$；我们不采用代数论证，而采用几何论证。
令 $X = \text{Proj}(\bigoplus_{d \geq 0} I^d)$。由《除子》引理 \ref{divisors-lemma-blowing-up-gives-effective-Cartier-divisor}，态射
$X \to \Spec(R)$ 在去心谱 $U = \Spec(R) \setminus \{\mathfrak m\}$ 上是同构。
因此可以而且将把 $U$ 看作 $X$ 的开子概形。由《除子》引理 \ref{divisors-lemma-blowing-up-projective}，态射 $X \to \Spec(R)$ 是射影的。
此外，例如由《除子》引理 \ref{divisors-lemma-blow-up-and-irreducible-components}，$X$ 的每个一般点都位于 $U$ 中。
于是由《代数簇》引理 \ref{varieties-lemma-finite-in-codim-1}，
$X \to \Spec(R)$ 是有限态射。因此 $X = \Spec(R')$ 是仿射的，且 $R \to R'$ 有限。
由于 $U = D(a)$，有 $R_a \cong R'_a$。故 $a$ 的某个幂零化有限 $R$-模 $R'/R$。
又因 $\mathfrak m = \sqrt{(a)}$，可知 $R'/R$ 被 $\mathfrak m$ 的某个幂零化。
由《除子》引理 \ref{divisors-lemma-blowing-up-gives-effective-Cartier-divisor}，
$IR'$ 是局部主理想。由于 $R'$ 是半局部环，$IR'$ 是主理想；见《交换代数》引理 \ref{algebra-lemma-locally-free-semi-local-free}。设 $IR' = (t)$，则
$a = a't$、$b = b't$，其余结论均显然。
\end{proof}

\begin{lemma}
\label{chow-lemma-not-infinitely-divisible}
设 $(R, \mathfrak m)$ 是维数为 $1$ 的诺特局部环。设 $a, b \in R$ 是非零因子，且
$a \in \mathfrak m$。存在整数 $n = n(R, a, b)$，使得对于任意有限环扩张
$R \subset R'$，如果对某个 $c \in R'$ 有 $b = a^m c$，那么 $m \leq n$。
\end{lemma}

\begin{proof}
取极小素理想 $\mathfrak q \subset R$。注意 $\dim(R/\mathfrak q) = 1$，特别地，
$R/\mathfrak q$ 不是域。可以选取一个与 $R/\mathfrak q$ 有相同分式域并支配它的离散赋值环 $A$；
见《交换代数》引理 \ref{algebra-lemma-dominate-by-dimension-1}。
由于 $R$ 中的非零因子不属于 $\mathfrak q$，$a$、$b$ 在 $A$ 中的像均非零。
设 $v$ 为 $A$ 上的离散赋值。因 $a \in \mathfrak m$，有 $v(a) > 0$。
我们断言 $n = v(b)/v(a)$ 可用。

\medskip\noindent
给定 $R \subset R'$，令 $A' = A \otimes_R R'$。由于
$\Spec(R') \to \Spec(R)$ 是满射
（《交换代数》引理 \ref{algebra-lemma-integral-overring-surjective}），
$\Spec(A') \to \Spec(A)$ 也为满射
（《交换代数》引理 \ref{algebra-lemma-surjective-spec-radical-ideal}）。
取位于 $(0) \subset A$ 上方的素理想 $\mathfrak q' \subset A'$。
那么 $A \subset A'' = A'/\mathfrak q'$ 是有限环扩张（并且仍诱导谱上的满射）。
取位于 $A$ 的极大理想上方的极大理想 $\mathfrak m'' \subset A''$，再取一个离散赋值环
$A'''$，它支配 $A''_{\mathfrak m''}$（见上引引理）。于是 $A \to A'''$ 是离散赋值环的扩张，并且等式 $b = a^m c$ 在 $A'''$ 中成立。因此 $v'''(b) \geq mv'''(a)$。由于 $v''' = ev$，其中 $e$ 是 $A'''/A$ 的分歧指数，便得到所需的 $m \leq n$。
\end{proof}

\begin{lemma}
\label{chow-lemma-prepare-tame-symbol}
设 $(A, \mathfrak m)$ 是维数为 $1$ 的诺特局部环。设 $r \geq 2$，并设
$a_1, \ldots, a_r \in A$ 是不全为单位的非零因子。那么存在
\begin{enumerate}
\item 有限环扩张 $A \subset B$，其中 $B/A$ 被 $\mathfrak m$ 的某个幂零化；
\item 对每个极大理想 $\mathfrak m_j \subset B$，存在非零因子
$\pi_j \in B_j = B_{\mathfrak m_j}$；并且
\item 在 $B_j$ 中存在分解 $a_i = u_{i, j} \pi_j^{e_{i, j}}$，
其中 $u_{i, j} \in B_j$ 为单位，$e_{i, j} \geq 0$。
\end{enumerate}
\end{lemma}

\begin{proof}
由于至少有一个 $a_i$ 不是单位，$\mathfrak m$ 不是 $A$ 的伴随素理想。
此外，对陈述中任意的 $A \subset B$，$\mathfrak m$ 不是 $B$ 的伴随素理想，
$\mathfrak m_j$ 也不是 $B_j$ 的伴随素理想。记住这一点有助于核查下面的论证。

\medskip\noindent
首先断言，只需对 $r = 2$ 证明引理。我们对 $r$ 作归纳；建议读者跳过这段证明。
设已经给定 $A \subset B$、$B_j = B_{\mathfrak m_j}$ 中的 $\pi_j$，以及在 $B_j$ 中的分解
$a_i = u_{i, j} \pi_j^{e_{i, j}}$（$i = 1, \ldots, r - 1$），
其中 $u_{i, j} \in B_j$ 为单位且 $e_{i, j} \geq 0$。
把 $r = 2$ 的情形用于 $B_j$ 中的 $\pi_j$ 与 $a_r$，便可找到扩张 $B_j \subset C_j$，
并且对 $C_j$ 的每个极大理想 $\mathfrak m_{j, k}$，找到非零因子
$\pi_{j, k} \in C_{j, k} = (C_j)_{\mathfrak m_{j, k}}$ 和引理所述的分解
$$
\pi_j = v_{j, k} \pi_{j, k}^{f_{j, k}}
\quad\text{且}\quad
a_r = w_{j, k} \pi_{j, k}^{g_{j, k}}
$$
由引理 \ref{chow-lemma-glue-at-max}，存在唯一的有限扩张 $B \subset C$，其中
$C/B$ 被 $\mathfrak m$ 的某个幂零化，并且对所有 $j$ 都有
$C_j \cong C_{\mathfrak m_j}$。$C$ 的极大理想与各局部化中的极大理想
$\mathfrak m_{j, k}$ 构成 $1$-对-$1$ 对应；在这些局部化中，对 $i \leq r - 1$ 有
$$
a_i = u_{i, j} \pi_j^{e_{i, j}} =
u_{i, j} v_{j, k}^{e_{i, j}} \pi_{j, k}^{e_{i, j}f_{j, k}}
$$
而 $a_r$ 也有正确的分解，故归纳步骤得证。

\medskip\noindent
下面证明 $r = 2$ 的情形。对
$$
\ell = \min(\text{length}_A(A/a_1A),\ \text{length}_A(A/a_2A)).
$$
作归纳。如果 $\ell = 0$，则 $a_1$ 或 $a_2$ 是单位，取 $A = B$ 即得引理。
因此可以而且将假设 $\ell > 0$。

\medskip\noindent
设有有限环扩张 $A \subset A'$，其中 $A'/A$ 被 $\mathfrak m$ 的某个幂零化，
且 $\mathfrak m$ 不是 $A'$ 的伴随素理想。设
$\mathfrak m_1, \ldots, \mathfrak m_r \subset A'$ 为极大理想，并令
$A'_i = A'_{\mathfrak m_i}$。如果能在每个 $A'_i$ 中解决 $a_1, a_2$ 的问题，
就可应用引理 \ref{chow-lemma-glue-at-max} 得到 $A$ 中 $a_1, a_2$ 的解。
选取 $x \in \{a_1, a_2\}$，使 $\ell = \text{length}_A(A/xA)$。
由引理 \ref{chow-lemma-compare-periodic-lengths} 
与 (\ref{chow-equation-multiplicity-coker-ker})，有
$\text{length}_A(A/xA) = \text{length}_A(A'/xA')$。另一方面，由《交换代数》引理 \ref{algebra-lemma-pushdown-module}，有
$$
\text{length}_A(A'/xA') =
\sum [\kappa(\mathfrak m_i) : \kappa(\mathfrak m)]
\text{length}_{A'_i}(A'_i/xA'_i)
$$
由于 $x \in \mathfrak m$，右边每一项都为正。因此，如果 $r > 1$，或者 $r = 1$ 且
$[\kappa(\mathfrak m_1) : \kappa(\mathfrak m)] > 1$，归纳假设可用于每个 $A'_i$ 中的
$a_1, a_2$。由此可以假设，上述每个 $A'$ 都是与 $A$ 有相同剩余域的局部环。

\medskip\noindent
应用上一段的讨论，可以把 $A$ 替换为引理 \ref{chow-lemma-Noetherian-domain-dim-1-two-elements} 对 $a_1, a_2 \in A$ 构造的环。
由于 $A$ 是局部环，必要时交换 $a_1$、$a_2$ 后，可设 $a_2 \in (a_1)$。
写成 $a_2 = a_1^m c$，其中 $m > 0$ 取最大值。事实上，由引理 \ref{chow-lemma-not-infinitely-divisible}，即使把 $A$ 换成上一段所述的任意有限扩张
$A \subset A'$，也可假设 $m$ 仍为最大值。如果 $c$ 是单位，证明结束；否则，把 $A$ 替换为
引理 \ref{chow-lemma-Noetherian-domain-dim-1-two-elements} 对 $a_1, c \in A$ 构造的环。
此时或者 (1) $c = a_1 c'$，或者 (2) $a_1 = c a'_1$。第一种情形不可能发生，
因为它给出 $a_2 = a_1^{m + 1} c'$，与 $m$ 的最大性矛盾。第二种情形给出
$a_1 = c a'_1$ 和 $a_2 = c^{m + 1} (a'_1)^m$。
因此只需对 $A$ 以及 $c, a'_1$ 证明引理。如果 $a'_1$ 是单位，证明结束；否则，
由于 $cA$ 严格大于 $a_1A$，有 $\text{length}_A(A/cA) < \ell$，
于是由归纳假设得到结论。
\end{proof}






\section{驯顺符号}
\label{chow-section-tame-symbol}

\noindent
考虑维数为 $1$ 的诺特局部环 $(A, \mathfrak m)$。以 $Q(A)$ 表示 $A$ 的全分式环；见
《交换代数》例 \ref{algebra-example-localize-at-prime}。
{\it 驯顺符号}将是一个映射
$$
\partial_A(-, -) : Q(A)^* \times Q(A)^* \longrightarrow \kappa(\mathfrak m)^*
$$
并满足以下性质：
\begin{enumerate}
\item $\partial_A(f, gh) = \partial_A(f, g) \partial_A(f, h)$
\label{chow-item-bilinear}
对 $f, g, h \in Q(A)^*$ 成立；
\item $\partial_A(f, g) \partial_A(g, f) = 1$
\label{chow-item-skew}
对 $f, g \in Q(A)^*$ 成立；
\item $\partial_A(f, 1 - f) = 1$
\label{chow-item-1-x}
对满足 $1 - f \in Q(A)^*$ 的 $f \in Q(A)^*$ 成立；
\item $\partial_A(aa', b) = \partial_A(a, b)\partial_A(a', b)$
\label{chow-item-bilinear-better}
和 $\partial_A(a, bb') = \partial_A(a, b)\partial_A(a, b')$
对非零因子 $a, a', b, b' \in A$ 成立；
\item $\partial_A(b, b) = (-1)^m$
\label{chow-item-skew-better}
对非零因子 $b \in A$ 成立，其中 $m = \text{length}_A(A/bA)$；
\item $\partial_A(u, b) = u^m \bmod \mathfrak m$
\label{chow-item-normalization}
对单位 $u \in A$ 和非零因子 $b \in A$ 成立，其中
$m = \text{length}_A(A/bA)$；并且
\item
\label{chow-item-1-x-better}
$\partial_A(a, b - a)\partial_A(b, b) = \partial_A(b, b - a)\partial_A(a, b)$
对使 $a, b, b - a$ 都是非零因子的 $a, b \in A$ 成立。
\end{enumerate}
由于处理 $A$ 的元素比较方便，我们常把 $\partial_A$ 看作定义在 $A$ 的非零因子对上的映射，
满足 (\ref{chow-item-bilinear-better})、(\ref{chow-item-skew-better})、
(\ref{chow-item-normalization}) 和 (\ref{chow-item-1-x-better})。
读者不难验证，令
$$
\partial_A(\frac{a}{b}, \frac{c}{d}) =
\partial_A(a, c) \partial_A(a, d)^{-1} \partial_A(b, c)^{-1} \partial_A(b, d)
$$
便得到良定义映射 $Q(A)^* \times Q(A)^* \to \kappa(\mathfrak m)^*$，
它满足 (\ref{chow-item-bilinear})、(\ref{chow-item-skew})、(\ref{chow-item-1-x}) 以及上述其他性质。

\medskip\noindent
我们并不声称满足这些性质的映射是唯一的，而是给出构造这种映射的方法。
具体而言，给定非零因子 $a_1, a_2 \in A$，选择如引理 \ref{chow-lemma-prepare-tame-symbol} 中的环扩张 $A \subset B$ 与局部分解。定义
\begin{equation}
\label{chow-equation-tame-symbol}
\partial_A(a_1, a_2) = \prod\nolimits_j
\text{Norm}_{\kappa(\mathfrak m_j)/\kappa(\mathfrak m)}
((-1)^{e_{1, j}e_{2, j}}u_{1, j}^{e_{2, j}}u_{2, j}^{-e_{1, j}}
\bmod \mathfrak m_j)^{m_j}
\end{equation}
其中 $m_j = \text{length}_{B_j}(B_j/\pi_j B_j)$，乘积遍历 $B$ 的极大理想
$\mathfrak m_1, \ldots, \mathfrak m_r$。

\begin{lemma}
\label{chow-lemma-well-defined-tame-symbol}
公式 (\ref{chow-equation-tame-symbol}) 确定 $\kappa(\mathfrak m)^*$ 中一个良定义的元素。
换言之，右边不依赖于局部分解的选择，也不依赖于 $B$ 的选择。
\end{lemma}

\begin{proof}
先证与分解选择无关。设 $B$ 是一个诺特 $1$-维局部环，$a_1, a_2 \in B$，
并设 $\pi, \theta$ 为非零因子，使得
$$
a_1 = u_1 \pi^{e_1} = v_1 \theta^{f_1},\quad
a_2 = u_2 \pi^{e_2} = v_2 \theta^{f_2}
$$
其中 $e_i, f_i \geq 0$ 是整数，$u_i, v_i$ 是 $B$ 中的单位。由此可得
$$
a_1^{e_2} = u_1^{e_2}u_2^{-e_1}a_2^{e_1},\quad
a_1^{f_2} = v_1^{f_2}v_2^{-f_1}a_2^{f_1}
$$
另一方面，令 $m = \text{length}_B(B/\pi B)$、$k = \text{length}_B(B/\theta B)$，
便有 $e_2 m = \text{length}_B(B/a_2 B) = f_2 k$。
用上式展开 $a_1^{e_2m} = a_1^{f_2 k}$，得到
$$
(u_1^{e_2}u_2^{-e_1})^m =  (v_1^{f_2}v_2^{-f_1})^k
$$
这已证明除符号外的所需等式。为了验证符号一致，只须证明 $me_1e_2$ 为偶数当且仅当
$kf_1f_2$ 为偶数。由于 $me_2 = kf_2$ 且 $me_1 = kf_1$（论证同上），结论成立。

\medskip\noindent
再证与 $B$ 的选择无关。设有引理 \ref{chow-lemma-prepare-tame-symbol} 中的两个扩张
$A \subset B$ 和 $A \subset B'$。那么
$$
C = (B \otimes_A B')/(\mathfrak m\text{-幂挠子})
$$
给出第三个这样的扩张。因此可以假设 $A \subset B \subset C$，并已在 $B$ 的各局部环中给定分解；
需要证明，在 $C$ 的各局部环中使用同一组分解会得到 $\kappa(\mathfrak m)$ 中的同一元素。
由范数的传递性（《域》引理 \ref{fields-lemma-trace-and-norm-tower}），
问题归结为：若 $B$ 是维数为 $1$ 的诺特局部环，且 $\pi \in B$ 是非零因子，则
$$
\lambda^m = \prod \text{Norm}_{\kappa_k/\kappa}(\lambda)^{m_k}
$$
这里使用如下记号：
(1) $\kappa$ 是 $B$ 的剩余域；
(2) $\lambda$ 是 $\kappa$ 的一个元素；
(3) $\mathfrak m_k \subset C$ 是 $C$ 的极大理想；
(4) $\kappa_k = \kappa(\mathfrak m_k)$ 是 $C_k = C_{\mathfrak m_k}$ 的剩余域；
(5) $m = \text{length}_B(B/\pi B)$；以及
(6) $m_k = \text{length}_{C_k}(C_k/\pi C_k)$。
上述等式成立，因为对 $\lambda \in \kappa$ 有
$\text{Norm}_{\kappa_k/\kappa}(\lambda) = \lambda^{[\kappa_k : \kappa]}$，
且 $m = \sum m_k[\kappa_k:\kappa]$。首先，由引理 \ref{chow-lemma-compare-periodic-lengths} 与 (\ref{chow-equation-multiplicity-coker-ker})，有
$m = \text{length}_B(B/xB) = \text{length}_B(C/\pi C)$。
最后，由《交换代数》引理 \ref{algebra-lemma-pushdown-module}，有
$\text{length}_B(C/\pi C) = \sum m_k[\kappa_k:\kappa]$。
\end{proof}

\begin{lemma}
\label{chow-lemma-tame-symbol}
驯顺符号 (\ref{chow-equation-tame-symbol}) 满足
(\ref{chow-item-bilinear-better})、(\ref{chow-item-skew-better})、
(\ref{chow-item-normalization}) 和 (\ref{chow-item-1-x-better})，因而给出映射
$\partial_A : Q(A)^* \times Q(A)^* \to \kappa(\mathfrak m)^*$，
它满足 (\ref{chow-item-bilinear})、(\ref{chow-item-skew}) 和 (\ref{chow-item-1-x})。
\end{lemma}

\begin{proof}
先证明 (\ref{chow-item-bilinear-better})。设非零因子 $a_1, a_2, a_3 \in A$。
为 $a_1, a_2, a_3$ 选择引理 \ref{chow-lemma-prepare-tame-symbol} 中的
$A \subset B$。等式
$$
\partial_A(a_1a_2, a_3) = \partial_A(a_1, a_3) \partial_A(a_2, a_3)
$$
由 $B_j$ 中的等式
$$
(-1)^{(e_{1, j} + e_{2, j})e_{3, j}}
(u_{1, j}u_{2, j})^{e_{3, j}}u_{3, j}^{-e_{1, j} - e_{2, j}} =
(-1)^{e_{1, j}e_{3, j}}
u_{1, j}^{e_{3, j}}u_{3, j}^{-e_{1, j}}
(-1)^{e_{2, j}e_{3, j}}
u_{2, j}^{e_{3, j}}u_{3, j}^{-e_{2, j}}
$$
推出。性质 (\ref{chow-item-skew-better}) 和 (\ref{chow-item-normalization}) 同样是立即可得的。

\medskip\noindent
下面证明 (\ref{chow-item-1-x-better})。设 $a_1, a_2, a_1 - a_2 \in A$ 为非零因子，
并令 $a_3 = a_1 - a_2$。为 $a_1, a_2, a_3$ 选择引理 \ref{chow-lemma-prepare-tame-symbol} 中的 $A \subset B$。只需证明
$$
(-1)^{e_{1, j}e_{2, j} + e_{1, j}e_{3, j} + e_{2, j}e_{3, j} + e_{2, j}}
u_{1, j}^{e_{2, j} - e_{3, j}}
u_{2, j}^{e_{3, j} - e_{1, j}}
u_{3, j}^{e_{1, j} - e_{2, j}} \bmod \mathfrak m_j = 1
$$
如果 $e_{1, j} = e_{2, j} = e_{3, j}$，这很显然。设 $e_{1, j} > e_{2, j}$。
由 $a_3 = a_1 - a_2$ 可知 $e_{3, j} = e_{2, j}$，且 $u_{3, j}$ 与
$-u_{2, j}$ 有相同的剩余类。因此公式成立，符号也恰好相符；
这项核查正是 (\ref{chow-equation-tame-symbol}) 中如此选择符号的原因。
其他情形的处理完全相同。
\end{proof}

\begin{lemma}
\label{chow-lemma-norm-down-tame-symbol}
设 $(A, \mathfrak m)$ 为维数为 $1$ 的诺特局部环。设 $A \subset B$ 是有限环扩张，
$B/A$ 被 $\mathfrak m$ 的某个幂零化，并且 $\mathfrak m$ 不是 $B$ 的伴随素理想。
对非零因子 $a, b \in A$，有
$$
\partial_A(a, b) = \prod
\text{Norm}_{\kappa(\mathfrak m_j)/\kappa(\mathfrak m)}(\partial_{B_j}(a, b))
$$
其中乘积遍历 $B$ 的极大理想 $\mathfrak m_j$，且 $B_j = B_{\mathfrak m_j}$。
\end{lemma}

\begin{proof}
为 $a, b$ 选择引理 \ref{chow-lemma-prepare-tame-symbol} 中的 $B_j \subset C_j$。
由引理 \ref{chow-lemma-glue-at-max}，可选取有限环扩张 $B \subset C$，
使得对所有 $j$ 都有 $C_j \cong C_{\mathfrak m_j}$。
设 $\mathfrak m_{j, k} \subset C$ 是位于 $\mathfrak m_j$ 上方的 $C$ 的极大理想。设
$$
a = u_{j, k}\pi_{j, k}^{f_{j, k}},\quad
b = v_{j, k}\pi_{j, k}^{g_{j, k}}
$$
为选择 $C_j \cong C_{\mathfrak m_j}$ 后得到的局部分解。由定义，
$$
\partial_A(a, b) = 
\prod\nolimits_{j, k}
\text{Norm}_{\kappa(\mathfrak m_{j, k})/\kappa(\mathfrak m)}
((-1)^{f_{j, k}g_{j, k}}u_{j, k}^{g_{j, k}}v_{j, k}^{-f_{j, k}}
\bmod \mathfrak m_{j, k})^{m_{j, k}}
$$
并且
$$
\partial_{B_j}(a, b) = 
\prod\nolimits_k
\text{Norm}_{\kappa(\mathfrak m_{j, k})/\kappa(\mathfrak m_j)}
((-1)^{f_{j, k}g_{j, k}}u_{j, k}^{g_{j, k}}v_{j, k}^{-f_{j, k}}
\bmod \mathfrak m_{j, k})^{m_{j, k}}
$$
对域塔 $\kappa(\mathfrak m_{j, k})/\kappa(\mathfrak m_j)/\kappa(\mathfrak m)$
应用范数的传递性即得结论；见《域》引理 \ref{fields-lemma-trace-and-norm-tower}。
\end{proof}

\begin{lemma}
\label{chow-lemma-tame-symbol-formally-smooth}
设 $(A, \mathfrak m, \kappa) \to (A', \mathfrak m', \kappa')$
是诺特局部环的局部同态。假设 $A \to A'$ 平坦，且 $\dim(A) = \dim(A') = 1$。
令 $m = \text{length}_{A'}(A'/\mathfrak mA')$。
对非零因子 $a_1, a_2 \in A$，$\partial_A(a_1, a_2)^m$ 经
$\kappa \to \kappa'$ 映为 $\partial_{A'}(a_1, a_2)$。
\end{lemma}

\begin{proof}
如果 $a_1, a_2$ 都是单位，那么 $\partial_A(a_1, a_2) = 1$ 且
$\partial_{A'}(a_1, a_2) = 1$，结论成立。否则，可以选择引理 \ref{chow-lemma-prepare-tame-symbol} 中的环扩张 $A \subset B$ 与局部分解。
以 $\mathfrak m_1, \ldots, \mathfrak m_m$ 表示 $B$ 的极大理想。
设 $\mathfrak m_1, \ldots, \mathfrak m_m$ 是 $B$ 的极大理想，其剩余域为
$\kappa_1, \ldots, \kappa_m$。对每个 $j \in \{1, \ldots, m\}$，
以 $\pi_j \in B_j = B_{\mathfrak m_j}$ 表示一个非零因子，使得有引理中的分解
$a_i = u_{i, j}\pi_j^{e_{i, j}}$。由定义，
$$
\partial_A(a_1, a_2) = \prod\nolimits_j
\text{Norm}_{\kappa_j/\kappa}
((-1)^{e_{1, j}e_{2, j}}u_{1, j}^{e_{2, j}}u_{2, j}^{-e_{1, j}}
\bmod \mathfrak m_j)^{m_j}
$$
其中 $m_j = \text{length}_{B_j}(B_j/\pi_j B_j)$。

\medskip\noindent
令 $B' = A' \otimes_A B$。由于 $A'$ 在 $A$ 上平坦，$A' \subset B'$ 是环扩张，
且 $B'/A'$ 被 $\mathfrak m'$ 的某个幂零化。设
$$
\mathfrak m'_{j, l},\quad l = 1, \ldots, n_j
$$
为 $B'$ 中位于 $\mathfrak m_j$ 上方的极大理想。以 $\kappa'_{j, l}$ 表示
$\mathfrak m'_{j, l}$ 的剩余域，以 $B'_{j, l}$ 表示 $B'$ 在 $\mathfrak m'_{j, l}$ 处的局部化。
在 $B'_{j, l}$ 中分解 $a_1$、$a_2$ 时，采用 $B_j$ 中给定分解
$a_i = u_{i, j} \pi_j^{e_{i, j}}$ 的像。由定义，
$$
\partial_{A'}(a_1, a_2) = \prod\nolimits_{j, l}
\text{Norm}_{\kappa'_{j, l}/\kappa'}
((-1)^{e_{1, j}e_{2, j}}u_{1, j}^{e_{2, j}}u_{2, j}^{-e_{1, j}}
\bmod \mathfrak m'_{j, l})^{m'_{j, l}}
$$
其中 $m'_{j, l} = \text{length}_{B'_{j, l}}(B'_{j, l}/\pi_j B'_{j, l})$。

\medskip\noindent
比较上述公式可知，只需证明：对每个 $j$ 和任意单位 $u \in B_j$，都有
\begin{equation}
\label{chow-equation-to-prove}
\left(\text{Norm}_{\kappa_j/\kappa}(u \bmod \mathfrak m_j)^{m_j}\right)^m
=
\prod\nolimits_l
\text{Norm}_{\kappa'_{j, l}/\kappa'}(u \bmod \mathfrak m'_{j, l})^{m'_{j, l}}
\end{equation}
在 $\kappa'$ 中成立。我们用《代数进阶》第 \ref{more-algebra-section-determinants-finite-length} 节中有限长度模的自同态行列式构造证明这一点。
令 $M = B_j/\pi_j B_j$。由《代数进阶》引理 \ref{more-algebra-lemma-multiplication}，有
$$
\text{Norm}_{\kappa_j/\kappa}(u \bmod \mathfrak m_j)^{m_j} =
\det\nolimits_\kappa(u : M \to M)
$$
因此，由《代数进阶》引理 \ref{more-algebra-lemma-flat-base-change-det}，
(\ref{chow-equation-to-prove}) 的左边等于
$\det_{\kappa'}(u : M \otimes_A A' \to M \otimes_A A')$。
有 $A'$-模同构
$$
M \otimes_A A' = (B_j/\pi_j B_j) \otimes_A A' =
\bigoplus\nolimits_l B'_{j, l}/\pi_j B'_{j, l}
$$
令 $M'_l = B'_{j, l}/\pi_j B'_{j, l}$。再次由《代数进阶》引理 \ref{more-algebra-lemma-multiplication}，可得
$\text{Norm}_{\kappa'_{j, l}/\kappa'}(u \bmod \mathfrak m'_{j, l})^{m'_{j, l}}
= \det_{\kappa'}(u_j : M'_l \to M'_l)$。
所以 (\ref{chow-equation-to-prove}) 由行列式构造的乘法性成立；见《代数进阶》引理 \ref{more-algebra-lemma-ses}。
\end{proof}






\section{一个关键引理}
\label{chow-section-key-lemma}

\noindent
本节应用上述结果证明引理 \ref{chow-lemma-milnor-gersten-low-degree}。
该引理是如下陈述的低次数情形：Milnor K-理论存在一个类似于 Quillen K-理论中
Gersten--Quillen 复形的复形。见注 \ref{chow-remark-gersten-complex-milnor}。

\begin{lemma}
\label{chow-lemma-perpare-key}
设 $(A, \mathfrak m)$ 是一个 $2$-维诺特局部环。设 $t \in \mathfrak m$ 为非零因子，并写成
$V(t) = \{\mathfrak m, \mathfrak q_1, \ldots, \mathfrak q_r\}$。
设 $A_{\mathfrak q_i} \subset B_i$ 是有限环扩张，其中
$B_i/A_{\mathfrak q_i}$ 被 $t$ 的某个幂零化。那么存在保持剩余域相同的有限局部环扩张
$A \subset B$，使得 $B_i \cong B_{\mathfrak q_i}$，并且 $B/A$ 被 $t$ 的某个幂零化。
\end{lemma}

\begin{proof}
选取 $n > 0$，使 $B_i \subset t^{-n}A_{\mathfrak q_i}$。
令 $M \subset t^{-n}A$ 和 $M' \subset t^{-2n}A$ 分别为\ 由如下元素组成的 $A$-子模：
它们在 $t^{-n}A_{\mathfrak q_i}$ 和 $t^{-2n}A_{\mathfrak q_i}$ 中分别\ 映入 $B_i$。
由于 $A$ 诺特，$M \subset M'$ 是有限 $A$-模；又因局部化正合，
$M_{\mathfrak q_i} = M'_{\mathfrak q_i} = B_i$。因此对某个 $c > 0$，
$M'/M$ 被 $\mathfrak m^c$ 零化。注意，在乘法
$t^{-n}A \times t^{-n}A \to t^{-2n}A$ 下，有 $M \cdot M \subset M'$。
故 $B = A + \mathfrak m^{c + 1}M$ 是有限 $A$-代数，并具有所需的局部化。
我们略去如下验证：$B$ 是局部环，其极大理想为
$\mathfrak m + \mathfrak m^{c + 1}M$。
\end{proof}

\begin{lemma}
\label{chow-lemma-key-nonzerodivisors}
设 $(A, \mathfrak m)$ 是一个 $2$-维诺特局部环，$a, b \in A$ 是非零因子。那么
$$
\sum
\text{ord}_{A/\mathfrak q}(\partial_{A_{\mathfrak q}}(a, b))
=
0
$$
其中求和遍历 $A$ 中所有高度为 $1$ 的素理想 $\mathfrak q$。
\end{lemma}

\begin{proof}
若 $\mathfrak q$ 是 $A$ 中高度为 $1$ 的素理想，并且 $a, b$ 在 $A_\mathfrak q$ 中的像是单位，
则 $\partial_{A_\mathfrak q}(a, b) = 1$。因此该和为有限和。事实上，若
$V(ab) = \{\mathfrak m, \mathfrak q_1, \ldots, \mathfrak q_r\}$，
则求和遍历 $i = 1, \ldots, r$。对每个 $i$，为 $a, b$ 选取引理 \ref{chow-lemma-prepare-tame-symbol} 中的扩张 $A_{\mathfrak q_i} \subset B_i$。
把引理 \ref{chow-lemma-perpare-key} 用于 $t = ab$ 与给定的素理想列表，
可假设存在有限局部扩张 $A \subset B$，其中 $B/A$ 被 $ab$ 的某个幂零化，
并且对每个 $i$ 都有 $B_{\mathfrak q_i} \cong B_i$。
注意，若 $\mathfrak q_{i, j}$ 是 $B$ 中位于 $\mathfrak q_i$ 上方的素理想，则由
引理 \ref{chow-lemma-norm-down-tame-symbol} 和《交换代数》引理 \ref{algebra-lemma-finite-extension-dim-1}，有
$$
\text{ord}_{A/\mathfrak q_i}(\partial_{A_{\mathfrak q_i}}(a, b))
=
\sum\nolimits_j
\text{ord}_{B/\mathfrak q_{i, j}}(\partial_{B_{\mathfrak q_{i, j}}}(a, b))
$$
因此可以用 $B$ 代替 $A$，并约化到下一段讨论的情形。

\medskip\noindent
假设对每个 $i$，存在非零因子 $\pi_i \in A_{\mathfrak q_i}$ 和单位
$u_i, v_i \in A_{\mathfrak q_i}$，使得对某些整数 $e_i, f_i \geq 0$，在
$A_{\mathfrak q_i}$ 中有
$$
a = u_i \pi_i^{e_i},\quad b = v_i \pi_i^{f_i}
$$
令 $m_i = \text{length}_{A_{\mathfrak q_i}}(A_{\mathfrak q_i}/\pi_i)$。
根据定义，
$\partial_{A_{\mathfrak q_i}}(a, b) =
((-1)^{e_if_i}u_i^{f_i}v_i^{-e_i})^{m_i}$。
由于 $a, b$ 是非零因子，$(2, 1)$-周期复形 $(A/(ab), a, b)$ 的上同调为零。
以 $M_i$ 表示 $A/(ab)$ 在 $A_{\mathfrak q_i}/(ab)$ 中的像。于是有映射
$$
A/(ab) \longrightarrow \bigoplus M_i
$$
其核与余核均支撑在 $\{\mathfrak m\}$ 上，因而具有有限长度。因此，由引理 \ref{chow-lemma-compare-periodic-lengths}，有
$$
\sum e_A(M_i, a, b) = 0
$$
故只须证明
$e_A(M_i, a, b) =
- \text{ord}_{A/\mathfrak q_i}(\partial_{A_{\mathfrak q_i}}(a, b))$。

\medskip\noindent
先证明 $\pi_i, u_i, v_i$ 是某些元素 $\pi_i, u_i, v_i \in A$ 的像这一情形
（使用相同符号不会引起混淆）。此时
\begin{align*}
e_A(M_i, a, b) & =
e_A(M_i, u_i\pi_i^{e_i}, v_i\pi_i^{f_i}) \\
& =
e_A(M_i, \pi_i^{e_i}, \pi_i^{f_i}) -
e_A(\pi_i^{e_i}M_i, 0, u_i) +
e_A(\pi_i^{f_i}M_i, 0, v_i) \\
& =
0 -
f_im_i\text{ord}_{A/\mathfrak q_i}(u_i) +
e_im_i\text{ord}_{A/\mathfrak q_i}(v_i) \\
& =
-m_i\text{ord}_{A/\mathfrak q_i}(u_i^{f_i}v_i^{-e_i}) =
-\text{ord}_{A/\mathfrak q_i}(\partial_{A_{\mathfrak q_i}}(a, b))
\end{align*}
第二个等式来自引理 \ref{chow-lemma-multiply-period-length}。注意
$M_i \subset (M_i)_{\mathfrak q_i} = A_{\mathfrak q_i}/(\pi_i^{e_i + f_i})$，并且
$(\pi_i^{e_i}M_i)_{\mathfrak q_i} \cong A_{\mathfrak q_i}/\pi_i^{f_i}$、
$(\pi_i^{f_i}M_i)_{\mathfrak q_i} \cong A_{\mathfrak q_i}/\pi_i^{e_i}$。
第三个等式中的 $0$ 来自引理 \ref{chow-lemma-powers-period-length-zero}，
另外两项来自引理 \ref{chow-lemma-length-multiplication}。
最后两个等式由阶函数的乘法性和驯顺符号的定义得到。

\medskip\noindent
在一般情形下，先选取 $c \in A$、$c \not \in \mathfrak q_i$，使 $c\pi_i \in A$。
以 $c\pi_i$ 代替 $\pi_i$，以 $c^{-e_i}u_i$ 代替 $u_i$，并以
$c^{-f_i}v_i$ 代替 $v_i$ 后，可以而且将假设 $\pi_i$ 属于 $A$。
接着选取 $c \in A$、$c \not \in \mathfrak q_i$，使 $cu_i, cv_i \in A$。
由引理 \ref{chow-lemma-length-multiplication}，有
$$
e_A(M_i, ca, cb) = e_A(M_i, a, b) - e_A(aM_i, 0, c) + e_A(bM_i, 0, c)
$$
另一方面，由于 $c$ 是 $A_{\mathfrak q_i}$ 中的单位，在 $\kappa(\mathfrak q_i)^*$ 中有
$$
\partial_{A_{\mathfrak q_i}}(ca, cb) =
c^{m_i(f_i - e_i)}\partial_{A_{\mathfrak q_i}}(a, b)
$$
上一段的论证表明
$e_A(M_i, ca, cb) = -
\text{ord}_{A/\mathfrak q_i}(\partial_{A_{\mathfrak q_i}}(ca, cb))$。
因此只须证明
$$
e_A(aM_i, 0, c) = \text{ord}_{A/\mathfrak q_i}(c^{m_if_i})
\quad\text{且}\quad
e_A(bM_i, 0, c) = \text{ord}_{A/\mathfrak q_i}(c^{m_ie_i})
$$
这由引理 \ref{chow-lemma-length-multiplication} 以及上文对在 $\mathfrak q_i$ 处局部化所得情形的描述推出。
\end{proof}

\begin{lemma}[关键引理]
\label{chow-lemma-milnor-gersten-low-degree}
\begin{reference}
当 $A$ 是优秀环时，这是 \cite[Proposition 1]{Kato-Milnor-K}。
\end{reference}
设 $A$ 是一个 $2$-维诺特局部整环，其分式域为 $K$。设 $f, g \in K^*$。
设 $\mathfrak q_1, \ldots, \mathfrak q_t$ 是 $A$ 中满足如下条件的高度为 $1$ 的素理想
$\mathfrak q$：$f$ 或 $g$ 至少有一个不属于 $A^*_{\mathfrak q}$。那么
$$
\sum\nolimits_{i = 1, \ldots, t}
\text{ord}_{A/\mathfrak q_i}(\partial_{A_{\mathfrak q_i}}(f, g))
=
0
$$
也可以写成
$$
\sum\nolimits_{\text{高度}(\mathfrak q) = 1}
\text{ord}_{A/\mathfrak q}(\partial_{A_{\mathfrak q}}(f, g))
=
0
$$
因为对 $A$ 中任意高度为 $1$ 的素理想 $\mathfrak q$，若
$f, g \in A^*_{\mathfrak q}$，则 $\partial_{A_{\mathfrak q}}(f, g) = 1$。
\end{lemma}

\begin{proof}
驯顺符号 $\partial_{A_{\mathfrak q}}(f, g)$ 是双线性的，而阶函数
$\text{ord}_{A/\mathfrak q}$ 是可加的，所以只需在 $f$、$g$ 都属于 $A$ 时证明该公式。
这一情形已在引理 \ref{chow-lemma-key-nonzerodivisors} 中证明。
\end{proof}

\begin{remark}[Milnor K-理论]
\label{chow-remark-gersten-complex-milnor}
对域 $k$，以 $K^M_*(k)$ 表示 $k^*$ 上的张量代数对如下双边理想的商：
该理想由所有 $x \in k \setminus \{0, 1\}$ 对应的元素
$x \otimes 1 - x$ 生成。因此 $K^M_0(k) = \mathbf{Z}$、$K_1^M(k) = k^*$，并且
$$
K^M_2(k) = k^* \otimes_\mathbf{Z} k^* / \langle x \otimes 1 - x \rangle
$$
若 $A$ 是分式域为 $F = \text{Frac}(A)$、剩余域为 $\kappa$ 的离散赋值环，
则有如第 \ref{chow-section-tame-symbol} 节定义的驯顺符号
$$
\partial_A : K_{i + 1}^M(F) \to K_i^M(\kappa)
$$
见 \cite{Kato-Milnor-K}。更一般地，利用正规化以及 $K_i^M$ 上的范数映射，可以把这一映射推广到
$A$ 为维数为 $1$ 的优秀局部整环的情形；见 \cite{Kato-Milnor-K}。
可以推测，第 \ref{chow-section-tame-symbol} 节的方法能够把驯顺符号 $\partial_A$ 的构造推广到任意维数为
$1$ 的诺特局部整环 $A$。接着，设 $X$ 是带有维数函数 $\delta$ 的诺特概形。
利用这些驯顺符号，可以得到下列箭头：
$$
\bigoplus\nolimits_{\delta(x) = j + 1} K^M_{i + 1}(\kappa(x))
\longrightarrow
\bigoplus\nolimits_{\delta(x) = j} K^M_i(\kappa(x))
\longrightarrow
\bigoplus\nolimits_{\delta(x) = j - 1} K^M_{i - 1}(\kappa(x))
$$
但尚不清楚该复合是否为零，亦即是否由此得到一个阿贝尔群复形。
对优秀概形 $X$，这在 \cite{Kato-Milnor-K} 中得到证明。
当 $i = 1$ 而 $j$ 任意时，结论来自引理 \ref{chow-lemma-milnor-gersten-low-degree}。
\end{remark}





















\section{基本设置}
\label{chow-section-setup}

\noindent
在本章中，我们始终在一个配有维数函数 $\delta$ 的局部诺特、泛链状基概形 $S$ 上工作。
尽管本章的讨论（至少一部分）很可能可以推广，但这似乎是一个良好的第一步近似；
它恰好就是 \cite{Thorup} 所讨论的一般程度。通常不假设所研究的概形是分离的或拟紧的。
许多有趣的代数栈并不分离和/或并不拟紧，因此这里也是研究如何为这些对象发展合理理论的一个好例子。
为了便于引用这些假设，我们为其编号。

\begin{situation}
\label{chow-situation-setup}
这里 $S$ 是一个局部诺特且泛链状的概形。此外，假设 $S$ 配有维数函数
$\delta : S \longrightarrow \mathbf{Z}$。
\end{situation}

\noindent
泛链状概形的概念见《概形态射》定义 \ref{morphisms-definition-universally-catenary}；维数函数的概念见《拓扑学》定义 \ref{topology-definition-dimension-function}。回忆，任意局部诺特链状概形在局部都有维数函数；
见《概形的性质》引理 \ref{properties-lemma-catenary-dimension-function}。
此外，泛链状概形相当常见；见《概形态射》引理 \ref{morphisms-lemma-ubiquity-uc}。

\medskip\noindent
设 $(S, \delta)$ 如情形 \ref{chow-situation-setup}。任意在 $S$ 上局部有限型的概形 $X$
都是局部诺特且链状的。事实上，$X$ 有一个与 $(f : X \to S, \delta)$ 相伴的典范维数函数
$$
\delta = \delta_{X/S} : X \longrightarrow \mathbf{Z}
$$
其规则为
$\delta_{X/S}(x) = \delta(f(x)) + \text{trdeg}_{\kappa(f(x))}\kappa(x)$。
见《概形态射》引理 \ref{morphisms-lemma-dimension-function-propagates}。
此外，若 $h : X \to Y$ 是 $S$ 上局部有限型概形之间的态射，且
$x \in X$、$y = h(x)$，则显然
$\delta_{X/S}(x) = \delta_{Y/S}(y) + \text{trdeg}_{\kappa(y)}\kappa(x)$。
下文将自由使用这个函数及其性质。

\medskip\noindent
下面给出上述设置的基本例子。事实上，最值得关注的是基为一个域的谱，
或者基为 Dedekind 环的谱（例如\ $\mathbf{Z}$ 或离散赋值环）的情形。

\begin{example}
\label{chow-example-field}
这里 $S = \Spec(k)$，其中 $k$ 是域。令 $\delta(pt) = 0$，其中 $pt$ 表示 $S$ 的唯一点。
$(S, \delta)$ 是情形 \ref{chow-situation-setup} 中设置的一个例子；
这来自《概形态射》引理 \ref{morphisms-lemma-ubiquity-uc}。
\end{example}

\begin{example}
\label{chow-example-domain-dimension-1}
这里 $S = \Spec(A)$，其中 $A$ 是维数为 $1$ 的诺特整环。例如可以取 $A = \mathbf{Z}$。
若 $\mathfrak p$ 是极大理想，令 $\delta(\mathfrak p) = 0$；若
$\mathfrak p = (0)$ 对应一般点，令 $\delta(\mathfrak p) = 1$。
这是情形 \ref{chow-situation-setup} 的一个例子；
结论来自《概形态射》引理 \ref{morphisms-lemma-ubiquity-uc}。
\end{example}

\begin{example}
\label{chow-example-CM-irreducible}
这里 $S$ 是 Cohen--Macaulay 概形。由《概形态射》引理 \ref{morphisms-lemma-ubiquity-uc}，$S$ 泛链状。令
$\delta(s) = -\dim(\mathcal{O}_{S, s})$。若 $s' \leadsto s$ 是 $S$ 中点的非平凡特殊化，
则 $\mathcal{O}_{S, s'}$ 是 $\mathcal{O}_{S, s}$ 在非极大素理想
$\mathfrak p \subset \mathcal{O}_{S, s}$ 处的局部化；见《概形》引理 \ref{schemes-lemma-specialize-points}。因此由《交换代数》引理 \ref{algebra-lemma-CM-dim-formula}，有
$\dim(\mathcal{O}_{S, s}) =
\dim(\mathcal{O}_{S, s'}) + \dim(\mathcal{O}_{S, s}/\mathfrak p) >
\dim(\mathcal{O}_{S, s'})$。故 $\delta(s') > \delta(s)$。
若 $s' \leadsto s$ 是紧邻特殊化，则 $\mathfrak p$ 与 $\mathfrak m_s$ 之间没有严格介于二者的素理想，
从而 $\delta(s') = \delta(s) + 1$。所以 $\delta$ 是维数函数。
换言之，$(S, \delta)$ 是情形 \ref{chow-situation-setup} 的一个例子。
\end{example}

\noindent
如果 $S$ 是 Jacobson 概形，而 $\delta$ 把闭点映为零，那么 $\delta$ 就是把每个点映到其闭包维数的函数。

\begin{lemma}
\label{chow-lemma-delta-is-dimension}
设 $(S, \delta)$ 如情形 \ref{chow-situation-setup}。进一步假设 $S$ 是 Jacobson 概形，
并且对 $S$ 的每个闭点 $s$ 都有 $\delta(s) = 0$。设 $X$ 在 $S$ 上局部有限型。
设 $Z \subset X$ 是整闭子概形，$\xi \in Z$ 是它的一般点。下列整数相同：
\begin{enumerate}
\item $\delta_{X/S}(\xi)$；
\item $\dim(Z)$；以及
\item $\dim(\mathcal{O}_{Z, z})$，其中 $z$ 是 $Z$ 的闭点。
\end{enumerate}
\end{lemma}

\begin{proof}
设 $X \to S$、$\xi \in Z \subset X$ 如引理所述。由于 $X$ 在 $S$ 上局部有限型，
$X$ 是 Jacobson 概形；见《概形态射》引理 \ref{morphisms-lemma-Jacobson-universally-Jacobson}。因此，$X$ 的闭点在 $Z$ 的每个闭子集中稠密，
并映到 $S$ 的闭点。所以，给定 $Z$ 中不可约闭子集的任意链，可以用 $Z$ 的某个闭点将其终结。
从而，当 $z \in Z$ 遍历 $Z$ 的闭点时，
$\dim(Z) = \sup_z(\dim(\mathcal{O}_{Z, z})$；
见《概形的性质》引理 \ref{properties-lemma-codimension-local-ring}。
由维数函数的性质，$\dim(\mathcal{O}_{Z, z}) = \delta(\xi) - \delta(z)$。
对每个闭点 $z \in Z$，域扩张 $\kappa(z)/\kappa(f(z))$ 有限；见《概形态射》引理 \ref{morphisms-lemma-jacobson-finite-type-points}。
因此，对闭点 $z \in Z$，有 $\delta_{X/S}(z) = \delta(f(z)) = 0$。
由此可知三个整数全都相等。
\end{proof}

\noindent
在上述引理的情形下，$\delta$ 在一个闭不可约子集的一般点处的值，就是该不可约闭子集的维数。
但一般来说不能期待这一等式成立。例如，若 $S = \Spec(\mathbf{C}[[t]])$ 且
$X = \Spec(\mathbf{C}((t)))$，则对 $X$ 的唯一点有 $\delta(x) = 1$，但
$\dim(X) = 0$。尽管如此，我们仍希望把 $\delta_{X/S}$ 视作给出不可约闭子概形的维数。
因此引入以下术语。

\begin{definition}
\label{chow-definition-delta-dimension}
设 $(S, \delta)$ 如情形 \ref{chow-situation-setup}。对任意在 $S$ 上局部有限型的概形 $X$
和任意不可约闭子集 $Z \subset X$，定义
$$
\dim_\delta(Z) = \delta(\xi)
$$
其中 $\xi \in Z$ 是 $Z$ 的一般点。称其为{\it $Z$ 的 $\delta$-维数}。
如果 $Z$ 是 $X$ 的闭子概形，则定义 $\dim_\delta(Z)$ 为其不可约分支的
$\delta$-维数的上确界。
\end{definition}







\section{循环}
\label{chow-section-cycles}

\noindent
由于不假设所研究的概形拟紧，定义循环时必须稍加谨慎。例如，有理函数可能有无穷多个极点，
因此必须允许无穷和。不过，如果 $X$ 拟紧，那么循环仍像通常一样是有限和。

\begin{definition}
\label{chow-definition-cycles}
设 $(S, \delta)$ 如情形 \ref{chow-situation-setup}，设 $X$ 在 $S$ 上局部有限型，并设
$k \in \mathbf{Z}$。
\begin{enumerate}
\item $X$ 上的一个{\it 循环}是形式和
$$
\alpha = \sum n_Z [Z]
$$
其中求和遍历整闭子概形 $Z \subset X$，每个 $n_Z \in \mathbf{Z}$，并且集合
$\{Z; n_Z \not = 0\}$ 局部有限
（《拓扑学》定义 \ref{topology-definition-locally-finite}）。
\item $X$ 上的一个 {\it $k$-循环}是满足下述条件的循环
$$
\alpha = \sum n_Z [Z]
$$
即 $n_Z \not = 0 \Rightarrow \dim_\delta(Z) = k$。
\item $X$ 上所有 $k$-循环构成的阿贝尔群记为 $Z_k(X)$。
\end{enumerate}
\end{definition}

\noindent
换言之，$X$ 上的 $k$-循环是 $\delta$-维数为 $k$ 的整闭子概形的局部有限形式
$\mathbf{Z}$-线性组合。$k$-循环 $\alpha = \sum n_Z[Z]$ 与
$\beta = \sum m_Z[Z]$ 的加法定义为
$$
\alpha + \beta = \sum (n_Z + m_Z)[Z],
$$
也就是把系数相加。

\begin{remark}
\label{chow-remark-cycles-pointwise}
设 $(S, \delta)$ 如情形 \ref{chow-situation-setup}，设 $X$ 在 $S$ 上局部有限型，并设
$k \in \mathbf{Z}$。可以写成
$$
Z_k(X) = \bigoplus\nolimits_{\delta(x) = k}' K_0^M(\kappa(x))
\quad\subset\quad
\bigoplus\nolimits_{\delta(x) = k} K_0^M(\kappa(x))
$$
其中采用以下记号与约定：
\begin{enumerate}
\item $K_0^M(\kappa(x)) = \mathbf{Z}$ 是点 $x \in X$ 的剩余域 $\kappa(x)$ 的
Milnor K-理论的次数 $0$ 部分（见注 \ref{chow-remark-gersten-complex-milnor}）；
\item 右边的直和遍历满足 $\delta(x) = k$ 的所有点 $x \in X$；
\item 记号 $\bigoplus'_x$ 表示只取由局部有限元素构成的子群；也就是取满足如下条件的元素
$\sum_x n_x$：对每个拟紧开集 $U \subset X$，满足 $n_x \not = 0$ 的点 $x \in U$ 只有有限多个。
\end{enumerate}
\end{remark}

\begin{definition}
\label{chow-definition-support-cycle}
设 $(S, \delta)$ 如情形 \ref{chow-situation-setup}，并设 $X$ 在 $S$ 上局部有限型。
$X$ 上循环 $\alpha = \sum n_Z [Z]$ 的{\it 支撑}为
$$
\text{Supp}(\alpha) = \bigcup\nolimits_{n_Z \not = 0} Z \subset X
$$
\end{definition}

\noindent
由于集合 $\{Z; n_Z \not = 0\}$ 局部有限，$\text{Supp}(\alpha)$ 是 $X$ 的闭子集。
若 $\alpha$ 是 $k$-循环，则 $\text{Supp}(\alpha)$ 的每个不可约分支 $Z$ 的
$\delta$-维数都为 $k$。

\begin{definition}
\label{chow-definition-effective-cycle}
设 $(S, \delta)$ 如情形 \ref{chow-situation-setup}，并设 $X$ 在 $S$ 上局部有限型。
$X$ 上的循环 $\alpha$ 如果可以写成 $\alpha =\sum n_Z [Z]$，且对所有 $Z$ 都有 $n_Z \geq 0$，
就称该循环是{\it 有效的}。
\end{definition}

\noindent
所有有效循环构成一个幺半群，因为两个有效循环之和仍有效；但它不是群
（除非 $X = \emptyset$）。




\section{与闭子概形相伴的循环}
\label{chow-section-cycle-of-closed-subscheme}

\begin{lemma}
\label{chow-lemma-multiplicity-finite}
设 $(S, \delta)$ 如情形 \ref{chow-situation-setup}，设 $X$ 在 $S$ 上局部有限型，
并设 $Z \subset X$ 是闭子概形。
\begin{enumerate}
\item 设 $Z' \subset Z$ 是一个不可约分支，$\xi \in Z'$ 是它的一般点。那么
$$
\text{length}_{\mathcal{O}_{X, \xi}} \mathcal{O}_{Z, \xi} < \infty
$$
\item 如果 $\dim_\delta(Z) \leq k$，并且 $\xi \in Z$ 满足
$\delta(\xi) = k$，那么 $\xi$ 是 $Z$ 的某个不可约分支的一般点。
\end{enumerate}
\end{lemma}

\begin{proof}
设 $Z' \subset Z$、$\xi \in Z'$ 如 (1) 所述。那么
$\dim(\mathcal{O}_{Z, \xi}) = 0$；例如见《概形的性质》引理 \ref{properties-lemma-codimension-local-ring}。因此 $\mathcal{O}_{Z, \xi}$ 是维数为零的诺特局部环，
从而作为自身的模具有有限长度（见《交换代数》命题 \ref{algebra-proposition-dimension-zero-ring}）。所以它作为 $\mathcal{O}_{X, \xi}$-模也具有有限长度；
见《交换代数》引理 \ref{algebra-lemma-length-independent}。

\medskip\noindent
假设 $\xi \in Z$ 且 $\delta(\xi) = k$。考虑闭包
$Z' = \overline{\{\xi\}}$。根据定义，它是满足 $\dim_\delta(Z') = k$ 的不可约闭子概形。
由于 $\dim_\delta(Z) = k$，它必为 $Z$ 的一个不可约分支。因此 (2) 成立。
\end{proof}

\begin{definition}
\label{chow-definition-cycle-associated-to-closed-subscheme}
设 $(S, \delta)$ 如情形 \ref{chow-situation-setup}，设 $X$ 在 $S$ 上局部有限型，
并设 $Z \subset X$ 是闭子概形。
\begin{enumerate}
\item 对任意一般点为 $\xi$ 的不可约分支 $Z' \subset Z$，称整数
$m_{Z', Z} = \text{length}_{\mathcal{O}_{X, \xi}} \mathcal{O}_{Z, \xi}$
（引理 \ref{chow-lemma-multiplicity-finite}）为 {\it $Z'$ 在 $Z$ 中的重数}。
\item 假设 $\dim_\delta(Z) \leq k$。定义{\it 与 $Z$ 相伴的 $k$-循环}为
$$
[Z]_k
=
\sum m_{Z', Z}[Z']
$$
其中求和遍历 $Z$ 中 $\delta$-维数为 $k$ 的不可约分支。
（由《除子》引理 \ref{divisors-lemma-components-locally-finite}，这是一个 $k$-循环。）
\end{enumerate}
\end{definition}

\noindent
必须注意，只有在 $Z$ 的 $\delta$-维数不超过 $k$ 时才定义 $[Z]_k$。
换言之，按约定，只要写下 $[Z]_k$，就隐含着 $\dim_\delta(Z) \leq k$。




\section{与凝聚层相伴的循环}
\label{chow-section-cycle-of-coherent-sheaf}



\begin{lemma}
\label{chow-lemma-length-finite}
设 $(S, \delta)$ 如情形 \ref{chow-situation-setup}，设 $X$ 在 $S$ 上局部有限型，
并设 $\mathcal{F}$ 是凝聚 $\mathcal{O}_X$-模。
\begin{enumerate}
\item $\mathcal{F}$ 的支撑的不可约分支所成集合局部有限。
\item 设 $Z' \subset \text{Supp}(\mathcal{F})$ 是一个不可约分支，
$\xi \in Z'$ 是它的一般点。那么
$$
\text{length}_{\mathcal{O}_{X, \xi}} \mathcal{F}_\xi < \infty
$$
\item 如果 $\dim_\delta(\text{Supp}(\mathcal{F})) \leq k$，并且
$\xi \in \text{Supp}(\mathcal{F})$ 满足 $\delta(\xi) = k$，那么 $\xi$ 是
$\text{Supp}(\mathcal{F})$ 的某个不可约分支的一般点。
\end{enumerate}
\end{lemma}

\begin{proof}
由《概形上同调》引理 \ref{coherent-lemma-coherent-support-closed}，
$\mathcal{F}$ 的支撑 $Z$ 是 $X$ 的闭子集。可以把 $Z$ 看成 $X$ 的约化闭子概形
（《概形》引理 \ref{schemes-lemma-reduced-closed-subscheme}）。
因此，把《除子》引理 \ref{divisors-lemma-components-locally-finite} 用于 $Z$ 即得 (1)，
把引理 \ref{chow-lemma-multiplicity-finite} 用于 $Z$ 即得 (3)。

\medskip\noindent
设 $\xi \in Z'$ 如 (2) 所述。在此情形下，对 $X$ 中任意特殊化
$\xi' \leadsto \xi$，都有 $\mathcal{F}_{\xi'} = 0$。回忆，
$\mathcal{O}_{X, \xi}$ 的非极大素理想对应于 $X$ 中特殊化到 $\xi$ 的点
（《概形》引理 \ref{schemes-lemma-specialize-points}）。
因此 $\mathcal{F}_\xi$ 是有限 $\mathcal{O}_{X, \xi}$-模，其支撑为
$\{\mathfrak m_\xi\}$。所以由《交换代数》引理 \ref{algebra-lemma-support-point}，
它具有有限长度。
\end{proof}

\begin{definition}
\label{chow-definition-cycle-associated-to-coherent-sheaf}
设 $(S, \delta)$ 如情形 \ref{chow-situation-setup}，设 $X$ 在 $S$ 上局部有限型，
并设 $\mathcal{F}$ 是凝聚 $\mathcal{O}_X$-模。
\begin{enumerate}
\item 对任意一般点为 $\xi$ 的不可约分支
$Z' \subset \text{Supp}(\mathcal{F})$，称整数
$m_{Z', \mathcal{F}} = \text{length}_{\mathcal{O}_{X, \xi}} \mathcal{F}_\xi$
（引理 \ref{chow-lemma-length-finite}）为 {\it $Z'$ 在 $\mathcal{F}$ 中的重数}。
\item 假设 $\dim_\delta(\text{Supp}(\mathcal{F})) \leq k$。定义
{\it 与 $\mathcal{F}$ 相伴的 $k$-循环}为
$$
[\mathcal{F}]_k
=
\sum m_{Z', \mathcal{F}}[Z']
$$
其中求和遍历 $\text{Supp}(\mathcal{F})$ 中 $\delta$-维数为 $k$ 的不可约分支。
（由引理 \ref{chow-lemma-length-finite}，这是一个 $k$-循环。）
\end{enumerate}
\end{definition}

\noindent
必须注意，只有当 $\mathcal{F}$ 凝聚，且 $\text{Supp}(\mathcal{F})$ 的
$\delta$-维数不超过 $k$ 时，才定义 $[\mathcal{F}]_k$。换言之，按约定，只要写下
$[\mathcal{F}]_k$，就隐含着 $\mathcal{F}$ 在 $X$ 上凝聚，且
$\dim_\delta(\text{Supp}(\mathcal{F})) \leq k$。

\begin{lemma}
\label{chow-lemma-cycle-closed-coherent}
设 $(S, \delta)$ 如情形 \ref{chow-situation-setup}，设 $X$ 在 $S$ 上局部有限型，
并设 $Z \subset X$ 是闭子概形。如果 $\dim_\delta(Z) \leq k$，那么
$[Z]_k = [{\mathcal O}_Z]_k$。
\end{lemma}

\begin{proof}
这是因为在此情形下，重数 $m_{Z', Z}$ 与 $m_{Z', \mathcal{O}_Z}$ 依定义相等。
\end{proof}

\begin{lemma}
\label{chow-lemma-additivity-sheaf-cycle}
设 $(S, \delta)$ 如情形 \ref{chow-situation-setup}，设 $X$ 在 $S$ 上局部有限型，
并设 $0 \to \mathcal{F} \to \mathcal{G} \to \mathcal{H} \to 0$
是 $X$ 上凝聚层的短正合列。假设 $\mathcal{F}$、$\mathcal{G}$ 与 $\mathcal{H}$ 的支撑的
$\delta$-维数都满足 $\leq k$。那么
$[\mathcal{G}]_k = [\mathcal{F}]_k + [\mathcal{H}]_k$。
\end{lemma}

\begin{proof}
这立即来自长度的可加性；见《交换代数》引理 \ref{algebra-lemma-length-additive}。
\end{proof}









\section{固有推前的准备工作}
\label{chow-section-preparation-pushforward}

\begin{lemma}
\label{chow-lemma-equal-dimension}
设 $(S, \delta)$ 如情形 \ref{chow-situation-setup}，设 $X$、$Y$ 在 $S$ 上局部有限型，
并设 $f : X \to Y$ 是态射。假设 $X$、$Y$ 都是整概形，且
$\dim_\delta(X) = \dim_\delta(Y)$。那么，或者 $f(X)$ 包含在 $Y$ 的某个真闭子概形中，
或者 $f$ 是支配态射，并且函数域扩张 $R(X)/R(Y)$ 有限。
\end{lemma}

\begin{proof}
由于 $X$ 不可约，闭包 $\overline{f(X)} \subset Y$ 不可约
（《拓扑学》引理 \ref{topology-lemma-image-irreducible-space} 和
\ref{topology-lemma-irreducible}）。如果 $\overline{f(X)} \not = Y$，结论已得。
如果 $\overline{f(X)} = Y$，则 $f$ 支配；由《概形态射》引理 \ref{morphisms-lemma-dominant-finite-number-irreducible-components}，
$Y$ 的一般点 $\eta_Y$ 位于 $f$ 的像中。这当然蕴含 $f(\eta_X) = \eta_Y$，
其中 $\eta_X \in X$ 是 $X$ 的一般点。由于 $\delta(\eta_X) = \delta(\eta_Y)$，
$R(Y) = \kappa(\eta_Y) \subset \kappa(\eta_X) = R(X)$ 是超越次数为 $0$ 的扩张。
所以由《概形态射》引理 \ref{morphisms-lemma-finite-degree}，
$R(Y) \subset R(X)$ 是有限扩张；该引理适用是由《概形态射》引理 \ref{morphisms-lemma-permanence-finite-type} 保证的。
\end{proof}

\begin{lemma}
\label{chow-lemma-quasi-compact-locally-finite}
设 $(S, \delta)$ 如情形 \ref{chow-situation-setup}，设 $X$、$Y$ 在 $S$ 上局部有限型，
并设 $f : X \to Y$ 是态射。假设 $f$ 拟紧，且 $\{Z_i\}_{i \in I}$ 是 $X$ 中闭子集的局部有限族。
那么 $\{\overline{f(Z_i)}\}_{i \in I}$ 是 $Y$ 中闭子集的局部有限族。
\end{lemma}

\begin{proof}
设 $V \subset Y$ 是拟紧开子集。由于 $f$ 拟紧，开集 $f^{-1}(V)$ 拟紧。
因此集合 $\{i \in I \mid Z_i \cap f^{-1}(V) \not = \emptyset \}$ 有限；
这里略去一个简单的拓扑论证。这个集合与
$$
\{i \in I \mid f(Z_i) \cap V \not = \emptyset \} =
\{i \in I \mid \overline{f(Z_i)} \cap V \not = \emptyset \}
$$
相同，故引理得证。
\end{proof}









\section{固有推前}
\label{chow-section-proper-pushforward}

\begin{definition}
\label{chow-definition-proper-pushforward}
设 $(S, \delta)$ 如情形 \ref{chow-situation-setup}，设 $X$、$Y$ 在 $S$ 上局部有限型，
并设 $f : X \to Y$ 是态射。假设 $f$ 固有。
\begin{enumerate}
\item 设 $Z \subset X$ 是满足 $\dim_\delta(Z) = k$ 的整闭子概形。定义
$$
f_*[Z] =
\left\{
\begin{matrix}
0 & \text{若} & \dim_\delta(f(Z))< k, \\
\deg(Z/f(Z)) [f(Z)] & \text{若} & \dim_\delta(f(Z)) = k.
\end{matrix}
\right.
$$
这里把 $f(Z) \subset Y$ 看作整闭子概形。若
$\dim_\delta(f(Z)) = \dim_\delta(Z)$，则由引理 \ref{chow-lemma-equal-dimension}，
$Z$ 在 $f(Z)$ 上的次数有限。
\item 设 $\alpha = \sum n_Z [Z]$ 是 $X$ 上的 $k$-循环。定义 $\alpha$ 的{\it 推前}为和
$$
f_* \alpha = \sum n_Z f_*[Z]
$$
其中每个 $f_*[Z]$ 如上定义。由引理 \ref{chow-lemma-quasi-compact-locally-finite}，这个和局部有限。
\end{enumerate}
\end{definition}

\noindent
依定义，循环的固有推前
$$
f_* : Z_k(X) \longrightarrow Z_k(Y)
$$
是阿贝尔群同态。它把 $X \mapsto Z_k(X)$ 变成如下范畴上的协变函子：
对象是在 $S$ 上局部有限型的概形，态射为固有态射。

\begin{lemma}
\label{chow-lemma-compose-pushforward}
设 $(S, \delta)$ 如情形 \ref{chow-situation-setup}，设 $X$、$Y$、$Z$ 在 $S$ 上局部有限型，
并设 $f : X \to Y$、$g : Y \to Z$ 为固有态射。那么作为映射
$Z_k(X) \to Z_k(Z)$，有 $g_* \circ f_* = (g \circ f)_*$。
\end{lemma}

\begin{proof}
设 $W \subset X$ 是维数为 $k$ 的整闭子概形。考虑
$W' = f(W) \subset Y$ 与 $W'' = g(f(W)) \subset Z$。
由于 $f$、$g$ 固有，$W'$（以及\ $W''$）分别是 $Y$（以及\ $Z$）的整闭子概形。
需要证明 $g_*(f_*[W]) = (g \circ f)_*[W]$。若 $\dim_\delta(W'') < k$，
则两边均为零。若 $\dim_\delta(W'') = k$，则所诱导的两个态射
$$
W \longrightarrow
W' \longrightarrow
W''
$$
都满足引理 \ref{chow-lemma-equal-dimension} 的假设。因此
$$
g_*(f_*[W]) = \deg(W/W')\deg(W'/W'')[W''],
\quad
(g \circ f)_*[W] = \deg(W/W'')[W''].
$$
再应用《概形态射》引理 \ref{morphisms-lemma-degree-composition} 即得结论。
\end{proof}

\noindent
闭浸入是固有态射。如果 $i : Z \to X$ 是闭浸入，那么映射
$$
i_* : Z_k(Z) \longrightarrow Z_k(X)
$$
全都是{\it 单射}。

\begin{lemma}
\label{chow-lemma-exact-sequence-closed}
设 $(S, \delta)$ 如情形 \ref{chow-situation-setup}。设 $X$ 在 $S$ 上局部有限型，并设
$X_1, X_2 \subset X$ 是闭子概形，使得在集合意义下 $X = X_1 \cup X_2$。
对每个 $k \in \mathbf{Z}$，阿贝尔群序列
$$
\xymatrix{
Z_k(X_1 \cap X_2) \ar[r] &
Z_k(X_1) \oplus Z_k(X_2) \ar[r] &
Z_k(X) \ar[r] &
0
}
$$
正合。这里 $X_1 \cap X_2$ 是概形论交，映射都是推前映射，其中一个乘以 $-1$。
\end{lemma}

\begin{proof}
先假设 $X$ 拟紧。那么 $Z_k(X)$ 是自由 $\mathbf{Z}$-模；它的一组基为元素 $[Z]$，
其中 $Z \subset X$ 是 $\delta$-维数为 $k$ 的整闭子概形。群
$Z_k(X_1)$、$Z_k(X_2)$、$Z_k(X_1 \cap X_2)$ 分别由这些 $Z$ 中满足
$Z \subset X_1$、$Z \subset X_2$、$Z \subset X_1 \cap X_2$ 的子集自由生成。
这立即证明了此情形下的引理。一般情形相似，略去证明。
\end{proof}

\begin{lemma}
\label{chow-lemma-cycle-push-sheaf}
设 $(S, \delta)$ 如情形 \ref{chow-situation-setup}，并设 $f : X \to Y$ 是在 $S$ 上局部有限型的概形之间的固有态射。
\begin{enumerate}
\item 设 $Z \subset X$ 是满足 $\dim_\delta(Z) \leq k$ 的闭子概形。那么
$$
f_*[Z]_k = [f_*{\mathcal O}_Z]_k.
$$
\item 设 $\mathcal{F}$ 是 $X$ 上的凝聚层，且
$\dim_\delta(\text{Supp}(\mathcal{F})) \leq k$。那么
$$
f_*[\mathcal{F}]_k = [f_*{\mathcal F}]_k.
$$
\end{enumerate}
该陈述有意义，因为由《概形上同调》命题 \ref{coherent-proposition-proper-pushforward-coherent}，$f_*\mathcal{F}$ 和
$f_*\mathcal{O}_Z$ 都是凝聚 $\mathcal{O}_Y$-模。
\end{lemma}

\begin{proof}
(1) 由 (2) 和引理 \ref{chow-lemma-cycle-closed-coherent} 推出。
设 $\mathcal{F}$ 是 $X$ 上的凝聚层，并假设
$\dim_\delta(\text{Supp}(\mathcal{F})) \leq k$。
由《概形上同调》引理 \ref{coherent-lemma-coherent-support-closed}，
存在闭子概形 $i : Z \to X$ 和凝聚 $\mathcal{O}_Z$-模 $\mathcal{G}$，使得
$i_*\mathcal{G} \cong \mathcal{F}$，并且 $\mathcal{F}$ 的支撑为 $Z$。
设 $Z' \subset Y$ 是 $f|_Z : Z \to Y$ 的概形论像。考虑概形的交换图
$$
\xymatrix{
Z \ar[r]_i \ar[d]_{f|_Z} &
X \ar[d]^f \\
Z' \ar[r]^{i'} & Y
}
$$
沿该图的两条路径可得
$f_*\mathcal{F} = f_*i_*\mathcal{G} = i'_*(f|_Z)_*\mathcal{G}$。
假设已知结果对闭浸入以及 $f|_Z$ 成立，则
$$
f_*[\mathcal{F}]_k = f_*i_*[\mathcal{G}]_k
= (i')_*(f|_Z)_*[\mathcal{G}]_k =
(i')_*[(f|_Z)_*\mathcal{G}]_k =
[(i')_*(f|_Z)_*\mathcal{G}]_k = [f_*\mathcal{F}]_k
$$
如所需。闭浸入的情形直接可得，略去证明。注意
$f|_Z : Z \to Z'$ 是支配态射；见《概形态射》引理 \ref{morphisms-lemma-quasi-compact-scheme-theoretic-image}。
因此已经约化到 $\dim_\delta(X) \leq k$ 且 $f : X \to Y$ 固有并支配的情形。

\medskip\noindent
假设 $\dim_\delta(X) \leq k$，且 $f : X \to Y$ 固有并支配。
由于 $f$ 支配，对一般点为 $\eta$ 的每个不可约分支 $Z \subset Y$，
都存在点 $\xi \in X$ 使 $f(\xi) = \eta$。因此
$\delta(\eta) \leq \delta(\xi) \leq k$。所以在表达式
$$
f_*[\mathcal{F}]_k = \sum n_Z[Z],
\quad
\text{且}
\quad
[f_*\mathcal{F}]_k = \sum m_Z[Z].
$$
中，只要 $n_Z \not = 0$ 或 $m_Z \not = 0$，整闭子概形 $Z$ 实际上就是 $Y$ 中
$\delta$-维数为 $k$ 的一个不可约分支。选取这样的整闭子概形 $Z \subset Y$，
并以 $\eta$ 表示其一般点。注意，对任何满足 $f(\xi) = \eta$ 的 $\xi \in X$，
都有 $\delta(\xi) \geq k$，因而 $\xi$ 也是 $X$ 中某个 $\delta$-维数为 $k$ 的不可约分支的一般点
（见引理 \ref{chow-lemma-multiplicity-finite}）。由于 $f$ 拟紧且 $X$ 局部诺特，
这样的点只有有限多个，故 $f^{-1}(\{\eta\})$ 有限。
由《概形态射》引理 \ref{morphisms-lemma-generically-finite}，存在开邻域
$\eta \in V \subset Y$，使 $f^{-1}(V) \to V$ 有限。以 $V$ 代替 $Y$，
以 $f^{-1}(V)$ 代替 $X$，便约化到 $Y$ 仿射且 $f$ 有限的情形。

\medskip\noindent
写成 $Y = \Spec(R)$、$X = \Spec(A)$；这是可能的，因为有限态射是仿射的。
那么 $R$、$A$ 是诺特环，且 $A$ 在 $R$ 上有限。此外，对某个有限 $A$-模 $M$，有
$\mathcal{F} = \widetilde{M}$。注意，$f_*\mathcal{F}$ 对应于被视作 $R$-模的 $M$。
设 $\mathfrak p \subset R$ 是对应于 $\eta \in Y$ 的极小素理想。
$Z$ 在 $[f_*\mathcal{F}]_k$ 中的系数显然为
$\text{length}_{R_{\mathfrak p}}(M_{\mathfrak p})$。
设 $\mathfrak q_i$（$i = 1, \ldots, t$）是 $A$ 中位于 $\mathfrak p$ 上方的素理想。
由于 $A_{\mathfrak p}$ 在维数为零的诺特局部环 $R_{\mathfrak p}$ 上有限，
它是 Artin 环，因而 $A_{\mathfrak p} = \prod A_{\mathfrak q_i}$。
$Z$ 在 $f_*[\mathcal{F}]_k$ 中的系数显然为
$$
\sum\nolimits_{i = 1, \ldots, t}
[\kappa(\mathfrak q_i) : \kappa(\mathfrak p)]
\text{length}_{A_{\mathfrak q_i}}(M_{\mathfrak q_i})
$$
所需等式于是来自《交换代数》引理 \ref{algebra-lemma-pushdown-module}。
\end{proof}




















\section{平坦拉回的准备工作}
\label{chow-section-preparation-flat-pullback}


\noindent
回忆，若局部有限型态射 $f : X \to Y$ 的每个非空纤维都是维数为 $r$ 的等维概形，
就称它的相对维数为 $r$。见《概形态射》定义 \ref{morphisms-definition-relative-dimension-d}。

\begin{lemma}
\label{chow-lemma-flat-inverse-image-dimension}
设 $(S, \delta)$ 如情形 \ref{chow-situation-setup}，设 $X$、$Y$ 在 $S$ 上局部有限型，
并设 $f : X \to Y$ 是态射。假设 $f$ 是相对维数为 $r$ 的平坦态射。
对任何闭子集 $Z \subset Y$，只要 $f^{-1}(Z)$ 非空，就有
$$
\dim_\delta(f^{-1}(Z)) = \dim_\delta(Z) + r.
$$
若 $Z$ 不可约，而 $Z' \subset f^{-1}(Z)$ 是一个不可约分支，
则 $Z'$ 支配 $Z$，并且 $\dim_\delta(Z') = \dim_\delta(Z) + r$。
\end{lemma}

\begin{proof}
只需证明最后一个断言。可以用整闭子概形 $Z$ 代替 $Y$，
并用概形论逆像 $f^{-1}(Z) = Z \times_Y X$ 代替 $X$。
因此可以假设 $Z = Y$ 为整概形，且 $f$ 是相对维数为 $r$ 的平坦态射。
由于 $Y$ 局部诺特，局部有限型态射 $f$ 实际上局部有限呈示。
故《概形态射》引理 \ref{morphisms-lemma-fppf-open} 适用，并表明 $f$ 是开态射。
设 $\xi \in X$ 是 $X$ 的某个不可约分支的一般点。由 $f$ 的开性，
$f(\xi)$ 是 $Z = Y$ 的一般点 $\eta$。因为 $f$ 的相对维数为 $r$，有
$\dim_\xi(X_\eta) = r$。另一方面，由于 $\xi$ 是 $X$ 的一般点，
$\mathcal{O}_{X, \xi} = \mathcal{O}_{X_\eta, \xi}$ 只有一个素理想，因而维数为 $0$。
所以由《概形态射》引理 \ref{morphisms-lemma-dimension-fibre-at-a-point}，
$\kappa(\xi)$ 在 $\kappa(\eta)$ 上的超越次数为 $r$。
换言之，$\delta(\xi) = \delta(\eta) + r$，如所需。
\end{proof}

\noindent
下面的引理将用于证明：闭子概形的局部有限族经平坦拉回后仍局部有限。

\begin{lemma}
\label{chow-lemma-inverse-image-locally-finite}
设 $(S, \delta)$ 如情形 \ref{chow-situation-setup}，设 $X$、$Y$ 在 $S$ 上局部有限型，
并设 $f : X \to Y$ 是态射。假设 $\{Z_i\}_{i \in I}$ 是 $Y$ 中闭子集的局部有限族，
那么 $\{f^{-1}(Z_i)\}_{i \in I}$ 是 $X$ 中闭子集的局部有限族。
\end{lemma}

\begin{proof}
设 $U \subset X$ 是拟紧开子集。由于像 $f(U) \subset Y$ 是拟紧子集，
存在拟紧开集 $V \subset Y$ 使 $f(U) \subset V$。注意
$$
\{i \in I \mid f^{-1}(Z_i) \cap U \not = \emptyset \}
\subset
\{i \in I \mid Z_i \cap V \not = \emptyset \}.
$$
右边依假设有限，故结论成立。
\end{proof}



\section{平坦拉回}
\label{chow-section-flat-pullback}

\noindent
下文以 $f^{-1}(Z)$ 表示如下概形论逆像：对概形态射 $f : X \to Y$ 与闭子概形
$Z \subset Y$，取该闭子概形的{\it 概形论逆像}。回忆，概形论逆像是纤维积
$$
\xymatrix{
f^{-1}(Z) \ar[r] \ar[d] & X \ar[d] \\
Z \ar[r] & Y
}
$$
如果 $\mathcal{I} \subset \mathcal{O}_Y$ 是对应于 $Y$ 中 $Z$ 的拟凝聚理想层，
它也就是由拟凝聚理想层 $f^{-1}(\mathcal{I})\mathcal{O}_X$ 截出的 $X$ 的闭子概形。
对此的讨论见《概形》第 \ref{schemes-section-closed-immersion} 节、
引理 \ref{schemes-lemma-fibre-product-immersion} 和定义 \ref{schemes-definition-inverse-image-closed-subscheme}。

\begin{definition}
\label{chow-definition-flat-pullback}
设 $(S, \delta)$ 如情形 \ref{chow-situation-setup}，设 $X$、$Y$ 在 $S$ 上局部有限型，
并设 $f : X \to Y$ 是态射。假设 $f$ 是相对维数为 $r$ 的平坦态射。
\begin{enumerate}
\item 设 $Z \subset Y$ 是 $\delta$-维数为 $k$ 的整闭子概形。定义 $f^*[Z]$ 为
$X$ 上与概形论逆像相伴的 $(k+r)$-循环
$$
f^*[Z] = [f^{-1}(Z)]_{k+r}.
$$
这是有意义的，因为由引理 \ref{chow-lemma-flat-inverse-image-dimension}，
$\dim_\delta(f^{-1}(Z)) = k + r$。
\item 设 $\alpha = \sum n_i [Z_i]$ 是 $Y$ 上的 $k$-循环。定义
{\it $\alpha$ 沿 $f$ 的平坦拉回}为和
$$
f^* \alpha = \sum n_i f^*[Z_i]
$$
其中每个 $f^*[Z_i]$ 如上定义。由引理 \ref{chow-lemma-inverse-image-locally-finite}，
该和局部有限。
\item 将如此得到的阿贝尔群映射记为 $f^* : Z_k(Y) \to Z_{k + r}(X)$。
\end{enumerate}
\end{definition}

\noindent
开浸入是平坦的。这是平坦态射的一个重要但平凡的特殊情形。如果 $U \subset X$ 是开集，
有时把循环沿 $j : U \to X$ 的拉回称为该循环在 $U$ 上的{\it 限制}。注意，此时映射
$$
j^* : Z_k(X) \longrightarrow Z_k(U)
$$
全都是{\it 满射}。理由如下：给定任意整闭子概形 $Z' \subset U$，
可以在 $X$ 中取 $Z'$ 的闭包 $Z$，并把它看成 $X$ 的约化闭子概形
（见《概形》引理 \ref{schemes-lemma-reduced-closed-subscheme}）。
显然 $Z \cap U = Z'$，换言之 $j^*[Z] = [Z']$，故得满射性。事实上还有稍强的结论。

\begin{lemma}
\label{chow-lemma-exact-sequence-open}
设 $(S, \delta)$ 如情形 \ref{chow-situation-setup}，设 $X$ 在 $S$ 上局部有限型。
设 $U \subset X$ 是开子概形，并把 $i : Y = X \setminus U \to X$ 视作 $X$ 的约化闭子概形。
对每个 $k \in \mathbf{Z}$，阿贝尔群复形
$$
\xymatrix{
Z_k(Y) \ar[r]^{i_*} & Z_k(X) \ar[r]^{j^*} & Z_k(U) \ar[r] & 0
}
$$
正合。
\end{lemma}

\begin{proof}
先假设 $X$ 拟紧。那么 $Z_k(X)$ 是自由 $\mathbf{Z}$-模；它的一组基由元素 $[Z]$ 给出，
其中 $Z \subset X$ 是 $\delta$-维数为 $k$ 的整闭子概形。
这样的基元素映到基元素 $[Z \cap U]$，或者在 $Z \subset Y$ 时映为零。
故此情形下引理显然。一般情形相似，略去证明。
\end{proof}

\begin{lemma}
\label{chow-lemma-compose-flat-pullback}
设 $(S, \delta)$ 如情形 \ref{chow-situation-setup}，设 $X, Y, Z$ 在 $S$ 上局部有限型。
设 $f : X \to Y$、$g : Y \to Z$ 分别是相对维数为 $r$、$s$ 的平坦态射。
那么 $g \circ f$ 是相对维数为 $r + s$ 的平坦态射，并且
$$
f^* \circ g^* = (g \circ f)^*
$$
作为映射 $Z_k(Z) \to Z_{k + r + s}(X)$ 成立。
\end{lemma}

\begin{proof}
由《概形态射》引理 \ref{morphisms-lemma-composition-relative-dimension-d}，
复合是相对维数为 $r + s$ 的平坦态射。假设
\begin{enumerate}
\item $W \subset Z$ 是 $\delta$-维数为 $k$ 的闭整子概形；
\item $W' \subset Y$ 是 $\delta$-维数为 $k + s$ 的闭整子概形，且
$W' \subset g^{-1}(W)$；以及
\item $W'' \subset Y$ 是 $\delta$-维数为 $k + s + r$ 的闭整子概形，且
$W'' \subset f^{-1}(W')$。
\end{enumerate}
需要证明，$[W'']$ 在 $(g \circ f)^*[W]$ 中的系数 $n$，等于
$[W'']$ 在 $f^*(g^*[W])$ 中的系数 $m$。由引理 \ref{chow-lemma-flat-inverse-image-dimension}，只核查这些情形已经足够。
设 $\xi'' \in W''$、$\xi' \in W'$、$\xi \in W$ 分别为一般点。考虑局部环
$A = \mathcal{O}_{Z, \xi}$、$B = \mathcal{O}_{Y, \xi'}$ 和
$C = \mathcal{O}_{X, \xi''}$。于是有平坦局部环同态 $A \to B$、$B \to C$，并且
$$
n = \text{length}_C(C/\mathfrak m_AC),
\quad
\text{且}
\quad
m = \text{length}_C(C/\mathfrak m_BC) \text{length}_B(B/\mathfrak m_AB)
$$
所需等式由《交换代数》引理 \ref{algebra-lemma-pullback-transitive} 推出。
\end{proof}

\begin{lemma}
\label{chow-lemma-pullback-coherent}
设 $(S, \delta)$ 如情形 \ref{chow-situation-setup}，设 $X, Y$ 在 $S$ 上局部有限型，
并设 $f : X \to Y$ 是相对维数为 $r$ 的平坦态射。
\begin{enumerate}
\item 设 $Z \subset Y$ 是满足 $\dim_\delta(Z) \leq k$ 的闭子概形。那么
$\dim_\delta(f^{-1}(Z)) \leq k + r$，并且在 $Z_{k + r}(X)$ 中有
$[f^{-1}(Z)]_{k + r} = f^*[Z]_k$。
\item 设 $\mathcal{F}$ 是 $Y$ 上满足
$\dim_\delta(\text{Supp}(\mathcal{F})) \leq k$ 的凝聚层。那么
$\dim_\delta(\text{Supp}(f^*\mathcal{F})) \leq k + r$，并且
$$
f^*[{\mathcal F}]_k = [f^*{\mathcal F}]_{k+r}
$$
在 $Z_{k + r}(X)$ 中成立。
\end{enumerate}
\end{lemma}

\begin{proof}
关于维数的断言立即来自引理 \ref{chow-lemma-flat-inverse-image-dimension}。
(1) 由 (2)、引理 \ref{chow-lemma-cycle-closed-coherent} 以及等式
$f^*\mathcal{O}_Z = \mathcal{O}_{f^{-1}(Z)}$ 推出。

\medskip\noindent
下面证明 (2)。由于 $X$、$Y$ 局部诺特，可以应用《概形上同调》引理 \ref{coherent-lemma-coherent-Noetherian}，得知 $\mathcal{F}$ 有限型；所以
$f^*\mathcal{F}$ 有限型（《模层》引理 \ref{modules-lemma-pullback-finite-type}），
进而 $f^*\mathcal{F}$ 凝聚（再次应用《概形上同调》引理 \ref{coherent-lemma-coherent-Noetherian}）。因此该引理有意义。
设 $W \subset Y$ 是 $\delta$-维数为 $k$ 的整闭子概形，并设 $W' \subset X$ 是维数为
$k + r$、经 $f$ 映入 $W$ 的整闭子概形。需要证明，$[W']$ 在
$f^*[{\mathcal F}]_k$ 中的系数 $n$ 等于 $[W']$ 在 $[f^*{\mathcal F}]_{k+r}$ 中的系数 $m$。
设 $\xi \in W$、$\xi' \in W'$ 为一般点。令
$A = \mathcal{O}_{Y, \xi}$、$B = \mathcal{O}_{X, \xi'}$，并把
$M = \mathcal{F}_\xi$ 看作 $A$-模。注意，由维数假设，$M$ 具有有限长度，
但实际上无需验证这一点；见引理 \ref{chow-lemma-length-finite}。
有 $f^*\mathcal{F}_{\xi'} = B \otimes_A M$，从而
$$
n = \text{length}_B(B \otimes_A M)
\quad
\text{且}
\quad
m = \text{length}_A(M) \text{length}_B(B/\mathfrak m_AB)
$$
所需等式来自《交换代数》引理 \ref{algebra-lemma-pullback-module}。
\end{proof}



\section{推前与拉回}
\label{chow-section-push-pull}

\noindent
本节验证，在有意义时，固有推前与平坦拉回彼此相容。由上面的工作，这归结为上同调与基变换。

\begin{lemma}
\label{chow-lemma-flat-pullback-proper-pushforward}
设 $(S, \delta)$ 如情形 \ref{chow-situation-setup}。设
$$
\xymatrix{
X' \ar[r]_{g'} \ar[d]_{f'} & X \ar[d]^f \\
Y' \ar[r]^g & Y
}
$$
是在 $S$ 上局部有限型的概形的纤维积图。假设 $f : X \to Y$ 固有，
而 $g : Y' \to Y$ 是相对维数为 $r$ 的平坦态射。那么 $f'$ 也固有，
$g'$ 也是相对维数为 $r$ 的平坦态射。对 $X$ 上任意 $k$-循环 $\alpha$，有
$$
g^*f_*\alpha = f'_*(g')^*\alpha
$$
在 $Z_{k + r}(Y')$ 中成立。
\end{lemma}

\begin{proof}
$f'$ 固有的断言来自《概形态射》引理 \ref{morphisms-lemma-base-change-proper}。
$g'$ 平坦且相对维数为 $r$ 的断言来自《概形态射》引理 \ref{morphisms-lemma-base-change-relative-dimension-d} 和
\ref{morphisms-lemma-base-change-flat}。只须在 $\alpha = [W]$ 时证明循环等式，
其中 $W \subset X$ 是 $\delta$-维数为 $k$ 的整闭子概形。
注意，此时有 $\alpha = [\mathcal{O}_W]_k$；见引理 \ref{chow-lemma-cycle-closed-coherent}。所以由引理 \ref{chow-lemma-cycle-push-sheaf} 和 \ref{chow-lemma-pullback-coherent}，
只须证明 $f'_*(g')^*\mathcal{O}_W$ 同构于 $g^*f_*\mathcal{O}_W$。
这来自上同调与基变换；见《概形上同调》引理 \ref{coherent-lemma-flat-base-change-cohomology}。
\end{proof}

\begin{lemma}
\label{chow-lemma-finite-flat}
设 $(S, \delta)$ 如情形 \ref{chow-situation-setup}，设 $X$、$Y$ 在 $S$ 上局部有限型。
设 $f : X \to Y$ 是次数为 $d$ 的有限局部自由态射
（见《概形态射》定义 \ref{morphisms-definition-finite-locally-free}）。
那么 $f$ 既固有，又是相对维数为 $0$ 的平坦态射，并且对每个
$\alpha \in Z_k(Y)$，有
$$
f_*f^*\alpha = d\alpha
$$
\end{lemma}

\begin{proof}
由《概形态射》引理 \ref{morphisms-lemma-finite-flat}，有限局部自由态射平坦且有限；
由《概形态射》引理 \ref{morphisms-lemma-finite-proper}，有限态射固有。
我们略去有限态射的相对维数为 $0$ 的证明。因此该公式有意义。
为证明它，设 $Z \subset Y$ 是 $\delta$-维数为 $k$ 的整闭子概形。
只须对 $\alpha = [Z]$ 证明公式。有限局部自由态射的基变换仍有限局部自由
（《概形态射》引理 \ref{morphisms-lemma-base-change-finite-locally-free}），
故 $f_*f^*\mathcal{O}_Z$ 是 $Z$ 上秩为 $d$ 的有限局部自由层。因此
$$
f_*f^*[Z] = f_*f^*[\mathcal{O}_Z]_k =
[f_*f^*\mathcal{O}_Z]_k = d[Z]
$$
这里使用了引理 \ref{chow-lemma-pullback-coherent} 和
\ref{chow-lemma-cycle-push-sheaf}。
\end{proof}








\section{主除子的准备}
\label{chow-section-preparation-principal-divisors}

\noindent
本节的一部分材料与《除子》第 \ref{divisors-section-Weil-divisors} 节
中的讨论有所重叠。

\begin{lemma}
\label{chow-lemma-divisor-delta-dimension}
设 $(S, \delta)$ 如情形 \ref{chow-situation-setup}。
设 $X$ 在 $S$ 上局部有限型，并假设 $X$ 是整概形。
\begin{enumerate}
\item 若 $Z \subset X$ 是整闭子概形，则下列条件等价：
\begin{enumerate}
\item $Z$ 是素除子；
\item $Z$ 在 $X$ 中的余维数为 $1$；
\item $\dim_\delta(Z) = \dim_\delta(X) - 1$。
\end{enumerate}
\item 若 $Z$ 是 $X$ 上某个有效 Cartier 除子的不可约分支，则
$\dim_\delta(Z) = \dim_\delta(X) - 1$。
\end{enumerate}
\end{lemma}

\begin{proof}
(1) 由素除子的定义（《除子》定义 \ref{divisors-definition-Weil-divisor}）和维数函数的定义（《拓扑学》定义 \ref{topology-definition-dimension-function}）立即得到。
设 $\xi \in Z$ 是有效 Cartier 除子 $D \subset X$ 的某个不可约分支 $Z$ 的一般点。
那么 $\dim(\mathcal{O}_{D, \xi}) = 0$，并且对某个非零因子
$f \in \mathcal{O}_{X, \xi}$，有
$\mathcal{O}_{D, \xi} = \mathcal{O}_{X, \xi}/(f)$（《除子》引理 \ref{divisors-lemma-effective-Cartier-in-points}）。
由《交换代数》引理 \ref{algebra-lemma-one-equation}，
$\dim(\mathcal{O}_{X, \xi}) = 1$。于是《概形的性质》引理 \ref{properties-lemma-codimension-local-ring} 表明 $Z$ 满足 (1)，证明完毕。
\end{proof}

\begin{lemma}
\label{chow-lemma-finite-in-codimension-one}
设 $f : X \to Y$ 是概形态射，$\xi \in Y$ 是一点。假设
\begin{enumerate}
\item $X$、$Y$ 都是整概形；
\item $Y$ 局部诺特；
\item $f$ 固有且支配，并且 $R(Y) \subset R(X)$ 是有限扩张；
\item $\dim(\mathcal{O}_{Y, \xi}) = 1$。
\end{enumerate}
那么存在 $\xi$ 的开邻域 $V \subset Y$，使得
$f|_{f^{-1}(V)} : f^{-1}(V) \to V$ 是有限态射。
\end{lemma}

\begin{proof}
本引理是《代数簇》引理 \ref{varieties-lemma-finite-in-codim-1} 的特例。
这里给出本情形下的直接论证。由《概形上同调》引理 \ref{coherent-lemma-proper-finite-fibre-finite-in-neighbourhood}，
只需证明 $f^{-1}(\{\xi\})$ 是有限集。用一个仿射开集替换 $Y$，记为
$Y = \Spec(R)$。由于假设 $Y$ 局部诺特，$R$ 是诺特环。
由于 $f$ 固有，它是拟紧的；因而可以取有限仿射开覆盖
$X = U_1 \cup \ldots \cup U_n$，其中每个 $U_i = \Spec(A_i)$。
注意，$R \to A_i$ 是整环之间有限型的单同态，且其诱导的分式域扩张有限。
因此该引理由《交换代数》引理 \ref{algebra-lemma-finite-in-codim-1} 得出。
\end{proof}


\section{主除子}
\label{chow-section-principal-divisors}

\noindent
下面的定义是《除子》定义 \ref{divisors-definition-principal-divisor} 在当前设定下的对应形式。

\begin{definition}
\label{chow-definition-principal-divisor}
设 $(S, \delta)$ 如情形 \ref{chow-situation-setup}。
设 $X$ 在 $S$ 上局部有限型，并假设 $X$ 是整概形且
$\dim_\delta(X) = n$。设 $f \in R(X)^*$。定义{\it 与 $f$ 相伴的主除子}为
$(n - 1)$-循环
$$
\text{div}(f) = \text{div}_X(f) = \sum \text{ord}_Z(f) [Z]
$$
其中右端定义见《除子》定义 \ref{divisors-definition-principal-divisor}。由引理 \ref{chow-lemma-divisor-delta-dimension}，素除子的 $\delta$-维数为 $n - 1$，
故此定义有意义。
\end{definition}

\noindent
在该定义的情形下，对 $f, g \in R(X)^*$，在 $Z_{n - 1}(X)$ 中有
$$
\text{div}_X(fg) = \text{div}_X(f) + \text{div}_X(g)
$$
见《除子》引理 \ref{divisors-lemma-div-additive}。
下面的引理将被更一般的引理 \ref{chow-lemma-flat-pullback-rational-equivalence} 所涵盖。

\begin{lemma}
\label{chow-lemma-flat-pullback-principal-divisor}
设 $(S, \delta)$ 如情形 \ref{chow-situation-setup}。
设 $X$、$Y$ 在 $S$ 上局部有限型，并假设 $X$、$Y$ 都是整概形，且
$n = \dim_\delta(Y)$。设 $f : X \to Y$ 是相对维数为 $r$ 的平坦态射，
$g \in R(Y)^*$。那么在 $Z_{n + r - 1}(X)$ 中有
$$
f^*(\text{div}_Y(g)) = \text{div}_X(g)
$$
\end{lemma}

\begin{proof}
由于 $f$ 平坦，它是支配态射，故 $f$ 诱导嵌入
$R(Y) \subset R(X)$；因而可以把 $g$ 看成 $R(X)^*$ 的元素。
设 $Z \subset X$ 是 $\delta$-维数为 $n + r - 1$ 的整闭子概形，
$\xi \in Z$ 是其一般点。若 $\dim_\delta(f(Z)) > n - 1$，
则方程两端 $[Z]$ 的系数都为零。因此可以假设
$Z' = \overline{f(Z)}$ 是 $Y$ 中 $\delta$-维数为 $n - 1$ 的整闭子概形。
令 $\xi' = f(\xi)$；它是 $Z'$ 的一般点。置
$A = \mathcal{O}_{Y, \xi'}$、$B = \mathcal{O}_{X, \xi}$。
环同态 $A \to B$ 是维数为 $1$ 的诺特局部整环之间的平坦局部同态。
$g$ 属于 $A$ 的分式域。所需证明的是
$$
\text{ord}_A(g) \text{length}_B(B/\mathfrak m_AB)
=
\text{ord}_B(g).
$$
这由《交换代数》引理 \ref{algebra-lemma-pullback-module} 得出
（细节从略）。
\end{proof}











\section{主除子与推前}
\label{chow-section-two-fun}

\noindent
第一个引理表明，沿一般有限态射推前主除子，所得仍是主除子。

\begin{lemma}
\label{chow-lemma-proper-pushforward-alteration}
设 $(S, \delta)$ 如情形 \ref{chow-situation-setup}。
设 $X$、$Y$ 在 $S$ 上局部有限型，并假设 $X$、$Y$ 都是整概形，且
$n = \dim_\delta(X) = \dim_\delta(Y)$。设 $p : X \to Y$ 是支配的固有态射，
$f \in R(X)^*$。置
$$
g = \text{Nm}_{R(X)/R(Y)}(f).
$$
那么
$p_*\text{div}(f) = \text{div}(g)$。
\end{lemma}

\begin{proof}
设 $Z \subset Y$ 是 $\delta$-维数为 $n - 1$ 的整闭子概形。
需要证明 $[Z]$ 在 $p_*\text{div}(f)$ 与 $\text{div}(g)$ 中的系数相等。
可以把引理 \ref{chow-lemma-finite-in-codimension-one} 用于态射
$p : X \to Y$ 和一般点 $\xi \in Z$。因此可以用 $\xi$ 的仿射开邻域替换 $Y$，
并假设 $p : X \to Y$ 是有限态射。写成 $Y = \Spec(R)$、$X = \Spec(A)$，
其中 $p$ 由诺特整环间的有限同态 $R \to A$ 诱导；该同态诱导分式域的有限扩张
$L/K$。现在有 $f \in L$、$g = \text{Nm}(f) \in K$，以及满足
$\dim(R_{\mathfrak p}) = 1$ 的素理想 $\mathfrak p \subset R$。
$[Z]$ 在 $\text{div}_Y(g)$ 中的系数为 $\text{ord}_{R_\mathfrak p}(g)$，
而 $[Z]$ 在 $p_*\text{div}_X(f)$ 中的系数为
$$
\sum\nolimits_{\mathfrak q\text{ 位于 }\mathfrak p\text{ 之上}}
[\kappa(\mathfrak q) : \kappa(\mathfrak p)]
\text{ord}_{A_{\mathfrak q}}(f)
$$
所需等式于是由《交换代数》引理 \ref{algebra-lemma-finite-extension-dim-1} 得出。
\end{proof}

\noindent
在主除子的讨论中，$\mathbf{P}^1_S$ 上的“泛”主除子
$[0] - [\infty]$ 起着重要作用。为更精确地表述这一点，以
\begin{equation}
\label{chow-equation-zero-infty}
D_0, D_\infty \subset
\mathbf{P}^1_S = \underline{\text{Proj}}_S(\mathcal{O}_S[T_0, T_1])
\end{equation}
表示分别由 $\mathcal{O}(1)$ 的截面 $T_1$ 与\ $T_0$ 截出的闭子概形。
它们是有效 Cartier 除子；见《除子》定义 \ref{divisors-definition-effective-Cartier-divisor} 和引理 \ref{divisors-lemma-characterize-OD}。下面的引理说，粗略而言，
“$\text{div}(T_1/T_0) = [D_0] - [D_1]$”，并且这是泛主除子。

\begin{lemma}
\label{chow-lemma-rational-function}
设 $(S, \delta)$ 如情形 \ref{chow-situation-setup}。
设 $X$ 在 $S$ 上局部有限型，并假设 $X$ 是整概形，且
$n = \dim_\delta(X)$。设 $f \in R(X)^*$。取非空开集 $U \subset X$，
使 $f$ 对应于一个截面 $f \in \Gamma(U, \mathcal{O}_X^*)$。
令 $Y \subset X \times_S \mathbf{P}^1_S$ 为
$f : U \to \mathbf{P}^1_S$ 的图的闭包。那么
\begin{enumerate}
\item 投影态射 $p : Y \to X$ 是固有的；
\item $p|_{p^{-1}(U)} : p^{-1}(U) \to U$ 是同构；
\item 沿态射 $q : Y \to \mathbf{P}^1_S$ 的拉回
$Y_0 = q^{-1}D_0$ 和 $Y_\infty = q^{-1}D_\infty$ 有定义
（《除子》定义 \ref{divisors-definition-pullback-effective-Cartier-divisor}）；
\item 有
$$
\text{div}_Y(f) = [Y_0]_{n - 1} - [Y_\infty]_{n - 1}
$$
\item 有
$$
\text{div}_X(f) = p_*\text{div}_Y(f)
$$
\item 若通过态射 $p$ 把 $Y_0$ 和 $Y_\infty$ 看成 $X$ 的闭子概形，则有
$$
\text{div}_X(f) = [Y_0]_{n - 1} - [Y_\infty]_{n - 1}
$$
\end{enumerate}
\end{lemma}

\begin{proof}
由于 $X$ 是整概形，$U$ 也是整概形。因此 $Y$ 是整概形，且
$(1, f)(U) \subset Y$ 是稠密开子概形。还要注意，闭子概形
$Y \subset X \times_S \mathbf{P}^1_S$ 与开集 $U$ 的选择无关，
因为它终究是单点集 $\{\eta'\} = \{(1, f)(\eta)\}$ 的闭包，
其中 $\eta \in X$ 是一般点。下面证明引理的各个断言。

\medskip\noindent
对于 (1)，注意 $p$ 是闭浸入
$Y \to X \times_S \mathbf{P}^1_S = \mathbf{P}^1_X$ 与固有态射
$\mathbf{P}^1_X \to X$ 的复合。固有态射的复合仍固有
（《概形态射》引理 \ref{morphisms-lemma-composition-proper}），故结论成立。

\medskip\noindent
显然 $Y \cap U \times_S \mathbf{P}^1_S = (1, f)(U)$，所以 (2) 成立；
同时也得到 $\dim_\delta(Y) = n$。

\medskip\noindent
由于 $f \in R(X)^*$，$q(\eta') = f(\eta)$ 不在 $D_0$ 或 $D_\infty$ 中。
因此 (3) 由《除子》引理 \ref{divisors-lemma-pullback-effective-Cartier-defined} 得到。
再由引理 \ref{chow-lemma-divisor-delta-dimension} 得到
$\dim_\delta(Y_0) = n - 1$ 和 $\dim_\delta(Y_\infty) = n - 1$。

\medskip\noindent
考虑有效 Cartier 除子 $Y_0$。在每一点 $\xi \in Y_0$，都有
$f \in \mathcal{O}_{Y, \xi}$，且 $Y_0$ 的局部方程由 $f$ 给出。
特别地，若 $\delta(\xi) = n - 1$，从而 $\xi$ 是某个
$\delta$-维数为 $n - 1$ 的整闭子概形 $Z$ 的一般点，则 $[Z]$ 在
$\text{div}_Y(f)$ 中的系数为
$$
\text{ord}_Z(f) =
\text{length}_{\mathcal{O}_{Y, \xi}}
(\mathcal{O}_{Y, \xi}/f\mathcal{O}_{Y, \xi}) =
\text{length}_{\mathcal{O}_{Y, \xi}}
(\mathcal{O}_{Y_0, \xi})
$$
这正是 $[Z]$ 在 $[Y_0]_{n - 1}$ 中的系数。对有理函数 $1/f$ 作类似论证，
可知 $-[Y_\infty]$ 对应于 $\text{div}_Y(f)$ 表达式中系数为负的各项。
故 (4) 成立。

\medskip\noindent
注意 $D_0 \to S$ 是同构，因而 $X \times_S D_0 \to X$ 也是同构。
显然，在 $X \times_S \mathbf{P}^1_S$ 中有
$Y_0 = Y \cap X \times_S D_0$（概形论交），故 $Y_0 \to X$ 确实是闭浸入。
于是
$$
p_*\mathcal{O}_{Y_0} = \mathcal{O}_{Y'_0}
$$
其中 $Y'_0 \subset X$ 是 $Y_0 \to X$ 的像。由引理 \ref{chow-lemma-cycle-push-sheaf}，有
$p_*[Y_0]_{n - 1} = [Y'_0]_{n - 1}$。对 $D_\infty$ 和 $Y_\infty$ 同样成立，
故 (6) 是 (5) 的推论。最后，(5) 立即由引理 \ref{chow-lemma-proper-pushforward-alteration} 得到。
\end{proof}

\noindent
下面的引理表明，固有曲线上主除子的次数为零。

\begin{lemma}
\label{chow-lemma-curve-principal-divisor}
设 $K$ 是任意域。设 $X$ 是一个 $1$-维整概形，并带有固有态射
$c : X \to \Spec(K)$。设 $f \in K(X)^*$ 是可逆有理函数。那么
$$
\sum\nolimits_{x \in X \text{ 为闭点}}
[\kappa(x) : K] \text{ord}_{\mathcal{O}_{X, x}}(f)
=
0
$$
其中 $\text{ord}$ 如《交换代数》定义 \ref{algebra-definition-ord} 所定义。
换言之，$c_*\text{div}(f) = 0$。
\end{lemma}

\begin{proof}
考虑图表
$$
\xymatrix{
Y \ar[r]_p \ar[d]_q & X \ar[d]^c \\
\mathbf{P}^1_K \ar[r]^-{c'} & \Spec(K)
}
$$
这是在引理 \ref{chow-lemma-rational-function} 中，从 $S = \Spec(K)$ 上的
$X$ 和有理函数 $f$ 出发构造的。下文不再说明而直接使用该引理的全部结论。
需要证明 $c_*\text{div}_X(f) = c_*p_*\text{div}_Y(f) = 0$，
这等价于证明 $c'_*q_*\text{div}_Y(f) = 0$。
若 $q(Y)$ 是 $\mathbf{P}^1_K$ 的闭点，则 $\text{div}_X(f) = 0$，引理成立。
所以可以假设 $q$ 是支配态射。假设能证明
$q : Y \to \mathbf{P}^1_K$ 是次数为 $d$ 的有限局部自由态射
（见《概形态射》定义 \ref{morphisms-definition-finite-locally-free}）。
由于 $\text{div}_Y(f) = [q^{-1}D_0]_0 - [q^{-1}D_\infty]_0$，
由平坦拉回的定义可知
$\text{div}_Y(f) = q^*([D_0]_0 - [D_\infty]_0)$。
再由引理 \ref{chow-lemma-finite-flat} 得到
$q_*\text{div}_Y(f) = d([D_0]_0 - [D_\infty]_0)$。
而显然 $c'_*[D_0]_0 = c'_*[D_\infty]_0$，结论即得。

\medskip\noindent
还须证明 $q$ 是有限局部自由态射。（由于 $\mathbf{P}^1_K$ 连通，
它会自动具有某个固定次数。）因为 $\dim(\mathbf{P}^1_K) = 1$，
例如由引理 \ref{chow-lemma-finite-in-codimension-one} 可知 $q$ 是有限态射。
$\mathbf{P}^1_K$ 在各闭点处的局部环都是维数为 $1$ 的正则局部环
（即离散赋值环），因为它们是 $K[T]$ 的局部化；见《交换代数》引理 \ref{algebra-lemma-dim-affine-space}。因此，对闭点 $y\in Y$，只要
$\mathcal{O}_{Y, y}$ 无挠，它在 $\mathcal{O}_{\mathbf{P}^1_K, q(y)}$ 上就是平坦的
（《代数进阶》引理 \ref{more-algebra-lemma-dedekind-torsion-free-flat}）。
这显然成立，因为 $\mathcal{O}_{Y, y}$ 是整环且 $q$ 支配。
所以 $q$ 平坦。最后由《概形态射》引理 \ref{morphisms-lemma-finite-flat}，$q$ 是有限局部自由态射。
\end{proof}





\section{有理等价}
\label{chow-section-rational-equivalence}

\noindent
本节定义 $k$-循环上的{\it 有理等价}。我们将允许主除子的像
（沿闭浸入）的局部有限和。这会产生一些颇为奇特的现象；见例 \ref{chow-example-weird}。然而，若不允许这种和，我们便不知道如何证明线丛陈类的帽积
可下降到有理等价类。

\begin{definition}
\label{chow-definition-rational-equivalence}
设 $(S, \delta)$ 如情形 \ref{chow-situation-setup}。
设概形 $X$ 在 $S$ 上局部有限型，并设 $k \in \mathbf{Z}$。
\begin{enumerate}
\item 给定任意局部有限族 $\{W_j \subset X\}$，其中各成员是满足
$\dim_\delta(W_j) = k + 1$ 的整闭子概形，并给定任意
$f_j \in R(W_j)^*$，可以考虑
$$
\sum (i_j)_*\text{div}(f_j) \in Z_k(X)
$$
其中 $i_j : W_j \to X$ 是包含态射。该表达式有意义，因为态射
$\coprod i_j : \coprod W_j \to X$ 是固有的。
\item 若 $\alpha$ 是上述形式的循环，就称 $\alpha \in Z_k(X)$
{\it 有理等价于零}。
\item 若 $\alpha - \beta$ 有理等价于零，就称
$\alpha, \beta \in Z_k(X)$ {\it 有理等价}，并记作
$\alpha \sim_{rat} \beta$。
\item 定义
$$
\CH_k(X) = Z_k(X) / \sim_{rat}
$$
为{\it $X$ 上 $k$-循环的 Chow 群}。它有时也称为
{\it $X$ 上模有理等价的 $k$-循环 Chow 群}。
\end{enumerate}
\end{definition}

\noindent
还有许多其他有趣的（适当的）等价关系；有理等价是其中最粗的一个。

\begin{remark}
\label{chow-remark-chow-group-pointwise}
设 $(S, \delta)$ 如情形 \ref{chow-situation-setup}。
设 $X$ 在 $S$ 上局部有限型，并设 $k \in \mathbf{Z}$。下面说明存在表现
$$
\bigoplus\nolimits_{\delta(x) = k + 1}' K_1^M(\kappa(x))
\xrightarrow{\partial}
\bigoplus\nolimits_{\delta(x) = k}' K_0^M(\kappa(x)) \to
\CH_k(X) \to 0
$$
这里采用注 \ref{chow-remark-cycles-pointwise} 中引入的记号与约定，此外：
\begin{enumerate}
\item $K_1^M(\kappa(x)) = \kappa(x)^*$ 是点 $x \in X$ 的剩余域
$\kappa(x)$ 的 Milnor K-理论的次数 $1$ 部分；见注 \ref{chow-remark-gersten-complex-milnor}；
\item 微分 $\partial$ 定义如下：给定元素 $\xi = \sum_x f_x$，
以 $W_x = \overline{x}$ 表示 $X$ 中以 $x$ 为一般点的整闭子概形，并置
$$
\partial(\xi) = \sum (W_x \to X)_*\text{div}(f_x)
$$
其属于 $Z_k(X)$。这是有意义的，因为注 \ref{chow-remark-cycles-pointwise} 已说明该复形的第二项等于 $Z_k(X)$。
\end{enumerate}
与定义 \ref{chow-definition-rational-equivalence} 比较，立即可知这确实给出
$\CH_k(X)$ 的一个表现。
\end{remark}

\noindent
下面是一个非常简单但很重要的引理。

\begin{lemma}
\label{chow-lemma-restrict-to-open}
设 $(S, \delta)$ 如情形 \ref{chow-situation-setup}。
设概形 $X$ 在 $S$ 上局部有限型。设 $U \subset X$ 是开子概形，并把
$i : Y = X \setminus U \to X$ 记为 $X$ 的约化闭子概形。设
$k \in \mathbf{Z}$，并设 $\alpha, \beta \in Z_k(X)$。
若 $\alpha|_U \sim_{rat} \beta|_U$，则存在循环
$\gamma \in Z_k(Y)$，使得
$$
\alpha \sim_{rat} \beta + i_*\gamma.
$$
换言之，Abel 群序列
$$
\xymatrix{
\CH_k(Y) \ar[r]^{i_*} & \CH_k(X) \ar[r]^{j^*} & \CH_k(U) \ar[r] & 0
}
$$
正合。
\end{lemma}

\begin{proof}
设 $\{W_j\}_{j \in J}$ 是 $U$ 中 $\delta$-维数为 $k + 1$ 的整闭子概形的
局部有限族，并设 $f_j \in R(W_j)^*$ 满足定义中的等式
$(\alpha - \beta)|_U = \sum (i_j)_*\text{div}(f_j)$。
令 $W_j' \subset X$ 为 $W_j$ 的闭包。设 $V \subset X$ 是拟紧开集。
由于 $V$ 是诺特的，$V \cap U$ 也是 $U$ 中的拟紧开集。因此集合
$\{j \in J \mid W_j \cap V \not = \emptyset\}
= \{j \in J \mid W'_j \cap V \not = \emptyset\}$
是有限的，因为 $\{W_j\}$ 局部有限。换言之，$\{W'_j\}$ 也局部有限。
由于 $R(W_j) = R(W'_j)$，可知
$$
\alpha - \beta - \sum (i'_j)_*\text{div}(f_j)
$$
是支撑在 $Y$ 上的循环，引理即得；见引理 \ref{chow-lemma-exact-sequence-open}。
\end{proof}

\begin{lemma}
\label{chow-lemma-exact-sequence-closed-chow}
设 $(S, \delta)$ 如情形 \ref{chow-situation-setup}。
设 $X$ 在 $S$ 上局部有限型。设 $X_1, X_2 \subset X$ 为闭子概形，
并且集合论上 $X = X_1 \cup X_2$。对每个 $k \in \mathbf{Z}$，Abel 群序列
$$
\xymatrix{
\CH_k(X_1 \cap X_2) \ar[r] &
\CH_k(X_1) \oplus \CH_k(X_2) \ar[r] &
\CH_k(X) \ar[r] &
0
}
$$
正合。这里 $X_1 \cap X_2$ 是概形论交；各映射均为推前映射，其中一个乘以
$-1$。
\end{lemma}

\begin{proof}
由引理 \ref{chow-lemma-exact-sequence-closed}，箭头
$\CH_k(X_1) \oplus \CH_k(X_2) \to \CH_k(X)$ 是满射。
假设 $(\alpha_1, \alpha_2)$ 在此映射下为零。写成
$\alpha_1 = \sum n_{1, i}[W_{1, i}]$ 和
$\alpha_2 = \sum n_{2, i}[W_{2, i}]$。于是得到 $X$ 中
$\delta$-维数为 $k + 1$ 的整闭子概形的局部有限族
$\{W_j\}_{j \in J}$，以及 $f_j \in R(W_j)^*$，使得
$$
\sum n_{1, i}[W_{1, i}] + \sum n_{2, i}[W_{2, i}] = \sum (i_j)_*\text{div}(f_j)
$$
作为 $X$ 上的循环成立，其中 $i_j : W_j \to X$ 为包含态射。
选择不交并分解 $J = J_1 \amalg J_2$，使得当 $j \in J_1$ 时
$W_j \subset X_1$，当 $j \in J_2$ 时 $W_j \subset X_2$。
这是可能的，因为各 $W_j$ 都是整概形。于是上述等式可写成
$$
\sum n_{1, i}[W_{1, i}] - \sum\nolimits_{j \in J_1} (i_j)_*\text{div}(f_j) =
- \sum n_{2, i}[W_{2, i}] + \sum\nolimits_{j \in J_2} (i_j)_*\text{div}(f_j)
$$
因此这个表达式是 $X_1 \cap X_2$ 上的一个循环（！）。换言之，
元素 $(\alpha_1, \alpha_2)$ 属于第一个箭头的像，证明完毕。
\end{proof}

\begin{example}
\label{chow-example-weird}
下面是一个“奇特”的例子。设 $S$ 是域 $k$ 的谱，$\delta$ 如例 \ref{chow-example-field}。设 $X = C_1 \cup C_2 \cup \ldots$ 是曲线
$C_j \cong \mathbf{P}^1_k$ 的无穷并，并按如下方式粘合：对
$j = 1, 2, 3, \ldots$，把点 $\infty \in C_j$ 与点
$0 \in C_{j + 1}$ 横截地粘合。取点 $0 \in C_1$，它给出零维循环
$[0] \in Z_0(X)$。本情形的“奇特”之处在于，事实上
$[0] \sim_{rat} 0$！具体地，可以选择有理函数 $f_j \in R(C_j)$，
使其在 $0$ 处有一个一阶零点，在 $\infty$ 处有一个一阶极点，且没有其他零点或极点。
于是和 $\sum (i_j)_*\text{div}(f_j)$ 恰好就是 $0$-循环 $[0]$。
事实上，在此例中 $\CH_0(X) = 0$。如果你觉得这过于离奇，只要始终要求所用空间拟紧即可
（这样 $X$ 对你而言根本不会出现）。
\end{example}

\begin{remark}
\label{chow-remark-infinite-sums-rational-equivalences}
设 $(S, \delta)$ 如情形 \ref{chow-situation-setup}。
设概形 $X$ 在 $S$ 上局部有限型。假设有 $X$ 上 $k$-循环的无穷族
$\alpha_i, \beta_i \in Z_k(X)$，$i \in I$。假设 $\alpha_i$ 与
$\beta_i$ 的支撑构成 $X$ 的闭子集的局部有限族，从而
$\sum \alpha_i$ 和 $\sum \beta_i$ 都定义了循环；再假设对每个 $i$，
$\alpha_i \sim_{rat} \beta_i$。此时并不能直接断定
$\sum \alpha_i \sim_{rat} \sum \beta_i$。问题在于，这些有理等价可能由局部有限族
$\{W_{i, j}, f_{i, j} \in R(W_{i, j})^*\}_{j \in J_i}$ 给出，
但它们的并 $\{W_{i, j}\}_{i \in I, j\in J_i}$ 未必局部有限。

\medskip\noindent
在许多实际情形中，存在闭子集的局部有限族 $\{T_i\}_{i \in I}$，
使 $\alpha_i, \beta_i$ 都支撑在 $T_i$ 上，并且在 $\CH_k(T_i)$ 中
$\alpha_i = \beta_i$；换言之，各族
$\{W_{i, j}, f_{i, j} \in R(W_{i, j})^*\}_{j \in J_i}$
由子概形 $W_{i, j} \subset T_i$ 组成。在这种情况下，$X$ 上确实有
$\sum \alpha_i \sim_{rat} \sum \beta_i$，因为此时族
$\{W_{i, j}\}_{i \in I, j\in J_i}$ 自动局部有限。
\end{remark}






\section{有理等价与推前、拉回}
\label{chow-section-properties-rational-equivalence}

\noindent
本节证明平坦拉回与固有推前都保持有理等价。

\begin{lemma}
\label{chow-lemma-prepare-flat-pullback-rational-equivalence}
设 $(S, \delta)$ 如情形 \ref{chow-situation-setup}。
设概形 $X$、$Y$ 在 $S$ 上局部有限型。假设 $Y$ 是整概形且
$\dim_\delta(Y) = k$。设 $f : X \to Y$ 是相对维数为 $r$ 的平坦态射。
那么对 $g \in R(Y)^*$，有
$$
f^*\text{div}_Y(g) =
\sum n_j i_{j, *}\text{div}_{X_j}(g \circ f|_{X_j})
$$
作为 $X$ 上的 $(k + r - 1)$-循环成立；其中求和遍历 $X$ 的不可约分支
$X_j$，而 $n_j$ 是 $X_j$ 在 $X$ 中的重数。
\end{lemma}

\begin{proof}
设 $Z \subset X$ 是 $\delta$-维数为 $k + r - 1$ 的整闭子概形。
需要证明 $[Z]$ 在 $f^*\text{div}(g)$ 中的系数 $n$，等于 $[Z]$ 在
$\sum i_{j, *} \text{div}(g \circ f|_{X_j})$ 中的系数 $m$。
设 $Z'$ 是 $f(Z)$ 的闭包；它是 $Y$ 的整闭子概形。由引理 \ref{chow-lemma-flat-inverse-image-dimension}，有
$\dim_\delta(Z') \geq k - 1$。因此，要么 $Z' = Y$，要么 $Z'$ 是
$Y$ 上的素除子。若 $Z' = Y$，则系数 $n$ 和 $m$ 都为零：
对 $n$，这由 $f^*$ 的定义显然；对 $m$，这是因为
$g \circ f|_{X_j}$ 在 $X_j$ 中任何映到 $Y$ 一般点的点处都是单位。
因此可以假设 $Z' \subset Y$ 是素除子。

\medskip\noindent
下面把 $n$ 与 $m$ 相等这一断言转化为代数等式。具体地，令 $\xi' \in Z'$ 和
$\xi \in Z$ 为一般点，置 $A = \mathcal{O}_{Y, \xi'}$、
$B = \mathcal{O}_{X, \xi}$。注意 $A$、$B$ 都是诺特环，$A \to B$
是平坦局部同态，$A$ 是整环，而 $\mathfrak m_AB$ 是局部环 $B$ 的定义理想。
有理函数 $g$ 是 $A$ 的分式域 $Q(A)$ 的元素。按构造，与 $\xi$ 相交的闭子概形
$X_j$ 以 $1$ 对 $1$ 的方式对应于极小素理想
$$
\mathfrak q_1, \ldots, \mathfrak q_s \subset B
$$
相应的整数 $n_j$ 是长度
$$
n_i = \text{length}_{B_{\mathfrak q_i}}(B_{\mathfrak q_i})
$$
有理函数 $g \circ f|_{X_j}$ 对应于 $g \in Q(A)$ 的像
$g_i \in \kappa(\mathfrak q_i)^*$。综合起来，有
$$
n = \text{ord}_A(g) \text{length}_B(B/\mathfrak m_AB)
$$
并且
$$
m = \sum \text{ord}_{B/\mathfrak q_i}(g_i)
\text{length}_{B_{\mathfrak q_i}}(B_{\mathfrak q_i})
$$
将该有理函数写成 $g = x/y$，其中非零的 $x, y \in A$。于是只需证明
$$
\text{length}_A(A/(x)) \text{length}_B(B/\mathfrak m_AB) =
\text{length}_B(B/xB)
$$
（此等式使用《交换代数》引理 \ref{algebra-lemma-pullback-module}）等于
$$
\sum\nolimits_{i = 1, \ldots, s}
\text{length}_{B/\mathfrak q_i}(B/(x, \mathfrak q_i))
\text{length}_{B_{\mathfrak q_i}}(B_{\mathfrak q_i})
$$
并对 $y$ 证明类似等式。由于 $A \to B$ 平坦，$x$ 是 $B$ 中的非零因子。
故所需等式由引理 \ref{chow-lemma-additivity-divisors-restricted} 得出。
\end{proof}

\begin{lemma}
\label{chow-lemma-flat-pullback-rational-equivalence}
设 $(S, \delta)$ 如情形 \ref{chow-situation-setup}。
设概形 $X$、$Y$ 在 $S$ 上局部有限型。设 $f : X \to Y$ 是相对维数为
$r$ 的平坦态射。设 $\alpha \sim_{rat} \beta$ 是 $Y$ 上有理等价的
$k$-循环。那么，作为 $X$ 上的 $(k + r)$-循环，
$f^*\alpha \sim_{rat} f^*\beta$。
\end{lemma}

\begin{proof}
需要证明什么？假设给定一族闭浸入
$$
i_j : W_j \longrightarrow Y
$$
其中每个 $W_j$ 都是 $\delta$-维数为 $k + 1$ 的整概形，并给定有理函数
$g_j \in R(W_j)^*$。再假设族 $\{i_j(W_j)\}_{j \in J}$ 在 $Y$ 上局部有限。
那么需要证明
$$
f^*(\sum i_{j, *}\text{div}(g_j)) = \sum f^*i_{j, *}\text{div}(g_j)
$$
在 $X$ 上有理等价于零。由引理 \ref{chow-lemma-inverse-image-locally-finite}，$\{W_j\}$ 在 $X$ 中局部有限，
故右端的和有意义。

\medskip\noindent
考虑纤维积
$$
i'_j : W'_j = W_j \times_Y X \longrightarrow X.
$$
并以 $f_j : W'_j \to W_j$ 表示第一投影。由引理 \ref{chow-lemma-flat-pullback-proper-pushforward}，上述和可写为
$$
\sum i'_{j, *}(f_j^*\text{div}(g_j))
$$
由引理 \ref{chow-lemma-prepare-flat-pullback-rational-equivalence}，
每个 $f_j^*\text{div}(g_j)$ 在 $W'_j$ 上有理等价于零。
因此每个 $i'_{j, *}(f_j^*\text{div}(g_j))$ 都有理等价于零。
由注 \ref{chow-remark-infinite-sums-rational-equivalences} 中的讨论，
展示的和也有同样的性质。
\end{proof}

\begin{lemma}
\label{chow-lemma-proper-pushforward-rational-equivalence}
设 $(S, \delta)$ 如情形 \ref{chow-situation-setup}。
设概形 $X$、$Y$ 在 $S$ 上局部有限型，并设 $p : X \to Y$ 是固有态射。
假设 $\alpha, \beta \in Z_k(X)$ 有理等价。那么 $p_*\alpha$ 与
$p_*\beta$ 有理等价。
\end{lemma}

\begin{proof}
需要证明什么？假设给定一族闭浸入
$$
i_j : W_j \longrightarrow X
$$
其中每个 $W_j$ 都是 $\delta$-维数为 $k + 1$ 的整概形，并给定有理函数
$f_j \in R(W_j)^*$。再假设族 $\{i_j(W_j)\}_{j \in J}$ 在 $X$ 上局部有限。
那么需要证明
$$
p_*\left(\sum i_{j, *}\text{div}(f_j)\right)
$$
在 $X$ 上有理等价于零。

\medskip\noindent
注意该和等于
$$
\sum p_*i_{j, *}\text{div}(f_j).
$$
设 $W'_j \subset Y$ 是 $p \circ i_j$ 的像这一整闭子概形。
由引理 \ref{chow-lemma-quasi-compact-locally-finite}，族 $\{W'_j\}$ 在 $Y$
中局部有限。因此，只须对给定的 $j$ 证明：要么
$p_*i_{j, *}\text{div}(f_j) = 0$，要么它等于
$i'_{j, *}\text{div}(g_j)$，其中某个 $g_j \in R(W'_j)^*$。

\medskip\noindent
上述论证把问题约化到一个 $\delta$-维数为 $k + 1$ 的整闭子概形
$W \subset X$。设 $f \in R(W)^*$，并如上令 $W' = p(W)$。
得到态射的交换图
$$
\xymatrix{
W \ar[r]_i \ar[d]_{p'} & X \ar[d]^p \\
W' \ar[r]^{i'} & Y
}
$$
由引理 \ref{chow-lemma-compose-pushforward}，
$p_*i_*\text{div}(f) = i'_*(p')_*\text{div}(f)$。如上所述，
需要证明 $(p')_*\text{div}(f)$ 要么是 $W'$ 上某个有理函数的除子，要么为零。
分三种情形。

\medskip\noindent
情形 $\dim_\delta(W') < k$。此时自动有 $(p')_*\text{div}(f) = 0$，
无须再证。

\medskip\noindent
情形 $\dim_\delta(W') = k$。下面证明此时 $(p')_*\text{div}(f) = 0$。
设 $\eta \in W'$ 为一般点。注意 $c : W_\eta \to \Spec(K)$ 是域
$K = \kappa(\eta)$ 上的固有整曲线，其函数域 $K(W_\eta)$ 与 $R(W)$ 等同。
有图表
$$
\xymatrix{
W_\eta \ar[r] \ar[d]_c & W \ar[d]^{p'} \\
\Spec(K) \ar[r] & W'
}
$$
以 $f_\eta \in K(W_\eta)^*$ 表示对应于 $f \in R(W)^*$ 的有理函数。
此外，$W_\eta$ 的闭点 $\xi$ 按 $1 - 1$ 对应于其中满足
$p'(Z) = W'$ 的 $\delta$-维数为 $k$ 的闭整子概形
$Z = Z_\xi \subset W$。注意 $Z_\xi$ 在 $\text{div}(f)$ 中的重数等于
$\text{ord}_{\mathcal{O}_{W_\eta, \xi}}(f_\eta)$，这只是因为局部环
$\mathcal{O}_{W_\eta, \xi}$ 与 $\mathcal{O}_{W, \xi}$ 作为各自分式域的子环彼此等同。
因此 $[W']$ 在 $(p')_*\text{div}(f)$ 中的重数等于
$[\Spec(K)]$ 在 $c_*\text{div}(f_\eta)$ 中的重数。
由引理 \ref{chow-lemma-curve-principal-divisor}，该重数为零。

\medskip\noindent
情形 $\dim_\delta(W') = k + 1$。此时可应用引理 \ref{chow-lemma-proper-pushforward-alteration}，从而确有
$p'_*\text{div}(f) = \text{div}(g)$，其中某个 $g \in R(W')^*$，正合所需。
\end{proof}















\section{有理等价与射影直线}
\label{chow-section-different-rational-equivalence}

\noindent
设 $(S, \delta)$ 如情形 \ref{chow-situation-setup}，并设概形 $X$ 在 $S$
上局部有限型。给定任意闭子概形
$Z \subset X \times_S \mathbf{P}^1_S = X \times \mathbf{P}^1$，
分别令 $Z_0$、\ $Z_\infty$ 为概形论闭子概形
$Z_0 = \text{pr}_2^{-1}(D_0)$、\ $Z_\infty = \text{pr}_2^{-1}(D_\infty)$。
这里 $D_0$、$D_\infty$ 如 (\ref{chow-equation-zero-infty}) 中所定义。

\begin{lemma}
\label{chow-lemma-rational-equivalence-family}
设 $(S, \delta)$ 如情形 \ref{chow-situation-setup}。
设概形 $X$ 在 $S$ 上局部有限型。设
$W \subset X \times_S \mathbf{P}^1_S$ 是 $\delta$-维数为 $k + 1$ 的整闭子概形。
假设 $W \not = W_0$ 且 $W \not = W_\infty$。那么
\begin{enumerate}
\item $W_0$、$W_\infty$ 是 $W$ 的有效 Cartier 除子；
\item 可以把 $W_0$、$W_\infty$ 看成 $X$ 的闭子概形，并且
$$
[W_0]_k \sim_{rat} [W_\infty]_k,
$$
\item 对任意局部有限族的整闭子概形
$W_i \subset X \times_S \mathbf{P}^1_S$，若其 $\delta$-维数为 $k + 1$，
且 $W_i \not = (W_i)_0$、$W_i \not = (W_i)_\infty$，则在 $X$ 上有
$\sum ([(W_i)_0]_k - [(W_i)_\infty]_k) \sim_{rat} 0$；
\item 对任意满足 $\alpha \sim_{rat} 0$ 的 $\alpha \in Z_k(X)$，
存在如上的整闭子概形局部有限族
$W_i \subset X \times_S \mathbf{P}^1_S$，使得
$\alpha = \sum ([(W_i)_0]_k - [(W_i)_\infty]_k)$。
\end{enumerate}
\end{lemma}

\begin{proof}
由假设，$W$ 的一般点在投影
$X \times_S \mathbf{P}^1_S \to \mathbf{P}^1_S$ 下不映入
$D_0$ 或 $D_\infty$，所以 (1) 由《除子》引理 \ref{divisors-lemma-pullback-effective-Cartier-defined} 得出。

\medskip\noindent
由于 $X \times_S D_0 \to X$ 是闭浸入，$W_0$ 同构于 $X$ 的某个闭子概形；
$W_\infty$ 同理。态射 $p : W \to X$ 是闭浸入
$W \to X \times_S \mathbf{P}^1_S$ 与固有态射
$X \times_S \mathbf{P}^1_S \to X$ 的复合，故它固有。
由引理 \ref{chow-lemma-rational-function}，作为 $W$ 上的循环，有
$[W_0]_k \sim_{rat} [W_\infty]_k$。而显然
$p_*[W_0]_k = [W_0]_k$，对 $W_\infty$ 也一样，所以 (2) 由引理 \ref{chow-lemma-proper-pushforward-rational-equivalence} 得出。

\medskip\noindent
在已知 (1)、(2) 后，断言 (3) 唯一需要说明的是：族
$\{(W_i)_0, (W_i)_\infty\}$ 是 $X$ 的闭子概形的局部有限族。这是显然的。

\medskip\noindent
假设 $\alpha \sim_{rat} 0$。按定义，存在 $\delta$-维数为 $k + 1$ 的整闭子概形
$V_i \subset X$ 和有理函数 $f_i \in R(V_i)^*$，使得族
$\{V_i\}_{i \in I}$ 在 $X$ 中局部有限，并且
$\alpha = \sum (V_i \to X)_*\text{div}(f_i)$。令
$$
W_i \subset V_i \times_S \mathbf{P}^1_S \subset X \times_S \mathbf{P}^1_S
$$
为引理 \ref{chow-lemma-rational-function} 中有理映射 $f_i$ 的图的闭包。
由同一引理，$(V_i \to X)_*\text{div}(f_i)$ 等于
$[(W_i)_0]_k - [(W_i)_\infty]_k$，结论显然。
\end{proof}

\begin{lemma}
\label{chow-lemma-closed-subscheme-cross-p1}
设 $(S, \delta)$ 如情形 \ref{chow-situation-setup}。
设概形 $X$ 在 $S$ 上局部有限型，并设 $Z$ 是
$X \times \mathbf{P}^1$ 的闭子概形。假设
\begin{enumerate}
\item $\dim_\delta(Z) \leq k + 1$；
\item $\dim_\delta(Z_0) \leq k$、$\dim_\delta(Z_\infty) \leq k$；
\item 对 $Z$ 的任意嵌入点 $\xi$（《除子》定义 \ref{divisors-definition-embedded}），要么
$\xi \not \in Z_0 \cup Z_\infty$，要么 $\delta(\xi) < k$。
\end{enumerate}
那么，作为 $X$ 上的 $k$-循环，有
$[Z_0]_k \sim_{rat} [Z_\infty]_k$。
\end{lemma}

\begin{proof}
设 $\{W_i\}_{i \in I}$ 是 $Z$ 中 $\delta$-维数为 $k + 1$ 的不可约分支的集合。
按定义写成
$$
[Z]_{k + 1} = \sum n_i[W_i]
$$
其中 $n_i > 0$。由《除子》引理 \ref{divisors-lemma-components-locally-finite}，$\{W_i\}$ 是
$X \times_S \mathbf{P}^1_S$ 的闭子集的局部有限族。我们断言
$$
[Z_0]_k = \sum n_i[(W_i)_0]_k
$$
并且对 $[Z_\infty]_k$ 也有类似等式。一旦证明此点，引理即由引理 \ref{chow-lemma-rational-equivalence-family} 得出。

\medskip\noindent
设 $Z' \subset X$ 是 $\delta$-维数为 $k$ 的整闭子概形。
为证明上述等式，只须证明 $[Z']$ 在 $[Z_0]_k$ 中的系数 $n$，
等于 $[Z']$ 在 $\sum n_i[(W_i)_0]_k$ 中的系数 $m$。
设 $\xi' \in Z'$ 为一般点，并置
$\xi = (\xi', 0) \in  X \times_S \mathbf{P}^1_S$。
考虑局部环 $A = \mathcal{O}_{X \times_S \mathbf{P}^1_S, \xi}$。
设 $I \subset A$ 是截出 $Z$ 的理想，即 $A/I = \mathcal{O}_{Z, \xi}$。
设 $t \in A$ 是截出 $X \times_S D_0$ 的元素（即把 $\mathbf{P}^1$
在零点处的坐标拉回）。由 $\xi' \in Z'$ 的选择，
$\delta(\xi) = k$，故 $\dim(A/I) = 1$。由假设 (3)，$\xi$ 不是嵌入点，
所以 $A/I$ 是 Cohen--Macaulay 环。由于 $\dim_\delta(Z_0)
= k$，
有 $\dim(A/(t, I)) = 0$，这蕴含 $t$ 是 $A/I$ 上的非零因子。
最后，经过 $\xi$ 的不可约闭子概形 $W_i$ 对应于位于 $I$ 上方的极小素理想
$I \subset \mathfrak q_i$。重数 $n_i$ 对应于长度
$\text{length}_{A_{\mathfrak q_i}}(A/I)_{\mathfrak q_i}$。因此
$$
n = \text{length}_A(A/(t, I))
$$
并且
$$
m = \sum
\text{length}_A(A/(t, \mathfrak q_i))
\text{length}_{A_{\mathfrak q_i}}(A/I)_{\mathfrak q_i}
$$
结论于是由引理 \ref{chow-lemma-additivity-divisors-restricted} 得出。
\end{proof}

\begin{lemma}
\label{chow-lemma-coherent-sheaf-cross-p1}
设 $(S, \delta)$ 如情形 \ref{chow-situation-setup}。
设概形 $X$ 在 $S$ 上局部有限型。设 $\mathcal{F}$ 是
$X \times \mathbf{P}^1$ 上的凝聚层。设
$i_0, i_\infty : X \to X \times \mathbf{P}^1$ 是满足
$i_t(x) = (x, t)$ 的闭浸入。记
$\mathcal{F}_0 = i_0^*\mathcal{F}$、$\mathcal{F}_\infty = i_\infty^*\mathcal{F}$。
假设
\begin{enumerate}
\item $\dim_\delta(\text{Supp}(\mathcal{F})) \leq k + 1$；
\item $\dim_\delta(\text{Supp}(\mathcal{F}_0)) \leq k$、
$\dim_\delta(\text{Supp}(\mathcal{F}_\infty)) \leq k$；
\item 对 $\mathcal{F}$ 的任意嵌入伴随点 $\xi$，要么
$\xi \not \in (X \times \mathbf{P}^1)_0 \cup (X \times \mathbf{P}^1)_\infty$，
要么 $\delta(\xi) < k$。
\end{enumerate}
那么，作为 $X$ 上的 $k$-循环，有
$[\mathcal{F}_0]_k \sim_{rat} [\mathcal{F}_\infty]_k$。
\end{lemma}

\begin{proof}
设 $\{W_i\}_{i \in I}$ 是 $\text{Supp}(\mathcal{F})$ 中
$\delta$-维数为 $k + 1$ 的不可约分支的集合。按定义写成
$$
[\mathcal{F}]_{k + 1} = \sum n_i[W_i]
$$
其中 $n_i > 0$。由引理 \ref{chow-lemma-length-finite}，$\{W_i\}$ 是
$X \times_S \mathbf{P}^1_S$ 的闭子集的局部有限族。我们断言
$$
[\mathcal{F}_0]_k = \sum n_i[(W_i)_0]_k
$$
并且对 $[\mathcal{F}_\infty]_k$ 也有类似等式。一旦证明此点，引理即由引理 \ref{chow-lemma-rational-equivalence-family} 得出。

\medskip\noindent
设 $Z' \subset X$ 是 $\delta$-维数为 $k$ 的整闭子概形。
为证明上述等式，只须证明 $[Z']$ 在 $[\mathcal{F}_0]_k$ 中的系数 $n$，
等于 $[Z']$ 在 $\sum n_i[(W_i)_0]_k$ 中的系数 $m$。
设 $\xi' \in Z'$ 为一般点，并置
$\xi = (\xi', 0) \in  X \times_S \mathbf{P}^1_S$。
考虑局部环 $A = \mathcal{O}_{X \times_S \mathbf{P}^1_S, \xi}$，
并把 $M = \mathcal{F}_\xi$ 看成 $A$-模。设 $t \in A$ 是截出
$X \times_S D_0$ 的元素（即把 $\mathbf{P}^1$ 在零点处的坐标拉回）。
由 $\xi' \in Z'$ 的选择，$\delta(\xi) = k$，故
$\dim(\text{Supp}(M)) = 1$。由假设 (3)，$\xi$ 不是 $\mathcal{F}$ 的伴随点，
所以 $M$ 是 Cohen--Macaulay 模。由于
$\dim_\delta(\text{Supp}(\mathcal{F}_0)) = k$，有
$\dim(\text{Supp}(M/tM)) = 0$，这蕴含 $t$ 是 $M$ 上的非零因子。
最后，经过 $\xi$ 的不可约闭子概形 $W_i$ 对应于
$\text{Ass}(M)$ 的极小素理想 $\mathfrak q_i$。重数 $n_i$ 对应于长度
$\text{length}_{A_{\mathfrak q_i}}M_{\mathfrak q_i}$。因此
$$
n = \text{length}_A(M/tM)
$$
并且
$$
m = \sum
\text{length}_A(A/(t, \mathfrak q_i)A)
\text{length}_{A_{\mathfrak q_i}}M_{\mathfrak q_i}
$$
结论于是由引理 \ref{chow-lemma-additivity-divisors-restricted} 得出。
\end{proof}







\section{Chow 群与包络}
\label{chow-section-envelopes}

\noindent
定义如下。

\begin{definition}
\label{chow-definition-envelope}
\begin{reference}
\cite[Definition 18.3]{F}
\end{reference}
设 $X$ 为概形。所谓{\it 包络}，是指完全分解的固有态射 $f : Y \to X$
（《态射进阶》定义 \ref{more-morphisms-definition-cd-morphism}）。
\end{definition}

\noindent
引理 \ref{chow-lemma-envelope} 的正合序列是采用该定义的主要动机。

\begin{lemma}
\label{chow-lemma-composition-envelope}
设 $(S, \delta)$ 如情形 \ref{chow-situation-setup}。
设概形 $X$ 在 $S$ 上局部有限型。若 $f : Y \to X$ 和
$g : Z \to Y$ 都是包络，则 $f \circ g$ 也是包络。
\end{lemma}

\begin{proof}
这由《概形态射》引理 \ref{morphisms-lemma-composition-proper} 和
《态射进阶》引理 \ref{more-morphisms-lemma-composition-cd} 得出。
\end{proof}

\begin{lemma}
\label{chow-lemma-base-change-envelope}
设 $(S, \delta)$ 如情形 \ref{chow-situation-setup}。
设 $X' \to X$ 是在 $S$ 上局部有限型的概形之间的态射。
若 $f : Y \to X$ 是包络，则 $f$ 的基变换 $f' : Y' \to X'$ 也是包络。
\end{lemma}

\begin{proof}
这由《概形态射》引理 \ref{morphisms-lemma-base-change-proper} 和
《态射进阶》引理 \ref{more-morphisms-lemma-base-change-cd} 得出。
\end{proof}

\begin{lemma}
\label{chow-lemma-envelope}
设 $(S, \delta)$ 如情形 \ref{chow-situation-setup}。
设概形 $X$ 在 $S$ 上局部有限型，并设 $f : Y \to X$ 是包络。
那么对所有 $k \in \mathbf{Z}$，有正合序列
$$
\CH_k(Y \times_X Y) \xrightarrow{p_* - q_*}
\CH_k(Y) \xrightarrow{f_*}
\CH_k(X) \to 0
$$
其中 $p, q : Y \times_X Y \to Y$ 是投影。
\end{lemma}

\begin{proof}
由于 $f$ 是包络，$f$ 是固有态射，故循环及循环类的推前都有定义；见第 \ref{chow-section-proper-pushforward} 节和 \ref{chow-section-push-pull}。
类似地，态射 $p$ 和 $q$ 是 $f$ 的基变换，所以也固有。两个箭头的复合为零，因为
$f_* \circ p_* = (p \circ f)_* = (q \circ f)_* = f_* \circ q_*$；
见引理 \ref{chow-lemma-compose-pushforward}。

\medskip\noindent
先证明 $f_* : Z_k(Y) \to Z_k(X)$ 是满射。设
$\alpha = \sum n_i[Z_i] \in Z_k(X)$，其中 $Z_i \subset X$ 是整闭子概形的
局部有限族。令 $x_i \in Z_i$ 为一般点。由于 $f$ 是包络，因而完全分解，
存在点 $y_i \in Y$，满足 $f(y_i) = x_i$，且
$\kappa(y_i)/\kappa(x_i)$ 是平凡扩张。令 $W_i \subset Y$ 为以 $y_i$ 为一般点的
整闭子概形。由于 $f$ 是闭映射，$f(W_i) = Z_i$。由此，闭子概形族 $W_i$
在 $Y$ 上局部有限。由于 $\kappa(y_i)/\kappa(x_i)$ 平凡，有
$\dim_\delta(W_i) = \dim_\delta(Z_i) = k$。因此
$\beta = \sum n_i[W_i]$ 属于 $Z_k(Y)$。最后，因为
$\kappa(y_i)/\kappa(x_i)$ 平凡，支配态射
$f|_{W_i} : W_i \to Z_i$ 的次数为 $1$，故 $f_*\beta = \alpha$。

\medskip\noindent
既然 $f_* : Z_k(Y) \to Z_k(X)$ 是满射，映射
$f_* : \CH_k(Y) \to \CH_k(X)$ 自然更是满射。

\medskip\noindent
设 $\beta \in Z_k(Y)$，且 $f_*\beta$ 在 $\CH_k(X)$ 中为零。
这意味着可以找到整闭子概形的局部有限族 $Z_j \subset X$，满足
$\dim_\delta(Z_j) = k + 1$，以及 $f_j \in R(Z_j)^*$，使得作为循环有
$$
f_*\beta = \sum (Z_j \to X)_*\text{div}(f_j)
$$
其中 $i_j : Z_j \to X$ 是给定的闭浸入。完全照上述方式，可以找到整闭子概形的
局部有限族 $W_j \subset Y$，使得 $f(W_j) = Z_j$ 且
$W_j \to Z_j$ 是双有理态射，即诱导同构 $R(Z_j) = R(W_j)$。
以 $g_j \in R(W_j)^*$ 表示对应于 $f_j$ 的元素。注意 $W_j \to Z_j$ 固有，
且在 $Z_j$ 上作为循环有
$(W_j \to Z_j)_*\text{div}(g_j) = \text{div}(f_j)$。
由此，若把 $\beta$ 替换为与其有理等价的循环
$$
\beta' = \beta - \sum (W_j \to Y)_*\text{div}(g_j)
$$
则 $f_*\beta' = 0$。（这里使用引理 \ref{chow-lemma-compose-pushforward}。）
因此，要完成引理的证明，只须证明下一段中的断言。

\medskip\noindent
断言：若 $\beta \in Z_k(Y)$ 且 $f_*\beta = 0$，则对某个
$\gamma \in Z_k(Y \times_X Y)$，在 $Z_k(Y)$ 中有
$\beta = \delta + p_*\gamma - q_*\gamma$。具体地，写成
$\beta = \sum_{j \in J} n_j[W_j]$，其中 $\{W_j\}_{j \in J}$ 是 $Y$ 中
满足 $\dim_\delta(W_j) = k$ 的整闭子概形的局部有限族。
固定一个整闭子概形 $Z \subset X$，并考虑子集
$J_Z = \{j \in J : f(W_j) = Z\}$。这是有限集。分三种情形：
\begin{enumerate}
\item $J_Z = \emptyset$。此时置 $\gamma_Z = 0$。
\item $J_Z \not = \emptyset$ 且 $\dim_\delta(Z) = k$。
考察 $Z$ 的系数，条件 $f_*\beta = 0$ 蕴含
$\sum_{j \in J_Z} n_j\deg(W_j/Z) = 0$。此时选择双有理地映到 $Z$ 的整闭子概形
$W \subset Y$（见上文）。考察一般点可知，$W_j \times_Z W$ 有唯一一个
不可约分支 $W'_j \subset W_j \times_Z W \subset Y \times_X Y$，
它双有理地映到 $W_j$。于是 $W'_j \to W$ 支配，且
$\deg(W'_j/W) = \deg(W_j/W)$。因此，若置
$\gamma_Z = \sum_{j \in J_Z} n_j[W'_j]$，则
$p_*\gamma_Z = \sum_{j \in J_Z} n_j[W_j]$，并且
$q_*\gamma_Z = \sum_{j \in J_Z} n_j\deg(W'_j/W)[W] = 0$。
\item $J_Z \not = \emptyset$ 且 $\dim_\delta(Z) < k$。
此时选择双有理地映到 $Z$ 的整闭子概形 $W \subset Y$（见上文）。
考察一般点可知，$W_j \times_Z W$ 有唯一一个不可约分支
$W'_j \subset W_j \times_Z W \subset Y \times_X Y$，它双有理地映到 $W_j$。
于是 $W'_j \to W$ 支配，且
$k = \dim_\delta(W'_j) > \dim_\delta(W) = \dim_\delta(Z)$。
因此，若置 $\gamma_Z = \sum_{j \in J_Z} n_j[W'_j]$，则
$p_*\gamma_Z = \sum_{j \in J_Z} n_j[W_j]$，并且 $q_*\gamma_Z = 0$。
\end{enumerate}
由于整闭子概形族 $\{f(W_j)\}$ 在 $X$ 上局部有限
（引理 \ref{chow-lemma-quasi-compact-locally-finite}），$Y \times_X Y$ 上的
$k$-循环
$$
\gamma = \sum\nolimits_{Z \subset X\text{ 为闭整子概形}} \gamma_Z
$$
有良好定义。由上述计算，$p_*\gamma_Z = \beta$ 且 $q_*\gamma_Z = 0$，
这就蕴含所需结论。
\end{proof}













\section{Chow 群与 K-群}
\label{chow-section-chow-and-K}

\noindent
本节比较凝聚层范畴的 $K_0$ 与 Chow 群。

\medskip\noindent
设 $(S, \delta)$ 如情形 \ref{chow-situation-setup}。
设概形 $X$ 在 $S$ 上局部有限型。以
$\textit{Coh}(X) = \textit{Coh}(\mathcal{O}_X)$ 表示 $X$ 上凝聚层的范畴。
这是 Abel 范畴；见《概形上同调》引理 \ref{coherent-lemma-coherent-abelian-Noetherian}。对任意
$k \in \mathbf{Z}$，以 $\textit{Coh}_{\leq k}(X)$ 表示
$\textit{Coh}(X)$ 的全子范畴，其对象是满足
$\dim_\delta(\text{Supp}(\mathcal{F})) \leq k$ 的凝聚层 $\mathcal{F}$。

\begin{lemma}
\label{chow-lemma-Serre-subcategories}
设 $(S, \delta)$ 如情形 \ref{chow-situation-setup}。
设概形 $X$ 在 $S$ 上局部有限型。范畴
$\textit{Coh}_{\leq k}(X)$ 是 Abel 范畴 $\textit{Coh}(X)$ 的 Serre 子范畴。
\end{lemma}

\begin{proof}
Serre 子范畴的定义见《同调代数》定义 \ref{homology-definition-serre-subcategory}。该引理的证明直接，故略去。
\end{proof}

\begin{lemma}
\label{chow-lemma-cycles-k-group}
设 $(S, \delta)$ 如情形 \ref{chow-situation-setup}。
设概形 $X$ 在 $S$ 上局部有限型。映射
$$
Z_k(X)
\longrightarrow
K_0(\textit{Coh}_{\leq k}(X)/\textit{Coh}_{\leq k - 1}(X)),
\quad
\sum n_Z[Z] \mapsto
\left[\bigoplus\nolimits_{n_Z > 0} \mathcal{O}_Z^{\oplus n_Z}\right]
-
\left[\bigoplus\nolimits_{n_Z < 0} \mathcal{O}_Z^{\oplus -n_Z}\right]
$$
与映射
$$
K_0(\textit{Coh}_{\leq k}(X)/\textit{Coh}_{\leq k - 1}(X))
\longrightarrow
Z_k(X),\quad
\mathcal{F} \longmapsto [\mathcal{F}]_k
$$
是互逆同构。
\end{lemma}

\begin{proof}
注意，若 $\sum n_Z[Z]$ 属于 $Z_k(X)$，则直和
$\bigoplus\nolimits_{n_Z > 0} \mathcal{O}_Z^{\oplus n_Z}$ 与
$\bigoplus\nolimits_{n_Z < 0} \mathcal{O}_Z^{\oplus -n_Z}$
都是 $X$ 上的凝聚层，因为族 $\{Z \mid n_Z > 0\}$ 在 $X$ 上局部有限。
映射 $\mathcal{F} \to [\mathcal{F}]_k$ 在
$\textit{Coh}_{\leq k}(X)$ 上可加；见引理 \ref{chow-lemma-additivity-sheaf-cycle}。而若
$\mathcal{F} \in \textit{Coh}_{\leq k - 1}(X)$，则
$[\mathcal{F}]_k = 0$。由《同调代数》引理 \ref{homology-lemma-serre-subcategory-K-groups} 的 (1)，第二个映射也有良好定义。
显然，先应用第一个映射再应用第二个映射所得复合是恒等映射。

\medskip\noindent
反过来，设从 $X$ 上的凝聚层 $\mathcal{F}$ 出发。写成
$[\mathcal{F}]_k = \sum_{i \in I} n_i[Z_i]$，其中 $n_i > 0$，而
$Z_i \subset X$、$i \in I$ 是两两不同、$\delta$-维数为 $k$ 的整闭子概形。
需要证明
$$
[\mathcal{F}] = [\bigoplus\nolimits_{i \in I} \mathcal{O}_{Z_i}^{\oplus n_i}]
$$
在 $K_0(\textit{Coh}_{\leq k}(X)/\textit{Coh}_{\leq k - 1}(X))$ 中成立。
以 $\xi_i \in Z_i$ 表示一般点。若置
$$
\mathcal{F}' = \Ker(\mathcal{F} \to \bigoplus \xi_{i, *}\mathcal{F}_{\xi_i})
$$
则 $\mathcal{F}'$ 是 $\mathcal{F}$ 的、支撑维数 $\leq k - 1$ 的最大凝聚子模。
特别地，$\mathcal{F}$ 与 $\mathcal{F}/\mathcal{F}'$ 在
$K_0(\textit{Coh}_{\leq k}(X)/\textit{Coh}_{\leq k - 1}(X))$ 中具有同一类。
因此，以 $\mathcal{F}/\mathcal{F}'$ 替换 $\mathcal{F}$ 后，可以并且将假设
上面展示的核 $\mathcal{F}'$ 为零。

\medskip\noindent
对每个 $i \in I$，选择滤过
$$
\mathcal{F}_{\xi_i} = \mathcal{F}_i^0 \supset \mathcal{F}_i^1 \supset
\ldots \supset \mathcal{F}_i^{n_i} = 0
$$
使各相继商在 $\xi_i$ 处的剩余域上维数为 $1$。这是可能的，因为
$\mathcal{F}_{\xi_i}$ 作为 $\mathcal{O}_{X, \xi_i}$-模的长度为 $n_i$。
对 $p > n_i$，置 $\mathcal{F}_i^p = 0$。对 $p \geq 0$，记
$$
\mathcal{F}^p =
\Ker\left(\mathcal{F} \longrightarrow \bigoplus
\xi_{i, *}(\mathcal{F}_{\xi_i}/\mathcal{F}_i^p)\right)
$$
那么 $\mathcal{F}^p$ 凝聚，$\mathcal{F}^0 = \mathcal{F}$；并且在
$\xi_i$ 的某个开邻域内，$\mathcal{F}^p/\mathcal{F}^{p + 1}$ 同构于秩为
$1$ 的自由 $\mathcal{O}_{Z_i}$-模（若 $n_i > p$），或同构于 $0$
（若 $n_i \leq p$）。此外，$\mathcal{F}' = \bigcap \mathcal{F}^p = 0$。
由于每个拟紧开集 $U \subset X$ 只包含有限多个 $\xi_i$，可知当
$p \gg 0$ 时 $\mathcal{F}^p|_U$ 为零。因此
$\bigoplus_{p \geq 0} \mathcal{F}^p$ 是凝聚 $\mathcal{O}_X$-模。
考虑短正合列
$$
0 \to
\bigoplus\nolimits_{p > 0} \mathcal{F}^p \to
\bigoplus\nolimits_{p \geq 0} \mathcal{F}^p \to
\bigoplus\nolimits_{p > 0} \mathcal{F}^p/\mathcal{F}^{p + 1} \to 0
$$
和
$$
0 \to
\bigoplus\nolimits_{p > 0} \mathcal{F}^p \to
\bigoplus\nolimits_{p \geq 0} \mathcal{F}^p \to
\mathcal{F} \to 0
$$
它们都是凝聚 $\mathcal{O}_X$-模的短正合列。这已经表明
$$
[\mathcal{F}] = [\bigoplus \mathcal{F}^p/\mathcal{F}^{p + 1}]
$$
在 $K_0(\textit{Coh}_{\leq k}(X)/\textit{Coh}_{\leq k - 1}(X))$ 中成立。
接下来，对每个满足 $n_i > p$ 的 $p \geq 0$ 与 $i \in I$，
选择非零理想层 $\mathcal{I}_{i, p} \subset \mathcal{O}_{Z_i}$，并选择 $X$ 上的映射
$\mathcal{I}_{i, p} \to \mathcal{F}^p/\mathcal{F}^{p + 1}$，使其在上述
$\xi_i$ 的开邻域上为同构。《概形上同调》引理 \ref{coherent-lemma-extend-coherent} 保证可以这样选择。然后考虑短正合列
$$
0 \to
\bigoplus\nolimits_{p \geq 0, i \in I, n_i > p} \mathcal{I}_{i, p}
\to
\bigoplus \mathcal{F}^p/\mathcal{F}^{p + 1} \to
\mathcal{Q} \to 0
$$
以及短正合列
$$
0 \to
\bigoplus\nolimits_{p \geq 0, i \in I, n_i > p} \mathcal{I}_{i, p}
\to
\bigoplus\nolimits_{p \geq 0, i \in I, n_i > p} \mathcal{O}_{Z_i}
\to
\mathcal{Q}' \to 0
$$
注意 $\mathcal{Q}$ 和 $\mathcal{Q}'$ 在各点 $\xi_i$ 的某个邻域内都为零，
且都支撑在 $\bigcup Z_i$ 上。因此 $\mathcal{Q}$ 和 $\mathcal{Q}'$ 属于
$\textit{Coh}_{\leq k - 1}(X)$。由于
$$
\bigoplus\nolimits_{i \in I} \mathcal{O}_{Z_i}^{\oplus n_i} \cong
\bigoplus\nolimits_{p \geq 0, i \in I, n_i > p} \mathcal{O}_{Z_i}
$$
证明完成。
\end{proof}

\begin{lemma}
\label{chow-lemma-finite-cycles-k-group}
设 $\pi : X \to Y$ 是在情形 \ref{chow-situation-setup} 中
$(S, \delta)$ 上局部有限型的概形之间的有限态射。那么
$\pi_* : \textit{Coh}(X) \to \textit{Coh}(Y)$ 是正合函子；它把
$\textit{Coh}_{\leq k}(X)$ 送入 $\textit{Coh}_{\leq k}(Y)$，并在这些范畴及其商范畴的
$K_0$ 上诱导同态。引理 \ref{chow-lemma-cycles-k-group} 的映射构成交换图
$$
\xymatrix{
Z_k(X) \ar[d]^{\pi_*} \ar[r] &
K_0(\textit{Coh}_{\leq k}(X)/\textit{Coh}_{\leq k - 1}(X))
\ar[d]^{\pi_*} \ar[r] &
Z_k(X) \ar[d]^{\pi_*} \\
Z_k(Y) \ar[r] &
K_0(\textit{Coh}_{\leq k}(Y)/\textit{Coh}_{\leq k - 1}(Y)) \ar[r] &
Z_k(Y)
}
$$
\end{lemma}

\begin{proof}
有限态射是仿射态射，所以由《概形上同调》引理 \ref{coherent-lemma-relative-affine-vanishing}，沿 $\pi$ 推前拟凝聚模是正合函子。
有限态射是固有态射，所以 $\pi_*$ 把凝聚层送到凝聚层；见《概形上同调》命题 \ref{coherent-proposition-proper-pushforward-coherent}。
关于支撑维数的断言显然。右侧的交换性立即由引理 \ref{chow-lemma-cycle-push-sheaf} 得到。由于各水平箭头是双射，左侧也交换。
\end{proof}

\begin{lemma}
\label{chow-lemma-from-chow-to-K}
设 $X$ 是情形 \ref{chow-situation-setup} 中在 $(S, \delta)$ 上局部有限型的概形。
存在典范映射
$$
\CH_k(X)
\longrightarrow
K_0(\textit{Coh}_{\leq k + 1}(X)/\textit{Coh}_{\leq k - 1}(X))
$$
它由引理 \ref{chow-lemma-cycles-k-group} 中的映射
$Z_k(X) \to K_0(\textit{Coh}_{\leq k}(X)/\textit{Coh}_{\leq k - 1}(X))$
诱导。
\end{lemma}

\begin{proof}
需要证明：$Z_k(X)$ 中有理等价于零的元素 $\alpha$，在
$K_0(\textit{Coh}_{\leq k + 1}(X)/\textit{Coh}_{\leq k - 1}(X))$
中映到零。按照定义 \ref{chow-definition-rational-equivalence} 写成
$\alpha = \sum (i_j)_*\text{div}(f_j)$。注意
$$
\pi = \coprod i_j : W = \coprod W_j \longrightarrow X
$$
是有限态射，因为每个 $i_j : W_j \to X$ 都是闭浸入，且 $W_j$ 的族在
$X$ 中局部有限。因此可以使用引理 \ref{chow-lemma-finite-cycles-k-group}
把问题约化到 $W$ 的情形。由于 $W$ 是整概形的不交并，进一步约化到下一段讨论的情形。

\medskip\noindent
假设 $X$ 是 $\delta$-维数为 $k + 1$ 的整概形。设 $f$ 是 $X$ 上的非零有理函数，
并令 $\alpha = \text{div}(f)$。需要证明 $\alpha$ 在
$K_0(\textit{Coh}_{\leq k + 1}(X)/\textit{Coh}_{\leq k - 1}(X))$
中映到零。设 $\mathcal{I} \subset \mathcal{O}_X$ 是 $f$ 的分母理想；见《除子》定义 \ref{divisors-definition-regular-meromorphic-ideal-denominators}。
于是有短正合列
$$
0 \to \mathcal{I} \to \mathcal{O}_X \to \mathcal{O}_X/\mathcal{I} \to 0
$$
和
$$
0 \to \mathcal{I} \xrightarrow{f} \mathcal{O}_X \to
\mathcal{O}_X/f\mathcal{I} \to 0
$$
见《除子》引理 \ref{divisors-lemma-regular-meromorphic-ideal-denominators}。我们断言
$$
[\mathcal{O}_X/\mathcal{I}]_k - [\mathcal{O}_X/f\mathcal{I}]_k =
\text{div}(f)
$$
该断言表明，元素 $\alpha = \text{div}(f)$ 在
$K_0(\textit{Coh}_{\leq k}(X)/\textit{Coh}_{\leq k - 1}(X))$ 中由
$[\mathcal{O}_X/\mathcal{I}] - [\mathcal{O}_X/f\mathcal{I}]$ 表示。
上述短正合列于是表明，该元素在
$K_0(\textit{Coh}_{\leq k + 1}(X)/\textit{Coh}_{\leq k - 1}(X))$ 中映到零。

\medskip\noindent
为证明该断言，设 $Z \subset X$ 是 $\delta$-维数为 $k$ 的整闭子概形，
$\xi \in Z$ 是其一般点。那么
$I = \mathcal{I}_\xi \subset A = \mathcal{O}_{X, \xi}$ 是满足
$fI \subset A$ 的理想。$[Z]$ 在 $\text{div}(f)$ 中的系数是
$\text{ord}_A(f)$。（当然，与往常一样，把 $X$ 的函数域与 $A$ 的分式域等同。）
另一方面，$[Z]$ 在
$[\mathcal{O}_X/\mathcal{I}] - [\mathcal{O}_X/f\mathcal{I}]$ 中的系数为
$$
\text{length}_A(A/I) - \text{length}_A(A/fI)
$$
利用《交换代数》定义 \ref{algebra-definition-distance} 的距离函数，
可将其改写为
$$
d(A, I) - d(A, fI) = d(I, fI) = \text{ord}_A(f)
$$
这些等式由《交换代数》引理 \ref{algebra-lemma-properties-distance-function} 和
\ref{algebra-lemma-order-vanishing-determinant} 得出。
使用这些引理并非必要，但很方便。
\end{proof}

\begin{remark}
\label{chow-remark-good-cases-K-A}
设 $(S, \delta)$ 如情形 \ref{chow-situation-setup}。
设概形 $X$ 在 $S$ 上局部有限型。稍后将在引理 \ref{chow-lemma-cycles-rational-equivalence-K-group} 中看到，引理 \ref{chow-lemma-from-chow-to-K} 的映射
$$
\CH_k(X)
\longrightarrow
K_0(\textit{Coh}_{k + 1}(X)/\textit{Coh}_{\leq k - 1}(X))
$$
是单射。将它与典范映射
$$
K_0(\textit{Coh}_{k + 1}(X)/\textit{Coh}_{\leq k - 1}(X))
\longrightarrow
K_0(\textit{Coh}(X)/\textit{Coh}_{\leq k - 1}(X))
$$
复合，得到典范映射
$$
\CH_k(X)
\longrightarrow
K_0(\textit{Coh}(X)/\textit{Coh}_{\leq k - 1}(X)).
$$
关于这个映射是否应为单射，我们未能在文献中找到相关断言或猜想。
合理的预期似乎是其核为挠群。稍后将回到这个问题（插入未来引用）。
\end{remark}

\begin{lemma}
\label{chow-lemma-K-coherent-supported-on-closed}
设 $X$ 是局部诺特概形，$Z \subset X$ 是闭子概形。以
$\textit{Coh}_Z(X) \subset \textit{Coh}(X)$ 表示这样的凝聚
$\mathcal{O}_X$-模构成的 Serre 子范畴：其集合论支撑包含在 $Z$ 中。
那么正合包含函子 $\textit{Coh}(Z) \to \textit{Coh}_Z(X)$ 诱导同构
$$
K'_0(Z) = K_0(\textit{Coh}(Z)) \longrightarrow K_0(\textit{Coh}_Z(X))
$$
\end{lemma}

\begin{proof}
设 $\mathcal{F}$ 是 $\textit{Coh}_Z(X)$ 的对象。设
$\mathcal{I} \subset \mathcal{O}_X$ 是 $Z$ 的拟凝聚理想层。考虑降滤过
$$
\ldots \subset
\mathcal{F}^p = \mathcal{I}^p \mathcal{F} \subset
\mathcal{F}^{p - 1} \subset \ldots \subset \mathcal{F}^0 = \mathcal{F}
$$
完全如引理 \ref{chow-lemma-from-chow-to-K} 的证明，这个滤过局部有限，因而
$\bigoplus_{p \geq 0} \mathcal{F}^p$、
$\bigoplus_{p \geq 1} \mathcal{F}^p$ 与
$\bigoplus_{p \geq 0} \mathcal{F}^p/\mathcal{F}^{p + 1}$
都是支撑在 $Z$ 上的凝聚 $\mathcal{O}_X$-模。因此，与引理 \ref{chow-lemma-from-chow-to-K} 的证明完全一样，在 $K_0(\textit{Coh}_Z(X))$ 中有
$$
[\mathcal{F}] =
[\bigoplus\nolimits_{p \geq 0} \mathcal{F}^p/\mathcal{F}^{p + 1}]
$$
凝聚模 $\bigoplus_{p \geq 0} \mathcal{F}^p/\mathcal{F}^{p + 1}$ 被
$\mathcal{I}$ 零化，所以 $[\mathcal{F}]$ 属于该像。事实上，我们断言映射
$$
\mathcal{F} \longmapsto 
c(\mathcal{F}) =
[\bigoplus\nolimits_{p \geq 0} \mathcal{F}^p/\mathcal{F}^{p + 1}]
$$
可通过 $K_0(\textit{Coh}_Z(X))$ 分解，并且是引理所述映射的逆。
为此只须证明：若
$$
0 \to \mathcal{F} \to \mathcal{G} \to \mathcal{H} \to 0
$$
是 $\textit{Coh}_Z(X)$ 中的短正合列，则
$c(\mathcal{G}) = c(\mathcal{F}) + c(\mathcal{H})$。
注意，对所有 $q \geq 0$，有短正合列
$$
0 \to
(\mathcal{F} \cap \mathcal{I}^q\mathcal{G})/
(\mathcal{F} \cap \mathcal{I}^{q + 1}\mathcal{G}) \to
\mathcal{G}^q/\mathcal{G}^{q + 1} \to
\mathcal{H}^q/\mathcal{H}^{q + 1} \to 0
$$
对 $p, q \geq 0$，考虑凝聚子模
$$
\mathcal{F}^{p, q} = \mathcal{I}^p\mathcal{F} \cap \mathcal{I}^q\mathcal{G}
$$
完全照上述方式，并利用滤过 $\mathcal{F}^p = \mathcal{I}^p\mathcal{F}$ 与
$\mathcal{F} \cap \mathcal{I}^q\mathcal{G}$ 都局部有限，得到
$$
[\bigoplus\nolimits_{p \geq 0} \mathcal{F}^p/\mathcal{F}^{p + 1}] =
[\bigoplus\nolimits_{p, q \geq 0}
\mathcal{F}^{p, q}/(\mathcal{F}^{p + 1, q} + \mathcal{F}^{p, q + 1})] =
[\bigoplus\nolimits_{q \geq 0}
(\mathcal{F} \cap \mathcal{I}^q\mathcal{G})/
(\mathcal{F} \cap \mathcal{I}^{q + 1}\mathcal{G})]
$$
其位于 $K_0(\textit{Coh}(Z))$ 中。结合上述正合列即可得到所需结论。
若干细节从略。
\end{proof}













\section{与可逆层相伴的除子}
\label{chow-section-divisor-invertible-sheaf}

\noindent
下面的定义是《除子》定义 \ref{divisors-definition-divisor-invertible-sheaf} 在当前设定下的对应形式。

\begin{definition}
\label{chow-definition-divisor-invertible-sheaf}
设 $(S, \delta)$ 如情形 \ref{chow-situation-setup}。
设 $X$ 在 $S$ 上局部有限型，并假设 $X$ 是整概形且
$n = \dim_\delta(X)$。设 $\mathcal{L}$ 是可逆 $\mathcal{O}_X$-模。
\begin{enumerate}
\item 对 $\mathcal{L}$ 的任意非零亚纯截面 $s$，定义{\it 与 $s$ 相伴的 Weil 除子}为
$(n - 1)$-循环
$$
\text{div}_\mathcal{L}(s) =
\sum \text{ord}_{Z, \mathcal{L}}(s) [Z]
$$
其定义见《除子》定义 \ref{divisors-definition-divisor-invertible-sheaf}。由引理 \ref{chow-lemma-divisor-delta-dimension}，Weil 除子的 $\delta$-维数为 $n - 1$，
故此定义有意义。
\item 定义{\it 与 $\mathcal{L}$ 相伴的 Weil 除子}为
$$
c_1(\mathcal{L}) \cap [X] =
\text{div}_\mathcal{L}(s)\text{ 的类} \in \CH_{n - 1}(X)
$$
其中 $s$ 是 $X$ 上 $\mathcal{L}$ 的任意非零亚纯截面。由《除子》引理 \ref{divisors-lemma-divisor-meromorphic-well-defined}，此定义与选择无关。
\end{enumerate}
\end{definition}

\noindent
设 $X$、$S$ 如上面的定义 \ref{chow-definition-divisor-invertible-sheaf}，
并置 $n = \dim_\delta(X)$。由定义显然有
$Cl(X) = \CH_{n - 1}(X)$，其中 $Cl(X)$ 是《除子》定义 \ref{divisors-definition-class-group} 中定义的 $X$ 的 Weil 除子类群。映射
$$
\Pic(X) \longrightarrow \CH_{n - 1}(X), \quad
\mathcal{L} \longmapsto c_1(\mathcal{L}) \cap [X]
$$
与 Divisors, Equation (\ref{divisors-equation-c1}) 对任意局部诺特整概形构造的映射
$\Pic(X) \to Cl(X)$ 相同。特别地，该映射是 Abel 群同态；若 $X$ 正规，则它是单射；
若 $X$ 的所有局部环都是 UFD，则它是同构。见《除子》引理 \ref{divisors-lemma-normal-c1-injective} 和
\ref{divisors-lemma-local-rings-UFD-c1-bijective}。
在某些情形下，与可逆层相伴的 Weil 除子很容易计算。

\begin{lemma}
\label{chow-lemma-compute-c1}
设 $(S, \delta)$ 如情形 \ref{chow-situation-setup}。
设 $X$ 在 $S$ 上局部有限型，并假设 $X$ 是整概形且
$n = \dim_\delta(X)$。设 $\mathcal{L}$ 是可逆 $\mathcal{O}_X$-模，
$s \in \Gamma(X, \mathcal{L})$ 是非零整体截面。那么在 $Z_{n - 1}(X)$ 中
$$
\text{div}_\mathcal{L}(s) = [Z(s)]_{n - 1}
$$
并且在 $\CH_{n - 1}(X)$ 中
$$
c_1(\mathcal{L}) \cap [X] = [Z(s)]_{n - 1}
$$
\end{lemma}

\begin{proof}
设 $Z \subset X$ 是 $\delta$-维数为 $n - 1$ 的整闭子概形，
$\xi \in Z$ 是其一般点。选择生成元 $s_\xi \in \mathcal{L}_\xi$，并对某个
$f \in \mathcal{O}_{X, \xi}$ 写成 $s = fs_\xi$。由 $Z(s)$ 的定义
（见《除子》定义 \ref{divisors-definition-zero-scheme-s}），
$Z(s)$ 由拟凝聚理想层 $\mathcal{I} \subset \mathcal{O}_X$ 截出，且
$\mathcal{I}_\xi = (f)$。因此
$\text{length}_{\mathcal{O}_{X, \xi}}(\mathcal{O}_{Z(s), \xi})
=
\text{length}_{\mathcal{O}_{X, \xi}}(\mathcal{O}_{X, \xi}/(f))
=
\text{ord}_{\mathcal{O}_{X, \xi}}(f)$，正合所需。
\end{proof}

\noindent
下面的引理将被更一般的引理 \ref{chow-lemma-flat-pullback-cap-c1} 所涵盖。

\begin{lemma}
\label{chow-lemma-flat-pullback-divisor-invertible-sheaf}
设 $(S, \delta)$ 如情形 \ref{chow-situation-setup}。
设 $X$、$Y$ 在 $S$ 上局部有限型，并假设 $X$、$Y$ 都是整概形且
$n = \dim_\delta(Y)$。设 $\mathcal{L}$ 是可逆 $\mathcal{O}_Y$-模，
$f : X \to Y$ 是相对维数为 $r$ 的平坦态射。那么在
$\CH_{n + r - 1}(X)$ 中
$$
f^*(c_1(\mathcal{L}) \cap [Y]) = c_1(f^*\mathcal{L}) \cap [X]
$$
\end{lemma}

\begin{proof}
设 $s$ 是 $\mathcal{L}$ 的非零亚纯截面。我们将证明实际上有
$f^*\text{div}_\mathcal{L}(s) = \text{div}_{f^*\mathcal{L}}(f^*s)$，
从而引理成立。为此，设 $\xi \in Y$ 是一点，且
$s_\xi \in \mathcal{L}_\xi$ 是生成元。对某个 $g \in R(Y)^*$ 写成
$s = gs_\xi$。存在 $\xi$ 的开邻域 $V \subset Y$，使得
$s_\xi \in \mathcal{L}(V)$，且 $s_\xi$ 生成 $\mathcal{L}|_V$。因此
$$
\text{div}_\mathcal{L}(s)|_V = \text{div}_Y(g)|_V.
$$
完全类似地，由于 $f^*s_\xi$ 在 $f^{-1}(V)$ 上生成 $f^*\mathcal{L}$，
且 $f^*s = g f^*s_\xi$，还有
$$
\text{div}_\mathcal{L}(f^*s)|_{f^{-1}(V)}
=
\text{div}_X(g)|_{f^{-1}(V)}.
$$
于是，$f^{-1}(V)$ 上所需的循环等式由主除子拉回的相应结论得到；见引理 \ref{chow-lemma-flat-pullback-principal-divisor}。
\end{proof}



\section{与可逆层相交}
\label{chow-section-intersecting-with-divisors}

\noindent
本节研究如下构造。

\begin{definition}
\label{chow-definition-cap-c1}
设 $(S, \delta)$ 如情形 \ref{chow-situation-setup}。
设 $X$ 在 $S$ 上局部有限型，并设 $\mathcal{L}$ 是可逆
$\mathcal{O}_X$-模。对每个整数 $k$，定义运算
$$
c_1(\mathcal{L}) \cap - :
Z_{k + 1}(X) \to \CH_k(X)
$$
称为{\it 与 $\mathcal{L}$ 的第一陈类相交}。
\begin{enumerate}
\item 给定满足 $\dim_\delta(W) = k + 1$ 的整闭子概形
$i : W \to X$，定义
$$
c_1(\mathcal{L}) \cap [W] = i_*(c_1({i^*\mathcal{L}}) \cap [W])
$$
其中右端由定义 \ref{chow-definition-divisor-invertible-sheaf} 定义。
\item 对一般的 $(k + 1)$-循环 $\alpha = \sum n_i [W_i]$，置
$$
c_1(\mathcal{L}) \cap \alpha = \sum n_i c_1(\mathcal{L}) \cap [W_i]
$$
\end{enumerate}
\end{definition}

\noindent
将每个 $c_1(\mathcal{L}) \cap W_i = \sum_j n_{i, j} [Z_{i, j}]$ 写成上述形式，
其中 $\{Z_{i, j}\}_j$ 是 $W_i$ 的整闭子概形的局部有限和。
由于 $\{W_i\}$ 是 $X$ 上整闭子概形的局部有限族，容易看出
$\{Z_{i, j}\}_{i, j}$ 是 $X$ 的闭子概形的局部有限族。因此
$c_1(\mathcal{L}) \cap \alpha = \sum n_in_{i, j}[Z_{i, j}]$ 是循环。
另一种更方便的理解方式是注意到态射 $\coprod W_i \to X$ 固有。
因此可将 $c_1(\mathcal{L}) \cap \alpha$ 看作
$\CH_k(\coprod W_i) = \prod \CH_k(W_i)$ 中某个类的推前。
这也说明所得结果在 $X$ 上模有理等价有良好定义。

\medskip\noindent
接下来几节的主要目标，是证明与 $c_1(\mathcal{L})$ 相交可下降到有理等价类。
这并非平凡事实。

\begin{lemma}
\label{chow-lemma-c1-cap-additive}
设 $(S, \delta)$ 如情形 \ref{chow-situation-setup}。
设 $X$ 在 $S$ 上局部有限型。设 $\mathcal{L}$、$\mathcal{N}$ 是 $X$ 上的可逆层。
那么对每个 $\alpha \in Z_{k + 1}(X)$，在 $\CH_k(X)$ 中有
$$
c_1(\mathcal{L}) \cap \alpha  + c_1(\mathcal{N}) \cap \alpha =
c_1(\mathcal{L} \otimes_{\mathcal{O}_X} \mathcal{N}) \cap \alpha
$$
此外，对所有 $\alpha$，有 $c_1(\mathcal{O}_X) \cap \alpha = 0$。
\end{lemma}

\begin{proof}
可加性直接由《除子》引理 \ref{divisors-lemma-c1-additive} 与定义得到。
为说明 $c_1(\mathcal{O}_X) \cap \alpha = 0$，考虑截面
$1 \in \Gamma(X, \mathcal{O}_X)$。它在任意整闭子概形
$W \subset X$ 上限制为处处非零的截面。因此
$c_1(\mathcal{O}_X) \cap [W] = 0$，正合所需。
\end{proof}

\noindent
回忆，$Z(s) \subset X$ 表示概形 $X$ 上某个可逆层的整体截面 $s$ 的零点概形；
见《除子》定义 \ref{divisors-definition-zero-scheme-s}。

\begin{lemma}
\label{chow-lemma-prepare-geometric-cap}
设 $(S, \delta)$ 如情形 \ref{chow-situation-setup}。
设 $Y$ 在 $S$ 上局部有限型。设 $\mathcal{L}$ 是可逆
$\mathcal{O}_Y$-模，并设 $s \in \Gamma(Y, \mathcal{L})$。假设
\begin{enumerate}
\item $\dim_\delta(Y) \leq k + 1$；
\item $\dim_\delta(Z(s)) \leq k$；
\item 对 $Z(s)$ 中 $\delta$-维数为 $k$ 的任意不可约分支的一般点 $\xi$，
乘以 $s$ 诱导单射 $\mathcal{O}_{Y, \xi} \to \mathcal{L}_\xi$。
\end{enumerate}
写成 $[Y]_{k + 1} = \sum n_i[Y_i]$，其中 $Y_i \subset Y$ 是 $Y$ 中
$\delta$-维数为 $k + 1$ 的不可约分支。置
$s_i = s|_{Y_i} \in \Gamma(Y_i, \mathcal{L}|_{Y_i})$。那么，作为 $Y$ 上的
$k$-循环，有
\begin{equation}
\label{chow-equation-equal-as-cycles}
[Z(s)]_k =  \sum n_i[Z(s_i)]_k
\end{equation}
\end{lemma}

\begin{proof}
设 $Z \subset Y$ 是 $\delta$-维数为 $k$ 的整闭子概形，
$\xi \in Z$ 是其一般点。要比较 $[Z]$ 在表达式
$\sum n_i[Z(s_i)]_k$ 中的系数 $n$ 与 $[Z]$ 在表达式 $[Z(s)]_k$ 中的系数 $m$。
选择生成元 $s_\xi \in \mathcal{L}_\xi$。写成
$A = \mathcal{O}_{Y, \xi}$、$L = \mathcal{L}_\xi$。那么
$L = As_\xi$。对某个唯一的 $f \in A$ 写成 $s = f s_\xi$。
假设 (3) 意味着 $f : A \to A$ 是单射。由于
$\dim_\delta(Y) \leq k + 1$ 且 $\dim_\delta(Z) = k$，有
$\dim(A) = 0$ 或 $1$。并且
$$
m = \text{length}_A(A/(f))
$$
在两种情形下都是有限的。

\medskip\noindent
若 $\dim(A) = 0$，则 $f : A \to A$ 为单射蕴含 $f \in A^*$。
所以此时 $m$ 为零。此外，条件 $\dim(A) = 0$ 意味着 $\xi$ 不位于任何
$\delta$-维数为 $k + 1$ 的不可约分支上，即 $n = 0$。

\medskip\noindent
现在设 $\dim(A) = 1$。由于 $A$ 是诺特局部环，它只有有限多个极小素理想
$\mathfrak q_1, \ldots, \mathfrak q_t$。它们与经过 $\xi'$ 的 $Y_i$ 一一对应。
此外，$n_i = \text{length}_{A_{\mathfrak q_i}}(A_{\mathfrak q_i})$。
$[Z]$ 在 $[Z(s_i)]_k$ 中的重数为
$\text{length}_A(A/(f, \mathfrak q_i))$。因此，此时需要证明的等式是
$$
\text{length}_A(A/(f))
=
\sum \text{length}_{A_{\mathfrak q_i}}(A_{\mathfrak q_i})
\text{length}_A(A/(f, \mathfrak q_i))
$$
它由引理 \ref{chow-lemma-additivity-divisors-restricted} 得出。
\end{proof}

\noindent
下面的引理可用于计算可逆层的 $c_1$ 与闭子概形所对应循环的交积。
回忆，$Z(s) \subset X$ 表示概形 $X$ 上某个可逆层的整体截面 $s$ 的零点概形；
见《除子》定义 \ref{divisors-definition-zero-scheme-s}。

\begin{lemma}
\label{chow-lemma-geometric-cap}
设 $(S, \delta)$ 如情形 \ref{chow-situation-setup}。
设 $X$ 在 $S$ 上局部有限型。设 $\mathcal{L}$ 是可逆
$\mathcal{O}_X$-模，$Y \subset X$ 是闭子概形，且
$s \in \Gamma(Y, \mathcal{L}|_Y)$。假设
\begin{enumerate}
\item $\dim_\delta(Y) \leq k + 1$；
\item $\dim_\delta(Z(s)) \leq k$；
\item 对 $Z(s)$ 中 $\delta$-维数为 $k$ 的任意不可约分支的一般点 $\xi$，
乘以 $s$ 诱导单射
$\mathcal{O}_{Y, \xi} \to (\mathcal{L}|_Y)_\xi$\footnote{例如，当
$s$ 是 $\mathcal{L}|_Y$ 的正则截面时，这一条件成立。}。
\end{enumerate}
那么在 $\CH_k(X)$ 中有
$$
c_1(\mathcal{L}) \cap [Y]_{k + 1} = [Z(s)]_k
$$
\end{lemma}

\begin{proof}
写成
$$
[Y]_{k + 1} = \sum n_i[Y_i]
$$
其中 $Y_i \subset Y$ 是 $Y$ 中 $\delta$-维数为 $k + 1$ 的不可约分支，
且 $n_i > 0$。由假设，限制
$s_i = s|_{Y_i} \in \Gamma(Y_i, \mathcal{L}|_{Y_i})$ 非零，因而是正则截面。
由引理 \ref{chow-lemma-compute-c1}，$[Z(s_i)]_k$ 表示
$c_1(\mathcal{L}|_{Y_i})$。所以按定义
$$
c_1(\mathcal{L}) \cap [Y]_{k + 1} = \sum n_i[Z(s_i)]_k
$$
结论于是由引理 \ref{chow-lemma-prepare-geometric-cap} 得出。
\end{proof}




\section{与可逆层相交及推前、拉回}
\label{chow-section-intersecting-with-divisors-push-pull}

\noindent
本节证明运算 $c_1(\mathcal{L}) \cap -$ 与平坦拉回、固有推前交换。

\begin{lemma}
\label{chow-lemma-prepare-flat-pullback-cap-c1}
设 $(S, \delta)$ 如情形 \ref{chow-situation-setup}。
设 $X$、$Y$ 在 $S$ 上局部有限型。设 $f : X \to Y$ 是相对维数为
$r$ 的平坦态射，$\mathcal{L}$ 是 $Y$ 上的可逆层。假设 $Y$ 是整概形且
$n = \dim_\delta(Y)$。设 $s$ 是 $\mathcal{L}$ 的非零亚纯截面。那么
$$
f^*\text{div}_\mathcal{L}(s) = \sum n_i\text{div}_{f^*\mathcal{L}|_{X_i}}(s_i)
$$
在 $Z_{n + r - 1}(X)$ 中成立。这里求和遍历其中
$\delta$-维数为 $n + r$ 的不可约分支 $X_i \subset X$；截面
$s_i = f|_{X_i}^*(s)$ 是 $s$ 的拉回；而
$n_i = m_{X_i, X}$ 是 $X_i$ 在 $X$ 中的重数。
\end{lemma}

\begin{proof}
为证明这个循环等式，可以在 $Y$ 上局部工作。因此可以假设 $Y$ 仿射，且对某些非零截面
$p \in \Gamma(Y, \mathcal{L})$ 和 $q \in \Gamma(Y, \mathcal{O})$，有
$s = p/q$。若能同时证明
$$
f^*\text{div}_\mathcal{L}(p) =
\sum n_i\text{div}_{f^*\mathcal{L}|_{X_i}}(p_i)
\quad\text{且}\quad
f^*\text{div}_\mathcal{O}(q) =
\sum n_i\text{div}_{\mathcal{O}_{X_i}}(q_i)
$$
（记号含义显然），则由可加性即得结论；见《除子》引理 \ref{divisors-lemma-c1-additive}。所以可以假设
$s \in \Gamma(Y, \mathcal{L})$。此时应用等式
(\ref{chow-equation-equal-as-cycles})，得到
$$
[Z(f^*(s))]_{k + r - 1} =
\sum n_i\text{div}_{f^*\mathcal{L}|_{X_i}}(s_i)
$$
其中 $f^*(s) \in f^*\mathcal{L}$ 表示 $s$ 到 $X$ 的拉回。另一方面，有
$$
f^*\text{div}_\mathcal{L}(s) = f^*[Z(s)]_{k - 1}
= [f^{-1}(Z(s))]_{k + r - 1},
$$
这由引理 \ref{chow-lemma-compute-c1} 和 \ref{chow-lemma-pullback-coherent} 得出。
由于 $Z(f^*(s)) = f^{-1}(Z(s))$，结论成立。
\end{proof}

\begin{lemma}
\label{chow-lemma-flat-pullback-cap-c1}
设 $(S, \delta)$ 如情形 \ref{chow-situation-setup}。
设 $X$、$Y$ 在 $S$ 上局部有限型。设 $f : X \to Y$ 是相对维数为
$r$ 的平坦态射，$\mathcal{L}$ 是 $Y$ 上的可逆层，$\alpha$ 是 $Y$ 上的
$k$-循环。那么在 $\CH_{k + r - 1}(X)$ 中
$$
f^*(c_1(\mathcal{L}) \cap \alpha) = c_1(f^*\mathcal{L}) \cap f^*\alpha
$$
\end{lemma}

\begin{proof}
写成 $\alpha = \sum n_i[W_i]$。我们将通过在 $X$ 的闭子概形
$f^{-1}(W_i)$ 上构造有理等价，证明在 $\CH_{k + r - 1}(X)$ 中
$$
f^*(c_1(\mathcal{L}) \cap [W_i]) = c_1(f^*\mathcal{L}) \cap f^*[W_i]
$$
注 \ref{chow-remark-infinite-sums-rational-equivalences} 中的讨论于是表明，
引理的等式成立。

\medskip\noindent
设 $W \subset Y$ 是 $\delta$-维数为 $k$ 的整闭子概形。
考虑闭子概形 $W' = f^{-1}(W) = W \times_Y X$，从而有纤维积图
$$
\xymatrix{
W' \ar[r] \ar[d]_h & X \ar[d]^f \\
W \ar[r] & Y
}
$$
需要证明
$f^*(c_1(\mathcal{L}) \cap [W]) = c_1(f^*\mathcal{L}) \cap f^*[W]$。
选择 $\mathcal{L}|_W$ 的非零亚纯截面 $s$。设
$W'_i \subset W'$ 是 $\delta$-维数为 $k + r$ 的不可约分支。
按定义写成 $[W']_{k + r} = \sum n_i[W'_i]$，其中 $n_i$ 是
$W'_i$ 在 $W'$ 中的重数。因此在 $Z_{k + r}(X)$ 中有
$f^*[W] = \sum n_i[W'_i]$。由于每个 $W'_i \to W$ 都支配，
对每个 $i$，$s_i = s|_{W'_i}$ 都是非零亚纯截面。由引理 \ref{chow-lemma-prepare-flat-pullback-cap-c1}，有循环等式
$$
h^*\text{div}_{\mathcal{L}|_W}(s) =
\sum n_i\text{div}_{f^*\mathcal{L}|_{W'_i}}(s_i)
$$
其位于 $Z_{k + r - 1}(W')$ 中。证明由此完成：左端是 $W'$ 上的循环，
其在 $\CH_{k + r - 1}(X)$ 中推前为
$f^*(c_1(\mathcal{L}) \cap [W])$；右端也是 $W'$ 上的循环，其在
$\CH_{k + r - 1}(X)$ 中推前为
$c_1(f^*\mathcal{L}) \cap f^*[W]$。
\end{proof}

\begin{lemma}
\label{chow-lemma-equal-c1-as-cycles}
设 $(S, \delta)$ 如情形 \ref{chow-situation-setup}。
设 $X$、$Y$ 在 $S$ 上局部有限型。设 $f : X \to Y$ 是固有态射，
$\mathcal{L}$ 是 $Y$ 上的可逆层。设 $s$ 是 $Y$ 上 $\mathcal{L}$ 的非零亚纯截面，
并将该截面记为 $s$。假设 $X$、$Y$ 都是整概形，$f$ 支配，且
$\dim_\delta(X) = \dim_\delta(Y)$。那么，作为 $Y$ 上的循环，
$$
f_*\left(\text{div}_{f^*\mathcal{L}}(f^*s)\right) =
[R(X) : R(Y)]\text{div}_\mathcal{L}(s).
$$
特别地，
$$
f_*(c_1(f^*\mathcal{L}) \cap [X]) =
[R(X) : R(Y)] c_1(\mathcal{L}) \cap [Y] =
c_1(\mathcal{L}) \cap f_*[X]
$$
\end{lemma}

\begin{proof}
由定义，$f_*[X] = [R(X) : R(Y)][Y]$，所以最后一个等式由第一个等式得出。
可以再次使用引理 \ref{chow-lemma-proper-pushforward-alteration} 来证明第一个等式。
因为所证的是循环等式，可以在 $Y$ 上局部工作。因此可以假设
$\mathcal{L} = \mathcal{O}_Y$。此时 $s$ 对应于有理函数 $g \in R(Y)$，
而所需证明的只是
$$
f_*\left(\text{div}_X(g)\right) =
[R(X) : R(Y)]\text{div}_Y(g).
$$
与上述引理 \ref{chow-lemma-proper-pushforward-alteration} 的结果比较，
由 $g \in R(Y)^*$ 时
$\text{Nm}_{R(X)/R(Y)}(g) = g^{[R(X) : R(Y)]}$，即得该等式。
\end{proof}

\begin{lemma}
\label{chow-lemma-pushforward-cap-c1}
设 $(S, \delta)$ 如情形 \ref{chow-situation-setup}。
设 $X$、$Y$ 在 $S$ 上局部有限型。设 $p : X \to Y$ 是固有态射，
$\alpha \in Z_{k + 1}(X)$，而 $\mathcal{L}$ 是 $Y$ 上的可逆层。
那么在 $\CH_k(Y)$ 中有
$$
p_*(c_1(p^*\mathcal{L}) \cap \alpha) = c_1(\mathcal{L}) \cap p_*\alpha
$$
\end{lemma}

\begin{proof}
先假设 $p$ 具有如下性质：对每个整闭子概形 $W \subset X$，映射
$p|_W : W \to Y$ 都是闭浸入。那么由与 $c_1(\mathcal{L})$ 作帽积的定义，
引理成立。

\medskip\noindent
利用这一观察约化到特殊情形。写成 $\alpha = \sum n_i[W_i]$，其中
$n_i \not = 0$，且各 $W_i$ 两两不同。设 $W'_i \subset Y$ 是 $W_i$ 的像
（视为整闭子概形）。考虑图表
$$
\xymatrix{
X' = \coprod W_i \ar[r]_-q \ar[d]_{p'} & X \ar[d]^p \\
Y' = \coprod W'_i \ar[r]^-{q'} & Y.
}
$$
由于 $\{W_i\}$ 在 $X$ 上局部有限，且 $p$ 固有，可知
$\{W'_i\}$ 在 $Y$ 上局部有限，并且 $q, q', p'$ 也都是固有态射。
还可以把 $\sum n_i[W_i]$ 看成 $k$-循环
$\alpha' \in Z_k(X')$。显然 $q_*\alpha' = \alpha$。由证明开头的观察，有
$q_*(c_1(q^*p^*\mathcal{L}) \cap \alpha')
= c_1(p^*\mathcal{L}) \cap q_*\alpha'$
以及
$(q')_*(c_1((q')^*\mathcal{L}) \cap p'_*\alpha') =
c_1(\mathcal{L}) \cap q'_*p'_*\alpha'$。因此只须对态射 $p'$ 和循环
$\sum n_i[W_i]$ 证明引理。显然，这意味着可以假设 $X$、$Y$ 都是整概形，
$f : X \to Y$ 支配，且 $\alpha = [X]$。此时结论由引理 \ref{chow-lemma-equal-c1-as-cycles} 得出。
\end{proof}






\section{关键公式}
\label{chow-section-key}

\noindent
设 $(S, \delta)$ 如情形 \ref{chow-situation-setup}。
设 $X$ 在 $S$ 上局部有限型，并假设
$X$ 是整概形且 $\dim_\delta(X) = n$。
设 $\mathcal{L}$ 和 $\mathcal{N}$ 是 $X$ 上的可逆层。
设 $s$ 是 $\mathcal{L}$ 的非零亚纯截面，而
$t$ 是 $\mathcal{N}$ 的非零亚纯截面。
设 $Z_i \subset X$（$i \in I$）是一族局部有限的余维 $1$
不可约闭子集，并具有如下性质：
若 $Z \not \in \{Z_i\}$ 且 $\xi$ 是其一般点，则 $s$ 是
$\mathcal{L}_\xi$ 的生成元，而 $t$ 是 $\mathcal{N}_\xi$ 的生成元。
由《除子》引理 \ref{divisors-lemma-divisor-meromorphic-locally-finite}，这样的集合存在。
于是
$$
\text{div}_\mathcal{L}(s) = \sum \text{ord}_{Z_i, \mathcal{L}}(s) [Z_i]
$$
并且类似地，
$$
\text{div}_\mathcal{N}(t) = \sum \text{ord}_{Z_i, \mathcal{N}}(t) [Z_i]
$$
进一步展开定义：对每个 $i$ 选取生成元
$s_i \in \mathcal{L}_{\xi_i}$ 和 $t_i \in \mathcal{N}_{\xi_i}$，
其中 $\xi_i$ 是 $Z_i$ 的一般点。于是可以写成
$$
s = f_i s_i
\quad\text{且}\quad
t = g_i t_i
$$
置 $B_i = \mathcal{O}_{X, \xi_i}$。由定义，
$$
\text{ord}_{Z_i, \mathcal{L}}(s) = \text{ord}_{B_i}(f_i)
\quad\text{且}\quad
\text{ord}_{Z_i, \mathcal{N}}(t) = \text{ord}_{B_i}(g_i)
$$
因为 $t_i$ 是 $\mathcal{N}_{\xi_i}$ 的生成元，它在纤维
$\mathcal{N}_{\xi_i} \otimes \kappa(\xi_i)$ 中的像是
$\mathcal{N}|_{Z_i}$ 的非零亚纯截面。将此像记为
$t_i|_{Z_i}$。由定义可得
$$
c_1(\mathcal{N}) \cap \text{div}_\mathcal{L}(s) =
\sum \text{ord}_{B_i}(f_i)
(Z_i \to X)_*\text{div}_{\mathcal{N}|_{Z_i}}(t_i|_{Z_i})
$$
以及类似的等式
$$
c_1(\mathcal{L}) \cap \text{div}_\mathcal{N}(t) =
\sum \text{ord}_{B_i}(g_i)
(Z_i \to X)_*\text{div}_{\mathcal{L}|_{Z_i}}(s_i|_{Z_i})
$$
它们位于 $\CH_{n - 2}(X)$ 中。下面要找出这两个循环之间的有理等价。
为此，考虑驯顺符号
$$
\partial_{B_i}(f_i, g_i) \in \kappa(\xi_i)^*
$$
见第 \ref{chow-section-tame-symbol} 节。

\begin{lemma}[关键公式]
\label{chow-lemma-key-formula}
在上述情形中，循环
$$
\sum
(Z_i \to X)_*\left(
\text{ord}_{B_i}(f_i) \text{div}_{\mathcal{N}|_{Z_i}}(t_i|_{Z_i}) -
\text{ord}_{B_i}(g_i) \text{div}_{\mathcal{L}|_{Z_i}}(s_i|_{Z_i}) \right)
$$
等于循环
$$
\sum (Z_i \to X)_*\text{div}(\partial_{B_i}(f_i, g_i))
$$
\end{lemma}

\begin{proof}
先考察用 $us_i$ 替换 $s_i$ 时会发生什么，其中 $u$ 是 $B_i$ 中的单位。
此时 $f_i$ 被 $u^{-1} f_i$ 替换。因此，引理第一个表达式的第一部分不变，
而在第二部分中增加
$$
-\text{ord}_{B_i}(g_i)\text{div}(u|_{Z_i})
$$
（其中 $u|_{Z_i}$ 是 $u$ 在剩余域中的像）；这是
《除子》引理 \ref{divisors-lemma-divisor-meromorphic-well-defined}
的结论。第二个表达式中则增加
$$
\text{div}(\partial_{B_i}(u^{-1}, g_i))
$$
这是驯顺符号双线性的结果。由驯顺符号的性质
(\ref{chow-item-normalization})，这两项相等。

\medskip\noindent
设 $Z \subset X$ 是满足 $\dim_\delta(Z) = n - 2$ 的不可约闭子集。
为证明引理中两个循环的 $Z$ 系数相同，可以像上一段那样作替换
$s_i \mapsto us_i$。完全类似地，也可以用 $t_i \mapsto vt_i$
作替换，其中 $v$ 是 $B_i$ 的单位。

\medskip\noindent
由于所证的是循环相等，可以逐个系数论证。因此，选取满足
$\dim_\delta(Z) = n - 2$ 的不可约闭子集 $Z \subset X$，并比较其系数。
设 $\xi \in Z$ 为一般点，并置 $A = \mathcal{O}_{X, \xi}$。
这是一个维数为 $2$ 的 Noether 局部整环。选取
$\mathcal{L}_\xi$ 和 $\mathcal{N}_\xi$ 的生成元 $\sigma$ 和 $\tau$。
缩小 $X$ 后，可以并且确实假设 $\sigma$ 和 $\tau$ 在整个 $X$ 上
给出可逆层 $\mathcal{L}$ 和 $\mathcal{N}$ 的平凡化。
由于缩小 $X$ 后 $Z_i$ 局部有限，可以假设对所有 $i \in I$ 都有
$Z \subset Z_i$，并且 $I$ 有限。此时 $\xi_i$ 对应于高度为 $1$ 的素理想
$\mathfrak q_i \subset A$。可以写成 $s_i = a_i \sigma$ 和
$t_i = b_i \tau$，其中 $a_i$ 和 $b_i$ 是 $A_{\mathfrak q_i}$ 中的单位。
由上述讨论，只须在所有 $i$ 都满足 $a_i = b_i = 1$ 时证明引理。

\medskip\noindent
假设所有 $i$ 都有 $a_i = b_i = 1$。由于所选 $\sigma$ 和 $\tau$
是平凡化截面，引理的第一个表达式为零。写
$s = f\sigma$ 和 $t = g \tau$，其中 $f$、$g$ 属于 $A$ 的分式域。
由上一段，问题已约化到对所有 $i$ 都有 $f_i = f$ 和 $g_i = g$ 的情形。
此外，若 $A$ 的高度 $1$ 素理想 $\mathfrak q$ 不在
$\{\mathfrak q_i\}$ 中，则 $f$、$g$ 都是 $A_\mathfrak q$ 中的单位
（这是由引理之前讨论中对族 $\{Z_i\}$ 的选取所得）。因此，引理第二个
表达式中 $Z$ 的系数为
$$
\sum\nolimits_i \text{ord}_{A/\mathfrak q_i}(\partial_{B_i}(f, g))
$$
而关键引理引理 \ref{chow-lemma-milnor-gersten-low-degree} 说明它为零。
\end{proof}

\begin{remark}
\label{chow-remark-higher-chow-pointwise}
设 $(S, \delta)$ 如情形 \ref{chow-situation-setup}。
设 $X$ 在 $S$ 上局部有限型，并设 $k \in \mathbf{Z}$。
我们断言存在复形
$$
\bigoplus\nolimits_{\delta(x) = k + 2}' K_2^M(\kappa(x))
\xrightarrow{\partial}
\bigoplus\nolimits_{\delta(x) = k + 1}' K_1^M(\kappa(x))
\xrightarrow{\partial}
\bigoplus\nolimits_{\delta(x) = k}' K_0^M(\kappa(x))
$$
这里采用注 \ref{chow-remark-chow-group-pointwise} 中引入的记号和约定，
并另作如下约定：
\begin{enumerate}
\item $K_2^M(\kappa(x))$ 是点 $x \in X$ 的剩余域 $\kappa(x)$ 的
Milnor K-理论的次数 $2$ 部分（见注 \ref{chow-remark-gersten-complex-milnor}）；它是
$\kappa(x)^* \otimes_\mathbf{Z} \kappa(x)^*$ 关于由形如
$\lambda \otimes (1 - \lambda)$ 的元素生成的子群所取的商，
其中 $\lambda \in \kappa(x) \setminus \{0, 1\}$；以及
\item 第一个微分 $\partial$ 定义如下：给定第一项中的元素
$\xi = \sum_x \alpha_x$，置
$$
\partial(\xi) = \sum\nolimits_{x \leadsto x',\ \delta(x') = k + 1}
\partial_{\mathcal{O}_{W_x, x'}}(\alpha_x)
$$
其中
$\partial_{\mathcal{O}_{W_x, x'}} : K_2^M(\kappa(x)) \to K_1^M(\kappa(x))$
是第 \ref{chow-section-tame-symbol} 节中构造的驯顺符号。
\end{enumerate}
我们断言所得确为复形，即 $\partial \circ \partial = 0$。
为此，只须取上述元素 $\xi$ 以及满足 $\delta(x'') = k$ 的点
$x'' \in X$，并验证元素 $\partial(\partial(\xi))$ 中 $x''$ 的系数为零。
由于 $\xi = \sum \alpha_x$ 是局部有限和，利用可加性，实际上可以假设
$\xi = \alpha_x$，其中某个 $x \in X$ 满足 $\delta(x) = k + 2$，且
$\alpha_x \in K_2^M(\kappa(x))$。再由线性性，可以假设
$\alpha_x = f \otimes g$，其中 $f, g \in \kappa(x)^*$。
以 $W \subset X$ 表示一般点为 $x$ 的整闭子概形。若
$x'' \not \in W$，则立即可见，$\partial(\partial(\xi))$ 中 $x$ 的系数为零。
若 $x'' \in W$，则 $\partial(\partial(x))$ 中 $x''$ 的系数等于
$$
\sum\nolimits_{x \leadsto x' \leadsto x'',\ \delta(x') = k + 1}
\text{ord}_{\mathcal{O}_{\overline{\{x'\}}, x''}}(
\partial_{\mathcal{O}_{W, x'}}(f, g))
$$
关键代数引理引理 \ref{chow-lemma-milnor-gersten-low-degree}
恰好说明这个和为零。
\end{remark}

\begin{remark}
\label{chow-remark-higher-chow}
设 $(S, \delta)$ 如情形 \ref{chow-situation-setup}。
设 $X$ 在 $S$ 上局部有限型，并设 $k \in \mathbf{Z}$。
注 \ref{chow-remark-higher-chow-pointwise} 中的复形以及注 \ref{chow-remark-chow-group-pointwise} 中 $\CH_k(X)$ 的表示提示我们，
可以定义第一高阶 Chow 群
$$
\CH^M_k(X, 1) =
H_1(\text{the complex of Remark \ref{chow-remark-higher-chow-pointwise}})
$$
我们使用上标 ${}^M$，以区别于文献中定义的高阶 Chow 群，
例如 Spencer Bloch 的论文（\cite{Bloch} 和 \cite{Bloch-moving}）中的定义。
设 $U \subset X$ 是开子集，其补集为 $Y \subset X$（视为既约闭子概形）。
于是对 $i = 2, 1, 0$ 有分裂短正合列
$$
0 \to
\bigoplus\nolimits_{y \in Y, \delta(y) = k + i}' K_i^M(\kappa(y)) \to
\bigoplus\nolimits_{x \in X, \delta(x) = k + i}' K_i^M(\kappa(x)) \to
\bigoplus\nolimits_{u \in U, \delta(u) = k + i}' K_i^M(\kappa(u)) \to 0
$$
它与注 \ref{chow-remark-higher-chow-pointwise} 中各复形的边界映射相容。
应用蛇形引理（见《同调代数》引理 \ref{homology-lemma-long-exact-sequence-chain}），得到六项正合列
$$
\CH^M_k(Y, 1) \to \CH^M_k(X, 1) \to \CH^M_k(U, 1) \to
\CH_k(Y) \to \CH_k(X) \to \CH_k(U) \to 0
$$
它延拓了引理 \ref{chow-lemma-restrict-to-open} 的典范正合列。
稍作进一步工作，还可以为第一高阶 Chow 群
$\CH^M_k(X, 1)$ 定义平坦拉回与固有推前。我们将在后文回到这一点
（此处插入未来引用）。
\end{remark}






\section{与可逆层相交及有理等价}
\label{chow-section-commutativity}

\noindent
应用关键引理，可以得到与可逆层相交的基本性质。特别地，我们将看到
$c_1(\mathcal{L}) \cap -$ 下降到有理等价类，并且不同可逆层所给出的
这类运算彼此交换。

\begin{lemma}
\label{chow-lemma-commutativity-on-integral}
设 $(S, \delta)$ 如情形 \ref{chow-situation-setup}。
设 $X$ 在 $S$ 上局部有限型，并假设 $X$ 是整概形且
$\dim_\delta(X) = n$。设 $\mathcal{L}$、$\mathcal{N}$ 是 $X$ 上的可逆层。
选取 $\mathcal{L}$ 的非零亚纯截面 $s$ 和 $\mathcal{N}$ 的非零亚纯截面
$t$。置 $\alpha = \text{div}_\mathcal{L}(s)$ 及
$\beta = \text{div}_\mathcal{N}(t)$。则
$$
c_1(\mathcal{N}) \cap \alpha
=
c_1(\mathcal{L}) \cap \beta
$$
在 $\CH_{n - 2}(X)$ 中成立。
\end{lemma}

\begin{proof}
这立即由关键引理引理 \ref{chow-lemma-key-formula} 及其之前的讨论得出。
\end{proof}

\begin{lemma}
\label{chow-lemma-factors}
设 $(S, \delta)$ 如情形 \ref{chow-situation-setup}。
设 $X$ 在 $S$ 上局部有限型，$\mathcal{L}$ 是 $X$ 上的可逆层。
运算 $\alpha \mapsto c_1(\mathcal{L}) \cap \alpha$ 下降到有理等价类，
从而给出运算
$$
c_1(\mathcal{L}) \cap - : \CH_{k + 1}(X) \to \CH_k(X)
$$
\end{lemma}

\begin{proof}
设 $\alpha \in Z_{k + 1}(X)$ 且 $\alpha \sim_{rat} 0$。
需要证明定义 \ref{chow-definition-cap-c1} 所定义的
$c_1(\mathcal{L}) \cap \alpha$ 为零。
由定义 \ref{chow-definition-rational-equivalence}，存在局部有限族
$\{W_j\}$，其成员都是满足 $\dim_\delta(W_j) = k + 2$
的整闭子概形；还存在有理函数 $f_j \in R(W_j)^*$，使得
$$
\alpha = \sum (i_j)_*\text{div}_{W_j}(f_j)
$$
注意，$p : \coprod W_j \to X$ 是固有态射。因此
$\alpha = p_*\alpha'$，其中 $\alpha' \in Z_{k + 1}(\coprod W_j)$
是主除子 $\text{div}_{W_j}(f_j)$ 之和。
由引理 \ref{chow-lemma-pushforward-cap-c1}，
$c_1(\mathcal{L}) \cap \alpha = p_*(c_1(p^*\mathcal{L}) \cap \alpha')$。
故只须证明每个
$c_1(\mathcal{L}|_{W_j}) \cap \text{div}_{W_j}(f_j)$ 都为零。
换言之，可以假设 $X$ 是整概形，并且对某个 $f \in R(X)^*$ 有
$\alpha = \text{div}_X(f)$。

\medskip\noindent
假设 $X$ 是整概形，且对某个 $f \in R(X)^*$ 有
$\alpha = \text{div}_X(f)$。可以把 $f$ 看成可逆层
$\mathcal{N} = \mathcal{O}_X$ 的正则亚纯截面。选取 $\mathcal{L}$ 的
一个亚纯截面 $s$，并记 $\beta = \text{div}_\mathcal{L}(s)$。
由引理 \ref{chow-lemma-commutativity-on-integral} 可得
$$
c_1(\mathcal{L}) \cap \alpha = c_1(\mathcal{O}_X) \cap \beta.
$$
然而引理 \ref{chow-lemma-c1-cap-additive} 说明右边在 $\CH_k(X)$ 中为零，
正合所需。
\end{proof}

\noindent
设 $(S, \delta)$ 如情形 \ref{chow-situation-setup}。
设 $X$ 在 $S$ 上局部有限型，$\mathcal{L}$ 是 $X$ 上的可逆层。
将运算 $c_1(\mathcal{L}) \cap - $ 记为
$$
c_1(\mathcal{L}) \cap - : \CH_{k + 1}(X) \to \CH_k(X)
$$
由引理 \ref{chow-lemma-factors}，这个记法有意义。对所有 $s \geq 0$，
以 $c_1(\mathcal{L})^s \cap -$ 表示该运算的 $s$ 次迭代。

\begin{lemma}
\label{chow-lemma-cap-commutative}
设 $(S, \delta)$ 如情形 \ref{chow-situation-setup}。
设 $X$ 在 $S$ 上局部有限型，$\mathcal{L}$、$\mathcal{N}$ 是 $X$ 上的
可逆层。对任意 $\alpha \in \CH_{k + 2}(X)$，都有
$$
c_1(\mathcal{L}) \cap c_1(\mathcal{N}) \cap \alpha
=
c_1(\mathcal{N}) \cap c_1(\mathcal{L}) \cap \alpha
$$
作为 $\CH_k(X)$ 中的元素。
\end{lemma}

\begin{proof}
将其写成 $\alpha = \sum m_j[Z_j]$，其中
$Z_j \subset X$ 是满足 $\dim_\delta(Z_j) = k + 2$ 的一族局部有限的
整闭子概形。考虑固有态射 $p : \coprod Z_j \to X$。
把 $\alpha' = \sum m_j[Z_j]$ 视为 $\coprod Z_j$ 上的
$(k + 2)$-循环。反复应用引理 \ref{chow-lemma-pushforward-cap-c1}，得到
$c_1(\mathcal{L}) \cap c_1(\mathcal{N}) \cap \alpha
= p_*(c_1(p^*\mathcal{L}) \cap c_1(p^*\mathcal{N}) \cap \alpha')$
以及
$c_1(\mathcal{N}) \cap c_1(\mathcal{L}) \cap \alpha
= p_*(c_1(p^*\mathcal{N}) \cap c_1(p^*\mathcal{L}) \cap \alpha')$。
因此，只须在 $X$ 为整概形且 $\alpha = [X]$ 的情形证明该公式。
此时结论由引理 \ref{chow-lemma-commutativity-on-integral} 和定义得出。
\end{proof}






\section{Gysin 同态}
\label{chow-section-intersecting-effective-Cartier}

\noindent
本节为可逆层的一个截面的零点概形 $D$ 定义 Gysin 映射。
当 $D$ 是有效 Cartier 除子时会出现一个重要情形；但把定义推广到任意
$D$，可以更灵活地表述各种相容性，见注 \ref{chow-remark-pullback-pairs} 以及引理 \ref{chow-lemma-closed-in-X-gysin}、\ref{chow-lemma-gysin-flat-pullback} 和
\ref{chow-lemma-gysin-commutes-gysin}。
这些结果还可以推广到带有虚法丛的局部主闭子概形
（注 \ref{chow-remark-generalize-to-virtual}），或推广到伪除子
（注 \ref{chow-remark-generalize-to-pseudo-divisor}）。

\medskip\noindent
回忆，有效 Cartier 除子与配对 $(\mathcal{L}, s)$ 的同构类按 $1$ 对 $1$ 对应，
其中 $\mathcal{L}$ 是可逆层，而 $s$ 是正则整体截面；见
《除子》引理 \ref{divisors-lemma-characterize-OD}。
若 $D$ 对应于 $(\mathcal{L}, s)$，则
$\mathcal{L} = \mathcal{O}_X(D)$。阅读本节时请记住这一点。

\begin{definition}
\label{chow-definition-gysin-homomorphism}
设 $(S, \delta)$ 如情形 \ref{chow-situation-setup}。
设 $X$ 在 $S$ 上局部有限型。设 $(\mathcal{L}, s)$ 是由一个可逆层
和一个整体截面 $s \in \Gamma(X, \mathcal{L})$ 组成的配对。
设 $D = Z(s)$ 是 $s$ 的零点概形，并以
$i : D \to X$ 表示闭浸入。对每个整数 $k$，按如下规则定义一个
{\it Gysin 同态}
$$
i^* : Z_{k + 1}(X) \to \CH_k(D).
$$
\begin{enumerate}
\item 给定满足 $\dim_\delta(W) = k + 1$ 的整闭子概形
$W \subset X$，定义
\begin{enumerate}
\item 若 $W \not \subset D$，则令 $i^*[W] = [D \cap W]_k$，
把它视为 $D$ 上的 $k$-循环；以及
\item 若 $W \subset D$，则令
$i^*[W] = i'_*(c_1(\mathcal{L}|_W) \cap [W])$，
其中 $i' : W \to D$ 是诱导的闭浸入。
\end{enumerate}
\item 对一般的 $(k + 1)$-循环 $\alpha = \sum n_j[W_j]$，置
$$
i^*\alpha = \sum n_j i^*[W_j]
$$
\item 若 $D$ 是有效 Cartier 除子，则以
$D \cdot \alpha = i_*i^*\alpha$ 表示类 $i^*\alpha$ 推前后所得的
$X$ 上的类。
\end{enumerate}
\end{definition}

\noindent
事实上，正如下文将看到的，这个 Gysin 同态 $i^*$ 可以视为非平坦拉回
的一个例子。因此，有时也非正式地把类 $i^*\alpha$ 称为类 $\alpha$ 的
{\it 拉回}。

\begin{remark}
\label{chow-remark-generalize-to-virtual}
设 $X$ 是情形 \ref{chow-situation-setup} 中那样在 $S$ 上局部有限型
的概形。设 $(D, \mathcal{N}, \sigma)$ 是如下三元组：局部主闭子概形
（《除子》定义 \ref{divisors-definition-effective-Cartier-divisor}）
$i : D \to X$、可逆 $\mathcal{O}_D$-模 $\mathcal{N}$，以及
$\mathcal{O}_D$-模的满射
$\sigma : \mathcal{N}^{\otimes -1} \to i^*\mathcal{I}_D$
\footnote{这个条件保证：若 $D$ 是有效 Cartier 除子，则
$\mathcal{N} = \mathcal{O}_X(D)|_D$。}。
应把 $\mathcal{N}$ 看作{\it $D$ 在 $X$ 中的虚法丛}。
定义 \ref{chow-definition-gysin-homomorphism} 中的构造
$i^* : Z_{k + 1}(X) \to \CH_k(D)$ 可以推广到这类三元组，见
第 \ref{chow-section-gysin-higher-codimension} 节。
\end{remark}

\begin{remark}
\label{chow-remark-generalize-to-pseudo-divisor}
设 $X$ 是情形 \ref{chow-situation-setup} 中那样在 $S$ 上局部有限型
的概形。在 \cite{F} 中，$X$ 上的{\it 伪除子}定义为三元组
$D = (\mathcal{L}, Z, s)$，其中 $\mathcal{L}$ 是可逆
$\mathcal{O}_X$-模，$Z \subset X$ 是闭子集，而
$s \in \Gamma(X \setminus Z, \mathcal{L})$ 是处处非零截面。
与上述情形类似，对 $\CH_{k + 1}(X)$ 中每个 $\alpha$，
可以在 $\CH_k(Z \cap |\alpha|)$ 中定义乘积 $D \cdot \alpha$，
其中 $|\alpha|$ 是 $\alpha$ 的支撑。
\end{remark}

\begin{lemma}
\label{chow-lemma-support-cap-effective-Cartier}
设 $(S, \delta)$ 如情形 \ref{chow-situation-setup}，并设 $X$ 在 $S$ 上
局部有限型。设 $(\mathcal{L}, s, i : D \to X)$ 如定义 \ref{chow-definition-gysin-homomorphism}，而 $\alpha$ 是 $X$ 上的
$(k + 1)$-循环。那么在 $\CH_k(X)$ 中，
$i_*i^*\alpha = c_1(\mathcal{L}) \cap \alpha$。
特别地，若 $D$ 是有效 Cartier 除子，则
$D \cdot \alpha = c_1(\mathcal{O}_X(D)) \cap \alpha$。
\end{lemma}

\begin{proof}
写 $\alpha = \sum n_j[W_j]$，其中 $i_j : W_j \to X$ 是满足
$\dim_\delta(W_j) = k$ 的整闭子概形。
由于 $D$ 是 $s$ 的零点概形，$D \cap W_j$ 是限制 $s|_{W_j}$ 的零点概形。
因此，对每个满足 $W_j \not \subset D$ 的 $j$，由引理 \ref{chow-lemma-geometric-cap} 有
$c_1(\mathcal{L}) \cap [W_j] = [D \cap W_j]_k$。于是，由定义 \ref{chow-definition-cap-c1}，在 $\CH_k(X)$ 中有
$$
c_1(\mathcal{L}) \cap \alpha
=
\sum\nolimits_{W_j \not \subset D} n_j[D \cap W_j]_k
+
\sum\nolimits_{W_j \subset D}
n_j i_{j, *}(c_1(\mathcal{L})|_{W_j}) \cap [W_j])
$$
右边逐项等于定义 \ref{chow-definition-gysin-homomorphism} 中
$D$ 上的类 $i^*\alpha$ 的推前。因此结论成立。
\end{proof}

\begin{lemma}
\label{chow-lemma-easy-gysin}
设 $(S, \delta)$ 如情形 \ref{chow-situation-setup}。
设 $X$ 在 $S$ 上局部有限型。设
$(\mathcal{L}, s, i : D \to X)$ 如定义 \ref{chow-definition-gysin-homomorphism}。
\begin{enumerate}
\item 设 $Z \subset X$ 是满足 $\dim_\delta(Z) \leq k + 1$ 的闭子概形，
并且 $D \cap Z$ 是 $Z$ 上的有效 Cartier 除子。则
$i^*[Z]_{k + 1} = [D \cap Z]_k$。
\item 设 $\mathcal{F}$ 是 $X$ 上的凝聚层，满足
$\dim_\delta(\text{Supp}(\mathcal{F})) \leq k + 1$，并且
$s : \mathcal{F} \to \mathcal{F} \otimes_{\mathcal{O}_X} \mathcal{L}$
是单射。则
$$
i^*[\mathcal{F}]_{k + 1} = [i^*\mathcal{F}]_k
$$
在 $\CH_k(D)$ 中成立。
\end{enumerate}
\end{lemma}

\begin{proof}
设 $Z \subset X$ 如 (1)，并置 $\mathcal{F} = \mathcal{O}_Z$。
$D \cap Z$ 是有效 Cartier 除子的假设等价于
$s : \mathcal{F} \to \mathcal{F} \otimes_{\mathcal{O}_X} \mathcal{L}$
为单射的假设。此外，
$[Z]_{k + 1} = [\mathcal{F}]_{k + 1}]$，并且
$[D \cap Z]_k = [\mathcal{O}_{D \cap Z}]_k = [i^*\mathcal{F}]_k$。
见引理 \ref{chow-lemma-cycle-closed-coherent}。因此 (1) 由 (2) 得出。

\medskip\noindent
写 $[\mathcal{F}]_{k + 1} = \sum m_j[W_j]$，其中 $m_j > 0$，
而 $W_j \subset X$ 是两两不同、$\delta$-维数为 $k + 1$ 的整闭子概形。
$s : \mathcal{F} \to \mathcal{F} \otimes_{\mathcal{O}_X} \mathcal{L}$
为单射这一假设蕴含对所有 $j$ 都有 $W_j \not \subset D$。
由定义，
$$
i^*[\mathcal{F}]_{k + 1} = \sum m_j [D \cap W_j]_k.
$$
我们断言循环等式
$$
\sum [D \cap W_j]_k = [i^*\mathcal{F}]_k
$$
成立。
设 $Z \subset D$ 是 $\delta$-维数为 $k$ 的整闭子概形。
设 $\xi \in Z$ 为其一般点，$A = \mathcal{O}_{X, \xi}$，
$M = \mathcal{F}_\xi$。设 $f \in A$ 生成 $D$ 的理想，即满足
$\mathcal{O}_{D, \xi} = A/fA$。根据假设，
$\dim(\text{Supp}(M)) = 1$，映射 $f : M \to M$ 是单射，并且
$\text{length}_A(M/fM) < \infty$。此外，$\text{length}_A(M/fM)$
是 $[i^*\mathcal{F}]_k$ 中 $[Z]$ 的系数。另一方面，设
$\mathfrak q_1, \ldots, \mathfrak q_t$ 是 $M$ 的支撑中的极小素理想。
那么
$$
\sum
\text{length}_{A_{\mathfrak q_i}}(M_{\mathfrak q_i})
\text{ord}_{A/\mathfrak q_i}(f)
$$
是 $\sum [D \cap W_j]_k$ 中 $[Z]$ 的系数。
因此，该等式由引理 \ref{chow-lemma-additivity-divisors-restricted} 得出。
\end{proof}

\begin{remark}
\label{chow-remark-gysin-on-cycles}
设 $X \to S$、$\mathcal{L}$、$s$、$i : D \to X$ 如定义 \ref{chow-definition-gysin-homomorphism}，并假设
$\mathcal{L}|_D \cong \mathcal{O}_D$。在这种情形，可以在循环上定义
典范映射 $i^* : Z_{k + 1}(X) \to Z_k(D)$：每当
$W \subset D$ 是整闭子概形时，要求 $i^*[W] = 0$。
这种做法将在后文有用。
\end{remark}

\begin{remark}
\label{chow-remark-pullback-pairs}
设 $f : X' \to X$ 是情形 \ref{chow-situation-setup} 中那样的
$S$ 上局部有限型概形间的态射。设
$(\mathcal{L}, s, i : D \to X)$ 是定义 \ref{chow-definition-gysin-homomorphism} 中的三元组。可以置
$\mathcal{L}' = f^*\mathcal{L}$、$s' = f^*s$ 以及
$D' = X' \times_X D = Z(s')$。这给出交换图
$$
\xymatrix{
D' \ar[d]_g \ar[r]_{i'} & X' \ar[d]^f \\
D \ar[r]^i & X
}
$$
于是可以考察 $i^*$ 与 $(i')^*$ 之间的各种相容性。
\end{remark}

\begin{lemma}
\label{chow-lemma-closed-in-X-gysin}
设 $(S, \delta)$ 如情形 \ref{chow-situation-setup}。
设 $f : X' \to X$ 是 $S$ 上局部有限型概形间的固有态射。
设 $(\mathcal{L}, s, i : D \to X)$ 如定义 \ref{chow-definition-gysin-homomorphism}。依注 \ref{chow-remark-pullback-pairs} 构成图表
$$
\xymatrix{
D' \ar[d]_g \ar[r]_{i'} & X' \ar[d]^f \\
D \ar[r]^i & X
}
$$
对 $X'$ 上任意 $(k + 1)$-循环 $\alpha'$，在 $\CH_k(D)$ 中有
$i^*f_*\alpha' = g_*(i')^*\alpha'$
（这是有意义的，因为 $f_*$ 已在循环层次上定义）。
\end{lemma}

\begin{proof}
设对某个整闭子概形 $W' \subset X'$ 有 $\alpha = [W']$。
令 $W = f(W') \subset X$。若 $W' \not \subset D'$，则
$W \not \subset D$，并且
$$
[W' \cap D']_k = \text{div}_{\mathcal{L}'|_{W'}}({s'|_{W'}})
\quad\text{且}\quad
[W \cap D]_k = \text{div}_{\mathcal{L}|_W}(s|_W)
$$
因此，由引理 \ref{chow-lemma-equal-c1-as-cycles}，第一个循环沿
$f_*$ 的像等于第二个循环，故等式作为循环成立。
若 $W' \subset D'$，则 $W \subset D$；由引理 \ref{chow-lemma-equal-c1-as-cycles} 的第二个断言，
$f_*(c_1(\mathcal{L}|_{W'}) \cap [W'])$ 在 $\CH_k(W)$ 中等于
$c_1(\mathcal{L}|_W) \cap [W]$。
由注 \ref{chow-remark-infinite-sums-rational-equivalences}，
对一般的 $\alpha'$ 也得到结论。
\end{proof}

\begin{lemma}
\label{chow-lemma-gysin-flat-pullback}
设 $(S, \delta)$ 如情形 \ref{chow-situation-setup}。设
$f : X' \to X$ 是 $S$ 上局部有限型概形间相对维数为 $r$ 的平坦态射。
设 $(\mathcal{L}, s, i : D \to X)$ 如定义 \ref{chow-definition-gysin-homomorphism}。依注 \ref{chow-remark-pullback-pairs} 构成图表
$$
\xymatrix{
D' \ar[d]_g \ar[r]_{i'} & X' \ar[d]^f \\
D \ar[r]^i & X
}
$$
对 $X$ 上任意 $(k + 1)$-循环 $\alpha$，在 $\CH_{k + r}(D')$ 中有
$(i')^*f^*\alpha = g^*i^*\alpha$
（这是有意义的，因为 $f^*$ 已在循环层次上定义）。
\end{lemma}

\begin{proof}
设对某个整闭子概形 $W \subset X$ 有 $\alpha = [W]$。
令 $W' = f^{-1}(W) \subset X'$。若 $W \not \subset D$，则
$W' \not \subset D'$，并且
$$
W' \cap D' = g^{-1}(W \cap D)
$$
作为 $D'$ 的闭子概形。故等式作为循环成立；见
引理 \ref{chow-lemma-pullback-coherent}。
若 $W \subset D$，则 $W' \subset D'$ 且 $W' = g^{-1}(W)$，
并有 $[W']_{k + 1 + r} = g^*[W]$；此时等式由引理 \ref{chow-lemma-flat-pullback-cap-c1} 在 $\CH_{k + r}(D')$ 中成立。
由注 \ref{chow-remark-infinite-sums-rational-equivalences}，
对一般的 $\alpha'$ 也得到结论。
\end{proof}








\section{Gysin 同态与有理等价}
\label{chow-section-gysin}

\noindent
本节使用关键公式证明 Gysin 同态下降到有理等价类。还将证明一个重要的交换性。

\begin{lemma}
\label{chow-lemma-gysin-factors-general}
设 $(S, \delta)$ 如情形 \ref{chow-situation-setup}。
设 $X$ 在 $S$ 上局部有限型。设 $X$ 为整概形，并且
$n = \dim_\delta(X)$。设 $i : D \to X$ 是有效 Cartier 除子。
设 $\mathcal{N}$ 是可逆 $\mathcal{O}_X$-模，而 $t$ 是
$\mathcal{N}$ 的非零亚纯截面。则
$i^*\text{div}_\mathcal{N}(t) = c_1(\mathcal{N}|_D) \cap [D]_{n - 1}$
在 $\CH_{n - 2}(D)$ 中成立。
\end{lemma}

\begin{proof}
对某些 $\delta$-维数为 $n - 1$ 的整闭子概形 $Z_i \subset X$，写
$\text{div}_\mathcal{N}(t) = \sum \text{ord}_{Z_i, \mathcal{N}}(t)[Z_i]$。
可以假设族 $\{Z_i\}$ 局部有限；在
$U = X \setminus \bigcup Z_i$ 上，$t \in \Gamma(U, \mathcal{N}|_U)$
是生成元；而且 $D$ 的每个不可约分支都是某个 $Z_i$。见
《除子》引理 \ref{divisors-lemma-components-locally-finite}、
\ref{divisors-lemma-divisor-locally-finite} 和
\ref{divisors-lemma-divisor-meromorphic-locally-finite}。

\medskip\noindent
置 $\mathcal{L} = \mathcal{O}_X(D)$。以
$s \in \Gamma(X, \mathcal{O}_X(D)) = \Gamma(X, \mathcal{L})$
表示典范截面。将把第 \ref{chow-section-key} 节中的讨论应用于当前情形。
对每个 $i$，设 $\xi_i \in Z_i$ 是其一般点，并置
$B_i = \mathcal{O}_{X, \xi_i}$。对每个 $i$，选取 $B_i$ 上的生成元
$s_i \in \mathcal{L}_{\xi_i}$ 和 $t_i \in \mathcal{N}_{\xi_i}$；
但若 $Z_i \not \subset D$，则坚持选取 $s_i = s$。
写 $s = f_i s_i$ 及 $t = g_i t_i$，其中 $f_i, g_i \in B_i$。
于是 $\text{ord}_{Z_i, \mathcal{N}}(t) = \text{ord}_{B_i}(g_i)$。
另一方面，$f_i \in B_i$，并且由于对 $s_i$ 的选取，有
$$
[D]_{n - 1} = \sum \text{ord}_{B_i}(f_i)[Z_i]
$$
我们断言循环等式
$$
i^*\text{div}_\mathcal{N}(t) =
\sum \text{ord}_{B_i}(g_i) \text{div}_{\mathcal{L}|_{Z_i}}(s_i|_{Z_i})
$$
成立。更准确地说，右边是表示左边的循环。事实上，这由
$\text{div}_\mathcal{N}(t)$ 的公式及如下事实直接得到：当
$Z_i \not \subset D$ 时，
$\text{div}_{\mathcal{L}|_{Z_i}}(s_i|_{Z_i}) = [Z(s_i|_{Z_i})]_{n - 2} =
[Z_i \cap D]_{n - 2}$，因为在这种情形，
$s_i|_{Z_i} = s|_{Z_i}$ 是正则截面；见引理 \ref{chow-lemma-compute-c1}。类似地，
$$
c_1(\mathcal{N}) \cap [D]_{n - 1} =
\sum \text{ord}_{B_i}(f_i) \text{div}_{\mathcal{N}|_{Z_i}}(t_i|_{Z_i})
$$
关键公式（引理 \ref{chow-lemma-key-formula}）给出循环等式
$$
\sum \left(
\text{ord}_{B_i}(f_i) \text{div}_{\mathcal{N}|_{Z_i}}(t_i|_{Z_i}) -
\text{ord}_{B_i}(g_i) \text{div}_{\mathcal{L}|_{Z_i}}(s_i|_{Z_i}) \right) =
\sum \text{div}_{Z_i}(\partial_{B_i}(f_i, g_i))
$$
若 $Z_i \not \subset D$，则 $f_i = 1$，从而
$\text{div}_{Z_i}(\partial_{B_i}(f_i, g_i)) = 0$。因此，在 $D$ 上，
表示 $i^*\text{div}_\mathcal{N}(t)$ 和
$c_1(\mathcal{N}) \cap [D]_{n - 1}$ 的上述具体循环有理等价。
证明完毕。
\end{proof}

\begin{lemma}
\label{chow-lemma-gysin-factors}
设 $(S, \delta)$ 如情形 \ref{chow-situation-setup}。
设 $X$ 在 $S$ 上局部有限型。设
$(\mathcal{L}, s, i : D \to X)$ 如定义 \ref{chow-definition-gysin-homomorphism}。Gysin 同态下降到有理等价类，
从而给出映射 $i^* : \CH_{k + 1}(X) \to \CH_k(D)$。
\end{lemma}

\begin{proof}
设 $\alpha \in Z_{k + 1}(X)$，并假设 $\alpha \sim_{rat} 0$。
这意味着存在一族局部有限的整闭子概形 $W_j \subset X$，
其 $\delta$-维数为 $k + 2$；并存在 $f_j \in R(W_j)^*$，使得
$\alpha = \sum i_{j, *}\text{div}_{W_j}(f_j)$。
置 $X' = \coprod W_i$，并考虑注 \ref{chow-remark-pullback-pairs} 中的图表
$$
\xymatrix{
D' \ar[d]_q \ar[r]_{i'} & X' \ar[d]^p \\
D \ar[r]^i & X
}
$$
由于 $X' \to X$ 固有，由引理 \ref{chow-lemma-closed-in-X-gysin} 有
$i^*p_* = q_*(i')^*$。又因 $q_*$ 下降到有理等价类
（引理 \ref{chow-lemma-proper-pushforward-rational-equivalence}），
只须对 $X'$ 上的 $\alpha' = \sum \text{div}_{W_j}(f_j)$ 证明结论。
显然，这把问题约化到 $X$ 为整概形，且对某个 $f \in R(X)^*$ 有
$\alpha = \text{div}(f)$ 的情形。

\medskip\noindent
假设 $X$ 是整概形，且对某个 $f \in R(X)^*$ 有
$\alpha = \text{div}(f)$。若 $X = D$，则 $i^*\alpha$ 等于
$c_1(\mathcal{L}) \cap \alpha$，而由引理 \ref{chow-lemma-factors}，
它有理等价于零。若 $D \not = X$，则由引理 \ref{chow-lemma-gysin-factors-general}，$i^*\text{div}_X(f)$ 在
$\CH_{n - 2}(D)$ 中等于
$c_1(\mathcal{O}_D) \cap [D]_{n - 1}$。当然，与
$c_1(\mathcal{O}_D)$ 作帽积是零映射
（引理 \ref{chow-lemma-c1-cap-additive}）。
\end{proof}

\begin{lemma}
\label{chow-lemma-gysin-back}
设 $(S, \delta)$ 如情形 \ref{chow-situation-setup}，并设 $X$ 在 $S$ 上
局部有限型。设 $(\mathcal{L}, s, i : D \to X)$ 如定义 \ref{chow-definition-gysin-homomorphism}。则
$i^*i_* : \CH_k(D) \to \CH_{k - 1}(D)$ 把 $\alpha$ 映为
$c_1(\mathcal{L}|_D) \cap \alpha$。
\end{lemma}

\begin{proof}
这立即由循环上 $i_*$ 的定义以及定义 \ref{chow-definition-gysin-homomorphism} 中 $i^*$ 的定义得出。
\end{proof}

\begin{lemma}
\label{chow-lemma-gysin-commutes-cap-c1}
设 $(S, \delta)$ 如情形 \ref{chow-situation-setup}，并设 $X$ 在 $S$ 上
局部有限型。设 $(\mathcal{L}, s, i : D \to X)$ 是定义 \ref{chow-definition-gysin-homomorphism} 中的三元组。
设 $\mathcal{N}$ 是可逆 $\mathcal{O}_X$-模。则对所有
$\alpha \in \CH_k(X)$，在 $\CH_{k - 2}(D)$ 中有
$i^*(c_1(\mathcal{N}) \cap \alpha) = c_1(i^*\mathcal{N}) \cap i^*\alpha$。
\end{lemma}

\begin{proof}
采用与引理 \ref{chow-lemma-gysin-factors} 完全相同的证明，
结论由引理 \ref{chow-lemma-pushforward-cap-c1}、
\ref{chow-lemma-cap-commutative} 和
\ref{chow-lemma-gysin-factors-general} 得出。
\end{proof}

\begin{lemma}
\label{chow-lemma-gysin-commutes-gysin}
设 $(S, \delta)$ 如情形 \ref{chow-situation-setup}，并设 $X$ 在 $S$ 上
局部有限型。设 $(\mathcal{L}, s, i : D \to X)$ 和
$(\mathcal{L}', s', i' : D' \to X)$ 是定义 \ref{chow-definition-gysin-homomorphism} 中的两个三元组。则图表
$$
\xymatrix{
\CH_k(X) \ar[r]_{i^*} \ar[d]_{(i')^*} & \CH_{k - 1}(D) \ar[d]^{j^*} \\
\CH_{k - 1}(D') \ar[r]^{(j')^*} & \CH_{k - 2}(D \cap D')
}
$$
交换，其中每个映射都是 Gysin 映射。
\end{lemma}

\begin{proof}
以 $j : D \cap D' \to D$ 和 $j' : D \cap D' \to D'$ 表示分别对应于
$(\mathcal{L}|_{D'}, s|_{D'}$ 和 $(\mathcal{L}'_D, s|_D)$ 的闭浸入。
需要证明对所有 $\alpha \in \CH_k(X)$ 都有
$(j')^*i^*\alpha = j^* (i')^*\alpha$。
设 $W \subset X$ 是维数为 $k$ 的整闭子概形。先在
$\alpha = [W]$ 的情形证明等式，并由关键公式推出它。

\medskip\noindent
令 $\sigma$ 为 $\mathcal{L}|_W$ 的一个非零亚纯截面；若
$W \not \subset D$，则要求它等于 $s|_W$。
令 $\sigma'$ 为 $\mathcal{L}'|_W$ 的一个非零亚纯截面；若
$W \not \subset D'$，则要求它等于 $s'|_W$。写
$$
\text{div}_{\mathcal{L}|_W}(\sigma) =
\sum \text{ord}_{Z_i, \mathcal{L}|_W}(\sigma)[Z_i] = \sum n_i[Z_i]
$$
以及类似的等式
$$
\text{div}_{\mathcal{L}'|_W}(\sigma') =
\sum \text{ord}_{Z_i, \mathcal{L}'|_W}(\sigma')[Z_i] = \sum n'_i[Z_i]
$$
如第 \ref{chow-section-key} 节的讨论。于是，若 $n_i \not = 0$，
则 $Z_i \subset D$；若 $n'_i \not = 0$，则 $Z'_i \subset D'$。
对每个 $i$，设 $\xi_i \in Z_i$ 为一般点。与第 \ref{chow-section-key} 节中一样，对每个 $i$ 选取元素
$\sigma_i \in \mathcal{L}_{\xi_i}$，以及相应的\ 
$\sigma'_i \in \mathcal{L}'_{\xi_i}$；它们分别生成
$B_i = \mathcal{O}_{W, \xi_i}$ 上的相应模，并且若
$Z_i \not \subset D$，则前者等于 $s$ 的像；\ 相应地，若
$Z_i \not \subset D'$，则\ 后者等于 $s'$ 的像。
写 $\sigma = f_i \sigma_i$ 及 $\sigma' = f'_i\sigma'_i$，使得
$n_i = \text{ord}_{B_i}(f_i)$ 且
$n'_i = \text{ord}_{B_i}(f'_i)$。由定义，作为循环有
$$
(j')^*i^*[W] =
\sum \text{ord}_{B_i}(f_i) \text{div}_{\mathcal{L}'|_{Z_i}}(\sigma'_i|_{Z_i})
$$
以及
$$
j^*(i')^*[W] =
\sum \text{ord}_{B_i}(f'_i) \text{div}_{\mathcal{L}|_{Z_i}}(\sigma_i|_{Z_i})
$$
此时关键公式（引理 \ref{chow-lemma-key-formula}）给出循环等式
$$
\sum \left(
\text{ord}_{B_i}(f_i) \text{div}_{\mathcal{L}'|_{Z_i}}(\sigma'_i|_{Z_i}) -
\text{ord}_{B_i}(f'_i) \text{div}_{\mathcal{L}|_{Z_i}}(\sigma_i|_{Z_i})
\right) =
\sum \text{div}_{Z_i}(\partial_{B_i}(f_i, f'_i))
$$
注意，若 $Z_i \not \subset D \cap D'$，则 $f_i = 1$ 或
$f'_i = 1$，从而
$\text{div}_{Z_i}(\partial_{B_i}(f_i, f'_i)) = 0$。
因此，在 $D \cap D' \cap W$ 上，表示
$(j')^*i^*[W]$ 与 $j^*(i')^*[W]$ 的上述具体循环有理等价。
由注 \ref{chow-remark-infinite-sums-rational-equivalences}，
对一般的 $\alpha$ 也得到结论。
\end{proof}
















\section{相对有效 Cartier 除子}
\label{chow-section-relative-effective-cartier}

\noindent
相对有效 Cartier 除子在《除子》第 \ref{divisors-section-effective-Cartier-morphisms} 节中定义并研究。
为建立向量丛陈类的基本结果，只需考虑周围概形和有效 Cartier 除子
都在基底上平坦的情形。

\begin{lemma}
\label{chow-lemma-relative-effective-cartier}
设 $(S, \delta)$ 如情形 \ref{chow-situation-setup}。
设 $X$、$Y$ 在 $S$ 上局部有限型。设 $p : X \to Y$ 是相对维数为
$r$ 的平坦态射。设 $i : D \to X$ 是相对有效 Cartier 除子
（《除子》定义 \ref{divisors-definition-relative-effective-Cartier-divisor}）。
置 $\mathcal{L} = \mathcal{O}_X(D)$。对任意
$\alpha \in \CH_{k + 1}(Y)$，在 $\CH_{k + r}(D)$ 中有
$$
i^*p^*\alpha = (p|_D)^*\alpha
$$
并且在 $\CH_{k + r}(X)$ 中有
$$
c_1(\mathcal{L}) \cap p^*\alpha = i_* ((p|_D)^*\alpha)
$$
\end{lemma}

\begin{proof}
设 $W \subset Y$ 是 $\delta$-维数为 $k + 1$ 的整闭子概形。
由《除子》引理 \ref{divisors-lemma-relative-Cartier}，
$D \cap p^{-1}W$ 是 $p^{-1}W$ 上的有效 Cartier 除子。
由引理 \ref{chow-lemma-easy-gysin}，下式中的第一个等式成立：
$$
i^*[p^{-1}W]_{k + r + 1} =
[D \cap p^{-1}W]_{k + r} =
[(p|_D)^{-1}(W)]_{k + r}.
$$
第二个等式成立，是因为作为概形有
$D \cap p^{-1}(W) = (p|_D)^{-1}(W)$。
又因定义给出 $p^*[W] = [p^{-1}W]_{k + r + 1}$，可知作为循环有
$i^*p^*[W] = (p|_D)^*[W]$。若
$\alpha = \sum m_j[W_j]$ 是一般的 $k + 1$ 循环，则作为循环有
$i^*\alpha = \sum m_j i^*p^*[W_j] = \sum m_j(p|_D)^*[W_j]$。
这就证明了第一个等式。再应用引理 \ref{chow-lemma-support-cap-effective-Cartier}，即可由第一个等式推出第二个。
\end{proof}








\section{仿射丛}
\label{chow-section-affine-vector}

\noindent
对仿射丛而言，Chow 群上的拉回映射是满射。

\begin{lemma}
\label{chow-lemma-pullback-affine-fibres-surjective}
设 $(S, \delta)$ 如情形 \ref{chow-situation-setup}。
设 $X$、$Y$ 在 $S$ 上局部有限型。设 $f : X \to Y$ 是相对维数为
$r$ 的平坦态射。假设对每个 $y \in Y$，都存在开邻域
$U \subset Y$，使得 $f|_{f^{-1}(U)} : f^{-1}(U) \to U$
可与态射 $U \times \mathbf{A}^r \to U$ 等同。则对所有
$k \in \mathbf{Z}$，映射
$f^* : \CH_k(Y) \to \CH_{k + r}(X)$ 都是满射。
\end{lemma}

\begin{proof}
设 $\alpha \in \CH_{k + r}(X)$。写
$\alpha = \sum m_j[W_j]$，其中 $m_j \not = 0$，而 $W_j$
是两两不同、$\delta$-维数为 $k + r$ 的整闭子概形。
于是族 $\{W_j\}$ 在 $X$ 中局部有限。对任意拟紧开集
$V \subset Y$，只有有限多个 $j$ 使
$f^{-1}(V) \cap W_j$ 非空。因此，像的闭包所成的族
$Z_j = \overline{f(W_j)}$ 是 $Y$ 中局部有限的一族整闭子概形。

\medskip\noindent
考虑纤维积图表
$$
\xymatrix{
f^{-1}(Z_j) \ar[r] \ar[d]_{f_j} & X \ar[d]^f \\
Z_j \ar[r] & Y
}
$$
假设对某个 $k$-循环 $\beta_j \in \CH_k(Z_j)$，
$[W_j] \in Z_{k + r}(f^{-1}(Z_j))$ 与 $f_j^*\beta_j$ 有理等价。
那么 $\beta = \sum m_j \beta_j$ 是 $Y$ 上的 $k$-循环，而
$f^*\beta = \sum m_j f_j^*\beta_j$ 与 $\alpha$ 有理等价
（见注 \ref{chow-remark-infinite-sums-rational-equivalences}）。
这把问题约化到如下情形：$Y$ 是整概形，并且
$\alpha = [W]$，其中所指的是支配 $Y$ 的 $X$ 中整闭子概形。
特别地，可以假设 $d = \dim_\delta(Y) < \infty$。

\medskip\noindent
于是可以对 $d = \dim_\delta(Y)$ 作归纳。若 $d < k$，则
$\CH_{k + r}(X) = 0$，引理成立。由假设，存在稠密开集
$V \subset Y$，使得作为 $V$ 上的概形有
$f^{-1}(V) \cong V \times \mathbf{A}^r$。
假设能够证明对某个 $\beta \in Z_k(V)$ 有
$\alpha|_{f^{-1}(V)} = f^*\beta$。由引理 \ref{chow-lemma-exact-sequence-open}，对某个 $\beta' \in Z_k(Y)$ 有
$\beta = \beta'|_V$。由引理 \ref{chow-lemma-restrict-to-open} 的正合列
$\CH_k(f^{-1}(Y \setminus V)) \to \CH_k(X) \to \CH_k(f^{-1}(V))$，
可知 $\alpha - f^*\beta'$ 来自某个循环
$\alpha' \in \CH_{k + r}(f^{-1}(Y \setminus V))$。
由于 $\dim_\delta(Y \setminus V) < d$，对 $d$ 使用归纳即可。

\medskip\noindent
所以可以假设 $X = Y \times \mathbf{A}^r$。此时可以把 $f$ 分解为
$$
X = Y \times \mathbf{A}^r \to
Y \times \mathbf{A}^{r - 1} \to \ldots \to
Y \times \mathbf{A}^1 \to Y.
$$
故只须处理 $r = 1$ 的情形。由证明第二段中的论证，问题约化到
$\alpha = [W]$、$Y$ 为整概形且 $W \to Y$ 支配的情形。
仍可对 $d = \dim_\delta(Y)$ 作归纳。若
$W = Y \times \mathbf{A}^1$，则 $[W] = f^*[Y]$。
最后，若 $W \subset Y \times \mathbf{A}^1$ 是真包含，则
$W \to Y$ 诱导有限域扩张 $R(W)/R(Y)$。
设 $P(T) \in R(Y)[T]$ 是首一不可约多项式，使得 $W \to Y$
的一般纤维在 $\mathbf{A}^1_{R(Y)}$ 中由 $P$ 截出。
选取非空开集 $V \subset Y$，使 $P \in \Gamma(V, \mathcal{O}_Y)[T]$，
并且 $W \cap f^{-1}(V)$ 仍由 $P$ 截出。于是
$\alpha|_{f^{-1}(V)} \sim_{rat} 0$，故对
$(Y \setminus V) \times \mathbf{A}^1$ 上某个循环 $\alpha'$ 有
$\alpha \sim_{rat} \alpha'$。对维数应用归纳即可。
\end{proof}

\begin{lemma}
\label{chow-lemma-linebundle}
设 $(S, \delta)$ 如情形 \ref{chow-situation-setup}。
设 $X$ 在 $S$ 上局部有限型，$\mathcal{L}$ 是可逆
$\mathcal{O}_X$-模。设
$$
p :
L = \underline{\Spec}(\text{Sym}^*(\mathcal{L}))
\longrightarrow
X
$$
是 $X$ 上相应的向量丛。则对所有 $k$，
$p^* : \CH_k(X) \to \CH_{k + 1}(L)$ 都是同构。
\end{lemma}

\begin{proof}
满射性见引理 \ref{chow-lemma-pullback-affine-fibres-surjective}。
设 $o : X \to L$ 是 $L \to X$ 的零截面，即对应于满射
$\text{Sym}^*(\mathcal{L}) \to \mathcal{O}_X$ 的态射；该满射对所有
$n > 0$ 都把 $\mathcal{L}^{\otimes n}$ 映为零。
于是 $p \circ o = \text{id}_X$，且 $o(X)$ 是 $L$ 上的有效
Cartier 除子。故由引理 \ref{chow-lemma-relative-effective-cartier}，
$o^* \circ p^* = \text{id}$，从而 $p^*$ 也是单射。
\end{proof}

\begin{remark}
\label{chow-remark-when-isomorphism}
后文将看到（引理 \ref{chow-lemma-vectorbundle}），若 $X$ 是 $Y$ 上秩为
$r$ 的向量丛，则拉回映射 $\CH_k(Y) \to \CH_{k + r}(X)$ 是同构。
只要 $X \to Y$ 满足引理 \ref{chow-lemma-pullback-affine-fibres-surjective} 的假设，这一结论就成立；
见 \cite[Lemma 2.2]{Totaro-group}。在注 \ref{chow-remark-higher-chow-isomorphism} 中，我们将用高阶 Chow 群勾勒一个证明。
\end{remark}

\begin{lemma}
\label{chow-lemma-linebundle-formulae}
在引理 \ref{chow-lemma-linebundle} 的情形中，以 $o : X \to L$
表示零截面（见该引理的证明）。则有
\begin{enumerate}
\item $o(X)$ 是 $p^*\mathcal{L}^{\otimes -1}$ 的某个正则整体截面的
零点概形；
\item 因为 $o$ 是闭浸入，有 $o_* : \CH_k(X) \to \CH_k(L)$；
\item 因为 $o(X)$ 是有效 Cartier 除子，有
$o^* : \CH_{k + 1}(L) \to \CH_k(X)$；
\item $o^* p^* : \CH_k(X) \to \CH_k(X)$ 是恒等映射；
\item 对任意 $\alpha \in \CH_k(X)$，有
$o_*\alpha = - p^*(c_1(\mathcal{L}) \cap \alpha)$；以及
\item $o^* o_* : \CH_k(X) \to \CH_{k - 1}(X)$ 等于映射
$\alpha \mapsto - c_1(\mathcal{L}) \cap \alpha$。
\end{enumerate}
\end{lemma}

\begin{proof}
由于 $p_*\mathcal{O}_L = \text{Sym}^*(\mathcal{L})$，投影公式
（《层上同调》引理 \ref{cohomology-lemma-projection-formula}）给出
$p_*(p^*\mathcal{L}^{\otimes -1}) =
\text{Sym}^*(\mathcal{L}) \otimes_{\mathcal{O}_X} \mathcal{L}^{\otimes -1}$，
而 (1) 中提到的截面是典范平凡化
$\mathcal{O}_X \to
\mathcal{L} \otimes_{\mathcal{O}_X} \mathcal{L}^{\otimes -1}$。
这里略去证明该截面的零点轨迹恰为 $o(X)$。这证明了 (1)。

\medskip\noindent
(2)、(3) 和 (4) 已在引理 \ref{chow-lemma-linebundle} 的证明过程中看到。
当然，(4) 正是引理 \ref{chow-lemma-relative-effective-cartier} 中的第一个公式。

\medskip\noindent
(5) 由引理 \ref{chow-lemma-relative-effective-cartier} 中的第二个公式、
与 $c_1$ 作帽积的可加性（引理 \ref{chow-lemma-c1-cap-additive}），
以及与 $c_1$ 作帽积和平坦拉回交换这一事实
（引理 \ref{chow-lemma-flat-pullback-cap-c1}）得出。

\medskip\noindent
(6) 由引理 \ref{chow-lemma-gysin-back} 以及
$o^*p^*\mathcal{L} = \mathcal{L}$ 得出。
\end{proof}

\begin{lemma}
\label{chow-lemma-decompose-section}
设 $Y$ 是概形。设 $\mathcal{L}_i$（$i = 1, 2$）是可逆
$\mathcal{O}_Y$-模。设 $s$ 是
$\mathcal{L}_1 \otimes_{\mathcal{O}_Y} \mathcal{L}_2$ 的整体截面，
并以 $i : D \to Y$ 表示 $s$ 的零点概形。则存在交换图
$$
\xymatrix{
D_1 \ar[r]_{i_1} \ar[d]_{p_1} &
L \ar[d]^p &
D_2 \ar[l]^{i_2} \ar[d]^{p_2} \\
D \ar[r]^i &
Y &
D \ar[l]_i
}
$$
以及 $p^*\mathcal{L}_i$ 的截面 $s_i$，使得：
\begin{enumerate}
\item $p^*s = s_1 \otimes s_2$；
\item $p$ 有限型、平坦且相对维数为 $1$；
\item $D_i$ 是 $s_i$ 的零点概形；
\item 对 $i = 1, 2$，在 $D$ 上有
$D_i \cong
\underline{\Spec}(\text{Sym}^*(\mathcal{L}_{3 - i}^{\otimes -1})|_D))$；
\item $p^{-1}D = D_1 \cup D_2$（概形论并）；
\item $D_1 \cap D_2$（概形论交）同构地映到 $D$；以及
\item 对 $i = 1, 2$，$D_1 \cap D_2 \to D_i$ 是线丛
$D_i \to D$ 的零截面。
\end{enumerate}
此外，该图表和各截面 $s_i$ 的形成与任意基变换交换。
\end{lemma}

\begin{proof}
设 $p : X \to Y$ 是如下拟凝聚 $\mathcal{O}_Y$-代数层的相对谱：
$$
\mathcal{A} =
\left(\bigoplus\nolimits_{a_1, a_2 \geq 0}
\mathcal{L}_1^{\otimes -a_1} \otimes_{\mathcal{O}_Y}
\mathcal{L}_2^{\otimes -a_2}\right)/\mathcal{J}
$$
其中 $\mathcal{J}$ 是由形如 $st - t$ 的局部截面生成的理想；
这里 $t$ 是满足 $a_1, a_2 > 0$ 的任意直和项
$\mathcal{L}_1^{\otimes -a_1} \otimes \mathcal{L}_2^{\otimes -a_2}$
的局部截面。把截面 $s_i$ 看作映射
$p^*\mathcal{L}_i^{\otimes -1} \to \mathcal{O}_X$，则它们定义为映射
$\mathcal{L}_i^{\otimes -1} \to \mathcal{A} = p_*\mathcal{O}_X$
的伴随映射。对任意 $y \in Y$，可以选取包含 $y$ 的仿射开集
$V \subset Y$，设 $V = \Spec(A)$，并选取平凡化
$z_i : \mathcal{O}_V \to \mathcal{L}_i^{\otimes -1}|_V$。
注意，$f = s(z_1z_2) \in A$ 截出闭子概形 $D$。显然
$$
p^{-1}(V) = \Spec(B[z_1, z_2]/(z_1 z_2 - f))
$$
由于 $D_i$ 由 $z_i$ 截出，其余各点都很清楚。
\end{proof}

\begin{lemma}
\label{chow-lemma-decompose-section-formulae}
在引理 \ref{chow-lemma-decompose-section} 的情形中，假设 $Y$ 是
情形 \ref{chow-situation-setup} 中那样在 $(S, \delta)$ 上局部有限型。
则对所有 $\alpha \in \CH_k(Y)$，在 $\CH_k(D_1)$ 中有
$i_1^*p^*\alpha = p_1^*i^*\alpha$。
\end{lemma}

\begin{proof}
设 $W \subset Y$ 是 $\delta$-维数为 $k$ 的整闭子概形。
分两种情形讨论。

\medskip\noindent
假设 $W \subset D$。由 Gysin 同态的定义以及引理 \ref{chow-lemma-c1-cap-additive} 的可加性，在 $\CH_{k - 1}(D)$ 中有
$i^*[W] = c_1(\mathcal{L}_1) \cap [W] + c_1(\mathcal{L}_2) \cap [W]$。
因此
$p_1^*i^*[W] =
p_1^*(c_1(\mathcal{L}_1) \cap [W]) + p_1^*(c_1(\mathcal{L}_2) \cap [W])$。
另一方面，平坦拉回的构造给出
$p^*[W] = [p^{-1}(W)]_{k + 1}$。由该引理
$p^{-1}(W) = W_1 \cup W_2$（概形论地），其中
$W_i = p_i^{-1}(W)$ 是 $W$ 上的线丛
（因为图表的形成与基变换交换）。于是
$[p^{-1}(W)]_{k + 1} = [W_1] + [W_2]$，因为 $W_i$ 是 $L$ 中
$\delta$-维数为 $k + 1$ 的整闭子概形。因此
\begin{align*}
i_1^*p^*[W]
& =
i_1^*[p^{-1}(W)]_{k + 1} \\
& =
i_1^*([W_1] + [W_2]) \\
& =
c_1(p_1^*\mathcal{L}_1) \cap [W_1] + [W_1 \cap W_2]_k \\
& =
c_1(p_1^*\mathcal{L}_1) \cap p_1^*[W] + [W_1 \cap W_2]_k \\
& =
p_1^*(c_1(\mathcal{L}_1) \cap [W]) + [W_1 \cap W_2]_k
\end{align*}
这里使用了 Gysin 同态的构造、平坦拉回的定义（用于第二个等式），
以及 $c_1 \cap -$ 与平坦拉回的相容性
（引理 \ref{chow-lemma-flat-pullback-cap-c1}）。
由于 $W_1 \cap W_2$ 是线丛 $W_1 \to W$ 的零截面，引理 \ref{chow-lemma-linebundle-formulae} 给出
$[W_1 \cap W_2]_k = p_1^*(c_1(\mathcal{L}_2) \cap [W])$。
注意，这里使用了 $D_1$ 是 $\mathcal{L}_2$ 的逆的相对谱所给出的线丛
这一事实。因此所得结果与前式相同。

\medskip\noindent
假设 $W \not \subset D$。此时，$i_1^*p^*[W]$ 与
$p_1^*i^*[W]$ 都由 $W$ 在 $D_1$ 中的概形论逆像所关联的
$k - 1$ 循环表示。
\end{proof}

\begin{lemma}
\label{chow-lemma-normal-cone-effective-Cartier}
在情形 \ref{chow-situation-setup} 中，设 $X$ 是 $S$ 上局部有限型的概形。
设 $(\mathcal{L}, s, i : D \to X)$ 是定义 \ref{chow-definition-gysin-homomorphism} 中的三元组。则存在交换图
$$
\xymatrix{
D' \ar[r]_{i'} \ar[d]_p & X' \ar[d]^g \\
D \ar[r]^i & X
}
$$
满足
\begin{enumerate}
\item $p$ 和 $g$ 有限型、平坦且相对维数为 $1$；
\item 对所有 $k$，$p^* : \CH_k(D) \to \CH_{k + 1}(D')$ 都是单射；
\item $D' \subset X'$ 是某个整体截面
$s' \in \Gamma(X', \mathcal{O}_{X'})$ 的零点概形；
\item 作为映射 $\CH_k(X) \to \CH_k(D')$，有
$p^*i^* = (i')^*g^*$。
\end{enumerate}
此外，经由局部有限型态射 $Y \to X$ 作任意基变换后，这些性质仍成立。
\end{lemma}

\begin{proof}
注意，$(i')^*$ 已有定义，因为存在定义 \ref{chow-definition-gysin-homomorphism} 中那样的三元组
$(\mathcal{O}_{X'}, s', i' : D' \to X')$。因此该陈述有意义。

\medskip\noindent
置 $\mathcal{L}_1 = \mathcal{O}_X$、$\mathcal{L}_2 = \mathcal{L}$，
并把引理 \ref{chow-lemma-decompose-section} 应用于
$\mathcal{L} = \mathcal{L}_1 \otimes_{\mathcal{O}_X} \mathcal{L}_2$
的截面 $s$。取 $D' = D_1$。此时结论由该引理和引理 \ref{chow-lemma-decompose-section-formulae} 得出，而单射性由引理 \ref{chow-lemma-linebundle} 得出。
\end{proof}

\begin{remark}
\label{chow-remark-higher-chow-isomorphism}
设 $(S, \delta)$ 如情形 \ref{chow-situation-setup}。
设 $Y$ 在 $S$ 上局部有限型，$r \geq 0$，而
$f : X \to Y$ 是概形态射。假设每个 $y \in Y$ 都包含于某个开集
$V \subset Y$，并且作为 $V$ 上的概形有
$f^{-1}(V) \cong V \times \mathbf{A}^r$。
本注记勾勒事实
$f^* : \CH_k(Y) \to \CH_{k + r}(X)$ 为同构的一个证明。
首先，由引理 \ref{chow-lemma-pullback-affine-fibres-surjective}，
该映射是满射。设 $\alpha \in \CH_k(Y)$ 且 $f^*\alpha = 0$。
下面证明 $\alpha = 0$。

\medskip\noindent
第 1 步。可以假设 $\dim_\delta(Y) < \infty$。（实际遇到的所有情形中，
这一点都是立即成立的，所以建议读者略过本步。）事实上，在 $X$ 上见证
$f^*\alpha = 0$ 的任一有理等价会用到一族局部有限、维数为
$k + r + 1$ 的整闭子概形。取它们在 $Y$ 中各像的闭包之并，
得到满足 $\dim_\delta(Y') \leq k + r + 1$ 的闭子概形
$Y' \subset Y$，使得 $\alpha$ 是某个 $\alpha' \in \CH_k(Y')$ 的像，
并且 $(f')^*\alpha = 0$，其中 $f'$ 是 $f$ 到 $Y'$ 的基变换。

\medskip\noindent
第 2 步。假设 $d = \dim_\delta(Y) < \infty$，并对 $d$ 作归纳。
若 $d < k$，则 $\alpha = 0$，结论成立；这是归纳的初始情形。
一般地，对 $f$ 的假设表明，可以选取稠密开集
$V \subset Y$，使 $U = f^{-1}(V) = \mathbf{A}^r_V$。
以 $Y' \subset Y$ 表示 $V$ 的补集，视为既约闭子概形，并置
$X' = f^{-1}(Y')$。考虑
$$
\xymatrix{
\CH^M_{k + r}(U, 1) \ar[r] &
\CH_{k + r}(X') \ar[r] &
\CH_{k + r}(X) \ar[r] &
\CH_{k + r}(U) \ar[r] &
0 \\
\CH^M_k(V, 1) \ar[r] \ar[u] &
\CH_k(Y') \ar[r] \ar[u] &
\CH_k(Y) \ar[r] \ar[u] &
\CH_k(V) \ar[r] \ar[u] &
0
}
$$
这里使用了 $V$ 和 $U$ 的第一高阶 Chow 群、注 \ref{chow-remark-higher-chow} 中构造的六项正合列，以及这些高阶 Chow 群
上的平坦拉回和平坦拉回与六项正合列的相容性。
由于 $U = \mathbf{A}^r_V$，右侧竖直映射是同构。
根据对 $d$ 的归纳，映射 $\CH_k(Y') \to \CH_{k + r}(X')$ 是双射。
所以只须证明
$$
\CH^M_k(V, 1) \longrightarrow \CH^M_{k + r}(U, 1)
$$
为满射即可完成论证。仿照引理 \ref{chow-lemma-pullback-affine-fibres-surjective} 的证明进行论证，
这又约化到下面的第 3 步。

\medskip\noindent
第 3 步。设 $F$ 是域。则 $\CH^M_0(\mathbf{A}^1_F, 1) = 0$。
（在上述引理的证明中，我们类似地证明了
$\CH_0(\mathbf{A}^1_F) = 0$。）有
$$
\CH^M_0(\mathbf{A}^1_F, 1) = \Coker\left(
\partial : K^M_2(F(T)) \longrightarrow
\bigoplus\nolimits_{\mathfrak p \subset F[T]\text{ 为极大理想}}
\kappa(\mathfrak p)^*\right)
$$
证明该余核为零的经典论证，是对 $\kappa(\mathfrak p)/F$ 的次数作归纳，
证明对应于 $\mathfrak p$ 的直和项在像中。
若 $\mathfrak p$ 由首一不可约多项式 $P(T) \in F[T]$ 生成，而
$u \in \kappa(x)^*$ 是某个满足 $\deg(Q) < \deg(P)$ 的
$Q(T) \in F[T]$ 的剩余类，则可以证明 $\partial(Q, P)$ 在
$\mathfrak p$ 处给出元素 $u$，并且可能还在整除 $Q$ 的一些次数更低
的素理想处给出其他单位。这就完成了证明概要。
\end{remark}











\section{双变相交理论}
\label{chow-section-bivariant}

\noindent
为妥善讨论向量丛的高阶陈类，我们依照 \cite{F} 引入双变 Chow 类。
我们的定义与 \cite{F} 有两点不同：(1) 所处的框架不同；(2) 仅要求
双变类与可逆模截面的零点概形所给出的 Gysin 同态交换
（第 \ref{chow-section-intersecting-effective-Cartier} 节）。
后文将在引理 \ref{chow-lemma-gysin-commutes} 中看到，我们的双变类
与所有高余维 Gysin 同态交换，因而满足 \cite{F} 中要求的全部性质；
另见 \cite[Theorem 17.1]{F}。

\begin{definition}
\label{chow-definition-bivariant-class}
\begin{reference}
类似于 \cite[Definition 17.1]{F}
\end{reference}
设 $(S, \delta)$ 如情形 \ref{chow-situation-setup}。
设 $f : X \to Y$ 是 $S$ 上局部有限型概形间的态射，并设
$p \in \mathbf{Z}$。{\it $f$ 的次数为 $p$ 的双变类 $c$}由如下规则给出：
对每个局部有限型态射 $Y' \to Y$ 及每个 $k$，指定映射
$$
c \cap - : \CH_k(Y') \longrightarrow \CH_{k - p}(X')
$$
其中 $X' = Y' \times_Y X$，并满足以下条件：
\begin{enumerate}
\item 若 $Y'' \to Y'$ 是固有态射，则对 $Y''$ 上所有 $\alpha''$，
有 $c \cap (Y'' \to Y')_*\alpha'' = (X'' \to X')_*(c \cap \alpha'')$，
其中 $X'' = Y'' \times_Y X$；
\item 若 $Y'' \to Y'$ 是相对维数固定、局部有限型的平坦态射，
则对 $Y'$ 上所有 $\alpha'$，有
$c \cap (Y'' \to Y')^*\alpha' = (X'' \to X')^*(c \cap \alpha')$；以及
\item 若 $(\mathcal{L}', s', i' : D' \to Y')$ 如定义 \ref{chow-definition-gysin-homomorphism}，且其到 $X'$ 的拉回为
$(\mathcal{N}', t', j' : E' \to X')$，则对 $Y'$ 上所有 $\alpha'$，有
$c \cap (i')^*\alpha' = (j')^*(c \cap \alpha')$。
\end{enumerate}
$f$ 的所有次数为 $p$ 的双变类所成的集合记为 $A^p(X \to Y)$。
\end{definition}

\noindent
设 $(S, \delta)$ 如情形 \ref{chow-situation-setup}。设 $X \to Y$
和 $Y \to Z$ 是 $S$ 上局部有限型概形间的态射，并设
$p \in \mathbf{Z}$。显然，$A^p(X \to Y)$ 是 Abel 群。
此外，显然存在结合的双线性复合
$$
A^p(X \to Y) \times A^q(Y \to Z) \to A^{p + q}(X \to Z)
$$

\begin{lemma}
\label{chow-lemma-flat-pullback-bivariant}
设 $(S, \delta)$ 如情形 \ref{chow-situation-setup}。
设 $f : X \to Y$ 是 $S$ 上局部有限型概形间相对维数为 $r$ 的平坦态射。
对 $Y' \to Y$ 指定
$(f')^* : \CH_k(Y') \to \CH_{k + r}(X')$ 的规则是次数为 $-r$ 的双变类，
其中 $X' = X \times_Y Y'$。
\end{lemma}

\begin{proof}
这由引理 \ref{chow-lemma-flat-pullback-rational-equivalence}、
\ref{chow-lemma-compose-flat-pullback}、
\ref{chow-lemma-flat-pullback-proper-pushforward} 和
\ref{chow-lemma-gysin-flat-pullback} 得出。
\end{proof}

\begin{lemma}
\label{chow-lemma-gysin-bivariant}
设 $(S, \delta)$ 如情形 \ref{chow-situation-setup}。
设 $X$ 在 $S$ 上局部有限型。设
$(\mathcal{L}, s, i : D \to X)$ 是定义 \ref{chow-definition-gysin-homomorphism} 中的三元组。
对 $f : X' \to X$ 指定
$(i')^* : \CH_k(X') \to \CH_{k - 1}(D')$ 的规则是次数为 $1$ 的双变类，
其中 $D' = D \times_X X'$。
\end{lemma}

\begin{proof}
这由引理 \ref{chow-lemma-gysin-factors}、
\ref{chow-lemma-closed-in-X-gysin}、\ref{chow-lemma-gysin-flat-pullback} 和
\ref{chow-lemma-gysin-commutes-gysin} 得出。
\end{proof}

\begin{lemma}
\label{chow-lemma-push-proper-bivariant}
设 $(S, \delta)$ 如情形 \ref{chow-situation-setup}。
设 $f : X \to Y$ 和 $g : Y \to Z$ 是 $S$ 上局部有限型概形间的态射。
设 $c \in A^p(X \to Z)$，并假设 $f$ 固有。
对 $Z' \to Z$ 指定
$\alpha \longmapsto f'_*(c \cap \alpha)$ 的规则是双变类，
记为 $f_* \circ c \in A^p(Y \to Z)$。
\end{lemma}

\begin{proof}
这由引理 \ref{chow-lemma-compose-pushforward}、
\ref{chow-lemma-flat-pullback-proper-pushforward} 和
\ref{chow-lemma-closed-in-X-gysin} 得出。
\end{proof}

\begin{remark}
\label{chow-remark-restriction-bivariant}
设 $(S, \delta)$ 如情形 \ref{chow-situation-setup}。设 $X \to Y$
和 $Y' \to Y$ 是 $S$ 上局部有限型概形间的态射，并置
$X' = Y' \times_Y X$。则有明显的限制映射
$$
A^p(X \to Y) \longrightarrow A^p(X' \to Y'),\quad
c \longmapsto res(c)
$$
其构造如下：把 $Y'$ 上局部有限型的概形 $Y''$ 视为 $Y$ 上局部有限型
的概形，并对任意 $\alpha'' \in \CH_k(Y'')$ 置
$res(c) \cap \alpha'' = c \cap \alpha''$。这个限制运算显然与复合相容。
\end{remark}

\begin{remark}
\label{chow-remark-bivariant-commute}
设 $(S, \delta)$ 如情形 \ref{chow-situation-setup}，并设 $X$ 在 $S$ 上
局部有限型。对 $i = 1, 2$，设 $Z_i \to X$ 是局部有限型概形的态射。
设 $c_i \in A^{p_i}(Z_i \to X)$（$i = 1, 2$）是双变类。
对任意 $\alpha \in \CH_k(X)$，可以问在
$\CH_{k - p_1 - p_2}(Z_1 \times_X Z_2)$ 中是否有
$$
c_1 \cap c_2 \cap \alpha = c_2 \cap c_1 \cap \alpha
$$
若该等式成立，并且经任意局部有限型基变换 $X' \to X$ 后仍成立，
就称 $c_1$ 和 $c_2$ {\it 交换}。当然，这等价于在
$A^{p_1 + p_2}(Z_1 \times_X Z_2 \to X)$ 中有
$$
res(c_1) \circ c_2 = res(c_2) \circ c_1
$$
这里 $res(c_1) \in A^{p_1}(Z_1 \times_X Z_2 \to Z_2)$ 是注 \ref{chow-remark-restriction-bivariant} 中那样的 $c_1$ 的限制；
$res(c_2)$ 类似。
\end{remark}

\begin{example}
\label{chow-example-gysin-commute}
设 $(S, \delta)$ 如情形 \ref{chow-situation-setup}，并设 $X$ 在 $S$ 上
局部有限型。设 $(\mathcal{L}, s, i : D \to X)$ 是定义 \ref{chow-definition-gysin-homomorphism} 中的三元组。
设 $Z \to X$ 是局部有限型概形的态射，而 $c \in A^p(Z \to X)$
是双变类。则引理 \ref{chow-lemma-gysin-bivariant} 的双变 Gysin 类
$c' \in A^1(D \to X)$ 按注 \ref{chow-remark-bivariant-commute} 的意义与 $c$ 交换。
这正是定义 \ref{chow-definition-bivariant-class} 中条件 (3) 的重述。
\end{example}

\begin{remark}
\label{chow-remark-more-general-bivariant}
文献似乎没有考虑如下更一般类型的双变类。设给定图表
$$
X \longrightarrow Z \longleftarrow Y
$$
其中各概形如情形 \ref{chow-situation-setup} 中那样在
$(S, \delta)$ 上局部有限型。设 $p \in \mathbf{Z}$。
可以考虑如下规则 $c$：对每个局部有限型态射 $Z' \to Z$
及所有 $k \in \mathbf{Z}$，指定映射
$$
c \cap - : \CH_k(Y') \longrightarrow \CH_{k - p}(X')
$$
其中 $X' = X \times_Z Z'$ 且 $Y' = Z' \times_Z Y$，并与以下运算相容：
\begin{enumerate}
\item 给定固有态射 $Z'' \to Z'$ 时的固有推前；
\item 给定相对维数固定的平坦态射 $Z'' \to Z'$ 时的平坦拉回；以及
\item 给定定义 \ref{chow-definition-gysin-homomorphism} 中那样的
$D' \subset Z'$ 时的 Gysin 映射。
\end{enumerate}
这里略去详细表述。把所有这类运算所成的集合记为
$A^p(X \to Z \leftarrow Y)$。这个概念有用的一个简单例子，是存在固有态射
$f : X_2 \to X_1$ 的情形。此时，按定义 \ref{chow-definition-bivariant-class} 的意义，$f_*$ 不是双变运算；
但按上述推广意义它是，即
$f_* \in A^0(X_1 \to X_1 \leftarrow X_2)$。
\end{remark}





\section{Chow 上同调与第一陈类}
\label{chow-section-chow-cohomology}

\noindent
我们最关心的是 $A^p(X) = A^p(X \to X)$；它始终指
$\text{id}_X$ 的双变上同调类。陈类正是位于这里。

\begin{definition}
\label{chow-definition-chow-cohomology}
设 $(S, \delta)$ 如情形 \ref{chow-situation-setup}。
设 $X$ 在 $S$ 上局部有限型。$X$ 的{\it Chow 上同调}是分次
$\mathbf{Z}$-代数 $A^*(X)$，其次数 $p$ 分量为 $A^p(X \to X)$。
\end{definition}

\noindent
警告：尚不清楚 $A^*(X)$ 上的 $\mathbf{Z}$-代数结构是否交换；
但我们将看到陈类位于其中心。

\begin{remark}
\label{chow-remark-pullback-cohomology}
设 $(S, \delta)$ 如情形 \ref{chow-situation-setup}。
设 $f : Y' \to Y$ 是 $S$ 上局部有限型概形间的态射。
作为注 \ref{chow-remark-restriction-bivariant} 的一个特例，
存在典范 $\mathbf{Z}$-代数映射 $res : A^*(Y) \to A^*(Y')$。
文献中常把这个映射记为 $f^*$。
\end{remark}

\begin{lemma}
\label{chow-lemma-cap-c1-bivariant}
设 $(S, \delta)$ 如情形 \ref{chow-situation-setup}。
设 $X$ 在 $S$ 上局部有限型，$\mathcal{L}$ 是可逆
$\mathcal{O}_X$-模。对 $f : X' \to X$ 指定
$c_1(f^*\mathcal{L}) \cap - : \CH_k(X') \to \CH_{k - 1}(X')$
的规则是次数为 $1$ 的双变类。
\end{lemma}

\begin{proof}
这由引理 \ref{chow-lemma-factors}、\ref{chow-lemma-pushforward-cap-c1}、
\ref{chow-lemma-flat-pullback-cap-c1} 和
\ref{chow-lemma-gysin-commutes-cap-c1} 得出。
\end{proof}

\noindent
上述引理终于使我们能够作出如下定义。

\begin{definition}
\label{chow-definition-first-chern-class}
设 $(S, \delta)$ 如情形 \ref{chow-situation-setup}。设 $X$ 在 $S$ 上
局部有限型，$\mathcal{L}$ 是可逆 $\mathcal{O}_X$-模。
$\mathcal{L}$ 的{\it 第一陈类}
$c_1(\mathcal{L}) \in A^1(X)$ 是引理 \ref{chow-lemma-cap-c1-bivariant} 的双变类。
\end{definition}

\noindent
有限局部自由模的陈类将在第 \ref{chow-section-intersecting-chern-classes} 节中构造。下面证明
$c_1(\mathcal{L})$ 位于 $A^*(X)$ 的中心。

\begin{lemma}
\label{chow-lemma-c1-center}
设 $(S, \delta)$ 如情形 \ref{chow-situation-setup}。
设 $X$ 在 $S$ 上局部有限型，$\mathcal{L}$ 是可逆
$\mathcal{O}_X$-模。则
\begin{enumerate}
\item $c_1(\mathcal{L}) \in A^1(X)$ 位于 $A^*(X)$ 的中心；以及
\item 若 $f : X' \to X$ 局部有限型，且 $c \in A^*(X' \to X)$，
则 $c \circ c_1(\mathcal{L}) = c_1(f^*\mathcal{L}) \circ c$。
\end{enumerate}
\end{lemma}

\begin{proof}
当然，(2) 蕴含 (1)。设 $p : L \to X$ 如引理 \ref{chow-lemma-linebundle}，并设 $o : X \to L$ 为零截面。
以 $p' : L' \to X'$ 和 $o' : X' \to L'$ 表示它们的基变换。
由引理 \ref{chow-lemma-linebundle-formulae}，有
$$
p^*(c_1(\mathcal{L}) \cap \alpha) = - o_* \alpha
\quad\text{且}\quad
(p')^*(c_1(f^*\mathcal{L}) \cap \alpha') = - o'_* \alpha'
$$
由于 $c$ 是双变类，有
\begin{align*}
(p')^*(c \cap c_1(\mathcal{L}) \cap \alpha)
& =
c \cap p^*(c_1(\mathcal{L}) \cap \alpha) \\
& =
- c \cap o_* \alpha \\
& =
- o'_*(c \cap \alpha) \\
& =
(p')^*(c_1(f^*\mathcal{L}) \cap c \cap \alpha)
\end{align*}
由上述某个引理，$(p')^*$ 是单射，故
$c \cap c_1(\mathcal{L}) \cap \alpha =
c_1(f^*\mathcal{L}) \cap c \cap \alpha$。
经任意局部有限型基变换 $Y \to X$ 后，同一结论仍成立，
所以得到 (2) 所述双变类的等式。
\end{proof}

\begin{lemma}
\label{chow-lemma-vanish-above-dimension}
设 $(S, \delta)$ 如情形 \ref{chow-situation-setup}。
设 $X$ 是 $S$ 上有限型且具有充足可逆层的概形。假设
$d = \dim(X) < \infty$（这里确实指维数，而不是 $\delta$-维数）。
则对 $X$ 上任意可逆层
$\mathcal{L}_1, \ldots, \mathcal{L}_{d + 1}$，在 $A^{d + 1}(X)$ 中有
$c_1(\mathcal{L}_1) \circ \ldots \circ c_1(\mathcal{L}_{d + 1}) = 0$。
\end{lemma}

\begin{proof}
对 $d$ 作归纳。初始情形 $d = 0$ 成立，因为此时 $X$ 是有限个闭点所成的
集合，所以每个可逆模都平凡。假设 $d > 0$。由《除子》引理 \ref{divisors-lemma-quasi-projective-Noetherian-pic-effective-Cartier}，
对某些有效 Cartier 除子 $D, D' \subset X$，可以写成
$\mathcal{L}_{d + 1} \cong \mathcal{O}_X(D) \otimes
\mathcal{O}_X(D')^{\otimes -1}$。于是 $c_1(\mathcal{L}_{d + 1})$ 是
$c_1(\mathcal{O}_X(D))$ 与 $c_1(\mathcal{O}_X(D'))$ 之差，故可以假设
对某个有效 Cartier 除子有
$\mathcal{L}_{d + 1} = \mathcal{O}_X(D)$。

\medskip\noindent
以 $i : D \to X$ 表示包含态射，并以 $i^* \in A^1(D \to X)$
表示引理 \ref{chow-lemma-gysin-bivariant} 中由 Gysin 同态给出的双变类。
由引理 \ref{chow-lemma-support-cap-effective-Cartier}，在 $A^1(X)$ 中有
$i_* \circ i^* = c_1(\mathcal{L}_{d + 1})$
（左边的意义由引理 \ref{chow-lemma-push-proper-bivariant} 保证）。
由于 $c_1(\mathcal{L}_i)$ 同时与 $i_*$ 和 $i^*$ 交换
（根据双变类的定义），可得
$$
c_1(\mathcal{L}_1) \circ \ldots \circ c_1(\mathcal{L}_{d + 1}) =
i_* \circ c_1(\mathcal{L}_1) \circ \ldots \circ c_1(\mathcal{L}_d) \circ i^* =
i_* \circ c_1(\mathcal{L}_1|_D) \circ \ldots \circ c_1(\mathcal{L}_d|_D)
\circ i^*
$$
于是由对 $d$ 的归纳得到结论。事实上，因为 $X$ 的一般点都不在 $D$ 中，
有 $\dim(D) < d$。
\end{proof}

\begin{remark}
\label{chow-remark-ring-loc-classes}
设 $(S, \delta)$ 如情形 \ref{chow-situation-setup}。
设 $Z \to X$ 是 $S$ 上局部有限型概形间的闭浸入，且 $p \geq 0$。
在这种框架下定义
$$
A^{(p)}(Z \to X) =
\prod\nolimits_{i \leq p - 1} A^i(X) \times
\prod\nolimits_{i \geq p} A^i(Z \to X).
$$
于是 $A^{(p)}(Z \to X)$ 典范地带有分次代数结构。
事实上，更一般地存在乘法
$$
A^{(p)}(Z \to X) \times A^{(q)}(Z \to X)
\longrightarrow A^{(\max(p, q))}(Z \to X)
$$
为定义这些乘法，定义映射
\begin{align*}
A^i(Z \to X) \times A^j(X) & \to A^{i + j}(Z \to X) \\
A^i(X) \times A^j(Z \to X) & \to A^{i + j}(Z \to X) \\
A^i(Z \to X) \times A^j(Z \to X) & \to A^{i + j}(Z \to X)
\end{align*}
对第一个映射，使用双变类的复合。对第二个映射，使用限制
$A^i(X) \to A^i(Z)$（注 \ref{chow-remark-restriction-bivariant}）
以及复合 $A^i(Z) \times A^j(Z \to X) \to A^{i + j}(Z \to X)$。
对第三个映射，把 $(c, c')$ 映为 $res(c) \circ c'$，其中
$res : A^i(Z \to X) \to A^i(Z)$ 是限制映射
（见注 \ref{chow-remark-restriction-bivariant}）。
这里略去验证这些乘法在适当意义下具有结合性。
\end{remark}

\begin{remark}
\label{chow-remark-res-push}
设 $(S, \delta)$ 如情形 \ref{chow-situation-setup}。
设 $Z \to X$ 是 $S$ 上局部有限型概形间的闭浸入。
以 $res : A^p(Z \to X) \to A^p(Z)$ 表示注 \ref{chow-remark-restriction-bivariant} 的限制映射。
对 $c \in A^p(Z \to X)$ 和 $\alpha \in \CH_*(Z)$，有
$res(c) \cap \alpha = c \cap i_*\alpha$。事实上，
$res(c) \cap \alpha = c \cap \alpha$，而 $c$ 与固有推前的相容性给出
$(Z \to Z)_*(c \cap \alpha) = c \cap (Z \to X)_*\alpha$。
\end{remark}






\section{关于双变类的引理}
\label{chow-section-bivariant-II}

\noindent
本节证明若干关于双变类的初等结果。先给出判断一个运算能否下降到
有理等价类的判据。

\begin{lemma}
\label{chow-lemma-factors-through-rational-equivalence}
\begin{reference}
\cite[Theorem 17.1]{F} 的一个很弱的形式
\end{reference}
设 $(S, \delta)$ 如情形 \ref{chow-situation-setup}。
设 $f : X \to Y$ 是 $S$ 上局部有限型概形间的态射，并设
$p \in \mathbf{Z}$。假设给定如下规则：对每个局部有限型态射
$Y' \to Y$ 及每个 $k$，指定映射
$$
c \cap - : Z_k(Y') \longrightarrow \CH_{k - p}(X')
$$
其中 $Y' = X' \times_X Y$；并且每当
$\mathcal{L}'|_{D'} \cong \mathcal{O}_{D'}$ 时，该规则满足定义 \ref{chow-definition-bivariant-class} 的条件 (3)。则
$c \cap -$ 下降到有理等价类。
\end{lemma}

\begin{proof}
该陈述有意义：若给定定义 \ref{chow-definition-gysin-homomorphism} 中那样的三元组
$(\mathcal{L}, s, i : D \to X)$，且
$\mathcal{L}|_D \cong \mathcal{O}_D$，则运算 $i^*$
已在循环层次上定义；见注 \ref{chow-remark-gysin-on-cycles}。
设 $\alpha \in Z_k(X')$ 是有理等价于零的循环。需要证明
$c \cap \alpha = 0$。由引理 \ref{chow-lemma-rational-equivalence-family}，
存在循环 $\beta \in Z_{k + 1}(X' \times \mathbf{P}^1)$，使得
$\alpha = i_0^*\beta - i_\infty^*\beta$，其中
$i_0, i_\infty : X' \to X' \times \mathbf{P}^1$ 是位于
$0, \infty$ 上的 $X'$ 的闭浸入。它们是具有平凡法丛的有效
Cartier 除子的例子，所以
$c \cap i_0^*\beta = j_0^*(c \cap \beta)$ 且
$c \cap i_\infty^*\beta = j_\infty^*(c \cap \beta)$，其中
$j_0, j_\infty : Y' \to Y' \times \mathbf{P}^1$ 是与前述相同的闭浸入。
又因 $j_0^*(c \cap \beta) \sim_{rat} j_\infty^*(c \cap \beta)$
（由引理 \ref{chow-lemma-rational-equivalence-family}），结论成立。
\end{proof}

\begin{lemma}
\label{chow-lemma-bivariant-weaker}
\begin{reference}
\cite[Theorem 17.1]{F} 的弱形式
\end{reference}
设 $(S, \delta)$ 如情形 \ref{chow-situation-setup}。
设 $f : X \to Y$ 是 $S$ 上局部有限型概形间的态射，并设
$p \in \mathbf{Z}$。假设给定如下规则：对每个局部有限型态射
$Y' \to Y$ 及每个 $k$，指定映射
$$
c \cap - : \CH_k(Y') \longrightarrow \CH_{k - p}(X')
$$
其中 $Y' = X' \times_X Y$；该规则满足定义 \ref{chow-definition-bivariant-class} 的条件 (1)、(2)，并且每当
$\mathcal{L}'|_{D'} \cong \mathcal{O}_{D'}$ 时满足条件 (3)。
则 $c \cap -$ 是双变类。
\end{lemma}

\begin{proof}
设 $Y' \to Y$ 是局部有限型概形态射。设
$(\mathcal{L}', s', i' : D' \to Y')$ 如定义 \ref{chow-definition-gysin-homomorphism}，其到 $X'$ 的拉回为
$(\mathcal{N}', t', j' : E' \to X')$。
需要证明对所有 $\alpha' \in \CH_k(Y')$，都有
$c \cap (i')^*\alpha' = (j')^*(c \cap \alpha')$。

\medskip\noindent
以 $g : Y'' \to Y'$ 表示引理 \ref{chow-lemma-normal-cone-effective-Cartier} 中构造的相对维数为 $1$
的光滑态射，并同时采用其中的 $i'' : D'' \to Y''$ 和
$p : D'' \to D'$。（警告：$D''$ 不是 $D'$ 的完整逆像。）
以 $f : X'' \to X'$ 和 $E'' \subset X''$ 表示它们经
$X' \to Y'$ 的基变换。图示为
$$
\xymatrix{
& X'' \ar[rr] \ar'[d][dd]_h & & Y'' \ar[dd]^g \\
E'' \ar[rr] \ar[dd]_q \ar[ru]^{j''} & & D'' \ar[dd]^p \ar[ru]^{i''} & \\
& X' \ar'[r][rr] & & Y' \\
E' \ar[rr] \ar[ru]^{j'} & & D' \ar[ru]^{i'}
}
$$
由该引理所给性质，$\beta' = (i')^*\alpha'$ 是满足
$p^*\beta' = (i'')^*g^*\alpha'$ 的 $\CH_{k - 1}(D')$ 中唯一元素。
类似地，$\gamma' = (j')^*(c \cap \alpha')$ 是满足
$q^*\gamma' = (j'')^*h^*(c \cap \alpha')$ 的
$\CH_{k - 1 - p}(E')$ 中唯一元素。由对 $c$ 的假设，已知
$$
(j'')^*h^*(c \cap \alpha') =
(j'')^*(c \cap g^*\alpha') =
c \cap (i'')^*g^*\alpha'
$$
注意，引理陈述中所假设的条件 (3) 的修改版适用于 $i''$
及其基变换 $j''$。类似地，已知
$$
q^*(c \cap \beta') = c \cap p^*\beta'
$$
由上述唯一性，得到 $\gamma' = c \cap \beta'$。
\end{proof}

\noindent
下面给出双变类为零的判据。

\begin{lemma}
\label{chow-lemma-bivariant-zero}
设 $(S, \delta)$ 如情形 \ref{chow-situation-setup}。
设 $f : X \to Y$ 是 $S$ 上局部有限型概形间的态射。
设 $c \in A^p(X \to Y)$。对 $Y'' \to Y' \to Y$，置
$X'' = Y'' \times_Y X$ 及 $X' = Y' \times_Y X$。以下条件等价：
\begin{enumerate}
\item $c$ 为零；
\item 对每个在 $Y$ 上局部有限型的整概形 $Y'$，在
$\CH_*(X')$ 中有 $c \cap [Y'] = 0$；以及
\item 对每个在 $Y$ 上局部有限型的整概形 $Y'$，存在固有双有理态射
$Y'' \to Y'$，使得在 $\CH_*(X'')$ 中有 $c \cap [Y''] = 0$。
\end{enumerate}
\end{lemma}

\begin{proof}
蕴含关系 (1) $\Rightarrow$ (2) $\Rightarrow$ (3) 是显然的。
条件 (3) 蕴含 (2)，因为 $(Y'' \to Y')_*[Y''] = [Y']$，并且由于
$c$ 是双变类，有
$c \cap [Y'] = (X'' \to X')_*(c \cap [Y''])$。现在假设 (2)。
设 $Y' \to Y$ 局部有限型，并设 $\alpha \in \CH_k(Y')$。
写 $\alpha = \sum n_i [Y'_i]$，其中 $Y'_i \subset Y'$ 是一族局部有限、
$\delta$-维数为 $k$ 的整闭子概形。于是 $\alpha$ 是循环
$\alpha' = \sum n_i[Y'_i]$ 在 $Y'' = \coprod Y'_i$ 上沿固有态射
$Y'' \to Y'$ 的推前。由双变类的性质，只须证明
$c \cap \alpha' = 0$ 在 $\CH_{k - p}(X'')$ 中成立。
有 $\CH_{k - p}(X'') = \prod \CH_{k - p}(X'_i)$，其中
$X'_i = Y'_i \times_Y X$；这立即由定义得出。投影映射
$\CH_{k - p}(X'') \to \CH_{k - p}(X'_i)$ 由平坦拉回给出。
由于与 $c$ 作帽积和平坦拉回交换，只须证明
$c \cap [Y'_i]$ 在 $\CH_{k - p}(X'_i)$ 中为零，而这由假设成立。
\end{proof}

\begin{lemma}
\label{chow-lemma-disjoint-decomposition-bivariant}
设 $(S, \delta)$ 如情形 \ref{chow-situation-setup}。
设 $f : X \to Y$ 是 $S$ 上局部有限型概形间的态射。
假设由开闭子概形给出不交并分解
$X = \coprod_{i \in I} X_i$ 和 $Y = \coprod_{j \in J} Y_j$，
并且存在集合映射 $a : I \to J$，使得
$f(X_i) \subset Y_{a(i)}$。则
$$
A^p(X \to Y) = \prod\nolimits_{i \in I} A^p(X_i \to Y_{a(i)})
$$
\end{lemma}

\begin{proof}
设给定元素 $(c_i) \in \prod_i A^p(X_i \to Y_{a(i)})$。
对 $\beta \in \CH_k(Y)$，可把它映到 $\CH_{k - p}(X)$ 中的如下元素：
其到 $X_i$ 的限制为 $c_i \cap \beta|_{Y_{a(i)}}$。
这是可行的，因为 $\CH_{k - p}(X) = \prod_i \CH_{k - p}(X_i)$。
经任意局部有限型基变换 $Y' \to Y$ 后，同一构造仍可进行，
从而得到 $c \in A^p(X \to Y)$。于是得到从公式右边到左边的映射
$\Psi$。反之，给定 $c \in A^p(X \to Y)$ 和元素
$\beta_i \in \CH_k(Y_{a(i)})$，可以考虑 $\CH_{k - p}(X_i)$ 中的元素
$(c \cap (Y_{a(i)} \to Y)_*\beta_i)|_{X_i}$。
经任意局部有限型基变换 $Y' \to Y$ 后，同样可行，从而得到
$c_i \in A^p(X_i \to Y_{a(i)})$。于是得到从公式左边到右边的映射
$\Phi$。显然，$\Phi \circ \Psi = \text{id}$。
对反向复合，设 $c \in A^p(X \to Y)$ 且 $\beta \in \CH_k(Y)$。
设 $\Phi(c) = (c_i)$，并取 $j \in J$。因为 $c$ 与平坦拉回交换，有
$$
(c \cap \beta)|_{\coprod_{a(i) = j} X_i} =
c \cap \beta|_{Y_j}
$$
又因为 $c$ 与固有推前交换，有
$$
(\coprod\nolimits_{a(i) = j} X_i \to X)_*
((c \cap \beta)|_{\coprod_{a(i) = j} X_i})
=
c \cap (Y_j \to Y)_*\beta|_{Y_j}
$$
左边是 $X$ 上的循环：对满足 $a(i) = j$ 的 $i \in I$，
其到 $X_i$ 的限制为 $(c \cap \beta)|_{X_i}$，其余分量上为 $0$。
右边是 $X$ 上的循环：对满足 $a(i) = j$ 的 $i \in I$，
其到 $X_i$ 的限制为 $c_i \cap \beta|_{Y_j}$。
所以如所需有 $c \cap \beta = \Psi((c_i))$。
\end{proof}

\begin{remark}
\label{chow-remark-completion-bivariant}
设 $(S, \delta)$ 如情形 \ref{chow-situation-setup}。
设 $f : X \to Y$ 是 $S$ 上局部有限型概形间的态射。
设 $X = \coprod_{i \in I} X_i$ 和 $Y = \coprod_{j \in J} Y_j$
分别是 $X$、$Y$ 到各连通分支的分解
（因为 $X$、$Y$ 局部 Noether，连通分支是开的；见《拓扑学》引理 \ref{topology-lemma-locally-Noetherian-locally-connected} 和
《概形的性质》引理 \ref{properties-lemma-Noetherian-topology}）。
设 $a(i) \in J$ 是满足 $f(X_i) \subset Y_{a(i)}$ 的指标。
由引理 \ref{chow-lemma-disjoint-decomposition-bivariant}，
$A^p(X \to Y) = \prod A^p(X_i \to Y_{a(i)})$。
在此框架下，置
$$
A^*(X \to Y)^\wedge = \prod\nolimits_i A^*(X_i \to Y_{a(i)})
$$
较为方便。这个“完备化”的双变群是子集
$$
A^*(X \to Y)^\wedge \quad\subset\quad \prod\nolimits_{p \geq 0} A^p(X \to Y)
$$
它由如下元素组成：$c = (c_0, c_1, c_2, \ldots)$，且对每个连通分支
$X_i$，$c_p$ 在 $A^p(X_i \to Y_{a(i)})$ 中的像仅对有限多个 $p$
非零。若 $Y \to Z$ 是第二个态射，则复合
$A^*(X \to Y) \times A^*(Y \to Z) \to A^*(X \to Z)$
延拓为完备化之间的复合
$A^*(X \to Y)^\wedge \times A^*(Y \to Z)^\wedge \to A^*(X \to Z)^\wedge$。
有时把 $A^*(X)^\wedge = A^*(X \to X)^\wedge$ 称为 $X$ 的
{\it 完备双变上同调环}。
\end{remark}

\begin{lemma}
\label{chow-lemma-envelope-bivariant}
设 $(S, \delta)$ 如情形 \ref{chow-situation-setup}。
设 $f : X \to Y$ 是 $S$ 上局部有限型概形间的态射。
设 $g : Y' \to Y$ 是包络（定义 \ref{chow-definition-envelope}），
并记 $X' = Y' \times_Y X$。设 $p \in \mathbf{Z}$ 且
$c' \in A^p(X' \to Y')$。若两个限制
$$
res_1(c') = res_2(c') \in A^p(X' \times_X X' \to Y' \times_Y Y')
$$
相等（见证明），则存在唯一的 $c \in A^p(X \to Y)$，使其在
$A^p(X' \to Y')$ 中的限制满足 $res(c) = c'$。
\end{lemma}

\begin{proof}
有交换图
$$
\xymatrix{
X' \times_X X' \ar[d]^{f''} \ar@<1ex>[r]^-a \ar@<-1ex>[r]_-b &
X' \ar[d]^{f'} \ar[r]_h &
X \ar[d]^f \\
Y' \times_Y Y' \ar@<1ex>[r]^-p \ar@<-1ex>[r]_-q &
Y' \ar[r]^g &
Y
}
$$
元素 $res_1(c')$ 是 $c'$ 关于由态射 $a, f', p, f''$ 组成的 Cartesian
方块的限制（见注 \ref{chow-remark-restriction-bivariant}）；
元素 $res_2(c')$ 是 $c'$ 关于由态射 $b, f', q, f''$ 组成的 Cartesian
方块的限制。假设 $res_1(c') = res_2(c')$，并设
$\beta \in \CH_k(Y)$。由引理 \ref{chow-lemma-envelope}，可以找到
$\beta' \in \CH_k(Y')$，使 $g_*\beta' = \beta$。于是置
$$
c \cap \beta = h_*(c' \cap \beta')
$$
为证明这与 $\beta'$ 的选取无关，只须证明对
$\gamma \in \CH_k(Y' \times_Y Y')$，
$h_*(c' \cap (p_*\gamma - q_*\gamma))$ 为零。
因为 $c'$ 是双变类，有
$$
h_*(c' \cap (p_*\gamma - q_*\gamma)) =
h_*(a_*(c' \cap \gamma) - b_*(c' \cap \gamma)) = 0
$$
最后一个等式成立，是因为 $h \circ a = h \circ b$，从而
$h_* \circ a_* = h_* \circ b_*$。

\medskip\noindent
注意，要求 $res(c) = c'$ 以及双变类与固有推前相容，已经唯一地强制
了对 $c \cap \beta$ 的上述选取。

\medskip\noindent
当然，为定义双变类 $c$，还需要对任意局部有限型态射
$Y_1 \to Y$ 构造满足定义 \ref{chow-definition-bivariant-class}
所列条件的映射
$c \cap -: \CH_k(Y_1) \to \CH_{k + p}(Y_1 \times_Y X)$。
记 $Y'_1 = Y' \times_Y Y_1$、$X_1 = X \times_Y Y_1$。
由引理 \ref{chow-lemma-base-change-envelope}，态射
$Y'_1 \to Y_1$ 是包络。因此可以使用基变换后的图表
$$
\xymatrix{
X'_1 \times_{X_1} X'_1 \ar[d]^{f''_1} \ar@<1ex>[r]^-{a_1} \ar@<-1ex>[r]_-{b_1} &
X'_1 \ar[d]^{f'_1} \ar[r]_{h_1} &
X_1 \ar[d]^{f_1} \\
Y'_1 \times_{Y_1} Y'_1 \ar@<1ex>[r]^-{p_1} \ar@<-1ex>[r]_-{q_1} &
Y'_1 \ar[r]^{g_1} &
Y_1
}
$$
并用同样的论证得到如前良定义的映射
$c \cap - : \CH_k(Y_1) \to \CH_{k + p}(X_1)$。

\medskip\noindent
接下来，需要验证 $c$ 满足定义 \ref{chow-definition-bivariant-class} 的条件 (1)、(2)、(3)。
例如，假设 $t : Y_2 \to Y_1$ 是 $Y$ 上局部有限型概形间的固有态射。
像上面一样，以下标 ${}_1$、相应地\ ${}_2$ 表示第一幅图到
$Y_1$、相应地\ $Y_2$ 的基变换。以 $t' : Y'_2 \to Y'_1$、
$s : X_2 \to X_1$ 和 $s' : X'_2 \to X'_1$ 表示 $t$ 到
$Y'$、$X$ 和 $X'$ 的基变换。需要证明
$$
s_*(c \cap \beta_2) = c \cap t_*\beta_2
$$
对 $\beta_2 \in \CH_k(Y_2)$ 成立。选取
$\beta'_2 \in \CH_k(Y'_2)$，使 $g_{2, *}\beta'_2 = \beta_2$。
由于 $c'$ 是双变类，且图表
$$
\vcenter{
\xymatrix{
X'_2 \ar[d]_{s'} \ar[r]_{h_2} & X_2 \ar[d]^s \\
X'_1 \ar[r]^{h_1} & X_1
}
}
\quad\text{且}\quad
\vcenter{
\xymatrix{
X'_2 \ar[d]_{s'} \ar[r]_{f'_2} & Y'_2 \ar[d]^{t'} \\
X'_2 \ar[r]^{f'_1} & Y'_1
}
}
$$
都是 Cartesian 的，故有
$$
s_*(c \cap \beta_2) =
s_*(h_{2, *}(c' \cap \beta'_2)) =
h_{1, *}s'_*(c' \cap \beta'_2) =
h_{1, *}(c' \cap (t'_*\beta'_2))
$$
按构造，最后一个表达式计算 $c \cap t_*\beta_2$：类
$t'_*\beta'_2 \in \CH_k(Y'_1)$ 沿 $g_{1, *}$ 的像是
$t_*\beta_2$。这证明了条件 (1)。其余条件用同样方法证明，
这里略去详细论证。
\end{proof}















\section{射影空间丛公式}
\label{chow-section-projective-space-bundle-formula}

\noindent
设 $(S, \delta)$ 如情形 \ref{chow-situation-setup}。
设 $X$ 在 $S$ 上局部有限型。考虑秩为 $r$ 的有限局部自由
$\mathcal{O}_X$-模 $\mathcal{E}$。我们的约定是，{\it 与
$\mathcal{E}$ 关联的射影丛}是 $X$ 上的态射
$$
\xymatrix{
\mathbf{P}(\mathcal{E}) =
\underline{\text{Proj}}_X(\text{Sym}^*(\mathcal{E}))
\ar[r]^-\pi
& X
}
$$
其中 $\mathcal{O}_{\mathbf{P}(\mathcal{E})}(1)$ 规范化为满足
$\pi_*(\mathcal{O}_{\mathbf{P}(\mathcal{E})}(1)) = \mathcal{E}$。
特别地，存在满射
$\pi^*\mathcal{E} \to \mathcal{O}_{\mathbf{P}(\mathcal{E})}(1)$。
下文将非正式地说“设 $(\pi : P \to X, \mathcal{O}_P(1))$
是与 $\mathcal{E}$ 关联的射影丛”，以表示
$P = \mathbf{P}(\mathcal{E})$ 且
$\mathcal{O}_P(1) = \mathcal{O}_{\mathbf{P}(\mathcal{E})}(1)$ 的情形。

\begin{lemma}
\label{chow-lemma-cap-projective-bundle}
设 $(S, \delta)$ 如情形 \ref{chow-situation-setup}。
设 $X$ 在 $S$ 上局部有限型。设 $\mathcal{E}$ 是秩为 $r$ 的有限局部自由
$\mathcal{O}_X$-模 $\mathcal{E}$。设
$(\pi : P \to X, \mathcal{O}_P(1))$ 是与 $\mathcal{E}$ 关联的射影丛。
对任意 $\alpha \in \CH_k(X)$，元素
$$
\pi_*\left(
c_1(\mathcal{O}_P(1))^s \cap \pi^*\alpha
\right)
\in
\CH_{k + r - 1 - s}(X)
$$
在 $s < r - 1$ 时为 $0$，在 $s = r - 1$ 时等于 $\alpha$。
\end{lemma}

\begin{proof}
设 $Z \subset X$ 是 $\delta$-维数为 $k$ 的整闭子概形。
注意，$\pi^*[Z] = [\pi^{-1}(Z)]$，因为 $\pi^{-1}(Z)$ 是
$\delta$-维数为 $r - 1$ 的整概形。若 $s < r - 1$，则根据构造，
$c_1(\mathcal{O}_P(1))^s \cap \pi^*[Z]$ 由一个支撑在
$\pi^{-1}(Z)$ 上的 $(k + r - 1 - s)$-循环表示。
因此，由维数原因，该循环的推前为零。

\medskip\noindent
设 $s = r - 1$。由上述论证，对某个 $n \in \mathbf{Z}$ 有
$\pi_*(c_1(\mathcal{O}_P(1))^s \cap \pi^*\alpha) = n [Z]$。
需要证明 $n = 1$。由同样的维数原因，只须在用
$X \setminus T$ 替换 $X$ 后证明这一结果，其中
$T \subset Z$ 是真闭子集。设 $\xi$ 是 $Z$ 的一般点。
可以选取元素 $e_1, \ldots, e_{r - 1} \in \mathcal{E}_\xi$，
使它们构成 $\mathcal{E}_\xi$ 的某个基的一部分。
它们给出 $\mathcal{O}_P(1)|_{\pi^{-1}(Z)}$ 的有理截面
$s_1, \ldots, s_{r - 1}$；这些截面的公共零点集是态射
$\mathbf{P}(\mathcal{E}|_Z) \to Z$ 的一个有理截面之像的闭包，
并上一个支撑映到 $Z$ 的真闭子集 $T$ 的闭子集。
从 $X$ 中移除 $T$（相应地从 $P$ 中移除 $\pi^{-1}(T)$）后，
$s_1, \ldots, s_n$ 构成一列整体截面
$s_i \in \Gamma(\pi^{-1}(Z), \mathcal{O}_{\pi^{-1}(Z)}(1))$，
其公共零点集是某个截面 $Z \to \pi^{-1}(Z)$ 的像。
因此，反复应用引理 \ref{chow-lemma-geometric-cap}，依次得到
\begin{eqnarray*}
\pi^*[Z] & = & [\pi^{-1}(Z)] \\
c_1(\mathcal{O}_P(1)) \cap \pi^*[Z] & = & [Z(s_1)] \\
c_1(\mathcal{O}_P(1))^2 \cap \pi^*[Z] & = & [Z(s_1) \cap Z(s_2)] \\
\ldots & = & \ldots \\
c_1(\mathcal{O}_P(1))^{r - 1} \cap \pi^*[Z] & = &
[Z(s_1) \cap \ldots \cap Z(s_{r - 1})]
\end{eqnarray*}
由于 $\pi$ 在 $Z$ 上的一个截面的像沿 $\pi$ 的推前显然为 $[Z]$，
在 $\alpha = [Z]$ 时便得到结论。这里略去验证这些论证对一般循环
$\alpha = \sum n_j [Z_j]$ 也给出结论。
\end{proof}

\begin{lemma}[射影空间丛公式]
\label{chow-lemma-chow-ring-projective-bundle}
设 $(S, \delta)$ 如情形 \ref{chow-situation-setup}。
设 $X$ 在 $S$ 上局部有限型。
设 $\mathcal{E}$ 是秩为 $r$ 的有限局部自由 $\mathcal{O}_X$-模
$\mathcal{E}$。设 $(\pi : P \to X, \mathcal{O}_P(1))$
是与 $\mathcal{E}$ 关联的射影丛。映射
$$
\bigoplus\nolimits_{i = 0}^{r - 1}
\CH_{k + i}(X)
\longrightarrow
\CH_{k + r - 1}(P),
$$
$$
(\alpha_0, \ldots, \alpha_{r-1})
\longmapsto
\pi^*\alpha_0 +
c_1(\mathcal{O}_P(1)) \cap \pi^*\alpha_1
+ \ldots +
c_1(\mathcal{O}_P(1))^{r - 1} \cap \pi^*\alpha_{r-1}
$$
是同构。
\end{lemma}

\begin{proof}
固定 $k \in \mathbf{Z}$。先证该映射为单射。
设 $(\alpha_0, \ldots, \alpha_{r - 1})$ 是左端的一个元素，
且映到零。由引理 \ref{chow-lemma-cap-projective-bundle} 可知
$$
0 = \pi_*(\pi^*\alpha_0 +
c_1(\mathcal{O}_P(1)) \cap \pi^*\alpha_1
+ \ldots +
c_1(\mathcal{O}_P(1))^{r - 1} \cap \pi^*\alpha_{r-1})
= \alpha_{r - 1}
$$
接着可知
$$
0 = \pi_*(c_1(\mathcal{O}_P(1)) \cap (\pi^*\alpha_0 +
c_1(\mathcal{O}_P(1)) \cap \pi^*\alpha_1
+ \ldots +
c_1(\mathcal{O}_P(1))^{r - 2} \cap \pi^*\alpha_{r - 2}))
= \alpha_{r - 2}
$$
依此类推。因此该映射为单射。

\medskip\noindent
还须证明该映射为满射。
设 $X_i$（$i \in I$）为 $X$ 的不可约分支。
则 $P_i = \mathbf{P}(\mathcal{E}|_{X_i})$（$i \in I$）
是 $P$ 的不可约分支。考虑交换图
$$
\xymatrix{
\coprod P_i \ar[d]_{\coprod \pi_i} \ar[r]_p & P \ar[d]^\pi \\
\coprod X_i \ar[r]^q & X
}
$$
注意 $p_*$ 为满射。若 $\beta \in \CH_k(\coprod X_i)$，则
$\pi^* q_* \beta = p_*(\coprod \pi_i)^* \beta$，参见
引理 \ref{chow-lemma-flat-pullback-proper-pushforward}。由
引理 \ref{chow-lemma-pushforward-cap-c1}，与
$c_1(\mathcal{O}(1))$ 作帽积时也有类似结论。
因此，若引理中的映射对每个态射 $\pi_i : P_i \to X_i$ 都是满射，
则它对 $\pi : P \to X$ 也是满射。故可设 $X$ 不可约。
于是 $\dim_\delta(X) < \infty$，特别地，可以对
$\dim_\delta(X)$ 作归纳。

\medskip\noindent
若 $\dim_\delta(X) < k$，结论显然成立。
设 $\alpha \in \CH_{k + r - 1}(P)$。
对任意局部闭子概形 $T \subset X$，记
$\gamma_T : \bigoplus \CH_{k + i}(T) \to \CH_{k + r - 1}(\pi^{-1}(T))$
为映射
$$
\gamma_T(\alpha_0, \ldots, \alpha_{r - 1})
= \pi^*\alpha_0 + \ldots +
c_1(\mathcal{O}_{\pi^{-1}(T)}(1))^{r - 1} \cap \pi^*\alpha_{r - 1}.
$$
设对某个非空开集 $U \subset X$ 有
$\alpha|_{\pi^{-1}(U)} = \gamma_U(\alpha_0, \ldots, \alpha_{r - 1})$。
此时可以选取提升 $\alpha'_i \in \CH_{k + i}(X)$，并由
引理 \ref{chow-lemma-restrict-to-open} 看出，
$\alpha - \gamma_X(\alpha'_0, \ldots, \alpha'_{r - 1})$
有理等价于 $P_Y = \mathbf{P}(\mathcal{E}|_Y)$ 上的一个 $k$-循环，
其中 $Y = X \setminus U$ 赋予既约闭子概形结构。
注意 $\dim_\delta(Y) < \dim_\delta(X)$。
由归纳假设，结论对 $P_Y \to Y$ 成立，从而对 $\alpha$ 也成立。
因此，可以用 $X$ 的任意非空开集替换 $X$。

\medskip\noindent
特别地，可以设 $\mathcal{E} \cong \mathcal{O}_X^{\oplus r}$。
此时 $\mathbf{P}(\mathcal{E}) = X \times \mathbf{P}^{r - 1}$。
采用分层
$$
\mathbf{P}^{r - 1} = \mathbf{A}^{r - 1}
\amalg \mathbf{A}^{r - 2}
\amalg \ldots
\amalg \mathbf{A}^0
$$
每一层的闭包都是一个 $\mathbf{P}^{r - 1 - i}$，它表示
$c_1(\mathcal{O}(1))^i \cap [\mathbf{P}^{r - 1}]$。
因此 $P$ 有类似的分层
$$
P = U^{r - 1} \amalg U^{r - 2} \amalg \ldots \amalg U^0
$$
设 $P^i$ 为 $U^i$ 的闭包，设 $\pi^i : P^i \to X$
为 $\pi$ 在 $P^i$ 上的限制。
设 $\alpha \in \CH_{k + r - 1}(P)$。由
引理 \ref{chow-lemma-pullback-affine-fibres-surjective}，
可对某个 $\alpha_0 \in \CH_k(X)$ 写成
$\alpha|_{U^{r - 1}} = \pi^*\alpha_0|_{U^{r - 1}}$。
因此差 $\alpha - \pi^*\alpha_0$ 是某个
$\alpha' \in \CH_{k + r - 1}(P^{r - 2})$ 的像。
再次应用引理 \ref{chow-lemma-pullback-affine-fibres-surjective}，
可对某个 $\alpha_1 \in \CH_{k + 1}(X)$ 写成
$\alpha'|_{U^{r - 2}} = (\pi^{r - 2})^*\alpha_1|_{U^{r - 2}}$。
由引理 \ref{chow-lemma-relative-effective-cartier}，
$(\pi^{r - 2})^*\alpha_1$ 的像表示
$c_1(\mathcal{O}_P(1)) \cap \pi^*\alpha_1$。
还可看出，
$\alpha - \pi^*\alpha_0 - c_1(\mathcal{O}_P(1)) \cap \pi^*\alpha_1$
是某个 $\alpha'' \in \CH_{k + r - 1}(P^{r - 3})$ 的像。
依此类推。
\end{proof}

\begin{lemma}
\label{chow-lemma-vectorbundle}
设 $(S, \delta)$ 如情形 \ref{chow-situation-setup}。
设 $X$ 在 $S$ 上局部有限型。
设 $\mathcal{E}$ 是 $X$ 上秩为 $r$ 的有限局部自由层。令
$$
p :
E = \underline{\Spec}(\text{Sym}^*(\mathcal{E}))
\longrightarrow
X
$$
为 $X$ 上相应的向量丛。则对所有 $k$，
$p^* : \CH_k(X) \to \CH_{k + r}(E)$ 都是同构。
\end{lemma}

\begin{proof}
（秩一向量丛的情形见引理 \ref{chow-lemma-linebundle}。）
满射性见引理 \ref{chow-lemma-pullback-affine-fibres-surjective}。
设 $(\pi  : P \to X, \mathcal{O}_P(1))$ 是与有限局部自由层
$\mathcal{E} \oplus \mathcal{O}_X$ 关联的射影空间丛。
设 $s \in \Gamma(P, \mathcal{O}_P(1))$ 对应于整体截面
$(0, 1) \in \Gamma(X, \mathcal{E} \oplus \mathcal{O}_X)$。
令 $D = Z(s) \subset P$。注意，
$(\pi|_D : D \to X , \mathcal{O}_P(1)|_D)$
是与 $\mathcal{E}$ 关联的射影空间丛。记
$\pi_D = \pi|_D$ 且 $\mathcal{O}_D(1) = \mathcal{O}_P(1)|_D$。
此外，$D$ 是 $P$ 上的有效 Cartier 除子。因此
$\mathcal{O}_P(D) = \mathcal{O}_P(1)$
（见《除子》引理 \ref{divisors-lemma-characterize-OD}）。
还有一个同构 $E \cong P \setminus D$。记
$j : E \to P$ 为相应的开浸入。
为证明单射性，使用如下事实：
$$
j^* :
\CH_{k + r}(P)
\longrightarrow
\CH_{k + r}(E)
$$
的核由支撑在有效 Cartier 除子 $D$ 上的循环组成，见
引理 \ref{chow-lemma-restrict-to-open}。所以，若 $p^*\alpha = 0$，
则对某个 $\beta \in \CH_{k + r}(D)$ 有
$\pi^*\alpha = i_*\beta$。由
引理 \ref{chow-lemma-chow-ring-projective-bundle}，可写成
$$
\beta = \pi_D^*\beta_0 +
\ldots + c_1(\mathcal{O}_D(1))^{r - 1} \cap \pi_D^* \beta_{r - 1}.
$$
其中 $\beta_i \in \CH_{k + i}(X)$。由引理 \ref{chow-lemma-relative-effective-cartier} 和
\ref{chow-lemma-pushforward-cap-c1}，这推出
$$
\pi^*\alpha = i_*\beta =
c_1(\mathcal{O}_P(1)) \cap \pi^*\beta_0 +
\ldots +
c_1(\mathcal{O}_D(1))^r \cap \pi^*\beta_{r - 1}.
$$
由于 $\mathcal{E} \oplus \mathcal{O}_X$ 的秩为 $r + 1$，
除非 $\alpha$ 以及所有 $\beta_i$ 都为零，否则这与
引理 \ref{chow-lemma-pushforward-cap-c1} 矛盾。
\end{proof}








\section{向量丛的陈类}
\label{chow-section-chern-classes-vector-bundles}

\noindent
可以利用射影空间丛公式，通过把 $c_1(\mathcal{O}(1))^r$
展开成较低次幂来定义秩为 $r$ 的向量丛的陈类，见公式
(\ref{chow-equation-chern-classes})。其中符号的缘由稍后说明。

\begin{definition}
\label{chow-definition-chern-classes}
设 $(S, \delta)$ 如情形 \ref{chow-situation-setup}。
设 $X$ 在 $S$ 上局部有限型。假设 $X$ 整且
$n = \dim_\delta(X)$。设 $\mathcal{E}$ 是 $X$ 上秩为 $r$
的有限局部自由层。设 $(\pi : P \to X, \mathcal{O}_P(1))$
是与 $\mathcal{E}$ 关联的射影空间丛。
\begin{enumerate}
\item 由引理 \ref{chow-lemma-chow-ring-projective-bundle}，存在元素
$c_i \in \CH_{n - i}(X)$（$i = 0, \ldots, r$），
使得 $c_0 = [X]$，且
\begin{equation}
\label{chow-equation-chern-classes}
\sum\nolimits_{i = 0}^r
(-1)^i c_1(\mathcal{O}_P(1))^i \cap \pi^*c_{r - i}
= 0.
\end{equation}
\item 采用上述记号，令
$c_i(\mathcal{E}) \cap [X] = c_i$，
视为 $\CH_{n - i}(X)$ 的元素。称它们为
{\it $X$ 上 $\mathcal{E}$ 的陈类}。
\item {\it $X$ 上 $\mathcal{E}$ 的全陈类}是组合
$$
c({\mathcal E}) \cap [X] =
c_0({\mathcal E}) \cap [X]
+ c_1({\mathcal E}) \cap [X] + \ldots
+ c_r({\mathcal E}) \cap [X]
$$
它是 $\CH_*(X) = \bigoplus_{k \in \mathbf{Z}} \CH_k(X)$
的元素。
\end{enumerate}
\end{definition}

\noindent
下面验证，当向量丛秩为 $1$ 时，这并未给出新的概念。

\begin{lemma}
\label{chow-lemma-first-chern-class}
设 $(S, \delta)$ 如情形 \ref{chow-situation-setup}。
设 $X$ 在 $S$ 上局部有限型。假设 $X$ 整且
$n = \dim_\delta(X)$。设 $\mathcal{L}$ 是可逆
$\mathcal{O}_X$-模。定义 \ref{chow-definition-chern-classes}
所定义的 $X$ 上 $\mathcal{L}$ 的第一陈类，等于
定义 \ref{chow-definition-divisor-invertible-sheaf}
赋予 $\mathcal{L}$ 的 Weil 除子。
\end{lemma}

\begin{proof}
本证明用 $c_1(\mathcal{L}) \cap [X]$ 表示
定义 \ref{chow-definition-divisor-invertible-sheaf} 中的构造。
由于 $\mathcal{L}$ 的秩为 $1$，根据我们的规范化有
$\mathbf{P}(\mathcal{L}) = X$ 且
$\mathcal{O}_{\mathbf{P}(\mathcal{L})}(1) = \mathcal{L}$。
因此 (\ref{chow-equation-chern-classes}) 成为
$$
(-1)^1 c_1(\mathcal{L}) \cap c_0 + (-1)^0 c_1 = 0
$$
由于 $c_0 = [X]$，可得所需的
$c_1 = c_1(\mathcal{L}) \cap [X]$。
\end{proof}

\begin{remark}
\label{chow-remark-equation-signs}
也可以把方程 \ref{chow-equation-chern-classes} 改写为
\begin{equation}
\label{chow-equation-signs}
\sum\nolimits_{i = 0}^r
c_1(\mathcal{O}_P(-1))^i \cap \pi^*c_{r - i}
= 0.
\end{equation}
不过，相比其对偶，我们认为使用重言商层
$\mathcal{O}_P(1)$ 更为方便。
\end{remark}




\section{与陈类相交}
\label{chow-section-intersecting-chern-classes}

\noindent
本节把 $X$ 上向量丛的陈类定义为 $X$ 上的双变类，参见
引理 \ref{chow-lemma-cap-cp-bivariant} 及该引理之后的讨论。
我们的构造遵循熟知的模式：先在素循环上定义运算，再求和。
在引理 \ref{chow-lemma-determine-intersections} 中，将证明结果由
相应射影丛上的通常公式所确定。随后证明，与陈类作帽积可下降到
有理等价，并与正常推前、平坦拉回以及有效 Cartier 除子嵌入的
Gysin 映射交换。也可以直接使用
引理 \ref{chow-lemma-determine-intersections} 中的刻画和
引理 \ref{chow-lemma-push-proper-bivariant}，从而避开这些引理；
希望看到具体推导的读者可参阅《空间的 Chow 群》引理 \ref{spaces-chow-lemma-segre-classes}。

\begin{definition}
\label{chow-definition-cap-chern-classes}
设 $(S, \delta)$ 如情形 \ref{chow-situation-setup}。
设 $X$ 在 $S$ 上局部有限型。设 $\mathcal{E}$ 是 $X$ 上秩为
$r$ 的有限局部自由层。对每个整数 $k$ 及任意
$0 \leq j \leq r$，定义运算
$$
c_j(\mathcal{E}) \cap - : Z_k(X) \to \CH_{k - j}(X)
$$
称为{\it 与 $\mathcal{E}$ 的第 $j$ 陈类相交}。
\begin{enumerate}
\item 给定一个 $\delta$-维数为 $k$ 的整闭子概形
$i : W \to X$，定义
$$
c_j(\mathcal{E}) \cap [W] = i_*(c_j({i^*\mathcal{E}}) \cap [W])
\in
\CH_{k - j}(X)
$$
其中 $c_j({i^*\mathcal{E}}) \cap [W]$ 如
定义 \ref{chow-definition-chern-classes} 中所定义。
\item 对一般的 $k$-循环 $\alpha = \sum n_i [W_i]$，令
$$
c_j(\mathcal{E}) \cap \alpha = \sum n_i c_j(\mathcal{E}) \cap [W_i]
$$
\end{enumerate}
\end{definition}

\noindent
若 $\mathcal{E}$ 的秩为 $1$，则这与此前的定义
（定义 \ref{chow-definition-cap-c1}）一致，这是
引理 \ref{chow-lemma-first-chern-class} 的结论。

\begin{lemma}
\label{chow-lemma-determine-intersections}
设 $(S, \delta)$ 如情形 \ref{chow-situation-setup}。
设 $X$ 在 $S$ 上局部有限型。设 $\mathcal{E}$ 是 $X$ 上秩为
$r$ 的有限局部自由层。设
$(\pi : P \to X, \mathcal{O}_P(1))$ 是与 $\mathcal{E}$ 关联的射影丛。
对 $\alpha \in Z_k(X)$，元素
$c_j(\mathcal{E}) \cap \alpha$ 是 $\CH_{k - j}(X)$ 中满足
$\alpha_0 = \alpha$ 以及
$$
\sum\nolimits_{i = 0}^r
(-1)^i c_1(\mathcal{O}_P(1))^i \cap
\pi^*(\alpha_{r - i}) = 0
$$
在 $P$ 的 Chow 群中成立的唯一元素 $\alpha_j$。
\end{lemma}

\begin{proof}
满足 $\alpha_0 = \alpha$ 及所示方程的
$\alpha_0, \ldots, \alpha_r$ 的唯一性，来自射影空间丛公式
引理 \ref{chow-lemma-chow-ring-projective-bundle}。
当 $\alpha = [W]$ 且 $W$ 是 $X$ 的整闭子概形时，该恒等式依定义成立。
对 $X$ 上一般的 $k$-循环 $\alpha$，写成
$\alpha = \sum n_a[W_a]$，其中 $n_a \not = 0$，且
$i_a : W_a \to X$ 是两两不同的整闭子概形。于是族
$\{W_a\}$ 在 $X$ 上局部有限。令
$P_a = \pi^{-1}(W_a) = \mathbf{P}(\mathcal{E}|_{W_a})$，
并记 $i'_a : P_a \to P$ 为相应的闭浸入。考虑纤维积图
$$
\xymatrix{
P' \ar@{=}[r] \ar[d]_{\pi'} &
\coprod P_a \ar[d]_{\coprod \pi_a} \ar[r]_{\coprod i'_a} &
P \ar[d]^\pi \\
X' \ar@{=}[r] &
\coprod W_a \ar[r]^{\coprod i_a} &
X
}
$$
态射 $p : X' \to X$ 是正常的。此外，
$\pi' : P' \to X'$ 连同可逆层
$\mathcal{O}_{P'}(1) = \coprod \mathcal{O}_{P_a}(1)$
是与 $\mathcal{E}' = p^*\mathcal{E}$ 关联的射影丛；该可逆层也正是
$\mathcal{O}_P(1)$ 的拉回。依定义，
$$
c_j(\mathcal{E}) \cap [\alpha]
=
\sum i_{a, *}(c_j(\mathcal{E}|_{W_a}) \cap [W_a]).
$$
记 $\beta_{a, j} = c_j(\mathcal{E}|_{W_a}) \cap [W_a]$，
它是 $\CH_{k - j}(W_a)$ 的元素。根据定义，对每个 $a$ 有
$$
\sum\nolimits_{i = 0}^r
(-1)^i c_1(\mathcal{O}_{P_a}(1))^i \cap \pi_a^*(\beta_{a, r - i}) = 0
$$
因而显然有
$$
\sum\nolimits_{i = 0}^r
(-1)^i c_1(\mathcal{O}_{P'}(1))^i \cap (\pi')^*(\beta_{r - i}) = 0
$$
其中 $\beta_j = \sum n_a\beta_{a, j} \in \CH_{k - j}(X')$。
记 $p' : P' \to P$ 为态射 $\coprod i'_a$。由
引理 \ref{chow-lemma-flat-pullback-proper-pushforward}，有
$\pi^*p_*\beta_j = p'_*(\pi')^*\beta_j$。
再由引理 \ref{chow-lemma-pushforward-cap-c1} 的投影公式，得到
$$
\sum\nolimits_{i = 0}^r
(-1)^i c_1(\mathcal{O}_P(1))^i \cap \pi^*(p_*\beta_j) = 0
$$
由于 $p_*\beta_j$ 表示 $c_j(\mathcal{E}) \cap \alpha$，故得证。
\end{proof}

\noindent
下文将始终使用陈类的这一刻画来证明更多性质。

\begin{lemma}
\label{chow-lemma-cap-chern-class-factors-rational-equivalence}
设 $(S, \delta)$ 如情形 \ref{chow-situation-setup}。
设 $X$ 在 $S$ 上局部有限型。设 $\mathcal{E}$ 是 $X$ 上秩为
$r$ 的有限局部自由层。若 $\alpha \sim_{rat} \beta$ 是 $X$ 上
有理等价的 $k$-循环，则在 $\CH_{k - j}(X)$ 中有
$c_j(\mathcal{E}) \cap \alpha = c_j(\mathcal{E}) \cap \beta$。
\end{lemma}

\begin{proof}
由引理 \ref{chow-lemma-determine-intersections}，元素
$\alpha_j = c_j(\mathcal{E}) \cap \alpha$（$j \geq 1$）与
$\beta_j = c_j(\mathcal{E}) \cap \beta$（$j \geq 1$），由与
$\mathcal{E}$ 关联的射影丛的 Chow 群中的{\it 同一个}方程唯一确定。
（这当然用到了平坦拉回与有理等价相容这一事实，见
引理 \ref{chow-lemma-flat-pullback-rational-equivalence}。）因此它们相等。
\end{proof}

\noindent
换言之，与有限局部自由层的陈类作帽积可下降到有理等价，从而给出映射
$$
c_j(\mathcal{E}) \cap - : \CH_k(X) \to \CH_{k - j}(X).
$$
下一步是证明陈类是双变类，见
定义 \ref{chow-definition-bivariant-class}。

\begin{lemma}
\label{chow-lemma-pushforward-cap-cj}
设 $(S, \delta)$ 如情形 \ref{chow-situation-setup}。
设 $X$、$Y$ 在 $S$ 上局部有限型。设 $\mathcal{E}$ 是 $X$ 上秩为
$r$ 的有限局部自由层。设 $p : X \to Y$ 是正常态射，
设 $\alpha$ 是 $X$ 上的 $k$-循环，并设 $\mathcal{E}$ 是 $Y$ 上
的有限局部自由层。则
$$
p_*(c_j(p^*\mathcal{E}) \cap \alpha) = c_j(\mathcal{E}) \cap p_*\alpha
$$
\end{lemma}

\begin{proof}
设 $(\pi : P \to Y, \mathcal{O}_P(1))$ 是与 $\mathcal{E}$ 关联的射影丛。
则 $P_X = X \times_Y P$ 是与 $p^*\mathcal{E}$ 关联的射影丛，且
$\mathcal{O}_{P_X}(1)$ 是 $\mathcal{O}_P(1)$ 的拉回。记
$\alpha_j = c_j(p^*\mathcal{E}) \cap \alpha$，于是
$\alpha_0 = \alpha$。由引理 \ref{chow-lemma-determine-intersections}，有
$$
\sum\nolimits_{i = 0}^r
(-1)^i c_1(\mathcal{O}_P(1))^i \cap
\pi_X^*(\alpha_{r - i}) = 0
$$
在 $P_X$ 的 Chow 群中成立。考虑纤维积图
$$
\xymatrix{
P_X \ar[r]_-{p'} \ar[d]_{\pi_X} & P \ar[d]^\pi \\
X \ar[r]^p & Y
}
$$
对上述等式施以正常推前 $p'_*$
（引理 \ref{chow-lemma-proper-pushforward-rational-equivalence}）。
利用引理 \ref{chow-lemma-pushforward-cap-c1} 和
\ref{chow-lemma-flat-pullback-proper-pushforward}，得到
$$
\sum\nolimits_{i = 0}^r
(-1)^i c_1(\mathcal{O}_P(1))^i \cap
\pi^*(p_*\alpha_{r - i}) = 0
$$
在 $P$ 的 Chow 群中成立。由
引理 \ref{chow-lemma-determine-intersections} 的刻画即得结论。
\end{proof}

\begin{lemma}
\label{chow-lemma-flat-pullback-cap-cj}
设 $(S, \delta)$ 如情形 \ref{chow-situation-setup}。
设 $X$、$Y$ 在 $S$ 上局部有限型。设 $\mathcal{E}$ 是 $Y$ 上秩为
$r$ 的有限局部自由层。设 $f : X \to Y$ 是相对维数为 $r$
的平坦态射，设 $\alpha$ 是 $Y$ 上的 $k$-循环。则
$$
f^*(c_j(\mathcal{E}) \cap \alpha) = c_j(f^*\mathcal{E}) \cap f^*\alpha
$$
\end{lemma}

\begin{proof}
记 $\alpha_j = c_j(\mathcal{E}) \cap \alpha$，于是
$\alpha_0 = \alpha$。由引理 \ref{chow-lemma-determine-intersections}，有
$$
\sum\nolimits_{i = 0}^r
(-1)^i c_1(\mathcal{O}_P(1))^i \cap
\pi^*(\alpha_{r - i}) = 0
$$
在与 $\mathcal{E}$ 关联的射影丛
$(\pi : P \to Y, \mathcal{O}_P(1))$ 的 Chow 群中成立。
考虑纤维积图
$$
\xymatrix{
P_X = \mathbf{P}(f^*\mathcal{E}) \ar[r]_-{f'} \ar[d]_{\pi_X} &
P \ar[d]^\pi \\
X \ar[r]^f & Y
}
$$
注意 $\mathcal{O}_{P_X}(1)$ 是 $\mathcal{O}_P(1)$ 的拉回。
对上述方程施以平坦拉回 $(f')^*$
（引理 \ref{chow-lemma-flat-pullback-rational-equivalence}）。由
引理 \ref{chow-lemma-flat-pullback-cap-c1} 和
\ref{chow-lemma-compose-flat-pullback} 可知
$$
\sum\nolimits_{i = 0}^r
(-1)^i c_1(\mathcal{O}_{P_X}(1))^i \cap
\pi_X^*(f^*\alpha_{r - i}) = 0
$$
在 $P_X$ 的 Chow 群中成立。由
引理 \ref{chow-lemma-determine-intersections} 的刻画即得结论。
\end{proof}

\begin{lemma}
\label{chow-lemma-cap-chern-class-commutes-with-gysin}
设 $(S, \delta)$ 如情形 \ref{chow-situation-setup}。
设 $X$ 在 $S$ 上局部有限型。设 $\mathcal{E}$ 是 $X$ 上秩为
$r$ 的有限局部自由层。设 $(\mathcal{L}, s, i : D \to X)$
如定义 \ref{chow-definition-gysin-homomorphism}。
则对所有 $\alpha \in \CH_k(X)$ 有
$c_j(\mathcal{E}|_D) \cap i^*\alpha = i^*(c_j(\mathcal{E}) \cap \alpha)$。
\end{lemma}

\begin{proof}
记 $\alpha_j = c_j(\mathcal{E}) \cap \alpha$，于是
$\alpha_0 = \alpha$。由引理 \ref{chow-lemma-determine-intersections}，有
$$
\sum\nolimits_{i = 0}^r
(-1)^i c_1(\mathcal{O}_P(1))^i \cap
\pi^*(\alpha_{r - i}) = 0
$$
在与 $\mathcal{E}$ 关联的射影丛
$(\pi : P \to X, \mathcal{O}_P(1))$ 的 Chow 群中成立。
考虑纤维积图
$$
\xymatrix{
P_D = \mathbf{P}(\mathcal{E}|_D) \ar[r]_-{i'} \ar[d]_{\pi_D} &
P \ar[d]^\pi \\
D \ar[r]^i & X
}
$$
注意 $\mathcal{O}_{P_D}(1)$ 是 $\mathcal{O}_P(1)$ 的拉回。
对上述方程施以 Gysin 映射 $(i')^*$
（引理 \ref{chow-lemma-gysin-factors}）。应用引理 \ref{chow-lemma-gysin-commutes-cap-c1} 和
\ref{chow-lemma-gysin-flat-pullback}，得到
$$
\sum\nolimits_{i = 0}^r
(-1)^i c_1(\mathcal{O}_{P_D}(1))^i \cap
\pi_D^*(i^*\alpha_{r - i}) = 0
$$
在 $P_D$ 的 Chow 群中成立。由
引理 \ref{chow-lemma-determine-intersections} 的刻画即得结论。
\end{proof}

\noindent
至此已有足够材料证明，与陈类作帽积定义了一个双变类。

\begin{lemma}
\label{chow-lemma-cap-cp-bivariant}
设 $(S, \delta)$ 如情形 \ref{chow-situation-setup}。
设 $X$ 在 $S$ 上局部有限型。设 $\mathcal{E}$ 是秩为 $r$ 的局部自由
$\mathcal{O}_X$-模，并设 $0 \leq p \leq r$。则把
$f : X' \to X$ 映为
$c_p(f^*\mathcal{E}) \cap - : \CH_k(X') \to \CH_{k - p}(X')$
的规则，是一个次数为 $p$ 的双变类。
\end{lemma}

\begin{proof}
这立即来自引理 \ref{chow-lemma-cap-chern-class-factors-rational-equivalence}、
\ref{chow-lemma-pushforward-cap-cj}、
\ref{chow-lemma-flat-pullback-cap-cj}、
\ref{chow-lemma-cap-chern-class-commutes-with-gysin}
以及定义 \ref{chow-definition-bivariant-class}。
\end{proof}

\noindent
借助这个引理，可以如下定义有限局部自由模的陈类。

\begin{definition}
\label{chow-definition-chern-classes-final}
设 $(S, \delta)$ 如情形 \ref{chow-situation-setup}。
设 $X$ 在 $S$ 上局部有限型。设 $\mathcal{E}$ 是秩为 $r$ 的局部自由
$\mathcal{O}_X$-模。对 $i = 0, \ldots, r$，$\mathcal{E}$ 的
{\it 第 $i$ 陈类}是引理 \ref{chow-lemma-cap-cp-bivariant}
所构造的次数为 $i$ 的双变类
$c_i(\mathcal{E}) \in A^i(X)$。$\mathcal{E}$ 的
{\it 全陈类}是形式和
$$
c(\mathcal{E}) = 
c_0(\mathcal{E}) + c_1(\mathcal{E}) + \ldots + c_r(\mathcal{E})
$$
并将它视为 $X$ 上的非齐次双变类。
\end{definition}

\noindent
由定义 \ref{chow-definition-cap-chern-classes} 后的注记，
若 $\mathcal{E}$ 可逆，则这个定义与
定义 \ref{chow-definition-first-chern-class} 一致。
下面说明陈类位于双变 Chow 上同调环 $A^*(X)$ 的中心。

\begin{lemma}
\label{chow-lemma-cap-commutative-chern}
设 $(S, \delta)$ 如情形 \ref{chow-situation-setup}。
设 $X$ 在 $S$ 上局部有限型。设 $\mathcal{E}$ 是秩为 $r$ 的局部自由
$\mathcal{O}_X$-模。则
\begin{enumerate}
\item $c_j(\mathcal{E}) \in A^j(X)$ 位于 $A^*(X)$ 的中心；
\item 若 $f : X' \to X$ 局部有限型且 $c \in A^*(X' \to X)$，则
$c \circ c_j(\mathcal{E}) = c_j(f^*\mathcal{E}) \circ c$。
\end{enumerate}
特别地，若 $\mathcal{F}$ 是 $X$ 上另一个秩为 $s$ 的局部自由
$\mathcal{O}_X$-模，则对所有 $\alpha \in \CH_k(X)$，
$$
c_i(\mathcal{E}) \cap c_j(\mathcal{F}) \cap \alpha
=
c_j(\mathcal{F}) \cap c_i(\mathcal{E}) \cap \alpha
$$
作为 $\CH_{k - i - j}(X)$ 的元素成立。
\end{lemma}

\begin{proof}
(2) 推出 (1) 是立即的。设 $\alpha \in \CH_k(X)$，并记
$\alpha_j = c_j(\mathcal{E}) \cap \alpha$，于是
$\alpha_0 = \alpha$。由引理 \ref{chow-lemma-determine-intersections}，有
$$
\sum\nolimits_{i = 0}^r
(-1)^i c_1(\mathcal{O}_P(1))^i \cap
\pi^*(\alpha_{r - i}) = 0
$$
在与 $\mathcal{E}$ 关联的射影丛
$(\pi : P \to Y, \mathcal{O}_P(1))$ 的 Chow 群中成立。
记 $\pi' : P' \to X'$ 为 $\pi$ 沿 $f$ 的基变换。利用
引理 \ref{chow-lemma-c1-center} 及双变类的性质，得到
\begin{align*}
0 & = c \cap \left(\sum\nolimits_{i = 0}^r
(-1)^i c_1(\mathcal{O}_P(1))^i \cap
\pi^*(\alpha_{r - i})\right) \\
& =
\sum\nolimits_{i = 0}^r
(-1)^i c_1(\mathcal{O}_{P'}(1))^i \cap
(\pi')^*(c \cap \alpha_{r - i})
\end{align*}
在 $P'$ 的 Chow 群中成立（略去计算）。因此，由
引理 \ref{chow-lemma-determine-intersections} 的刻画，
$c \cap \alpha_j$ 等于
$c_j(f^*\mathcal{E}) \cap (c \cap \alpha)$。引理得证。
\end{proof}

\begin{remark}
\label{chow-remark-extend-to-finite-locally-free}
设 $(S, \delta)$ 如情形 \ref{chow-situation-setup}。
设 $X$ 在 $S$ 上局部有限型。设 $\mathcal{E}$ 是有限局部自由
$\mathcal{O}_X$-模。即使 $\mathcal{E}$ 的秩不是常数，仍可定义
$\mathcal{E}$ 的陈类。
事实上，此时可写成
$$
X = X_0 \amalg X_1 \amalg X_2 \amalg \ldots
$$
其中 $X_r \subset X$ 是 $\mathcal{E}$ 秩为 $r$ 的开闭子空间。
由引理 \ref{chow-lemma-disjoint-decomposition-bivariant}，有
$A^p(X) = \prod A^p(X_r)$。因此，可以把 $c_p(\mathcal{E})$
定义为各 $A^p(X_r)$ 中的类 $c_p(\mathcal{E}|_{X_r})$ 之积。
明确地说，若 $X' \to X$ 是局部有限型态射，则通过拉回得到
$X'$ 的相应分解，并由定义有
$$
\CH_*(X') = \prod\nolimits_{r \geq 0} \CH_*(X'_r)
$$
此时 $c_p(\mathcal{E}) \in A^p(X)$ 是保持这些直积分解，
并在各因子上由已经定义的运算
$c_i(\mathcal{E}|_{X_r}) \cap -$ 作用的双变类。
注意，在这种情形下，可能有无穷多个 $p$ 使
$c_p(\mathcal{E})$ 非零。因此，全陈类是完备化双变上同调环中的元素
$$
c(\mathcal{E}) =
c_0(\mathcal{E}) + c_1(\mathcal{E}) + c_2(\mathcal{E}) + \ldots
\in A^*(X)^\wedge
$$
参见注 \ref{chow-remark-completion-bivariant}。
在这里，把 $\mathcal{E}$ 的“秩”定义为元素
$r(\mathcal{E}) \in A^0(X)$，即把
$(\alpha_r) \in \prod \CH_*(X'_r)$ 映到
$(r\alpha_r) \in \prod \CH_*(X'_r)$ 的双变运算。
仍然有 $c_p(\mathcal{E})$ 和 $r(\mathcal{E})$ 位于
$A^*(X)$ 的中心。
\end{remark}

\begin{remark}
\label{chow-remark-top-chern-class}
设 $(S, \delta)$ 如情形 \ref{chow-situation-setup}，设 $X$
在 $S$ 上局部有限型，并设 $\mathcal{E}$ 是有限局部自由
$\mathcal{O}_X$-模。一般地，如
注 \ref{chow-remark-extend-to-finite-locally-free} 写成
$X = \coprod X_r$。若仅有有限多个 $X_r$ 非空，则可令
$$
c_{top}(\mathcal{E}) = \sum\nolimits_r c_r(\mathcal{E}|_{X_r})
\in A^*(X) = \bigoplus A^*(X_r)
$$
其中等式来自引理 \ref{chow-lemma-disjoint-decomposition-bivariant}。
若有无穷多个 $X_r$ 非空，则用同一记号表示
$$
c_{top}(\mathcal{E}) = \prod c_r(\mathcal{E}|_{X_r})
\in \prod A^r(X_r) \subset A^*(X)^\wedge
$$
记号参见注 \ref{chow-remark-completion-bivariant}。
\end{remark}











\section{陈类之间的多项式关系}
\label{chow-section-relations-chern-classes}

\noindent
设 $(S, \delta)$ 如情形 \ref{chow-situation-setup}，设 $X$ 在 $S$
上局部有限型，并设 $\mathcal{E}_i$ 是 $X$ 上有限局部自由层的有限族。
由引理 \ref{chow-lemma-cap-commutative-chern}，陈类
$$
c_j(\mathcal{E}_i) \in A^*(X)
$$
生成 Chow 上同调代数 $A^*(X)$ 的一个交换（甚至中心的）
$\mathbf{Z}$-子代数。因此，可以说明这些陈类的多项式为零，
或者两个多项式相等，具体是什么意思。例如，等式
$c_1(\mathcal{E}_1)^5 + c_2(\mathcal{E}_2)c_3(\mathcal{E}_3) = 0$
表示：对所有局部有限型态射 $f : Y \to X$，运算
$$
\CH_k(Y) \longrightarrow \CH_{k - 5}(Y), \quad
\alpha \longmapsto
c_1(\mathcal{E}_1)^5 \cap \alpha +
c_2(\mathcal{E}_2) \cap c_3(\mathcal{E}_3) \cap \alpha
$$
都为零。由引理 \ref{chow-lemma-bivariant-zero}，这等价于如下要求：
对任意态射 $f : Y \to X$，其中 $Y$ 是 $S$ 上局部有限型的整概形，循环
$$
c_1(\mathcal{E}_1)^5 \cap [Y] +
c_2(\mathcal{E}_2) \cap c_3(\mathcal{E}_3) \cap [Y]
$$
在 $\CH_{\dim(Y) - 5}(Y)$ 中为零。

\medskip\noindent
一个具体例子是引理 \ref{chow-lemma-c1-cap-additive} 中证明的关系
$$
c_1(\mathcal{L} \otimes_{\mathcal{O}_X} \mathcal{N})
=
c_1(\mathcal{L}) + c_1(\mathcal{N})
$$
更一般地，下面说明任意局部自由层与可逆层张量时的情形。

\begin{lemma}
\label{chow-lemma-chern-classes-E-tensor-L}
设 $(S, \delta)$ 如情形 \ref{chow-situation-setup}。
设 $X$ 在 $S$ 上局部有限型。设 $\mathcal{E}$ 是 $X$ 上秩为 $r$
的有限局部自由层，并设 $\mathcal{L}$ 是 $X$ 上的可逆层。则在
$A^*(X)$ 中有
\begin{equation}
\label{chow-equation-twist}
c_i({\mathcal E} \otimes {\mathcal L})
=
\sum\nolimits_{j = 0}^i
\binom{r - i + j}{j} c_{i - j}({\mathcal E}) c_1({\mathcal L})^j
\end{equation}
\end{lemma}

\begin{proof}
该式应对任意三元组 $(X, \mathcal{E}, \mathcal{L})$ 成立。
特别地，它应在 $X$ 整时成立；由引理 \ref{chow-lemma-bivariant-zero}，
只须对这样的 $X$ 证明与 $[X]$ 作帽积后的等式。因此设 $X$ 整。
令 $(\pi : P \to X, \mathcal{O}_P(1))$，resp.\ $(\pi' : P' \to X, \mathcal{O}_{P'}(1))$
分别为与
$\mathcal{E}$，resp.\ $\mathcal{E} \otimes \mathcal{L}$
关联的射影空间丛。考虑规范态射
$$
\xymatrix{
P \ar[rd]_\pi \ar[rr]_g & & P' \ar[ld]^{\pi'} \\
& X &
}
$$
参见《概形的构造》引理 \ref{constructions-lemma-twisting-and-proj}。
它满足
$g^*\mathcal{O}_{P'}(1)
= \mathcal{O}_P(1) \otimes \pi^* {\mathcal L}$。
这意味着
$$
\sum\nolimits_{i = 0}^r
(-1)^i
(\xi + x)^i \cap \pi^*(c_{r - i}(\mathcal{E} \otimes \mathcal{L}) \cap [X])
=
0
$$
在 $\CH_*(P)$ 中成立，其中 $\xi$ 表示
$c_1(\mathcal{O}_P(1))$，而 $x$ 表示
$c_1(\pi^*\mathcal{L})$。由简单的代数运算，这等价于
$$
\sum\nolimits_{i = 0}^r
(-1)^i \xi^i \left(
\sum\nolimits_{j = i}^r
(-1)^{j - i}
\binom{j}{i}
x^{j - i} \cap
\pi^*(c_{r - j}(\mathcal{E} \otimes \mathcal{L}) \cap [X])
\right)
=
0
$$
与 Equation (\ref{chow-equation-chern-classes}) 比较，可得
$$
c_{r - i}(\mathcal{E}) \cap [X] =
\sum\nolimits_{j = i}^r
\binom{j}{i}
(-c_1(\mathcal{L}))^{j - i} \cap
c_{r - j}(\mathcal{E} \otimes \mathcal{L}) \cap [X]
$$
重新整理（消去负号并重新编号）便得到所需关系。
\end{proof}

\noindent
(\ref{chow-equation-twist}) 的一些具体情形为
\begin{align*}
c_1(\mathcal{E} \otimes \mathcal{L})
& =
c_1(\mathcal{E}) +
r c_1(\mathcal{L}) \\
c_2(\mathcal{E} \otimes \mathcal{L})
& =
c_2(\mathcal{E}) +
(r - 1) c_1(\mathcal{E}) c_1(\mathcal{L}) +
\binom{r}{2} c_1(\mathcal{L})^2 \\
c_3(\mathcal{E} \otimes \mathcal{L})
& =
c_3(\mathcal{E}) +
(r - 2) c_2(\mathcal{E})c_1(\mathcal{L}) +
\binom{r - 1}{2} c_1(\mathcal{E})c_1(\mathcal{L})^2 +
\binom{r}{3} c_1(\mathcal{L})^3
\end{align*}








\section{陈类的可加性}
\label{chow-section-additivity-chern-classes}

\noindent
所有预备引理都从最终结果立即推出。

\begin{lemma}
\label{chow-lemma-get-rid-of-trivial-subbundle}
设 $(S, \delta)$ 如情形 \ref{chow-situation-setup}。
设 $X$ 在 $S$ 上局部有限型。设 $\mathcal{E}$、$\mathcal{F}$ 是
$X$ 上秩分别为 $r$、$r - 1$ 的有限局部自由层，且嵌入短正合列
$$
0 \to \mathcal{O}_X \to \mathcal{E} \to \mathcal{F} \to 0
$$
则在 $A^*(X)$ 中有
$$
c_r(\mathcal{E}) = 0, \quad
c_j(\mathcal{E}) = c_j(\mathcal{F}), \quad j = 0, \ldots, r - 1
$$
\end{lemma}

\begin{proof}
由引理 \ref{chow-lemma-bivariant-zero}，只须证明：若 $X$ 整，则
$c_j(\mathcal{E}) \cap [X] = c_j(\mathcal{F}) \cap [X]$。
令 $(\pi : P \to X, \mathcal{O}_P(1))$，resp.\ $(\pi' : P' \to X, \mathcal{O}_{P'}(1))$
分别表示与 $\mathcal{E}$，resp.\ $\mathcal{F}$ 关联的射影空间丛。
满射 $\mathcal{E} \to \mathcal{F}$ 给出 $X$ 上的闭浸入
$$
i : P' \longrightarrow P
$$
此外，元素
$1 \in \Gamma(X, \mathcal{O}_X) \subset \Gamma(X, \mathcal{E})$
给出整体截面 $s \in \Gamma(P, \mathcal{O}_P(1))$，其零点集恰为
$P'$。因此 $P'$ 是 $P$ 上的有效 Cartier 除子，并满足
$\mathcal{O}_P(P') \cong \mathcal{O}_P(1)$。故由
引理 \ref{chow-lemma-relative-effective-cartier}，对 $X$ 上任意循环类
$\alpha$ 有
$$
c_1(\mathcal{O}_P(1)) \cap \pi^*\alpha = i_*((\pi')^*\alpha)
$$
由引理 \ref{chow-lemma-determine-intersections}，
$\alpha_j = c_j(\mathcal{F}) \cap [X]$（$j = 0, \ldots, r - 1$）满足
$$
\sum\nolimits_{j = 0}^{r - 1} (-1)^jc_1(\mathcal{O}_{P'}(1))^j
\cap (\pi')^*\alpha_j = 0
$$
把它推前到 $P$，并使用上述注记及
引理 \ref{chow-lemma-pushforward-cap-c1}，得到
$$
\sum\nolimits_{j = 0}^{r - 1}
(-1)^j c_1(\mathcal{O}_P(1))^{j + 1}
\cap \pi^*\alpha_j = 0
$$
由引理 \ref{chow-lemma-determine-intersections} 中的唯一性，可得
$c_r(\mathcal{E}) \cap [X] = 0$，并且
$c_j(\mathcal{E}) \cap [X] = \alpha_j = c_j(\mathcal{F}) \cap [X]$
对 $j = 0, \ldots, r - 1$ 成立。引理得证。
\end{proof}

\begin{lemma}
\label{chow-lemma-additivity-invertible-subsheaf}
设 $(S, \delta)$ 如情形 \ref{chow-situation-setup}。
设 $X$ 在 $S$ 上局部有限型。设 $\mathcal{E}$、$\mathcal{F}$ 是
$X$ 上秩分别为 $r$、$r - 1$ 的有限局部自由层，且嵌入短正合列
$$
0 \to \mathcal{L} \to \mathcal{E} \to \mathcal{F} \to 0
$$
其中 $\mathcal{L}$ 是可逆层。则在 $A^*(X)$ 中有
$$
c(\mathcal{E}) = c(\mathcal{L}) c(\mathcal{F})
$$
\end{lemma}

\begin{proof}
这个关系实际上就是
$c_i(\mathcal{E}) = c_i(\mathcal{F}) + c_1(\mathcal{L})c_{i - 1}(\mathcal{F})$。
由引理 \ref{chow-lemma-get-rid-of-trivial-subbundle}，对
$j = 0, \ldots, r$ 有
$c_j(\mathcal{E} \otimes \mathcal{L}^{\otimes -1})
= c_j(\mathcal{F} \otimes \mathcal{L}^{\otimes -1})$，其中约定
$c_r(\mathcal{F} \otimes \mathcal{L}^{-1}) = 0$。
应用引理 \ref{chow-lemma-chern-classes-E-tensor-L}，得到
$$
\sum_{j = 0}^i
\binom{r - i + j}{j} (-1)^j c_{i - j}({\mathcal E}) c_1({\mathcal L})^j
=
\sum_{j = 0}^i
\binom{r - 1 - i + j}{j} (-1)^j c_{i - j}({\mathcal F}) c_1({\mathcal L})^j
$$
令 $c_i(\mathcal{E}) = c_i(\mathcal{F}) + c_1(\mathcal{L})c_{i - 1}(\mathcal{F})$，
便给出这个方程的一个“解”。若证明这是唯一可能的解，引理即得。
这里略去验证。
\end{proof}

\begin{lemma}
\label{chow-lemma-additivity-chern-classes}
设 $(S, \delta)$ 如情形 \ref{chow-situation-setup}。
设 $X$ 是 $S$ 上局部有限型概形。假设 ${\mathcal E}$ 位于秩为
$r_i$ 的有限局部自由层 $\mathcal{E}_i$ 的正合列
$$
0
\to
{\mathcal E}_1
\to
{\mathcal E}
\to
{\mathcal E}_2
\to
0
$$
中。则全陈类在 $A^*(X)$ 中满足
$$
c({\mathcal E}) = c({\mathcal E}_1) c({\mathcal E}_2)
$$
\end{lemma}

\begin{proof}
由引理 \ref{chow-lemma-bivariant-zero}，可设 $X$ 整，并只须证明与
$[X]$ 作帽积后的恒等式。对 $r_1$ 归纳。当 $r_1 = 1$ 时，结论是
引理 \ref{chow-lemma-additivity-invertible-subsheaf}。设 $r_1 > 1$。
令 $(\pi : P \to X, \mathcal{O}_P(1))$ 表示与
$\mathcal{E}_1$ 关联的射影空间丛。注意
\begin{enumerate}
\item $\pi^* : \CH_*(X) \to \CH_*(P)$ 是单射；
\item $\pi^*\mathcal{E}_1$ 位于短正合列
$0 \to \mathcal{F} \to \pi^*\mathcal{E}_1 \to \mathcal{L} \to 0$
中，其中 $\mathcal{L}$ 可逆。
\end{enumerate}
第一点来自射影空间丛公式，第二点来自射影空间丛的定义。
（事实上 $\mathcal{L} = \mathcal{O}_P(1)$。）令
$Q = \pi^*\mathcal{E}/\mathcal{F}$，它位于正合列
$0 \to \mathcal{L} \to Q \to \pi^*\mathcal{E}_2 \to 0$ 中。
由归纳假设，
\begin{eqnarray*}
c(\pi^*\mathcal{E}) \cap [P]
& = &
c(\mathcal{F}) \cap c(\pi^*\mathcal{E}/\mathcal{F}) \cap [P] \\
& = &
c(\mathcal{F}) \cap c(\mathcal{L}) \cap c(\pi^*\mathcal{E}_2) \cap [P] \\
& = &
c(\pi^*\mathcal{E}_1) \cap c(\pi^*\mathcal{E}_2) \cap [P]
\end{eqnarray*}
由于 $[P] = \pi^*[X]$，由
引理 \ref{chow-lemma-flat-pullback-cap-cj} 即得结论。
\end{proof}

\begin{lemma}
\label{chow-lemma-chern-filter-by-linebundles}
设 $(S, \delta)$ 如情形 \ref{chow-situation-setup}。
设 $X$ 在 $S$ 上局部有限型。设 ${\mathcal L}_i$（$i = 1, \ldots, r$）
是 $X$ 上的可逆 $\mathcal{O}_X$-模。设 $\mathcal{E}$ 是赋有滤过
$$
0 = \mathcal{E}_0 \subset \mathcal{E}_1 \subset \mathcal{E}_2
\subset \ldots \subset \mathcal{E}_r = \mathcal{E}
$$
的局部自由秩 $\mathcal{O}_X$-模，且
$\mathcal{E}_i/\mathcal{E}_{i - 1} \cong \mathcal{L}_i$。
令 $c_1({\mathcal L}_i) = x_i$。则在 $A^*(X)$ 中有
$$
c(\mathcal{E})
=
\prod\nolimits_{i = 1}^r (1 + x_i)
$$
\end{lemma}

\begin{proof}
应用引理 \ref{chow-lemma-additivity-invertible-subsheaf} 并作归纳。
\end{proof}




\section{零维循环的次数}
\label{chow-section-degree-zero-cycles}

\noindent
先定义域上的正常概形上零维循环的次数。一种做法是如
引理 \ref{chow-lemma-spell-out-degree-zero-cycle} 直接定义，
再由引理 \ref{chow-lemma-curve-principal-divisor} 证明它良定义。
这里改用如下定义。

\begin{definition}
\label{chow-definition-degree-zero-cycle}
设 $k$ 是域（例 \ref{chow-example-field}），并设
$p : X \to \Spec(k)$ 正常。$X$ 上{\it 零维循环的次数}由正常推前
$$
p_* : \CH_0(X) \to \CH_0(\Spec(k))
$$
（引理 \ref{chow-lemma-proper-pushforward-rational-equivalence}）
与把 $[\Spec(k)]$ 映到 $1$ 的自然同构
$\CH_0(\Spec(k)) = \mathbf{Z}$ 复合给出。记作 $\deg(\alpha)$。
\end{definition}

\noindent
进一步把它写明如下。

\begin{lemma}
\label{chow-lemma-spell-out-degree-zero-cycle}
设 $k$ 是域，$X$ 在 $k$ 上正常，并设
$\alpha = \sum n_i[Z_i]$ 属于 $Z_0(X)$。则
$$
\deg(\alpha) = \sum n_i\deg(Z_i)
$$
其中 $\deg(Z_i)$ 是 $Z_i \to \Spec(k)$ 的次数，即
$\deg(Z_i) = \dim_k \Gamma(Z_i, \mathcal{O}_{Z_i})$。
\end{lemma}

\begin{proof}
这正是正常推前的定义
（定义 \ref{chow-definition-proper-pushforward}）。
\end{proof}

\noindent
下面把它与《代数簇》第 \ref{varieties-section-divisors-curves} 节
所定义的域上 $1$-维正常概形上的向量丛次数联系起来。

\begin{lemma}
\label{chow-lemma-degree-vector-bundle}
设 $k$ 是域。设 $X$ 是 $k$ 上维数 $\leq 1$ 的正常概形。
设 $\mathcal{E}$ 是常秩有限局部自由 $\mathcal{O}_X$-模。则
$$
\deg(\mathcal{E}) = \deg(c_1(\mathcal{E}) \cap [X]_1)
$$
其中左端由《代数簇》定义 \ref{varieties-definition-degree-invertible-sheaf} 定义。
\end{lemma}

\begin{proof}
设 $C_i \subset X$（$i = 1, \ldots, t$）是赋予诱导既约概形结构的
$1$ 维不可约分支，并设 $m_i$ 是 $C_i$ 在 $X$ 中的重数。则
$[X]_1 = \sum m_i[C_i]$，且 $c_1(\mathcal{E}) \cap [X]_1$
是循环 $m_i c_1(\mathcal{E}|_{C_i}) \cap [C_i]$ 的推前之和。
由《代数簇》引理 \ref{varieties-lemma-degree-in-terms-of-components}，
$\mathcal{E}$ 的次数也有类似分解，故只须证明 $X$ 是 $k$ 上正常曲线的情形。

\medskip\noindent
假设 $X$ 是 $k$ 上正常曲线。由《除子》引理 \ref{divisors-lemma-filter-after-modification}，存在修改
$f : X' \to X$，使 $f^*\mathcal{E}$ 有一个滤过，其逐次商是可逆
$\mathcal{O}_{X'}$-模。由于 $f_*[X']_1 = [X]_1$，由
引理 \ref{chow-lemma-pushforward-cap-cj} 可得
$$
\deg(c_1(\mathcal{E}) \cap [X]_1) = \deg(c_1(f^*\mathcal{E}) \cap [X']_1)
$$
由《代数簇》引理 \ref{varieties-lemma-degree-birational-pullback}，
次数也满足类似关系，因此归约到 $\mathcal{E}$ 有逐次商为可逆
$\mathcal{O}_X$-模的滤过这一情形。此时利用次数的可加性
（《代数簇》引理 \ref{varieties-lemma-degree-additive}）以及第一陈类的可加性
（引理 \ref{chow-lemma-additivity-chern-classes}），可归约到下一段的情形。

\medskip\noindent
假设 $X$ 是 $k$ 上正常曲线，且 $\mathcal{E}$ 是可逆
$\mathcal{O}_X$-模。由《除子》引理 \ref{divisors-lemma-quasi-projective-Noetherian-pic-effective-Cartier}，
$\mathcal{E}$ 同构于
$\mathcal{O}_X(D) \otimes \mathcal{O}_X(D')^{\otimes -1}$，
其中 $D, D'$ 是 $X$ 上的有效 Cartier 除子（这里还用到 $X$ 射影；
例如参见《代数簇》引理 \ref{varieties-lemma-dim-1-proper-projective}）。
利用可逆层张量积下次数的可加性
（《代数簇》引理 \ref{varieties-lemma-degree-tensor-product}）
以及可逆层张量积下 $c_1$ 的可加性
（引理 \ref{chow-lemma-c1-cap-additive} 或 \ref{chow-lemma-chern-classes-E-tensor-L}），
归约到 $\mathcal{E} = \mathcal{O}_X(D)$ 的情形。此时左端给出
$\deg(D)$（《代数簇》引理 \ref{varieties-lemma-degree-effective-Cartier-divisor}），而由
引理 \ref{chow-lemma-geometric-cap}，右端给出 $\deg([D]_0)$。
根据定义，
$$
[D]_0 = \sum\nolimits_{x \in D}
\text{length}_{\mathcal{O}_{X, x}}(\mathcal{O}_{D, x}) [x] =
\sum\nolimits_{x \in D}
\text{length}_{\mathcal{O}_{D, x}}(\mathcal{O}_{D, x}) [x]
$$
因此
$$
\deg([D]_0) = \sum\nolimits_{x \in D}
\text{length}_{\mathcal{O}_{D, x}}(\mathcal{O}_{D, x}) [\kappa(x) : k] =
\dim_k \Gamma(D, \mathcal{O}_D) = \deg(D)
$$
倒数第二个等式由《交换代数》引理 \ref{algebra-lemma-pushdown-module} 以及 $D$ 为仿射这一事实得到。
\end{proof}

\noindent
最后，把以上内容与《代数簇》第 \ref{varieties-section-num} 节所定义的数值交数联系起来。

\begin{lemma}
\label{chow-lemma-degrees-and-numerical-intersections}
设 $k$ 是域，$X$ 是 $k$ 上正常概形，$Z \subset X$ 是维数为
$d$ 的闭子概形，并设 $\mathcal{L}_1, \ldots, \mathcal{L}_d$
是可逆 $\mathcal{O}_X$-模。则
$$
(\mathcal{L}_1 \cdots \mathcal{L}_d \cdot Z) =
\deg(
c_1(\mathcal{L}_1) \cap \ldots \cap c_1(\mathcal{L}_d) \cap [Z]_d)
$$
其中左端由《代数簇》定义 \ref{varieties-definition-intersection-number} 定义。特别地，
若 $\mathcal{L}$ 是充足可逆 $\mathcal{O}_X$-模，则
$$
\deg_\mathcal{L}(Z) = \deg(c_1(\mathcal{L})^d \cap [Z]_d)
$$
\end{lemma}

\begin{proof}
对 $d$ 归纳。若 $d = 0$，结论由《代数簇》引理 \ref{varieties-lemma-chi-tensor-finite} 成立。设 $d > 0$。

\medskip\noindent
设 $Z_i \subset Z$（$i = 1, \ldots, t$）是赋予诱导既约概形结构的
$d$ 维不可约分支，并设 $m_i$ 是 $Z_i$ 在 $Z$ 中的重数。则
$[Z]_d = \sum m_i[Z_i]$，且
$c_1(\mathcal{L}_1) \cap \ldots \cap c_1(\mathcal{L}_d) \cap [Z]_d$
是循环
$m_i c_1(\mathcal{L}_1) \cap \ldots \cap c_1(\mathcal{L}_d) \cap [Z_i]$
之和。由《代数簇》引理 \ref{varieties-lemma-numerical-polynomial-leading-term}，
$(\mathcal{L}_1 \cdots \mathcal{L}_d \cdot Z)$ 也有类似分解，
故只须证明 $Z = X$ 是 $k$ 上 $d$ 维正常簇的情形。

\medskip\noindent
由 Chow 引理，存在双有理正常态射 $f : Y \to X$，其中 $Y$
在 $k$ 上 H-射影。参见《概形上同调》引理 \ref{coherent-lemma-chow-Noetherian} 和注 \ref{coherent-remark-chow-Noetherian}。此时，由《代数簇》引理 \ref{varieties-lemma-intersection-number-and-pullback} 有
$$
(f^*\mathcal{L}_1 \cdots f^*\mathcal{L}_d \cdot Y) =
(\mathcal{L}_1 \cdots \mathcal{L}_d \cdot X)
$$
并由引理 \ref{chow-lemma-pushforward-cap-c1} 有
$$
f_*(c_1(f^*\mathcal{L}_1) \cap \ldots \cap c_1(f^*\mathcal{L}_d) \cap [Y]) =
c_1(\mathcal{L}_1) \cap \ldots \cap c_1(\mathcal{L}_d) \cap [X]
$$
因此可用 $Y$ 替换 $X$，并设 $X$ 在 $k$ 上射影。

\medskip\noindent
若 $X$ 是正常的 $d$ 维射影簇，则由《除子》引理 \ref{divisors-lemma-quasi-projective-Noetherian-pic-effective-Cartier}，
可对某些有效 Cartier 除子 $D, D' \subset X$ 写成
$\mathcal{L}_1 = \mathcal{O}_X(D) \otimes \mathcal{O}_X(D')^{\otimes -1}$。
利用等式两端的可加性（《代数簇》引理 \ref{varieties-lemma-intersection-number-additive} 和
引理 \ref{chow-lemma-c1-cap-additive}），可归约到
$\mathcal{L}_1 = \mathcal{O}_X(D)$，其中 $D$ 是有效 Cartier 除子。
由《代数簇》引理 \ref{varieties-lemma-numerical-intersection-effective-Cartier-divisor}，有
$$
(\mathcal{L}_1 \cdots \mathcal{L}_d \cdot X) =
(\mathcal{L}_2 \cdots \mathcal{L}_d \cdot D)
$$
而由引理 \ref{chow-lemma-geometric-cap}，有
$$
c_1(\mathcal{L}_1) \cap \ldots \cap c_1(\mathcal{L}_d) \cap [X] =
c_1(\mathcal{L}_2) \cap \ldots \cap c_1(\mathcal{L}_d) \cap [D]_{d - 1}
$$
归纳假设即给出结论。
\end{proof}








\section{给定余维数的循环}
\label{chow-section-cycles-codimension}

\noindent
在某些情形下，所有循环的阿贝尔群还有一个由余维数给出的分次。

\begin{lemma}
\label{chow-lemma-locally-equidimensional}
设 $(S, \delta)$ 如情形 \ref{chow-situation-setup}。
设 $X$ 在 $S$ 上局部有限型。如第 \ref{chow-section-setup} 节，记
$\delta = \delta_{X/S}$。下列条件等价：
\begin{enumerate}
\item 存在分解 $X = \coprod_{n \in \mathbf{Z}} X_n$ 为开闭子概形，
使得只要 $\xi \in X_n$ 是 $X_n$ 某个不可约分支的一般点，
就有 $\delta(\xi) = n$。
\item 对所有 $x \in X$，存在 $x$ 的开邻域 $U \subset X$ 和整数
$n$，使得只要 $\xi \in U$ 是 $U$ 某个不可约分支的一般点，
就有 $\delta(\xi) = n$。
\item 对所有 $x \in X$，存在整数 $n_x$，使得对任意包含 $x$ 的
$X$ 的不可约分支的一般点 $\xi$，都有 $\delta(\xi) = n_x$。
\end{enumerate}
若 $X$ 正规或 Cohen--Macaulay，则这些条件成立\footnote{事实上，
只须 $X$ 是 $(S_2)$。可与《局部上同调》引理 \ref{local-cohomology-lemma-catenary-S2-equidimensional} 比较。}。
\end{lemma}

\begin{proof}
(1) $\Rightarrow$ (2) $\Rightarrow$ (3) 显然。反之，若 (3) 成立，
令 $X_n = \{x \in X \mid n_x = n\}$，便得到 (1) 中的分解。
事实上，$X_n$ 是开的：给定 $x$，所有经过 $x$ 的 $X$ 的不可约分支之并，
减去所有不经过 $x$ 的 $X$ 的不可约分支之并，是 $x$ 的一个开邻域。
若 $X$ 正规，则 $X$ 是整概形的不交并
（《概形的性质》引理 \ref{properties-lemma-normal-locally-Noetherian}），
故这些性质成立。若 $X$ 是 Cohen--Macaulay，则
$\delta' : X \to \mathbf{Z}$，$x \mapsto -\dim(\mathcal{O}_{X, x})$
是 $X$ 上的维数函数（见例 \ref{chow-example-CM-irreducible}）。
由于 $\delta - \delta'$ 局部常值
（《拓扑学》引理 \ref{topology-lemma-dimension-function-unique}），
且对 $X$ 的每个一般点 $\xi$ 都有 $\delta'(\xi) = 0$，故 (2) 成立。
\end{proof}

\noindent
设 $(S, \delta)$ 如情形 \ref{chow-situation-setup}。设 $X$ 在 $S$
上局部有限型，并满足引理 \ref{chow-lemma-locally-equidimensional}
中的等价条件。对整闭子概形 $Z \subset X$，有 $Z$ 在 $X$ 中的
余维数 $\text{codim}(Z, X)$，见《拓扑学》定义 \ref{topology-definition-codimension}。定义一个{\it 余维数 $p$ 的循环}
为 $X$ 上满足
$n_Z \not = 0 \Rightarrow \text{codim}(Z, X) = p$ 的循环
$\alpha = \sum n_Z[Z]$。所有余维数 $p$ 的循环所成的阿贝尔群记作
$Z^p(X)$。令 $X = \coprod X_n$ 为
引理 \ref{chow-lemma-locally-equidimensional} (1) 中的分解。
回忆循环被定义为局部有限和，显然
$$
Z^p(X) = \prod\nolimits_n Z_{n - p}(X_n)
$$
此外，$\prod_p Z^p(X) = \prod_k Z_k(X)$。现在可以完全照此前的方式
定义 $X$ 上余维数 $p$ 循环的有理等价；事实上，对满足
引理 \ref{chow-lemma-locally-equidimensional} 的 $X$，可以从头重新发展
给定余维数循环的整套理论。不过，这里直接把{\it 余维数 $p$ 循环的
Chow 群}定义为
$$
\CH^p(X) = \prod\nolimits_n \CH_{n - p}(X_n)
$$
如前，有 $\prod_p \CH^p(X) = \prod_k \CH_k(X)$。
若 $X$ 拟紧，则公式中的积为有限积（因而也是直和），且有
$\bigoplus_p \CH^p(X) = \bigoplus_k \CH_k(X)$。若 $X$ 还为有限维，
则这些群中仅有限多个非零。

\medskip\noindent
上文对 Chow 群证明的许多构造与结果，都有关于 Chow 群
$\CH^*(X)$ 的自然对应版本。证明方法是如
引理 \ref{chow-lemma-locally-equidimensional} 把有关概形分解成“等维”部分，
再对上述积分解中的各因子应用已经证明的结果。列举如下：
\begin{enumerate}
\item 若 $f : X \to Y$ 是 $S$ 上局部有限型概形之间的平坦态射，
且 $X$、$Y$ 满足引理 \ref{chow-lemma-locally-equidimensional} 的等价条件，
则平坦拉回给出映射
$$
f^* : \CH^p(Y) \to \CH^p(X)
$$
\item 若 $f : X \to Y$ 是 $S$ 上局部有限型概形之间的态射，
且 $X$、$Y$ 满足引理 \ref{chow-lemma-locally-equidimensional} 的等价条件，
则称 $f$ 的{\it 余维数}为 $r \in \mathbf{Z}$，如果对所有满足
$f(Z) \subset W$ 的不可约分支对 $Z \subset X$、$W \subset Y$，都有
$\dim_\delta(W) - \dim_\delta(Z) = r$。
\item 若 $f : X \to Y$ 是 $S$ 上局部有限型概形之间的正常态射，
且 $X$、$Y$ 满足引理 \ref{chow-lemma-locally-equidimensional} 的等价条件，
并且 $f$ 的余维数为 $r$，则正常推前是映射
$$
f_* : \CH^p(X) \to \CH^{p + r}(Y)
$$
\item 若 $f : X \to Y$ 是 $S$ 上局部有限型概形之间的态射，
且 $X$、$Y$ 满足引理 \ref{chow-lemma-locally-equidimensional} 的等价条件，
$f$ 的余维数为 $r$，并且 $c \in A^q(X \to Y)$，则 $c$ 诱导映射
$$
c \cap -  : \CH^p(Y) \to \CH^{p + q - r}(X)
$$
\item 若 $X$ 是 $S$ 上局部有限型概形，满足
引理 \ref{chow-lemma-locally-equidimensional} 的等价条件，且
$\mathcal{L}$ 是可逆 $\mathcal{O}_X$-模，则
$$
c_1(\mathcal{L}) \cap - : \CH^p(X) \to \CH^{p + 1}(X)
$$
\item 若 $X$ 是 $S$ 上局部有限型概形，满足
引理 \ref{chow-lemma-locally-equidimensional} 的等价条件，且
$\mathcal{E}$ 是有限局部自由 $\mathcal{O}_X$-模，则
$$
c_i(\mathcal{E}) \cap - : \CH^p(X) \to \CH^{p + i}(X)
$$
\end{enumerate}
警告：一个态射具有余维数 $r$ 的性质在基变换下并不保持。

\begin{remark}
\label{chow-remark-fundamental-class}
设 $(S, \delta)$ 如情形 \ref{chow-situation-setup}。
设 $X$ 在 $S$ 上局部有限型，并满足
引理 \ref{chow-lemma-locally-equidimensional} 的等价条件。
设 $X = \coprod X_n$ 是开闭子概形的分解，使 $X_n$ 的每个不可约分支
都具有 $\delta$-维数 $n$。在此情形下，有时令
$$
[X] = \sum\nolimits_n [X_n]_n \in \CH^0(X)
$$
这个类是 Chow 理论中 $X$ 的一种“基本类”。
\end{remark}




\section{分裂原理}
\label{chow-section-splitting-principle}

\noindent
在我们的设定中，不太容易精确说明分裂原理的内容。下面给出一种表述。

\begin{lemma}
\label{chow-lemma-splitting-principle}
设 $(S, \delta)$ 如情形 \ref{chow-situation-setup}。设 $X$ 在 $S$
上局部有限型，并设 $\mathcal{E}_i$ 是秩为 $r_i$ 的局部自由
$\mathcal{O}_X$-模的有限族。存在相对维数为 $d$ 的射影平坦态射
$\pi : P \to X$，使得
\begin{enumerate}
\item 对任意态射 $f : Y \to X$，映射
$\pi_Y^* : \CH_*(Y) \to \CH_{* + d}(Y \times_X P)$ 是单射；
\item 每个 $\pi^*\mathcal{E}_i$ 都有一个滤过，其逐次商
$\mathcal{L}_{i, 1}, \ldots, \mathcal{L}_{i, r_i}$ 是可逆
${\mathcal O}_P$-模。
\end{enumerate}
此外，当 (1) 成立时，限制映射 $A^*(X) \to A^*(P)$
（注 \ref{chow-remark-pullback-cohomology}）是单射。
\end{lemma}

\begin{proof}
可以设对所有 $i$ 都有 $r_i \geq 1$。对 $\sum (r_i - 1)$ 归纳。
若这个整数为 $0$，则对所有 $i$，$\mathcal{E}_i$ 都可逆，取
$\pi = \text{id}_X$ 即得结论。否则，可选取满足 $r_i > 1$ 的
$i$，并考虑态射 $\pi_i : P_i = \mathbf{P}(\mathcal{E}_i) \to X$。
有有限局部自由 $\mathcal{O}_{P_i}$-模的短正合列
$$
0 \to \mathcal{F} \to \pi_i^*\mathcal{E}_i \to \mathcal{O}_{P_i}(1) \to 0
$$
其中三项的秩分别为 $r_i - 1$、$r_i$、$1$。注意，由射影丛公式
（引理 \ref{chow-lemma-chow-ring-projective-bundle}），在任意基变换后，
$\pi_i^*$ 在 Chow 群上仍为单射。对有限局部自由
$\mathcal{O}_{P_i}$-模 $\mathcal{F}$ 以及
$\pi_{i'}^*\mathcal{E}_{i'}$（$i' \not = i$）应用归纳假设，
得到具有引理所述性质的态射 $\pi : P \to P_i$。
于是复合 $\pi_i \circ \pi : P \to X$ 满足要求。略去一些细节。
\end{proof}

\begin{remark}
\label{chow-remark-the-proof-shows-more}
引理 \ref{chow-lemma-splitting-principle} 的证明还表明态射
$\pi : P \to X$ 具有下列性质：
\begin{enumerate}
\item $\pi$ 是与有限常秩局部自由模关联的射影空间丛的有限复合；
\item 对每个 $\alpha \in \CH_k(X)$，都有
$\alpha = \pi_*(\xi_1 \cap \ldots \cap \xi_d \cap \pi^*\alpha)$，
其中 $\xi_i$ 是某个可逆 $\mathcal{O}_P$-模的第一陈类。
\end{enumerate}
第二点由第一点及引理 \ref{chow-lemma-cap-projective-bundle} 得出。
以后可按需要在此添加更多说明。
\end{remark}

\noindent
设 $(S, \delta)$、$X$ 和 $\mathcal{E}_i$ 如
引理 \ref{chow-lemma-splitting-principle}。所谓{\it 分裂原理}，是指形式地写成
$$
c(\mathcal{E}_i) = \prod (1 + x_{i, j})
$$
符号 $x_{i, 1}, \ldots, x_{i, r_i}$ 称为 $\mathcal{E}_i$ 的
{\it 陈根}。换言之，$\mathcal{E}_i$ 的第 $p$ 陈类是陈根的第 $p$
初等对称函数。分裂原理的用处在于：要证明 $\mathcal{E}_i$ 的陈类之间
的多项式关系，只须证明陈根之间对应的关系。

\medskip\noindent
具体地，设 $\pi : P \to X$ 如引理 \ref{chow-lemma-splitting-principle}。
回忆有规范的 $\mathbf{Z}$-代数映射
$\pi^* : A^*(X) \to A^*(P)$，见
注 \ref{chow-remark-pullback-cohomology}。对 $X$ 上每个 $Y$，
$\pi_Y^*$ 在 Chow 群上的单射性推出映射
$\pi^* : A^*(X) \to A^*(P)$ 是单射（略去细节）。由
引理 \ref{chow-lemma-chern-filter-by-linebundles}，有
$$
\pi^*c(\mathcal{E}_i) = \prod (1 + c_1(\mathcal{L}_{i, j}))
$$
因此，可以把陈根 $x_{i, j}$ 视为元素
$c_1(\mathcal{L}_{i, j}) \in A^*(P)$，并把所示方程理解为：
对方程左端应用单射 $\pi^* : A^*(X) \to A^*(P)$ 后，
方程在 $A^*(P)$ 中成立。

\medskip\noindent
要看清这一方法，最好给出若干例子。

\begin{lemma}
\label{chow-lemma-chern-classes-dual}
在情形 \ref{chow-situation-setup} 中，设 $X$ 在 $S$ 上局部有限型。
设 $\mathcal{E}$ 是有限局部自由 $\mathcal{O}_X$-模，其对偶为
$\mathcal{E}^\vee$。则在 $A^i(X)$ 中有
$$
c_i(\mathcal{E}^\vee) = (-1)^i c_i(\mathcal{E})
$$
\end{lemma}

\begin{proof}
选取引理 \ref{chow-lemma-splitting-principle} 中的态射
$\pi : P \to X$。由于 $\pi^*$ 在任意基变换后都是单射，只须在拉回到
$P$ 后证明 $\mathcal{E}$ 与 $\mathcal{E}^\vee$ 的陈类之间的关系。
因此可以设存在可逆 $\mathcal{O}_X$-模
${\mathcal L}_i$（$i = 1, \ldots, r$）以及滤过
$$
0 = \mathcal{E}_0 \subset \mathcal{E}_1 \subset \mathcal{E}_2
\subset \ldots \subset \mathcal{E}_r = \mathcal{E}
$$
使得 $\mathcal{E}_i/\mathcal{E}_{i - 1} \cong \mathcal{L}_i$。
于是得到对偶滤过
$$
0 = \mathcal{E}_r^\perp \subset \mathcal{E}_1^\perp \subset \mathcal{E}_2^\perp
\subset \ldots \subset \mathcal{E}_0^\perp = \mathcal{E}^\vee
$$
并且 $\mathcal{E}_{i - 1}^\perp/\mathcal{E}_i^\perp \cong
\mathcal{L}_i^{\otimes -1}$。令 $x_i = c_1(\mathcal{L}_i)$。
由引理 \ref{chow-lemma-c1-cap-additive}，有
$c_1(\mathcal{L}_i^{\otimes -1}) = - x_i$。由
引理 \ref{chow-lemma-chern-filter-by-linebundles}，在 $A^*(X)$ 中有
$$
c(\mathcal{E}) = \prod\nolimits_{i = 1}^r (1 + x_i)
\quad\text{且}\quad
c(\mathcal{E}^\vee) = \prod\nolimits_{i = 1}^r (1 - x_i)
$$
结论由一个形式计算得出，这里略去。
\end{proof}

\begin{lemma}
\label{chow-lemma-chern-classes-tensor-product}
在情形 \ref{chow-situation-setup} 中，设 $X$ 在 $S$ 上局部有限型。
设 $\mathcal{E}$ 和 $\mathcal{F}$ 是秩分别为 $r$ 和 $s$ 的有限局部自由
$\mathcal{O}_X$-模。则在 $A^*(X)$ 中有
$$
c_1(\mathcal{E} \otimes \mathcal{F})
=
r c_1(\mathcal{F}) + s c_1(\mathcal{E})
$$
$$
c_2(\mathcal{E} \otimes \mathcal{F})
=
r c_2(\mathcal{F}) + s c_2(\mathcal{E}) +
{r \choose 2} c_1(\mathcal{F})^2 +
(rs - 1) c_1(\mathcal{F})c_1(\mathcal{E}) +
{s \choose 2} c_1(\mathcal{E})^2
$$
依此类推。
\end{lemma}

\begin{proof}
完全仿照引理 \ref{chow-lemma-chern-classes-dual} 的证明，可以设有可逆
$\mathcal{O}_X$-模 ${\mathcal L}_i$（$i = 1, \ldots, r$）、
${\mathcal N}_i$（$i = 1, \ldots, s$）以及滤过
$$
0 = \mathcal{E}_0 \subset \mathcal{E}_1 \subset \mathcal{E}_2
\subset \ldots \subset \mathcal{E}_r = \mathcal{E}
\quad\text{且}\quad
0 = \mathcal{F}_0 \subset \mathcal{F}_1 \subset \mathcal{F}_2
\subset \ldots \subset \mathcal{F}_s = \mathcal{F}
$$
使得 $\mathcal{E}_i/\mathcal{E}_{i - 1} \cong \mathcal{L}_i$，且
$\mathcal{F}_j/\mathcal{F}_{j - 1} \cong \mathcal{N}_j$。
按字典序排列二元组 $(i, j)$，得到滤过
$$
0 \subset \ldots \subset
\mathcal{E}_i \otimes \mathcal{F}_j
+
\mathcal{E}_{i - 1} \otimes \mathcal{F}
\subset \ldots \subset \mathcal{E} \otimes \mathcal{F}
$$
其逐次商为
$$
\mathcal{L}_1 \otimes \mathcal{N}_1,
\mathcal{L}_1 \otimes \mathcal{N}_2,
\ldots,
\mathcal{L}_1 \otimes \mathcal{N}_s,
\mathcal{L}_2 \otimes \mathcal{N}_1,
\ldots,
\mathcal{L}_r \otimes \mathcal{N}_s
$$
由引理 \ref{chow-lemma-chern-filter-by-linebundles}，在 $A^*(X)$ 中有
$$
c(\mathcal{E}) = \prod (1 + x_i),
\quad
c(\mathcal{F}) = \prod (1 + y_j),
\quad\text{且}\quad
c(\mathcal{E} \otimes \mathcal{F}) = \prod (1 + x_i + y_j),
$$
结论由一个形式计算得出，这里略去。
\end{proof}

\begin{remark}
\label{chow-remark-equalities-nonconstant-rank}
即使考虑秩可以不是常数的有限局部自由 $\mathcal{O}_X$-模，
上面证明的等式仍然成立。事实上，可以如下使用秩和陈类的多项式。
考虑分次多项式环 $\mathbf{Z}[r, c_1, c_2, c_3, \ldots]$，
其中 $r$ 的次数为 $0$，$c_i$ 的次数为 $i$。设
$$
P \in \mathbf{Z}[r, c_1, c_2, c_3, \ldots]
$$
是次数为 $p$ 的齐次多项式。则对 $X$ 上任意有限局部自由
$\mathcal{O}_X$-模 $\mathcal{E}$，可以考虑
$$
P(\mathcal{E}) =
P(r(\mathcal{E}), c_1(\mathcal{E}), c_2(\mathcal{E}), c_3(\mathcal{E}), \ldots)
\in A^p(X)
$$
记号与约定见注 \ref{chow-remark-extend-to-finite-locally-free}。
要证明这些多项式之间的关系（涉及多个有限局部自由模时亦然），
可以在 $X$ 上局部工作，并使用上述分裂原理。例如，断言
$$
c_2(\SheafHom_{\mathcal{O}_X}(\mathcal{E}, \mathcal{E})) =
P(\mathcal{E})
$$
其中 $P = 2rc_2 - (r - 1)c_1^2$。事实上，由于
$\SheafHom_{\mathcal{O}_X}(\mathcal{E}, \mathcal{E}) =
\mathcal{E} \otimes \mathcal{E}^\vee$，该式容易从上述引理 \ref{chow-lemma-chern-classes-dual} 和
\ref{chow-lemma-chern-classes-tensor-product} 得出：只须把 $X$ 如
注 \ref{chow-remark-extend-to-finite-locally-free} 分解成
$\mathcal{E}$ 的秩为常数的部分。
\end{remark}

\begin{example}
\label{chow-example-power-sum}
对每个 $p \geq 1$，存在唯一的次数为 $p$ 的齐次多项式
$P_p \in \mathbf{Z}[c_1, c_2, c_3, \ldots]$，使得对任意
$n \geq p$，在 $\mathbf{Z}[x_1, \ldots, x_n]$ 中有
$$
P_p(s_1, s_2, \ldots, s_p) = \sum x_i^p
$$
其中 $s_1, \ldots, s_p$ 是 $x_1, \ldots, x_n$ 的初等对称多项式，即
$$
s_i = \sum\nolimits_{1 \leq j_1 < \ldots < j_i \leq n}
x_{j_1}x_{j_2} \ldots x_{j_i}
$$
$P_p$ 的存在性来自熟知的事实：初等对称函数生成整数上的全部对称函数环。
刻画 $P_p \in \mathbf{Z}[c_1, c_2, c_3, \ldots]$ 的另一种方式是，
作为形式幂级数有
$$
\log(1 + c_1 + c_2 + c_3 + \ldots) =
\sum\nolimits_{p \geq 1} (-1)^{p - 1}\frac{P_p}{p}
$$
这可由写成 $1 + c_1 + c_2 + \ldots = \prod (1 + x_i)$ 并应用
对数函数的幂级数立即看出。展开左端，得到
\begin{align*}
& (c_1 + c_2 + \ldots) - (1/2)(c_1 + c_2 + \ldots)^2 +
(1/3)(c_1 + c_2 + \ldots)^3 - \ldots \\
& =
c_1 + (c_2 - (1/2)c_1^2) + (c_3 - c_1c_2 + (1/3)c_1^3) + \ldots
\end{align*}
由此得到
\begin{align*}
P_1 & = c_1, \\
P_2 & = c_1^2 - 2c_2, \\
P_3 & = c_1^3 - 3c_1c_2 + 3c_3, \\
P_4 & = c_1^4 - 4c_1^2c_2 + 4c_1c_3 + 2c_2^2 - 4c_4,
\end{align*}
依此类推。由于有限局部自由 $\mathcal{O}_X$-模 $\mathcal{E}$ 的陈类
是陈根 $x_i$ 的初等对称多项式，可见
$$
P_p(\mathcal{E}) = \sum x_i^p
$$
为方便起见，在 $\mathbf{Z}[r, c_1, c_2, c_3, \ldots]$ 中令
$P_0 = r$，使得作为双变类有 $P_0(\mathcal{E}) = r(\mathcal{E})$
（如注 \ref{chow-remark-extend-to-finite-locally-free} 和
\ref{chow-remark-equalities-nonconstant-rank}）。
\end{example}






\section{陈类与截面}
\label{chow-section-top-chern-class}

\noindent
本节很短，主要结果是：可以用一个正则截面的零点轨迹计算有限局部自由模的顶陈类。

\medskip\noindent
设 $(S, \delta)$ 如情形 \ref{chow-situation-setup}。设 $X$ 是 $S$
上局部有限型概形，$\mathcal{E}$ 是有限局部自由 $\mathcal{O}_X$-模，
且 $f : X' \to X$ 局部有限型。设
$$
s \in \Gamma(X', f^*\mathcal{E})
$$
是 $\mathcal{E}$ 到 $X'$ 的拉回的整体截面。设 $Z(s) \subset X'$
为 $s$ 的零点概形。更精确地，定义 $Z(s)$ 为如下闭子概形：
其拟凝聚理想层是映射
$s : f^*\mathcal{E}^\vee \to \mathcal{O}_{X'}$ 的像。

\begin{lemma}
\label{chow-lemma-top-chern-class}
在上述情形中，假设 $\dim_\delta(X') = n$，$f^*\mathcal{E}$ 的常秩为
$r$，$\dim_\delta(Z(s)) \leq n - r$，并且对每个满足
$\delta(\xi) = n - r$ 的一般点 $\xi \in Z(s)$，$Z(s)$ 在
$\mathcal{O}_{X', \xi}$ 中的理想由长度为 $r$ 的正则序列生成。则在
$\CH_*(X')$ 中有
$$
c_r(\mathcal{E}) \cap [X']_n = [Z(s)]_{n - r}
$$
\end{lemma}

\begin{proof}
由于 $c_r(\mathcal{E})$ 是双变类
（引理 \ref{chow-lemma-cap-cp-bivariant}），可以设 $X = X'$，
并须证明 $\CH_{n - r}(X)$ 中的等式
$c_r(\mathcal{E}) \cap [X]_n = [Z(s)]_{n - r}$。
对 $r \geq 0$ 归纳。（$r = 0$ 的情形显然。）当 $r = 1$ 时，
结论由引理 \ref{chow-lemma-geometric-cap} 给出。设 $r > 1$。

\medskip\noindent
设 $\pi : P \to X$ 是与 $\mathcal{E}$ 关联的射影空间丛，
并考虑短正合列
$$
0 \to \mathcal{E}' \to \pi^*\mathcal{E} \to \mathcal{O}_P(1) \to 0
$$
由射影空间丛公式（引理 \ref{chow-lemma-chow-ring-projective-bundle}），
只须在沿 $\pi$ 拉回后证明该等式。注意，
$\pi^{-1}Z(s) = Z(\pi^*s)$ 的 $\delta$-维数 $\leq n - 1$，
且在 $\delta$-维数为 $n - 1$ 的一般点处关于正则序列的假设，
由平坦拉回仍成立；见《交换代数》引理 \ref{algebra-lemma-flat-increases-depth}。设
$t \in \Gamma(P, \mathcal{O}_P(1))$ 是 $\pi^*s$ 的像。断言
$$
[Z(t)]_{n + r - 2} = c_1(\mathcal{O}_P(1)) \cap [P]_{n + r - 1}
$$
假设该断言成立，则如下完成证明。限制 $\pi^*s|_{Z(t)}$ 在
$\mathcal{O}_P(1)|_{Z(t)}$ 中映到零，故来自唯一元素
$s' \in \Gamma(Z(t), \mathcal{E}'|_{Z(t)})$。注意，作为 $P$
的闭子概形，有 $Z(s') = Z(\pi^*s)$。若
$\xi \in Z(s')$ 是满足 $\delta(\xi) = n - 1$ 的一般点，则
$Z(s')$ 在 $\mathcal{O}_{Z(t), \xi}$ 中的理想可由长度为
$r - 1$ 的正则序列生成：它由 $\mathcal{O}_{P, \xi}$ 中
$r - 1$ 个元素的像生成，而这 $r - 1$ 个元素连同
$Z(t)$ 在 $\mathcal{O}_{P, \xi}$ 中理想的一个生成元，构成
$\mathcal{O}_{P, \xi}$ 中长度为 $r$ 的正则序列。因此可对
$Z(t)$ 上的 $s'$ 应用归纳假设，得到
$c_{r - 1}(\mathcal{E}') \cap [Z(t)]_{n + r - 2} = [Z(s')]_{n - 1}$。
综合以上各式，得到
\begin{align*}
c_r(\pi^*\mathcal{E}) \cap [P]_{n + r - 1}
& =
c_{r - 1}(\mathcal{E}') \cap c_1(\mathcal{O}_P(1)) \cap [P]_{n + r - 1} \\
& =
c_{r - 1}(\mathcal{E}') \cap [Z(t)]_{n + r - 2} \\
& =
[Z(s')]_{n - 1} \\
& = [Z(\pi^*s)]_{n - 1}
\end{align*}
这正是所需结论。

\medskip\noindent
证明断言。它来自对已经使用过的
引理 \ref{chow-lemma-geometric-cap} 的一次应用。我们有
$\pi^{-1}(Z(s)) = Z(\pi^*s) \subset Z(t)$。另一方面，对
$x \in X$，若 $P_x \subset Z(t)$，则 $t|_{P_x} = 0$；这推出
$s$ 在纤维 $\mathcal{E} \otimes \kappa(x)$ 中为零，进而推出
$x \in Z(s)$。因此 $\dim_\delta(Z(t)) \leq n + (r - 1) - 1$。
最后，设 $\xi \in Z(t)$ 是满足
$\delta(\xi) = n + r - 2$ 的一般点。若 $\xi$ 不是 $P \to X$
某条纤维的一般点，则 $Z(t)$ 的一个局部方程在
$\mathcal{O}_{P, \xi}$ 中是非零因子，这是立即的（因为由
《交换代数》引理 \ref{algebra-lemma-grothendieck}，可在纤维上检验）。
若 $\xi$ 是某条纤维的一般点，则 $x = \pi(\xi) \in Z(s)$，且
$\delta(x) = n + r - 2 - (r - 1) = n - 1$。由于 $r > 1$，
这与 $\dim_\delta(Z(s)) \leq n - r$ 矛盾，所以这种情形不会发生。
\end{proof}

\begin{lemma}
\label{chow-lemma-easy-virtual-class}
设 $(S, \delta)$ 如情形 \ref{chow-situation-setup}。设 $X$ 是 $S$
上局部有限型概形，并设
$$
0 \to \mathcal{N}' \to \mathcal{N} \to \mathcal{E} \to 0
$$
是有限局部自由 $\mathcal{O}_X$-模的短正合列。考虑闭嵌入
$$
i :
N' = \underline{\Spec}_X(\text{Sym}((\mathcal{N}')^\vee))
\longrightarrow
N = \underline{\Spec}_X(\text{Sym}(\mathcal{N}^\vee))
$$
对 $\alpha \in \CH_k(X)$，有
$$
i_*(p')^*\alpha = p^*(c_{top}(\mathcal{E}) \cap \alpha)
$$
其中 $p' : N' \to X$ 和 $p : N \to X$ 是结构态射。
\end{lemma}

\begin{proof}
这里 $c_{top}(\mathcal{E})$ 是注 \ref{chow-remark-top-chern-class}
中定义的双变类。由其定义，为验证该公式，可以设 $\mathcal{E}$
为常秩。类似地，可设 $\mathcal{N}'$ 和 $\mathcal{N}$ 的常秩分别为
$r'$ 和 $r$，于是 $\mathcal{E}$ 的秩为 $r - r'$，且
$c_{top}(\mathcal{E}) = c_{r - r'}(\mathcal{E})$。注意，
$p^*\mathcal{E}$ 有规范截面
$$
s \in \Gamma(N, p^*\mathcal{E}) = \Gamma(X, p_*p^*\mathcal{E}) =
\Gamma(X, \mathcal{E} \otimes_{\mathcal{O}_X} \text{Sym}(\mathcal{N}^\vee)
\supset \Gamma(X, \SheafHom(\mathcal{N}, \mathcal{E}))
$$
它对应于引理陈述中的满射 $\mathcal{N} \to \mathcal{E}$。
该截面的零点概形恰为 $N' \subset N$。设 $Y \subset X$ 是
$\delta$-维数为 $n$ 的整闭子概形。于是
\begin{enumerate}
\item $p^*[Y] = [p^{-1}(Y)]$，因为 $p^{-1}(Y)$ 是
$\delta$-维数为 $n + r$ 的整概形；
\item $(p')^*[Y] = [(p')^{-1}(Y)]$，因为 $(p')^{-1}(Y)$ 是
$\delta$-维数为 $n + r'$ 的整概形；
\item $s$ 在 $p^{-1}Y$ 上的限制以 $(p')^{-1}Y$ 为零点概形，
且闭浸入 $(p')^{-1}Y \to p^{-1}Y$ 是正则浸入（局部由正则序列截出）。
\end{enumerate}
由引理 \ref{chow-lemma-top-chern-class}，得到
$$
(p')^*[Y] = c_{r - r'}(p^*\mathcal{E}) \cap p^*[Y]
\quad\text{于}\quad \CH_*(N)
$$
引理得证。
\end{proof}










\section{陈特征与张量积}
\label{chow-section-chern-classes-tensor}

\noindent
设 $(S, \delta)$ 如情形 \ref{chow-situation-setup}，设 $X$ 在 $S$
上局部有限型。若 $x_i$ 是 ${\mathcal E}$ 的陈根，则把有限局部自由
$\mathcal{O}_X$-模的{\it 陈特征}定义为形式表达式
$$
ch({\mathcal E}) = \sum\nolimits_{i=1}^r e^{x_i}
$$
把它写成陈类的多项式，得到
\begin{align*}
ch(\mathcal{E})
& =
r(\mathcal{E})
+
c_1(\mathcal{E}) +
\frac{1}{2}(c_1(\mathcal{E})^2 - 2c_2(\mathcal{E}))
+
\frac{1}{6}(c_1(\mathcal{E})^3 - 3c_1(\mathcal{E})c_2(\mathcal{E}) + 3c_3(\mathcal{E})) \\
& \quad\quad +
\frac{1}{24}(c_1(\mathcal{E})^4 - 4c_1(\mathcal{E})^2c_2(\mathcal{E}) + 4c_1(\mathcal{E})c_3(\mathcal{E}) + 2c_2(\mathcal{E})^2 - 4c_4(\mathcal{E}))
+
\ldots \\
& =
\sum\nolimits_{p = 0, 1, 2, \ldots} \frac{P_p(\mathcal{E})}{p!}
\end{align*}
其中 $P_p$ 是例 \ref{chow-example-power-sum} 中陈类的多项式。
上述表达式的次数 $p$ 分量是
$$
ch_p(\mathcal{E}) = \frac{P_p(\mathcal{E})}{p!} \in A^p(X) \otimes \mathbf{Q}
$$
系数是有理数是什么意思？这只是说，把
$ch_p(\mathcal{E})$ 视为 $A^p(X) \otimes \mathbf{Q}$ 的元素。

\begin{remark}
\label{chow-remark-extend-chern-character-to-finite-locally-free}
在上述讨论中，即使 $\mathcal{E}$ 的秩不是常数，也已经定义了
$\mathcal{E}$ 的陈特征分量
$ch_p(\mathcal{E}) \in A^p(X) \otimes \mathbf{Q}$。
见注 \ref{chow-remark-extend-to-finite-locally-free} 和
\ref{chow-remark-equalities-nonconstant-rank}。因此 $\mathcal{E}$ 的完整陈特征
是 $\prod_{p \geq 0} (A^p(X) \otimes \mathbf{Q})$ 的元素。若 $X$
拟紧且 $\dim(X) < \infty$（通常维数），则利用
引理 \ref{chow-lemma-vanish-above-dimension} 和分裂原理可证明
$ch(\mathcal{E}) \in A^*(X) \otimes \mathbf{Q}$。
\end{remark}

\begin{lemma}
\label{chow-lemma-chern-character-additive}
设 $(S, \delta)$ 如情形 \ref{chow-situation-setup}。设 $X$ 在 $S$
上局部有限型，并设
$
0 \to \mathcal{E}_1 \to \mathcal{E} \to \mathcal{E}_2 \to 0
$
是有限局部自由 $\mathcal{O}_X$-模的短正合列。则有等式
$$
ch(\mathcal{E}) = ch(\mathcal{E}_1) + ch(\mathcal{E}_2)
$$
更精确地，在 $A^p(X)$ 中有
$P_p(\mathcal{E}) = P_p(\mathcal{E}_1) + P_p(\mathcal{E}_2)$，
其中 $P_p$ 如例 \ref{chow-example-power-sum}。
\end{lemma}

\begin{proof}
只须证明更精确的陈述。由第 \ref{chow-section-splitting-principle} 节，
这是因为：若 $x_{1, i}$（$i = 1, \ldots, r_1$）与
$x_{2, i}$（$i = 1, \ldots, r_2$）分别是
$\mathcal{E}_1$ 与 $\mathcal{E}_2$ 的陈根，则
$x_{1, 1}, \ldots, x_{1, r_1}, x_{2, 1}, \ldots, x_{2, r_2}$
是 $\mathcal{E}$ 的陈根。因此，结论来自例 \ref{chow-example-power-sum} 中对 $P_p$ 的选取。
\end{proof}

\begin{lemma}
\label{chow-lemma-chern-character-multiplicative}
设 $(S, \delta)$ 如情形 \ref{chow-situation-setup}。设 $X$ 在 $S$
上局部有限型，并设 $\mathcal{E}_1$ 和 $\mathcal{E}_2$ 是有限局部自由
$\mathcal{O}_X$-模。则有等式
$$
ch(\mathcal{E}_1 \otimes_{\mathcal{O}_X} \mathcal{E}_2) =
ch(\mathcal{E}_1) ch(\mathcal{E}_2)
$$
更精确地，在 $A^p(X)$ 中有
$$
P_p(\mathcal{E}_1 \otimes_{\mathcal{O}_X} \mathcal{E}_2) =
\sum\nolimits_{p_1 + p_2 = p}
{p \choose p_1} P_{p_1}(\mathcal{E}_1) P_{p_2}(\mathcal{E}_2)
$$
其中 $P_p$ 如例 \ref{chow-example-power-sum}。
\end{lemma}

\begin{proof}
只须证明更精确的陈述。由第 \ref{chow-section-splitting-principle} 节，
这是因为：若 $x_{1, i}$（$i = 1, \ldots, r_1$）与
$x_{2, i}$（$i = 1, \ldots, r_2$）分别是
$\mathcal{E}_1$ 与 $\mathcal{E}_2$ 的陈根，则
$x_{1, i} + x_{2, j}$（$1 \leq i \leq r_1$，$1 \leq j \leq r_2$）
是 $\mathcal{E}_1 \otimes \mathcal{E}_2$ 的陈根。因此，结论来自
$(x_{1, i} + x_{2, j})^p$ 的二项式公式以及例 \ref{chow-example-power-sum} 中多项式 $P_p$ 的形式。
\end{proof}

\begin{lemma}
\label{chow-lemma-chern-character-dual}
在情形 \ref{chow-situation-setup} 中，设 $X$ 在 $S$ 上局部有限型。
设 $\mathcal{E}$ 是有限局部自由 $\mathcal{O}_X$-模，其对偶为
$\mathcal{E}^\vee$。则在 $A^i(X) \otimes \mathbf{Q}$ 中有
$ch_i(\mathcal{E}^\vee) = (-1)^i ch_i(\mathcal{E})$。
\end{lemma}

\begin{proof}
结论来自陈类的相应结果
（引理 \ref{chow-lemma-chern-classes-dual}）。
\end{proof}











\section{陈类与导出范畴}
\label{chow-section-pre-derived}

\noindent
本节在如下假设下定义概形 $X$ 的导出范畴中完美对象 $E$ 的全陈类：
$E$ 可在 $X$ 的一个包络上由有限局部自由模的有限复形表示。

\medskip\noindent
设 $(S, \delta)$ 如情形 \ref{chow-situation-setup}，设 $X$ 在 $S$
上局部有限型，并设
$$
\mathcal{E}^a \to \mathcal{E}^{a + 1} \to \ldots \to \mathcal{E}^b
$$
是常秩有限局部自由 $\mathcal{O}_X$-模的有界复形。
用公式
$$
c(\mathcal{E}^\bullet) = \prod\nolimits_{n = a, \ldots, b}
c(\mathcal{E}^n)^{(-1)^n} \in \prod\nolimits_{p \geq 0} A^p(X)
$$
定义{\it 复形的全陈类}。这里的逆是形式逆，即
$$
(1 + c_1 + c_2 + c_3 + \ldots)^{-1} =
1 - c_1 + c_1^2 - c_2 - c_1^3 + 2c_1 c_2 - c_3 + \ldots
$$
用 $c_p(\mathcal{E}^\bullet) \in A^p(X)$ 表示
$c(\mathcal{E}^\bullet)$ 的次数 $p$ 部分。类似地，用公式
$$
ch(\mathcal{E}^\bullet) = \sum\nolimits_{n = a, \ldots, b}
(-1)^n ch(\mathcal{E}^n) \in
\prod\nolimits_{p \geq 0} (A^p(X) \otimes \mathbf{Q})
$$
定义{\it 复形的陈特征}。用
$ch_p(\mathcal{E}^\bullet) \in A^p(X) \otimes \mathbf{Q}$ 表示
$ch(\mathcal{E}^\bullet)$ 的次数 $p$ 部分。最后，对例 \ref{chow-example-power-sum} 中的
$P_p \in \mathbf{Z}[r, c_1, c_2, c_3, \ldots]$，定义
$$
P_p(\mathcal{E}^\bullet) = \sum\nolimits_{n = a, \ldots, b}
(-1)^n P_p(\mathcal{E}^n)
$$
它属于 $A^p(X)$。如常，有
$ch_p(\mathcal{E}^\bullet) = (1/p!)P_p(\mathcal{E}^\bullet)$。
下一个引理表明，这些构造只依赖于复形在导出范畴中的像。

\begin{lemma}
\label{chow-lemma-pre-derived-chern-class}
设 $(S, \delta)$ 如情形 \ref{chow-situation-setup}。
设 $X$ 在 $S$ 上局部有限型。设 $E \in D(\mathcal{O}_X)$ 是一个对象，
并存在表示 $E$ 的有限局部自由 $\mathcal{O}_X$-模的局部有界复形
$\mathcal{E}^\bullet$。则上述构造的轻微推广
$$
c(\mathcal{E}^\bullet) \in \prod\nolimits_{p \geq 0} A^p(X),\quad
ch(\mathcal{E}^\bullet) \in
\prod\nolimits_{p \geq 0} A^p(X) \otimes \mathbf{Q},\quad
P_p(\mathcal{E}^\bullet) \in A^p(X)
$$
与复形 $\mathcal{E}^\bullet$ 的选择无关。
\end{lemma}

\begin{proof}
先对全陈类证明；另两种情形用同样的论证，只须以
引理 \ref{chow-lemma-chern-character-additive} 代替
引理 \ref{chow-lemma-additivity-chern-classes}。

\medskip\noindent
如注 \ref{chow-remark-extend-to-finite-locally-free}，为定义全陈类
$c(\mathcal{E}^\bullet)$，把 $X$ 分解为开闭子概形
$$
X = \coprod\nolimits_{i \in I} X_i
$$
使对所有 $n$、$i$，$\mathcal{E}^n$ 在 $X_i$ 上的秩为常数。
（这些秩是 $X$ 上的局部常值函数，所以可以这样做。）由于
$\mathcal{E}^\bullet$ 局部有界，对固定的 $i$，层
$\mathcal{E}^n|_{X_i}$ 中仅有限多个非零。因此可以如上定义
$$
c(\mathcal{E}^\bullet|_{X_i}) =
\prod\nolimits_n c(\mathcal{E}^n|_{X_i})^{(-1)^n}
\in \prod\nolimits_{p \geq 0} A^p(X_i)
$$
由引理 \ref{chow-lemma-disjoint-decomposition-bivariant}，有
$A^p(X) = \prod_i A^p(X_i)$。因此，对每个 $p \in \mathbf{Z}$，
存在唯一元素 $c_p(\mathcal{E}^\bullet) \in A^p(X)$，使得对所有
$i$，其在 $X_i$ 上的限制为 $c_p(\mathcal{E}^\bullet|_{X_i})$。

\medskip\noindent
设有限局部自由 $\mathcal{O}_X$-模的另一个局部有界复形
$\mathcal{F}^\bullet$ 也表示 $E$。设 $g : Y \to X$ 局部有限型且
$Y$ 整。由引理 \ref{chow-lemma-bivariant-zero}，只须证明
$c(g^*\mathcal{E}^\bullet) \cap [Y]$ 与
$c(g^*\mathcal{F}^\bullet) \cap [Y]$ 相同；甚至只须在把 $Y$
替换为 $Y$ 上整、正常且双有理的概形后证明。首先可得
$g^*\mathcal{E}^\bullet$ 和 $g^*\mathcal{F}^\bullet$ 是常秩有限局部自由
$\mathcal{O}_Y$-模的有界复形。接着，由《平坦性进阶》引理 \ref{flat-lemma-blowup-complex-integral}，可设对所有
$i \in \mathbf{Z}$，$H^i(Lg^*E)$ 都是 Tor 维数 $\leq 1$ 的完美对象。
这就归约到下一段的情形。

\medskip\noindent
假设 $X$ 整，$\mathcal{E}^\bullet$ 和 $\mathcal{F}^\bullet$ 是常秩有限局部自由模的
有界复形，且对所有 $i \in \mathbf{Z}$，$H^i(E)$ 都是 Tor 维数
$\leq 1$ 的完美 $\mathcal{O}_X$-模。须证明
$c(\mathcal{E}^\bullet) \cap [X]$ 与
$c(\mathcal{F}^\bullet) \cap [X]$ 相同。记
$d_\mathcal{E}^i : \mathcal{E}^i \to \mathcal{E}^{i + 1}$ 和
$d_\mathcal{F}^i : \mathcal{F}^i \to \mathcal{F}^{i + 1}$
为这些复形的微分。由《平坦性进阶》注 \ref{flat-remark-when-you-have-a-complex}，对所有 $i$，
$\Im(d_\mathcal{E}^i)$、$\Ker(d_\mathcal{E}^i)$、
$\Im(d_\mathcal{F}^i)$ 和 $\Ker(d_\mathcal{F}^i)$ 都是有限局部自由
$\mathcal{O}_X$-模。由可加性（引理 \ref{chow-lemma-additivity-chern-classes}），
$$
c(\mathcal{E}^\bullet) = \prod\nolimits_i
c(\Ker(d_\mathcal{E}^i))^{(-1)^i} c(\Im(d_\mathcal{E}^i))^{(-1)^i}
$$
$\mathcal{F}^\bullet$ 也类似。由于有短正合列
$$
0 \to \Im(d_\mathcal{E}^i) \to \Ker(d_\mathcal{E}^i) \to H^i(E) \to 0
\quad\text{且}\quad
0 \to \Im(d_\mathcal{F}^i) \to \Ker(d_\mathcal{F}^i) \to H^i(E) \to 0
$$
便归约到下一段陈述并解决的问题。

\medskip\noindent
假设 $X$ 整，且有两个短正合列
$$
0 \to \mathcal{E}' \to \mathcal{E} \to \mathcal{Q} \to 0
\quad\text{且}\quad
0 \to \mathcal{F}' \to \mathcal{F} \to \mathcal{Q} \to 0
$$
其中 $\mathcal{E}$、$\mathcal{E}'$、$\mathcal{F}$、$\mathcal{F}'$
有限局部自由。问题是证明
$c(\mathcal{E})c(\mathcal{E}')^{-1} \cap [X] =
c(\mathcal{F})c(\mathcal{F}')^{-1} \cap [X]$。为此，考虑定义
$\mathcal{G}$ 的短正合列
$$
0 \to \mathcal{G} \to \mathcal{E} \oplus \mathcal{F} \to \mathcal{Q} \to 0
$$
由于 $\mathcal{Q}$ 的 Tor 维数 $\leq 1$，$\mathcal{G}$ 有限局部自由。
追图表明，满射 $\mathcal{G} \to \mathcal{F}$ 的核同构地映到
$\mathcal{E}$ 中的 $\mathcal{E}'$，而满射
$\mathcal{G} \to \mathcal{E}$ 的核同构地映到 $\mathcal{F}$ 中的
$\mathcal{F}'$。（在仿射局部，这来自 Schanuel 引理或等价于该引理，
见《交换代数》引理 \ref{algebra-lemma-Schanuel}。）因此
$$
c(\mathcal{E})c(\mathcal{F}') = c(\mathcal{G}) =
c(\mathcal{F})c(\mathcal{E}')
$$
正是所需结论。
\end{proof}

\begin{lemma}
\label{chow-lemma-defined-by-envelope}
设 $(S, \delta)$ 如情形 \ref{chow-situation-setup}。
设 $X$ 在 $S$ 上局部有限型，并设 $E \in D(\mathcal{O}_X)$
是完美对象。假设存在一个包络
$f : Y \to X$（定义 \ref{chow-definition-envelope}），
使得 $Lf^*E$ 在 $D(\mathcal{O}_Y)$ 中同构于有限局部自由
$\mathcal{O}_Y$-模组成的局部有界复形 $\mathcal{E}^\bullet$。
则存在唯一的双变类
$c(E) \in \prod_{p \geq 0} A^p(X)$、
$ch(E) \in \prod_{p \geq 0} A^p(X) \otimes \mathbf{Q}$ 以及
$P_p(E) \in A^p(X)$；它们与 $f : Y \to X$
及 $\mathcal{E}^\bullet$ 的选择无关，并且这些类在
$Y$ 上的限制分别等于 $c(\mathcal{E}^\bullet)$、
$ch(\mathcal{E}^\bullet)$ 和 $P_p(\mathcal{E}^\bullet)$。
\end{lemma}

\begin{proof}
固定 $p \in \mathbf{Z}$。我们对陈类
$c_p(E) \in A^p(X)$ 证明此引理；其余情形的论证从略。

\medskip\noindent
设 $g : T \to X$ 是局部有限型态射，并且存在有限局部自由
$\mathcal{O}_T$-模组成的局部有界复形 $\mathcal{E}^\bullet$，
它在 $D(\mathcal{O}_T)$ 中表示 $Lg^*E$。由
引理 \ref{chow-lemma-pre-derived-chern-class}，双变类
$c_p(\mathcal{E}^\bullet) \in A^p(T)$ 与
$\mathcal{E}^\bullet$ 的选择无关。将此类记为
$c_p(Lg^*E) \in A^p(T)$。对任意进一步的局部有限型态射
$h : T' \to T$，令 $g' = g \circ h$，则
$L(g')^*E = L(g \circ h)^*E = Lh^*Lg^*E$ 在
$D(\mathcal{O}_{T'})$ 中由 $h^*\mathcal{E}^\bullet$ 表示。
因此 $c_p(L(g')^*E)$ 有定义，并且等于
$c_p(Lg^*E)$ 在 $T'$ 上的限制
（注 \ref{chow-remark-restriction-bivariant}；严格来说，这里需要应用
引理 \ref{chow-lemma-cap-cp-bivariant}）。

\medskip\noindent
取引理陈述中的 $f : Y \to X$ 和 $\mathcal{E}^\bullet$。
把引理 \ref{chow-lemma-envelope-bivariant} 应用于
$f : Y \to X$ 以及上一段构造的类
$c' = c_p(Lf^*E)$，便得到双变类 $c_p(E) \in A^p(X)$。
该引理的假设成立，因为上一段的讨论给出
$res_1(c') = c_p(Lg^*E) = res_2(c')$，其中
$g = f \circ p = f \circ q : Y \times_X Y \to X$。

\medskip\noindent
最后，假设 $f' : Y' \to X$ 是另一个包络，并且
$L(f')^*E$ 由有限局部自由 $\mathcal{O}_{Y'}$-模组成的有界复形表示。
于是 $c_p(Lf^*E)$ 与 $c_p(L(f')^*E)$ 在
$Y \times_X Y'$ 上的限制相等。由于
$Y \times_X Y' \to X$ 是包络（引理 \ref{chow-lemma-base-change-envelope} 和
\ref{chow-lemma-composition-envelope}），由
引理 \ref{chow-lemma-envelope-bivariant} 中的唯一性可知，
$c_p(E)$ 的这两个候选者相同。
\end{proof}

\begin{definition}
\label{chow-definition-defined-on-perfect}
设 $(S, \delta)$ 如情形 \ref{chow-situation-setup}。
设 $X$ 在 $S$ 上局部有限型，并设 $E \in D(\mathcal{O}_X)$
是完美对象。
\begin{enumerate}
\item 若存在包络 $f : Y \to X$，使得 $Lf^*E$ 在
$D(\mathcal{O}_Y)$ 中同构于有限局部自由
$\mathcal{O}_Y$-模组成的局部有界复形，则称
{\it $E$ 的陈类已定义}\footnote{若干判据见
引理 \ref{chow-lemma-chern-classes-defined}。}。
\item 若 $E$ 的陈类已定义，则应用
引理 \ref{chow-lemma-defined-by-envelope} 定义
$$
c(E) \in \prod\nolimits_{p \geq 0} A^p(X),\quad
ch(E) \in
\prod\nolimits_{p \geq 0} A^p(X) \otimes \mathbf{Q},\quad
P_p(E) \in A^p(X)
$$
\end{enumerate}
\end{definition}

\noindent
此定义适用于本章所设想的许多情形，但并非全部情形，见
引理 \ref{chow-lemma-chern-classes-defined}。对于定义中一般的
$E/X/(S,\delta)$，也许存在这些双变类的初等构造；我们尚不知道。

\begin{lemma}
\label{chow-lemma-chern-classes-defined}
设 $(S, \delta)$ 如情形 \ref{chow-situation-setup}。
设 $X$ 在 $S$ 上局部有限型，并设 $E \in D(\mathcal{O}_X)$
是完美对象。若下列条件之一成立，则 $E$ 的陈类已定义：
\begin{enumerate}
\item 存在包络 $f : Y \to X$，使得 $Lf^*E$ 在
$D(\mathcal{O}_Y)$ 中同构于有限局部自由
$\mathcal{O}_Y$-模组成的局部有界复形；
\item $E$ 可由有限局部自由 $\mathcal{O}_X$-模组成的有界复形表示；
\item $X$ 的不可约分支都是拟紧的；
\item $X$ 是拟紧的；
\item 存在 $S$ 上局部有限型概形的态射 $X \to X'$，
使得 $E$ 是 $X'$ 上某个陈类已定义的完美对象 $E'$ 的拉回；或者
\item 此处增补更多情形。
\end{enumerate}
\end{lemma}

\begin{proof}
条件 (1) 正是定义 \ref{chow-definition-defined-on-perfect} 的第 (1) 部分。
条件 (2) 蕴含 (1)。

\medskip\noindent
在 (3) 的情形，假设 $X$ 的不可约分支 $X_i$ 都是拟紧的。
把 $X_i$ 视为 $X$ 上的既约整闭子概形。态射
$\coprod X_i \to X$ 是包络。对每个 $i$，存在包络
$X'_i \to X_i$，使得 $X'_i$ 具有充足的可逆模族；见《态射进阶》命题 \ref{more-morphisms-proposition-envelope-with-resolution-property}。
注意 $f : Y = \coprod X'_i \to X$ 是包络；此处略去一个小细节。
由《概形的导出范畴》引理 \ref{perfect-lemma-resolution-property-ample-family}，
每个 $X'_i$ 都具有消解性质。因此 $D(\mathcal{O}_{X'_i})$
中的完美对象 $L(f|_{X'_i})^*E$ 可由有限局部自由
$\mathcal{O}_{X'_i}$-模组成的有界复形表示；见
《概形的导出范畴》引理 \ref{perfect-lemma-resolution-property-perfect-complex}。
这证明了 (3) 蕴含 (1)。

\medskip\noindent
第 (4) 部分蕴含 (3)。

\medskip\noindent
设 $g : X \to X'$ 和 $E'$ 如第 (5) 部分。则存在包络
$f' : Y' \to X'$，使得 $L(f')^*E'$ 由一个局部有界的
$\mathcal{O}_{Y'}$-模复形 $(\mathcal{E}')^\bullet$ 表示。
由引理 \ref{chow-lemma-base-change-envelope}，基变换
$f : Y \to X$ 是包络。此外，拉回复形
$\mathcal{E}^\bullet = g^*(\mathcal{E}')^\bullet$ 表示 $Lf^*E$，
从而 $E$ 的陈类已定义。
\end{proof}

\begin{lemma}
\label{chow-lemma-chern-classes-computed}
设 $(S, \delta)$ 如情形 \ref{chow-situation-setup}。
设 $X$ 在 $S$ 上局部有限型，并设 $E \in D(\mathcal{O}_X)$
是完美对象。假设 $E$ 的陈类已定义。对局部有限型态射
$g : W \to X$，若 $W$ 整，则存在交换图
$$
\xymatrix{
W' \ar[rd]_{g'} \ar[rr]_b & & W \ar[ld]^g \\
& X
}
$$
其中 $W'$ 整，$b : W' \to W$ 是真双有理态射，
$L(g')^*E$ 由常秩局部自由 $\mathcal{O}_{W'}$-模组成的有界复形
$\mathcal{E}^\bullet$ 表示，并且在 $A^p(W')$ 中有
$res(c_p(E)) = c_p(\mathcal{E}^\bullet)$。
\end{lemma}

\begin{proof}
选择包络 $f : Y \to X$，使得 $Lf^*E$ 在
$D(\mathcal{O}_Y)$ 中同构于有限局部自由
$\mathcal{O}_Y$-模组成的局部有界复形 $\mathcal{E}^\bullet$。
由引理 \ref{chow-lemma-base-change-envelope}，$f$ 的基变换
$Y \times_X W \to W$ 是包络。选择一点
$\xi \in Y \times_X W$，它映到 $W$ 的泛点且剩余域相同。
考虑以 $\xi$ 为泛点的整闭子概形
$W' \subset Y \times_X W$。投影
$Y \times_X W \to W$ 在 $W'$ 上的限制是真双有理态射
$b : W' \to W$。令 $g' = g \circ b$。最后，考虑拉回
$(W' \to Y)^*\mathcal{E}^\bullet$。这是 $W'$ 上有限局部自由模
组成的局部有界复形。由于 $W'$ 整，它是有界的，且各项具有常秩。
最后，按构造，$(W' \to Y)^*\mathcal{E}^\bullet$ 表示
$L(g')^*E$，而其第 $p$ 个陈类按构造给出 $c_p(E)$
经 $W' \to X$ 的限制。证明完毕。
\end{proof}

\begin{lemma}
\label{chow-lemma-commutative-chern-perfect}
设 $(S, \delta)$ 如情形 \ref{chow-situation-setup}。
设 $X$ 在 $S$ 上局部有限型。设 $E \in D(\mathcal{O}_X)$
是完美对象。若 $E$ 的陈类已定义，则
\begin{enumerate}
\item $c_p(E)$ 属于代数 $A^*(X)$ 的中心；并且
\item 若 $g : X' \to X$ 局部有限型且 $c \in A^*(X' \to X)$，则
$c \circ c_p(E) = c_p(Lg^*E) \circ c$。
\end{enumerate}
\end{lemma}

\begin{proof}
第 (1) 部分立即由第 (2) 部分推出。设 $g : X' \to X$ 和
$c \in A^*(X' \to X)$ 如 (2)。为证明
$c \circ c_p(E) - c_p(Lg^*E) \circ c = 0$，使用
引理 \ref{chow-lemma-bivariant-zero} 的判据。因此可以假设 $X$ 整；由
引理 \ref{chow-lemma-chern-classes-computed}，甚至可以假设
$E$ 由常秩有限局部自由 $\mathcal{O}_X$-模组成的有界复形
$\mathcal{E}^\bullet$ 表示。此时只须证明
$$
c \cap c_p(\mathcal{E}^\bullet) \cap [X] =
c_p(\mathcal{E}^\bullet) \cap c \cap [X] 
$$
在 $\CH_*(X')$ 中成立。这立即由
引理 \ref{chow-lemma-cap-commutative-chern} 以及
$c_p(\mathcal{E}^\bullet)$ 作为局部自由模
$\mathcal{E}^n$ 的陈类之多项式的构造得出。
\end{proof}

\begin{lemma}
\label{chow-lemma-additivity-on-perfect}
设 $(S, \delta)$ 如情形 \ref{chow-situation-setup}。
设 $X$ 在 $S$ 上局部有限型，并设
$$
E_1 \to E_2 \to E_3 \to E_1[1]
$$
是 $D(\mathcal{O}_X)$ 中完美对象的可区别三角。若下列条件之一成立：
\begin{enumerate}
\item 存在包络 $f : Y \to X$，使得
$Lf^*E_1 \to Lf^*E_2$ 可由有限局部自由
$\mathcal{O}_Y$-模组成的局部有界复形之间的映射表示；
\item $E_1 \to E_2$ 可由有限局部自由 $\mathcal{O}_X$-模组成的
局部有界复形之间的映射表示；
\item $X$ 的不可约分支都是拟紧的；
\item $X$ 是拟紧的；或者
\item 此处增补更多情形；
\end{enumerate}
则 $E_1$、$E_2$、$E_3$ 的陈类均已定义，并且
$c(E_2) = c(E_1) c(E_3)$、$ch(E_2) = ch(E_1) + ch(E_3)$，以及
$P_p(E_2) = P_p(E_1) + P_p(E_3)$。
\end{lemma}

\begin{proof}
设 $f : Y \to X$ 是包络，并设
$\alpha^\bullet : \mathcal{E}_1^\bullet \to \mathcal{E}_2^\bullet$
是有限局部自由 $\mathcal{O}_Y$-模组成的局部有界复形之间的映射，
它表示 $Lf^*E_1 \to Lf^*E_2$。于是锥
$C(\alpha)^\bullet$ 表示 $Lf^*E_3$。由于
$C(\alpha)^n = \mathcal{E}_2^n \oplus \mathcal{E}_1^{n + 1}$，
$C(\alpha)^\bullet$ 是有限局部自由 $\mathcal{O}_Y$-模组成的
局部有界复形。因此 $E_1$、$E_2$、$E_3$ 的陈类均已定义。
此外，回忆 $c_p(E_1)$ 定义为 $A^p(X)$ 中唯一的元素，
其在 $A^p(Y)$ 中的限制是 $c_p(\mathcal{E}_1^\bullet)$；
$E_2$ 和 $E_3$ 也同理。因此只须在
$\prod_{p \geq 0} A^p(Y)$ 中证明
$c(\mathcal{E}_2^\bullet) =
c(\mathcal{E}_1^\bullet) c(C(\alpha)^\bullet)$。
进一步，只须在限制到 $Y$ 的每个连通分支后证明它。
因而可以假设复形 $\mathcal{E}_1^\bullet$ 和
$\mathcal{E}_2^\bullet$ 都是常秩有限局部自由
$\mathcal{O}_Y$-模组成的有界复形。此时所需等式由
引理 \ref{chow-lemma-additivity-chern-classes} 的乘法性推出。
对 $ch$ 或 $P_p$，使用引理 \ref{chow-lemma-chern-character-additive}。

\medskip\noindent
上一段已经说明，在条件 (1) 成立时引理成立。
由于 (2) 蕴含 (1)，第二种情形也已处理。现假设 (3)。
完全仿照引理 \ref{chow-lemma-chern-classes-defined} 的证明，
可找到包络 $f : Y \to X$，使得 $Y$ 是不交并
$Y = \coprod Y_i$，其中每个分支都是拟紧（且拟分离）的
具有消解性质的概形。于是 $Lf^*E_1 \to Lf^*E_2$
在 $Y_i$ 上的限制可由有限局部自由模组成的有界复形之间的映射表示；见
《概形的导出范畴》命题 \ref{perfect-proposition-perfect-resolution-property}。
由此可见条件 (3) 蕴含条件 (1)。当然，条件 (4) 蕴含条件 (3)，
证明完毕。
\end{proof}

\begin{remark}
\label{chow-remark-splitting-principle-perfect}
完美复形的陈类一旦有定义，就满足一种分裂原理。具体地，假设
$(S, \delta), X, E$ 如定义 \ref{chow-definition-defined-on-perfect}，并且 $E$ 的陈类已定义。
假设要证明双变类 $c_p(E)$、$P_p(E)$ 与 $ch_p(E)$ 之间的某个关系。
为此，可以选择包络 $f : Y \to X$ 以及表示 $E$ 的、由有限局部自由
$\mathcal{O}_X$-模组成的局部有界复形 $\mathcal{E}^\bullet$。
由引理 \ref{chow-lemma-defined-by-envelope} 中的唯一性，
只须证明双变类 $c_p(\mathcal{E}^\bullet)$、
$P_p(\mathcal{E}^\bullet)$ 与 $ch_p(\mathcal{E}^\bullet)$
之间所需的关系。因此可将 $X$ 替换为 $Y$ 的一个连通分支，
并假设 $E$ 由常秩有限局部自由模组成的有界复形
$\mathcal{E}^\bullet$ 表示。利用分裂原理
（引理 \ref{chow-lemma-splitting-principle}），可以假设每个
$\mathcal{E}^i$ 都有一个滤过，其逐次商
$\mathcal{L}_{i, j}$ 是可逆模。令
$x_{i, j} = c_1(\mathcal{L}_{i, j})$，则
$$
c(E) =
\prod\nolimits_{i\text{ 为偶数}} (1 + x_{i, j})
\prod\nolimits_{i\text{ 为奇数}} (1 + x_{i, j})^{-1}
$$
并且
$$
P_p(E) =  \sum\nolimits_{i\text{ 为偶数}} (x_{i, j})^p -
\sum\nolimits_{i\text{ 为奇数}} (x_{i, j})^p
$$
对上面 $c(E)$ 的表达式形式地取对数，得到
$$
\log(c(E)) = \sum (-1)^{p - 1}\frac{P_p(E)}{p}
$$
考察例 \ref{chow-example-power-sum} 中多项式
$P_p$ 的构造可知，$P_p(E)$ 关于 $E$ 的陈类的表达式
与向量丛情形完全相同。换言之，
\begin{align*}
P_1(E) & = c_1(E), \\
P_2(E) & = c_1(E)^2 - 2c_2(E), \\
P_3(E) & = c_1(E)^3 - 3c_1(E)c_2(E) + 3c_3(E), \\
P_4(E) & = c_1(E)^4 - 4c_1(E)^2c_2(E) + 4c_1(E)c_3(E) + 2c_2(E)^2 - 4c_4(E),
\end{align*}
等等。另一方面，双变类 $P_0(E) = r(E) = ch_0(E)$
不能从 $E$ 的陈类 $c(E)$ 恢复；陈类并不记录复形的秩。
\end{remark}

\begin{lemma}
\label{chow-lemma-chern-classes-perfect-dual}
在情形 \ref{chow-situation-setup} 中，设 $X$ 在 $S$ 上局部有限型。
设 $E \in D(\mathcal{O}_X)$ 是陈类已定义的完美对象。则在
$A^i(X)$ 中，$c_i(E^\vee) = (-1)^i c_i(E)$、
$P_i(E^\vee) = (-1)^iP_i(E)$，且
$ch_i(E^\vee) = (-1)^ich_i(E)$。
\end{lemma}

\begin{proof}
第一种证明：仿照引理 \ref{chow-lemma-commutative-chern-perfect}
的证明，把问题约化到 $E$ 由常秩有限局部自由模组成的有界复形表示的情形，
再应用引理 \ref{chow-lemma-chern-classes-dual}。
第二种证明：使用注 \ref{chow-remark-splitting-principle-perfect}
中讨论的分裂原理，以及 $E^\vee$ 的陈根是 $E$ 的陈根的相反数这一事实。
\end{proof}

\begin{lemma}
\label{chow-lemma-chern-class-perfect-tensor-invertible}
在情形 \ref{chow-situation-setup} 中，设 $X$ 在 $S$ 上局部有限型。
设 $E$ 是 $D(\mathcal{O}_X)$ 中陈类已定义的完美对象。
设 $\mathcal{L}$ 是可逆 $\mathcal{O}_X$-模。则
$$
c_i(E \otimes \mathcal{L}) =
\sum\nolimits_{j = 0}^i
\binom{r - i + j}{j} c_{i - j}(E) c_1(\mathcal{L})^j
$$
只要 $E$ 具有常秩 $r \in \mathbf{Z}$。
\end{lemma}

\begin{proof}
当 $E$ 是秩为 $r$ 的局部自由对象时，这就是
引理 \ref{chow-lemma-chern-classes-E-tensor-L}。读者可由这个特殊情形
通过形式计算推出本引理。另一种方法是使用
注 \ref{chow-remark-splitting-principle-perfect} 的分裂原理。
在这种方法中，最终需要证明如下代数事实：若形式地写成
$$
\frac{\prod_{a = 1, \ldots, n} (1 + x_a)}{\prod_{n = 1, \ldots, m} (1 + y_b)}
= 1 + c_1 + c_2 + c_3 + \ldots
$$
其中 $c_i$ 是 $\mathbf{Z}[x_i, y_j]$ 中次数为 $i$ 的齐次多项式，
则有
$$
\frac{\prod_{a = 1, \ldots, n} (1 + x_a + t)}{\prod_{b = 1, \ldots, m} (1 + y_b + t)}
= \sum\nolimits_{i \geq 0} \sum\nolimits_{j = 0}^i
\binom{r - i + j}{j} c_{i - j} t^j
$$
其中 $r = n - m$。细节从略。
\end{proof}

\begin{lemma}
\label{chow-lemma-chern-classes-perfect-tensor-product}
在情形 \ref{chow-situation-setup} 中，设 $X$ 在 $S$ 上局部有限型。
设 $E$ 与 $F$ 是 $D(\mathcal{O}_X)$ 中陈类已定义的完美对象。则
$$
c_1(E \otimes_{\mathcal{O}_X}^\mathbf{L} F) =
r(E) c_1(\mathcal{F}) + r(F) c_1(\mathcal{E})
$$
而 $c_2(E \otimes_{\mathcal{O}_X}^\mathbf{L} F)$ 的表达式为
$$
r(E) c_2(F) + r(F) c_2(E) + {r(E) \choose 2} c_1(F)^2 +
(r(E)r(F) - 1) c_1(F)c_1(E) + {r(F) \choose 2} c_1(E)^2
$$
$A^*(X)$ 中的更高陈类也有类似表达式。相应地，在
$A^*(X) \otimes \mathbf{Q}$ 中有
$ch(E \otimes_{\mathcal{O}_X}^\mathbf{L} F) = ch(E) ch(F)$。
更精确地，在 $A^p(X)$ 中有
$$
P_p(E \otimes_{\mathcal{O}_X}^\mathbf{L} F) = \sum\nolimits_{p_1 + p_2 = p}
{p \choose p_1} P_{p_1}(E) P_{p_2}(F)
$$
。
\end{lemma}

\begin{proof}
选择包络 $f : Y \to X$，使得 $Lf^*E$ 与 $Lf^*F$
可由有限局部自由 $\mathcal{O}_X$-模组成的局部有界复形表示后，
由向量丛的相应结论作计算即可得到本结论；见引理 \ref{chow-lemma-chern-classes-tensor-product} 和
\ref{chow-lemma-chern-character-multiplicative}。
更好的证明也许是像注 \ref{chow-remark-splitting-principle-perfect}
那样使用分裂原理，把本引理约化为多项式环中的计算；
下一段将说明这些计算。

\medskip\noindent
设 $A$ 是交换环（对我们而言，它是 $X$ 的双变 Chow 环中
由陈类生成的子环）。设 $S$ 是有限集，并给定映射
$\epsilon : S \to \{\pm 1\}$ 和 $f : S \to A$。定义
$$
P_p(S, f , \epsilon) = \sum\nolimits_{s \in S} \epsilon(s) f(s)^p
$$
为 $A$ 中的元素。给定第二个三元组 $(S', \epsilon', f')$，
关于 $P_p$ 所须证明的等式是
$$
P_p(S \times S', f + f' , \epsilon \epsilon') = 
\sum\nolimits_{p_1 + p_2 = p}
{p \choose p_1} P_{p_1}(S, f, \epsilon) P_{p_2}(S', f', \epsilon')
$$
为说明它成立，可约化到以 $S \amalg S'$ 为变量的多项式环，
并验证每一项 $f(s)^if'(s')^j$ 在等式两边的系数相同。
为验证 $c_1(E \otimes_{\mathcal{O}_X}^\mathbf{L} F)$ 和
$c_2(E \otimes_{\mathcal{O}_X}^\mathbf{L} F)$ 的公式，
利用分裂原理把问题约化为在无挠环中检验这些公式。
然后使用注 \ref{chow-remark-splitting-principle-perfect}
中证明的 $P_j(E)$ 与 $c_i(E)$ 之间的关系。例如
$$
c_1(E \otimes F) = P_1(E \otimes F) = r(F)P_1(E) + r(E)P_1(F) =
r(F)c_1(E) + r(E)c_1(F)
$$
其中中间的等式使用了按定义 $r(E) = P_0(E)$。类似地，
\begin{align*}
& 2c_2(E \otimes F) \\
& = c_1(E \otimes F)^2 - P_2(E \otimes F) \\
& =
(r(F)c_1(E) + r(E)c_1(F))^2 -
r(F)P_2(E) - P_1(E)P_1(F) - r(E)P_2(F) \\
& =
(r(F)c_1(E) + r(E)c_1(F))^2 -
r(F)(c_1(E)^2 - 2c_2(E)) - c_1(E)c_1(F) - \\
& \quad r(E)(c_1(F)^2 - 2c_2(F))
\end{align*}
读者可验证，这与引理陈述中的公式相差因子 $2$ 后一致。
\end{proof}







\section{局部化陈类的一个简单情形}
\label{chow-section-preparation-localized-chern}

\noindent
本节讨论下一引理所构造的双变类的一些性质；其中大多数性质
立即由引理给出的刻画推出。我们建议读者跳过本节其余部分。

\begin{lemma}
\label{chow-lemma-silly}
设 $(S, \delta)$ 如情形 \ref{chow-situation-setup}。
设 $X$ 在 $S$ 上局部有限型。设 $i_j : X_j \to X$（$j = 1, 2$）
是闭浸入，并且在集合论意义下 $X = X_1 \cup X_2$。
设 $E_2 \in D(\mathcal{O}_{X_2})$ 是完美对象。假设
\begin{enumerate}
\item $E_2$ 的陈类已定义；
\item 限制 $E_2|_{X_1 \cap X_2}$ 为零，
分别地\ 同构于置于上同调次数 $0$ 的、秩 $< p$ 的有限局部自由
$\mathcal{O}_{X_1 \cap X_2}$-模。
\end{enumerate}
则存在典范双变类
$$
P'_p(E_2),\text{ 分别 }c'_p(E_2) \in A^p(X_2 \to X)
$$
它由下列性质刻画：
$$
P'_p(E_2) \cap i_{2, *} \alpha_2 = P_p(E_2) \cap \alpha_2
\quad\text{且}\quad
P'_p(E_2) \cap i_{1, *} \alpha_1 = 0,
$$
相应地，
$$
c'_p(E_2) \cap i_{2, *} \alpha_2 = c_p(E_2) \cap \alpha_2
\quad\text{且}\quad
c'_p(E_2) \cap i_{1, *} \alpha_1 = 0
$$
对 $\alpha_i \in \CH_k(X_i)$ 成立，并且在任意局部有限型基变换
$X' \to X$ 后也同样成立。
\end{lemma}

\begin{proof}
下文将直接使用第 \ref{chow-section-pre-derived} 节的内容，不再另行说明。

\medskip\noindent
假设 $E_2|_{X_1 \cap X_2}$ 为零。考虑局部有限型概形态射
$X' \to X$，并记 $i'_j : X'_j \to X'$ 为 $i_j$ 的基变换。
由引理 \ref{chow-lemma-exact-sequence-closed-chow}，任意元素
$\alpha' \in \CH_k(X')$ 都可写成
$i'_{1, *}\alpha'_1 + i'_{2, *}\alpha'_2$，其中
$\alpha'_2 \in \CH_k(X'_2)$ 的不确定性至多相差一个由
$X'_1 \cap X'_2 \to X'_2$ 推前的像中的元素。于是可以令
$P'_p(E_2) \cap \alpha' = P_p(E_2) \cap \alpha'_2 \in \CH_{k - p}(X'_2)$。
由于假设 $E_2$ 在 $X_1 \cap X_2$ 上的限制为零，
这个定义是良定的。

\medskip\noindent
若 $E_2|_{X_1 \cap X_2}$ 同构于置于上同调次数 $0$ 的、秩
$< p$ 的有限局部自由 $\mathcal{O}_{X_1 \cap X_2}$-模，
则由秩的考量有 $c_p(E_2|_{X_1 \cap X_2}) = 0$，
因而可以完全相同的方式论证。
\end{proof}

\begin{lemma}
\label{chow-lemma-silly-independent}
引理 \ref{chow-lemma-silly} 中 $A^p(X_2 \to X)$ 内的双变类
$P'_p(E_2)$，分别地\ $c'_p(E_2)$，与 $X_1$ 的选择无关。
\end{lemma}

\begin{proof}
设 $X_1' \subset X$ 是另一个闭子概形，使得在集合论意义下
$X = X'_1 \cup X_2$，并且限制 $E_2|_{X'_1 \cap X_2}$ 为零，
分别地\ 同构于置于上同调次数 $0$ 的、秩 $< p$ 的有限局部自由
$\mathcal{O}_{X'_1 \cap X_2}$-模。于是
$X = (X_1 \cap X'_1) \cup X_2$。因此任意元素
$\alpha \in \CH_k(X)$ 都可写成 $i_*\beta + i_{2, *}\alpha_2$，其中
$\alpha_2 \in \CH_k(X'_2)$ 且 $\beta \in \CH_k(X_1 \cap X'_1)$。
由此显然，
$P'_p(E_2) \cap \alpha = P_p(E_2) \cap \alpha_2 \in \CH_{k - p}(X_2)$，
分别地\  $c'_p(E_2) \cap \alpha = c_p(E_2) \cap \alpha_2 \in \CH_{k - p}(X_2)$，
与使用 $X_1$ 还是 $X'_1$ 无关。任意基变换后也同样成立。
\end{proof}

\begin{lemma}
\label{chow-lemma-base-change-silly}
在引理 \ref{chow-lemma-silly} 中，设 $X' \to X$ 是局部有限型态射。
记 $X' = X'_1 \cup X'_2$ 以及
$E'_2 \in D(\mathcal{O}_{X'_2})$ 为拉回到 $X'$ 的对象。
则使用 $X' = X'_1 \cup X'_2$ 和 $E_2'$，由
引理 \ref{chow-lemma-silly} 构造的 $A^p(X_2' \to X')$ 中的类
$P'_p(E_2')$，分别地\ $c'_p(E_2')$，是
$A^p(X_2 \to X)$ 中的类 $P'_p(E_2)$，分别地\ $c'_p(E_2)$，
的限制（注 \ref{chow-remark-restriction-bivariant}）。
\end{lemma}

\begin{proof}
立即由引理 \ref{chow-lemma-silly} 对这些类的刻画得出。
\end{proof}

\begin{lemma}
\label{chow-lemma-silly-silly}
在引理 \ref{chow-lemma-silly} 中，设 $E_2$ 是某个完美对象
$E \in D(\mathcal{O}_X)$ 的限制，并且 $E|_{X_1}$ 为零，
分别地\ 同构于置于上同调次数 $0$ 的、秩 $< p$ 的有限局部自由
$\mathcal{O}_{X_1}$-模。若 $E$ 的陈类已定义，则
$i_{2, *} \circ P'_p(E_2) = P_p(E)$，
分别地\ $i_{2, *} \circ c'_p(E_2) = c_p(E)$
（其中 $\circ$ 如引理 \ref{chow-lemma-push-proper-bivariant}）。
\end{lemma}

\begin{proof}
先假设 $E|_{X_1}$ 为零。采用引理 \ref{chow-lemma-silly}
证明中的记号，此情形下本引理由下式推出：
\begin{align*}
P_p(E) \cap \alpha'
& =
i'_{1, *}(P_p(E) \cap \alpha'_1) +
i'_{2, *}(P_p(E) \cap \alpha'_2) \\
& =
i'_{1, *}(P_p(E|_{X_1}) \cap \alpha'_1) +
i'_{2, *}(P'_p(E_2) \cap \alpha') \\
& =
i'_{2, *}(P'_p(E_2) \cap \alpha')
\end{align*}
当 $E|_{X_1}$ 同构于置于上同调次数 $0$ 的、秩 $< p$ 的
有限局部自由 $\mathcal{O}_{X_1}$-模时，论证类似。
\end{proof}

\begin{lemma}
\label{chow-lemma-silly-shrink}
在引理 \ref{chow-lemma-silly} 中，假设存在闭子概形
$X'_2 \subset X_2$ 和 $X_1 \subset X'_1 \subset X$，使得在集合论意义下
$X = X'_1 \cup X'_2$。假设 $E_2|_{X'_1 \cap X_2}$ 为零，
分别地\ 同构于置于次数 $0$ 的、秩 $< p$ 的有限局部自由模。则
$(X'_2 \to X_2)_* \circ P'_p(E_2|_{X'_2}) = P'_p(E_2)$，
分别地\  $(X'_2 \to X_2)_* \circ c'_p(E_2|_{X'_2}) = c_p(E_2)$
（其中 $\circ$ 如引理 \ref{chow-lemma-push-proper-bivariant}）。
\end{lemma}

\begin{proof}
立即由引理 \ref{chow-lemma-silly} 对这些类的刻画得出。
\end{proof}

\begin{lemma}
\label{chow-lemma-silly-commutes}
在引理 \ref{chow-lemma-silly} 中，设 $f : Y \to X$ 局部有限型，
并设 $c \in A^*(Y \to X)$。则在 $A^*(Y_2 \to Y)$ 中有
$$
c \circ P'_p(E_2) = P'_p(Lf_2^*E_2) \circ c
\quad\text{分别}\quad
c \circ c'_p(E_2) = c'_p(Lf_2^*E_2) \circ c
$$
其中 $f_2 : Y_2 \to X_2$ 是 $f$ 的基变换。
\end{lemma}

\begin{proof}
设 $\alpha \in \CH_k(X)$。可写成
$$
\alpha = \alpha_1 + \alpha_2
$$
其中 $\alpha_i \in \CH_k(X_i)$；这里省略了闭浸入
$X_i \to X$ 的推前。读者于是可检验
$c'_p(E_2) \cap \alpha = c_p(E_2) \cap \alpha_2$、
$c \cap c'_p(E_2) \cap \alpha = c \cap c_p(E_2) \cap \alpha_2$、
$c \cap \alpha = c \cap \alpha_1 + c \cap \alpha_2$，以及
$c'_p(Lf_2^*E_2) \cap c \cap \alpha = c_p(Lf_2^*E_2) \cap c \cap \alpha_2$。
最后应用引理 \ref{chow-lemma-commutative-chern-perfect}。
\end{proof}

\begin{lemma}
\label{chow-lemma-silly-compose}
在引理 \ref{chow-lemma-silly} 中，假设 $E_2|_{X_1 \cap X_2}$ 为零。则
\begin{align*}
P'_1(E_2) & = c'_1(E_2), \\
P'_2(E_2) & = c'_1(E_2)^2 - 2c'_2(E_2), \\
P'_3(E_2) & = c'_1(E_2)^3 - 3c'_1(E_2)c'_2(E_2) + 3c'_3(E_2), \\
P'_4(E_2) & = c'_1(E_2)^4 - 4c'_1(E_2)^2c'_2(E_2) +
4c'_1(E_2)c'_3(E_2) + 2c'_2(E_2)^2 - 4c'_4(E_2),
\end{align*}
等等，其中乘法如注 \ref{chow-remark-ring-loc-classes}。
\end{lemma}

\begin{proof}
零层的秩 $< 1$，所以类 $c'_p(E_2)$ 对所有 $p \geq 1$
都有定义，因而该陈述有意义。这些等式立即由
引理 \ref{chow-lemma-silly} 所构造类的刻画，以及
注 \ref{chow-remark-splitting-principle-perfect} 给出的
与 $E_2$ 的陈类取帽积的相应结论推出。
\end{proof}

\begin{lemma}
\label{chow-lemma-silly-sum-c}
设 $(S, \delta)$ 如情形 \ref{chow-situation-setup}。
设 $X$ 在 $S$ 上局部有限型。设 $i_j : X_j \to X$（$j = 1, 2$）
是闭浸入，使得在集合论意义下 $X = X_1 \cup X_2$。
设 $E, F \in D(\mathcal{O}_X)$ 是完美对象。假设
\begin{enumerate}
\item $E$ 与 $F$ 的陈类已定义；
\item 限制 $E|_{X_1 \cap X_2}$ 与 $F|_{X_1 \cap X_2}$
同构于置于上同调次数 $0$ 的有限局部自由
$\mathcal{O}_{X_1}$-模，其秩分别 $< p$ 和 $< q$。
\end{enumerate}
采用注 \ref{chow-remark-ring-loc-classes} 的记号，令
$$
c^{(p)}(E) = 1 + c_1(E) + \ldots + c_{p - 1}(E) +
c'_p(E|_{X_2}) + c'_{p + 1}(E|_{X_2}) + \ldots \in A^{(p)}(X_2 \to X)
$$
其中 $c'_p(E|_{X_2})$ 如引理 \ref{chow-lemma-silly}。
$c^{(q)}(F)$ 与 $c^{(p + q)}(E \oplus F)$ 也类似定义。
则在 $A^{(p + q)}(X_2 \to X)$ 中有
$c^{(p + q)}(E \oplus F) = c^{(p)}(E)c^{(q)}(F)$。
\end{lemma}

\begin{proof}
立即由引理 \ref{chow-lemma-silly} 对这些类的刻画以及
引理 \ref{chow-lemma-additivity-on-perfect} 中的可加性得出。
\end{proof}

\begin{lemma}
\label{chow-lemma-silly-sum-P}
设 $(S, \delta)$ 如情形 \ref{chow-situation-setup}。
设 $X$ 在 $S$ 上局部有限型。设 $i_j : X_j \to X$（$j = 1, 2$）
是闭浸入，使得在集合论意义下 $X = X_1 \cup X_2$。
设 $E, F \in D(\mathcal{O}_{X_2})$ 是完美对象。假设
\begin{enumerate}
\item $E$ 与 $F$ 的陈类已定义；
\item 限制 $E|_{X_1 \cap X_2}$ 与 $F|_{X_1 \cap X_2}$ 均为零。
\end{enumerate}
对 $p \geq 0$，记由引理 \ref{chow-lemma-silly} 构造的类为
$P'_p(E), P'_p(F), P'_p(E \oplus F) \in A^p(X_2 \to X)$。则
$P'_p(E \oplus F) = P'_p(E) + P'_p(F)$。
\end{lemma}

\begin{proof}
立即由引理 \ref{chow-lemma-silly} 对这些类的刻画以及
引理 \ref{chow-lemma-additivity-on-perfect} 中的可加性得出。
\end{proof}

\begin{lemma}
\label{chow-lemma-silly-tensor-invertible}
在引理 \ref{chow-lemma-silly} 中，假设 $E_2$ 具有常秩 $0$。
设 $\mathcal{L}$ 是可逆 $\mathcal{O}_X$-模。则
$$
c'_i(E_2 \otimes \mathcal{L}) =
\sum\nolimits_{j = 0}^i
\binom{- i + j}{j} c'_{i - j}(E_2) c_1(\mathcal{L})^j
$$
\end{lemma}

\begin{proof}
关于秩的假设蕴含 $E_2|_{X_1 \cap X_2}$ 为零。
因此 $c'_i(E_2)$ 对所有 $i \geq 1$ 都有定义，陈述有意义。
所述等式立即由引理 \ref{chow-lemma-chern-class-perfect-tensor-invertible}
以及引理 \ref{chow-lemma-silly} 对 $c'_i$ 的刻画得出。
\end{proof}

\begin{lemma}
\label{chow-lemma-silly-tensor-product}
在情形 \ref{chow-situation-setup} 中，设 $X$ 在 $S$ 上局部有限型。设
$$
X = X_1 \cup X_2 = X'_1 \cup X'_2
$$
是在集合论意义下把 $X$ 写成闭子概形之并的两种方式。
设 $E$、$E'$ 是 $D(\mathcal{O}_X)$ 中陈类已定义的完美对象。
假设 $E|_{X_1}$ 和 $E'|_{X'_1}$ 为零\footnote{大概存在本引理的一个变体，
其中只假设这些限制同构于秩 $< p$ 和 $< p'$ 的有限局部自由模。}，
对 $i = 1, 2$，记
\begin{enumerate}
\item $r = P'_0(E) \in A^0(X_2 \to X)$ 以及
$r' = P'_0(E') \in A^0(X'_2 \to X)$；
\item $\gamma_p = c'_p(E|_{X_2}) \in A^p(X_2 \to X)$ 以及
$\gamma'_p = c'_p(E'|_{X'_2}) \in A^p(X'_2 \to X)$；
\item $\chi_p = P'_p(E|_{X_2}) \in A^p(X_2 \to X)$ 以及
$\chi'_p = P'_p(E'|_{X'_2}) \in A^p(X'_2 \to X)$
\end{enumerate}
为引理 \ref{chow-lemma-silly} 所构造的类。则有
$$
c'_1((E \otimes_{\mathcal{O}_X}^\mathbf{L} E')|_{X_2 \cap X'_2}) =
r \gamma'_1 + r' \gamma_1
$$
在 $A^1(X_2 \cap X'_2 \to X)$ 中成立，并且
$$
c'_2((E \otimes_{\mathcal{O}_X}^\mathbf{L} E')|_{X_2 \cap X'_2}) =
r \gamma'_2 + r' \gamma_2 + {r \choose 2} (\gamma'_1)^2 +
(rr' - 1) \gamma'_1\gamma_1 + {r' \choose 2} \gamma_1^2
$$
在 $A^2(X_2 \cap X'_2 \to X)$ 中成立；更高陈类也类似。
同样有
$$
P'_p((E \otimes_{\mathcal{O}_X}^\mathbf{L} E')|_{X_2 \cap X'_2}) =
\sum\nolimits_{p_1 + p_2 = p}
{p \choose p_1} \chi_{p_1} \chi'_{p_2}
$$
在 $A^p(X_2 \cap X'_2 \to X)$ 中成立。
\end{lemma}

\begin{proof}
先说明该陈述有意义。事实上，
$X = (X_2 \cap X'_2) \cup Y$，其中
$Y = (X_1 \cap X'_1) \cup (X_1 \cap X'_2) \cup (X_2 \cap X'_1)$，
而对象 $E \otimes_{\mathcal{O}_X}^\mathbf{L} E'$ 在 $Y$ 上的限制为零。
所述等式由引理 \ref{chow-lemma-silly} 对这些类的刻画以及
引理 \ref{chow-lemma-chern-classes-perfect-tensor-product} 的等式推出。
细节从略。
\end{proof}







\section{无穷远处的 Gysin 类}
\label{chow-section-gysin-at-infty}

\noindent
本节讨论下一引理所构造的双变类。我们建议读者跳过本节其余部分。

\begin{lemma}
\label{chow-lemma-gysin-at-infty}
设 $(S, \delta)$ 如情形 \ref{chow-situation-setup}。
设 $X$ 在 $S$ 上局部有限型。设
$b : W \to \mathbf{P}^1_X$ 是真概形态射，并且在
$\mathbf{A}^1_X$ 上是同构。记
$i_\infty : W_\infty \to W$ 为除子
$D_\infty \subset \mathbf{P}^1_X$ 的逆像，其中补集为
$\mathbf{A}^1_X$。则存在典范双变类
$$
C \in A^0(W_\infty \to X)
$$
具有如下性质：对 $\alpha \in \CH_k(X)$，
$i_{\infty, *}(C \cap \alpha) = i_{0, *}\alpha$；
任意局部有限型基变换 $X' \to X$ 后也同样成立。
\end{lemma}

\begin{proof}
给定 $\alpha \in \CH_k(X)$，存在 $\beta \in \CH_{k + 1}(W)$，
其在 $b^{-1}(\mathbf{A}^1_X)$ 上的限制是 $\alpha$ 的平坦拉回；见
引理 \ref{chow-lemma-exact-sequence-open}。$\beta$ 的另一个选择与
$\beta$ 相差一个支撑在 $W_\infty$ 上的循环；见
引理 \ref{chow-lemma-restrict-to-open}。由于
(\ref{chow-equation-zero-infty}) 中的有效 Cartier 除子
$D_\infty \subset \mathbf{P}^1_X$ 的正规丛平凡，
Gysin 同态 $i_\infty^*$ 消去支撑在 $W_\infty$ 上的循环类；见
注 \ref{chow-remark-gysin-on-cycles}。因此令
$C \cap \alpha = i_\infty^*\beta$ 是良定的。

\medskip\noindent
由于 $W_\infty$ 与 $W_0 = X \times \{0\}$ 分别是
$\mathbf{P}^1_X$ 中有理等价的有效 Cartier 除子
$D_0, D_\infty$ 的拉回，$i_\infty^*\beta$ 与
$i_0^*\beta$ 映到 $W$ 上同一个循环类；确切地说，由
引理 \ref{chow-lemma-support-cap-effective-Cartier}，二者都表示类
$c_1(\mathcal{O}_{\mathbf{P}^1_X}(1)) \cap \beta$。
由 $\beta$ 的选择，在 $W_0 = X \times \{0\}$ 上作为循环有
$i_0^*\beta = \alpha$；例如见
引理 \ref{chow-lemma-relative-effective-cartier}。
因此得到引理所述的
$i_{\infty, *}(C \cap \alpha) = i_{0, *}\alpha$。

\medskip\noindent
注意，对任意局部有限型基变换 $X' \to X$，关于 $b$ 的假设仍然成立。
因此由上述同一构造得到运算
$C \cap - : \CH_k(X') \to \CH_k(W'_\infty)$。
为说明这族运算定义了双变类，考虑图
$$
\xymatrix{
& & & \CH_*(X) \ar[d]^{\text{平坦拉回}} \\
\CH_{* + 1}(W_\infty) \ar[r] \ar[rd]^0 &
\CH_{* + 1}(W) \ar[d]^{i_\infty^*} \ar[rr]^{\text{平坦拉回}} & &
\CH_{* + 1}(\mathbf{A}^1_X) \ar[r] \ar@{..>}[lld]^{C \cap -} &
0 \\
& \CH_*(W_\infty)
}
$$
这里取图示的 $X$，并考虑对任意 $X' \to X$ 的基变换图。已知平坦拉回和
$i_\infty^*$ 都是双变运算；见引理 \ref{chow-lemma-flat-pullback-bivariant} 和 \ref{chow-lemma-gysin-bivariant}。
于是一个形式论证（其中涉及概形及其 Chow 群的巨大交换图）
表明虚线箭头是双变运算。
\end{proof}

\begin{lemma}
\label{chow-lemma-base-change-gysin-at-infty}
在引理 \ref{chow-lemma-gysin-at-infty} 中，设 $X' \to X$ 是局部有限型态射。
记 $b' : W' \to \mathbf{P}^1_{X'}$ 和
$i'_\infty : W'_\infty \to W'$ 分别为 $b$ 与 $i_\infty$ 的基变换。
则使用 $b'$ 按引理 \ref{chow-lemma-gysin-at-infty} 构造的类
$C' \in A^0(W'_\infty \to X')$ 是 $C$ 的限制
（注 \ref{chow-remark-restriction-bivariant}）。
\end{lemma}

\begin{proof}
立即由构造以及平坦拉回和 $i_\infty^*$ 也满足类似陈述这一事实得出。
\end{proof}

\begin{lemma}
\label{chow-lemma-gysin-at-infty-independent}
在引理 \ref{chow-lemma-gysin-at-infty} 中，设 $g : W' \to W$ 是真态射，
并且在 $\mathbf{A}^1_X$ 上是同构。设
$C' \in A^0(W'_\infty \to X)$ 和 $C \in A^0(W_\infty \to X)$
是引理 \ref{chow-lemma-gysin-at-infty} 所构造的类。则在
$A^0(W_\infty \to X)$ 中有 $g_{\infty, *} \circ C' = C$。
\end{lemma}

\begin{proof}
令 $b' = b \circ g : W' \to \mathbf{P}^1_X$。
记 $i'_\infty : W'_\infty \to W'$ 为包含态射，并记
$g_\infty : W'_\infty \to W_\infty$ 为 $g$ 的限制。
给定 $\alpha \in \CH_k(X)$，选择 $\beta' \in \CH_{k + 1}(W')$，
使其在 $(b')^{-1}\mathbf{A}^1_X$ 上的限制为 $\alpha$ 的平坦拉回。
于是 $\beta = g_*\beta' \in \CH_{k + 1}(W)$ 在
$b^{-1}\mathbf{A}^1_X$ 上的限制为 $\alpha$ 的平坦拉回。由
引理 \ref{chow-lemma-closed-in-X-gysin}，
$i_\infty^*\beta = g_{\infty, *}(i'_\infty)^*\beta'$。
此等式及其经任意局部有限型态射 $X' \to X$ 基变换后的相应事实，
正给出引理的断言。
\end{proof}

\begin{lemma}
\label{chow-lemma-homomorphism-pre}
在引理 \ref{chow-lemma-gysin-at-infty} 中有
$C \circ (W_\infty \to X)_* \circ i_\infty^* = i_\infty^*$。
\end{lemma}

\begin{proof}
设 $\beta \in \CH_{k + 1}(W)$。记
$i_0 : X = X \times \{0\} \to W$ 为 $\mathbf{P}^1$ 中位于
$0$ 上方的纤维的闭浸入。则在 $\CH_k(X)$ 中有
$(W_\infty \to X)_* i_\infty^* \beta = i_0^*\beta$，因为
$i_{\infty, *}i_\infty^*\beta$ 与 $i_{0, *}i_0^*\beta$
表示 $W$ 上的同一个类（例如由
引理 \ref{chow-lemma-support-cap-effective-Cartier}），
所以推前到 $X$ 上也得到同一个类。
$\beta$ 在 $b^{-1}(\mathbf{A}^1_X)$ 上的限制进一步限制为
$i_0^*\beta = (W_\infty \to X)_* i_\infty^* \beta$ 的平坦拉回，
因为可在由 $i_0$ 拉回后检验这一点；见引理 \ref{chow-lemma-linebundle} 和 \ref{chow-lemma-linebundle-formulae}。
因此在计算 $(W_\infty \to X)_*i_\infty^*\beta$ 在 $C$
下的像时可以使用 $\beta$，从而得到所需结论。
\end{proof}

\begin{lemma}
\label{chow-lemma-gysin-at-infty-commutes}
在引理 \ref{chow-lemma-gysin-at-infty} 中，设 $f : Y \to X$
是局部有限型态射且 $c \in A^*(Y \to X)$。则在
$A^*(W_\infty \times_X Y \to X)$ 中有 $C \circ c = c \circ C$。
\end{lemma}

\begin{proof}
考虑交换图
$$
\xymatrix{
W_\infty \times_X Y \ar@{=}[r] &
W_{Y, \infty} \ar[r]_{i_{Y, \infty}} \ar[d] &
W_Y \ar[r]_{b_Y} \ar[d] &
\mathbf{P}^1_Y \ar[r]_{p_Y} \ar[d] &
Y \ar[d]^f \\
& W_\infty \ar[r]^{i_\infty} &
W \ar[r]^b &
\mathbf{P}^1_X \ar[r]^p &
X
}
$$
其中各方块均为笛卡尔方块。对元素 $\alpha \in \CH_k(X)$，
选择 $\beta \in \CH_{k + 1}(W)$，使其在
$b^{-1}(\mathbf{A}^1_X)$ 上的限制为 $\alpha$ 的平坦拉回。
于是 $c \cap \beta$ 是 $\CH_*(W_Y)$ 中的类，其在
$b_Y^{-1}(\mathbf{A}^1_Y)$ 上的限制为 $c \cap \alpha$ 的平坦拉回。
其次，由于 $c$ 是双变类，有
$$
i_{Y, \infty}^*(c \cap \beta) = c \cap i_\infty^*\beta
$$
这正说明 $C \cap c \cap \alpha = c \cap C \cap \alpha$。
对任意局部有限型基变换 $X' \to X$，同一论证仍然适用。
这就证明了引理。
\end{proof}





\section{局部化陈类的准备}
\label{chow-section-preparation-localized-chern-II}

\noindent
本节讨论下一引理所构造的双变类的一些性质。
我们建议读者跳过本节其余部分。

\begin{lemma}
\label{chow-lemma-localized-chern-pre}
设 $(S, \delta)$ 如情形 \ref{chow-situation-setup}。
设 $X$ 在 $S$ 上局部有限型，并设 $Z \subset X$ 是闭子概形。设
$$
b : W \longrightarrow \mathbf{P}^1_X
$$
是真概形态射。设 $Q \in D(\mathcal{O}_W)$ 是完美对象。
记 $W_\infty \subset W$ 为除子
$D_\infty \subset \mathbf{P}^1_X$ 的逆像，其中补集为
$\mathbf{A}^1_X$。假设
\begin{enumerate}
\item[(A0)] $Q$ 的陈类已定义
（第 \ref{chow-section-pre-derived} 节）；
\item[(A1)] $b$ 在 $\mathbf{A}^1_X$ 上是同构；
\item[(A2)] 存在闭子概形 $T \subset W_\infty$，
它包含 $W_\infty$ 中所有位于 $X \setminus Z$ 上方的点，
并且 $Q|_T$ 为零，分别地\ 同构于置于上同调次数 $0$ 的、
秩 $< p$ 的有限局部自由 $\mathcal{O}_T$-模。
\end{enumerate}
则存在典范双变类
$$
P'_p(Q),\text{ 分别 }c'_p(Q) \in A^p(Z \to X)
$$
使得
$(Z \to X)_* \circ P'_p(Q) = P_p(Q|_{X \times \{0\}})$，
分别地\ $(Z \to X)_* \circ c'_p(Q) = c_p(Q|_{X \times \{0\}})$。
\end{lemma}

\begin{proof}
记 $E \subset W_\infty$ 为 $Z$ 的逆像。于是
$W_\infty = T \cup E$，且 $b$ 诱导真态射 $E \to Z$。
记 $C \in A^0(W_\infty \to X)$ 为
引理 \ref{chow-lemma-gysin-at-infty} 所构造的双变类。记
$A^p(E \to W_\infty)$ 中的 $P'_p(Q|_E)$，分别地\ $c'_p(Q|_E)$，
为引理 \ref{chow-lemma-silly} 所构造的双变类。这是有意义的，
因为由假设 (A2)，$(Q|_E)|_{T \cap E}$ 为零，分别地\ 同构于
置于上同调次数 $0$ 的、秩 $< p$ 的有限局部自由
$\mathcal{O}_{E \cap T}$-模。定义
$$
P'_p(Q) = (E \to Z)_* \circ P'_p(Q|_E) \circ C,\text{ 分别 }
c'_p(Q) = (E \to Z)_* \circ c'_p(Q|_E) \circ C
$$
这是双变类；见引理 \ref{chow-lemma-push-proper-bivariant}。
由于 $E \to Z \to X$ 等于 $E \to W_\infty \to W \to X$，有
\begin{align*}
(Z \to X)_* \circ c'_p(Q)
& =
(W \to X)_* \circ i_{\infty, *} \circ (E \to W_\infty)_*
\circ c'_p(Q|_E) \circ C \\
& =
(W \to X)_* \circ i_{\infty, *} \circ c_p(Q|_{W_\infty}) \circ C \\
& =
(W \to X)_* \circ c_p(Q) \circ i_{\infty, *} \circ C \\
& =
(W \to X)_*\circ c_p(Q) \circ i_{0, *} \\
& =
(W \to X)_* \circ i_{0, *} \circ c_p(Q|_{X \times \{0\}}) \\
& =
c_p(Q|_{X \times \{0\}})
\end{align*}
第二个等式由引理 \ref{chow-lemma-silly-silly} 得出。
第三个等式使用了 $c_p(Q)$ 是双变类。
第四个等式由引理 \ref{chow-lemma-gysin-at-infty} 得出。
第五个等式再次使用了 $c_p(Q)$ 是双变类。
最后一个等式是因为按上述方式将 $W_0$ 与 $X$ 识别后，
$(W_0 \to W) \circ (W \to X)$ 是 $X$ 上的恒等态射。
完全相同的一串等式证明关于 $P'_p(Q)$ 的性质。
\end{proof}

\begin{lemma}
\label{chow-lemma-base-change-localized-chern-pre}
在引理 \ref{chow-lemma-localized-chern-pre} 中，设
$X' \to X$ 是局部有限型态射。记
$Z'$、$b' : W' \to \mathbf{P}^1_{X'}$ 与
$T' \subset W'_\infty$ 分别为 $Z$、
$b : W \to \mathbf{P}^1_X$ 与 $T \subset W_\infty$ 的基变换。
令 $Q' = (W' \to W)^*Q$。则使用 $b'$、$Q'$ 与 $T'$，
按引理 \ref{chow-lemma-localized-chern-pre} 构造的
$A^p(Z' \to X')$ 中的类 $P'_p(Q')$，分别地\ $c'_p(Q')$，
是 $A^p(Z \to X)$ 中的类 $P'_p(Q)$，分别地\ $c'_p(Q)$，
的限制（注 \ref{chow-remark-restriction-bivariant}）。
\end{lemma}

\begin{proof}
回忆该构造为
$$
P'_p(Q) = (E \to Z)_* \circ P'_p(Q|_E) \circ C,\text{ 分别 }
c'_p(Q) = (E \to Z)_* \circ c'_p(Q|_E) \circ C
$$
因此本引理由 $C$ 的相应基变换性质
（引理 \ref{chow-lemma-base-change-gysin-at-infty}），以及
引理 \ref{chow-lemma-silly} 所构造的类也具有同样基变换性质这一事实推出
（省略一个小细节）。
\end{proof}

\begin{lemma}
\label{chow-lemma-localized-chern-pre-independent}
引理 \ref{chow-lemma-localized-chern-pre} 中的双变类
$P'_p(Q)$，分别地\ $c'_p(Q)$，与闭子概形 $T$ 的选择无关。
此外，给定在 $\mathbf{A}^1_X$ 上为同构的真态射
$g : W' \to W$，令 $Q' = g^*Q$，则有
$P'_p(Q) = P'_p(Q')$，分别地\ $c'_p(Q) = c'_p(Q')$。
\end{lemma}

\begin{proof}
关于 $T$ 的独立性立即由引理 \ref{chow-lemma-silly-independent} 得出。

\medskip\noindent
设 $g : W' \to W$ 是在 $\mathbf{A}^1_X$ 上为同构的真态射。
注意，取 $T' = g^{-1}(T) \subset W'_\infty$ 得到满足 (A2) 的闭子概形，
所以与 $b' = b \circ g : W' \to \mathbf{P}^1_X$ 和 $Q'$ 对应的
$A^p(Z \to X)$ 中的运算子 $P'_p(Q')$，分别地\ $c'_p(Q')$，
已有定义。记 $E' \subset W'_\infty$ 为 $Z$ 在
$W'_\infty$ 中的逆像。回忆
$$
c'_p(Q') = (E' \to Z)_* \circ c'_p(Q'|_{E'}) \circ C'
$$
其中 $C' \in A^0(W'_\infty \to X)$ 且
$c'_p(Q'|_{E'}) \in A^p(E' \to W'_\infty)$。
由引理 \ref{chow-lemma-gysin-at-infty-independent}，
$g_{\infty, *} \circ C' = C$。还要注意，$E'$ 也是 $E$ 经
$g_\infty$ 在 $W'_\infty$ 中的逆像。又因 $Q' = g^*Q$，
$c'_p(Q'|_{E'})$ 正是 $c'_p(Q|_E)$ 在位于
$W'_\infty$ 上方的概形上的限制；见
注 \ref{chow-remark-restriction-bivariant}。因此
\begin{align*}
c'_p(Q')
& = 
(E' \to Z)_* \circ c'_p(Q'|_{E'}) \circ C' \\
& =
(E \to Z)_* \circ (E' \to E)_* \circ c'_p(Q|_E) \circ C' \\
& =
(E \to Z)_* \circ c'_p(Q|_E) \circ g_{\infty, *} \circ C' \\
& =
(E \to Z)_* \circ c'_p(Q|_E) \circ C \\
& =
c'_p(Q)
\end{align*}
第三个等式使用了 $c'_p(Q|_E)$ 是双变类，因而与真推前交换。
等式 $P'_p(Q) = P'_p(Q')$ 完全同理证明。
\end{proof}

\begin{lemma}
\label{chow-lemma-homomorphism}
在引理 \ref{chow-lemma-localized-chern-pre} 中，假设 $Q|_T$
同构于秩 $< p$ 的有限局部自由 $\mathcal{O}_T$-模。
记 $C \in A^0(W_\infty \to X)$ 为
引理 \ref{chow-lemma-gysin-at-infty} 的类。则
$$
C \circ c_p(Q|_{X \times \{0\}}) =
C \circ (Z \to X)_* \circ c'_p(Q) = c_p(Q|_{W_\infty}) \circ C
$$
\end{lemma}

\begin{proof}
第一个等式成立，因为由引理 \ref{chow-lemma-localized-chern-pre}，
$c_p(Q|_{X \times \{0\}}) =
(Z \to X)_* \circ c'_p(Q)$。第二个等式可逐个循环类证明
（见引理 \ref{chow-lemma-bivariant-zero}）。由于引理中的双变类构造
与基变换相容，可取某个 $\alpha \in \CH_k(X)$，并只须证明
$C \cap (Z \to X)_*(c'_p(Q) \cap \alpha) =
c_p(Q|_{W_\infty}) \cap C \cap \alpha$。注意
\begin{align*}
C \cap (Z \to X)_*(c'_p(Q) \cap \alpha)
& =
C \cap (Z \to X)_* (E \to Z)_*(c'_p(Q|_E) \cap C \cap \alpha) \\
& =
C \cap (W_\infty \to X)_*(E \to W_\infty)_*(c'_p(Q|_E) \cap C \cap \alpha) \\
& =
C \cap (W_\infty \to X)_*(E \to W_\infty)_*(c'_p(Q|_E) \cap i_\infty^*\beta) \\
& =
C \cap (W_\infty \to X)_*(c_p(Q|_{W_\infty}) \cap i_\infty^*\beta) \\
& =
C \cap (W_\infty \to X)_*i_\infty^*(c_p(Q) \cap \beta) \\
& =
i_\infty^*(c_p(Q) \cap \beta) \\
& =
c_p(Q|_{W_\infty}) \cap i_\infty^*\beta \\
& =
c_p(Q|_{W_\infty}) \cap C \cap \alpha
\end{align*}
这正是所需结论。第一个等式使用了
$c'_p(Q) = (E \to Z)_* \circ c'_p(Q|_E) \circ C$，其中
$E \subset W_\infty$ 是 $Z$ 的逆像，而 $c'_p(Q|_E)$ 是
引理 \ref{chow-lemma-silly} 所构造的类。第二个等式只是说
$E \to Z \to X$ 等于 $E \to W_\infty \to X$。
在第三个等式中，选择 $\beta \in \CH_{k + 1}(W)$，使其在
$b^{-1}(\mathbf{A}^1_X)$ 上的限制是 $\alpha$ 的平坦拉回，
从而按构造有 $C \cap \alpha = i_\infty^*\beta$。
第四个等式是引理 \ref{chow-lemma-silly-silly}。
第五个等式使用了 $c_p(Q)$ 是双变类，因而与 $i_\infty^*$ 交换。
第六个等式是引理 \ref{chow-lemma-homomorphism-pre}。
第七个等式再次使用了 $c_p(Q)$ 是双变类。
最后一个等式由 $C \cap \alpha = i_\infty^*\beta$ 得出。
\end{proof}

\begin{lemma}
\label{chow-lemma-homomorphism-commute}
在引理 \ref{chow-lemma-localized-chern-pre} 中，设 $Y \to X$
是局部有限型态射，并设 $c \in A^*(Y \to X)$ 是双变类。则在
$A^*(Y \times_X Z \to X)$ 中有
$$
P'_p(Q) \circ c = c \circ P'_p(Q)
\quad\text{分别}\quad
c'_p(Q) \circ c = c \circ c'_p(Q)
$$
\end{lemma}

\begin{proof}
设 $E \subset W_\infty$ 为 $Z$ 的逆像。回忆
$P'_p(Q) = (E \to Z)_* \circ P'_p(Q|_E) \circ C$，
分别地\ $c'_p(Q) = (E \to Z)_* \circ c'_p(Q|_E) \circ C$，
其中 $C$ 如引理 \ref{chow-lemma-gysin-at-infty}，而
$P'_p(Q|_E)$，分别地\ $c'_p(Q|_E)$，如引理 \ref{chow-lemma-silly}。
由引理 \ref{chow-lemma-gysin-at-infty-commutes}，$C$ 与 $c$ 交换；
由引理 \ref{chow-lemma-silly-commutes}，$P'_p(Q|_E)$，
分别地\ $c'_p(Q|_E)$，与 $c$ 交换。由于 $c$ 是双变类，
按定义它与沿 $E \to Z$ 的真推前交换。证明完毕。
\end{proof}

\begin{lemma}
\label{chow-lemma-localized-chern-pre-compose}
在引理 \ref{chow-lemma-localized-chern-pre} 中，假设 $Q|_T$ 为零。
在 $A^*(Z \to X)$ 中有
\begin{align*}
P'_1(Q) & = c'_1(Q), \\
P'_2(Q) & = c'_1(Q)^2 - 2c'_2(Q), \\
P'_3(Q) & = c'_1(Q)^3 - 3c'_1(Q)c'_2(Q) + 3c'_3(Q), \\
P'_4(Q) & = c'_1(Q)^4 - 4c'_1(Q)^2c'_2(Q) +
4c'_1(Q)c'_3(Q) + 2c'_2(Q)^2 - 4c'_4(Q),
\end{align*}
等等，其中乘法如注 \ref{chow-remark-ring-loc-classes}。
\end{lemma}

\begin{proof}
零层的秩 $< 1$，所以类 $c'_p(Q)$ 对所有 $p \geq 1$
都有定义，因而该陈述有意义。在
引理 \ref{chow-lemma-localized-chern-pre} 的证明中，我们用
引理 \ref{chow-lemma-gysin-at-infty} 的双变类
$C \in A^0(W_\infty \to X)$，以及引理 \ref{chow-lemma-silly}
给出的双变类 $P'_p(Q|_E)$ 和 $c'_p(Q|_E)$，构造了类
$P'_p(Q)$ 与 $c'_p(Q)$；这里 $Q|_E$ 是 $Q$ 在
$W_\infty$ 中 $Z$ 的逆像 $E$ 上的限制。由
引理 \ref{chow-lemma-silly-compose}，$P'_p(Q|_E)$ 与
$c'_p(Q|_E)$ 之间满足所需关系。回忆
$$
P'_p(Q) = (E \to Z)_* \circ P'_p(Q|_E) \circ C
\quad\text{且}\quad
c'_p(Q) = (E \to Z)_* \circ c'_p(Q|_E) \circ C
$$
为完成证明，只须说明注 \ref{chow-remark-ring-loc-classes}
在各类 $a_p = c'_p(Q)$ 上定义的乘法，与在各类
$b_p = c'_p(Q|_E)$ 上定义的乘法相容：
$$
a_{p_1}a_{p_2} \ldots a_{p_r} =
(E \to Z)_* \circ b_{p_1}b_{p_2} \ldots b_{p_r} \circ C
$$
省略一些细节。当 $r = 1$ 时这显然成立。若 $r > 1$，则由
注 \ref{chow-remark-res-push}，注 \ref{chow-remark-ring-loc-classes}
中的乘法是在乘积 $a_{p_1}a_{p_2} \ldots a_{p_r}$，分别地\ 
$b_{p_1}b_{p_2} \ldots b_{p_r}$ 的各因子之间插入
$(Z \to X)_*$，分别地\ $(E \to W_\infty)_*$，
然后取双变类的复合。由引理 \ref{chow-lemma-silly}，
$$
(E \to W_\infty)_* \circ b_{p_i} = c_{p_i}(Q|_{W_\infty})
$$
而由引理 \ref{chow-lemma-homomorphism}，
$$
C \circ (Z \to X)_* \circ a_{p_i} = c_{p_i}(Q|_{W_\infty}) \circ C
$$
对 $i = 2, \ldots, r$ 成立。计算表明，所需等式的两边均化简为
$$
(E \to Z)_* \circ c'_{p_1}(Q|_E) \circ
c_{p_2}(Q|_{W_\infty}) \circ \ldots \circ
c_{p_r}(Q|_{W_\infty}) \circ C
$$
证明完毕。
\end{proof}

\begin{lemma}
\label{chow-lemma-localized-chern-pre-sum-c}
在引理 \ref{chow-lemma-localized-chern-pre} 中，假设 $Q|_T$
同构于秩 $< p$ 的有限局部自由 $\mathcal{O}_T$-模。
再假设存在另一个陈类已定义的完美对象
$Q' \in D(\mathcal{O}_W)$，且 $Q'|_T$ 同构于置于上同调次数
$0$ 的、秩 $< p'$ 的有限局部自由 $\mathcal{O}_T$-模。
采用注 \ref{chow-remark-ring-loc-classes} 的记号，令
$$
c^{(p)}(Q) = 1 + c_1(Q|_{X \times \{0\}}) + \ldots +
c_{p - 1}(Q|_{X \times \{0\}}) +
c'_{p}(Q) + c'_{p + 1}(Q) + \ldots
$$
为 $A^{(p)}(Z \to X)$ 中的元素，其中对 $i \geq p$，
$c'_i(Q)$ 如引理 \ref{chow-lemma-localized-chern-pre}。
$c^{(p')}(Q')$ 与 $c^{(p + p')}(Q \oplus Q')$ 也类似定义。
则在 $A^{(p + p')}(Z \to X)$ 中有
$c^{(p + p')}(Q \oplus Q') = c^{(p)}(Q)c^{(p')}(Q')$。
\end{lemma}

\begin{proof}
回忆，对 $i \geq p$，$c'_i(Q)$ 在 $A^p(X)$ 中的像等于
$c_i(Q|_{X \times \{0\}})$；$Q'$ 和 $Q \oplus Q'$ 的情形也类似，
见引理 \ref{chow-lemma-localized-chern-pre}。因此次数
$< p + p'$ 的等式由引理 \ref{chow-lemma-additivity-on-perfect}
的可加性推出。

\medskip\noindent
取 $n \geq p + p'$。仿照引理 \ref{chow-lemma-localized-chern-pre}
的证明，记 $E \subset W_\infty$ 为 $Z$ 的逆像。注意，由
引理 \ref{chow-lemma-silly-sum-c}，在
$A^{(p + p')}(E \to W_\infty)$ 中有
$$
c^{(p + p')}(Q|_E \oplus Q'|_E) =
c^{(p)}(Q|_E)c^{(p')}(Q'|_E)
$$
又因按构造
$$
c'_p(Q \oplus Q') = (E \to Z)_* \circ c'_p(Q|_E \oplus Q'|_E) \circ C
$$
所以只须对所有 $i + j = n$ 证明
$$
(E \to Z)_* \circ c^{(p)}_i(Q|_E)c^{(p')}_j(Q'|_E) \circ C
=
c^{(p)}_i(Q)c^{(p')}_j(Q')
$$
在 $A^n(Z \to X)$ 中成立，其中两边的乘法均为
注 \ref{chow-remark-ring-loc-classes} 中的乘法。依照
$i \geq p$、$j \geq p'$ 是否成立，共有三种情形。

\medskip\noindent
假设 $i \geq p$ 且 $j \geq p'$。此时乘积的定义是在两个因子之间
插入 $(E \to W_\infty)_*$，分别地\ $(Z \to X)_*$，
再取双变类的复合；见注 \ref{chow-remark-res-push}。
换言之，只须证明
$$
(E \to Z)_* \circ c'_i(Q|_E) \circ
(E \to W_\infty)_* \circ c'_j(Q'|_E) \circ C =
c'_i(Q) \circ (Z \to X)_* \circ c'_j(Q')
$$
由引理 \ref{chow-lemma-silly}，左边等于
$$
(E \to Z)_* \circ c'_i(Q|_E) \circ c_j(Q'|_{W_\infty}) \circ C
$$
由于 $c'_i(Q) = (E \to Z)_* \circ c'_i(Q|_E) \circ C$，
右边等于
$$
(E \to Z)_* \circ c'_i(Q|_E) \circ C \circ (Z \to X)_* \circ c'_j(Q')
$$
由引理 \ref{chow-lemma-homomorphism}，这立即等于上述左边。

\medskip\noindent
假设 $i \geq p$ 且 $j < p$。展开此情形中的乘积，只须证明
$$
(E \to Z)_* \circ c'_i(Q|_E) \circ c_j(Q'|_{W_\infty}) \circ C =
c'_i(Q) \circ c_j(Q'|_{X \times \{0\}})
$$
再次使用 $c'_i(Q) = (E \to Z)_* \circ c'_i(Q|_E) \circ C$，
只须证明 $c_j(Q'|_{W_\infty}) \circ C =
C \circ c_j(Q'|_{X \times \{0\}})$；这是
引理 \ref{chow-lemma-homomorphism} 的一部分。

\medskip\noindent
假设 $i < p$ 且 $j \geq p'$。展开此情形中的乘积，只须证明
$$
(E \to Z)_* \circ c_i(Q|_E) \circ c'_j(Q'|_E) \circ C =
c_i(Q|_{Z \times \{0\}}) \circ c'_j(Q')
$$
然而，由于 $c'_j(Q|_E)$ 与 $c'_j(Q')$ 是双变类，
它们同与陈类取帽积交换
（引理 \ref{chow-lemma-cap-commutative-chern}）。因此只须证明
$$
(E \to Z)_* \circ c'_j(Q'|_E) \circ c_i(Q|_{W_\infty}) \circ C =
c'_j(Q') \circ c_i(Q|_{X \times \{0\}})
$$
这就约化到上一段讨论的情形。
\end{proof}

\begin{lemma}
\label{chow-lemma-localized-chern-pre-sum-P}
在引理 \ref{chow-lemma-localized-chern-pre} 中，假设 $Q|_T$ 为零。
再假设有另一个陈类已定义的完美对象
$Q' \in D(\mathcal{O}_W)$，且其限制 $Q'|_T$ 为零。
此时由引理 \ref{chow-lemma-localized-chern-pre} 构造的类
$P'_p(Q), P'_p(Q'), P'_p(Q \oplus Q') \in A^p(Z \to X)$
满足 $P'_p(Q \oplus Q') = P'_p(Q) + P'_p(Q')$。
\end{lemma}

\begin{proof}
立即由这些类的构造以及引理 \ref{chow-lemma-silly-sum-P} 得出。
\end{proof}









 
\section{局部化陈类}
\label{chow-section-localized-chern}

\noindent
先概述构造。设 $F$ 是域，$X$ 是 $F$ 上的簇，
$E$ 是 $D(\mathcal{O}_X)$ 中的完美对象，并设
$Z \subset X$ 是闭子概形，使得 $E|_{X \setminus Z} = 0$。
目标是构造元素
$$
c_p(Z \to X, E) \in A^p(Z \to X)
$$
为此将构造图
$$
\xymatrix{
W \ar[d]_f \ar[r]_q & X \\
\mathbf{P}^1_F
}
$$
以及 $D(\mathcal{O}_W)$ 中的完美对象 $Q$，使得
\begin{enumerate}
\item $f$ 平坦，且 $f$、$q$ 都是真态射；对
$t \in \mathbf{P}^1_F$，记 $W_t$ 为 $f$ 的纤维，
$q_t : W_t \to X$ 为 $q$ 的限制，并令 $Q_t = Q|_{W_t}$；
\item 对 $t \in \mathbf{A}^1_F$，$q_t : W_t \to X$
是同构且 $Q_t = q_t^*E$；
\item $q_\infty : W_\infty \to X$ 在 $X \setminus Z$ 上是同构；
\item 若 $T \subset W_\infty$ 是
$q_\infty^{-1}(X \setminus Z)$ 的闭包，则 $Q_\infty|_T$ 为零。
\end{enumerate}
其思想是把它看成由 $t \in \mathbf{P}^1$ 参数化的族
$\{(W_t, Q_t)\}$。当 $t \not = \infty$ 时，在
$Q_t = X$ 的 Chow 群上，$c_p(Q_t)$ 就是 $c_p(E)$。
而当 $t = \infty$ 时，由于 $Q_\infty|_T = 0$，
$c_p(Q_\infty)$ 把 $Q_\infty$ 上的类送到支撑在
$E = q_\infty^{-1}(Z)$ 上的类。把 $E$ 看作
$q_\infty : W_\infty \to X$ 的例外轨迹。由于任意
$\alpha \in \CH_*(X)$ 都产生一个循环“族”
$\alpha_t \in \CH_*(W_t)$，可将
$c_p(Z \to X, E) \cap \alpha$ 定义为推前
$(E \to Z)_*(c_p(Q_\infty) \cap \alpha_\infty)$。

\medskip\noindent
使该方案成立有两个主要要素：(1) $W$ 与 $Q$ 的构造是某种代数的
Macpherson 图构造；它在《平坦性进阶》第 \ref{flat-section-blowup-complexes-III} 节中完成。
(2) 给定 $W$ 与 $Q$ 后实际类的构造在第 \ref{chow-section-preparation-localized-chern-II} 节中完成，并依赖于第 \ref{chow-section-gysin-at-infty} 节和
\ref{chow-section-preparation-localized-chern}。

\begin{situation}
\label{chow-situation-loc-chern}
设 $(S, \delta)$ 如情形 \ref{chow-situation-setup}。
设 $X$ 是 $S$ 上局部有限型概形，$i : Z \to X$ 是闭浸入，
$E \in D(\mathcal{O}_X)$ 是对象，并设 $p \geq 0$。假设
\begin{enumerate}
\item $E$ 是 $D(\mathcal{O}_X)$ 中的完美对象；
\item 限制 $E|_{X \setminus Z}$ 为零，分别地\ 同构于置于上同调次数
$0$ 的、秩 $< p$ 的有限局部自由
$\mathcal{O}_{X \setminus Z}$-模；并且
\item 下列条件至少有一个\footnote{初次阅读时请忽略此技术条件；
相关讨论见注 \ref{chow-remark-loc-chern}。}成立：
(a) $X$ 拟紧；
(b) $X$ 的不可约分支拟紧；
(c) 存在表示 $E$ 的有限局部自由 $\mathcal{O}_X$-模组成的
局部有界复形；或者
(d) 存在 $S$ 上局部有限型概形的态射 $X \to X'$，使得
$E$ 是 $X'$ 上某个完美对象的拉回，且 $X'$ 的不可约分支拟紧。
\end{enumerate}
\end{situation}

\begin{lemma}
\label{chow-lemma-independent-loc-chern}
在情形 \ref{chow-situation-loc-chern} 中，存在典范双变类
$$
P_p(Z \to X, E) \in A^p(Z \to X),
\quad\text{分别}\quad
c_p(Z \to X, E) \in A^p(Z \to X)
$$
其性质为，在 $X$ 上作为双变类有
\begin{equation}
\label{chow-equation-defining-property-localized-classes}
i_* \circ P_p(Z \to X, E) = P_p(E),
\quad\text{分别}\quad
i_* \circ c_p(Z \to X, E) = c_p(E)
\end{equation}
（其中 $\circ$ 如引理 \ref{chow-lemma-push-proper-bivariant}）。
\end{lemma}

\begin{proof}
这些双变类的构造如下。令
$$
b : W \longrightarrow \mathbf{P}^1_X
\quad\text{且}\quad
T \longrightarrow W_\infty
\quad\text{且}\quad
Q
$$
分别为《平坦性进阶》第 \ref{flat-section-blowup-complexes-III} 节和引理 \ref{flat-lemma-graph-construction} 中构造的爆破、
$D(\mathcal{O}_W)$ 中的完美对象 $Q$ 以及闭浸入。
令 $T' \subset T$ 为开闭子概形，使得 $Q|_{T'}$ 为零，
分别地\ 同构于置于上同调次数 $0$ 的、秩 $< p$ 的有限局部自由
$\mathcal{O}_{T'}$-模。由情形 \ref{chow-situation-loc-chern}
的条件 (2)，态射
$$
T' \to T \to W_\infty \to X
$$
在 $X$ 的开子概形 $X \setminus Z$ 上都是概形同构。
下面检验 $Q$ 的陈类已定义。回忆按构造
$Q|_{X \times \{0\}} \cong E$，于是使用 $W, b, Q, T'$，
由引理 \ref{chow-lemma-localized-chern-pre} 构造的双变类给出
$$
P_p(Z \to X, E) = P'_p(Q) \in A^p(Z \to X)
$$
以及
$$
c_p(Z \to X, E) = c'_p(Q) \in A^p(Z \to X)
$$
并满足 (\ref{chow-equation-defining-property-localized-classes})。

\medskip\noindent
本段证明 $Q$ 的陈类已定义
（定义 \ref{chow-definition-defined-on-perfect}）；建议读者跳过。
若情形 \ref{chow-situation-loc-chern} 的假设 (3)(a) 或 (3)(b) 成立，
即 $X$ 具有拟紧不可约分支，则 $W$ 也如此（因为
$W \to X$ 是真态射）。所以由
引理 \ref{chow-lemma-chern-classes-defined}，$D(\mathcal{O}_W)$
中任意完美对象的陈类均已定义。若 (3)(c) 成立，即 $E$
可由有限局部自由模组成的局部有界复形表示，则由《平坦性进阶》引理 \ref{flat-lemma-graph-construction} 的第 (5) 部分，
对象 $Q$ 可由有限局部自由 $\mathcal{O}_W$-模组成的局部有界复形表示。
故 $Q$ 的陈类已定义。最后假设 (3)(d) 成立，即存在
$S$ 上局部有限型概形的态射 $X \to X'$，使得 $E$ 是
$X'$ 上完美对象 $E'$ 的拉回，且 $X'$ 的不可约分支拟紧。
令 $b' : W' \to \mathbf{P}^1_{X'}$ 和
$Q' \in D(\mathcal{O}_{W'})$ 为从三元组
$(\mathbf{P}^1_{X'}, (\mathbf{P}^1_{X'})_\infty, L(p')^*E')$
出发，按《平坦性进阶》第 \ref{flat-section-blowup-complexes-III} 节构造的态射与完美对象。
由上述讨论，$Q'$ 的陈类已定义。由于 $b$ 和 $b'$ 是应用
《平坦性进阶》引理 \ref{flat-lemma-complex-and-divisor-blowup-pre} 构造的，
由《平坦性进阶》引理 \ref{flat-lemma-complex-and-divisor-blowup-base-change}，
存在态射 $W \to W'$，使得 $Q = L(W \to W')^*Q'$。
于是由引理 \ref{chow-lemma-chern-classes-defined}，
$Q$ 的陈类已定义。
\end{proof}

\begin{definition}
\label{chow-definition-localized-chern}
设 $(S, \delta)$、$X$、$E \in D(\mathcal{O}_X)$ 以及
$i : Z \to X$ 如情形 \ref{chow-situation-loc-chern}。
\begin{enumerate}
\item 若限制 $E|_{X \setminus Z}$ 为零，则对所有 $p \geq 0$，
由引理 \ref{chow-lemma-independent-loc-chern} 的构造定义
$$
P_p(Z \to X, E) \in A^p(Z \to X)
$$
并用公式
$$
ch(Z \to X, E) =
\sum\nolimits_{p = 0, 1, 2, \ldots} \frac{P_p(Z \to X, E)}{p!}
\quad\text{于}\quad \prod\nolimits_{p \geq 0} A^p(Z \to X) \otimes \mathbf{Q}
$$
定义{\it 局部化陈特征}。
\item 若限制 $E|_{X \setminus Z}$ 同构于置于上同调次数
$0$ 的、秩 $< p$ 的有限局部自由
$\mathcal{O}_{X \setminus Z}$-模，则由
引理 \ref{chow-lemma-independent-loc-chern} 的构造定义
{\it 局部化第 $p$ 陈类} $c_p(Z \to X, E)$。
\end{enumerate}
\end{definition}

\noindent
在该定义的情形中，假设 $E|_{X \setminus Z}$ 为零。
那么确有等式
$$
i_* \circ ch(Z \to X, E) = ch(E)
$$
在 $A^*(X) \otimes \mathbf{Q}$ 中成立，因为上面已经证明了等式
(\ref{chow-equation-defining-property-localized-classes})。

\medskip\noindent
下面给出一个重要的一致性检验。

\begin{lemma}
\label{chow-lemma-base-change-loc-chern}
在情形 \ref{chow-situation-loc-chern} 中，设
$f : X' \to X$ 是局部有限型概形态射。记
$E' = f^*E$ 且 $Z' = f^{-1}(Z)$。则定义 \ref{chow-definition-localized-chern} 中的双变类
$$
P_p(Z' \to X', E') \in A^p(Z' \to X'),
\quad\text{分别}\quad
c_p(Z' \to X', E') \in A^p(Z' \to X')
$$
即使用 $X', Z', E'$ 按引理 \ref{chow-lemma-independent-loc-chern}
构造的类，是双变类 $P_p(Z \to X, E) \in A^p(Z \to X)$，分别地\ 
$c_p(Z \to X, E) \in A^p(Z \to X)$，的限制
（注 \ref{chow-remark-restriction-bivariant}）。
\end{lemma}

\begin{proof}
记 $p : \mathbf{P}^1_X \to X$ 与
$p' : \mathbf{P}^1_{X'} \to X'$ 为结构态射。回忆
$b : W \to \mathbf{P}^1_X$ 与
$b' : W' \to \mathbf{P}^1_{X'}$ 是在《平坦性进阶》引理 \ref{flat-lemma-complex-and-divisor-blowup-pre} 中分别从三元组
$(\mathbf{P}^1_X, (\mathbf{P}^1_X)\infty, p^*E)$ 与
$(\mathbf{P}^1_{X'}, (\mathbf{P}^1_{X'})\infty, (p')^*E')$
构造的态射。此外
$Q = L\eta_{\mathcal{I}_\infty}p^*E$ 且
$Q = L\eta_{\mathcal{I}'_\infty}(p')^*E'$，其中
$\mathcal{I}_\infty \subset \mathcal{O}_W$ 是 $W_\infty$ 的理想层，
$\mathcal{I}'_\infty \subset \mathcal{O}_{W'}$ 是
$W'_\infty$ 的理想层。其次，
$h : \mathbf{P}^1_{X'} \to \mathbf{P}^1_X$ 是概形态射，使得有效
Cartier 除子 $(\mathbf{P}^1_X)_\infty$ 的拉回是有效 Cartier 除子
$(\mathbf{P}^1_{X'})_\infty$，且 $h^*p^*E = (p')^*E'$。
由《平坦性进阶》引理 \ref{flat-lemma-complex-and-divisor-blowup-base-change}，得到交换图
$$
\xymatrix{
W' \ar[rd]_{b'} \ar[r]_-g &
\mathbf{P}^1_{X'} \times_{\mathbf{P}^1_X} W \ar[d]_r \ar[r]_-q &
W \ar[d]^b \\
&
\mathbf{P}^1_{X'} \ar[r] &
\mathbf{P}^1_X
}
$$
其中 $W'$ 是 $\mathbf{P}^1_{X'}$ 关于 $b$ 的“严格变换”，
且 $Q' = (q \circ g)^*Q$。回忆
$P_p(Z \to X, E) = P'_p(Q)$，分别地\ 
$c_p(Z \to X, E) = c'_p(Q)$，其中 $P'_p(Q)$，分别地\ $c'_p(Q)$，
是在引理 \ref{chow-lemma-localized-chern-pre} 中使用 $b, Q, T'$ 构造的；
这里 $T'$ 是闭子概形 $T' \subset W_\infty$，满足以下两个性质：
(a) $T'$ 包含 $W_\infty$ 中所有位于 $X \setminus Z$ 上方的点；
(b) $Q|_{T'}$ 为零，分别地\ 同构于置于次数 $0$ 的、秩
$< p$ 的有限局部自由模。在引理 \ref{chow-lemma-localized-chern-pre} 的构造中，我们选定了一个具有
(a)、(b) 性质的闭子概形 $T'$，但 $T'$ 的具体选择无关紧要；见
引理 \ref{chow-lemma-localized-chern-pre-independent}。

\medskip\noindent
其次，由引理 \ref{chow-lemma-base-change-localized-chern-pre}，
双变类 $P_p(Z \to X, E) = P'_p(Q)$，分别地\ 
$c_p(Z \to X, E) = c_p(Q')$，在 $X'$ 上的限制对应于类
$P'_p(q^*Q)$，分别地\ $c'_p(q^*Q)$；后者是在
引理 \ref{chow-lemma-localized-chern-pre} 中使用
$r : \mathbf{P}^1_{X'} \times_{\mathbf{P}^1_X} W \to \mathbf{P}^1_{X'}$、
复形 $q^*Q$ 及逆像 $q^{-1}(T')$ 构造的。

\medskip\noindent
现在由引理 \ref{chow-lemma-localized-chern-pre-independent}
的第二个陈述，$P'_p(Q') = P'_p(q^*Q)$，分别地\ 
$c'_p(q^*Q) = c'_p(Q')$。又因
$P_p(Z' \to X', E') = P'_p(Q')$，分别地\ 
$c_p(Z' \to X', E') = c'_p(Q')$，引理得证。
\end{proof}

\begin{remark}
\label{chow-remark-loc-chern}
在情形 \ref{chow-situation-loc-chern} 中，更自然的做法是把假设 (3)
替换为假设“$E$ 的陈类已定义”。事实上，把引理 \ref{chow-lemma-independent-loc-chern} 和
\ref{chow-lemma-base-change-loc-chern} 与
引理 \ref{chow-lemma-envelope-bivariant} 结合起来，很容易把定义推广到
这个略微更一般的情形。若以后需要，我们将在此完成该推广。
\end{remark}

\begin{lemma}
\label{chow-lemma-loc-chern-after-pushforward}
在情形 \ref{chow-situation-loc-chern} 中，对任意
$\alpha \in \CH_*(Z)$，在 $\CH_*(Z)$ 中有
$$
P_p(Z \to X, E) \cap i_*\alpha = P_p(E|_Z) \cap \alpha,
\quad\text{分别}\quad
c_p(Z \to X, E) \cap i_*\alpha = c_p(E|_Z) \cap \alpha
$$
\end{lemma}

\begin{proof}
只证明第二个等式，第一个的证明从略。由于
$c_p(Z \to X, E)$ 是双变类，且 $Z \to X$ 沿 $Z \to X$
的基变换是 $\text{id} : Z \to Z$，有
$c_p(Z \to X, E) \cap i_*\alpha = c_p(Z \to X, E) \cap \alpha$。
由引理 \ref{chow-lemma-base-change-loc-chern}，
$c_p(Z \to X, E)$ 在 $Z$ 上的限制（！）是
$\text{id} : Z \to Z$ 与 $E|_Z$ 的局部化陈类。
因此取 $X = Z$，结论由
(\ref{chow-equation-defining-property-localized-classes}) 得出。
\end{proof}

\begin{lemma}
\label{chow-lemma-loc-chern-disjoint}
在情形 \ref{chow-situation-loc-chern} 中，若
$\alpha \in \CH_k(X)$ 的支撑与 $Z$ 不交，则
$P_p(Z \to X, E) \cap \alpha = 0$，分别地\ 
$c_p(Z \to X, E) \cap \alpha = 0$。
\end{lemma}

\begin{proof}
这立即由局部化陈类的构造得出。也可先把
$c_p(Z \to X, E)$ 限制到 $\alpha$ 的支撑，再用
引理 \ref{chow-lemma-base-change-loc-chern} 看出该限制为零，
从而计算 $c_p(Z \to X, E) \cap \alpha$。
\end{proof}

\begin{lemma}
\label{chow-lemma-loc-chern-shrink-Z}
在情形 \ref{chow-situation-loc-chern} 中，假设
$Z \subset Z' \subset X$，其中 $Z'$ 是 $X$ 的闭子概形。则
$P_p(Z' \to X, E) = (Z \to Z')_* \circ P_p(Z \to X, E)$，
分别地\ $c_p(Z' \to X, E) = (Z \to Z')_* \circ c_p(Z \to X, E)$
（其中 $\circ$ 如引理 \ref{chow-lemma-push-proper-bivariant}）。
\end{lemma}

\begin{proof}
引理 \ref{chow-lemma-independent-loc-chern} 中
$P_p(Z' \to X, E)$，分别地\ $c_p(Z' \to X, E)$ 的构造，
使用完全相同的态射 $b : W \to \mathbf{P}^1_X$ 与
$D(\mathcal{O}_W)$ 中的完美对象 $Q$。于是应用
引理 \ref{chow-lemma-silly-shrink} 即得结论。省略一些细节。
\end{proof}

\begin{lemma}
\label{chow-lemma-loc-chern-agree}
在引理 \ref{chow-lemma-silly} 中，设 $E_2$ 是某个完美对象
$E \in D(\mathcal{O}_X)$ 的限制，且该对象在 $X_1$ 上的限制为零，
分别地\ 同构于置于上同调次数 $0$ 的、秩 $< p$ 的有限局部自由
$\mathcal{O}_{X_1}$-模。则引理 \ref{chow-lemma-silly} 的类
$P'_p(E_2)$，分别地\ $c'_p(E_2)$，与定义 \ref{chow-definition-localized-chern} 的
$P_p(X_2 \to X, E)$，分别地\ $c_p(X_2 \to X, E)$ 相同，
只要 $E$ 满足情形 \ref{chow-situation-loc-chern} 的假设 (3)。
\end{lemma}

\begin{proof}
关于 $E$ 的假设蕴含存在包含 $X_1$ 的开集 $U \subset X$，
使得 $E|_U$ 为零，分别地\ 同构于秩 $< p$ 的有限局部自由
$\mathcal{O}_U$-模。见《代数进阶》引理 \ref{more-algebra-lemma-lift-perfect-from-residue-field}。
令 $Z \subset X$ 为 $U$ 在 $X$ 中的补集，并赋予诱导的既约闭子概形结构。
则由引理 \ref{chow-lemma-loc-chern-shrink-Z}，
$P_p(X_2 \to X, E) = (Z \to X_2)_* \circ P_p(Z \to X, E)$，
分别地\ $c_p(X_2 \to X, E) = (Z \to X_2)_* \circ c_p(Z \to X, E)$。
现在可证明 $P_p(X_2 \to X, E)$，分别地\ $c_p(X_2 \to X, E)$，
满足引理 \ref{chow-lemma-silly} 给出的对 $P'_p(E_2)$，分别地\ 
$c'_p(E_2)$ 的刻画。事实上，由刚刚证明的关系
$P_p(X_2 \to X, E) = (Z \to X_2)_* \circ P_p(Z \to X, E)$，
分别地\ $c_p(X_2 \to X, E) = (Z \to X_2)_* \circ c_p(Z \to X, E)$，
以及 $X_1 \cap Z = \emptyset$，由
引理 \ref{chow-lemma-loc-chern-disjoint}，复合
$P_p(X_2 \to X, E) \circ i_{1, *}$，分别地\ 
$c_p(X_2 \to X, E) \circ i_{1, *}$ 为零。另一方面，由
引理 \ref{chow-lemma-loc-chern-after-pushforward}，
$P_p(X_2 \to X, E) \circ i_{2, *} = P_p(E_2)$，
分别地\ $c_p(X_2 \to X, E) \circ i_{2, *} = c_p(E_2)$。
\end{proof}





\section{两个技术引理}
\label{chow-section-tools-loc-chern}

\noindent
本节建立一些附加工具，以便更方便地处理局部化陈类。
下一个引理是我们在引理 \ref{chow-lemma-localized-chern-pre-compose}
和 \ref{chow-lemma-localized-chern-pre-sum-c} 的证明中已经遇到的
某个事实的更精确版本。

\begin{lemma}
\label{chow-lemma-homomorphism-final}
设 $(S, \delta)$ 如情形 \ref{chow-situation-setup}。
设 $X$ 在 $S$ 上局部有限型，并设
$b : W \longrightarrow \mathbf{P}^1_X$ 是真概形态射。
设 $n \geq 1$。对 $i = 1, \ldots, n$，设
$Z_i \subset X$ 是闭子概形，$Q_i \in D(\mathcal{O}_W)$
是完美对象，$p_i \geq 0$ 是整数，并设
$T_i \subset W_\infty$（$i = 1, \ldots, n$）是闭子概形。
记 $W_i = b^{-1}(\mathbf{P}^1_{Z_i})$。假设
\begin{enumerate}
\item 对 $i = 1, \ldots, n$，数据 $b, Z_i, Q_i, T_i, p_i$
满足引理 \ref{chow-lemma-localized-chern-pre} 的假设；
\item $Q_i|_{W \setminus W_i}$ 为零，分别地\ 同构于置于上同调次数
$0$ 的、秩 $< p_i$ 的有限局部自由模；
\item $W$ 上的 $Q_i$ 满足情形 \ref{chow-situation-loc-chern}
的假设 (3)。
\end{enumerate}
则 $P'_{p_n}(Q_n) \circ \ldots \circ P'_{p_1}(Q_1)$ 等于
$$
(W_{n, \infty} \cap \ldots \cap W_{1, \infty} \to
Z_n \cap \ldots \cap Z_1)_* \circ
P'_{p_n}(Q_n|_{W_{n, \infty}}) \circ \ldots \circ P'_{p_1}(Q_1|_{W_{1, \infty}})
\circ C
$$
这是 $A^{p_n + \ldots + p_1}(Z_n \cap \ldots \cap Z_1 \to X)$
中的等式；分别地\ $c'_{p_n}(Q_n) \circ \ldots \circ c'_{p_1}(Q_1)$ 等于
$$
(W_{n, \infty} \cap \ldots \cap W_{1, \infty} \to
Z_n \cap \ldots \cap Z_1)_* \circ
c'_{p_n}(Q_n|_{W_{n, \infty}}) \circ \ldots \circ c'_{p_1}(Q_1|_{W_{1, \infty}})
\circ C
$$
这是 $A^{p_n + \ldots + p_1}(Z_n \cap \ldots \cap Z_1 \to X)$
中的等式。
\end{lemma}

\begin{proof}
对 $n$ 归纳证明关于陈类的陈述；关于 $P_p(-)$ 的陈述完全同理证明。
当 $n = 1$ 时，陈述就是 $c'_{p_1}(Q_1)$ 的构造，因为
$W_{1, \infty}$ 是 $Z_1$ 在 $W_\infty$ 中的逆像。
当 $n > 1$ 时，由归纳假设，
$c'_{p_n}(Q_n) \circ \ldots \circ c'_{p_1}(Q_1)$ 等于
$$
c'_{p_n}(Q_n) \circ
(W_{n - 1, \infty} \cap \ldots \cap W_{1, \infty} \to
Z_{n - 1} \cap \ldots \cap Z_1)_* \circ
c'_{p_{n - 1}}(Q_{n - 1}|_{W_{n - 1}, \infty}) \circ \ldots \circ
c'_{p_1}(Q_1|_{W_{1, \infty}})
\circ C
$$
由引理 \ref{chow-lemma-base-change-localized-chern-pre}，
$c'_{p_n}(Q_n)$ 在 $Z_{n - 1} \cap \ldots \cap Z_1$ 上的限制，
由闭子集 $Z_n \cap \ldots \cap Z_1$、态射
$b' : W_{n - 1} \cap \ldots \cap W_1 \to
\mathbf{P}^1_{Z_{n - 1} \cap \ldots \cap Z_1}$，以及
$Q_n$ 在 $W_{n - 1} \cap \ldots \cap W_1$ 上的限制计算。
注意 $(b')^{-1}(Z_n) = W_n \cap \ldots \cap W_1$，并且
$(W_n \cap \ldots \cap W_1)_\infty =
W_{n, \infty} \cap \ldots \cap W_{1, \infty}$。
记 $C_{n - 1} \in A^0(W_{n - 1, \infty} \cap \ldots \cap W_{1, \infty} \to
Z_{n - 1} \cap \ldots \cap Z_1)$ 为引理 \ref{chow-lemma-gysin-at-infty}
的类。于是 $c'_{p_n}(Q_n)$ 在
$Z_{n - 1} \cap \ldots \cap Z_1$ 上的限制为
\begin{align*}
&
(W_{n, \infty} \cap \ldots \cap W_{1, \infty} \to Z_n \cap \ldots \cap Z_1)_*
\circ
c'_{p_n}(Q_n|_{(W_n \cap \ldots \cap W_1)_\infty})
\circ
C_{n - 1} \\
& =
(W_{n, \infty} \cap \ldots \cap W_{1, \infty} \to Z_n \cap \ldots \cap Z_1)_*
\circ
c'_{p_n}(Q_n|_{W_{n, \infty}})
\circ
C_{n - 1}
\end{align*}
其中等式由引理 \ref{chow-lemma-base-change-silly} 得出
（右边的限制记号从略）。因此上式化为
\begin{align*}
(W_{n, \infty} \cap \ldots \cap W_{1, \infty} \to
Z_n \cap \ldots \cap Z_1)_* \circ
c'_{p_n}(Q_n|_{W_n, \infty}) \circ \\
C_{n - 1} \circ
(W_{n - 1, \infty} \cap \ldots \cap W_{1, \infty} \to
Z_{n - 1} \cap \ldots \cap Z_1)_* \\
\circ
c'_{p_{n - 1}}(Q_{n - 1}|_{W_{n - 1}, \infty}) \circ \ldots \circ
c'_{p_1}(Q_1|_{W_{1, \infty}})
\circ C
\end{align*}
由引理 \ref{chow-lemma-homomorphism-pre}，复合
$C_{n - 1} \circ (W_{n - 1, \infty} \cap \ldots \cap W_{1, \infty} \to
Z_{n - 1} \cap \ldots \cap Z_1)_*$
在 Gysin 映射
$$
(W_{n - 1, \infty} \cap \ldots \cap W_{1, \infty} \to
W_{n - 1} \cap \ldots \cap W_1)^*
$$
的像中的元素上是恒等映射。因此只须证明
$c'_{p_{n - 1}}(Q_{n - 1}|_{W_{n - 1}, \infty}) \circ \ldots \circ
c'_{p_1}(Q_1|_{W_{1, \infty}}) \circ C$
的像中每个元素都在该 Gysin 映射的像中。由
引理 \ref{chow-lemma-loc-chern-agree} 以及引理陈述中关于 $Q_i$
的假设 (2)、(3)，可以写成
$$
c'_{p_i}(Q_i|_{W_{i, \infty}}) = c_{p_i}(W_i \to W, Q_i)\text{ 在 } W_{i, \infty}\text{ 上的限制}
$$
因此，若 $\beta \in \CH_{k + 1}(W)$ 在
$b^{-1}(\mathbf{A}^1_X)$ 上限制为 $\alpha$ 的平坦拉回，则
\begin{align*}
& c'_{p_{n - 1}}(Q_{n - 1}|_{W_{n - 1}, \infty}) \cap \ldots \cap
c'_{p_1}(Q_1|_{W_{1, \infty}})
\cap C \cap \alpha \\
& =
c'_{p_{n - 1}}(Q_{n - 1}|_{W_{n - 1}, \infty}) \cap \ldots \cap
c'_{p_1}(Q_1|_{W_{1, \infty}})
\cap i_\infty^* \beta \\
& =
c_{p_{n - 1}}(W_{n - 1} \to W, Q_{n - 1}) \cap \ldots \cap
c_{p_{n - 1}}(W_1 \to W, Q_1) \cap i_\infty^* \beta \\
& =
(W_{n - 1, \infty} \cap \ldots \cap W_{1, \infty} \to
W_{n - 1} \cap \ldots \cap W_1)^*
\left(c_{p_{n - 1}}(W_{n - 1} \to W, Q_{n - 1}) \cap \ldots \cap
c_{p_1}(W_1 \to W, Q_1) \cap \beta\right)
\end{align*}
如所需。最后一个等式使用了
$c_{p_i}(W_i \to W, Q_i)$ 是双变类，因而按定义与
$i_\infty^*$ 交换。
\end{proof}

\noindent
若要计算复形的局部化陈类，下面的引理提供了极大的灵活性。

\begin{lemma}
\label{chow-lemma-independent-loc-chern-bQ}
假设 $(S, \delta), X, Z, b : W \to \mathbf{P}^1_X, Q, T, p$
满足引理 \ref{chow-lemma-localized-chern-pre} 的假设。
设 $F \in D(\mathcal{O}_X)$ 是完美对象，使得
\begin{enumerate}
\item $Q$ 在 $b^{-1}(\mathbf{A}^1_X)$ 上的限制同构于 $F$ 的拉回；
\item $F|_{X \setminus Z}$ 为零，分别地\ 同构于置于上同调次数
$0$ 的、秩 $< p$ 的有限局部自由
$\mathcal{O}_{X \setminus Z}$-模；并且
\item $W$ 上的 $Q$ 与 $X$ 上的 $F$ 满足
情形 \ref{chow-situation-loc-chern} 的假设 (3)。
\end{enumerate}
则引理 \ref{chow-lemma-localized-chern-pre} 所构造的
$A^p(Z \to X)$ 中的类 $P'_p(Q)$，分别地\ $c'_p(Q)$，
等于定义 \ref{chow-definition-localized-chern} 中的
$P_p(Z \to X, F)$，分别地\ $c_p(Z \to X, F)$。
\end{lemma}

\begin{proof}
这些假设在沿局部有限型态射 $X' \to X$ 作基变换后仍然成立。
因此只须对任意 $\alpha \in \CH_k(X)$ 证明
$P_p(Z \to X, F) \cap \alpha = P'_p(Q) \cap \alpha$，
分别地\ $c_p(Z \to X, F) \cap \alpha = c'_p(Q) \cap \alpha$。
依照引理 \ref{chow-lemma-gysin-at-infty} 中 $C$ 的构造，
选择 $\beta \in \CH_{k + 1}(W)$，使其在
$b^{-1}(\mathbf{A}^1_X)$ 上的限制等于 $\alpha$ 的平坦拉回。
记 $W' = b^{-1}(Z)$，并记
$E = W'_\infty \subset W_\infty$ 为 $Z$ 经
$W_\infty \to X$ 的逆像。引理由以下等式链推出
（$P_p$ 的情形类似）：
\begin{align*}
c'_p(Q) \cap \alpha
& =
(E \to Z)_*(c'_p(Q|_E) \cap i_\infty^*\beta) \\
& =
(E \to Z)_*(c_p(E \to W_\infty, Q|_{W_\infty}) \cap i_\infty^*\beta) \\
& =
(W'_\infty \to Z)_*(c_p(W' \to W, Q) \cap i_\infty^*\beta) \\
& =
(W'_\infty \to Z)_*((i'_\infty)^*(c_p(W' \to W, Q) \cap \beta)) \\
& =
(W'_\infty \to Z)_*((i'_\infty)^*(c_p(Z' \to X, F) \cap \beta)) \\
& =
(W'_0 \to Z)_*((i'_0)^*(c_p(Z' \to X, F) \cap \beta)) \\
& =
(W'_0 \to Z)_*(c_p(Z' \to X, F) \cap i_0^*\beta)) \\
& =
c_p(Z \to X, F) \cap \alpha
\end{align*}
第一个等式是引理 \ref{chow-lemma-localized-chern-pre} 中
$c'_p(Q)$ 的构造。第二个等式是引理 \ref{chow-lemma-loc-chern-agree}。
$W' \to W$ 沿 $W_\infty \to W$ 的基变换是态射
$E = W'_\infty \to W_\infty$，所以第三个等式由
引理 \ref{chow-lemma-base-change-loc-chern} 得出。第四个等式中
$i'_\infty : W'_\infty \to W'$ 是包含态射；该等式由
$c_p(W' \to W, Q)$ 是双变类这一事实得出。对于第五个等式，
注意由引理的假设 (1)，$Q$ 与 $F$ 在
$D(\mathcal{O}_{b^{-1}(\mathbf{A}^1_X)})$ 中限制为同一个对象；
应用引理 \ref{chow-lemma-base-change-loc-chern}，
$c_p(W' \to W, Q)$ 与 $c_p(Z' \to X, F)$ 在
$A^p((b')^{-1} \to b^{-1}(\mathbf{A}^1_X))$ 中限制为同一个双变类。
由于 $(i'_\infty)^*$ 消去支撑在 $W'_\infty$ 上的循环
（见注 \ref{chow-remark-gysin-on-cycles}），第五个等式成立。
第六个等式成立，因为 $W'_\infty$ 与 $W'_0$ 是
$\mathbf{P}^1_Z$ 中有理等价的有效 Cartier 除子
$D_0, D_\infty$ 的拉回，所以 $i_\infty^*\beta$ 与
$i_0^*\beta$ 映到 $W'$ 上同一个循环类；确切地说，由
引理 \ref{chow-lemma-support-cap-effective-Cartier}，二者均表示类
$c_1(\mathcal{O}_{\mathbf{P}^1_Z}(1)) \cap c_p(Z \to X, F_) \cap \beta$。
第七个等式成立，因为 $c_p(Z \to X, F)$ 是双变类。
按构造 $W'_0 = Z$ 且 $i_0^*\beta = \alpha$，这说明最后一个等式成立。
\end{proof}







\section{局部化陈类的性质}
\label{chow-section-properties-loc-chern}

\noindent
本节的主要结果是局部化陈类的可加性和乘法性。

\begin{lemma}
\label{chow-lemma-loc-chern-character}
在情形 \ref{chow-situation-loc-chern} 中，假设
$E|_{X \setminus Z}$ 为零。则
\begin{align*}
P_1(Z \to X, E) & = c_1(Z \to X, E), \\
P_2(Z \to X, E) & = c_1(Z \to X, E)^2 - 2c_2(Z \to X, E), \\
P_3(Z \to X, E) & = c_1(Z \to X, E)^3 - 3c_1(Z \to X, E)c_2(Z \to X, E)
+ 3c_3(Z \to X, E),
\end{align*}
依此类推，其中乘积取自注 \ref{chow-remark-ring-loc-classes}
中的代数 $A^{(1)}(Z \to X)$。
\end{lemma}

\begin{proof}
该陈述是有意义的，因为零层的秩 $< 1$，所以对所有
$p \geq 1$，类 $c_p(Z \to X, E)$ 都有定义。
结论本身立即由更一般的引理 \ref{chow-lemma-localized-chern-pre-compose} 得出，因为局部化陈类
是使用引理 \ref{chow-lemma-localized-chern-pre} 的步骤、在
第 \ref{chow-section-localized-chern} 节中定义的。
\end{proof}

\begin{lemma}
\label{chow-lemma-loc-chern-classes-commute}
在情形 \ref{chow-situation-loc-chern} 中，设
$Y \to X$ 局部有限型，且 $c \in A^*(Y \to X)$。则
$$
P_p(Z \to X, E) \circ c = c \circ P_p(Z \to X, E),
$$
分别地
$$
c_p(Z \to X, E) \circ c = c \circ c_p(Z \to X, E)
$$
在 $A^*(Y \times_X Z \to X)$ 中成立。
\end{lemma}

\begin{proof}
这由引理 \ref{chow-lemma-homomorphism-commute} 得出。
更确切地，令
$$
b : W \to \mathbf{P}^1_X
\quad\text{且}\quad
Q
\quad\text{且}\quad
T' \subset T \subset W_\infty
$$
如引理 \ref{chow-lemma-independent-loc-chern} 的证明中所取。
按定义，作为双变运算，
$c_p(Z \to X, E) = c'_p(Q)$，其中右边是在
引理 \ref{chow-lemma-localized-chern-pre} 中使用
$W, b, Q, T'$ 构造的双变类。由
引理 \ref{chow-lemma-homomorphism-commute}，在
$A^*(Y \times_X Z \to X)$ 中有
$P'_p(Q) \circ c = c \circ P'_p(Q)$，分别地\ 
$c'_p(Q) \circ c = c \circ c'_p(Q)$，结论随之成立。
\end{proof}

\begin{remark}
\label{chow-remark-loc-chern-classes}
在情形 \ref{chow-situation-loc-chern} 中，方便起见定义
$$
c^{(p)}(Z \to X, E) = 1 + c_1(E) + \ldots + c_{p - 1}(E) +
c_p(Z \to X, E) + c_{p + 1}(Z \to X, E) + \ldots
$$
将其视为注 \ref{chow-remark-ring-loc-classes} 中所考虑的代数
$A^{(p)}(Z \to X)$ 的元素。
\end{remark}

\begin{lemma}
\label{chow-lemma-additivity-loc-chern-c}
设 $(S, \delta)$ 如情形 \ref{chow-situation-setup}。
设 $X$ 在 $S$ 上局部有限型，且 $Z \to X$ 是闭浸入。设
$$
E_1 \to E_2 \to E_3 \to E_1[1]
$$
是 $D(\mathcal{O}_X)$ 中完美对象组成的可区别三角。假设
\begin{enumerate}
\item 限制 $E_1|_{X \setminus Z}$ 和 $E_3|_{X \setminus Z}$
分别同构于置于次数 $0$ 的、秩 $< p_1$ 和 $< p_3$ 的有限局部自由
$\mathcal{O}_{X \setminus Z}$-模；并且
\item 下列条件至少有一个成立：
(a) $X$ 拟紧；
(b) $X$ 的不可约分支拟紧；
(c) $E_3 \to E_1[1]$ 可由有限局部自由
$\mathcal{O}_X$-模组成的局部有界复形之间的映射表示；或者
(d) 存在包络 $f : Y \to X$，使得
$Lf^*E_3 \to Lf^*E_1[1]$ 可由有限局部自由
$\mathcal{O}_Y$-模组成的局部有界复形之间的映射表示。
\end{enumerate}
采用注 \ref{chow-remark-loc-chern-classes} 中的记号，在
$A^{(p_1 + p_3)}(Z \to X)$ 中有
$$
c^{(p_1 + p_3)}(Z \to X, E_2) = c^{(p_1)}(Z \to X, E_1)c^{(p_3)}(Z \to X, E_3)
$$
\end{lemma}

\begin{proof}
注意，这些假设蕴含 $E_2|_{X \setminus Z}$ 为零，
分别地\ 同构于秩 $< p_1 + p_3$ 的有限局部自由
$\mathcal{O}_{X \setminus Z}$-模。因此该陈述是有意义的。

\medskip\noindent
设 $f : Y \to X$ 是包络。展开引理中公式的左右两边可见，
需要证明 $A^*(X)$ 和 $A^*(Z \to X)$ 中若干类的等式。
由引理 \ref{chow-lemma-envelope-bivariant} 中的唯一性，只须证明
$A^*(Y)$ 和 $A^*(Z \to Y)$ 中对应的关系。又因所涉及各类的构造
与基变换相容（引理 \ref{chow-lemma-base-change-loc-chern}），
可将 $X$ 替换为 $Y$，并将可区别三角替换为其拉回。

\medskip\noindent
在引理 \ref{chow-lemma-additivity-on-perfect} 的证明中已经看到，
条件 (2)(a)、(2)(b) 和 (2)(c) 都蕴含条件 (2)(d)。
结合上一段的讨论，问题化归为下一段所述的情形。

\medskip\noindent
设 $\varphi^\bullet : \mathcal{E}_3^\bullet[-1] \to \mathcal{E}_1^\bullet$
是有限局部自由 $\mathcal{O}_X$-模组成的局部有界复形之间的映射，
它在导出范畴中表示映射 $E_3[-1] \to E_1$。考虑概形
$X' = \mathbf{A}^1 \times X$ 及投影
$g : X' \to X$。令 $Z' = g^{-1}(Z) = \mathbf{A}^1 \times Z$。
以 $t$ 表示 $\mathbf{A}^1$ 上的坐标。考虑 $X'$ 上复形映射
$$
t g^*\varphi^\bullet :
g^*\mathcal{E}_3^\bullet[-1]
\longrightarrow
g^*\mathcal{E}_1^\bullet
$$
的锥 $\mathcal{C}^\bullet$。得到可区别三角
$$
g^*\mathcal{E}_1^\bullet \to \mathcal{C}^\bullet \to
g^*\mathcal{E}_3^\bullet \to g^*\mathcal{E}_1^\bullet[1]
$$
其中前三项逐项组成分裂短正合复形列。显然
$\mathcal{C}^\bullet$ 是有限局部自由
$\mathcal{O}_{X'}$-模组成的有界复形；它在
$X' \setminus Z'$ 上的限制同构于置于次数 $0$ 的、秩
$< p_1 + p_3$ 的有限局部自由
$\mathcal{O}_{X' \setminus Z'}$-模。因此，对
$p \geq p_1 + p_3$，局部化陈类
$$
c_p(Z' \to X', \mathcal{C}^\bullet) \in A^p(Z' \to X')
$$
已有定义。对任意 $\alpha \in \CH_k(X)$，考虑
$$
c_p(Z' \to X', \mathcal{C}^\bullet) \cap g^*\alpha
\in \CH_{k + 1 - p}(\mathbf{A}^1 \times X)
$$
若限制到 $t = 0$，则映射 $t g^*\varphi^\bullet$ 限制为零，
而 $\mathcal{C}^\bullet|_{t = 0}$ 是
$\mathcal{E}_1^\bullet$ 与 $\mathcal{E}_3^\bullet$ 的直和。
由局部化陈类与基变换的相容性
（引理 \ref{chow-lemma-base-change-loc-chern}），得到
$$
i_0^* \circ c^{(p_1 + p_3)}(Z' \to X', \mathcal{C}^\bullet) \circ g^* =
c^{(p_1 + p_2)}(Z \to X, E_1 \oplus E_3)
$$
作为 $A^{(p_1 + p_3)}(Z \to X)$ 中的等式。另一方面，若限制到
$t = 1$，则映射 $t g^*\varphi^\bullet$ 限制为 $\varphi$，
而 $\mathcal{C}^\bullet|_{t = 1}$ 是表示 $E_2$ 的有限局部自由模
组成的有界复形。因此
$$
i_1^* \circ c^{(p_1 + p_3)}(Z' \to X', \mathcal{C}^\bullet) \circ g^* =
c^{(p_1 + p_2)}(Z \to X, E_2)
$$
在 $A^{(p_1 + p_3)}(Z \to X)$ 中成立。按有理等价的定义，
$i_0^* = i_1^*$（更确切地，这由
引理 \ref{chow-lemma-linebundle-formulae} 中的公式得出），所以
$$
c^{(p_1 + p_2)}(Z \to X, E_2) = c^{(p_1 + p_2)}(Z \to X, E_1 \oplus E_3)
$$
这把问题化归为下一段所讨论的情形。

\medskip\noindent
假设 $E_2 = E_1 \oplus E_3$，并且三元组 $(X, Z, E_i)$
如情形 \ref{chow-situation-loc-chern}。对 $i = 1, 3$，令
$$
b_i : W_i \to \mathbf{P}^1_X
\quad\text{且}\quad
Q_i
\quad\text{且}\quad
T'_i \subset T_i \subset W_{i, \infty}
$$
如引理 \ref{chow-lemma-independent-loc-chern} 的证明中所取。
按定义
$$
c_p(Z \to X, E_i) = c'_p(Q_i)
$$
其中右边是在引理 \ref{chow-lemma-localized-chern-pre} 中使用
$W_i, b_i, Q_i, T'_i$ 构造的双变类。令
$W = W_1 \times_{b_1, \mathbf{P}^1_X, b_2} W_2$，并考虑笛卡儿图
$$
\xymatrix{
W \ar[d]_{g_1} \ar[rd]^b \ar[r]_{g_3} & W_3 \ar[d]^{b_3} \\
W_1 \ar[r]^{b_1} & \mathbf{P}^1_X
}
$$
当然，$b^{-1}(\mathbf{A}^1)$ 同构地映到 $\mathbf{A}^1_X$。
注意，$T' = g_1^{-1}(T'_1) \cap g_2^{-1}(T'_2)$ 仍包含
$W_\infty$ 中所有位于 $X \setminus Z$ 上方的点。
由引理 \ref{chow-lemma-localized-chern-pre-independent}，可以使用
$W$、$b$、$g_i^*\mathcal{Q}_i$ 与
$T'$ 构造 $i = 1, 3$ 时的 $c_p(Z \to X, E_i)$。
此外，由引理 \ref{chow-lemma-independent-loc-chern-bQ}
给出的更强独立性，可以使用
$W$、$b$、$g_1^*Q_1 \oplus g_3^*Q_3$ 与 $T'$
计算各类 $c_p(Z \to X, E_2)$。于是所需等式由
引理 \ref{chow-lemma-localized-chern-pre-sum-c} 得出。
\end{proof}

\begin{lemma}
\label{chow-lemma-additivity-loc-chern-P}
设 $(S, \delta)$ 如情形 \ref{chow-situation-setup}。
设 $X$ 在 $S$ 上局部有限型，且 $Z \to X$ 是闭浸入。设
$$
E_1 \to E_2 \to E_3 \to E_1[1]
$$
是 $D(\mathcal{O}_X)$ 中完美对象组成的可区别三角。假设
\begin{enumerate}
\item 限制 $E_1|_{X \setminus Z}$ 和 $E_3|_{X \setminus Z}$
为零；并且
\item 下列条件至少有一个成立：
(a) $X$ 拟紧；
(b) $X$ 的不可约分支拟紧；
(c) $E_3 \to E_1[1]$ 可由有限局部自由
$\mathcal{O}_X$-模组成的局部有界复形之间的映射表示；或者
(d) 存在包络 $f : Y \to X$，使得
$Lf^*E_3 \to Lf^*E_1[1]$ 可由有限局部自由
$\mathcal{O}_Y$-模组成的局部有界复形之间的映射表示。
\end{enumerate}
则对所有 $p \in \mathbf{Z}$ 都有
$$
P_p(Z \to X, E_2) = P_p(Z \to X, E_1) + P_p(Z \to X, E_3)
$$
从而
$ch(Z \to X, E_2) = ch(Z \to X, E_1) + ch(Z \to X, E_3)$。
\end{lemma}

\begin{proof}
证明与引理 \ref{chow-lemma-additivity-loc-chern-c} 的证明完全相同，
只是最后使用引理 \ref{chow-lemma-localized-chern-pre-sum-P}。
当 $p > 0$ 时，可令 $p_1 = p_3 = 1$，由
引理 \ref{chow-lemma-additivity-loc-chern-c} 以及
引理 \ref{chow-lemma-loc-chern-character} 中
$P_p(Z \to X, E)$ 与 $c_p(Z \to X, E)$ 的关系推出本引理。
$p = 0$ 的情形可直接证明（只有当 $X$ 有某个连通分支完全包含在
$Z$ 中时才有实际内容）。
\end{proof}

\begin{lemma}
\label{chow-lemma-loc-chern-tensor-product}
在情形 \ref{chow-situation-setup} 中，设 $X$ 在 $S$ 上局部有限型。
设 $Z_i \subset X$（$i = 1, 2$）是闭子概形，且
$F_i$（$i = 1, 2$）是 $D(\mathcal{O}_X)$ 中的完美对象。
假设对 $i = 1, 2$，限制 $F_i|_{X \setminus Z_i}$
为零\footnote{想必本引理还有一个变体，其中只假设
$F_i|_{X \setminus Z_i}$ 同构于秩 $< p_i$ 的有限局部自由
$\mathcal{O}_{X \setminus Z_i}$-模。}，并且 $X$ 上的 $F_i$
满足情形 \ref{chow-situation-loc-chern} 的假设 (3)。记
$r_i = P_0(Z_i \to X, F_i) \in A^0(Z_i \to X)$。则在
$A^1(Z_1 \cap Z_2 \to X)$ 中有
$$
c_1(Z_1 \cap Z_2 \to X, F_1 \otimes_{\mathcal{O}_X}^\mathbf{L} F_2) =
r_1 c_1(Z_2 \to X, F_2) + r_2 c_1(Z_1 \to X, F_1)
$$
并且在 $A^2(Z_1 \cap Z_2 \to X)$ 中有
\begin{align*}
c_2(Z_1 \cap Z_2 \to X, F_1 \otimes_{\mathcal{O}_X}^\mathbf{L} F_2)
& =
r_1 c_2(Z_2 \to X, F_2) +
r_2 c_2(Z_1 \to X, F_1) + \\
& {r_1 \choose 2} c_1(Z_2 \to X, F_2)^2 + \\
& (r_1r_2 - 1) c_1(Z_2 \to X, F_2)c_1(Z_1 \to X, F_1) + \\
& {r_2 \choose 2} c_1(Z_1 \to X, F_1)^2
\end{align*}
高阶陈类依此类推。类似地，在
$\prod_{p \geq 0} A^p(Z_1 \cap Z_2 \to X) \otimes \mathbf{Q}$
中有
$$
ch(Z_1 \cap Z_2 \to X, F_1 \otimes_{\mathcal{O}_X}^\mathbf{L} F_2) =
ch(Z_1 \to X, F_1) ch(Z_2 \to X, F_2)
$$
更确切地，在 $A^p(Z_1 \cap Z_2 \to X)$ 中有
$$
P_p(Z_1 \cap Z_2 \to X, F_1 \otimes_{\mathcal{O}_X}^\mathbf{L} F_2) =
\sum\nolimits_{p_1 + p_2 = p}
{p \choose p_1} P_{p_1}(Z_1 \to X, F_1) P_{p_2}(Z_2 \to X, F_2)
$$
\end{lemma}

\begin{proof}
选取真态射 $b_i : W_i \to \mathbf{P}^1_X$、对象
$Q_i \in D(\mathcal{O}_{W_i})$ 以及闭子概形
$T_i \subset W_{i, \infty}$，如 $F_i$ 的局部化陈类构造中所取，
或更一般地如引理 \ref{chow-lemma-independent-loc-chern-bQ} 所取。
选取交换图
$$
\xymatrix{
W \ar[d]^{g_1} \ar[rd]^b \ar[r]_{g_2} & W_2 \ar[d]^{b_2} \\
W_1 \ar[r]^{b_1} & \mathbf{P}^1_X
}
$$
其中所有态射都是真态射，并且在 $\mathbf{A}^1_X$ 上为同构。
例如，可取 $W$ 为 $b_1^{-1}(\mathbf{A}^1_X)$ 与
$b_2^{-1}(\mathbf{A}^1_X)$ 之间同构之图的闭包。
由引理 \ref{chow-lemma-independent-loc-chern-bQ}，可以使用
$W$、$b = b_i \circ g_i$、$Lg_i^*Q_i$ 和
$g_i^{-1}(T_i)$ 构造局部化陈类
$c_p(Z_i \to X, F_i)$。于是化归为下一段所述情形。

\medskip\noindent
假设已有
\begin{enumerate}
\item 真态射 $b : W \to \mathbf{P}^1_X$，它在
$\mathbf{A}^1_X$ 上为同构；
\item $E_i \subset W_\infty$ 是 $Z_i$ 的逆像；
\item 陈类已有定义的完美对象 $Q_i \in D(\mathcal{O}_W)$，且满足
\begin{enumerate}
\item $Q_i$ 在 $b^{-1}(\mathbf{A}^1_X)$ 上的限制是
$F_i$ 的拉回；并且
\item 存在闭子概形 $T_i \subset W_\infty$，它包含
$W_\infty$ 中所有位于 $X \setminus Z_i$ 上方的点，并使
$Q_i|_{T_i}$ 为零。
\end{enumerate}
\end{enumerate}
由引理 \ref{chow-lemma-independent-loc-chern-bQ}，对 $i = 1, 2$ 有
$$
c_p(Z_i \to X, F_i) = c'_p(Q_i) =
(E_i \to Z_i)_* \circ c'_p(Q_i|_{E_i}) \circ C
$$
以及
$$
P_p(Z_i \to X, F_i) = P'_p(Q_i) =
(E_i \to Z_i)_* \circ P'_p(Q_i|_{E_i}) \circ C
$$
其次，注意 $Q = Q_1 \otimes_{\mathcal{O}_W}^\mathbf{L} Q_2$
对 $F_1 \otimes_{\mathcal{O}_X}^\mathbf{L} F_2$ 和
$T_1 \cup T_2$ 满足 (3)(a) 与 (3)(b)。因此由同一个引理可知
$$
c_p(Z_1 \cap Z_2 \to X, F_1 \otimes_{\mathcal{O}_X}^\mathbf{L} F_2) =
(E_1 \cap E_2 \to Z_1 \cap Z_2)_* \circ
c'_p(Q|_{E_1 \cap E_2}) \circ C
$$
以及
$$
P_p(Z_1 \cap Z_2 \to X, F_1 \otimes_{\mathcal{O}_X}^\mathbf{L} F_2) =
(E_1 \cap E_2 \to Z_1 \cap Z_2)_* \circ
P'_p(Q|_{E_1 \cap E_2}) \circ C
$$
由引理 \ref{chow-lemma-silly-tensor-product}，各类
$c'_p(Q|_{E_1 \cap E_2})$ 与 $P'_p(Q|_{E_1 \cap E_2})$
可用各类 $c'_p(Q_i|_{E_i})$ 与 $P'_p(Q_i|_{E_i})$
按正确方式展开。最后，引理 \ref{chow-lemma-homomorphism-final}
说明 $c'_p(Q_i|_{E_i})$ 与 $P'_p(Q_i|_{E_i})$ 的多项式，
同所需的 $c'_p(Q_i)$ 与 $P'_p(Q_i)$ 的对应多项式一致。
\end{proof}






\section{无穷远处的爆破}
\label{chow-section-blowup-Z-first}

\noindent
设 $X$ 是概形，$Z \subset X$ 是由有限型拟凝聚理想层截出的闭子概形。
记 $X' \to X$ 为以 $Z$ 为中心的爆破。设
$b : W \to \mathbf{P}^1_X$ 为以 $\infty(Z)$ 为中心的爆破，
并记 $E \subset W$ 为例外除子。存在交换图
$$
\xymatrix{
X' \ar[r] \ar[d] & W \ar[d]^b \\
X \ar[r]^\infty & \mathbf{P}^1_X
}
$$
其中水平箭头均为闭浸入（《除子》引理 \ref{divisors-lemma-strict-transform}）。记 $E \subset W$
为例外除子，$W_\infty \subset W$ 为
$(\mathbf{P}^1_X)_\infty$ 的逆像。则有
\begin{enumerate}
\item $b$ 在 $\mathbf{A}^1_X \cup \mathbf{P}^1_{X \setminus Z}$
上为同构；
\item $X'$ 是 $W$ 上的有效 Cartier 除子；
\item $X' \cap E$ 是 $X' \to X$ 的例外除子；
\item 作为 $W$ 上的有效 Cartier 除子，
$W_\infty = X' + E$；
\item $E = \underline{\text{Proj}}_Z(\mathcal{C}_{Z/X, *}[S])$，
其中 $S$ 是置于次数 $1$ 的变量；
\item $X' \cap E = \underline{\text{Proj}}_Z(\mathcal{C}_{Z/X, *})$；
\item
\label{chow-item-cone-is-open}
$E \setminus X' = E \setminus (X' \cap E) =
\underline{\Spec}_Z(\mathcal{C}_{Z/X, *}) = C_ZX$；
\item
\label{chow-item-find-Z-in-blowup}
存在闭浸入 $\mathbf{P}^1_Z \to W$，它与 $b$ 的复合是包含态射
$\mathbf{P}^1_Z \to \mathbf{P}^1_X$，而沿 $\infty$
的基变换是复合 $Z \to C_ZX \to E \to W_\infty$；
其中第一个箭头是锥顶。
\end{enumerate}
回忆 $\mathcal{C}_{Z/X, *}$ 是 $Z$ 在 $X$ 中的余法代数，见
《除子》定义 \ref{divisors-definition-conormal-sheaf}；
而 $C_ZX$ 是 $Z$ 在 $X$ 中的法锥，见
《除子》定义 \ref{divisors-definition-normal-cone}。

\medskip\noindent
现在证明上述各项。我们强烈建议读者亲自推演若干例子，而不是阅读
这些证明。

\medskip\noindent
第 (1) 项由《除子》引理 \ref{divisors-lemma-blowing-up-gives-effective-Cartier-divisor}
中的相应陈述得出。由同一引理，$E \subset W$ 是有效 Cartier 除子。

\medskip\noindent
由《除子》引理 \ref{divisors-lemma-blow-up-pullback-effective-Cartier}，
$W_\infty$ 是有效 Cartier 除子。由于 $E \subset W_\infty$，
可对某个有效 Cartier 除子 $D$ 写成 $W_\infty = D + E$；见
《除子》引理 \ref{divisors-lemma-difference-effective-Cartier-divisors}。
下文将看到 $D = X'$，这便证明 (2) 和 (4)。

\medskip\noindent
由于 $X'$ 是闭浸入 $\infty : X \to \mathbf{P}^1_X$
的严格变换（见上文），$X' \to X$ 的例外除子等于交
$X' \cap E$（例如因为二者都由 $Z$ 的理想层在 $X'$ 上的拉回截出）。
这证明了 (3)。

\medskip\noindent
$\infty(Z)$ 与 $\mathbf{P}^1_Z$ 的交是有效 Cartier 除子
$(\mathbf{P}^1_Z)_\infty$，所以 $\mathbf{P}^1_Z$ 经爆破 $b$
的严格变换同构地映到 $\mathbf{P}^1_Z$；见《除子》引理 \ref{divisors-lemma-strict-transform} 和
\ref{divisors-lemma-blow-up-effective-Cartier-divisor}。
这给出 (8) 中提到的态射 $\mathbf{P}^1_Z \to W$。
由于 $b$ 分离，该态射是闭浸入；见《概形》引理 \ref{schemes-lemma-section-immersion}。

\medskip\noindent
假设 $\Spec(A) \subset X$ 是仿射开集，并且
$Z \cap \Spec(A)$ 对应于有限生成理想 $I \subset A$。
$\infty(Z \cap \Spec(A))$ 的一个仿射邻域是以
$s = T_0/T_1$ 为坐标的 $A$ 上仿射空间。记
$J = (I, s) \subset A[s]$ 为由 $I$ 和 $s$ 生成的理想。
令 $B = A[s] \oplus J \oplus J^2 \oplus \ldots$ 为
$(A[s], J)$ 的 Rees 代数。注意，对所有 $n \geq 0$，
作为 $A[s]$ 的 $A$-子模，有
$$
J^n =
I^n \oplus sI^{n - 1} \oplus s^2I^{n - 2} \ldots \oplus s^nA
\oplus s^{n + 1}A \oplus \ldots
$$
考虑开子概形
$$
\text{Proj}(B) = \text{Proj}(A[s] \oplus J \oplus J^2 \oplus \ldots)
\subset W
$$
最后，将元素 $s \in J$ 视为 $B$ 的次数 $1$ 元素，并记作 $S$。

\medskip\noindent
由于 $\text{Proj}$ 的形成与基变换交换（《概形的构造》引理 \ref{constructions-lemma-base-change-map-proj}），可见
$$
E = \text{Proj}(B \otimes_{A[s]} A/I) =
\text{Proj}((A/I \oplus I/I^2 \oplus I^2/I^3 \oplus \ldots)[S])
$$
由上面对 $J^n$ 各次幂的描述，立即验证
$B \otimes_{A[s]} A/I = \bigoplus J^n/J^{n + 1}$
确实如上所示。这证明了 (5)，因为由《除子》引理 \ref{divisors-lemma-affine-conormal-sheaf}，
$Z \cap \Spec(A)$ 在 $\Spec(A)$ 中的余法代数对应于分次
$A$-代数 $A/I \oplus I/I^2 \oplus I^2/I^3 \oplus \ldots$。

\medskip\noindent
回忆 $\text{Proj}(B)$ 由仿射开集 $D_+(S)$ 以及
$f \in I$ 时的 $D_+(f^{(1)})$ 覆盖；它们分别是仿射爆破代数
$A[s][\frac{J}{s}]$ 和 $A[s][\frac{J}{f}]$ 的谱，见
《除子》引理 \ref{divisors-lemma-blowing-up-affine} 和
《交换代数》定义 \ref{algebra-definition-blow-up}。
逐一描述这些仿射开集即可完成证明。

\medskip\noindent
先考虑开集 $D_+(S)$，即 $A[s][\frac{J}{s}]$ 的谱。
由上面对 $J$ 的各次幂的描述，
$$
A[s][\textstyle{\frac{J}{s}}] = \sum s^{-n}I^n[s] \subset A[s, s^{-1}]
$$
该环中的元素 $s$ 是非零因子，并且既定义例外除子 $E$，
也定义 $W_\infty$。因此 $D \cap D_+(S) = \emptyset$。
最后，$A[s][\frac{J}{s}]$ 对 $s$ 的商是余法代数
$$
A/I \oplus I/I^2 \oplus I^2/I^3 \oplus \ldots
$$
这证明了 (7)。

\medskip\noindent
再考虑开集 $D_+(f^{(1)})$，即 $A[s][\frac{J}{f}]$ 的谱。
由上面对 $J$ 的各次幂的描述，
$$
A[s][\textstyle{\frac{J}{f}}] =
A[\textstyle{\frac{I}{f}}][\textstyle{\frac{s}{f}}]
$$
其中 $\frac{s}{f}$ 是变量。元素 $f$ 在该环中是非零因子，
其零点概形定义例外除子 $E$。由于 $s$ 定义 $W_\infty$，且
$s = f \cdot \frac{s}{f}$，可知 $\frac{s}{f}$ 定义上面构造的除子
$D$。于是
$$
D \cap D_+(f^{(1)}) = \Spec(A[\textstyle{\frac{I}{f}}])
$$
这是爆破 $X'$ 在 $\Spec(A)$ 上的相应开集。具体而言，到
$(A, I)$ 的 Rees 代数的满分次 $A[s]$-代数映射
$B \to A \oplus I \oplus I^2 \oplus \ldots$
对应于 $\Spec(A[s])$ 上的闭浸入 $X' \to W$。
这便如需地证明了 $D = X'$。

\medskip\noindent
下面证明 (6)。注意，上一段中 $\frac{s}{f}$ 的零点概形，
是 $S$ 的零点概形在仿射开集 $D_+(f^{(1)})$ 上的限制。
因此 $S = 0$ 在 $E$ 上定义 $X' \cap E$。由 (5) 即得 (6)。

\medskip\noindent
最后还须证明 (8) 的最后一部分。这是清楚的，因为态射
$\mathbf{P}^1_Z \to W$ 在仿射局部由满射
$$
B \to B \otimes_{A[s]} A/I =
(A/I \oplus I/I^2 \oplus I^2/I^3 \oplus \ldots)[S] \to
A/I[S]
$$
以及同一化 $\text{Proj}(A/I[S]) = \Spec(A/I)$ 给出。
省略一些细节。






\section{高余维 Gysin 同态}
\label{chow-section-gysin-higher-codimension}

\noindent
设 $(S, \delta)$ 如情形 \ref{chow-situation-setup}，且
$X$ 是 $S$ 上局部有限型概形。本节考虑三元组
$$
(Z \to X, \mathcal{N}, \sigma : \mathcal{N}^\vee \to \mathcal{C}_{Z/X})
$$
其中 $Z \to X$ 是闭浸入，$\mathcal{N}$ 是局部自由
$\mathcal{O}_Z$-模，而
$\sigma : \mathcal{N}^\vee \to \mathcal{C}_{Z/X}$
是从 $\mathcal{N}$ 的对偶到 $Z$ 在 $X$ 中的余法层的满射；见
《概形态射》第 \ref{morphisms-section-conormal-sheaf} 节。
我们称 $\mathcal{N}$ 是{\it $Z$ 在 $X$ 中的虚拟法层}。

\begin{lemma}
\label{chow-lemma-pullback-virtual-normal-sheaf}
设 $(S, \delta)$ 如情形 \ref{chow-situation-setup}。设
$$
\xymatrix{
Z' \ar[r] \ar[d]_g & X' \ar[d]^f \\
Z \ar[r] & X
}
$$
是 $S$ 上局部有限型概形的笛卡儿图，其中水平箭头均为闭浸入。
若 $\mathcal{N}$ 是 $Z$ 在 $X$ 中的虚拟法层，则
$\mathcal{N}' = g^*\mathcal{N}$ 是 $Z'$ 在 $X'$ 中的虚拟法层。
\end{lemma}

\begin{proof}
这由《概形态射》引理 \ref{morphisms-lemma-conormal-functorial-flat} 所证明的映射
$g^*\mathcal{C}_{Z/X} \to \mathcal{C}_{Z'/X'}$ 的满性得出。
\end{proof}

\noindent
设 $(S, \delta)$ 如情形 \ref{chow-situation-setup}，且
$X$ 是 $S$ 上局部有限型概形。设 $\mathcal{N}$ 是闭浸入
$Z \to X$ 的虚拟法丛。在此情形下，令
$$
p : N = \underline{\Spec}_Z(\text{Sym}(\mathcal{N}^\vee)) \longrightarrow Z
$$
为 $Z$ 上的向量丛，其截面对应于 $\mathcal{N}$ 的截面。
此时存在典范闭浸入
$$
C_ZX \longrightarrow N_ZX \longrightarrow N
$$
第一个闭浸入是 Divisors, Equation
(\ref{divisors-equation-normal-cone-in-normal-bundle})，
第二个闭浸入对应于 $\sigma$ 诱导的满射
$\text{Sym}(\mathcal{N}^\vee) \to \text{Sym}(\mathcal{C}_{Z/X})$。
令
$$
b : W \longrightarrow \mathbf{P}^1_X
$$
为第 \ref{chow-section-blowup-Z-first} 节中构造的沿
$\infty(Z)$ 的爆破。由引理 \ref{chow-lemma-gysin-at-infty}，
存在典范双变类
$$
C \in A^0(W_\infty \to X)
$$
考虑 (\ref{chow-item-cone-is-open}) 中的开浸入
$j : C_ZX \to W_\infty$，以及上面构造的闭浸入
$i : C_ZX \to N$。由引理 \ref{chow-lemma-vectorbundle}，
对每个 $\alpha \in \CH_k(X)$，存在唯一的
$\beta \in \CH_*(Z)$，使得
$$
i_*j^*(C \cap \alpha) = p^*\beta
$$
定义 $c(Z \to X, \mathcal{N}) \cap \alpha = \beta$。

\begin{lemma}
\label{chow-lemma-construction-gysin}
上述构造定义了双变类\footnote{记号
$A^*(Z \to X)^\wedge$ 见注 \ref{chow-remark-completion-bivariant} 的讨论。若 $X$ 拟紧，则
$A^*(Z \to X)^\wedge = A^*(Z \to X)$。}
$$
c(Z \to X, \mathcal{N}) \in A^*(Z \to X)^\wedge
$$
此外，该构造与引理 \ref{chow-lemma-pullback-virtual-normal-sheaf} 中的基变换相容。
若 $\mathcal{N}$ 的秩恒为 $r$，则
$c(Z \to X, \mathcal{N}) \in A^r(Z \to X)$。
\end{lemma}

\begin{proof}
由于 $i_* \circ j^* \circ C$ 和 $p^*$ 都是双变类
（见引理 \ref{chow-lemma-flat-pullback-bivariant} 和
\ref{chow-lemma-push-proper-bivariant}），可以用等式
$$
i_* \circ j^* \circ C  = p^* \circ c(Z \to X, \mathcal{N})
$$
（作适当解释）将 $c(Z \to X, \mathcal{N})$ 定义为双变类。
这是可行的，因为由引理 \ref{chow-lemma-vectorbundle}，
$p^*$ 在 Chow 群上总是双射。

\medskip\noindent
设 $X' \to X$、$Z' \to X'$ 与 $\mathcal{N}'$ 如
引理 \ref{chow-lemma-pullback-virtual-normal-sheaf}。记
$c = c(Z \to X, \mathcal{N})$ 且
$c' = c(Z' \to X', \mathcal{N}')$。引理的第二个陈述是说，
$c'$ 是注 \ref{chow-remark-restriction-bivariant} 意义下
$c$ 的限制。由于该断言声称对所有局部有限型的 $X'/X$ 都成立，
一个形式论证表明，只须对
$\alpha' \in \CH_k(X')$ 检验
$c' \cap \alpha' = c \cap \alpha'$。为此，注意存在交换图
$$
\xymatrix{
C_{Z'}X' \ar[d] \ar[r] &
W'_\infty \ar[d] \ar[r] &
W' \ar[d] \ar[r] &
\mathbf{P}^1_{X'} \ar[d] \\
C_ZX \ar[r] &
W_\infty \ar[r] &
W \ar[r] &
\mathbf{P}^1_X
}
$$
它诱导闭浸入
$$
W' \to W \times_{\mathbf{P}^1_X} \mathbf{P}^1_{X'},\quad
W'_\infty \to W_\infty \times_X X',\quad
C_{Z'}X' \to C_ZX \times_Z Z'
$$
为得到 $c \cap \alpha'$，使用在引理 \ref{chow-lemma-gysin-at-infty} 中由态射
$W \times_{\mathbf{P}^1_X} \mathbf{P}^1_{X'} \to \mathbf{P}^1_{X'}$
定义的类 $C \cap \alpha'$。另一方面，为得到
$c' \cap \alpha'$，使用由态射
$W' \to \mathbf{P}^1_{X'}$ 定义的类 $C' \cap \alpha'$。
由引理 \ref{chow-lemma-gysin-at-infty-independent}，
$C' \cap \alpha'$ 沿闭浸入
$W'_\infty \to (W \times_{\mathbf{P}^1_X} \mathbf{P}^1_{X'})_\infty$
的推前等于 $C \cap \alpha'$。因此，由引理 \ref{chow-lemma-flat-pullback-proper-pushforward}，二者到开集
$$
C_{Z'}X' \subset W'_\infty,\quad
C_ZX \times_Z Z' \subset (W \times_{\mathbf{P}^1_X} \mathbf{P}^1_{X'})_\infty
$$
的拉回也相同。又有交换图
$$
\xymatrix{
C_{Z'} X' \ar[d] \ar[r] & N' \ar@{=}[d] \\
C_ZX \times_Z Z' \ar[r] & N \times_Z Z'
}
$$
这些类推前到 $N'$ 上得到同一个类，故在 $\CH_*(Z')$ 中得到
同一元素 $c \cap \alpha' = c' \cap \alpha'$。
\end{proof}

\begin{lemma}
\label{chow-lemma-gysin-decompose}
设 $(S, \delta)$ 如情形 \ref{chow-situation-setup}。
设 $X$ 是 $S$ 上局部有限型概形，且 $\mathcal{N}$
是 $X$ 的闭子概形 $Z$ 的虚拟法层。假设存在有限局部自由
$\mathcal{O}_Z$-模的短正合列
$0 \to \mathcal{N}' \to \mathcal{N} \to \mathcal{E} \to 0$，
使得给定满射
$\sigma : \mathcal{N}^\vee \to \mathcal{C}_{Z/X}$
分解经过映射
$\sigma' : (\mathcal{N}')^\vee \to \mathcal{C}_{Z/X}$。
则作为双变类，有
$$
c(Z \to X, \mathcal{N}) = c_{top}(\mathcal{E}) \circ c(Z \to X, \mathcal{N}')
$$
\end{lemma}

\begin{proof}
记 $N' \to N$ 为对应于满射
$\mathcal{N}^\vee \to (\mathcal{N}')^\vee$ 的向量丛闭浸入。
于是存在闭浸入
$$
C_ZX \to N' \to N
$$
所需双变类关系立即由引理 \ref{chow-lemma-easy-virtual-class} 得出。
\end{proof}

\begin{lemma}
\label{chow-lemma-gysin-excess}
设 $(S, \delta)$ 如情形 \ref{chow-situation-setup}。考虑
$S$ 上局部有限型概形的笛卡儿图
$$
\xymatrix{
Z' \ar[r] \ar[d]_g & X' \ar[d]^f \\
Z \ar[r] & X
}
$$
其中水平箭头均为闭浸入。设 $\mathcal{N}$，分别地\ 
$\mathcal{N}'$，是 $Z \subset X$，分别地\ $Z' \to X'$，
的虚拟法层。假设给定 $Z'$ 上有限局部自由模的短正合列
$0 \to \mathcal{N}' \to g^*\mathcal{N} \to \mathcal{E} \to 0$，
使得图
$$
\xymatrix{
g^*\mathcal{N}^\vee \ar[r] \ar[d] &
(\mathcal{N}')^\vee \ar[d] \\
g^*\mathcal{C}_{Z/X} \ar[r] &
\mathcal{C}_{Z'/X'}
}
$$
交换。则在 $A^*(Z' \to X')^\wedge$ 中有
$$
res(c(Z \to X, \mathcal{N})) =
c_{top}(\mathcal{E}) \circ c(Z' \to X', \mathcal{N}')
$$
\end{lemma}

\begin{proof}
由引理 \ref{chow-lemma-construction-gysin}，
$res(c(Z \to X, \mathcal{N})) = c(Z' \to X', g^*\mathcal{N})$；
于是等式由引理 \ref{chow-lemma-gysin-decompose} 得出。
\end{proof}

\noindent
设 $(S, \delta)$ 如情形 \ref{chow-situation-setup}，且
$X$ 是 $S$ 上局部有限型概形。设 $\mathcal{N}$ 是
$X$ 的闭子概形 $Z$ 的虚拟法层。设 $Y \to X$ 是局部有限型态射。
假设 $Z \times_X Y \to Y$ 是正则闭浸入；见《除子》第 \ref{divisors-section-regular-immersions} 节。此时余法层
$\mathcal{C}_{Z \times_X Y/Y}$ 是有限局部自由
$\mathcal{O}_{Z \times_X Y}$-模，并得到短正合列
$$
0 \to \mathcal{E}^\vee \to
\mathcal{N}^\vee|_{Z \times_X Y} \to \mathcal{C}_{Z \times_X Y/Y} \to 0
$$
商 $\mathcal{N}|_{Y \times_X Z} \to \mathcal{E}$
称为此情形的{\it 超额法层}。

\begin{lemma}
\label{chow-lemma-gysin-fundamental}
在上述情形中，假设 $\dim_\delta(Y) = n$，且
$\mathcal{C}_{Y \times_X Z/Z}$ 的秩恒为 $r$。则在
$\CH_*(Z \times_X Y)$ 中有
$$
c(Z \to X, \mathcal{N}) \cap [Y]_n =
c_{top}(\mathcal{E}) \cap [Z \times_X Y]_{n - r}
$$
\end{lemma}

\begin{proof}
双变类 $c_{top}(\mathcal{E}) \in A^*(Z \times_X Y)$
在注 \ref{chow-remark-top-chern-class} 中定义。
由引理 \ref{chow-lemma-construction-gysin}，可将 $X$ 替换为 $Y$。
因此可假设 $Z \to X$ 是余维 $r$ 的正则闭浸入，
$\dim_\delta(X) = n$，并须在 $\CH_*(Z)$ 中证明
$c(Z \to X, \mathcal{N}) \cap [X]_n =
c_{top}(\mathcal{E}) \cap [Z]_{n - r}$。
由引理 \ref{chow-lemma-gysin-decompose}，甚至可假设
$\mathcal{N}^\vee \to \mathcal{C}_{Z/X}$ 是同构。
换言之，须在 $\CH_*(Z)$ 中证明
$c(Z \to X, \mathcal{C}_{Z/X}^\vee) \cap [X]_n = [Z]_{n - r}$。

\medskip\noindent
下面逐步追踪
$c(Z \to X, \mathcal{C}_{Z/X}^\vee) \cap [X]_n$ 的定义。
设 $b : W \to \mathbf{P}^1_X$ 是沿 $\infty(Z)$ 的爆破。
首先计算 $C \cap [X]_n$，其中
$C \in A^0(W_\infty \to X)$ 是
引理 \ref{chow-lemma-gysin-at-infty} 中的类。注意
$[W]_{n + 1}$ 是 $W$ 上的周，它在 $\mathbf{A}^1_X$
上的限制等于 $[X]_n$ 的平坦拉回。因此
$C \cap [X]_n$ 等于 $i_\infty^*[W]_{n + 1}$。
由于 $W_\infty$ 是 $W$ 上的有效 Cartier 除子，
$i_\infty^*[W]_{n + 1} = [W_\infty]_n$；见
引理 \ref{chow-lemma-easy-gysin}。该类在开集
$C_ZX \subset W_\infty$ 上的限制当然正是 $[C_ZX]_n$。
因为 $Z \subset X$ 是正则嵌入，有
$$
\mathcal{C}_{Z/X, *} = \text{Sym}(\mathcal{C}_{Z/X})
$$
作为分次 $\mathcal{O}_Z$-代数；见《除子》引理 \ref{divisors-lemma-quasi-regular-immersion}。因此
$p : N = C_ZX \to Z$ 是与秩 $r$ 的有限局部自由模
$\mathcal{C}_{Z/X}$ 对应之向量丛的结构态射。显然
$p^*[Z]_{n - r} = [C_ZX]_n$，证明完毕。
\end{proof}

\begin{lemma}
\label{chow-lemma-gysin-easy}
设 $(S, \delta)$ 如情形 \ref{chow-situation-setup}。
设 $X$ 是 $S$ 上局部有限型概形，且 $\mathcal{N}$
是 $X$ 的闭子概形 $Z$ 的虚拟法层。设
$Y \to X$ 是局部有限型态射。给定整数 $r$、$n$，假设
\begin{enumerate}
\item $\mathcal{N}$ 局部自由且秩为 $r$；
\item $Y$ 的每个不可约分支的 $\delta$-维数均为 $n$；
\item $\dim_\delta(Z \times_X Y) \leq n - r$；并且
\item 对满足 $\delta(\xi) = n - r$ 的
$\xi \in Z \times_X Y$，局部环
$\mathcal{O}_{Y, \xi}$ 是 Cohen--Macaulay 环。
\end{enumerate}
则在 $\CH_{n - r}(Z \times_X Y)$ 中有
$c(Z \to X, \mathcal{N}) \cap [Y]_n = [Z \times_X Y]_{n - r}$。
\end{lemma}

\begin{proof}
由于 $Z \times_X Y$ 是 $Y$ 的闭子概形，该陈述是有意义的。
因为 $\mathcal{N}$ 的秩为 $r$，可知
$c(Z \to X, \mathcal{N}) \cap [Y]_n$ 属于
$\CH_{n - r}(Z \times_X Y)$。由于
$\dim_\delta(Z \cap Y) \leq n - r$，Chow 群
$\CH_{n - r}(Z \times_X Y)$ 由 $\delta$-维数为
$n - r$ 的各不可约分支 $W \subset Z \times_X Y$
的周类自由生成。设 $\xi \in W$ 是泛点。
由假设 (2)，$\dim(\mathcal{O}_{Y, \xi}) = r$。
另一方面，由于 $\mathcal{N}$ 的秩为 $r$，且
$\mathcal{N}^\vee \to \mathcal{C}_{Z/X}$ 是满射，
$Z$ 的理想层在局部由 $r$ 个方程截出。因此
$Z \times_X Y$ 在 $Y$ 中的拟凝聚理想层
$\mathcal{I} \subset \mathcal{O}_Y$ 在局部由 $r$ 个元素生成。
由于 $\mathcal{O}_{Y, \xi}$ 是维数 $r$ 的 Cohen--Macaulay 环，
且 $\mathcal{I}_\xi$ 是定义理想（因为 $\xi$ 是
$Z \times_X Y$ 的泛点），$\mathcal{I}_\xi$ 由正则序列生成
（《交换代数》引理 \ref{algebra-lemma-reformulate-CM}）。
由《除子》引理 \ref{divisors-lemma-Noetherian-scheme-regular-ideal}，
在 $\xi$ 的某个开邻域 $V \subset Y$ 上，$\mathcal{I}$
由正则序列生成。根据上面对
$\CH_{n - r}(Z \times_X Y)$ 的描述，只须在
$\CH_{n - r}(Z \times_X V)$ 中证明
$c(Z \to X, \mathcal{N}) \cap [V]_n = [Z \times_X V]_{n - r}$。
这由引理 \ref{chow-lemma-gysin-fundamental} 得出，因为超额法层在
$V$ 上为 $0$。
\end{proof}

\begin{lemma}
\label{chow-lemma-gysin-agrees}
设 $(S, \delta)$ 如情形 \ref{chow-situation-setup}，且
$X$ 是 $S$ 上局部有限型概形。设
$(\mathcal{L}, s, i : D \to X)$ 是定义 \ref{chow-definition-gysin-homomorphism} 中的三元组。
将 Gysin 同态 $i^*$ 视为 $A^1(D \to X)$ 的元素
（见引理 \ref{chow-lemma-gysin-bivariant}），则它与双变类
$c(D \to X, \mathcal{N}) \in A^1(D \to X)$ 相同；
后者把 $\mathcal{N} = i^*\mathcal{L}$ 视为
$D$ 在 $X$ 中的虚拟法层而构造。
\end{lemma}

\begin{proof}
使用引理 \ref{chow-lemma-bivariant-zero} 的判据。
因此可假设 $X$ 是整概形，并须证明
$i^*[X]$ 等于 $c \cap [X]$。令
$n = \dim_\delta(X)$。照例分两种情形。

\medskip\noindent
若 $X = D$，则两个类都由
$c_1(\mathcal{N}) \cap [X]_n$ 表示。见
引理 \ref{chow-lemma-gysin-fundamental} 和定义 \ref{chow-definition-gysin-homomorphism}。

\medskip\noindent
若 $D \not = X$，则 $D \to X$ 是有效 Cartier 除子，
特别地，是余维 $1$ 的正则闭浸入。再次由
引理 \ref{chow-lemma-gysin-fundamental} 得到
$c(D \to X, \mathcal{N}) \cap [X]_n = [D]_{n - 1}$。
根据定义，Gysin 同态也具有同一性质，故仍得到结论。
\end{proof}

\begin{lemma}
\label{chow-lemma-gysin-commutes}
设 $(S, \delta)$ 如情形 \ref{chow-situation-setup}，且
$X$ 是 $S$ 上局部有限型概形。设 $Z \subset X$ 是具有虚拟法层
$\mathcal{N}$ 的闭子概形。设 $Y \to X$ 局部有限型，且
$c \in A^*(Y \to X)$。则 $c$ 与 $c(Z \to X, \mathcal{N})$
交换（注 \ref{chow-remark-bivariant-commute}）。
\end{lemma}

\begin{proof}
可用引理 \ref{chow-lemma-bivariant-zero} 检验此事。
因此可假设 $X$ 是整概形，并须在 $\CH_*(Z \times_X Y)$ 中证明
$c \cap c(Z \to X, \mathcal{N}) \cap [X] =
c(Z \to X, \mathcal{N}) \cap c \cap [X]$。

\medskip\noindent
若 $Z = X$，则由引理 \ref{chow-lemma-gysin-fundamental}，
$c(Z \to X, \mathcal{N}) = c_{top}(\mathcal{N})$；
它与双变类 $c$ 交换，见引理 \ref{chow-lemma-cap-commutative-chern}。

\medskip\noindent
假设 $Z$ 不等于 $X$。由引理 \ref{chow-lemma-bivariant-zero}，
甚至只须在对 $X$ 沿某个非零理想爆破后证明结论。
沿 $Z$ 的理想层爆破 $X$。这化归为 $Z$ 是有效 Cartier 除子的情形；
见《除子》引理 \ref{divisors-lemma-blowing-up-gives-effective-Cartier-divisor}。

\medskip\noindent
若 $Z$ 是有效 Cartier 除子，则
$$
c(Z \to X, \mathcal{N}) =
c_{top}(\mathcal{E}) \circ i^*
$$
其中 $i^* \in A^1(Z \to X)$ 是与 $i : Z \to X$ 对应的
Gysin 同态（引理 \ref{chow-lemma-gysin-bivariant}），而
$\mathcal{E}$ 是 $\mathcal{N}^\vee \to \mathcal{C}_{Z/X}$
之核的对偶；见引理 \ref{chow-lemma-gysin-decompose} 和
\ref{chow-lemma-gysin-agrees}。陈类位于双变环的中心
（以引理 \ref{chow-lemma-cap-commutative-chern} 所表述的强意义），
而按双变类的定义，$c$ 与 Gysin 同态 $i^*$ 交换，故结论成立。
\end{proof}

\noindent
设 $(S, \delta)$ 如情形 \ref{chow-situation-setup}。
设 $X$ 是 $S$ 上局部有限型、$\delta$-维数为 $n$ 的整概形。
设 $Z \subset Y \subset X$ 是闭子概形，且二者均为
$X$ 中的有效 Cartier 除子。记 $o : Y \to C_Y X$
为 $Y$ 在 $X$ 中的法线锥之零截面。由于
$C_YX$ 是 $Y$ 上的线丛，得到双变类
$o^* \in A^1(Y \to C_YX)$；见引理 \ref{chow-lemma-gysin-bivariant}。

\begin{lemma}
\label{chow-lemma-relation-normal-cones}
采用上述记号，在 $\CH_{n - 1}(Y \times_{o, C_Y X} C_ZX)$ 中有
$$
o^*[C_ZX]_n = [C_Z Y]_{n - 1}
$$
\end{lemma}

\begin{proof}
记 $W \to \mathbf{P}^1_X$ 为第 \ref{chow-section-blowup-Z-first} 节中沿 $\infty(Z)$ 的爆破。
类似地，记 $W' \to \mathbf{P}^1_X$ 为沿 $\infty(Y)$ 的爆破。
由于 $\infty(Z) \subset \infty(Y)$，得到理想层的反向包含，
从而得到定义这些爆破的分次代数之间的映射。这产生从
$W$ 到 $W'$ 的有理态射；事实上，它有典范代表
$$
W \supset U \longrightarrow W'
$$
见《概形的构造》引理 \ref{constructions-lemma-morphism-relative-proj}。
一次局部计算（省略）表明，$U$ 至少包含 $W$ 中所有不位于
$\infty$ 上方的点，以及特殊纤维的开子概形 $C_Z X$。
缩小 $U$ 后，可假设 $U_\infty = C_Z X$ 且
$\mathbf{A}^1_X \subset U$。另一次局部计算（省略）表明，
态射 $U_\infty \to W'_\infty$ 诱导典范态射
$C_Z X \to C_Y X \subset W'_\infty$；这是由
$Z \subset Y$ 所给的理想层包含诱导的法锥态射。
记 $W'' \subset W$ 为 $\mathbf{P}^1_Y \subset \mathbf{P}^1_X$
在 $W$ 中的严格变换。由《除子》引理 \ref{divisors-lemma-strict-transform}，$W''$ 是
$\mathbf{P}^1_Y$ 沿 $\infty(Z)$ 的爆破，因而
$(W'' \cap U)_\infty = C_ZY$。

\medskip\noindent
考虑 (\ref{chow-item-find-Z-in-blowup}) 中的有效 Cartier 除子
$i : \mathbf{P}^1_Y \to W'$，以及由引理 \ref{chow-lemma-gysin-bivariant} 得到的相应双变类
$i^* \in A^1(\mathbf{P}^1_Y \to W')$。
类似地，记 $(i'_\infty)^* \in A^1(W'_\infty \to W')$
为无穷远处的 Gysin 映射。注意，$i'_\infty$ 在 $U$ 上的限制
（注 \ref{chow-remark-restriction-bivariant}）等于
$i_\infty^* \in A^1(W_\infty \to W)$ 在 $U$ 上的限制。
一方面，
$$
(i'_\infty)^* i^* [U]_{n + 1} =
i_\infty^* i^* [U]_{n + 1} =
i_\infty^* [(W'' \cap U)_\infty]_{n + 1} =
[C_ZY]_n
$$
这是因为 $i_\infty^*$ 消灭所有支撑在 $\infty$ 上方的类，
$i^*[U]$ 与 $[W'']$ 作为 $\mathbf{A}^1$ 上的周相同，
并且 $C_ZY$ 是 $W'' \cap U$ 在 $\infty$ 上的纤维。
另一方面，
$$
(i'_\infty)^* i^* [U]_{n + 1} =
i^* i_\infty^*[U]_{n + 1} =
i^* [U_\infty] =
o^*[C_YX]_n
$$
这是因为 $(i'_\infty)^*$ 与 $i^*$ 交换
（引理 \ref{chow-lemma-gysin-commutes-gysin}），并且
$i : \mathbf{P}^1_Y \to W'$ 在 $\infty$ 上的纤维分解为
$o : Y \to C_YX$ 与开浸入 $C_YX \to W'_\infty$。
引理得证。
\end{proof}

\begin{lemma}
\label{chow-lemma-gysin-composition}
设 $(S, \delta)$ 如情形 \ref{chow-situation-setup}。
设 $Z \subset Y \subset X$ 是某个 $S$ 上局部有限型概形的闭子概形。
设 $\mathcal{N}$ 是 $Z \subset X$ 的虚拟法层，
$\mathcal{N}'$ 是 $Z \subset Y$ 的虚拟法层，而
$\mathcal{N}''$ 是 $Y \subset X$ 的虚拟法层。假设存在交换图
$$
\xymatrix{
(\mathcal{N}'')^\vee|_Z \ar[r] \ar[d] &
\mathcal{N}^\vee \ar[r] \ar[d] &
(\mathcal{N}')^\vee \ar[d] \\
\mathcal{C}_{Y/X}|_Z \ar[r] &
\mathcal{C}_{Z/X} \ar[r] &
\mathcal{C}_{Z/Y}
}
$$
其中下方序列来自《态射进阶》引理 \ref{more-morphisms-lemma-transitivity-conormal}，而上方序列是短正合列。
则在 $A^*(Z \to X)^\wedge$ 中有
$$
c(Z \to X, \mathcal{N}) =
c(Z \to Y, \mathcal{N}') \circ c(Y \to X, \mathcal{N}'')
$$
\end{lemma}

\begin{proof}
注意，沿任意局部有限型态射 $X' \to X$ 作基变换后，
这些假设仍然成立（虚拟法层的短正合列在局部分裂，
故沿任意基变换仍然正合）。因此可用
引理 \ref{chow-lemma-bivariant-zero} 检验双变类的等式。
于是可假设 $X$ 是整概形，并须证明
$c(Z \to X, \mathcal{N}) \cap [X] =
c(Z \to Y, \mathcal{N}') \cap c(Y \to X, \mathcal{N}'') \cap [X]$。

\medskip\noindent
若 $Y = X$，则
\begin{align*}
c(Z \to Y, \mathcal{N}') \cap c(Y \to X, \mathcal{N}'') \cap [X]
& =
c(Z \to Y, \mathcal{N}') \cap c_{top}(\mathcal{N}'') \cap [Y] \\
& =
c_{top}(\mathcal{N}''|_Z) \cap c(Z \to Y, \mathcal{N}') \cap [Y] \\
& =
c(Z \to X, \mathcal{N}) \cap [X] 
\end{align*}
第一个等式由引理 \ref{chow-lemma-gysin-decompose} 得出。
第二个等式是因为陈类与双变类交换
（引理 \ref{chow-lemma-cap-commutative-chern}）。
第三个等式由引理 \ref{chow-lemma-gysin-decompose} 得出。

\medskip\noindent
假设 $Y \not = X$。由引理 \ref{chow-lemma-bivariant-zero}，
甚至只须在对 $X$ 沿某个非零理想爆破后证明结论。
沿 $Y$ 的理想层与 $Z$ 的理想层之积爆破 $X$。
这化归为 $Y$ 和 $Z$ 均为 $X$ 上有效 Cartier 除子的情形；见
《除子》引理 \ref{divisors-lemma-blowing-up-gives-effective-Cartier-divisor} 和
\ref{divisors-lemma-blowing-up-two-ideals}。

\medskip\noindent
记 $\mathcal{N}'' \to \mathcal{E}$ 为有限局部自由
$\mathcal{O}_Z$-模的满射，使得
$0 \to \mathcal{E}^\vee \to (\mathcal{N}'')^\vee \to \mathcal{C}_{Y/X} \to 0$
是短正合列。于是 $\mathcal{N} \to \mathcal{E}|_Z$
也是满射。记 $\mathcal{N}_1$ 为该映射的有限局部自由核，并注意
$\mathcal{N}^\vee \to \mathcal{C}_{Z/X}$ 分解经过
$\mathcal{N}_1$。由引理 \ref{chow-lemma-gysin-decompose}，
$$
c(Y \to X, \mathcal{N}'') = c_{top}(\mathcal{E}) \circ
c(Y \to X, \mathcal{C}_{Y/X}^\vee)
$$
且
$$
c(Z \to X, \mathcal{N}) = c_{top}(\mathcal{E}|_Z) \circ
c(Z \to X, \mathcal{N}_1)
$$
由于丛的陈类与双变类交换
（引理 \ref{chow-lemma-cap-commutative-chern}），只须在
$A^*(Z \to X)$ 中证明
$$
c(Z \to X, \mathcal{N}_1) =
c(Z \to Y, \mathcal{N}') \circ c(Y \to X, \mathcal{C}_{Y/X}^\vee)
$$
为此可假设 $\mathcal{N}'' = \mathcal{C}_{Y/X}$。
这化归为下一段所讨论的情形。

\medskip\noindent
在本段中，$Z$ 和 $Y$ 是维数为 $n$ 的整概形 $X$
上的有效 Cartier 除子，且
$\mathcal{N}'' = \mathcal{C}_{Y/X}$。此时由
引理 \ref{chow-lemma-gysin-fundamental}，
$c(Y \to X, \mathcal{C}_{Y/X}^\vee) \cap [X] = [Y]_{n - 1}$。
因此须证明
$c(Z \to X, \mathcal{N}) \cap [X] = c(Z \to Y, \mathcal{N}') \cap [Y]_{n - 1}$。
记 $N$ 和 $N'$ 为与 $\mathcal{N}$ 和 $\mathcal{N}'$
对应的 $Z$ 上向量丛。考虑 $Z$ 上锥和向量丛的交换图
$$
\xymatrix{
N' \ar[r]_i &
N \ar[r] &
(C_Y X) \times_Y Z \\
C_Z Y \ar[r] \ar[u] &
C_Z X \ar[u]
}
$$
注意，$N'$ 是 $N$ 中在 $Z$ 上相对的有效 Cartier 除子，且
$$
\xymatrix{
N' \ar[d] \ar[r]_i & N \ar[d] \\
Z \ar[r]^-o & (C_Y X) \times_Y Z
}
$$
是笛卡儿图，其中 $o$ 是 $Y$ 上线丛
$C_Y X$ 的零截面。由引理 \ref{chow-lemma-relation-normal-cones}，有
$o^*[C_ZX]_n = [C_Z Y]_{n - 1}$，作为
$$
\CH_{n - 1}(Y \times_{o, C_Y X} C_ZX) =
\CH_{n - 1}(Z \times_{o, (C_Y X) \times_Y Z} C_ZX)
$$
中的等式。由上方正方形的笛卡儿性质，这蕴含
$$
i^*[C_ZX]_n = [C_Z Y]_{n - 1}
$$
在 $\CH_{n - 1}(N')$ 中成立。现在注意
$\gamma = c(Z \to X, \mathcal{N}) \cap [X]$ 与
$\gamma' = c(Z \to Y, \mathcal{N}') \cap [Y]_{n - 1}$
分别由 $\CH_n(N)$ 中的 $p^*\gamma = [C_Z X]_n$，
以及 $\CH_{n - 1}(N')$ 中的
$(p')^*\gamma' = [C_Z Y]_{n - 1}$ 刻画。
由引理 \ref{chow-lemma-relative-effective-cartier}，
$i^* \circ p^* = (p')^*$，证明完毕。
\end{proof}

\begin{remark}[浸入情形的变体]
\label{chow-remark-gysin-for-immersion}
设 $(S, \delta)$ 如情形 \ref{chow-situation-setup}。
设 $X$ 是 $S$ 上局部有限型概形，且
$i : Z \to X$ 是概形的浸入。在此情形下，
\begin{enumerate}
\item $Z$ 在 $X$ 中的余法层 $\mathcal{C}_{Z/X}$ 已有定义
（《概形态射》定义 \ref{morphisms-definition-conormal-sheaf}）；
\item 由有限局部自由 $\mathcal{O}_Z$-模 $\mathcal{N}$
和满射 $\sigma : \mathcal{N}^\vee \to \mathcal{C}_{Z/X}$
组成的对，称为浸入 $Z \to X$ 的虚拟法丛；
\item 选取开子概形 $U \subset X$，使得 $Z \to X$
分解经过闭浸入 $Z \to U$，并定义
$c(Z \to X, \mathcal{N}) = c(Z \to U, \mathcal{N}) \circ (U \to X)^*$。
\end{enumerate}
双变类 $c(Z \to X, \mathcal{N})$ 与开子概形 $U$ 的选择无关。
所有引理对这一稍为一般的构造都有直接对应的版本。省略细节。
\end{remark}







\section{计算若干类}
\label{chow-section-calculate}

\noindent
为进一步推进，需要计算上面构造的若干类的取值。

\begin{lemma}
\label{chow-lemma-compute-koszul}
设 $(S, \delta)$ 如情形 \ref{chow-situation-setup}。
设 $X$ 是 $S$ 上局部有限型概形，且 $\mathcal{E}$
是秩 $r$ 的局部自由 $\mathcal{O}_X$-模。则
$$
\prod\nolimits_{n = 0, \ldots, r} c(\wedge^n \mathcal{E})^{(-1)^n} =
1 - (r - 1)! c_r(\mathcal{E}) + \ldots
$$
\end{lemma}

\begin{proof}
由分裂原理，可将其化为在以 $\mathcal{E}$ 的陈根
$x_1, \ldots, x_r$ 为变量的多项式环中计算。见
第 \ref{chow-section-splitting-principle} 节。注意
$$
c(\wedge^n \mathcal{E}) =
\prod\nolimits_{1 \leq i_1 < \ldots < i_n \leq r}
(1 + x_{i_1} + \ldots + x_{i_n})
$$
因此，引理等式左边的对数为
$$
-
\sum\nolimits_{p \geq 1}
\sum\nolimits_{n = 0}^r
\sum\nolimits_{1 \leq i_1 < \ldots < i_n \leq r}
\frac{(-1)^{p + n}}{p}(x_{i_1} + \ldots + x_{i_n})^p
$$
请注意前面的负号。不过，
$$
\sum\nolimits_{p \geq 0}
\sum\nolimits_{n = 0}^r
\sum\nolimits_{1 \leq i_1 < \ldots < i_n \leq r}
\frac{(-1)^{p + n}}{p!}(x_{i_1} + \ldots + x_{i_n})^p
=
\prod (1 - e^{-x_i})
$$
由此可见，所得陈类的第一个非零项在次数 $r$，且等于所预言的值。
\end{proof}

\begin{lemma}
\label{chow-lemma-compute-section}
设 $(S, \delta)$ 如情形 \ref{chow-situation-setup}。
设 $X$ 是 $S$ 上局部有限型概形，且 $\mathcal{C}$
是秩 $r$ 的局部自由 $\mathcal{O}_X$-模。考虑态射
$$
X = \underline{\text{Proj}}_X(\mathcal{O}_X[T])
\xrightarrow{i}
E = \underline{\text{Proj}}_X(\text{Sym}^*(\mathcal{C})[T])
\xrightarrow{\pi}
X
$$
对 $t = 1, \ldots, r - 1$ 有
$c_t(i_*\mathcal{O}_X) = 0$，并且在 $A^0(C \to E)$ 中有
$$
p^* \circ \pi_* \circ c_r(i_*\mathcal{O}_X) = (-1)^{r - 1}(r - 1)! j^*
$$
其中 $j : C \to E$ 和 $p : C \to X$
分别是向量丛
$C = \underline{\Spec}(\text{Sym}^*(\mathcal{C}))$
的包含态射和结构态射。
\end{lemma}

\begin{proof}
典范映射 $\pi^*\mathcal{C} \to \mathcal{O}_E(1)$
恰好沿 $i(X)$ 消失。因此，映射
$$
\pi^*\mathcal{C} \otimes \mathcal{O}_E(-1) \to \mathcal{O}_E
$$
上的 Koszul 复形是 $i_*\mathcal{O}_X$ 的消解。
特别地，$i_*\mathcal{O}_X$ 是 $D(\mathcal{O}_E)$ 中陈类已有定义的
完美对象。对 $t = 1, \ldots, t - 1$，
$c_t(i_*\mathcal{O}_X)$ 的消失由
引理 \ref{chow-lemma-compute-koszul} 得出。该引理还给出
$$
c_r(i_*\mathcal{O}_X) = - (r - 1)!
c_r(\pi^*\mathcal{C} \otimes \mathcal{O}_E(-1))
$$
另一方面，由引理 \ref{chow-lemma-chern-classes-dual}，
$$
c_r(\pi^*\mathcal{C} \otimes \mathcal{O}_E(-1)) =
(-1)^r c_r(\pi^*\mathcal{C}^\vee \otimes \mathcal{O}_E(1))
$$
而 $\pi^*\mathcal{C}^\vee \otimes \mathcal{O}_E(1)$
有一个恰好沿 $i(X)$ 消失的截面 $s$。

\medskip\noindent
将 $X$ 替换为 $X$ 上局部有限型概形后，只须证明等式两边
对元素 $\alpha \in \CH_*(E)$ 的作用相同。由于
$C \to X$ 是向量丛，$C$ 上每个周类都形如
$p^*\beta$，其中 $\beta \in \CH_*(X)$
（引理 \ref{chow-lemma-vectorbundle}）。因此由
引理 \ref{chow-lemma-restrict-to-open}，可写
$\alpha = \pi^*\beta + \gamma$，其中 $\gamma$
支撑在 $E \setminus C$ 上。使用上述等式，只须证明
$$
p^*(\pi_*(c_r(\pi^*\mathcal{C}^\vee \otimes \mathcal{O}_E(1)) \cap [W])) =
j^*[W]
$$
其中 $W \subset E$ 是整闭子概形，并且或者 (a) 与 $C$ 不交，
或者 (b) 形如 $W = \pi^{-1}Y$，其中
$Y \subset X$ 是整闭子概形。使用截面 $s$ 和
引理 \ref{chow-lemma-top-chern-class}，在情形 (a) 有
$c_r(\pi^*\mathcal{C}^\vee \otimes \mathcal{O}_E(1)) \cap [W] = 0$，
在情形 (b) 有
$c_r(\pi^*\mathcal{C}^\vee \otimes \mathcal{O}_E(1)) \cap [W] = [i(Y)]$。
结论容易由此得出；省略细节。
\end{proof}

\begin{lemma}
\label{chow-lemma-agreement-with-loc-chern}
设 $(S, \delta)$ 如情形 \ref{chow-situation-setup}。
设 $i : Z \to X$ 是 $S$ 上局部有限型概形之间余维 $r$
的正则闭浸入。令 $\mathcal{N} = \mathcal{C}_{Z/X}^\vee$
为法层。若 $X$ 拟紧（或其不可约分支拟紧），则对
$t = 1, \ldots, r - 1$ 有
$c_t(Z \to X, i_*\mathcal{O}_Z) = 0$，并且
$$
c_r(Z \to X, i_*\mathcal{O}_Z) = (-1)^{r - 1} (r - 1)! c(Z \to X, \mathcal{N})
\quad\text{于}\quad
A^r(Z \to X)
$$
其中 $c_t(Z \to X, i_*\mathcal{O}_Z)$ 是定义 \ref{chow-definition-localized-chern} 中的局部化陈类。
\end{lemma}

\begin{proof}
对任意 $x \in Z$，可选取仿射开邻域
$\Spec(A) \subset X$，使得
$Z \cap \Spec(A) = V(f_1, \ldots, f_r)$，其中
$f_1, \ldots, f_r \in A$ 是正则序列。见《除子》定义 \ref{divisors-definition-regular-immersion} 和引理 \ref{divisors-lemma-Noetherian-scheme-regular-ideal}。
于是由例如《代数进阶》引理 \ref{more-algebra-lemma-regular-koszul-regular}，
$f_1, \ldots, f_r$ 上的 Koszul 复形是
$A/(f_1, \ldots, f_r)$ 的消解。因此
$A/(f_1, \ldots, f_r)$ 作为 $A$-模是完美的。
由此 $F = i_*\mathcal{O}_Z$ 是 $D(\mathcal{O}_X)$ 中的完美对象，
其在 $X \setminus Z$ 上的限制为零。$X$ 拟紧
（或其不可约分支拟紧）的假设意味着局部化陈类
$c_t(Z \to X, i_*\mathcal{O}_Z)$ 已有定义；见
情形 \ref{chow-situation-loc-chern} 和定义 \ref{chow-definition-localized-chern}。综上，该陈述是有意义的。

\medskip\noindent
记 $b : W \to \mathbf{P}^1_X$ 为第 \ref{chow-section-blowup-Z-first} 节中沿 $\infty(Z)$ 的爆破。
由 (\ref{chow-item-find-Z-in-blowup})，存在闭浸入
$$
i' : \mathbf{P}^1_Z \longrightarrow W
$$
我们断言 $Q = i'_*\mathcal{O}_{\mathbf{P}^1_Z}$
是 $D(\mathcal{O}_W)$ 中的完美对象，并且 $F$、$Q$
满足引理 \ref{chow-lemma-independent-loc-chern-bQ} 的假设。

\medskip\noindent
假设该断言成立。引理 \ref{chow-lemma-independent-loc-chern-bQ} 给出：对所有
$p \geq 1$，
$$
c_p(Z \to X, F) = c'_p(Q) = (E \to Z)_* \circ c'_p(Q|_E) \circ C
$$
注意，$Q|_E$ 等于 $Z$ 的结构层沿态射 $Z \to E$ 的推前，
而该态射是 $i'$ 沿 $\infty$ 的基变换。
因此，将引理 \ref{chow-lemma-compute-section} 应用于 $E \to Z$，
便得到 $1 \leq t \leq r - 1$ 时
$c_t(Z \to X, F)$ 的消失。由于
$\mathcal{C}_{Z/X} = \mathcal{N}^\vee$ 局部自由，双变类
$c(Z \to X, \mathcal{N})$ 由关系
$$
j^* \circ C = p^* \circ c(Z \to X, \mathcal{N})
$$
刻画，其中 $j : C_ZX \to W_\infty$ 和 $p : C_ZX \to Z$
为给定映射。（回忆 $C \in A^0(W_\infty \to X)$ 是
引理 \ref{chow-lemma-gysin-at-infty} 中的类。）
因此，引理陈述中的显示等式由
引理 \ref{chow-lemma-compute-section} 中的相应等式得出。

\medskip\noindent
证明该断言。设 $A$ 和 $f_1, \ldots, f_r$ 如上。
考虑第 \ref{chow-section-blowup-Z-first} 节中的仿射开集
$\Spec(A[s]) \subset \mathbf{P}^1_X$。回忆在该开集上，
$s = 0$ 定义 $(\mathbf{P}^1_X)_\infty$。因此在
$\Spec(A[s])$ 上，我们沿正则序列
$s, f_1, \ldots, f_r$ 生成的理想爆破。由《代数进阶》引理 \ref{more-algebra-lemma-blowup-regular-sequence}，这 $r + 1$
个仿射图均为 $A[s]$ 上的整体完全交。与仿射爆破代数
$$
A[s][f_1/s, \ldots, f_r/s] = A[s, y_1, \ldots, y_r]/(sy_i - f_i)
$$
对应的图包含 $i'(Z \cap \Spec(A))$，后者是由
$y_1, \ldots, y_r$ 截出的闭子概形。由于
$y_1, \ldots, y_r, sy_1 - f_1, \ldots, sy_r - f_r$
是多项式环 $A[s, y_1, \ldots, y_r]$ 中的正则序列，
可知 $i'$ 是正则浸入。省略一些细节。
如上可得 $Q = i'_*\mathcal{O}_{\mathbf{P}^1_Z}$
是 $D(\mathcal{O}_W)$ 中的完美对象。$F$ 与 $Q$ 的其他假设在
引理 \ref{chow-lemma-independent-loc-chern-bQ}
（以及引理 \ref{chow-lemma-localized-chern-pre}）中均立即得到验证。
\end{proof}

\begin{lemma}
\label{chow-lemma-actual-computation}
在引理 \ref{chow-lemma-agreement-with-loc-chern} 的情形中，设
$\dim_\delta(X) = n$。则
\begin{enumerate}
\item 当 $t = 1, \ldots, r - 1$ 时，
$c_t(Z \to X, i_*\mathcal{O}_Z) \cap [X]_n = 0$；
\item $c_r(Z \to X, i_*\mathcal{O}_Z) \cap [X]_n =
(-1)^{r - 1}(r - 1)![Z]_{n - r}$；
\item 当 $t = 0, \ldots, r - 1$ 时，
$ch_t(Z \to X, i_*\mathcal{O}_Z) \cap [X]_n = 0$；并且
\item $ch_r(Z \to X, i_*\mathcal{O}_Z) \cap [X]_n = [Z]_{n - r}$。
\end{enumerate}
\end{lemma}

\begin{proof}
第 (1)、(2) 项立即由引理 \ref{chow-lemma-agreement-with-loc-chern} 与引理 \ref{chow-lemma-gysin-fundamental} 合并得出。
然后利用引理 \ref{chow-lemma-loc-chern-character} 中
$ch_p = (1/p!)P_p$ 与 $c_p$ 之间的关系推出第 (3)、(4) 项。
（具体而言，若 $c_1 = c_2 = \ldots = c_{r - 1} = 0$，则
$(-1)^{r - 1}(r - 1)!ch_r = c_r$。）
\end{proof}







\section{一个 Adams 算子}
\label{chow-section-adams}

\noindent
我们只做定义第二个 Adams 算子所需的最少工作。
设 $X$ 是概形。回忆 $\textit{Vect}(X)$ 表示有限局部自由
$\mathcal{O}_X$-模的范畴。此外，回忆在
《概形的导出范畴》第 \ref{perfect-section-K-groups} 节中，已经构造了与该范畴对应的第零
$K$-群 $K_0(\textit{Vect}(X))$。最后，
$K_0(\textit{Vect}(X))$ 是环；见
《概形的导出范畴》注 \ref{perfect-remark-K-ring}。

\begin{lemma}
\label{chow-lemma-second-adams-operator}
设 $X$ 是概形。存在环映射
$$
\psi^2 :
K_0(\textit{Vect}(X))
\longrightarrow
K_0(\textit{Vect}(X))
$$
当 $\mathcal{L}$ 可逆时，它将 $[\mathcal{L}]$ 映到
$[\mathcal{L}^{\otimes 2}]$，并且与拉回相容。
\end{lemma}

\begin{proof}
设 $X$ 是概形，并设 $\mathcal{E}$ 是有限局部自由
$\mathcal{O}_X$-模。考虑 $K_0(\textit{Vect}(X))$ 的元素
$$
\psi^2(\mathcal{E}) = [\text{Sym}^2(\mathcal{E})] - [\wedge^2(\mathcal{E})]
$$

\medskip\noindent
设 $X$ 是概形，并考虑有限局部自由 $\mathcal{O}_X$-模的短正合列
$$
0 \to \mathcal{E} \to \mathcal{F} \to \mathcal{G} \to 0
$$
将其视为 $\mathcal{F}$ 上具有 $2$ 步的滤过。
$\text{Sym}^2(\mathcal{F})$ 上诱导的滤过具有 $3$ 步，其分次片为
$\text{Sym}^2(\mathcal{E})$、$\mathcal{E} \otimes \mathcal{F}$
以及 $\text{Sym}^2(\mathcal{G})$。因此
$$
[\text{Sym}^2(\mathcal{F})] =
[\text{Sym}^2(\mathcal{E})] +
[\mathcal{E} \otimes \mathcal{F}] +
[\text{Sym}^2(\mathcal{G})]
$$
完全类似地可证
$$
[\wedge^2(\mathcal{F})] =
[\wedge^2(\mathcal{E})] +
[\mathcal{E} \otimes \mathcal{F}] +
[\wedge^2(\mathcal{G})]
$$
由此
$\psi^2(\mathcal{F}) = \psi^2(\mathcal{E}) + \psi^2(\mathcal{G})$。
因而得到良定义的加法映射
$\psi^2 : K_0(\textit{Vect}(X)) \to K_0(\textit{Vect}(X))$。

\medskip\noindent
显然，该映射与拉回交换。

\medskip\noindent
还须证明 $\psi^2$ 是环映射。设 $X$ 是概形，且
$\mathcal{E}$、$\mathcal{F}$ 是有限局部自由
$\mathcal{O}_X$-模。注意存在短正合列
$$
0 \to \wedge^2(\mathcal{E}) \otimes \wedge^2(\mathcal{F}) \to
\text{Sym}^2(\mathcal{E} \otimes \mathcal{F}) \to
\text{Sym}^2(\mathcal{E}) \otimes \text{Sym}^2(\mathcal{F}) \to 0
$$
其中第一个映射将
$(e \wedge e') \otimes (f \wedge f')$ 映到
$(e \otimes f)(e' \otimes f') - (e' \otimes f)(e \otimes f')$，
第二个映射将 $(e \otimes f) (e' \otimes f')$ 映到
$ee' \otimes ff'$。类似地，存在短正合列
$$
0 \to \text{Sym}^2(\mathcal{E}) \otimes \wedge^2(\mathcal{F}) \to
\wedge^2(\mathcal{E} \otimes \mathcal{F}) \to
\wedge^2(\mathcal{E}) \otimes \text{Sym}^2(\mathcal{F}) \to 0
$$
其中第一个映射将 $e e' \otimes f \wedge f'$ 映到
$(e \otimes f) \wedge (e' \otimes f') + (e' \otimes f) \wedge (e \otimes f')$，
第二个映射将
$(e \otimes f) \wedge (e' \otimes f')$ 映到
$(e \wedge e') \otimes (f f')$。
如上可知映射 $\psi^2$ 具有乘法性。又显然
$\psi^2(1) = 1$，证明完毕。
\end{proof}

\begin{remark}
\label{chow-remark-adams-derived}
设 $X$ 是概形，并且 $2$ 在 $X$ 上可逆。
则可以直接在 $K$-群
$K_0(X) = K_0(D_{perf}(\mathcal{O}_X))$
上定义 Adams 算子 $\psi^2$（《概形的导出范畴》定义 \ref{perfect-definition-K-group}）。
具体而言，给定 $X$ 上的完美复形 $L$，交换因子便给出群
$\{\pm 1\}$ 在 $L \otimes^\mathbf{L} L$ 上的作用。于是可以定义
$$
\psi^2(L) = [(L \otimes^\mathbf{L} L)^+] -
[(L \otimes^\mathbf{L} L)^-]
$$
其中 $(-)^+$ 表示取不变量，$(-)^-$ 表示取反不变量
（作适当定义）。利用取不变量和反不变量的正合性，可以仿照
引理 \ref{chow-lemma-second-adams-operator} 的证明说明它良定义。
当 $2$ 在 $X$ 上不可逆时，情形复杂得多，必须采用另一种方法。
\end{remark}

\begin{lemma}
\label{chow-lemma-minus-adams-operator}
设 $X$ 是概形。存在环映射
$\psi^{-1} : K_0(\textit{Vect}(X)) \to K_0(\textit{Vect}(X))$；
当 $\mathcal{E}$ 有限局部自由时，它将 $[\mathcal{E}]$ 映到
$[\mathcal{E}^\vee]$，并与拉回相容。
\end{lemma}

\begin{proof}
只须检验取对偶与短正合列及拉回相容。这是清楚的。
\end{proof}

\begin{remark}
\label{chow-remark-chern-classes-K}
设 $(S, \delta)$ 如情形 \ref{chow-situation-setup}，且
$X$ 在 $S$ 上局部有限型。陈类映射定义典范映射
$$
c : K_0(\textit{Vect}(X)) \longrightarrow \prod\nolimits_{i \geq 0} A^i(X)
$$
方法是将左边的生成元 $[\mathcal{E}]$ 映到
$c(\mathcal{E}) = 1 + c_1(\mathcal{E}) + c_2(\mathcal{E}) + \ldots$，
再作乘法延拓。因此 $-[\mathcal{E}]$ 被映到形式逆元
$c(\mathcal{E})^{-1}$；这也解释了右边为何取无限乘积。
由引理 \ref{chow-lemma-additivity-chern-classes}，该映射良定义。
\end{remark}

\begin{remark}
\label{chow-remark-chern-character-K}
设 $(S, \delta)$ 如情形 \ref{chow-situation-setup}，且
$X$ 在 $S$ 上局部有限型。陈特征映射定义典范环映射
$$
ch : K_0(\textit{Vect}(X)) \longrightarrow
\prod\nolimits_{i \geq 0} A^i(X) \otimes \mathbf{Q}
$$
方法是将左边的生成元 $[\mathcal{E}]$ 映到
$ch(\mathcal{E})$，再作加法延拓。由引理 \ref{chow-lemma-chern-character-additive}，该映射良定义；由
引理 \ref{chow-lemma-chern-character-multiplicative}，它是环同态。
\end{remark}

\begin{lemma}
\label{chow-lemma-adams-and-chern}
设 $(S, \delta)$ 如情形 \ref{chow-situation-setup}，且
$X$ 在 $S$ 上局部有限型。设 $\psi^2$ 如引理 \ref{chow-lemma-second-adams-operator}，而 $c$ 与 $ch$ 如注 \ref{chow-remark-chern-classes-K} 和 \ref{chow-remark-chern-character-K}。
则对所有 $\alpha \in K_0(\textit{Vect}(X))$，有
$c_i(\psi^2(\alpha)) = 2^i c_i(\alpha)$ 以及
$ch_i(\psi^2(\alpha)) = 2^i ch_i(\alpha)$。
\end{lemma}

\begin{proof}
注意，在 $A^i(X)$ 上乘以 $2^i$ 的映射
$\prod_{i \geq 0} A^i(X) \to \prod_{i \geq 0} A^i(X)$
是环映射。由于 $\psi^2$ 也是环映射，只须对
$K_0(\textit{Vect}(X))$ 的加法生成元证明这些公式。
因此可假设 $\alpha = [\mathcal{E}]$，其中 $\mathcal{E}$
是某个有限局部自由 $\mathcal{O}_X$-模。
根据 $\mathcal{E}$ 的陈类构造，立即化归为
$\mathcal{E}$ 的秩恒为 $r$ 的情形；见注 \ref{chow-remark-extend-to-finite-locally-free}。
此时可选取投影光滑态射 $p : P \to X$，使得限制
$A^*(X) \to A^*(P)$ 是单射，并且 $p^*\mathcal{E}$
有有限滤过，其各分次片是可逆 $\mathcal{O}_P$-模
$\mathcal{L}_j$；见引理 \ref{chow-lemma-splitting-principle}。于是
$[p^*\mathcal{E}] = \sum [\mathcal{L}_j]$，从而按
$\psi^2$ 的定义，
$\psi^2([p^\mathcal{E}]) = \sum [\mathcal{L}_j^{\otimes 2}]$。
令 $x_j  = c_1(\mathcal{L}_j)$，则在 $\prod A^i(P)$ 中有
$$
c(\alpha) = \prod (1 + x_j)
\quad\text{且}\quad
c(\psi^2(\alpha)) = \prod (1 + 2 x_j)
$$
并且在 $\prod A^i(P)$ 中有
$$
ch(\alpha) = \sum \exp(x_j)
\quad\text{且}\quad
ch(\psi^2(\alpha)) = \sum \exp(2 x_j)
$$
所需结论由这些公式得出。
\end{proof}

\begin{remark}
\label{chow-remark-perf-Z-cohomology-K}
设 $X$ 是局部 Noether 概形，且 $Z \subset X$ 是闭子概形。
考虑严格满、饱和的三角子范畴
$$
D_{Z, perf}(\mathcal{O}_X) \subset D(\mathcal{O}_X)
$$
其对象是由 $\mathcal{O}_X$-模组成的完美复形，并且其上同调层
在集合论意义下支撑在 $Z$ 上。记
$\textit{Coh}_Z(X) \subset \textit{Coh}(X)$ 为凝聚
$\mathcal{O}_X$-模的 Serre 子范畴，其中各模的集合论支撑包含在
$Z$ 中。注意，给定
$E \in D_{Z, perf}(\mathcal{O}_X)$，在 $X$ 上 Zariski 局部，
只有有限多个上同调层 $H^i(E)$ 非零
（并且它们在集合论意义下都支撑在 $Z$ 上）。因此可按规则
$$
E \longmapsto
[\bigoplus\nolimits_{i \in \mathbf{Z}} H^{2i}(E)] -
[\bigoplus\nolimits_{i \in \mathbf{Z}} H^{2i + 1}(E)]
$$
定义映射
$$
K_0(D_{Z, perf}(\mathcal{O}_X))
\longrightarrow
K_0(\textit{Coh}_Z(X)) = K'_0(Z)
$$
（等号由引理 \ref{chow-lemma-K-coherent-supported-on-closed} 得出）。
这是可行的，因为 $D_{Z, perf}(\mathcal{O}_X)$ 中每个可区别三角
都给出上同调层的长正合列。
\end{remark}

\begin{remark}
\label{chow-remark-perf-Z-regular}
设 $X$、$Z$、$D_{Z, perf}(\mathcal{O}_X)$ 如注 \ref{chow-remark-perf-Z-cohomology-K}，并假设 $X$ 正则。
则存在典范映射
$$
K_0(\textit{Coh}(Z)) \longrightarrow K_0(D_{Z, perf}(\mathcal{O}_X))
$$
定义如下。对任意凝聚 $\mathcal{O}_Z$-模 $\mathcal{F}$，
以 $\mathcal{F}[0]$ 表示 $D(\mathcal{O}_X)$ 中次数 $0$
处为 $\mathcal{F}$、其他次数均为零的对象。
由《概形的导出范畴》引理 \ref{perfect-lemma-perfect-on-regular}，$\mathcal{F}[0]$
是 $X$ 上的完美复形。因此 $\mathcal{F}[0]$ 是
$D_{Z, perf}(\mathcal{O}_X)$ 的对象。另一方面，给定凝聚
$\mathcal{O}_Z$-模的短正合列
$0 \to \mathcal{F} \to \mathcal{F}' \to \mathcal{F}'' \to 0$，
得到可区别三角
$\mathcal{F}[0] \to \mathcal{F}'[0] \to \mathcal{F}''[0] \to \mathcal{F}[1]$；
见《导出范畴》第 \ref{derived-section-canonical-delta-functor} 节。这表明，将
$[\mathcal{F}]$ 映到 $[\mathcal{F}[0]]$ 定义了映射
$K_0(\textit{Coh}(Z)) \to K_0(D_{Z, perf}(\mathcal{O}_X))$。
为这个可怕的记号致歉。
\end{remark}

\begin{lemma}
\label{chow-lemma-perf-Z-regular}
设 $X$ 是 Noether 正则概形，且 $Z \subset X$ 是闭子概形。
注 \ref{chow-remark-perf-Z-cohomology-K} 和
\ref{chow-remark-perf-Z-regular} 中构造的映射互为逆映射，从而
$K'_0(Z) = K_0(D_{Z, perf}(\mathcal{O}_X))$。
\end{lemma}

\begin{proof}
显然，复合
$$
K_0(\textit{Coh}(Z)) \longrightarrow
K_0(D_{Z, perf}(\mathcal{O}_X)) \longrightarrow
K_0(\textit{Coh}(Z))
$$
是恒等映射。因此只须证明第一个箭头为满射。
设 $E$ 是 $D_{Z, perf}(\mathcal{O}_X)$ 的对象。
回忆由《概形的导出范畴》引理 \ref{perfect-lemma-perfect-on-regular}，
$D_{perf}(\mathcal{O}_X) = D^b_{\textit{Coh}}(\mathcal{O}_X)$。
因此各上同调 $H^i(E)$ 都凝聚，可视为
$D_{Z, perf}(\mathcal{O}_X)$ 的对象，并且仅有有限多个非零。
利用典范截断的可区别三角可见
$$
[E] = \sum (-1)^i[H^i(E)[0]]
$$
在 $K_0(D_{Z, perf}(\mathcal{O}_X))$ 中成立。
于是只须对任意在集合论意义下支撑在 $Z$ 上的凝聚
$\mathcal{O}_X$-模证明 $[\mathcal{F}[0]]$ 位于该映射的像中。
由于可在 $\mathcal{F}$ 上找到有限滤过，其各次商都是
$\mathcal{O}_Z$-模，证明完毕。
\end{proof}

\begin{remark}
\label{chow-remark-localized-chern-classes-K}
设 $(S, \delta)$ 如情形 \ref{chow-situation-setup}。
设 $X$ 在 $S$ 上局部有限型，$Z \subset X$ 是闭子概形，并设
$D_{Z, perf}(\mathcal{O}_X)$ 如注 \ref{chow-remark-perf-Z-cohomology-K}。若 $X$ 拟紧
（或更一般地，$X$ 的不可约分支拟紧），则局部化陈类定义典范映射
$$
c(Z \to X, -) : K_0(D_{Z, perf}(\mathcal{O}_X)) \longrightarrow
A^0(X) \times \prod\nolimits_{i \geq 1} A^i(Z \to X)
$$
方法是将左边的生成元 $[E]$ 映到
$$
c(Z \to X, E) = 1 + c_1(Z \to X, E) + c_2(Z \to X, E) + \ldots
$$
再作乘法延拓（右边的乘积如注 \ref{chow-remark-ring-loc-classes}）。$X$ 的拟紧条件保证
局部化陈类已有定义（情形 \ref{chow-situation-loc-chern} 和
定义 \ref{chow-definition-localized-chern}），并且这些局部化陈类
将可区别三角转化为双变 Chow 环中的相应乘积
（引理 \ref{chow-lemma-additivity-loc-chern-c}）。
\end{remark}

\begin{remark}
\label{chow-remark-localized-chern-character-K}
设 $(S, \delta)$ 如情形 \ref{chow-situation-setup}。
设 $X$ 在 $S$ 上局部有限型，$Z \subset X$ 是闭子概形，并设
$D_{Z, perf}(\mathcal{O}_X)$ 如注 \ref{chow-remark-perf-Z-cohomology-K}。若 $X$ 的不可约分支拟紧，
则局部化陈特征定义典范的加法且乘法映射
$$
ch(Z \to X, -) : K_0(D_{Z, perf}(\mathcal{O}_X)) \longrightarrow
\prod\nolimits_{i \geq 0} A^i(Z \to X) \otimes \mathbf{Q}
$$
方法是将左边的生成元 $[E]$ 映到
$ch(Z \to X, E)$，再作加法延拓。具体而言，$X$
的不可约分支满足的条件保证局部化陈特征已有定义
（情形 \ref{chow-situation-loc-chern} 和定义 \ref{chow-definition-localized-chern}），并且这些局部化陈特征
将可区别三角转化为双变 Chow 环中的相应和
（引理 \ref{chow-lemma-additivity-loc-chern-P}）。
$K_0(D_{Z, perf}(X))$ 上的乘法用导出张量积定义
（《概形的导出范畴》注 \ref{perfect-remark-perf-Z}），所以由引理 \ref{chow-lemma-loc-chern-tensor-product}，
$ch(Z \to X, \alpha \beta) = ch(Z \to X, \alpha) ch(Z \to X, \beta)$。
若 $X$ 拟紧，则映射 $ch(Z \to X, -)$ 的像包含在
$A^*(Z \to X) \otimes \mathbf{Q}$ 中；省略细节。
\end{remark}

\begin{remark}
\label{chow-remark-chern-classes-agree}
设 $(S, \delta)$ 如情形 \ref{chow-situation-setup}。
设 $X$ 在 $S$ 上局部有限型，并假设 $X$ 拟紧
（或更一般地，$X$ 的不可约分支拟紧）。令 $Z = X$，并采用
注 \ref{chow-remark-localized-chern-classes-K} 和
\ref{chow-remark-localized-chern-character-K} 中的记号，则
$D_{Z, perf}(\mathcal{O}_X) = D_{perf}(\mathcal{O}_X)$，且
$$
K_0(D_{Z, perf}(\mathcal{O}_X)) = K_0(D_{perf}(\mathcal{O}_X)) = K_0(X)
$$
见《概形的导出范畴》定义 \ref{perfect-definition-K-group}。因此，作为注 \ref{chow-remark-localized-chern-classes-K} 和
\ref{chow-remark-localized-chern-character-K} 的特殊情形，得到
$$
c : K_0(X) \to \prod A^i(X)
\quad\text{且}\quad
ch : K_0(X) \to \prod A^i(X) \otimes \mathbf{Q}
$$
当然，也可以直接使用定义 \ref{chow-definition-defined-on-perfect}、引理 \ref{chow-lemma-additivity-on-perfect} 和
\ref{chow-lemma-chern-classes-perfect-tensor-product} 构造这些映射
（立即可见这会产生相同的类）。
回忆存在典范映射 $K_0(\textit{Vect}(X)) \to K_0(X)$，
它将有限局部自由模映到作为完美复形的自身
（置于次数 $0$）；见《概形的导出范畴》第 \ref{perfect-section-K-groups} 节。于是图
$$
\xymatrix{
K_0((\textit{Vect}(X)) \ar[rd]_c \ar[rr] & &
K_0(D_{perf}(\mathcal{O}_X)) = K_0(X) \ar[ld]^c \\
& \prod A^i(X)
}
$$
交换，其中东南方向的箭头是在注 \ref{chow-remark-chern-classes-K} 中构造的。类似地，图
$$
\xymatrix{
K_0((\textit{Vect}(X)) \ar[rd]_{ch} \ar[rr] & &
K_0(D_{perf}(\mathcal{O}_X)) = K_0(X) \ar[ld]^{ch} \\
& \prod A^i(X) \otimes \mathbf{Q}
}
$$
交换，其中东南方向的箭头是在注 \ref{chow-remark-chern-character-K} 中构造的。
\end{remark}











\section{重访 Chow 群与 K 群}
\label{chow-section-chow-and-K-II}

\noindent
本节是第 \ref{chow-section-chow-and-K} 节的继续。
设 $(S, \delta)$ 如情形 \ref{chow-situation-setup}，且
$X$ 在 $S$ 上局部有限型。$X$ 上凝聚层的 K 群
$K'_0(X) = K_0(\textit{Coh}(X))$ 具有典范递增滤过
$$
F_kK'_0(X) =
\Im\Big(K_0(\textit{Coh}_{\leq k}(X)) \to K_0(\textit{Coh}(X)\Big)
$$
这称为按支撑维数的滤过。注意
$$
\text{gr}_k K'_0(X) \subset K'_0(X)/F_{k - 1}K'_0(X) =
K_0(\textit{Coh}(X)/\textit{Coh}_{\leq k - 1}(X))
$$
其中等号由《同调代数》引理 \ref{homology-lemma-serre-subcategory-K-groups} 得出。
注 \ref{chow-remark-good-cases-K-A} 中的讨论表明，存在典范映射
$$
\CH_k(X) \longrightarrow \text{gr}_k K'_0(X)
$$
它将 $\delta$-维数为 $k$ 的整闭子概形
$Z \subset X$ 的类映到右边的类 $[\mathcal{O}_Z]$。

\begin{proposition}
\label{chow-proposition-K-tensor-Q}
设 $(S, \delta)$ 如情形 \ref{chow-situation-setup}。
假设给定 $S$ 上局部有限型概形之间的闭浸入
$X \to Y$，其中 $Y$ 正则且拟紧。考虑以下三个映射的复合：
注 \ref{chow-remark-perf-Z-regular} 中的映射
$\mathcal{F} \mapsto \mathcal{F}[0]$，注 \ref{chow-remark-localized-chern-character-K} 中的映射
$ch(X \to Y, -)$，以及映射 $c \mapsto c \cap [Y]$；即
$$
K'_0(X) \to
K_0(D_{X, perf}(\mathcal{O}_Y)) \to
A^*(X \to Y) \otimes \mathbf{Q} \to
\CH_*(X) \otimes \mathbf{Q}
$$
该复合诱导同构
$$
K'_0(X) \otimes \mathbf{Q}
\longrightarrow
\CH_*(X) \otimes \mathbf{Q}
$$
此同构依赖于 $Y$ 的选择。此外，上述典范映射
$$
\CH_k(X) \otimes \mathbf{Q}
\longrightarrow
\text{gr}_k K'_0(X) \otimes \mathbf{Q}
$$
对所有 $k \in \mathbf{Z}$ 都是 $\mathbf{Q}$-向量空间同构。
\end{proposition}

\begin{proof}
由于 $Y$ 正则，可以应用注 \ref{chow-remark-perf-Z-regular} 中的构造。由于 $Y$ 拟紧，
可以应用注 \ref{chow-remark-localized-chern-character-K} 中的构造。
$Y$ 局部等维（引理 \ref{chow-lemma-locally-equidimensional}），
因而“基本周” $[Y]$ 已定义为 $\CH_*(Y)$ 的元素；见注 \ref{chow-remark-fundamental-class}。结合上面构造的映射
$\CH_k(X) \to \text{gr}_kK'_0(X)$，只须证明：
\begin{enumerate}
\item 若 $\mathcal{F}$ 是凝聚 $\mathcal{O}_X$-模，
其支撑的 $\delta$-维数 $\leq k$，则上述复合将
$[\mathcal{F}]$ 映入
$\bigoplus_{k' \leq k} \CH_{k'}(X) \otimes \mathbf{Q}$；
\item 若 $Z \subset X$ 是 $\delta$-维数为 $k$ 的整闭子概形，
则上述复合将 $[\mathcal{O}_Z]$ 映到某个元素，其次数 $k$
部分是 $\CH_k(X) \otimes \mathbf{Q}$ 中的类 $[Z]$。
\end{enumerate}
事实上，若这些断言成立，则我们的映射诱导映射
$\text{gr}_kK'_0(X) \otimes \mathbf{Q} \to CH_k(X) \otimes \mathbf{Q}$，
它们与命题上方给出的典范映射
$\CH_k(X) \otimes \mathbf{Q} \to \text{gr}_k K'_0(X) \otimes \mathbf{Q}$
互为逆映射。

\medskip\noindent
给定凝聚 $\mathcal{O}_X$-模 $\mathcal{F}$，上述复合将
$[\mathcal{F}]$ 映到
$$
ch(X \to Y, \mathcal{F}[0]) \cap [Y] \in \CH_*(X) \otimes \mathbf{Q}
$$
若 $\mathcal{F}$ 在集合论意义下支撑于闭子概形
$Z \subset X$，则由引理 \ref{chow-lemma-loc-chern-shrink-Z}，
$$
ch(X \to Y, \mathcal{F}[0]) = (Z \to X)_* \circ ch(Z \to Y, \mathcal{F}[0])
$$
因此在此情形下，所得元素位于
$\CH_*(Z) \to \CH_*(X)$ 的像中。这证明了条件 (1)。

\medskip\noindent
设 $Z \subset X$ 是 $\delta$-维数为 $k$ 的整闭子概形。
由与上面相同的论证，上述复合将 $[\mathcal{O}_Z]$ 映到元素
$$
ch(X \to Y, \mathcal{O}_Z[0]) \cap [Y] =
(Z \to X)_* ch(Z \to Y, \mathcal{O}_Z[0]) \cap [Y]
$$
因此只须证明
$ch(Z \to Y, \mathcal{O}_Z[0]) \cap [Y] \in
\CH_*(Z) \otimes \mathbf{Q}$ 的次数 $k$ 部分为 $[Z]$。
由于 $\CH_k(Z) = \mathbf{Z}$，为证明此事，可将 $Y$
替换为 $Z$ 的泛点 $\xi$ 的一个开邻域。由于正则局部环
$\mathcal{O}_{X, \xi}$ 的极大理想由正则序列生成
（《交换代数》引理 \ref{algebra-lemma-regular-ring-CM}），
可假设 $Z$ 的理想由正则序列生成；见《除子》引理 \ref{divisors-lemma-Noetherian-scheme-regular-ideal}。
于是结论由引理 \ref{chow-lemma-actual-computation} 得出。
\end{proof}







\section{正则概形上的有理交乘积}
\label{chow-section-intersection-regular}

\noindent
我们将证明：若 $X$ 是具有仿射对角态射的 Noether 正则有限维概形，
则 $\CH_*(X) \otimes \mathbf{Q}$ 上存在交乘积。
该构造以如下结果为基础（它是上一节命题的推论）。

\begin{lemma}
\label{chow-lemma-K-tensor-Q}
设 $(S, \delta)$ 如情形 \ref{chow-situation-setup}。
设 $X$ 是 $S$ 上有限型的拟紧正则概形，具有仿射对角态射，
且 $\delta_{X/S} : X \to \mathbf{Z}$ 有界。则
注 \ref{chow-remark-chern-character-K} 中的映射 $ch$
与映射 $c \mapsto c \cap [X]$ 的复合
$$
K_0(\textit{Vect}(X)) \otimes \mathbf{Q}
\longrightarrow
A^*(X) \otimes \mathbf{Q}
\longrightarrow
\CH_*(X) \otimes \mathbf{Q}
$$
是同构。
\end{lemma}

\begin{proof}
由《概形的导出范畴》引理 \ref{perfect-lemma-Kprime-K}、
\ref{perfect-lemma-regular-resolution-property} 和
\ref{perfect-lemma-K-is-old-K}，有
$K'_0(X) = K_0(X) = K_0(\textit{Vect}(X))$。
由注 \ref{chow-remark-chern-classes-agree}，当 $X = Y$ 时，
给定复合与命题 \ref{chow-proposition-K-tensor-Q} 中的映射一致。
故结论由该命题得出。
\end{proof}

\noindent
设 $X, S, \delta$ 如引理 \ref{chow-lemma-K-tensor-Q}。
为简单起见，考虑给定余维的周；见第 \ref{chow-section-cycles-codimension} 节。令 $[X]$ 为 $X$ 的基本周；
见注 \ref{chow-remark-fundamental-class}。
选取 $\alpha \in CH^i(X)$ 和 $\beta \in CH^j(X)$。
由该引理，存在唯一的
$\alpha' \in K_0(\textit{Vect}(X)) \otimes \mathbf{Q}$，使得
$ch(\alpha') \cap [X]  = \alpha$。这当然意味着：若
$i' \not = i$，则 $ch_{i'}(\alpha') \cap [X] = 0$；并且
$ch_i(\alpha') \cap [X] = \alpha$。
由引理 \ref{chow-lemma-adams-and-chern}，
$\alpha'' = 2^{-i}\psi^2(\alpha')$ 是另一个解。
由唯一性，$\alpha'' = \alpha'$，从而当 $i' \not = i$ 时，
$ch_{i'}(\alpha) = 0$ 在
$A^{i'}(X) \otimes \mathbf{Q}$ 中成立。于是可以定义
$$
\alpha \cdot \beta = ch(\alpha') \cap \beta =
ch_i(\alpha') \cap \beta
$$
作为 $\CH^{i + j}(X) \otimes \mathbf{Q}$ 的元素；
这是由上面观察到的 $\alpha'$ 的性质保证的。
该配对是对称的：事实上，若选取提升 $\beta$ 的
$\beta' \in K_0(\textit{Vect}(X)) \otimes \mathbf{Q}$，则
$$
\alpha \cdot \beta = ch(\alpha') \cap \beta =
ch(\alpha') \cap ch(\beta') \cap [X]
$$
而我们知道陈类彼此交换。出于同样理由，交乘积满足结合律：
$$
(\alpha \cdot \beta) \cdot \gamma =
ch(\alpha') \cap ch(\beta') \cap ch(\gamma') \cap [X]
$$
因为双变类的复合满足结合律。也许更好的表述如下：
$\CH^*(X) \otimes \mathbf{Q}$ 上存在唯一的、与分次相容的交换结合交乘积，
使得同构
$K_0(\textit{Vect}(X)) \otimes \mathbf{Q} \to \CH^*(X) \otimes \mathbf{Q}$
是环同构。






\section{局部完全交态射的 Gysin 映射}
\label{chow-section-koszul}

\noindent
阅读本节前，建议读者先了解正则浸入
（《除子》第 \ref{divisors-section-regular-immersions} 节）
以及局部完全交态射
（《态射进阶》第 \ref{more-morphisms-section-lci} 节）。

\medskip\noindent
设 $(S, \delta)$ 如情形 \ref{chow-situation-setup}。
设 $i : X \to Y$ 是 $S$ 上局部有限型概形之间的正则浸入\footnote{见
《除子》定义 \ref{divisors-definition-regular-immersion}。
请注意，对于局部 Noether 概形，正则浸入与
Koszul-正则浸入或拟正则浸入是同一概念；见
《除子》引理 \ref{divisors-lemma-regular-immersion-noetherian}。
本节中将不再另行说明地使用这一事实。}。
特别地，余法层 $\mathcal{C}_{X/Y}$ 是有限局部自由的
（见《除子》引理 \ref{divisors-lemma-quasi-regular-immersion}）。
因此法层
$$
\mathcal{N}_{X/Y} = \SheafHom_{\mathcal{O}_X}(\mathcal{C}_{X/Y}, \mathcal{O}_X)
$$
也是有限局部自由的，并且存在满射
$\mathcal{N}_{X/Y}^\vee \to \mathcal{C}_{X/Y}$
（因为同构当然也是满射）。
第 \ref{chow-section-gysin-higher-codimension} 节中的构造给出规范双变类
$$
i^! = c(X \to Y, \mathcal{N}_{X/Y}) \in A^*(X \to Y)^\wedge
$$
我们需要关于这一概念的几个引理。

\begin{lemma}
\label{chow-lemma-composition-regular-immersion}
设 $(S, \delta)$ 如情形 \ref{chow-situation-setup}。
设 $i : X \to Y$ 和 $j : Y \to Z$ 是 $S$ 上局部有限型概形之间的正则浸入。
则 $j \circ i$ 是正则浸入，并且
$(j \circ i)^! = i^! \circ j^!$。
\end{lemma}

\begin{proof}
第一项断言即《除子》引理 \ref{divisors-lemma-composition-regular-immersion}。
由《除子》引理 \ref{divisors-lemma-transitivity-conormal-quasi-regular}，
存在短正合列
$$
0 \to
i^*(\mathcal{C}_{Y/Z}) \to
\mathcal{C}_{X/Z} \to
\mathcal{C}_{X/Y} \to 0
$$
因此结论由更一般的引理 \ref{chow-lemma-gysin-composition} 得出。
\end{proof}

\begin{lemma}
\label{chow-lemma-section-smooth}
设 $(S, \delta)$ 如情形 \ref{chow-situation-setup}。
设 $p : P \to X$ 是 $S$ 上局部有限型概形之间的光滑态射，
且 $s : X \to P$ 是一个截面。则 $s$ 是正则浸入，并且在
$A^*(X)^\wedge$ 中有 $1 = s^! \circ p^*$，其中
$p^* \in A^*(P \to X)^\wedge$ 是引理 \ref{chow-lemma-flat-pullback-bivariant} 的双变类。
\end{lemma}

\begin{proof}
第一项断言即《除子》引理 \ref{divisors-lemma-section-smooth-regular-immersion}。
由于这些假设在沿任意局部有限型态射 $X' \to X$ 作基变换时保持不变，
只需对任意整数闭子概形 $Z \subset X$ 证明
$\CH_*(X)$ 中 $s^! \cap p^*[Z] = [Z]$。
将 $P$ 替换为 $s(Z)$ 的一个开邻域后，可以假设
$P \to X$ 是固定相对维数 $r$ 的光滑态射。
记 $\dim_\delta(Z) = n$。则 $p^{-1}(Z)$ 的每个不可约分支的维数均为
$r + n$，且 $p^*[Z]$ 由 $[p^{-1}(Z)]_{n + r}$ 给出。
请注意，按概形论意义有 $s(X) \cap p^{-1}(Z) = s(Z)$。
因此，由上面所用的同一文献结果，$s(X) \cap p^{-1}(Z)$ 是
$\overline{p}^{-1}(Z)$ 中余维数为 $r$ 的正则嵌入闭子概形。
结论由引理 \ref{chow-lemma-gysin-fundamental} 得出。
\end{proof}

\noindent
设 $(S, \delta)$ 如情形 \ref{chow-situation-setup}。考虑 $S$ 上局部有限型概形的交换图
$$
\xymatrix{
X \ar[rd]_f \ar[rr]_i & & P \ar[ld]^g \\
& Y
}
$$
其中 $g$ 光滑而 $i$ 是正则浸入。将上面讨论的双变类 $i^!$ 与
引理 \ref{chow-lemma-flat-pullback-bivariant} 的双变类
$g^* \in A^*(P \to Y)^\wedge$ 复合，得到
$$
f^! = i^! \circ g^* \in A^*(X \to Y)
$$
请注意，态射 $f$ 是局部完全交态射；见《态射进阶》定义 \ref{more-morphisms-definition-lci}。
反之，若 $f : X \to Y$ 是局部 Noether 概形之间的局部完全交态射，
且 $f = g \circ i$、$g$ 光滑，则 $i$ 是正则浸入。
我们断言，对 $f^!$ 的这一构造仅依赖于态射 $f$，
而不依赖于分解 $f = g \circ i$ 的选择。

\begin{lemma}
\label{chow-lemma-lci-gysin-well-defined}
设 $(S, \delta)$ 如情形 \ref{chow-situation-setup}。
设 $f : X \to Y$ 是 $S$ 上局部有限型概形之间的局部完全交态射。
双变类 $f^!$ 与满足 $g$ 光滑的分解 $f = g \circ i$ 的选择无关
（只要这种分解存在）。
\end{lemma}

\begin{proof}
给定另一个这样的分解 $f = g' \circ i'$，可以考虑光滑态射
$g'' : P \times_Y P' \to Y$、浸入
$i'' : X \to P \times_Y P'$ 以及分解 $f = g'' \circ i''$。
因此可以假设有图
$$
\xymatrix{
& P' \ar[d]^p \ar[rd]^{g'} \\
X \ar[r]^i \ar[ru]^{i'} & P \ar[r]^g & Y
}
$$
其中 $p$ 是光滑态射。于是 $(g')^* = p^* \circ g^*$
（引理 \ref{chow-lemma-compose-flat-pullback}），故只需在
$A^*(X \to P)$ 中证明 $i^! = (i')^! \circ p^*$。
考虑交换图
$$
\xymatrix{
& X \times_P P' \ar[d]^{\overline{p}} \ar[r]_j & P' \ar[d]^p \\
X \ar[ru]^s \ar[r]^1 & X \ar[r]^i & P
}
$$
其中 $s =(1, i')$。则 $s$ 和 $j$ 都是正则浸入
（分别由《除子》引理 \ref{divisors-lemma-section-smooth-regular-immersion}
和《除子》引理 \ref{divisors-lemma-flat-base-change-regular-immersion}），
且 $i' = j \circ s$。由引理 \ref{chow-lemma-composition-regular-immersion}，
有 $(i')^! = s^! \circ j^!$。
由于该方块是笛卡尔的，双变类 $j^!$ 是 $i^!$ 在 $P'$ 上的限制
（注 \ref{chow-remark-restriction-bivariant}）；见
引理 \ref{chow-lemma-construction-gysin}。
由于双变类与平坦拉回可交换，得到
$j^! \circ p^* = \overline{p}^* \circ i^!$。
因此只需证明 $s^! \circ \overline{p}^* = \text{id}$，
这已由引理 \ref{chow-lemma-section-smooth} 完成。
\end{proof}

\begin{definition}
\label{chow-definition-lci-gysin}
设 $(S, \delta)$ 如情形 \ref{chow-situation-setup}。
设 $f : X \to Y$ 是 $S$ 上局部有限型概形之间的局部完全交态射。
若可以写成 $f = g \circ i$，其中 $g$ 光滑而 $i$ 是浸入，
则称{\it $f$ 的 Gysin 映射存在}。
在这种情形下，如上定义{\it Gysin 映射}
$f^! = i^! \circ g^* \in A^*(X \to Y)$。
\end{definition}

\noindent
由定义可知：对于正则浸入，该定义与前面的构造一致；
对于光滑态射，它与平坦拉回一致。事实上，
对所有 syntomic 态射，这种一致性都成立。

\begin{lemma}
\label{chow-lemma-lci-gysin-flat}
设 $(S, \delta)$ 如情形 \ref{chow-situation-setup}。
设 $f : X \to Y$ 是 $S$ 上局部有限型概形之间的局部完全交态射。
若 $f$ 的 Gysin 映射存在且 $f$ 平坦，则 $f^!$ 等于
引理 \ref{chow-lemma-flat-pullback-bivariant} 的双变类。
\end{lemma}

\begin{proof}
选取分解 $f = g \circ i$，其中 $i : X \to P$ 是浸入，
$g : P \to Y$ 光滑。请注意，对于任意局部有限型态射 $Y' \to Y$，
基变换 $f'$、$g'$、$i'$ 满足相同假设
（见《概形态射》引理 \ref{morphisms-lemma-base-change-smooth}、
\ref{morphisms-lemma-base-change-syntomic} 以及《态射进阶》引理 \ref{more-morphisms-lemma-flat-lci}）。
因此，由引理 \ref{chow-lemma-bivariant-zero}，归结为当 $Y$ 整时证明
$f^*[Y] = i^!(g^*[Y])$。令 $n = \dim_\delta(Y)$。
将 $X$ 和 $P$ 分解为连通分支后，可以假设 $f$ 是相对维数为 $r$ 的平坦态射，
而 $g$ 是相对维数为 $t$ 的光滑态射。
此时 $f^*[Y] = [X]_{n + s}$ 且 $g^*[Y] = [P]_{n + t}$。
另一方面，$i$ 是余维数为 $t - s$ 的正则浸入。
因此 $i^![P]_{n + t} = [X]_{n + s}$
（引理 \ref{chow-lemma-gysin-fundamental}），证明完成。
\end{proof}

\begin{lemma}
\label{chow-lemma-lci-gysin-composition}
设 $(S, \delta)$ 如情形 \ref{chow-situation-setup}。
设 $f : X \to Y$ 和 $g : Y \to Z$ 是 $S$ 上局部有限型概形之间的局部完全交态射。
假设 $g \circ f$ 和 $g$ 的 Gysin 映射存在。
则 $f$ 的 Gysin 映射存在，并且
$(g \circ f)^! = f^! \circ g^!$。
\end{lemma}

\begin{proof}
由《态射进阶》引理 \ref{more-morphisms-lemma-composition-lci}，$g \circ f$ 是局部完全交态射，
故引理的陈述是有意义的。
若 $X \to P$ 是 $X$ 到某个 $Z$ 上光滑概形 $P$ 中的浸入，
则 $X \to P \times_Z Y$ 是 $X$ 到某个 $Y$ 上光滑概形中的浸入。
这证明了引理的第一项断言。
设 $Y \to P'$ 是 $Y$ 到某个 $Z$ 上光滑概形 $P'$ 中的浸入。
考虑交换图
$$
\xymatrix{
X \ar[r] \ar[d] &
P \times_Z Y \ar[r]_a \ar[ld]^p &
P \times_Z P' \ar[ld]^q \\
Y \ar[r]_b \ar[d] &
P' \ar[ld] \\
Z
}
$$
其中水平箭头是正则浸入，西南方向的箭头光滑，且该方块是笛卡尔的。
由于双变类与平坦拉回可交换，因而
$a^! \circ q^* = p^* \circ b^!$。
将此事实与引理 \ref{chow-lemma-composition-regular-immersion} 和
\ref{chow-lemma-compose-flat-pullback} 结合，读者即可验证引理的陈述成立。
这里省略一个小细节。
\end{proof}

\begin{lemma}
\label{chow-lemma-lci-gysin-commutes}
设 $(S, \delta)$ 如情形 \ref{chow-situation-setup}。
考虑 $S$ 上局部有限型概形的交换图
$$
\xymatrix{
X'' \ar[d] \ar[r] &
X' \ar[d] \ar[r] &
X \ar[d]^f \\
Y'' \ar[r] &
Y' \ar[r] &
Y
}
$$
其中两个方块都是笛卡尔的。假设 $f : X \to Y$ 是局部完全交态射，
且 $f$ 的 Gysin 映射存在。令 $c \in A^*(Y'' \to Y')$。
以 $res(f^!) \in A^*(X' \to Y')$ 表示 $f^!$ 在 $Y'$ 上的限制
（注 \ref{chow-remark-restriction-bivariant}）。
则 $c$ 与 $res(f^!)$ 可交换（注 \ref{chow-remark-bivariant-commute}）。
\end{lemma}

\begin{proof}
选取分解 $f = g \circ i$，其中 $g$ 光滑而 $i$ 是浸入。
由于 $f^! = i^! \circ g^!$，只需分别对 $g^!$（它由平坦拉回给出）
和 $i^!$ 证明该引理。关于平坦拉回的结论是双变类定义的一部分；
关于 $i^!$ 的情形立即由引理 \ref{chow-lemma-gysin-commutes} 得出。
\end{proof}

\begin{lemma}
\label{chow-lemma-lci-gysin-easy}
设 $(S, \delta)$ 如情形 \ref{chow-situation-setup}。
考虑 $S$ 上局部有限型概形的笛卡尔图
$$
\xymatrix{
X' \ar[d]_{f'} \ar[r] &
X \ar[d]^f \\
Y' \ar[r] &
Y
}
$$
假设
\begin{enumerate}
\item $f$ 是局部完全交态射，且 $f$ 的 Gysin 映射存在；
\item $X$、$X'$、$Y$、$Y'$ 满足引理 \ref{chow-lemma-locally-equidimensional} 中的等价条件；
\item 对任意 $x' \in X'$，若其在 $X$、$Y'$、$Y$ 中的像分别为
$x$、$y'$、$y$，则 $n_{x'} - n_{y'} = n_x - n_y$，
其中 $n_{x'}$、$n_x$、$n_{y'}$、$n_y$ 如该引理所定义；以及
\item 对每个一般点 $\xi \in X'$，局部环
$\mathcal{O}_{Y', f'(\xi)}$ 都是 Cohen--Macaulay 的。
\end{enumerate}
则 $f^![Y'] = [X']$，其中 $[Y']$ 和 $[X']$ 如
注 \ref{chow-remark-fundamental-class} 所定义。
\end{lemma}

\begin{proof}
回忆 $n_{x'}$ 是 $\delta(\xi)$ 的共同值，这里 $\xi$ 遍历经过 $x'$ 的
不可约分支的一般点。此外，函数
$x' \mapsto n_{x'}$、$x \mapsto n_x$、$y' \mapsto n_{y'}$ 和
$y \mapsto n_y$ 都是局部常值的。分别以 $X'_n$、$X_n$、$Y'_n$、
$Y_n$ 表示 $X'$、$X$、$Y'$、$Y$ 中相应函数取值为 $n$ 的开闭子概形。
回忆 $[X'] = \sum [X'_n]_n$ 且 $[Y'] = \sum [Y'_n]_n$。
由此显然，为证明该引理，可以将 $X'$ 替换为它的一个连通分支，
并将 $X$、$Y'$、$Y$ 替换为该分支所映入的连通分支。
于是 $X'$、$X$、$Y'$、$Y$ 都是 $\delta$-等维的，
即各自的每个不可约分支都具有相同的 $\delta$-维数。
分别记 $X'$、$X$、$Y'$、$Y$ 的这一共同值为
$n'$、$n$、$m'$、$m$。最后一项假设意味着
$n' - m' = n - m$。

\medskip\noindent
选取分解 $f = g \circ i$，其中 $i : X \to P$ 是浸入，
$g : P \to Y$ 光滑。由于 $X$ 连通，$P \to Y$ 在 $i(X)$ 各点处的相对维数恒定。
因此，将 $P$ 替换为 $i(X)$ 的一个开邻域后，可以假设
$P \to Y$ 具有恒定相对维数，且 $i : X \to P$ 是闭浸入。
以 $g' : Y' \times_Y P \to Y'$ 表示 $g$ 的基变换，
以 $i' : X' \to Y' \times_Y P$ 表示 $i$ 的基变换。
显然 $g^*[Y] = [P]$ 且 $(g')^*[Y'] = [Y' \times_Y P]$。
最后，若 $\xi' \in X'$ 是一般点，则
$\mathcal{O}_{Y' \times_Y P, i'(\xi)}$ 是 Cohen--Macaulay 的。
事实上，局部环映射
$\mathcal{O}_{Y', f'(\xi)} \to \mathcal{O}_{Y' \times_Y P, i'(\xi)}$
平坦且具有正则纤维
（见《交换代数》第 \ref{algebra-section-smooth-overview} 节）；
正则局部环是 Cohen--Macaulay 的
（《交换代数》引理 \ref{algebra-lemma-regular-ring-CM}）；
由假设 (4)，$\mathcal{O}_{Y', f'(\xi)}$ 是 Cohen--Macaulay 的；
于是《交换代数》引理 \ref{algebra-lemma-CM-goes-up} 给出所需结论。
因此归结为下一段所讨论的情形。

\medskip\noindent
假设 $f$ 是正则闭浸入，并且 $X'$、$X$、$Y'$、$Y$ 分别是
$\delta$-维数为 $n'$、$n$、$m'$、$m$ 的 $\delta$-等维概形，且
$m' - n' = m - n$。在这种情形下，结论立即由
引理 \ref{chow-lemma-gysin-easy} 得出。
\end{proof}

\begin{remark}
\label{chow-remark-gysin-chern-classes}
设 $(S, \delta)$ 如情形 \ref{chow-situation-setup}。
设 $f : X \to Y$ 是 $S$ 上局部有限型概形之间的局部完全交态射，
并假设 $f$ 的 Gysin 映射存在。则
$f^! \circ c_i(\mathcal{E}) = c_i(f^*\mathcal{E}) \circ f^!$；
对于 Chern 特征标也有类似结论；见
引理 \ref{chow-lemma-lci-gysin-commutes}。
例如，若 $X$ 和 $Y$ 满足引理 \ref{chow-lemma-locally-equidimensional} 的等价条件，且 $Y$ 是
Cohen--Macaulay 的，则由引理 \ref{chow-lemma-lci-gysin-easy} 有
$f^![Y] = [X]$。在这种情形下还有
$f^!(c_i(\mathcal{E}) \cap [Y]) = c_i(f^*\mathcal{E}) \cap [X]$；
对于 Chern 特征标同样如此。
\end{remark}

\begin{lemma}
\label{chow-lemma-compare-gysin-base-change}
设 $(S, \delta)$ 如情形 \ref{chow-situation-setup}。
考虑 $S$ 上局部有限型概形的笛卡尔方块
$$
\xymatrix{
X' \ar[d]_{f'} \ar[r]_{g'} &
X \ar[d]^f \\
Y' \ar[r]^g &
Y
}
$$
假设
\begin{enumerate}
\item $f$ 和 $f'$ 都是局部完全交态射；以及
\item $f$ 的 Gysin 映射存在。
\end{enumerate}
则 $\mathcal{C} = \Ker(H^{-1}((g')^*\NL_{X/Y}) \to H^{-1}(\NL_{X'/Y'}))$
是有限局部自由 $\mathcal{O}_{X'}$-模，$f'$ 的 Gysin 映射存在，且
$$
res(f^!) = c_{top}(\mathcal{C}^\vee) \circ (f')^!
$$
在 $A^*(X' \to Y')$ 中成立。
\end{lemma}

\begin{proof}
$\mathcal{C}$ 有限局部自由这一事实立即由《代数进阶》引理 \ref{more-algebra-lemma-base-change-lci-bis} 得出。
选取分解 $f = h \circ i$，其中 $h : P \to Y$ 光滑而 $i$ 是浸入。
则可分解 $f' = h' \circ i'$，其中 $h' : P' \to Y'$ 和
$i' : X' \to P'$ 是相应的基变换。图示为
$$
\xymatrix{
X' \ar[r] \ar[d] &
P' \ar[r] \ar[d] &
Y' \ar[d] \\
X \ar[r]^i &
P \ar[r]^h &
Y
}
$$
特别地，$f'$ 的 Gysin 映射存在。由《态射进阶》引理 \ref{more-morphisms-lemma-get-NL}，有
$$
\NL_{X/Y} = \left( \mathcal{C}_{X/P} \to i^*\Omega_{P/Y} \right)
$$
其中 $\mathcal{C}_{X/P}$ 是嵌入 $i$ 的余法层；带撇版本同理。
由《概形态射》引理 \ref{morphisms-lemma-base-change-differentials}，
$\Omega_{P/Y}$ 拉回为 $\Omega_{P'/Y'}$，故
$(g')^*i^*\Omega_{P/Y} = (i')^*\Omega_{P'/Y'}$。
还要回忆，$(g')^*\mathcal{C}_{X/P} \to \mathcal{C}_{X'/P'}$ 是满射；见
《概形态射》引理 \ref{morphisms-lemma-conormal-functorial-flat}。
由此推出，层 $\mathcal{C}$ 规范同构于有限局部自由模之间映射
$(g')^*\mathcal{C}_{X/P} \to \mathcal{C}_{X'/P'}$ 的核。
回忆 $i^!$ 是利用 $\mathcal{N} = \mathcal{C}_{Z/X}^\vee$ 定义的，
$(i')^!$ 亦然。因此应用引理 \ref{chow-lemma-gysin-excess} 得到
$$
res(i^!) = c_{top}(\mathcal{C}^\vee) \circ (i')^!
$$
在 $A^*(X' \to P')$ 中成立。最后，由
$f^! = i^! \circ h^*$、$(f')^! = (i')^! \circ (h')^*$ 以及
$(h')^* = res(h^*)$，即得结论。
\end{proof}

\begin{lemma}[爆破公式]
\label{chow-lemma-blow-up-formula}
设 $(S, \delta)$ 如情形 \ref{chow-situation-setup}。
设 $i : Z \to X$ 是 $S$ 上局部有限型概形之间的正则闭浸入，
$b : X' \to X$ 是以 $Z$ 为中心的爆破。图示为
$$
\xymatrix{
E \ar[r]_j \ar[d]_\pi & X' \ar[d]^b \\
Z \ar[r]^i & X
}
$$
假设 $b$ 的 Gysin 映射存在。则
$$
res(b^!) = c_{top}(\mathcal{F}^\vee) \circ \pi^*
$$
在 $A^*(E \to Z)$ 中成立，其中 $\mathcal{F}$ 是规范映射
$\pi^*\mathcal{C}_{Z/X} \to \mathcal{C}_{E/X'}$ 的核。
\end{lemma}

\begin{proof}
由《代数进阶》引理 \ref{more-algebra-lemma-blowup-regular-sequence}，
态射 $b$ 是局部完全交态射，故该陈述有意义。
由于 $Z \to X$ 是正则浸入（因而更是拟正则浸入），
$\mathcal{C}_{Z/X}$ 有限局部自由，且映射
$\text{Sym}^*(\mathcal{C}_{Z/X}) \to \mathcal{C}_{Z/X, *}$
是同构；见《除子》引理 \ref{divisors-lemma-quasi-regular-immersion}。
由 $E = \text{Proj}(\mathcal{C}_{Z/X, *})$ 可知
$E = \mathbf{P}(\mathcal{C}_{Z/X})$ 是 $Z$ 上的射影空间丛。
因此 $E \to Z$ 光滑，当然也是局部完全交态射。
于是引理 \ref{chow-lemma-compare-gysin-base-change} 适用，并给出
$$
res(b^!) = c_{top}(\mathcal{C}^\vee) \circ \pi^!
$$
其中 $\mathcal{C}$ 如该引理的陈述所定义。
当然，由引理 \ref{chow-lemma-lci-gysin-flat}，有 $\pi^* = \pi^!$。
只需再证明 $\mathcal{F}$ 等于映射
$H^{-1}(j^*\NL_{X'/X}) \to H^{-1}(\NL_{E/Z})$ 的核 $\mathcal{C}$。

\medskip\noindent
由于 $E \to Z$ 光滑，有 $H^{-1}(\NL_{E/Z}) = 0$；见
《态射进阶》引理 \ref{more-morphisms-lemma-NL-smooth}。
因此只需证明 $\mathcal{F}$ 可与 $H^{-1}(j^*\NL_{X'/X})$ 等同。
由《态射进阶》引理 \ref{more-morphisms-lemma-get-triangle-NL} 和
\ref{more-morphisms-lemma-NL-immersion}，有正合列
$$
0 \to H^{-1}(j^*\NL_{X'/X}) \to H^{-1}(\NL_{E/X}) \to
\mathcal{C}_{E/X'} \to \ldots
$$
将同样的引理应用于 $E \to Z \to X$，得到同构
$\pi^*\mathcal{C}_{Z/X} = H^{-1}(\pi^*\NL_{Z/X}) \to H^{-1}(\NL_{E/X})$。
于是结论成立。
\end{proof}

\begin{lemma}
\label{chow-lemma-lci-gysin-product-regular}
设 $(S, \delta)$ 如情形 \ref{chow-situation-setup}。
设 $f : X \to Y$ 是 $S$ 上局部有限型概形之间的态射，
并且 $X$ 和 $Y$ 都拟紧、正则、具有仿射对角态射且维数有限。
则 $f$ 是局部完全交态射。再假设 $f$ 的 Gysin 映射存在，则
$$
f^!(\alpha \cdot \beta) = f^!\alpha \cdot f^!\beta
$$
在 $\CH^*(X) \otimes \mathbf{Q}$ 中成立，其中交乘积如
第 \ref{chow-section-intersection-regular} 节所定义。
\end{lemma}

\begin{proof}
第一项断言由《态射进阶》引理 \ref{more-morphisms-lemma-morphism-regular-schemes-is-lci} 得出。
请注意，由引理 \ref{chow-lemma-lci-gysin-easy}，有 $f^![Y] = [X]$。
如第 \ref{chow-section-intersection-regular} 节，写成
$\alpha = ch(\alpha') \cap [Y]$ 和 $\beta = ch(\beta') \cap [Y]$，
其中 $\alpha', \beta' \in K_0(\textit{Vect}(X)) \otimes \mathbf{Q}$。
令 $c = ch(\alpha')$、$c' = ch(\beta')$；由构造可知
$\alpha \cdot \beta = c \cap c' \cap [Y]$。
由引理 \ref{chow-lemma-lci-gysin-commutes}，$f^!$ 与 $c$、$c'$ 均可交换。因此
\begin{align*}
f^!(\alpha \cdot \beta)
& =
f^!(c \cap c' \cap [Y]) \\
& =
c \cap c' \cap f^![Y] \\
& =
c \cap c' \cap [X] \\
& =
(c \cap [X]) \cdot (c' \cap [X]) \\
& =
(c \cap f^![Y]) \cdot (c' \cap f^![Y]) \\
& =
f^!(\alpha) \cdot f^!(\beta)
\end{align*}
这正是所需结论。
\end{proof}

\begin{lemma}
\label{chow-lemma-projection-formula-regular}
设 $(S, \delta)$ 如情形 \ref{chow-situation-setup}。
设 $f : X \to Y$ 是 $S$ 上局部有限型概形之间的态射，
并且 $X$ 和 $Y$ 都拟紧、正则、具有仿射对角态射且维数有限。
则 $f$ 是局部完全交态射。再假设 $f$ 的 Gysin 映射存在且 $f$ 固有，则
$$
f_*(\alpha \cdot f^!\beta) = f_*\alpha \cdot \beta
$$
在 $\CH^*(Y) \otimes \mathbf{Q}$ 中成立，其中交乘积如
第 \ref{chow-section-intersection-regular} 节所定义。
\end{lemma}

\begin{proof}
第一项断言由《态射进阶》引理 \ref{more-morphisms-lemma-morphism-regular-schemes-is-lci} 得出。
请注意，由引理 \ref{chow-lemma-lci-gysin-easy}，有 $f^![Y] = [X]$。
如第 \ref{chow-section-intersection-regular} 节，写成
$\alpha = ch(\alpha') \cap [X]$ 和 $\beta = ch(\beta') \cap [Y]$，
其中 $\alpha' \in K_0(\textit{Vect}(X)) \otimes \mathbf{Q}$ 且
$\beta' \in K_0(\textit{Vect}(Y)) \otimes \mathbf{Q}$。
令 $c = ch(\alpha')$、$c' = ch(\beta')$。于是
\begin{align*}
f_*(\alpha \cdot f^!\beta)
& =
f_*(c \cap f^!(c' \cap [Y]_e)) \\
& =
f_*(c \cap c' \cap f^![Y]_e) \\
& =
f_*(c \cap c' \cap [X]_d) \\
& =
f_*(c' \cap c \cap [X]_d) \\
& =
c' \cap f_*(c \cap [X]_d) \\
& =
\beta \cdot f_*(\alpha)
\end{align*}
第一项等式来自交乘积的构造。由引理 \ref{chow-lemma-lci-gysin-commutes}，$f^!$ 与 $c'$ 可交换。
Chern 类位于双变环的中心，这使得可以交换以 $c$ 和 $c'$ 对 $[X]$ 作帽积的次序。
由于 $c'$ 是双变类，允许将 $c'$ 与 $f_*$ 交换。最后一项等式再次来自交乘积的构造。
\end{proof}













\section{对角态射的 Gysin 映射}
\label{chow-section-gysin-for-diagonal}

\noindent
设 $(S, \delta)$ 如情形 \ref{chow-situation-setup}。
设 $f : X \to Y$ 是 $S$ 上局部有限型概形之间的光滑态射。
则对角态射 $\Delta : X \longrightarrow X \times_Y X$ 是正则浸入；见
《态射进阶》引理 \ref{more-morphisms-lemma-smooth-diagonal-perfect}。
因此有第 \ref{chow-section-koszul} 节中构造的 Gysin 映射
$$
\Delta^! \in A^*(X \to X \times_Y X)^\wedge
$$
若 $X \to Y$ 具有恒定相对维数 $d$，则
$\Delta^! \in A^d(X \to X \times_Y X)$。

\begin{lemma}
\label{chow-lemma-diagonal-identity}
在上述情形中，$A^0(X)$ 内有 $\Delta^! \circ \text{pr}_i^! = 1$。
\end{lemma}

\begin{proof}
请注意，投影 $\text{pr}_i : X \times_Y X \to X$ 都是光滑的，
因而这些投影也有 Gysin 映射。因此该引理的陈述有意义，
并且它是引理 \ref{chow-lemma-lci-gysin-composition} 的特殊情形。
\end{proof}

\begin{proposition}
\label{chow-proposition-compute-bivariant}
\begin{reference}
\cite[Proposition 17.4.2]{F}
\end{reference}
设 $(S, \delta)$ 如情形 \ref{chow-situation-setup}。
设 $f : X \to Y$ 和 $g : Y \to Z$ 是 $S$ 上局部有限型概形之间的态射。
若 $g$ 是相对维数为 $d$ 的光滑态射，则
$A^p(X \to Y) = A^{p - d}(X \to Z)$。
\end{proposition}

\begin{proof}
我们将使用如下事实：光滑态射是 Gysin 映射存在的局部完全交态射
（见第 \ref{chow-section-koszul} 节）。特别地，有
$g^! \in A^{-d}(Y \to Z)$。于是可以将
$c \in A^p(X \to Y)$ 映到 $c \circ g^! \in A^{p - d}(X \to Z)$。

\medskip\noindent
反之，令 $c' \in A^{p - d}(X \to Z)$。
以 $res(c')$ 表示 $c'$ 沿态射 $Y \to Z$ 的限制
（注 \ref{chow-remark-restriction-bivariant}）。由于图
$$
\xymatrix{
X \times_Z Y \ar[r]_{\text{pr}_2} \ar[d]_{\text{pr}_1} & Y \ar[d]^g \\
X \ar[r]^f & Z
}
$$
是笛卡尔的，得到 $res(c') \in A^{p - d}(X \times_Z Y \to Y)$。
令 $\Delta : Y \to Y \times_Z Y$ 为对角态射，
并以 $res(\Delta^!)$ 表示 $\Delta^!$ 沿态射
$X \times_Z Y \to Y \times_Z Y$ 在 $X \times_Z Y$ 上的限制。
由于图
$$
\xymatrix{
X \ar[r] \ar[d] & X \times_Z Y \ar[d] \\
Y \ar[r]^-\Delta & Y \times_Z Y
}
$$
是笛卡尔的，可见 $res(\Delta^!) \in A^d(X \to X \times_Z Y)$。
将这两个限制复合，得到
$$
res(\Delta^!) \circ res(c') \in A^p(X \to Y)
$$
于是构造了映射 $A^p(X \to Y) \to A^{p - d}(X \to Z)$ 和
$A^{p - d}(X \to Z) \to A^p(X \to Y)$。
为完成证明，将说明这两个映射互为逆映射。

\medskip\noindent
先从 $c \in A^p(X \to Y)$ 开始。考虑图
$$
\xymatrix{
X \ar[d] \ar[r] & Y \ar[d] \\
X \times_Z Y \ar[r] \ar[d]^{\text{pr}_1} &
Y \times_Z Y \ar[r]_{p_2} \ar[d]^{p_1} &
Y \ar[d]^g \\
X \ar[r]^f &
Y \ar[r]^g &
Z
}
$$
其中各方块都是笛卡尔的。该图下方的两个方块表明
$res(c \circ g^!) = res(c) \cap p_2^!$；在此公式中，
$res(c)$ 表示 $c$ 沿 $p_1$ 的限制。
考察该图上方的方块并使用引理 \ref{chow-lemma-lci-gysin-commutes}，得到
$c \circ \Delta^! = res(\Delta^!) \circ res(c)$。计算得
\begin{align*}
res(\Delta^!) \circ res(c \circ g^!)
& =
res(\Delta^!) \circ res(c) \circ p_2^! \\
& =
c \circ \Delta^! \circ p_2^! \\
& =
c
\end{align*}
最后一项等式由引理 \ref{chow-lemma-diagonal-identity} 得出。

\medskip\noindent
反之，从 $c' \in A^{p - d}(X \to Z)$ 开始。
考察上图下方的矩形，得到
$res(c') \circ g^! = \text{pr}_1^! \circ c'$。计算得
\begin{align*}
res(\Delta^!) \circ res(c') \circ g^!
& =
res(\Delta^!) \circ \text{pr}_1^! \circ c' \\
& =
c'
\end{align*}
最后一项等式成立，是因为该图左侧两个方块表明
$\text{id} = res(\Delta^! \circ p_1^!) = res(\Delta^!) \circ \text{pr}_1^!$。
证明完成。
\end{proof}









\section{外积}
\label{chow-section-exterior-product}

\noindent
设 $k$ 是域。本节在 $S = \Spec(k)$ 上工作，其中
$\delta : S \to \mathbf{Z}$ 将唯一的点映到 $0$；见
例 \ref{chow-example-field}。

\medskip\noindent
考虑 $k$ 上局部有限型概形的笛卡尔方块
$$
\xymatrix{
X \times_k Y \ar[r] \ar[d] & Y \ar[d] \\
X \ar[r] & \Spec(k) = S
}
$$
则存在规范映射
$$
\times :
\CH_n(X) \otimes_{\mathbf{Z}} \CH_m(Y)
\longrightarrow
\CH_{n + m}(X \times_k Y)
$$
它由如下规则唯一确定：给定维数分别为 $n$ 和 $m$ 的整闭子概形
$X' \subset X$ 与 $Y' \subset Y$，有
$$
[X'] \times [Y'] = [X' \times_k Y']_{n + m}
$$
在 $\CH_{n + m}(X \times_k Y)$ 中成立。

\begin{lemma}
\label{chow-lemma-exterior-product-well-defined}
映射
$\times : \CH_n(X) \otimes_{\mathbf{Z}} \CH_m(Y) \to \CH_{n + m}(X \times_k Y)$
是良定义的。
\end{lemma}

\begin{proof}
先注意，若 $\alpha = \sum n_i[X_i]$ 且 $\beta = \sum m_j[Y_j]$，
其中 $X_i \subset X$ 和 $Y_j \subset Y$ 分别是维数为 $n$ 和 $m$ 的
整闭子概形的局部有限族，则 $X_i \times_k Y_j$ 构成
$X \times_k Y$ 中维数为 $n + m$ 的闭子概形的局部有限族，因而确实可以将
$$
\alpha \times \beta = \sum n_i m_j [X_i \times_k Y_j]_{n + m}
$$
视为 $X \times_k Y$ 上的一个 $(n + m)$-周。
由此得到加性映射
$\times : Z_n(X) \otimes_{\mathbf{Z}} Z_m(Y) \to Z_{n + m}(X \times_k Y)$。
问题在于证明这一过程与有理等价相容。

\medskip\noindent
设 $i : X' \to X$ 是一个维数为 $n$ 的整闭子概形的包含态射。
由引理 \ref{chow-lemma-flat-pullback-bivariant}，沿态射
$p' : X' \to \Spec(k)$ 的平坦拉回是元素
$(p')^* \in A^{-n}(X' \to \Spec(k))$；因而由
引理 \ref{chow-lemma-push-proper-bivariant}，有
$c' = i_* \circ (p')^* \in A^{-n}(X \to \Spec(k))$。
这给出映射
$$
c' \cap - : \CH_m(Y) \longrightarrow \CH_{m + n}(X \times_k Y)
$$
读者容易验证，对于任意维数为 $m$ 的整闭子概形 $Y' \subset Y$，
该映射把 $[Y']$ 映到 $[X' \times_k Y']_{n + m}$。
因此，构造 $([X'], [Y']) \mapsto [X' \times_k Y']_{n + m}$
在第二个变量中可通过有理等价下降，也就是说，它给出良定义映射
$Z_n(X) \otimes_{\mathbf{Z}} \CH_m(Y) \to \CH_{n + m}(X \times_k Y)$。
由对称性，另一个变量也同样如此，结论得证。
\end{proof}

\begin{lemma}
\label{chow-lemma-chow-cohomology-towards-point}
设 $k$ 是域，$X$ 是 $k$ 上局部有限型概形。
则对所有 $p \in \mathbf{Z}$，都有规范等同
$$
A^p(X \to \Spec(k)) = \CH_{-p}(X)
$$
\end{lemma}

\begin{proof}
考虑元素 $[\Spec(k)] \in \CH_0(\Spec(k))$。
把 $c$ 映到 $c \cap [\Spec(k)]$，即得到映射
$A^p(X \to \Spec(k)) \to \CH_{-p}(X)$。

\medskip\noindent
反之，设有 $\alpha \in \CH_{-p}(X)$。
可按如下方式定义 $c_\alpha \in A^p(X \to \Spec(k))$：
给定 $X' \to \Spec(k)$ 和 $\alpha' \in \CH_n(X')$，令
$$
c_\alpha \cap \alpha' = \alpha \times \alpha'
$$
作为 $\CH_{n - p}(X \times_k X')$ 中的元素。
为证明这是双变类，如定义 \ref{chow-definition-cycles} 写成
$\alpha = \sum n_i[X_i]$。考虑复合
$$
\coprod X_i \xrightarrow{g} X \to \Spec(k)
$$
并以 $f : \coprod X_i \to \Spec(k)$ 表示该复合。
则 $g$ 固有，而 $f$ 是相对维数为 $-p$ 的平坦态射。
由引理 \ref{chow-lemma-flat-pullback-bivariant}，沿 $f$ 的拉回是双变类
$f^* \in A^p(\coprod X_i \to \Spec(k))$。
以 $\nu \in A^0(\coprod X_i)$ 表示如下双变类：
它在第 $i$ 个分支上将周乘以 $n_i$。于是
$\nu \circ f^* \in A^p(\coprod X_i \to X)$。
最后，由引理 \ref{chow-lemma-push-proper-bivariant}，有双变类
$$
g_* \circ \nu \circ f^*
$$
读者容易验证 $c_\alpha$ 等于该双变类，因而它本身也是双变类。

\medskip\noindent
为完成证明，还须说明这两个构造互为逆映射。
由于 $c_\alpha \cap [\Spec(k)] = \alpha$，其中一个方向是显然的。
对另一个方向，令 $c \in A^p(X \to \Spec(k))$，并置
$\alpha = c \cap [\Spec(k)]$。由引理 \ref{chow-lemma-bivariant-zero}，
只需在 $X'$ 是 $\Spec(k)$ 上局部有限型整概形时证明
$$
c \cap [X'] = c_\alpha \cap [X']
$$
但此时 $p' : X' \to \Spec(k)$ 是相对维数为 $\dim(X')$ 的平坦态射，
故 $[X'] = (p')^*[\Spec(k)]$。
因此，双变类 $c$ 和 $c_\alpha$ 在 $[\Spec(k)]$ 上相同，
就蕴含它们与 $[X']$ 作帽积时也相同，证明完成。
\end{proof}

\begin{lemma}
\label{chow-lemma-chow-cohomology-towards-point-commutes}
设 $k$ 是域，$X$ 是 $k$ 上局部有限型概形。
设 $c \in A^p(X \to \Spec(k))$，$Y \to Z$ 是 $k$ 上局部有限型概形之间的态射，
且 $c' \in A^q(Y \to Z)$。则在 $A^{p + q}(X \times_k Y \to Z)$ 中有
$c \circ c' = c' \circ c$。
\end{lemma}

\begin{proof}
在引理 \ref{chow-lemma-chow-cohomology-towards-point} 的证明中已经看到，
$c$ 由固有前推、在各连通分支上乘以整数以及平坦拉回组合而成。
依双变类的定义，$c'$ 与这些运算中的每一个都可交换，故结论成立。
这里省略若干细节。
\end{proof}

\begin{remark}
\label{chow-remark-commuting-exterior}
引理 \ref{chow-lemma-chow-cohomology-towards-point} 和
\ref{chow-lemma-chow-cohomology-towards-point-commutes} 的要点如下。
设 $k$ 是域，$X$ 是 $k$ 上局部有限型概形。
设 $\alpha \in \CH_*(X)$，$Y \to Z$ 是 $k$ 上局部有限型概形之间的态射，
且 $c' \in A^q(Y \to Z)$。则对任意 $\beta \in \CH_*(Z)$，
$$
\alpha \times (c' \cap \beta) = c' \cap (\alpha \times \beta)
$$
在 $\CH_*(X \times_k Y)$ 中成立。事实上，只需取与 $\alpha$ 对应的双变类
$c = c_\alpha \in A^*(X \to \Spec(k))$；见引理 \ref{chow-lemma-chow-cohomology-towards-point} 的证明。
\end{remark}

\begin{lemma}
\label{chow-lemma-exterior-product-associative}
外积满足结合律。更确切地，设 $k$ 是域，
$X, Y, Z$ 是 $k$ 上局部有限型概形，并设
$\alpha \in \CH_*(X)$、$\beta \in \CH_*(Y)$、$\gamma \in \CH_*(Z)$。
则在 $\CH_*(X \times_k Y \times_k Z)$ 中有
$(\alpha \times \beta) \times \gamma =
\alpha \times (\beta \times \gamma)$。
\end{lemma}

\begin{proof}
从略。提示：概形纤维积的结合律。
\end{proof}







\section{交乘积}
\label{chow-section-intersection-product}

\noindent
设 $k$ 是域。本节在 $S = \Spec(k)$ 上工作，其中
$\delta : S \to \mathbf{Z}$ 将唯一的点映到 $0$；见
例 \ref{chow-example-field}。

\medskip\noindent
设 $X$ 是 $k$ 上光滑概形。第 \ref{chow-section-gysin-for-diagonal} 节的双变类
$\Delta^!$ 使我们能够在 $X$ 上局部有限型概形的 Chow 群上定义一种交乘积。
具体地，设 $Y \to X$ 和 $Z \to X$ 是局部有限型概形之间的态射。注意到
$$
Y \times_X Z =  (Y \times_k Z) \times_{X \times_k X, \Delta} X
$$
因此可以考虑映射序列
$$
\CH_n(Y) \otimes_\mathbf{Z} \CH_m(Z)
\xrightarrow{\times}
\CH_{n + m}(Y \times_k Z)
\xrightarrow{\Delta^!}
\CH_{n + m - *}(Y \times_X Z)
$$
其中第一支箭头是第 \ref{chow-section-exterior-product} 节中构造的外积，
第二支箭头是第 \ref{chow-section-gysin-for-diagonal} 节中研究的对角态射 Gysin 映射。
若 $X$ 是维数为 $d$ 的等维概形，则所得元素落在
$\CH_{n + m - d}(Y \times_X Z)$ 中；一般而言，可依照 $X$ 的开闭子概形
（其上 $X$ 具有给定维数），将其分解为位于这些部分之上的分量。
给定 $\alpha \in \CH_*(Y)$ 和 $\beta \in \CH_*(Z)$，记
$$
\alpha \cdot \beta = \Delta^!(\alpha \times \beta)
\in \CH_*(Y \times_X Z)
$$
在特殊情形 $X = Y = Z$ 下，得到乘法
$$
\CH_*(X) \times \CH_*(X) \to \CH_*(X),\quad
(\alpha, \beta) \mapsto \alpha \cdot \beta
$$
称为{\it 交乘积}。显然，这一乘积是对称的。
结合律由下一个引理给出。

\begin{lemma}
\label{chow-lemma-associative}
上面定义的乘积满足结合律。更确切地，设 $k$ 是域，
$X$ 在 $k$ 上光滑，$Y, Z, W$ 是 $X$ 上局部有限型概形，并设
$\alpha \in \CH_*(Y)$、$\beta \in \CH_*(Z)$、$\gamma \in \CH_*(W)$。
则在 $\CH_*(Y \times_X Z \times_X W)$ 中有
$(\alpha \cdot \beta) \cdot \gamma =
\alpha \cdot (\beta \cdot \gamma)$。
\end{lemma}

\begin{proof}
由引理 \ref{chow-lemma-exterior-product-associative}，在
$\CH_*(Y \times_k Z \times_k W)$ 中有
$(\alpha \times \beta) \times \gamma =
\alpha \times (\beta \times \gamma)$。考虑闭浸入
$$
\Delta_{12} : X \times_k X \longrightarrow X \times_k X \times_k X,
\quad (x, x') \mapsto (x, x, x')
$$
以及
$$
\Delta_{23} : X \times_k X \longrightarrow X \times_k X \times_k X,
\quad (x, x') \mapsto (x, x', x')
$$
以 $\Delta_{12}^!$ 和 $\Delta_{23}^!$ 表示相应的双变类。
注意，$\Delta_{12}^!$ 是 $\Delta^!$ 沿映射 $\text{pr}_{12}$ 在
$X \times_k X \times_k X$ 上的限制
（注 \ref{chow-remark-restriction-bivariant}），而
$\Delta_{23}^!$ 是 $\Delta^!$ 沿映射 $\text{pr}_{23}$ 在
$X \times_k X \times_k X$ 上的限制。
因此，$\Delta_{12}^!$ 沿 $\Delta_{23}$ 的限制显然是 $\Delta^!$，
而 $\Delta_{23}^!$ 沿 $\Delta_{12}$ 的限制同样是 $\Delta^!$。
于是由引理 \ref{chow-lemma-gysin-commutes}，有
$$
\Delta^! \circ \Delta_{12}^! =
\Delta^! \circ \Delta_{23}^! 
$$
现在用如下等式链证明该引理：
\begin{align*}
(\alpha \cdot \beta) \cdot \gamma
& =
\Delta^!(\Delta^!(\alpha \times \beta) \times \gamma) \\
& =
\Delta^!(\Delta_{12}^!((\alpha \times \beta) \times \gamma)) \\
& =
\Delta^!(\Delta_{23}^!((\alpha \times \beta) \times \gamma)) \\
& =
\Delta^!(\Delta_{23}^!(\alpha \times (\beta \times \gamma)) \\
& =
\Delta^!(\alpha \times \Delta^!(\beta \times \gamma)) \\
& =
\alpha \cdot (\beta \cdot \gamma)
\end{align*}
除第二项和倒数第二项等式或许需要说明外，其余等式均由上文显然成立。
等式
$\Delta_{23}^!(\alpha \times (\beta \times \gamma)) =
\alpha \times \Delta^!(\beta \times \gamma)$
由注 \ref{chow-remark-commuting-exterior} 成立；第二项等式同理。
\end{proof}

\begin{lemma}
\label{chow-lemma-identify-chow-for-smooth}
设 $k$ 是域，$X$ 是 $k$ 上维数为 $d$ 的等维光滑概形。映射
$$
A^p(X) \longrightarrow \CH_{d - p}(X),\quad
c \longmapsto c \cap [X]_d
$$
是同构。在这一同构下，双变类的复合变为上面定义的交乘积。
\end{lemma}

\begin{proof}
以 $g : X \to \Spec(k)$ 表示结构态射。该映射是同构的复合
$$
A^p(X) \to A^{p - d}(X \to \Spec(k)) \to \CH_{d - p}(X)
$$
第一支是命题 \ref{chow-proposition-compute-bivariant} 的同构
$c \mapsto c \circ g^*$，第二支是引理 \ref{chow-lemma-chow-cohomology-towards-point} 的同构
$c \mapsto c \cap [\Spec(k)]$。
由引理 \ref{chow-lemma-chow-cohomology-towards-point} 的证明，
第二支箭头的逆把 $\alpha \in \CH_{d - p}(X)$ 映到双变类
$c_\alpha$；对于 $k$ 上局部有限型的 $Y$ 和
$\beta \in \CH_*(Y)$，该双变类把它映到
$\CH_*(X \times_k Y)$ 中的 $\alpha \times \beta$。
再由命题 \ref{chow-proposition-compute-bivariant} 的证明，
第一支箭头的逆把 $c_\alpha$ 映到如下双变类：
对于局部有限型的 $Y \to X$ 和 $\beta \in \CH_*(Y)$，它把后者映到
$\Delta^!(\alpha \times \beta) = \alpha \cdot \beta$。
由此得到引理的最后一项结论。
\end{proof}

\begin{lemma}
\label{chow-lemma-lci-gysin-product}
设 $k$ 是域，$f : X \to Y$ 是 $k$ 上光滑概形之间的态射。
则 $f$ 的 Gysin 映射存在，并且
$f^!(\alpha \cdot \beta) = f^!\alpha \cdot f^!\beta$。
\end{lemma}

\begin{proof}
请注意，$X \to X \times_k Y$ 是 $X$ 到某个 $Y$ 上光滑概形中的浸入。
故 $f$ 的 Gysin 映射存在（定义 \ref{chow-definition-lci-gysin}）。
为证明该公式，可以把 $X$ 和 $Y$ 分解为连通分支；
因此可以假设 $X$ 是 $k$ 上维数为 $d$ 的等维光滑概形，
而 $Y$ 是 $k$ 上维数为 $e$ 的等维光滑概形。
请注意 $f^![Y]_e = [X]_d$（例如见引理 \ref{chow-lemma-lci-gysin-easy}）。
写成 $\alpha = c \cap [Y]_e$ 和 $\beta = c' \cap [Y]_e$，
于是 $\alpha \cdot \beta = c \cap c' \cap [Y]_e$；见
引理 \ref{chow-lemma-identify-chow-for-smooth}。
由引理 \ref{chow-lemma-lci-gysin-commutes}，$f^!$ 与 $c$、$c'$ 均可交换。因此
\begin{align*}
f^!(\alpha \cdot \beta)
& =
f^!(c \cap c' \cap [Y]_e) \\
& =
c \cap c' \cap f^![Y]_e \\
& =
c \cap c' \cap [X]_d \\
& =
(c \cap [X]_d) \cdot (c' \cap [X]_d) \\
& =
(c \cap f^![Y]_e) \cdot (c' \cap f^![Y]_e) \\
& =
f^!(\alpha) \cdot f^!(\beta)
\end{align*}
这正是所需结论；这里也对 $X$ 使用了引理 \ref{chow-lemma-identify-chow-for-smooth}。

\medskip\noindent
另一种证明方法是先证明
$(f \times f)^!(\alpha \times \beta) = f^!\alpha \times f^!\beta$，
再使用引理 \ref{chow-lemma-lci-gysin-composition}。
\end{proof}

\begin{lemma}
\label{chow-lemma-projection-formula}
设 $k$ 是域，$f : X \to Y$ 是 $k$ 上光滑概形之间的固有态射。
则 $f$ 的 Gysin 映射存在，并且
$f_*(\alpha \cdot f^!\beta) = f_*\alpha \cdot \beta$。
\end{lemma}

\begin{proof}
请注意，$X \to X \times_k Y$ 是 $X$ 到某个 $Y$ 上光滑概形中的浸入。
故 $f$ 的 Gysin 映射存在（定义 \ref{chow-definition-lci-gysin}）。
为证明该公式，可以把 $X$ 和 $Y$ 分解为连通分支；
因此可以假设 $X$ 是 $k$ 上维数为 $d$ 的等维光滑概形，
而 $Y$ 是 $k$ 上维数为 $e$ 的等维光滑概形。
请注意 $f^![Y]_e = [X]_d$（例如见引理 \ref{chow-lemma-lci-gysin-easy}）。
写成 $\alpha = c \cap [X]_d$ 和 $\beta = c' \cap [Y]_e$；见
引理 \ref{chow-lemma-identify-chow-for-smooth}。于是
\begin{align*}
f_*(\alpha \cdot f^!\beta)
& =
f_*(c \cap f^!(c' \cap [Y]_e)) \\
& =
f_*(c \cap c' \cap f^![Y]_e) \\
& =
f_*(c \cap c' \cap [X]_d) \\
& =
f_*(c' \cap c \cap [X]_d) \\
& =
c' \cap f_*(c \cap [X]_d) \\
& =
\beta \cdot f_*(\alpha)
\end{align*}
第一项等式来自对 $X$ 应用引理 \ref{chow-lemma-identify-chow-for-smooth} 的结果。
由引理 \ref{chow-lemma-lci-gysin-commutes}，$f^!$ 与 $c'$ 可交换。
交乘积的交换律使得可以交换以 $c$ 和 $c'$ 对 $[X]_d$ 作帽积的次序
（通过该引理）。由于 $c'$ 是双变类，允许将 $c'$ 与 $f_*$ 交换。
最后一项等式仍由该引理给出。
\end{proof}

\begin{lemma}
\label{chow-lemma-intersect-properly}
设 $k$ 是域，$X$ 是 $k$ 上光滑整概形。
设 $Y, Z \subset X$ 是整闭子概形，并置
$d = \dim(Y) + \dim(Z) - \dim(X)$。假设
\begin{enumerate}
\item $\dim(Y \cap Z) \leq d$；以及
\item 对每个满足 $\xi \in Y \cap Z$ 和 $\delta(\xi) = d$ 的点，
$\mathcal{O}_{Y, \xi}$ 与 $\mathcal{O}_{Z, \xi}$ 都是 Cohen--Macaulay 的。
\end{enumerate}
则在 $\CH_d(X)$ 中有 $[Y] \cdot [Z] = [Y \cap Z]_d$。
\end{lemma}

\begin{proof}
回忆 $[Y] \cdot [Z] = \Delta^!([Y \times Z])$，其中
$\Delta^! = c(\Delta : X \to X \times X, \mathcal{T}_{X/k})$
是高余维 Gysin 映射（第 \ref{chow-section-gysin-higher-codimension} 节），且
$\mathcal{T}_{X/k} = \SheafHom(\Omega_{X/k}, \mathcal{O}_X)$
是秩为 $\dim(X)$ 的局部自由层。概形之间有等式
$$
Y \cap Z = X \times_{\Delta, (X \times X)} (Y \times Z)
$$
并且 $\dim(Y \times Z) = \dim(Y) + \dim(Z)$，
故引理 \ref{chow-lemma-gysin-easy} 的条件 (1)、(2)、(3) 成立。
最后，若 $\xi \in Y \cap Z$，则有平坦局部同态
$$
\mathcal{O}_{Y, \xi} \longrightarrow
\mathcal{O}_{Y \times Z, \xi}
$$
其“纤维”为 $\mathcal{O}_{Z, \xi}$。因此，若
$\mathcal{O}_{Y, \xi}$ 和 $\mathcal{O}_{Z, \xi}$ 都是 Cohen--Macaulay 的，
则 $\mathcal{O}_{Y \times Z, \xi}$ 也是；见
《交换代数》引理 \ref{algebra-lemma-CM-goes-up}。
由此可见，引理 \ref{chow-lemma-gysin-easy} 的所有假设都满足，结论随即得出。
\end{proof}

\begin{lemma}
\label{chow-lemma-intersect-regularly-embedded}
设 $k$ 是域，$X$ 是 $k$ 上光滑概形，$i : Y \to X$ 是正则闭浸入，
且 $\alpha \in \CH_*(X)$。若 $Y$ 是维数为 $e$ 的等维概形，则
$\alpha \cdot [Y]_e = i_*(i^!(\alpha))$ 在 $\CH_*(X)$ 中成立。
\end{lemma}

\begin{proof}
将 $X$ 分解为连通分支后，可以并且确实假设 $X$ 是维数为 $d$ 的等维概形。
写成 $\alpha = c \cap [X]_n$，其中 $x \in A^*(X)$；见
引理 \ref{chow-lemma-identify-chow-for-smooth}。于是
$$
i_*(i^!(\alpha)) = i_*(i^!(c \cap [X]_n)) =
i_*(c \cap i^![X]_n) = i_*(c \cap [Y]_e) =
c \cap i_*[Y]_e = \alpha \cdot [Y]_e
$$
第一项等式来自 $c$ 的选择；第二项等式由
引理 \ref{chow-lemma-lci-gysin-commutes} 给出；第三项等式是因为
$\CH_*(Y)$ 中 $i^![X]_d = [Y]_e$
（引理 \ref{chow-lemma-lci-gysin-easy}）；第四项等式是因为双变类与固有前推可交换；
最后一项等式由引理 \ref{chow-lemma-identify-chow-for-smooth} 给出。
\end{proof}

\begin{lemma}
\label{chow-lemma-intersection-regular-smooth}
设 $k$ 是域，$X$ 是 $k$ 上拟紧且具有仿射对角态射的光滑概形。
则本节在 $\CH^*(X)$ 上构造的交乘积，在与 $\mathbf{Q}$ 张量后，
与第 \ref{chow-section-intersection-regular} 节中构造的交乘积一致。
\end{lemma}

\begin{proof}
设 $\alpha \in \CH^i(X)$ 且 $\beta \in \CH^j(X)$。
如第 \ref{chow-section-intersection-regular} 节，写成
$\alpha = ch(\alpha') \cap [X]$ 和 $\beta = ch(\beta') \cap [X]$，
其中 $\alpha', \beta' \in K_0(\textit{Vect}(X)) \otimes \mathbf{Q}$。
令 $c = ch(\alpha')$、$c' = ch(\beta')$。
于是第 \ref{chow-section-intersection-regular} 节中的交乘积给出
$c \cap c' \cap [X]$。由引理 \ref{chow-lemma-identify-chow-for-smooth}，
这与 $\alpha \cdot \beta$ 相同（更确切地说，是使用其推广：
对任意域 $k$ 上光滑概形 $X$，映射
$A^i(X) \to \CH^i(X)$、$c \mapsto c \cap [X]$ 是同构）。
\end{proof}








\section{Dedekind 整环上的外积}
\label{chow-section-exterior-product-dim-1}

\noindent
设 $S$ 是局部 Noether 概形，并且存在由 Dedekind 整环的谱组成的开覆盖。
对闭点 $s \in S$ 置 $\delta(s) = 0$，否则置 $\delta(s) = 1$。
则 $(S, \delta)$ 是一般情形 \ref{chow-situation-setup} 的特殊情形；
见例 \ref{chow-example-domain-dimension-1}。
请注意，$S$ 正规（《交换代数》引理 \ref{algebra-lemma-characterize-Dedekind}），
因而是正规整概形的不交并
（《概形的性质》引理 \ref{properties-lemma-normal-locally-Noetherian}）。
所以以下所有论证都归结为 $S$ 不可约的情形。
另一方面，我们允许 $S$ 非分离（例如，$S$ 可以是原点 $0$ 加倍的仿射直线）。

\medskip\noindent
考虑 $S$ 上局部有限型概形的笛卡尔方块
$$
\xymatrix{
X \times_S Y \ar[r] \ar[d] & Y \ar[d] \\
X \ar[r] & S
}
$$
我们断言存在规范映射
$$
\times :
\CH_n(X) \otimes_{\mathbf{Z}} \CH_m(Y)
\longrightarrow
\CH_{n + m - 1}(X \times_S Y)
$$
它由如下规则唯一确定：给定 $\delta$-维数分别为 $n$ 和 $m$ 的整闭子概形
$X' \subset X$ 与 $Y' \subset Y$，置
\begin{enumerate}
\item 若 $X'$ 或 $Y'$ 支配 $S$ 的一个不可约分支，则
$[X'] \times [Y'] = [X' \times_S Y']_{n + m - 1}$；
\item 若 $X'$ 和 $Y'$ 都不支配 $S$ 的任何不可约分支，则
$[X'] \times [Y'] = 0$。
\end{enumerate}

\begin{lemma}
\label{chow-lemma-exterior-product-well-defined-dim-1}
映射 $\times : \CH_n(X) \otimes_{\mathbf{Z}} \CH_m(Y) \to
\CH_{n + m - 1}(X \times_S Y)$ 是良定义的。
\end{lemma}

\begin{proof}
考虑 $n$-周和 $m$-周 $\alpha = \sum_{i \in I} n_i[X_i]$ 与
$\beta = \sum_{j \in J} m_j[Y_j]$，其中 $X_i \subset X$ 和 $Y_j \subset Y$
分别是 $\delta$-维数为 $n$ 和 $m$ 的整闭子概形的局部有限族。
令 $K \subset I \times J$ 为所有满足如下条件的序对
$(i, j) \in I \times J$ 的集合：$X_i$ 或 $Y_j$ 支配 $S$ 的一个不可约分支。
则 $\{X_i \times_S Y_j\}_{(i, j) \in K}$ 是 $X \times_S Y$ 中
$\delta$-维数为 $n + m - 1$ 的闭子概形的局部有限族。
这意味着确实可以将
$$
\alpha \times \beta =
\sum\nolimits_{(i, j) \in K} n_i m_j [X_i \times_S Y_j]_{n + m - 1}
$$
视为 $X \times_S Y$ 上的一个 $(n + m - 1)$-周。
由此得到加性映射
$\times : Z_n(X) \otimes_{\mathbf{Z}} Z_m(Y) \to Z_{n + m}(X \times_S Y)$。
问题在于证明这一过程与有理等价相容。

\medskip\noindent
设 $i : X' \to X$ 是一个 $\delta$-维数为 $n$、
并支配 $S$ 某个不可约分支的整闭子概形的包含态射。
则 $p' : X' \to S$ 是相对维数为 $n - 1$ 的平坦态射；见
《代数进阶》引理 \ref{more-algebra-lemma-dedekind-torsion-free-flat}。
因此，由引理 \ref{chow-lemma-flat-pullback-bivariant}，沿 $p'$ 的平坦拉回是元素
$(p')^* \in A^{-n + 1}(X' \to S)$；进而由
引理 \ref{chow-lemma-push-proper-bivariant}，有
$c' = i_* \circ (p')^* \in A^{-n + 1}(X \to S)$。
这给出映射
$$
c' \cap - : \CH_m(Y) \longrightarrow \CH_{m + n - 1}(X \times_S Y)
$$
对于任意 $\delta$-维数为 $m$ 的整闭子概形 $Y' \subset Y$，
该映射把 $[Y']$ 映到 $[X' \times_S Y']_{n + m - 1}$。

\medskip\noindent
设 $i : X' \to X$ 是一个 $\delta$-维数为 $n$ 的整闭子概形的包含态射，
并且复合 $X' \to X \to S$ 通过某个闭点 $s \in S$ 分解。
由于 $s$ 是 Dedekind 整环之谱的闭点，$s$ 是 $S$ 上法丛平凡的有效
Cartier 除子。以 $c \in A^1(s \to S)$ 表示 Gysin 同态；见
引理 \ref{chow-lemma-gysin-bivariant}。态射 $p' : X' \to s$
是相对维数为 $n$ 的平坦态射。因此，由
引理 \ref{chow-lemma-flat-pullback-bivariant}，沿 $p'$ 的平坦拉回是元素
$(p')^* \in A^{-n}(X' \to S)$。于是
$$
c' = i_* \circ (p')^* \circ c \in A^{-n}(X \to S)
$$
由引理 \ref{chow-lemma-push-proper-bivariant}，这给出映射
$$
c' \cap - : \CH_m(Y) \longrightarrow \CH_{m + n - 1}(X \times_S Y)
$$
对于任意 $\delta$-维数为 $m$ 的整闭子概形 $Y' \subset Y$，
若 $Y'$ 支配 $S$ 的一个不可约分支，该映射把 $[Y']$ 映到
$[X' \times_S Y']_{n + m - 1}$；否则映到 $0$。

\medskip\noindent
由前两段可知，构造
$([X'], [Y']) \mapsto [X' \times_S Y']_{n + m - 1}$
在第二个变量中可通过有理等价下降，也就是说，它给出良定义映射
$Z_n(X) \otimes_{\mathbf{Z}} \CH_m(Y) \to \CH_{n + m - 1}(X \times_S Y)$。
由对称性，另一个变量也同样如此，结论得证。
\end{proof}

\begin{lemma}
\label{chow-lemma-chow-cohomology-towards-base-dim-1}
设 $(S, \delta)$ 如上，$X$ 是 $S$ 上局部有限型概形。
则对所有 $p \in \mathbf{Z}$，都有规范等同
$$
A^p(X \to S) = \CH_{1 - p}(X)
$$
\end{lemma}

\begin{proof}
考虑元素 $[S]_1 \in \CH_1(S)$。
把 $c$ 映到 $c \cap [S]_1$，即得到映射
$A^p(X \to S) \to \CH_{1 - p}(X)$。

\medskip\noindent
反之，设有 $\alpha \in \CH_{1 - p}(X)$。
可按如下方式定义 $c_\alpha \in A^p(X \to S)$：
给定 $X' \to S$ 和 $\alpha' \in \CH_n(X')$，令
$$
c_\alpha \cap \alpha' = \alpha \times \alpha'
$$
作为 $\CH_{n - p}(X \times_S X')$ 中的元素。
为证明这是双变类，如定义 \ref{chow-definition-cycles} 写成
$\alpha = \sum_{i \in I} n_i[X_i]$。特别地，态射
$$
g : \coprod\nolimits_{i \in I} X_i \longrightarrow X
$$
是固有的。选取 $i \in I$。若 $X_i$ 支配 $S$ 的一个不可约分支，
则结构态射 $p_i : X_i \to S$ 平坦，并且有
$\xi_i = p_i^* \in A^p(X_i \to S)$。
另一方面，若 $p_i$ 分解为 $p'_i : X_i \to s_i$，随后是闭点的包含
$s_i \to S$，则有
$\xi_i = (p'_i)^* \circ c_i \in A^p(X_i \to S)$，
其中 $c_i \in A^1(s_i \to S)$ 是 Gysin 同态，而
$(p'_i)^*$ 是平坦拉回。注意
$$
A^p(\coprod\nolimits_{i \in I} X_i \to S) =
\prod\nolimits_{i \in I} A^p(X_i \to S)
$$
因此有
$$
\xi = \sum n_i \xi_i \in A^p(\coprod\nolimits_{i \in I} X_i \to S)
$$
最后，由于 $g$ 固有，引理 \ref{chow-lemma-push-proper-bivariant} 给出双变类
$$
g_* \circ \xi \in A^p(X \to S)
$$
读者容易验证 $c_\alpha$ 等于该双变类
（请与引理 \ref{chow-lemma-exterior-product-well-defined-dim-1} 的证明比较），
因而它本身也是双变类。

\medskip\noindent
为完成证明，还须说明这两个构造互为逆映射。
由于 $c_\alpha \cap [S]_1 = \alpha$，其中一个方向是显然的。
对另一个方向，令 $c \in A^p(X \to S)$，并置 $\alpha = c \cap [S]_1$。
由引理 \ref{chow-lemma-bivariant-zero}，只需在 $X'$ 是 $S$ 上局部有限型整概形时证明
$$
c \cap [X'] = c_\alpha \cap [X']
$$
但或者 $p' : X' \to S$ 是相对维数为 $\dim_\delta(X') - 1$ 的平坦态射，
因而 $[X'] = (p')^*[S]_1$；或者 $X' \to S$ 分解为
$X' \to s \to S$，因而 $[X'] = (p')^*(s \to S)^*[S]_1$。
所以，双变类 $c$ 和 $c_\alpha$ 在 $[S]_1$ 上相同，
就蕴含它们与 $[X']$ 作帽积时也相同
（因为双变类与平坦拉回和 Gysin 映射可交换），证明完成。
\end{proof}

\begin{lemma}
\label{chow-lemma-chow-cohomology-towards-base-dim-1-commutes}
设 $(S, \delta)$ 如上，$X$ 是 $S$ 上局部有限型概形。
设 $c \in A^p(X \to S)$，$Y \to Z$ 是 $S$ 上局部有限型概形之间的态射，
且 $c' \in A^q(Y \to Z)$。则在
$A^{p + q}(X \times_S Y \to X \times_S Z)$ 中有
$c \circ c' = c' \circ c$。
\end{lemma}

\begin{proof}
在引理 \ref{chow-lemma-chow-cohomology-towards-base-dim-1} 的证明中已经看到，
$c$ 由固有前推、在各连通分支上乘以整数、平坦拉回以及 Gysin 映射组合而成。
依双变类的定义，$c'$ 与这些运算中的每一个都可交换，故结论成立。
这里省略若干细节。
\end{proof}

\begin{remark}
\label{chow-remark-commuting-exterior-dim-1}
引理 \ref{chow-lemma-chow-cohomology-towards-base-dim-1} 和
\ref{chow-lemma-chow-cohomology-towards-base-dim-1-commutes} 的要点如下。
设 $(S, \delta)$ 如上，$X$ 是 $S$ 上局部有限型概形。
设 $\alpha \in \CH_*(X)$，$Y \to Z$ 是 $S$ 上局部有限型概形之间的态射，
且 $c' \in A^q(Y \to Z)$。则对任意 $\beta \in \CH_*(Z)$，
$$
\alpha \times (c' \cap \beta) = c' \cap (\alpha \times \beta)
$$
在 $\CH_*(X \times_S Y)$ 中成立。事实上，只需取与 $\alpha$ 对应的双变类
$c = c_\alpha \in A^*(X \to S)$；见引理 \ref{chow-lemma-chow-cohomology-towards-base-dim-1} 的证明。
\end{remark}

\begin{lemma}
\label{chow-lemma-exterior-product-associative-dim-1}
外积满足结合律。更确切地，设 $(S, \delta)$ 如上，
$X, Y, Z$ 是 $S$ 上局部有限型概形，并设
$\alpha \in \CH_*(X)$、$\beta \in \CH_*(Y)$、$\gamma \in \CH_*(Z)$。
则在 $\CH_*(X \times_S Y \times_S Z)$ 中有
$(\alpha \times \beta) \times \gamma =
\alpha \times (\beta \times \gamma)$。
\end{lemma}

\begin{proof}
从略。提示：概形纤维积的结合律。
\end{proof}















\section{Dedekind 整环上的交乘积}
\label{chow-section-intersection-product-dim-1}

\noindent
设 $S$ 是局部 Noether 概形，并且存在由 Dedekind 整环的谱组成的开覆盖。
对闭点 $s \in S$ 置 $\delta(s) = 0$，否则置 $\delta(s) = 1$。
则 $(S, \delta)$ 是一般情形 \ref{chow-situation-setup} 的特殊情形；
见例 \ref{chow-example-domain-dimension-1} 以及
第 \ref{chow-section-exterior-product-dim-1} 节中的讨论。

\medskip\noindent
设 $X$ 是 $S$ 上光滑概形。第 \ref{chow-section-gysin-for-diagonal} 节的双变类
$\Delta^!$ 使我们能够在 $X$ 上局部有限型概形的 Chow 群上定义一种交乘积。
具体地，设 $Y \to X$ 和 $Z \to X$ 是局部有限型概形之间的态射。注意到
$$
Y \times_X Z =  (Y \times_S Z) \times_{X \times_S X, \Delta} X
$$
因此可以考虑映射序列
$$
\CH_n(Y) \otimes_\mathbf{Z} \CH_m(Y)
\xrightarrow{\times}
\CH_{n + m - 1}(Y \times_S Z)
\xrightarrow{\Delta^!}
\CH_{n + m - *}(Y \times_X Z)
$$
其中第一支箭头是第 \ref{chow-section-exterior-product-dim-1} 节中构造的外积，
第二支箭头是第 \ref{chow-section-gysin-for-diagonal} 节中研究的对角态射 Gysin 映射。
若 $X$ 是维数为 $d$ 的等维概形，则 $X \to S$ 是相对维数为 $d - 1$ 的光滑态射，
故所得元素落在 $\CH_{n + m - d}(Y \times_X Z)$ 中。
一般而言，可依照 $X$ 的开闭子概形（其上 $X$ 具有给定维数），
将其分解为位于这些部分之上的分量。
给定 $\alpha \in \CH_*(Y)$ 和 $\beta \in \CH_*(Z)$，记
$$
\alpha \cdot \beta = \Delta^!(\alpha \times \beta)
\in \CH_*(Y \times_X Z)
$$
在特殊情形 $X = Y = Z$ 下，得到乘法
$$
\CH_*(X) \times \CH_*(X) \to \CH_*(X),\quad
(\alpha, \beta) \mapsto \alpha \cdot \beta
$$
称为{\it 交乘积}。显然，这一乘积是对称的。
结合律由下一个引理给出。

\begin{lemma}
\label{chow-lemma-associative-dim-1}
上面定义的乘积满足结合律。更确切地，设 $(S, \delta)$ 如上，
$X$ 在 $S$ 上光滑，$Y, Z, W$ 是 $X$ 上局部有限型概形，并设
$\alpha \in \CH_*(Y)$、$\beta \in \CH_*(Z)$、$\gamma \in \CH_*(W)$。
则在 $\CH_*(Y \times_X Z \times_X W)$ 中有
$(\alpha \cdot \beta) \cdot \gamma =
\alpha \cdot (\beta \cdot \gamma)$。
\end{lemma}

\begin{proof}
由引理 \ref{chow-lemma-exterior-product-associative-dim-1}，在
$\CH_*(Y \times_S Z \times_S W)$ 中有
$(\alpha \times \beta) \times \gamma =
\alpha \times (\beta \times \gamma)$。考虑闭浸入
$$
\Delta_{12} : X \times_S X \longrightarrow X \times_S X \times_S X,
\quad (x, x') \mapsto (x, x, x')
$$
以及
$$
\Delta_{23} : X \times_S X \longrightarrow X \times_S X \times_S X,
\quad (x, x') \mapsto (x, x', x')
$$
以 $\Delta_{12}^!$ 和 $\Delta_{23}^!$ 表示相应的双变类。
注意，$\Delta_{12}^!$ 是 $\Delta^!$ 沿映射 $\text{pr}_{12}$ 在
$X \times_S X \times_S X$ 上的限制
（注 \ref{chow-remark-restriction-bivariant}），而
$\Delta_{23}^!$ 是 $\Delta^!$ 沿映射 $\text{pr}_{23}$ 在
$X \times_S X \times_S X$ 上的限制。
因此，$\Delta_{12}^!$ 沿 $\Delta_{23}$ 的限制显然是 $\Delta^!$，
而 $\Delta_{23}^!$ 沿 $\Delta_{12}$ 的限制同样是 $\Delta^!$。
于是由引理 \ref{chow-lemma-gysin-commutes}，有
$$
\Delta^! \circ \Delta_{12}^! =
\Delta^! \circ \Delta_{23}^! 
$$
现在用如下等式链证明该引理：
\begin{align*}
(\alpha \cdot \beta) \cdot \gamma
& =
\Delta^!(\Delta^!(\alpha \times \beta) \times \gamma) \\
& =
\Delta^!(\Delta_{12}^!((\alpha \times \beta) \times \gamma)) \\
& =
\Delta^!(\Delta_{23}^!((\alpha \times \beta) \times \gamma)) \\
& =
\Delta^!(\Delta_{23}^!(\alpha \times (\beta \times \gamma)) \\
& =
\Delta^!(\alpha \times \Delta^!(\beta \times \gamma)) \\
& =
\alpha \cdot (\beta \cdot \gamma)
\end{align*}
除第二项和倒数第二项等式或许需要说明外，其余等式均由上文显然成立。
等式
$\Delta_{23}^!(\alpha \times (\beta \times \gamma)) =
\alpha \times \Delta^!(\beta \times \gamma)$
由注 \ref{chow-remark-commuting-exterior} 成立；第二项等式同理。
\end{proof}

\begin{lemma}
\label{chow-lemma-identify-chow-for-smooth-dim-1}
设 $(S, \delta)$ 如上，$X$ 是 $S$ 上维数为 $d$ 的等维光滑概形。映射
$$
A^p(X) \longrightarrow \CH_{d - p}(X),\quad
c \longmapsto c \cap [X]_d
$$
是同构。在这一同构下，双变类的复合变为上面定义的交乘积。
\end{lemma}

\begin{proof}
以 $g : X \to S$ 表示结构态射。该映射是同构的复合
$$
A^p(X) \to A^{p - d + 1}(X \to S) \to \CH_{d - p}(X)
$$
第一支是命题 \ref{chow-proposition-compute-bivariant} 的同构
$c \mapsto c \circ g^*$，第二支是引理 \ref{chow-lemma-chow-cohomology-towards-base-dim-1} 的同构
$c \mapsto c \cap [S]_1$。
由引理 \ref{chow-lemma-chow-cohomology-towards-base-dim-1} 的证明，
第二支箭头的逆把 $\alpha \in \CH_{d - p}(X)$ 映到双变类
$c_\alpha$；对于 $k$ 上局部有限型的 $Y$ 和
$\beta \in \CH_*(Y)$，该双变类把它映到
$\CH_*(X \times_k Y)$ 中的 $\alpha \times \beta$。
再由命题 \ref{chow-proposition-compute-bivariant} 的证明，
第一支箭头的逆把 $c_\alpha$ 映到如下双变类：
对于局部有限型的 $Y \to X$ 和 $\beta \in \CH_*(Y)$，它把后者映到
$\Delta^!(\alpha \times \beta) = \alpha \cdot \beta$。
由此得到引理的最后一项结论。
\end{proof}










\section{Todd 类}
\label{chow-section-todd-classes}

\noindent
与秩为 $r$ 的向量丛 $\mathcal{E}$ 相关联的最后一种类是其
{\it Todd 类} $Todd(\mathcal{E})$。用 Chern 根
$x_1, \ldots, x_r$ 表示时，它定义为
$$
Todd(\mathcal{E})
=
\prod\nolimits_{i = 1}^r
\frac{x_i}{1 - e^{-x_i}}
$$
用 Chern 类 $c_i = c_i(\mathcal{E})$ 表示时，有
$$
Todd(\mathcal{E})
=
1
+
\frac{1}{2}c_1
+
\frac{1}{12}(c_1^2 + c_2)
+
\frac{1}{24}c_1c_2
+
\frac{1}{720}(-c_1^4 + 4c_1^2c_2 + 3c_2^2 + c_1c_3 - c_4)
+
\ldots
$$
上一节已经对分母作了适当说明。对于任意正合列
$$
0
\to
{\mathcal E}_1
\to
{\mathcal E}
\to
{\mathcal E}_2
\to
0
$$
都有
$$
Todd({\mathcal E}) = Todd({\mathcal E}_1) Todd({\mathcal E}_2).
$$







\section{Grothendieck--Riemann--Roch}
\label{chow-section-grr}

\noindent
设 $(S, \delta)$ 如情形 \ref{chow-situation-setup}，
$X, Y$ 在 $S$ 上局部有限型。设 $\mathcal{E}$ 是有限局部自由层；
具体地，${\mathcal E}$ 是 $X$ 上秩为 $r$ 的层。
设 $f : X \to Y$ 是固有光滑态射，并假设
$R^if_*\mathcal{E}$ 都是 $Y$ 上有限秩局部自由层。
在这种情形下，Grothendieck--Riemann--Roch 定理断言
$$
f_*(Todd(T_{X/Y}) ch(\mathcal{E}))
=
\sum (-1)^i ch(R^if_*\mathcal{E})
$$
其中
$$
T_{X/Y} = \SheafHom_{\mathcal{O}_X}(\Omega_{X/Y}, \mathcal{O}_X)
$$
是 $X$ 在 $Y$ 上的相对切丛。若 $Y = \Spec(k)$，其中 $k$ 是域，
则可将其改述为
$$
\chi(X, \mathcal{E}) = \deg(Todd(T_{X/k}) ch(\mathcal{E}))
$$
该定理还有更一般的形式，并且在正确的一般性下表述时更容易证明。
我们将在别处回到这个问题（在此插入未来引用）。







\section{基概形的变换}
\label{chow-section-change-base}

\noindent
本节说明如何比较不同基概形上的理论。

\begin{situation}
\label{chow-situation-setup-base-change}
这里 $(S, \delta)$ 和 $(S', \delta')$ 如情形 \ref{chow-situation-setup}。
此外，$g : S' \to S$ 是概形的平坦态射，$c \in \mathbf{Z}$ 是满足如下条件的整数：
对每个 $s \in S$，以及作为 $g^{-1}(\{s\})$ 某个不可约分支一般点的
$s' \in S'$，都有 $\delta'(s') = \delta(s) + c$。
\end{situation}

\noindent
我们将看到，对于 $S$ 上局部有限型概形 $X$，存在 Chow 群之间的良定义映射
$\CH_k(X) \to \CH_{k + c}(X \times_S S')$；它大体上与本章定义的运算可交换。

\begin{lemma}
\label{chow-lemma-dimension-base-change}
在情形 \ref{chow-situation-setup-base-change} 中，设 $X \to S$ 局部有限型，
并以 $X' \to S'$ 表示沿 $S' \to S$ 的基变换。
若 $X$ 是满足 $\dim_\delta(X) = k$ 的整概形，则
$X'$ 的每个不可约分支 $Z'$ 都满足 $\dim_{\delta'}(Z') = k + c$。
\end{lemma}

\begin{proof}
投影 $X' \to X$ 是平坦态射 $S' \to S$ 的基变换，因而平坦
（《概形态射》引理 \ref{morphisms-lemma-base-change-flat}）。
所以 $X'$ 任一不可约分支的一般点 $x'$ 都映到一般点 $x \in X$
（因为由《概形态射》引理 \ref{morphisms-lemma-generalizations-lift-flat}，
一般化可沿 $X' \to X$ 提升）。令 $s \in S$ 为 $x$ 的像。
回忆概形 $S'_s = S' \times_S s$ 与 $g^{-1}(\{s\})$ 具有相同的底拓扑空间
（《概形》引理 \ref{schemes-lemma-fibre-topological}）。
可将 $x'$ 视为概形 $S'_s \times_s x$ 的一个点；该概形带有单态射
$S'_s \times_s x \to S' \times_S X$。当然，$x'$ 也是
$S'_s \times_s x$ 某个不可约分支的一般点。
利用 $\Spec(\kappa(x)) \to \Spec(\kappa(s)) = s$ 的平坦性并重复上述论证，
可见 $x'$ 映到 $g^{-1}(\{s\})$ 某个不可约分支的一般点 $s'$。
因此由假设，$\delta'(s') = \delta(s) + c$。
由《概形态射》引理 \ref{morphisms-lemma-dimension-fibre-after-base-change}，有
$\dim_x(X_s) = \dim_{x'}(X_{s'})$。
由于 $x$ 是不可约分支 $X_s$ 的一般点
（这是不可约概形，不过我们不需要这一点），而 $x'$ 是 $X'_{s'}$ 某个不可约分支的一般点，
故由《概形态射》引理 \ref{morphisms-lemma-dimension-fibre-at-a-point}，得到
$\text{trdeg}_{\kappa(s)}(\kappa(x)) = 
\text{trdeg}_{\kappa(s')}(\kappa(x'))$。于是
$$
\delta_{X'/S'}(x') = \delta(s') + \text{trdeg}_{\kappa(s')}(\kappa(x')) =
\delta(s) + c + \text{trdeg}_{\kappa(s)}(\kappa(x)) = \delta_{X/S}(x) + c
$$
由定义 \ref{chow-definition-delta-dimension}，这就证明了所需结论。
\end{proof}

\noindent
在情形 \ref{chow-situation-setup-base-change} 中，设 $X \to S$ 局部有限型，
并以 $X' \to S'$ 表示沿 $g : S' \to S$ 的基变换。
存在唯一同态
$$
g^* : Z_k(X) \longrightarrow Z_{k + c}(X')
$$
它把任意 $\delta$-维数为 $k$ 的整闭子概形 $Z \subset X$ 的 $[Z]$
映到 $[Z \times_S S']_{k + c}$。由引理 \ref{chow-lemma-dimension-base-change}，这是有意义的。

\begin{lemma}
\label{chow-lemma-pullback-coherent-base-change}
在情形 \ref{chow-situation-setup-base-change} 中，设 $X \to S$ 局部有限型，
且 $X' \to S$ 是沿 $S' \to S$ 的基变换。
\begin{enumerate}
\item 设 $Z \subset X$ 是满足 $\dim_\delta(Z) \leq k$ 的闭子概形，
其基变换为 $Z' \subset X'$。则 $\dim_{\delta'}(Z')) \leq k + c$，
且在 $Z_{k + c}(X')$ 中有 $[Z']_{k + c} = g^*[Z]_k$。
\item 设 $\mathcal{F}$ 是 $X$ 上满足
$\dim_\delta(\text{Supp}(\mathcal{F})) \leq k$ 的凝聚层，
其在 $X'$ 上的基变换为 $\mathcal{F}'$。
则 $\dim_\delta(\text{Supp}(\mathcal{F}')) \leq k + c$，
且在 $Z_{k + c}(X')$ 中有
$g^*[\mathcal{F}]_k = [\mathcal{F}']_{k + c}$。
\end{enumerate}
\end{lemma}

\begin{proof}
该证明与引理 \ref{chow-lemma-pullback-coherent} 的证明完全相同，
建议读者略过。

\medskip\noindent
关于维数的断言由引理 \ref{chow-lemma-dimension-base-change} 得出。
由引理 \ref{chow-lemma-cycle-closed-coherent}，以及凝聚模
$\mathcal{O}_Z$ 的基变换是 $\mathcal{O}_{Z'}$ 这一事实，
(1) 由 (2) 得出。

\medskip\noindent
证明 (2)。由于 $X$、$X'$ 局部 Noether，可应用
《概形上同调》引理 \ref{coherent-lemma-coherent-Noetherian} 得知
$\mathcal{F}$ 有限型；所以 $\mathcal{F}'$ 有限型
（《模层》引理 \ref{modules-lemma-pullback-finite-type}），
进而 $\mathcal{F}'$ 凝聚
（再次使用《概形上同调》引理 \ref{coherent-lemma-coherent-Noetherian}）。故引理的陈述有意义。
设 $W \subset X$ 是 $\delta$-维数为 $k$ 的整闭子概形，
$W' \subset X'$ 是 $\delta'$-维数为 $k + c$ 的整闭子概形，
并且在 $X' \to X$ 下映入 $W$。须证明 $g^*[\mathcal{F}]_k$ 中
$[W']$ 的系数 $n$ 等于 $[\mathcal{F}']_{k + c}$ 中 $[W']$ 的系数 $m$。
设 $\xi \in W$ 和 $\xi' \in W'$ 为一般点。令
$A = \mathcal{O}_{X, \xi}$、$B = \mathcal{O}_{X', \xi'}$，
并将 $M = \mathcal{F}_\xi$ 视为 $A$-模。
（由维数假设，$M$ 的长度有限，但实际上不必验证这一点；
见引理 \ref{chow-lemma-length-finite}。）有
$\mathcal{F}'_{\xi'} = B \otimes_A M$。因此
$$
n = \text{length}_B(B \otimes_A M)
\quad
\text{且}
\quad
m = \text{length}_A(M) \text{length}_B(B/\mathfrak m_AB)
$$
所需等式遂由《交换代数》引理 \ref{algebra-lemma-pullback-module} 得出。
\end{proof}

\begin{lemma}
\label{chow-lemma-pullback-base-change}
在情形 \ref{chow-situation-setup-base-change} 中，设 $X \to S$ 局部有限型，
且 $X' \to S'$ 是沿 $S' \to S$ 的基变换。
上述映射 $g^* : Z_k(X) \to Z_{k + c}(X')$ 可通过有理等价下降，
从而给出 Chow 群之间的映射
$$
g^* : \CH_k(X) \longrightarrow \CH_{k + c}(X')
$$
\end{lemma}

\begin{proof}
设 $\alpha \in Z_k(X)$ 是有理等价于零的 $k$-周。
由引理 \ref{chow-lemma-rational-equivalence-family}，存在一个整闭子概形的局部有限族
$W_i \subset X \times \mathbf{P}^1$，其 $\delta$-维数为 $k$，
且不包含在 $X \times \mathbf{P}^1$ 的除子
$(X \times \mathbf{P}^1)_0$ 或 $(X \times \mathbf{P}^1)_\infty$ 中，并满足
$\alpha = \sum ([(W_i)_0]_k - [(W_i)_\infty]_k)$。
因此，只需对 $\delta$-维数为 $k$ 的整闭子概形
$W \subset X \times \mathbf{P}^1$ 证明如下两点，
其中该子概形不包含在 $X \times \mathbf{P}^1$ 的除子
$(X \times \mathbf{P}^1)_0$ 或 $(X \times \mathbf{P}^1)_\infty$ 中：
\begin{enumerate}
\item 基变换 $W' \subset X' \times \mathbf{P}^1$ 满足
引理 \ref{chow-lemma-closed-subscheme-cross-p1} 的假设，其中以 $k + c$ 代替 $k$；以及
\item $g^*[W_0]_k = [(W')_0]_{k + c}$ 且
$g^*[W_\infty]_k = [(W')_\infty]_{k + c}$。
\end{enumerate}
(2) 立即由引理 \ref{chow-lemma-pullback-coherent-base-change} 以及
$(W')_0$ 是 $W_0$ 的基变换这一事实得出（利用纤维积的结合律）。
对于 (1)，维数断言先由引理 \ref{chow-lemma-dimension-base-change} 得出。
再令 $w' \in (W')_0$，其像为 $w \in W_0$，并令
$z \in \mathbf{P}^1_S$。以 $t \in \mathcal{O}_{\mathbf{P}^1_S, z}$
表示 $0 : S \to \mathbf{P}^1_S$ 的通常方程。
由于 $\mathcal{O}_{W, w} \to \mathcal{O}_{W', w'}$ 平坦，
且 $t$ 是 $\mathcal{O}_{W, w}$ 上的非零因子
（因为 $W$ 整且 $W \not = W_0$），故 $t$ 也是
$\mathcal{O}_{W', w'}$ 中的非零因子。因此 $W'$ 没有落在 $W'_0$ 上的伴随点。
\end{proof}

\begin{lemma}
\label{chow-lemma-pullback-base-change-pullback}
在情形 \ref{chow-situation-setup-base-change} 中，设
$Y \to X \to S$ 局部有限型，并以 $Y' \to X' \to S'$ 表示沿
$S' \to S$ 的基变换。假设 $f : Y \to X$ 是相对维数为 $r$ 的平坦态射。
则 $f' : Y' \to X'$ 是相对维数为 $r$ 的平坦态射，并且周群与 Chow 群的图
$$
\vcenter{
\xymatrix{
Z_{k + r}(Y) \ar[r]_{g^*} & Z_{k + c + r}(Y') \\
Z_k(X) \ar[r]^{g^*} \ar[u]^{f^*} & Z_{k + c}(X') \ar[u]_{(f')^*}
}
}
\quad\text{且}\quad
\vcenter{
\xymatrix{
\CH_{k + r}(Y) \ar[r]_{g^*} & \CH_{k + c + r}(Y') \\
\CH_k(X) \ar[r]^{g^*} \ar[u]^{f^*} & \CH_{k + c}(X') \ar[u]_{(f')^*}
}
}
$$
均交换。
\end{lemma}

\begin{proof}
只需证明第一个图交换。为此，设 $Z \subset X$ 是
$\delta$-维数为 $k$ 的整闭子概形，并以 $Z' \subset X'$ 表示其基变换。
由构造，有 $g^*[Z] = [Z']_{k + c}$。由引理 \ref{chow-lemma-pullback-coherent}，有
$(f')^*g^*[Z] = [Z' \times_{X'} Y']_{k + c + r}$。
反过来，由定义 \ref{chow-definition-flat-pullback}，有
$f^*[Z] = [Z \times_X Y]_{k + r}$；由
引理 \ref{chow-lemma-pullback-coherent-base-change}，有
$g^*f^*[Z] = [(Z \times_X Y)']_{k + r + c}$。
由于纤维积的结合律给出
$(Z \times_X Y)' = Z' \times_{X'} Y'$，结论成立。
\end{proof}

\begin{lemma}
\label{chow-lemma-pullback-base-change-pushforward}
在情形 \ref{chow-situation-setup-base-change} 中，设
$Y \to X \to S$ 局部有限型，并以 $Y' \to X' \to S'$ 表示沿
$S' \to S$ 的基变换。假设 $f : Y \to X$ 固有。
则 $f' : Y' \to X'$ 固有，并且周群与 Chow 群的图
$$
\vcenter{
\xymatrix{
Z_k(Y) \ar[r]_{g^*} \ar[d]_{f_*} & Z_{k + c}(Y') \ar[d]^{f'_*} \\
Z_k(X) \ar[r]^{g^*} & Z_{k + c}(X')
}
}
\quad\text{且}\quad
\vcenter{
\xymatrix{
\CH_k(Y) \ar[r]_{g^*} \ar[d]_{f_*} & \CH_{k + c}(Y') \ar[d]^{f'_*} \\
\CH_k(X) \ar[r]^{g^*} & \CH_{k + c}(X')
}
}
$$
均交换。
\end{lemma}

\begin{proof}
只需证明第一个图交换。为此，设 $Z \subset Y$ 是
$\delta$-维数为 $k$ 的整闭子概形，并以 $Z' \subset X'$ 表示其基变换。
由构造，有 $g^*[Z] = [Z']_{k + c}$。由引理 \ref{chow-lemma-cycle-push-sheaf}，有
$(f')_*g^*[Z] = [f'_*\mathcal{O}_{Z'}]_{k + c}$。
由同一引理，有 $f_*[Z] = [f_*\mathcal{O}_Z]_k$；由
引理 \ref{chow-lemma-pullback-coherent-base-change}，有
$g^*f_*[Z] = [(X' \to X)^*f_*\mathcal{O}_Z]_{k + r}$。
因此只需证明 $X'$ 上凝聚模之间的同构
$$
(X' \to X)^*f_*\mathcal{O}_Z \cong f'_*\mathcal{O}_{Z'}
$$
由于 $X' \to X$ 平坦且
$\mathcal{O}_{Z'} = (Y' \to Y)^*\mathcal{O}_Z$，
该同构由平坦基变换得出；见《概形上同调》引理 \ref{coherent-lemma-flat-base-change-cohomology}。
\end{proof}

\begin{lemma}
\label{chow-lemma-pullback-base-change-c1}
在情形 \ref{chow-situation-setup-base-change} 中，设 $X \to S$ 局部有限型，
并以 $X' \to S'$ 表示沿 $S' \to S$ 的基变换。
设 $\mathcal{L}$ 是可逆 $\mathcal{O}_X$-模，其在 $X'$ 上的基变换为
$\mathcal{L}'$。则 Chow 群的图
$$
\xymatrix{
\CH_k(X) \ar[r]_{g^*} \ar[d]_{c_1(\mathcal{L}) \cap -} &
\CH_{k + c}(X') \ar[d]^{c_1(\mathcal{L}') \cap -} \\
\CH_{k - 1}(X) \ar[r]^{g^*} & \CH_{k + c - 1}(X')
}
$$
交换。
\end{lemma}

\begin{proof}
设 $p : L \to X$ 是与 $\mathcal{L}$ 相关联的线丛，
其零截面为 $o : X \to L$。对 $\alpha \in CH_k(X)$，
我们知道 $\beta = c_1(\mathcal{L}) \cap \alpha$ 是满足
$o_*\alpha = - p^*\beta$ 的唯一元素 $\CH_{k - 1}(X)$；见
引理 \ref{chow-lemma-linebundle} 和 \ref{chow-lemma-linebundle-formulae}。
拉回后仍有相同的刻画。因此引理由引理 \ref{chow-lemma-pullback-base-change-pullback} 和
\ref{chow-lemma-pullback-base-change-pushforward} 得出。
\end{proof}

\begin{lemma}
\label{chow-lemma-pullback-base-change-chern-classes}
在情形 \ref{chow-situation-setup-base-change} 中，设 $X \to S$ 局部有限型，
并以 $X' \to S'$ 表示沿 $S' \to S$ 的基变换。
设 $\mathcal{E}$ 是秩为 $r$ 的有限局部自由 $\mathcal{O}_X$-模，
其在 $X'$ 上的基变换为 $\mathcal{E}'$。则对所有 $i$，Chow 群的图
$$
\xymatrix{
\CH_k(X) \ar[r]_{g^*} \ar[d]_{c_i(\mathcal{E}) \cap -} &
\CH_{k + c}(X') \ar[d]^{c_i(\mathcal{E}') \cap -} \\
\CH_{k - i}(X) \ar[r]^{g^*} & \CH_{k + c - i}(X')
}
$$
交换。
\end{lemma}

\begin{proof}
置 $P = \mathbf{P}(\mathcal{E})$。$P$ 的基变换 $P'$ 等于
$\mathbf{P}(\mathcal{E}')$。我们已经知道，平坦拉回以及与可逆模的
$c_1$ 作杯积都与基变换可交换（引理 \ref{chow-lemma-pullback-base-change-pullback} 和
\ref{chow-lemma-pullback-base-change-c1}）。因此，该引理由
引理 \ref{chow-lemma-determine-intersections} 中对与
$c_i(\mathcal{E})$ 作帽积的刻画得出。
\end{proof}

\begin{lemma}
\label{chow-lemma-compose-base-change}
设 $(S, \delta)$、$(S', \delta')$、$(S'', \delta'')$ 如
情形 \ref{chow-situation-setup}。设 $g : S' \to S$ 和 $g' : S'' \to S'$
是概形的平坦态射，并设 $c, c' \in \mathbf{Z}$ 是整数，使得
$S, \delta, S', \delta', g, c$ 以及 $S', \delta', S'', g', c'$ 均如
情形 \ref{chow-situation-setup-base-change}。
设 $X \to S$ 局部有限型，并以 $X' \to S'$ 和 $X'' \to S''$
表示分别沿 $S' \to S$ 和 $S'' \to S$ 的基变换。则
\begin{enumerate}
\item $S, \delta, S'', \delta'', g \circ g', c + c'$ 如
情形 \ref{chow-situation-setup-base-change}；
\item 映射 $g^* : Z_k(X) \to Z_{k + c}(X')$ 与
$(g')^* : Z_{k + c}(X') \to Z_{k + c + c'}(X'')$ 的复合给出映射
$(g \circ g')^* : Z_k(X) \to Z_{k + c + c'}(X'')$；以及
\item 引理 \ref{chow-lemma-pullback-base-change} 的映射
$g^* : \CH_k(X) \to \CH_{k + c}(X')$ 与
$(g')^* : \CH_{k + c}(X') \to \CH_{k + c + c'}(X'')$ 的复合，
给出引理 \ref{chow-lemma-pullback-base-change} 的映射
$(g \circ g')^* : \CH_k(X) \to \CH_{k + c + c'}(X'')$。
\end{enumerate}
\end{lemma}

\begin{proof}
设 $s \in S$，并设 $s'' \in S''$ 是
$(g \circ g')^{-1}(\{s\})$ 某个不可约分支的一般点。置 $s' = g'(s'')$。
显然，$s''$ 是 $(g')^{-1}(\{s'\})$ 某个不可约分支的一般点。
此外，由于 $g'$ 平坦，一般化可沿 $g'$ 提升
（《概形态射》引理 \ref{morphisms-lemma-base-change-flat}），
所以 $s'$ 也是 $g^{-1}(\{s\})$ 某个不可约分支的一般点。
因此由假设，$\delta'(s') = \delta(s) + c$ 且
$\delta''(s'') = \delta'(s') + c'$。于是
$\delta''(s'') = \delta(s) + c + c'$，第一项断言成立。

\medskip\noindent
对于第二项，设 $Z \subset X$ 是 $\delta$-维数为 $k$ 的整闭子概形，
并以 $Z' \subset X'$ 和 $Z'' \subset X''$ 表示其基变换。
由定义，$g^*[Z] = [Z']_{k + c}$。由引理 \ref{chow-lemma-pullback-coherent-base-change}，有
$(g')^*[Z']_{k + c} = [Z'']_{k + c + c'}$。这证明了最后一项断言。
\end{proof}

\begin{lemma}
\label{chow-lemma-chow-limit}
在情形 \ref{chow-situation-setup-base-change} 中，假设 $c = 0$，
并假设 $S' = \lim_{i \in I} S_i$ 是 $S$ 上仿射概形 $S_i$ 的滤过极限，且
\begin{enumerate}
\item 当 $\delta_i$ 等于 $S_i \to S \xrightarrow{\delta} \mathbf{Z}$ 时，
序对 $(S_i, \delta_i)$ 如情形 \ref{chow-situation-setup}；
\item $S_i, \delta_i, S, \delta, S \to S_i, c = 0$ 如
情形 \ref{chow-situation-setup-base-change}；
\item 对 $i \geq i'$，$S_i, \delta_i, S_{i'}, \delta_{i'}, S_i \to S_{i'}, c = 0$
如情形 \ref{chow-situation-setup-base-change}。
\end{enumerate}
若 $X$ 是 $S$ 上拟紧有限型概形，并以 $X'$ 和 $X_i$ 表示分别沿
$S' \to S$ 和 $S_i \to S$ 的基变换，则
$Z_k(X') = \colim Z_k(X_i)$ 且 $\CH_k(X') = \colim \CH_k(X_i)$。
\end{lemma}

\begin{proof}
由引理 \ref{chow-lemma-compose-base-change}，得到周群系 $Z_k(X_i)$、
Chow 群系 $\CH_k(X_i)$，以及映射
$\colim Z_k(X_i) \to Z_k(X')$ 和 $\colim \CH_i(X_i) \to \CH_k(X')$。
可以把 $S$ 替换为 $X \to S$ 所通过的某个拟紧开子概形；
因此可以并且确实假设本证明中出现的所有概形都是 Noether 的
（从而拟紧且拟分离）。

\medskip\noindent
先证明映射 $\colim Z_k(X_i) \to Z_k(X')$ 满射。
设 $Z' \subset X'$ 是 $\delta'$-维数为 $k$ 的整闭子概形。
由《概形的极限》引理 \ref{limits-lemma-descend-finite-presentation}，
可找到某个 $i$ 以及有限表现态射 $Z_i \to X_i$，其基变换为 $Z'$。
增大 $i$ 后，可以假设 $Z_i$ 是 $X_i$ 的闭子概形；见
《概形的极限》引理 \ref{limits-lemma-descend-closed-immersion-finite-presentation}。
于是 $Z' \to X_i$ 通过 $Z_i$ 分解，并且可以用 $Z' \to X_i$ 的概形论像替换
$Z_i$。由此可见，可以假设 $Z_i$ 是 $X_i$ 的整闭子概形。
由引理 \ref{chow-lemma-dimension-base-change}，得到
$\dim_{\delta_i}(Z_i) = \dim_{\delta'}(Z') = k$。
因此，$Z_k(X_i) \to Z_k(X')$ 把 $[Z_i]$ 映到 $[Z']$，
从而满射性成立。

\medskip\noindent
下面证明映射 $\colim Z_k(X_i) \to Z_k(X')$ 单射。
设 $\alpha_i = \sum n_j[Z_j] \in Z_k(X_i)$ 是一个在
$Z_k(X')$ 中的像为零的周。可以并且确实假设：若 $j \not = j'$，
则 $Z_j \not = Z_{j'}$；并且对所有 $j$ 都有 $n_j \not = 0$。
以 $Z'_j \subset X'$ 表示 $Z_j$ 的基变换。
由引理 \ref{chow-lemma-dimension-base-change}，$Z'_j$ 的每个不可约分支都具有
$\delta'$-维数 $k$。此外，由于 $Z_j$ 不可约且 $Z'_j \to Z_j$ 平坦
（它是 $S' \to S$ 的基变换），$Z'_j \to Z_j$ 支配。
所以，若 $Z'_j$ 非空，则 $Z'_j$ 的某个不可约分支（记为 $Z'$）支配 $Z_j$。
于是当 $j' \not = j$ 时，$Z'$ 不可能是 $Z'_{j'}$ 的不可约分支。
因此，若 $Z'_j$ 非空，则
$(S' \to S_i)^*\alpha_i
= \sum [Z'_j]_r$ 非零（因为 $Z'$ 的系数将非零）。
故对所有 $j$ 都有 $Z'_j = \emptyset$。然而，由《概形的极限》引理 \ref{limits-lemma-limit-nonempty}，这意味着存在某个过渡态射
$S_{i'} \to S_i$，使 $Z_j$ 的基变换为空。
因此 $\alpha_i$ 在余极限中消失，正合所需。

\medskip\noindent
映射 $\colim Z_k(X_i) \to Z_k(X')$ 的满射性蕴含
$\colim \CH_k(X_i) \to \CH_k(X')$ 满射。
为完成证明，须说明后一个映射单射。
设 $\alpha_i \in \CH_k(X_i)$ 是一个周，其像
$\alpha' \in \CH_k(X')$ 为零。则存在 $\delta"$-维数为 $k + 1$ 的整闭子概形
$W_l' \subset X'$、$l = 1, \ldots, r$，以及 $W'_l$ 上的非零有理函数
$f'_l$，使得
$\alpha' = \sum_{l = 1, \ldots, r} \text{div}_{W'_l}(f'_l)$。
重复上述论证，可找到某个 $i$ 以及 $\delta_i$-维数为 $k + 1$ 的整闭子概形
$W_{i, l} \subset X_i$，其基变换为 $W'_l$。
增大 $i$ 后，可以假设在 $W_{i, l}$ 上有有理函数 $f_{i, l}$。
具体地，可将 $f'_l$ 视为非空开子概形 $U'_l \subset W'_l$ 上的结构层截面；
由《概形的极限》引理 \ref{limits-lemma-descend-opens}，这些开子概形可以下降；
再增大 $i$ 后，由《概形的极限》引理 \ref{limits-lemma-descend-section}，
$f'_l$ 也可以下降。我们断言，可能再次增大 $i$ 后，有
$$
\alpha_i = \sum\nolimits_{l = 1, \ldots, r} \text{div}_{W_{i, l}}(f_{i, l})
$$

\medskip\noindent
为证明该断言，取 $\delta'$-维数为 $k$ 的整闭子概形的有限族
$Z'_{l, j} \subset W'_l$，使 $f'_l$ 在
$\bigcup_j Y'_{l, j}$ 之外是可逆正则函数。
增大 $i$ 后（利用上述论证），可以假设存在 $\delta_i$-维数为 $k$ 的整闭子概形
$Z_{i, l, j} \subset W_i$，使 $f_{i, l}$ 在
$\bigcup_j Z_{i, l, j}$ 之外是可逆正则函数。于是可写成
$$
\text{div}_{W'_l}(f'_l) = \sum n_{l, j} [Z'_{l, j}]
$$
以及
$$
\text{div}_{W_{i, l}}(f_{i, l}) = \sum n_{i, l, j} [Z_{i, l, j}]
$$
为证明该断言，只需说明 $n_{l, i} = n_{i, l, j}$。
事实上，这将蕴含
$\beta_i =
\alpha_i - \sum\nolimits_{l = 1, \ldots, r} \text{div}_{W_{i, l}}(f_{i, l})$
是 $X_i$ 上的周，并且其拉回在 $X'$ 上作为周为零！
由一个从略的简单论证可知，存在某个 $i' \geq i$，使得
$\beta_i$ 在 $X_{i'}$ 上的拉回作为周为零。

\medskip\noindent
为证明等式 $n_{l, i} = n_{i, l, j}$，选取一般点
$\xi' \in Z'_{l, j}$，并以 $\xi \in Z_{i, l, j}$ 表示其像；后者也是一般点。
此时局部环映射
$$
\mathcal{O}_{W_{i, l}, \xi}
\longrightarrow
\mathcal{O}_{W'_l, \xi'}
$$
平坦，因为 $W'_l \to W_{i, l}$ 是平坦态射 $S' \to S_i$ 的基变换。
还因为 $Z_{i, l, j}$ 拉回为 $Z'_{l, j}$，有
$\mathfrak m_\xi \mathcal{O}_{W'_l, \xi'} = \mathfrak m_{\xi'}$！
因此，等式
$$
n_{l, j} = \text{ord}_{Z'_{l, j}}(f'_l) =
\text{ord}_{\mathcal{O}_{W'_l, \xi'}}(f'_l)
\quad\text{且}\quad
n_{i, l, j} = \text{ord}_{Z_{i, l, j}}(f_{i, l}) =
\text{ord}_{\mathcal{O}_{W_{i, l}, \xi}}(f_{i, l})
$$
由《交换代数》引理 \ref{algebra-lemma-pullback-module} 以及
《交换代数》第 \ref{algebra-section-orders-of-vanishing} 节中
$\text{ord}$ 的构造得出。
\end{proof}












\section{附录 A：关键引理的另一种处理方法}
\label{chow-section-appendix-A}

\noindent
在本附录中，我们首先定义局部环 $(R, \mathfrak m, \kappa)$ 上有限长度模
$M$ 的行列式 $\det_\kappa(M)$，见小节
\ref{chow-subsection-determinants-finite-length}。
行列式 $\det_\kappa(M)$ 是 $1$ 维 $\kappa$-向量空间。
在小节 \ref{chow-subsection-periodic-complexes-determinants} 中，我们用它定义
有限长度模 $M$ 上正合 $(2, 1)$-周期复形 $(M, \varphi, \psi)$ 的行列式
$\det_\kappa(M, \varphi, \psi) \in \kappa^*$。
在小节 \ref{chow-subsection-symbols} 中，当 $R$ 是 $1$ 维诺特环时，我们利用这些
行列式为一对非零因子 $a, b \in R$ 构造驯符号
$d_R(a, b) = \det_\kappa(R/ab, a, b)$。
尽管无疑有
$$
d_R(a, b) = \partial_R(a, b)
$$
其中 $\partial_R$ 如第 \ref{chow-section-tame-symbol} 节所定义，但我们（尚）未加入
验证。本附录所构造的驯符号的优点是，它可以推广到有限 $R$-模 $M$ 的一对
单射自同态 $\varphi, \psi$（例如），其中该模的维数为 $1$，且
$\varphi(\psi(M)) = \psi(\varphi(M))$。在小节
\ref{chow-subsection-length-determinant} 中，我们把 Herbrand 商与行列式联系起来。
主结果（命题 \ref{chow-proposition-length-determinant-periodic-complex}）的一种
易于陈述的形式是公式
$$
-e_R(M, \varphi, \psi) =
\text{ord}_R(\det\nolimits_K(M_K, \varphi, \psi))
$$
这里 $(M, \varphi, \psi)$ 是定义在 $1$ 维诺特局部整环 $R$ 上的
$(2, 1)$-周期复形，其分式域为 $K$，且该复形的 Herbrand 商 $e_R$
（定义 \ref{chow-definition-periodic-length}）有定义。
我们用这一命题，对本附录所构造的驯符号给出关键引理
（引理 \ref{chow-lemma-milnor-gersten-low-degree}）的另一种证明；见引理
\ref{chow-lemma-secondary-ramification}。


\subsection{有限长度模的行列式}
\label{chow-subsection-determinants-finite-length}

\noindent
本节内容与论文 \cite{determinant} 及学位论文 \cite{Joe} 中的内容有关。

\medskip\noindent
设 $(R, \mathfrak m, \kappa)$ 为局部环，设
$\varphi : M \to M$ 是有限长度 $R$-模 $M$ 的 $R$-线性自同态。
在《More on Algebra》第
\ref{more-algebra-section-determinants-finite-length} 节中，我们已经定义了
$\varphi$ 关于 $\kappa$ 的行列式 $\det_\kappa(\varphi)$
（以及迹和特征多项式）。本节将构造一个典范的 $1$ 维 $\kappa$-向量空间
$\det_\kappa(M)$，使得
$\det_\kappa(\varphi : M \to M) : \det_\kappa(M) \to \det_\kappa(M)$
等于乘以 $\det_\kappa(\varphi)$ 的映射。
若 $M$ 被 $\mathfrak m$ 零化，则可把 $M$ 视为有限维 $\kappa$-向量空间，
此时有 $\det_\kappa(M) = \wedge^n_\kappa(M)$，其中
$n = \dim_\kappa(M)$。我们的构造将把这一点推广到 $R$ 上所有有限长度模；
若 $R$ 包含其剩余域，则在适当意义下，行列式 $\det_\kappa(M)$ 由通常的
行列式给出，见注 \ref{chow-remark-explain-determinant}。

\begin{definition}
\label{chow-definition-determinant}
设 $R$ 为局部环，其极大理想为 $\mathfrak m$，剩余域为 $\kappa$。
设 $M$ 为有限长度 $R$-模，并记 $l = \text{length}_R(M)$。
\begin{enumerate}
\item 给定元素 $x_1, \ldots, x_r \in M$，我们用
$\langle x_1, \ldots, x_r \rangle = Rx_1 + \ldots + Rx_r$ 表示由
$x_1, \ldots, x_r$ 生成的 $M$ 的 $R$-子模。
\item 若 $M$ 中元素的一个 $l$-元组 $(e_1, \ldots, e_l)$ 对每个
$i = 1, \ldots, l$ 都满足
$\mathfrak m e_i \subset \langle e_1, \ldots, e_{i - 1} \rangle$，
则称它是{\it 容许的}。
\item 一个{\it 符号} $[e_1, \ldots, e_l]$ 意指
$(e_1, \ldots, e_l)$ 是容许 $l$-元组。
\item 符号之间的{\it 容许关系}是下列关系之一：
\begin{enumerate}
\item 若 $(e_1, \ldots, e_l)$ 是容许序列，且对某个
$1 \leq a \leq l$ 有
$e_a \in \langle e_1, \ldots, e_{a - 1}\rangle$，则
$[e_1, \ldots, e_l] = 0$；
\item 若 $(e_1, \ldots, e_l)$ 是容许序列，且对某个
$1 \leq a \leq l$ 有 $e_a = \lambda e'_a + x$，其中
$\lambda \in R^*$，并且
$x \in \langle e_1, \ldots, e_{a - 1}\rangle$，则
$$
[e_1, \ldots, e_l] =
\overline{\lambda} [e_1, \ldots, e_{a - 1}, e'_a, e_{a + 1}, \ldots, e_l]
$$
其中 $\overline{\lambda} \in \kappa^*$ 是 $\lambda$ 在剩余域中的像；
\item 若 $(e_1, \ldots, e_l)$ 是容许序列，且
$\mathfrak m e_a \subset \langle e_1, \ldots, e_{a - 2}\rangle$，则
$$
[e_1, \ldots, e_l] =
- [e_1, \ldots, e_{a - 2}, e_a, e_{a - 1}, e_{a + 1}, \ldots, e_l].
$$
\end{enumerate}
\item
我们把有限长度 $R$-模 $M$ 的{\it 行列式}定义为
$$
\det\nolimits_\kappa(M) =
\left\{
\frac{\text{由符号生成的 }\kappa\text{-向量空间}}
{\text{可容许关系的 }\kappa\text{-线性组合}}
\right\}
$$
\end{enumerate}
\end{definition}

\noindent
我们强调，始终有 $l = \text{length}_R(M)$。还要强调的是，这并不推出
符号 $[e_1, \ldots, e_l]$ 对各个分量具有可加性（通常并非如此）。
在证明行列式 $\det_\kappa(M)$ 的维数确实为 $1$ 之前，我们必须先证明
它的维数至多为 $1$。

\begin{lemma}
\label{chow-lemma-dimension-at-most-one}
沿用上述记号，有 $\dim_\kappa(\det_\kappa(M)) \leq 1$。
\end{lemma}

\begin{proof}
固定 $M$ 的一个容许序列 $(f_1, \ldots, f_l)$，使得
$$
\text{length}_R(\langle f_1, \ldots, f_i\rangle) = i
$$
对 $i = 1, \ldots, l$ 成立。这样的容许序列存在，恰是因为 $M$ 的长度为
$l$。我们将证明，$\det_\kappa(M)$ 的任意元素都是符号
$[f_1, \ldots, f_l]$ 的一个 $\kappa$-倍数。这就证明了引理。

\medskip\noindent
设 $(e_1, \ldots, e_l)$ 是 $M$ 的容许序列。只需证明
$[e_1, \ldots, e_l]$ 是 $[f_1, \ldots, f_l]$ 的倍数。
先假设 $\langle e_1, \ldots, e_l\rangle \not = M$。于是存在
$i \in [1, \ldots, l]$，使得
$e_i \in \langle e_1, \ldots, e_{i - 1}\rangle$。由第一条容许关系立即得到
$[e_1, \ldots, e_n] = 0$ 在 $\det_\kappa(M)$ 中成立。
因此可假设 $\langle e_1, \ldots, e_l\rangle = M$。特别地，存在最小指标
$i \in \{1, \ldots, l\}$，使得
$f_1 \in \langle e_1, \ldots, e_i\rangle$。这意味着
$e_i = \lambda f_1 + x$，其中
$x \in \langle e_1, \ldots, e_{i - 1}\rangle$ 且 $\lambda \in R^*$。
由第二条容许关系，这意味着
$[e_1, \ldots, e_l] =
\overline{\lambda}[e_1, \ldots, e_{i - 1}, f_1, e_{i + 1}, \ldots, e_l]$。
注意 $\mathfrak m f_1 = 0$。因此，应用第三条容许关系 $i - 1$ 次可得
$$
[e_1, \ldots, e_l] =
(-1)^{i - 1}\overline{\lambda}
[f_1, e_1, \ldots, e_{i - 1}, e_{i + 1}, \ldots, e_l].
$$
还要注意，
$ \langle f_1, e_1, \ldots, e_{i - 1}, e_{i + 1}, \ldots, e_l\rangle = M$
也成立。归纳假设已经证明，原符号等于某个标量乘以
$$
[f_1, \ldots, f_j, e_{j + 1}, \ldots, e_l]
$$
其中 $(f_1, \ldots, f_j, e_{j + 1}, \ldots, e_l)$ 是元素生成 $M$ 的
某个容许序列，即，\ 满足
$\langle f_1, \ldots, f_j, e_{j + 1}, \ldots, e_l\rangle = M$。
然后取最小的 $i$，使得
$f_{j + 1} \in \langle f_1, \ldots, f_j, e_{j + 1}, \ldots, e_i\rangle$，
并重复上述过程，得到
$$
[f_1, \ldots, f_j, e_{j + 1}, \ldots, e_l]
=
(\text{标量}) [f_1, \ldots, f_j, f_{j + 1}, e_{j + 1},
\ldots, \hat{e_i}, \ldots, e_l]
$$
如此继续，便得到所需结论。
\end{proof}

\noindent
在证明 $\det_\kappa(M)$ 的维数总是 $1$ 之前，先证明当该模是
$\kappa$ 上的向量空间时，它与通常的最高外幂一致。

\begin{lemma}
\label{chow-lemma-compare-det}
设 $R$ 为局部环，其极大理想为 $\mathfrak m$，剩余域为 $\kappa$。
设 $M$ 是被 $\mathfrak m$ 零化的有限长度 $R$-模，并设
$l = \dim_\kappa(M)$。则映射
$$
\det\nolimits_\kappa(M) \longrightarrow \wedge^l_\kappa(M),
\quad
[e_1, \ldots, e_l] \longmapsto e_1 \wedge \ldots \wedge e_l
$$
是同构。
\end{lemma}

\begin{proof}
引理中所述规则显然给出一个 $\kappa$-线性映射，因为通常的符号
$e_1 \wedge \ldots \wedge e_l$ 满足全部容许关系。该映射也显然满射。
另一方面，由引理 \ref{chow-lemma-dimension-at-most-one}，左端的维数至多为一，
故此映射是同构。
\end{proof}

\begin{lemma}
\label{chow-lemma-determinant-dimension-one}
设 $R$ 为局部环，其极大理想为 $\mathfrak m$，剩余域为 $\kappa$。
设 $M$ 为有限长度 $R$-模。上面定义的行列式 $\det_\kappa(M)$ 是维数为
$1$ 的 $\kappa$-向量空间。对任意满足
$\langle f_1, \ldots f_l \rangle = M$ 的容许序列，符号
$[f_1, \ldots, f_l]$ 都生成该空间。
\end{lemma}

\begin{proof}
由引理 \ref{chow-lemma-dimension-at-most-one} 及其证明可知，
$\det_\kappa(M)$ 的维数至多为 $1$，而且它实际上由
$[f_1, \ldots, f_l]$ 生成。我们将对 $l = \text{length}(M)$ 归纳，
证明它非零。当 $l = 1$ 时，结论由引理 \ref{chow-lemma-compare-det} 得出。
选取非零元素 $f \in M$，使得 $\mathfrak m f = 0$。令
$\overline{M} = M /\langle f \rangle$，并把商映射记为
$x \mapsto \overline{x}$。我们将定义一个满射
$$
\psi : \det\nolimits_k(M) \to \det\nolimits_\kappa(\overline{M})
$$
由归纳假设，$\overline{M}$ 的行列式非零，因而这将证明引理。

\medskip\noindent
我们如下在符号上定义 $\psi$。设 $(e_1, \ldots, e_l)$ 是容许序列。
若 $f \not \in \langle e_1, \ldots, e_l \rangle$，就令
$\psi([e_1, \ldots, e_l]) = 0$。若
$f \in \langle e_1, \ldots, e_l \rangle$，则选取使
$f \in \langle e_1, \ldots, e_i \rangle$ 成立的最小 $i$。
可写成 $e_i = \lambda f + x$，其中 $\lambda \in R$ 是单位，且
$x \in \langle e_1, \ldots, e_{i - 1} \rangle$。在此情形令
$$
\psi([e_1, \ldots, e_l]) =
(-1)^i
\overline{\lambda}[\overline{e}_1, \ldots,
\overline{e}_{i - 1},
\overline{e}_{i + 1}, \ldots, \overline{e}_l].
$$
确实，$(\overline{e}_1, \ldots,
\overline{e}_{i - 1},
\overline{e}_{i + 1}, \ldots, \overline{e}_l)$ 是
$\overline{M}$ 中的容许序列，所以该定义有意义。下面证明，将此规则按
$\kappa$-线性方式延拓到符号的线性组合，确实会在行列式上给出映射。
为此必须证明容许关系都被映为零。

\medskip\noindent
类型 (a) 的关系。设 $(e_1, \ldots, e_l)$ 是容许序列，且对某个
$1 \leq a \leq l$ 有
$e_a \in \langle e_1, \ldots, e_{a - 1}\rangle$。
设 $i$ 是使 $f \in \langle e_1, \ldots, e_i\rangle$ 成立的最小指标。
则 $i \not = a$，而且当 $i < a$ 时，
$\overline{e}_a \in \langle \overline{e}_1, \ldots,
\hat{\overline{e}_i}, \ldots, \overline{e}_{a - 1}\rangle$；
当 $i > a$ 时，
$\overline{e}_a \in \langle \overline{e}_1, \ldots,
\overline{e}_{a - 1}\rangle$。因此，$\det_\kappa(\overline{M})$ 中的
相应容许关系迫使符号 $[\overline{e}_1, \ldots,
\overline{e}_{i - 1},
\overline{e}_{i + 1}, \ldots, \overline{e}_l]$ 为零，正如所需。

\medskip\noindent
类型 (b) 的关系。设 $(e_1, \ldots, e_l)$ 是容许序列，且对某个
$1 \leq a \leq l$ 有 $e_a = \lambda e'_a + x$，其中
$\lambda \in R^*$，并且
$x \in \langle e_1, \ldots, e_{a - 1}\rangle$。
设 $i$ 是使 $f \in \langle e_1, \ldots, e_i\rangle$ 成立的最小指标，
并写成 $e_i = \mu f + y$，其中
$y \in \langle e_1, \ldots, e_{i - 1}\rangle$。若 $i < a$，则所需等式为
$$
(-1)^i
\overline{\lambda}
[\overline{e}_1,
\ldots,
\overline{e}_{i - 1},
\overline{e}_{i + 1},
\ldots,
\overline{e}_l]
=
(-1)^i
\overline{\lambda}
[\overline{e}_1,
\ldots,
\overline{e}_{i - 1},
\overline{e}_{i + 1},
\ldots,
\overline{e}_{a - 1},
\overline{e}'_a,
\overline{e}_{a + 1},
\ldots,
\overline{e}_l]
$$
它由 $\overline{e}_a = \lambda \overline{e}'_a + \overline{x}$ 以及
$\det_\kappa(\overline{M})$ 中相应的容许关系得出。若 $i > a$，
则所需等式为
$$
(-1)^i
\overline{\lambda}
[\overline{e}_1,
\ldots,
\overline{e}_{i - 1},
\overline{e}_{i + 1},
\ldots,
\overline{e}_l]
=
(-1)^i
\overline{\lambda}
[\overline{e}_1,
\ldots,
\overline{e}_{a - 1},
\overline{e}'_a,
\overline{e}_{a + 1},
\ldots,
\overline{e}_{i - 1},
\overline{e}_{i + 1},
\ldots,
\overline{e}_l]
$$
它同样由 $\overline{e}_a = \lambda \overline{e}'_a + \overline{x}$ 以及
$\det_\kappa(\overline{M})$ 中相应的容许关系得出。需要特别处理的是
$i = a$ 的情形。此时
$e_a = \lambda e'_a + x = \mu f + y$，从而还有
$e'_a = \lambda^{-1}(\mu f + y - x)$。于是
$$
\psi([e_1, \ldots, e_l])
= (-1)^i \overline{\mu}
[\overline{e}_1, \ldots,
\overline{e}_{i - 1},
\overline{e}_{i + 1}, \ldots, \overline{e}_l]
=
\psi(
\overline{\lambda}
[e_1, \ldots, e_{a - 1}, e'_a, e_{a + 1}, \ldots, e_l]
)
$$
正如所需。

\medskip\noindent
类型 (c) 的关系。设 $(e_1, \ldots, e_l)$ 是容许序列，且
$\mathfrak m e_a \subset \langle e_1, \ldots, e_{a - 2}\rangle$。
设 $i$ 是使 $f \in \langle e_1, \ldots, e_i\rangle$ 成立的最小指标。
写成 $e_i = \lambda f + x$，其中
$x \in \langle e_1, \ldots, e_{i - 1}\rangle$。分成 $4$ 种情形：

\medskip\noindent
情形 1：$i < a - 1$。所需等式为
\begin{align*}
& (-1)^i
\overline{\lambda}
[\overline{e}_1,
\ldots,
\overline{e}_{i - 1},
\overline{e}_{i + 1},
\ldots,
\overline{e}_l] \\
& =
(-1)^{i + 1}
\overline{\lambda}
[\overline{e}_1,
\ldots,
\overline{e}_{i - 1},
\overline{e}_{i + 1},
\ldots,
\overline{e}_{a - 2},
\overline{e}_a,
\overline{e}_{a - 1},
\overline{e}_{a + 1},
\ldots,
\overline{e}_l]
\end{align*}
它由 $\det_\kappa(\overline{M})$ 中类型 (c) 的容许关系得出。

\medskip\noindent
情形 2：$i > a$。所需等式为
\begin{align*}
& (-1)^i
\overline{\lambda}
[\overline{e}_1,
\ldots,
\overline{e}_{i - 1},
\overline{e}_{i + 1},
\ldots,
\overline{e}_l] \\
& =
(-1)^{i + 1}
\overline{\lambda}
[\overline{e}_1,
\ldots,
\overline{e}_{a - 2},
\overline{e}_a,
\overline{e}_{a - 1},
\overline{e}_{a + 1},
\ldots,
\overline{e}_{i - 1},
\overline{e}_{i + 1},
\ldots,
\overline{e}_l]
\end{align*}
它由 $\det_\kappa(\overline{M})$ 中类型 (c) 的容许关系得出。

\medskip\noindent
情形 3：$i = a$。写成 $e_a = \lambda f + \mu e_{a - 1} + y$，其中
$y \in \langle e_1, \ldots, e_{a - 2}\rangle$。于是根据定义，
$$
\psi([e_1, \ldots, e_l]) =
(-1)^a
\overline{\lambda}
[\overline{e}_1,
\ldots,
\overline{e}_{a - 1},
\overline{e}_{a + 1},
\ldots,
\overline{e}_l]
$$
若 $\overline{\mu}$ 非零，则
$e_{a - 1} = - \mu^{-1} \lambda f + \mu^{-1}e_a - \mu^{-1} y$，
因而根据定义得到
$$
\psi(-[e_1, \ldots, e_{a - 2}, e_a, e_{a - 1}, e_{a + 1}, \ldots, e_l]) =
(-1)^a
\overline{\mu^{-1}\lambda}
[\overline{e}_1,
\ldots,
\overline{e}_{a - 2},
\overline{e}_a,
\overline{e}_{a + 1},
\ldots,
\overline{e}_l]
$$
由于在 $\overline{M}$ 中有
$\overline{e}_a = \mu \overline{e}_{a - 1} + \overline{y}$，由
$\det_\kappa(\overline{M})$ 的关系 (a) 可知两个结果相等。
另一方面，若 $\overline{\mu}$ 为零，则可写成
$e_a = \lambda f + y$，其中
$y \in \langle e_1, \ldots, e_{a - 2}\rangle$，而且
$$
\psi(-[e_1, \ldots, e_{a - 2}, e_a, e_{a - 1}, e_{a + 1}, \ldots, e_l]) =
(-1)^a
\overline{\lambda}
[\overline{e}_1,
\ldots,
\overline{e}_{a - 1},
\overline{e}_{a + 1},
\ldots,
\overline{e}_l]
$$
它等于 $\psi([e_1, \ldots, e_l])$。

\medskip\noindent
情形 4：$i = a - 1$。此时根据定义，
$$
\psi([e_1, \ldots, e_l]) =
(-1)^{a - 1}
\overline{\lambda}
[\overline{e}_1,
\ldots,
\overline{e}_{a - 2},
\overline{e}_a,
\ldots,
\overline{e}_l]
$$
若 $f \not \in \langle e_1, \ldots, e_{a - 2}, e_a \rangle$，则
$$
\psi(-[e_1, \ldots, e_{a - 2}, e_a, e_{a - 1}, e_{a + 1}, \ldots, e_l]) =
(-1)^{a + 1}\overline{\lambda}
[\overline{e}_1,
\ldots,
\overline{e}_{a - 2},
\overline{e}_a,
\ldots,
\overline{e}_l]
$$
由于 $(-1)^{a - 1} = (-1)^{a + 1}$，两个表达式相同。
最后假设 $f \in \langle e_1, \ldots, e_{a - 2}, e_a \rangle$。
此时可见 $e_{a - 1} = \lambda f + x$，其中
$x \in \langle e_1, \ldots, e_{a - 2}\rangle$，并且
$e_a = \mu f + y$，其中
$y \in \langle e_1, \ldots, e_{a - 2}\rangle$，这里
$\lambda, \mu \in R$ 均为单位。于是同时有
$e_a \in \langle e_1, \ldots, e_{a - 1} \rangle$ 和
$e_{a - 1} \in \langle e_1, \ldots, e_{a - 2}, e_a\rangle$。
在这种情形，类型 (a) 的关系同时适用于
$[e_1, \ldots, e_l]$ 与
$[e_1, \ldots, e_{a - 2}, e_a, e_{a - 1}, e_{a + 1}, \ldots, e_l]$。
结合上面已经证明的 $\psi$ 与这些关系的相容性，可见
$$
\psi([e_1, \ldots, e_l])
\quad\text{且}\quad
\psi([e_1, \ldots, e_{a - 2}, e_a, e_{a - 1}, e_{a + 1}, \ldots, e_l])
$$
二者都为零，正如所需。

\medskip\noindent
至此已经证明 $\psi$ 良定义，只剩证明它是满射。为此，设
$(\overline{f}_2, \ldots, \overline{f}_l)$ 是 $\overline{M}$ 中的
容许序列。可以选取提升 $f_2, \ldots, f_l \in M$，于是
$(f, f_2, \ldots, f_l)$ 是 $M$ 中的容许序列。由于
$\psi([f, f_2, \ldots, f_l]) = [f_2, \ldots, f_l]$，结论成立。
\end{proof}

\noindent
设 $R$ 为局部环，其极大理想为 $\mathfrak m$，剩余域为 $\kappa$。
注意，若 $\varphi : M \to N$ 是有限长度 $R$-模的同构，则由规则
$$
\det\nolimits_\kappa(\varphi) :
\det\nolimits_\kappa(M)
\to
\det\nolimits_\kappa(N)
$$
以及
$$
\det\nolimits_\kappa(\varphi)([e_1, \ldots, e_l])
=
[\varphi(e_1), \ldots, \varphi(e_l)]
$$
对 $M$ 的任意符号 $[e_1, \ldots, e_l]$，得到一个同构。
因此可见，$\det\nolimits_\kappa$ 是函子
\begin{equation}
\label{chow-equation-functor}
\left\{
\begin{matrix}
\text{有限长度 }R\text{-模}\\
\text{（以同构为态射）}
\end{matrix}
\right\}
\longrightarrow
\left\{
\begin{matrix}
1\text{-维 }\kappa\text{-向量空间}\\
\text{（以同构为态射）}
\end{matrix}
\right\}
\end{equation}
这是“行列式函子”的典型情形（见 \cite{Knudsen}），下面的可加性也是如此。

\begin{lemma}
\label{chow-lemma-det-exact-sequences}
设 $(R, \mathfrak m, \kappa)$ 为局部环。对有限长度 $R$-模的每个短正合列
$$
0 \to K \to L \to M \to 0
$$
都存在一个典范同构
$$
\gamma_{K \to L \to M} :
\det\nolimits_\kappa(K) \otimes_\kappa \det\nolimits_\kappa(M)
\longrightarrow
\det\nolimits_\kappa(L)
$$
它在非零符号上由规则
$$
[e_1, \ldots, e_k]
\otimes
[\overline{f}_1, \ldots, \overline{f}_m]
\longrightarrow
[e_1, \ldots, e_k, f_1, \ldots, f_m]
$$
定义，并具有下列性质：
\begin{enumerate}
\item 对短正合列的每个同构，即对每个交换图
$$
\xymatrix{
0 \ar[r] &
K \ar[r] \ar[d]^u &
L \ar[r] \ar[d]^v &
M \ar[r] \ar[d]^w &
0 \\
0 \ar[r] &
K' \ar[r] &
L' \ar[r] &
M' \ar[r] &
0
}
$$
若两行短正合且 $u, v, w$ 均为同构，则
$$
\gamma_{K' \to L' \to M'} \circ
(\det\nolimits_\kappa(u) \otimes \det\nolimits_\kappa(w))
=
\det\nolimits_\kappa(v) \circ
\gamma_{K \to L \to M},
$$
\item 对有限长度 $R$-模的每个行列均正合的交换方图
$$
\xymatrix{
& 0 \ar[d] & 0 \ar[d] & 0 \ar[d] & \\
0 \ar[r] & A \ar[r] \ar[d] & B \ar[r] \ar[d] & C \ar[r] \ar[d] & 0 \\
0 \ar[r] & D \ar[r] \ar[d] & E \ar[r] \ar[d] & F \ar[r] \ar[d] & 0 \\
0 \ar[r] & G \ar[r] \ar[d] & H \ar[r] \ar[d] & I \ar[r] \ar[d] & 0 \\
& 0  & 0  & 0  &
}
$$
下图交换：
$$
\xymatrix{
\det\nolimits_\kappa(A) \otimes
\det\nolimits_\kappa(C) \otimes
\det\nolimits_\kappa(G) \otimes
\det\nolimits_\kappa(I)
\ar[dd]_{\epsilon}
\ar[rrr]_-{\gamma_{A \to B \to C} \otimes \gamma_{G \to H \to I}}
& & &
\det\nolimits_\kappa(B) \otimes
\det\nolimits_\kappa(H)
\ar[d]^{\gamma_{B \to E \to H}}
\\
& & & \det\nolimits_\kappa(E)
\\
\det\nolimits_\kappa(A) \otimes
\det\nolimits_\kappa(G) \otimes
\det\nolimits_\kappa(C) \otimes
\det\nolimits_\kappa(I)
\ar[rrr]^-{\gamma_{A \to D \to G} \otimes \gamma_{C \to F \to I}}
& & &
\det\nolimits_\kappa(D) \otimes
\det\nolimits_\kappa(F)
\ar[u]_{\gamma_{D \to E \to F}}
}
$$
其中 $\epsilon$ 是交换张量积因子的映射再乘以 $(-1)^{cg}$，这里
$c = \text{length}_R(C)$ 且 $g = \text{length}_R(G)$；
\item 若 $0 \to K \to L \to M \to 0$ 实际上是 $\kappa$-向量空间的
短正合列，则映射 $\gamma_{K \to L \to M}$ 与通常的同构一致。
\end{enumerate}
\end{lemma}

\begin{proof}
在映射 $\gamma_{K \to L \to M}$ 的显式描述中采用非零符号，原因如下。
若 $(e_1, \ldots, e_l)$ 是 $K$ 中的容许序列，而
$(\overline{f}_1, \ldots, \overline{f}_m)$ 是 $M$ 中的容许序列，
则 $(e_1, \ldots, e_l, f_1, \ldots, f_m)$ 不一定是 $L$ 中的容许序列
（这里 $f_i \in L$ 当然表示 $\overline{f}_i$ 的一个提升）。然而，若符号
$[e_1, \ldots, e_l]$ 在 $\det_\kappa(K)$ 中非零，则必有
$K = \langle e_1, \ldots, e_k\rangle$（见引理
\ref{chow-lemma-dimension-at-most-one} 的证明），在这种情形，
$(e_1, \ldots, e_k, f_1, \ldots, f_m)$ 的确是容许序列。
此外，由 $\det_\kappa(L)$ 中类型 (b) 的容许关系可见，
$[e_1, \ldots, e_k, f_1, \ldots, f_m]$ 在 $\det_\kappa(L)$ 中的值
也与提升 $f_i$ 的选取无关。根据这一点，很明显，$K$ 中
$e_1, \ldots, e_k$ 之间的容许关系会转化为 $L$ 中
$e_1, \ldots, e_k, f_1, \ldots, f_m$ 之间的容许关系；
$\overline{f}_1, \ldots, \overline{f}_m$ 之间的容许关系亦然。
因此，$\gamma$ 按引理所述定义了向量空间的线性映射。

\medskip\noindent
由引理 \ref{chow-lemma-determinant-dimension-one}，$\det_\kappa(L)$ 由任意一个
符号 $[x_1, \ldots, x_{k + m}]$ 生成，只要
$(x_1, \ldots, x_{k + m})$ 是满足
$L = \langle x_1, \ldots, x_{k + m}\rangle$ 的容许序列。
所以映射 $\gamma_{K \to L \to M}$ 显然满射，因而是同构。

\medskip\noindent
性质 (1) 成立，因为
\begin{eqnarray*}
& & \det\nolimits_\kappa(v)([e_1, \ldots, e_k, f_1, \ldots, f_m]) \\
& = &
[v(e_1), \ldots, v(e_k), v(f_1), \ldots, v(f_m)] \\
& = &
\gamma_{K' \to L' \to M'}([u(e_1), \ldots, u(e_k)]
\otimes [w(f_1), \ldots, w(f_m)]).
\end{eqnarray*}
性质 (2) 意味着：给定生成 $\det_\kappa(A)$ 的符号
$[\alpha_1, \ldots, \alpha_a]$、生成 $\det_\kappa(C)$ 的符号
$[\gamma_1, \ldots, \gamma_c]$、生成 $\det_\kappa(G)$ 的符号
$[\zeta_1, \ldots, \zeta_g]$，以及生成 $\det_\kappa(I)$ 的符号
$[\iota_1, \ldots, \iota_i]$，有
\begin{eqnarray*}
& & [\alpha_1, \ldots, \alpha_a, \tilde\gamma_1, \ldots, \tilde\gamma_c,
\tilde\zeta_1, \ldots, \tilde\zeta_g, \tilde\iota_1, \ldots, \tilde\iota_i] \\
& = &
(-1)^{cg} [\alpha_1, \ldots, \alpha_a, \tilde\zeta_1, \ldots, \tilde\zeta_g,
\tilde\gamma_1, \ldots, \tilde\gamma_c, \tilde\iota_1, \ldots, \tilde\iota_i]
\end{eqnarray*}
在 $\det_\kappa(E)$ 中成立（其中 $\tilde{x}$ 是 $E$ 中适当的提升）。
这是因为可以按下列次序使用类型 (c) 的容许关系 $cg$ 次：先将
$\tilde\zeta_1$ 越过各元素
$\tilde\gamma_c, \ldots, \tilde\gamma_1$
（这是允许的，因为 $\mathfrak m\tilde\zeta_1 \subset A$），
再将 $\tilde\zeta_2$ 越过各元素
$\tilde\gamma_c, \ldots, \tilde\gamma_1$
（这是允许的，因为 $\mathfrak m\tilde\zeta_2 \subset A + R\tilde\zeta_1$），
依此类推。

\medskip\noindent
引理的第 (3) 部分显然成立。证明完毕。
\end{proof}

\noindent
可以利用引理中的映射 $\gamma$ 如下定义更一般的映射 $\gamma$。
设 $(R, \mathfrak m, \kappa)$ 为局部环，$M$ 为有限长度 $R$-模，并设给定
一个有限滤过（见《Homology》定义 \ref{homology-definition-filtered}）
$$
0 = F^m \subset F^{m - 1} \subset \ldots \subset F^{n + 1} \subset F^n = M
$$
则有一个良定义的典范同构
$$
\gamma_{(M, F)} :
\det\nolimits_\kappa(F^{m - 1}/F^m) \otimes_\kappa \ldots \otimes_k 
\det\nolimits_\kappa(F^n/F^{n + 1})
\longrightarrow
\det\nolimits_\kappa(M)
$$
为构造它，我们使用由短正合列
$0 \to F^{i - 1}/F^i \to M/F^i \to M/F^{i - 1} \to 0$
给出的引理 \ref{chow-lemma-det-exact-sequences} 中的同构。
引理 \ref{chow-lemma-det-exact-sequences} 的第 (2) 部分取 $G = 0$ 表明，若改用
短正合列 $0 \to F^i \to F^{i - 1} \to F^{i - 1}/F^i \to 0$，
得到的是同一个同构。

\medskip\noindent
下面是行列式函子的另一个典型结果。它不难证明；通常困难之处是证明
行列式函子的存在性。

\begin{lemma}
\label{chow-lemma-uniqueness-det}
设 $(R, \mathfrak m, \kappa)$ 为任意局部环。赋有映射
$\gamma_{K \to L \to M}$ 的函子
$$
\det\nolimits_\kappa :
\left\{
\begin{matrix}
\text{有限长度 }R\text{-模} \\
\text{（以同构为态射）}
\end{matrix}
\right\}
\longrightarrow
\left\{
\begin{matrix}
1\text{-维 }\kappa\text{-向量空间} \\
\text{（以同构为态射）}
\end{matrix}
\right\}
$$
由下列性质刻画：
\begin{enumerate}
\item 它限制到被 $\mathfrak m$ 零化的模所成子范畴后，与通常的行列式
函子同构（见引理 \ref{chow-lemma-compare-det}）；
\item 引理 \ref{chow-lemma-det-exact-sequences} 的 (1)、(2)、(3) 均成立。
\end{enumerate}
\end{lemma}

\begin{proof}
从略。
\end{proof}

\begin{lemma}
\label{chow-lemma-determinant-quotient-ring}
设 $(R', \mathfrak m') \to (R, \mathfrak m)$ 为局部环同态，并且它在
剩余域 $\kappa$ 上诱导同构。则对每个有限长度 $R$-模，限制
$M_{R'}$ 是有限长度 $R'$-模，而且存在典范同构
$$
\det\nolimits_{R, \kappa}(M)
\longrightarrow
\det\nolimits_{R', \kappa}(M_{R'})
$$
该同构关于 $M$ 函子性成立，并与分别为 $\det_{R, \kappa}$ 和
$\det_{R', \kappa}$ 定义的引理 \ref{chow-lemma-det-exact-sequences} 中的同构
$\gamma_{K \to L \to M}$ 相容。
\end{lemma}

\begin{proof}
若 $M$ 作为 $R$-模的长度为 $l$，则 $M$ 作为 $R'$-模（即 $M_{R'}$）
的长度也为 $l$，见《Algebra》引理 \ref{algebra-lemma-pushdown-module}。
注意，$M$ 在 $R$ 上的容许序列 $x_1, \ldots, x_l$ 也是 $M$ 在
$R'$ 上的容许序列，因为 $\mathfrak m'$ 映入 $\mathfrak m$。
所需同构把符号
$[x_1, \ldots, x_l] \in \det\nolimits_{R, \kappa}(M)$
映到相应的符号
$[x_1, \ldots, x_l] \in \det\nolimits_{R', \kappa}(M)$。
可以立即验证，它关于同构函子性成立，并与引理
\ref{chow-lemma-det-exact-sequences} 中的同构 $\gamma$ 相容。
\end{proof}

\begin{remark}
\label{chow-remark-explain-determinant}
设 $(R, \mathfrak m, \kappa)$ 为局部环，并假设 $\kappa$ 的特征为零，
或者其特征为 $p$ 且 $p R = 0$。设 $M_1, \ldots, M_n$ 是有限长度
$R$-模。下面将证明，存在一个理想 $I \subset \mathfrak m$，对
$i = 1, \ldots, n$ 零化 $M_i$，并且存在典范满射
$R/I \to \kappa$ 的一个截面 $\sigma : \kappa \to R/I$。
通过 $\sigma$ 限制得到的 $M_i$ 的模 $M_{i, \kappa}$ 是维数为
$l_i = \text{length}_R(M_i)$ 的 $\kappa$-向量空间；由引理
\ref{chow-lemma-determinant-quotient-ring} 可见
$$
\det\nolimits_\kappa(M_i) = \wedge_\kappa^{l_i}(M_{i, \kappa})
$$
对于被 $I$ 零化的有限长度 $R$-模所成的短正合列，这些同构与引理
\ref{chow-lemma-det-exact-sequences} 中的同构
$\gamma_{K \to M \to L}$ 相容。由此可知，验证 $\det_\kappa$ 的一个性质，
往往可以归结为在有限维向量空间范畴上验证通常行列式的相应性质。

\medskip\noindent
作为 $I$，可以取模 $M = \bigoplus M_i$ 的零化子
（《Algebra》定义 \ref{algebra-definition-annihilator}）。此时
$R/I \subset \text{End}_R(M)$，因而它具有有限长度。因此 $R/I$ 是
剩余域为 $\kappa$ 的阿廷局部环。由于阿廷局部环完备，由 Cohen 结构定理
（《Algebra》定理 \ref{algebra-theorem-cohen-structure-theorem}）可知
$R/I$ 有系数环；根据我们对 $R$ 的假设，该系数环是域。
\end{remark}

\noindent
下面给出一种能够计算线性映射行列式的情形。事实上，任何情形都没有
神秘之处；一个任取的例子见例 \ref{chow-example-determinant-map}。

\begin{lemma}
\label{chow-lemma-times-u-determinant}
设 $R$ 为剩余域为 $\kappa$ 的局部环，$u \in R^*$ 为单位，
$M$ 为 $R$ 上有限长度模。把乘以 $u$ 的映射记为 $u_M : M \to M$。
则
$$
\det\nolimits_\kappa(u_M) :
\det\nolimits_\kappa(M)
\longrightarrow
\det\nolimits_\kappa(M)
$$
等于乘以 $\overline{u}^l$ 的映射，其中 $l = \text{length}_R(M)$，而
$\overline{u} \in \kappa^*$ 是 $u$ 的像。
\end{lemma}

\begin{proof}
把满足
$\det\nolimits_\kappa(u_M) = f_M \text{id}_{\det\nolimits_\kappa(M)}$
的元素记为 $f_M \in \kappa^*$。设
$0 \to K \to L \to M \to 0$ 是有限 $R$-模的短正合列。
于是 $u_k$、$u_L$、$u_M$ 给出短正合列的一个同构。因此由引理
\ref{chow-lemma-det-exact-sequences} (1) 得到 $f_K f_M = f_L$。
这表明，对长度归纳后，只需证明长度为 $1$ 的情形，而该情形显然成立。
\end{proof}

\begin{example}
\label{chow-example-determinant-map}
考虑局部环 $R = \mathbf{Z}_p$。令
$M = \mathbf{Z}_p/(p^2) \oplus \mathbf{Z}_p/(p^3)$，并设
$u : M \to M$ 是由矩阵
$$
u =
\left(
\begin{matrix}
a & b \\
pc & d
\end{matrix}
\right)
$$
给出的映射，其中 $a, b, c, d \in \mathbf{Z}_p$，且
$a, d \in \mathbf{Z}_p^*$。在这种情形，$\det_\kappa(u)$ 等于乘以
$a^2d^3 \bmod p \in \mathbf{F}_p^*$。考察 $u$ 在符号
$[p^2e, pe, pf, e, f]$ 上的作用即可容易看出这一点，其中
$e = (0 , 1) \in M$ 且 $f = (1, 0) \in M$。
\end{example}





\subsection{周期复形与行列式}
\label{chow-subsection-periodic-complexes-determinants}

\noindent
设 $R$ 为剩余域为 $\kappa$ 的局部环，$(M, \varphi, \psi)$ 为 $R$ 上的
$(2, 1)$-周期复形。假设 $M$ 具有有限长度，且
$(M, \varphi, \psi)$ 正合。我们将利用行列式构造定义这种情形的一个不变量；
见小节 \ref{chow-subsection-determinants-finite-length}。简记
$K_\varphi = \Ker(\varphi)$、
$I_\varphi = \Im(\varphi)$、
$K_\psi = \Ker(\psi)$ 以及
$I_\psi = \Im(\psi)$。短正合列
$$
0 \to K_\varphi \to M \to I_\varphi \to 0, \quad
0 \to K_\psi \to M \to I_\psi \to 0
$$
给出同构
$$
\gamma_\varphi :
\det\nolimits_\kappa(K_\varphi)
\otimes
\det\nolimits_\kappa(I_\varphi)
\longrightarrow
\det\nolimits_\kappa(M), \quad
\gamma_\psi :
\det\nolimits_\kappa(K_\psi)
\otimes
\det\nolimits_\kappa(I_\psi)
\longrightarrow
\det\nolimits_\kappa(M),
$$
见引理 \ref{chow-lemma-det-exact-sequences}。另一方面，复形的正合性给出等式
$K_\varphi = I_\psi$ 和 $K_\psi = I_\varphi$，从而通过交换因子得到同构
$$
\sigma :
\det\nolimits_\kappa(K_\varphi)
\otimes
\det\nolimits_\kappa(I_\varphi)
\longrightarrow
\det\nolimits_\kappa(K_\psi)
\otimes
\det\nolimits_\kappa(I_\psi)
$$
利用这些记号，可以定义所需不变量。

\begin{definition}
\label{chow-definition-periodic-determinant}
设 $R$ 为剩余域为 $\kappa$ 的局部环，$(M, \varphi, \psi)$ 为 $R$ 上的
$(2, 1)$-周期复形。假设 $M$ 具有有限长度，且
$(M, \varphi, \psi)$ 正合。$(M, \varphi, \psi)$ 的{\it 行列式}是元素
$$
\det\nolimits_\kappa(M, \varphi, \psi) \in \kappa^*
$$
它满足：复合
$$
\det\nolimits_\kappa(M)
\xrightarrow{\gamma_\psi \circ \sigma \circ \gamma_\varphi^{-1}}
\det\nolimits_\kappa(M)
$$
等于乘以
$(-1)^{\text{length}_R(I_\varphi)\text{length}_R(I_\psi)}
\det\nolimits_\kappa(M, \varphi, \psi)$ 的映射。
\end{definition}

\begin{remark}
\label{chow-remark-more-elementary}
下面给出上述行列式更具体的描述。设 $R$ 为剩余域为 $\kappa$ 的局部环，
$(M, \varphi, \psi)$ 为 $R$ 上的 $(2, 1)$-周期复形。假设 $M$ 具有
有限长度，且 $(M, \varphi, \psi)$ 正合。仍简记
$I_\varphi = \Im(\varphi)$、$I_\psi = \Im(\psi)$。
假设 $\text{length}_R(I_\varphi) = a$ 且
$\text{length}_R(I_\psi) = b$；由正合性，
$a + b = \text{length}_R(M)$。选取容许序列
$x_1, \ldots, x_a \in I_\varphi$ 和 $y_1, \ldots, y_b \in I_\psi$，
使符号 $[x_1, \ldots, x_a]$ 生成 $\det_\kappa(I_\varphi)$，而符号
$[x_1, \ldots, x_b]$ 生成 $\det_\kappa(I_\psi)$。选取
$\tilde x_i \in M$ 使得 $\varphi(\tilde x_i) = x_i$，并选取
$\tilde y_j \in M$ 使得 $\psi(\tilde y_j) = y_j$。于是
$\det_\kappa(M, \varphi, \psi)$ 由等式
$$
[x_1, \ldots, x_a, \tilde y_1, \ldots, \tilde y_b]
=
(-1)^{ab} \det\nolimits_\kappa(M, \varphi, \psi)
[y_1, \ldots, y_b, \tilde x_1, \ldots, \tilde x_a]
$$
在 $\det_\kappa(M)$ 中成立这一条件刻画。这也解释了符号因子。
\end{remark}

\begin{lemma}
\label{chow-lemma-periodic-determinant-shift}
设 $R$ 为剩余域为 $\kappa$ 的局部环，$(M, \varphi, \psi)$ 为 $R$ 上的
$(2, 1)$-周期复形。假设 $M$ 具有有限长度，且
$(M, \varphi, \psi)$ 正合。则
$$
\det\nolimits_\kappa(M, \varphi, \psi)
\det\nolimits_\kappa(M, \psi, \varphi)
= 1.
$$
\end{lemma}

\begin{proof}
从略。
\end{proof}

\begin{lemma}
\label{chow-lemma-periodic-determinant-sign}
设 $R$ 为剩余域为 $\kappa$ 的局部环，$(M, \varphi, \varphi)$ 为 $R$ 上的
$(2, 1)$-周期复形。假设 $M$ 具有有限长度，且
$(M, \varphi, \varphi)$ 正合。则
$\text{length}_R(M) = 2 \text{length}_R(\Im(\varphi))$，并且
$$
\det\nolimits_\kappa(M, \varphi, \varphi)
=
(-1)^{\text{length}_R(\Im(\varphi))}
=
(-1)^{\frac{1}{2}\text{length}_R(M)}
$$
\end{lemma}

\begin{proof}
这直接由定义中的符号规则得出。
\end{proof}

\begin{lemma}
\label{chow-lemma-periodic-determinant-easy-case}
设 $R$ 为剩余域为 $\kappa$ 的局部环，$M$ 为有限长度 $R$-模。
\begin{enumerate}
\item 若 $\varphi : M \to M$ 是同构，则
$\det_\kappa(M, \varphi, 0) = \det_\kappa(\varphi)$；
\item 若 $\psi : M \to M$ 是同构，则
$\det_\kappa(M, 0, \psi) = \det_\kappa(\psi)^{-1}$。
\end{enumerate}
\end{lemma}

\begin{proof}
证明 (1)。令 $\psi = 0$。使用定义 \ref{chow-definition-periodic-determinant}
上方的记号，可作认同
$K_\varphi = I_\psi = 0$、$I_\varphi = K_\psi = M$。
在这些认同下，映射
$$
\gamma_\varphi :
\kappa \otimes \det\nolimits_\kappa(M)
=
\det\nolimits_\kappa(K_\varphi)
\otimes
\det\nolimits_\kappa(I_\varphi)
\longrightarrow
\det\nolimits_\kappa(M)
$$
被认同为 $\det_\kappa(\varphi^{-1})$。另一方面，映射
$\gamma_\psi$ 被认同为恒等映射。因此在这种情形，
$\gamma_\psi \circ \sigma \circ \gamma_\varphi^{-1}$ 等于
$\det_\kappa(\varphi)$，结论随即得到。省略 (2) 的证明。
\end{proof}

\begin{lemma}
\label{chow-lemma-periodic-determinant}
设 $R$ 为剩余域为 $\kappa$ 的局部环。假设有 $(2, 1)$-周期复形的短正合列
$$
0 \to (M_1, \varphi_1, \psi_1)
\to (M_2, \varphi_2, \psi_2)
\to (M_3, \varphi_3, \psi_3)
\to 0
$$
其中所有 $M_i$ 都具有有限长度，且每个 $(M_1, \varphi_1, \psi_1)$
都正合。则
$$
\det\nolimits_\kappa(M_2, \varphi_2, \psi_2) =
\det\nolimits_\kappa(M_1, \varphi_1, \psi_1)
\det\nolimits_\kappa(M_3, \varphi_3, \psi_3).
$$
在 $\kappa^*$ 中成立。
\end{lemma}

\begin{proof}
简记
$I_{\varphi, i} = \Im(\varphi_i)$、
$K_{\varphi, i} = \Ker(\varphi_i)$、
$I_{\psi, i} = \Im(\psi_i)$ 以及
$K_{\psi, i} = \Ker(\psi_i)$。注意，有一个有限长度 $R$-模的行列均正合的
交换方图
$$
\xymatrix{
& 0 \ar[d] & 0 \ar[d] & 0 \ar[d] & \\
0 \ar[r] &
K_{\varphi, 1} \ar[r] \ar[d] &
K_{\varphi, 2} \ar[r] \ar[d] &
K_{\varphi, 3} \ar[r] \ar[d] &
0 \\
0 \ar[r] &
M_1 \ar[r] \ar[d] &
M_2 \ar[r] \ar[d] &
M_3 \ar[r] \ar[d] &
0 \\
0 \ar[r] &
I_{\varphi, 1} \ar[r] \ar[d] &
I_{\varphi, 2} \ar[r] \ar[d] &
I_{\varphi, 3} \ar[r] \ar[d] &
0 \\
& 0  & 0  & 0  &
}
$$
顶行正合，因为它可以与像所成的序列
$I_{\psi, 1} \to I_{\psi, 2} \to I_{\psi, 3} \to 0$ 认同；
底行同理。还存在一个涉及模 $I_{\psi, i}$ 与 $K_{\psi, i}$ 的类似图。
按照定义，$\det_\kappa(M_2, \varphi_2, \psi_2)$ 在不计符号时对应于
下图左侧竖直映射的复合：
$$
\xymatrix{
\det_\kappa(M_1) \otimes
\det_\kappa(M_3) \ar[r]^\gamma
\ar[d]^{\gamma^{-1} \otimes \gamma^{-1}} &
\det_\kappa(M_2)
\ar[d]^{\gamma^{-1}} \\
\det\nolimits_\kappa(K_{\varphi, 1})
\otimes
\det\nolimits_\kappa(I_{\varphi, 1})
\otimes
\det\nolimits_\kappa(K_{\varphi, 3})
\otimes
\det\nolimits_\kappa(I_{\varphi, 3})
\ar[d]^{\sigma \otimes \sigma}
\ar[r]^-{\gamma \otimes \gamma} &
\det\nolimits_\kappa(K_{\varphi, 2})
\otimes
\det\nolimits_\kappa(I_{\varphi, 2})
\ar[d]^\sigma
\\
\det\nolimits_\kappa(K_{\psi, 1})
\otimes
\det\nolimits_\kappa(I_{\psi, 1})
\otimes
\det\nolimits_\kappa(K_{\psi, 3})
\otimes
\det\nolimits_\kappa(I_{\psi, 3})
\ar[d]^{\gamma \otimes \gamma}
\ar[r]^-{\gamma \otimes \gamma}
&
\det\nolimits_\kappa(K_{\psi, 2})
\otimes
\det\nolimits_\kappa(I_{\psi, 2})
\ar[d]^\gamma \\
\det_\kappa(M_1)
\otimes
\det_\kappa(M_3) \ar[r]^\gamma
&
\det_\kappa(M_2)
}
$$
应用引理 \ref{chow-lemma-det-exact-sequences} (2)，可知顶部和底部方块在一个
符号因子意义下交换。中间方块显然交换（只是交换因子）。因此可见
$\det\nolimits_\kappa(M_2, \varphi_2, \psi_2) =
\epsilon \det\nolimits_\kappa(M_1, \varphi_1, \psi_1)
\det\nolimits_\kappa(M_3, \varphi_3, \psi_3)
$
对某个符号 $\epsilon$ 成立。该符号可以算出；具体而言，上图外侧矩形交换时
相差
\begin{eqnarray*}
\epsilon & = &
(-1)^{\text{length}(I_{\varphi, 1})\text{length}(K_{\varphi, 3})
+ \text{length}(I_{\psi, 1})\text{length}(K_{\psi, 3})} \\
& = &
(-1)^{\text{length}(I_{\varphi, 1})\text{length}(I_{\psi, 3})
+ \text{length}(I_{\psi, 1})\text{length}(I_{\varphi, 3})}
\end{eqnarray*}
（证明从略）。由此容易看出符号也完全吻合。
\end{proof}

\begin{example}
\label{chow-example-dual-numbers}
设 $k$ 为域。考虑 $k$ 上的对偶数环 $R = k[T]/(T^2)$，并把 $T$ 在
$R$ 中的类记为 $t$。令 $M = R$，并取 $\varphi = ut$、$\psi = vt$，
其中 $u, v \in k^*$。在这种情形，$\det_k(M)$ 有生成元
$e = [t, 1]$。我们作认同
$I_\varphi = K_\varphi = I_\psi = K_\psi = (t)$。于是
$\gamma_\varphi(t \otimes t) = u^{-1}[t, 1]$
（因为 $u^{-1} \in M$ 是 $t \in I_\varphi$ 的一个提升），并且
$\gamma_\psi(t \otimes t) = v^{-1}[t, 1]$（理由相同）。因此
$\det_k(M, \varphi, \psi) = -u/v \in k^*$。
\end{example}

\begin{example}
\label{chow-example-Zp}
设 $R = \mathbf{Z}_p$ 且 $M = \mathbf{Z}_p/(p^l)$。令
$\varphi = p^b u$ 且 $\varphi = p^a v$，其中 $a, b \geq 0$、
$a + b = l$，而 $u, v \in \mathbf{Z}_p^*$。与例
\ref{chow-example-dual-numbers} 中相同的计算表明
\begin{eqnarray*}
\det\nolimits_{\mathbf{F}_p}(\mathbf{Z}_p/(p^l), p^bu, p^av) & = &
(-1)^{ab}u^a/v^b \bmod p \\
& = &
(-1)^{\text{ord}_p(\alpha)\text{ord}_p(\beta)}
\frac{\alpha^{\text{ord}_p(\beta)}}{\beta^{\text{ord}_p(\alpha)}} \bmod p
\end{eqnarray*}
其中 $\alpha = p^bu, \beta = p^av \in \mathbf{Z}_p$。
更一般的情形（及其证明）见引理 \ref{chow-lemma-symbol-is-usual-tame-symbol}。
\end{example}

\begin{example}
\label{chow-example-generic-vector-space}
设 $R = k$ 为域，设 $M = k^{\oplus a} \oplus k^{\oplus b}$ 的维数为
$l = a + b$。令 $\varphi$ 和 $\psi$ 为下列对角矩阵
$$
\varphi = \text{diag}(u_1, \ldots, u_a, 0, \ldots, 0),
\quad
\psi = \text{diag}(0, \ldots, 0, v_1, \ldots, v_b)
$$
其中 $u_i, v_j \in k^*$。此时
$$
\det\nolimits_k(M, \varphi, \psi)
=
\frac{u_1 \ldots u_a}{v_1 \ldots v_b}.
$$
这可以直接计算得到，也可以先计算 $l = 1$ 的情形，再使用引理
\ref{chow-lemma-periodic-determinant} 的可加性。
\end{example}

\begin{example}
\label{chow-example-special-vector-space}
设 $R = k$ 为域，设 $M = k^{\oplus a} \oplus k^{\oplus a}$ 的维数为
$l = 2a$。令 $\varphi$ 和 $\psi$ 为下列分块矩阵
$$
\varphi =
\left(
\begin{matrix}
0 & U \\
0 & 0
\end{matrix}
\right),
\quad
\psi =
\left(
\begin{matrix}
0 & V \\
0 & 0
\end{matrix}
\right),
$$
其中 $U, V \in \text{Mat}(a \times a, k)$ 可逆。此时
$$
\det\nolimits_k(M, \varphi, \psi)
=
(-1)^a\frac{\det(U)}{\det(V)}.
$$
这可以直接计算得到。当 $a = 1$ 时，该计算与例
\ref{chow-example-dual-numbers} 中的计算类似。
\end{example}

\begin{example}
\label{chow-example-a-la-oort}
设 $R = k$ 为域，$M = k^{\oplus 4}$。令
$$
\varphi =
\left(
\begin{matrix}
  0 &   0 &   0 &   0 \\
u_1 &   0 &   0 &   0 \\
  0 &   0 &   0 &   0 \\
  0 &   0 & u_2 &   0
\end{matrix}
\right)
\quad
\varphi =
\left(
\begin{matrix}
  0 &   0 &   0 &   0 \\
  0 &   0 & v_2 &   0 \\
  0 &   0 &   0 &   0 \\
v_1 &   0 &   0 &   0
\end{matrix}
\right)
\quad
$$
其中 $u_1, u_2, v_1, v_2 \in k^*$。则
$$
\det\nolimits_k(M, \varphi, \psi) = -\frac{u_1u_2}{v_1v_2}.
$$
\end{example}

\noindent
下面讨论线性自同态复合的行列式等于各行列式之积这一事实的对应结论。
为避免公式过长，对任意 $R$-模映射 $\varphi : M \to M$，记
$I_\varphi = \Im(\varphi)$，并记
$K_\varphi = \Ker(\varphi)$。对一对态射
$\varphi, \psi : M \to M$，还记
$\varphi\psi = \varphi \circ \psi$。

\begin{lemma}
\label{chow-lemma-multiplicativity-determinant}
设 $R$ 为剩余域为 $\kappa$ 的局部环，$M$ 为有限长度 $R$-模，
$\alpha, \beta, \gamma$ 为 $M$ 的自同态。假设
\begin{enumerate}
\item $I_\alpha = K_{\beta\gamma}$，而且对 $\alpha, \beta, \gamma$
的任意置换也有类似等式；
\item $K_\alpha = I_{\beta\gamma}$，而且对 $\alpha, \beta, \gamma$
的任意置换也有类似等式。
\end{enumerate}
则
\begin{enumerate}
\item 三元组 $(M, \alpha, \beta\gamma)$ 是正合 $(2, 1)$-周期复形；
\item 三元组 $(I_\gamma, \alpha, \beta)$ 是正合 $(2, 1)$-周期复形；
\item 三元组 $(M/K_\beta, \alpha, \gamma)$ 是正合 $(2, 1)$-周期复形；
\item 有
$$
\det\nolimits_\kappa(M, \alpha, \beta\gamma)
=
\det\nolimits_\kappa(I_\gamma, \alpha, \beta)
\det\nolimits_\kappa(M/K_\beta, \alpha, \gamma).
$$
\end{enumerate}
\end{lemma}

\begin{proof}
假设显然蕴含引理的第 (1) 部分。

\medskip\noindent
为看出第 (1) 部分，注意假设蕴含
$I_{\gamma\alpha} = I_{\alpha\gamma}$；对核及任意其他一对态射也类似。
此外，$I_{\gamma\beta} =I_{\beta\gamma} = K_\alpha \subset I_\gamma$；
对任意其他一对也类似。特别地，得到短正合列
$$
0 \to I_{\beta\gamma} \to I_\gamma \xrightarrow{\alpha} I_{\alpha\gamma} \to 0
$$
以及类似的短正合列
$$
0 \to I_{\alpha\gamma} \to I_\gamma \xrightarrow{\beta} I_{\beta\gamma} \to 0.
$$
这证明 $(I_\gamma, \alpha, \beta)$ 是正合 $(2, 1)$-周期复形，故引理的
第 (2) 部分成立。

\medskip\noindent
为看出 $\alpha$、$\gamma$ 在 $M/K_\beta$ 上给出良定义的自同态，必须验证
$\alpha(K_\beta) \subset K_\beta$ 和 $\gamma(K_\beta) \subset K_\beta$。
这确实成立，因为
$\alpha(K_\beta) = \alpha(I_{\gamma\alpha}) = I_{\alpha\gamma\alpha}
\subset I_{\alpha\gamma} = K_\beta$；另一情形类似。映射
$\alpha : M/K_\beta \to M/K_\beta$ 的核为
$K_{\beta\alpha}/K_\beta = I_\gamma/K_\beta$。类似地，映射
$\gamma : M/K_\beta \to M/K_\beta$ 的核等于 $I_\alpha/K_\beta$。
因此 (3) 成立。

\medskip\noindent
记 $r = \text{length}_R(K_\alpha)$、
$s = \text{length}_R(K_\beta)$、$t = \text{length}_R(K_\gamma)$。
由上述正合列及假设，有
$\text{length}_R(I_\alpha) = s + t$、$\text{length}_R(I_\beta) = r + t$、
$\text{length}_R(I_\gamma) = r + s$，以及
$\text{length}(M) = r + s + t$。选取
\begin{enumerate}
\item 生成 $K_\alpha$ 的容许序列 $x_1, \ldots, x_r \in K_\alpha$；
\item 生成 $K_\beta$ 的容许序列 $y_1, \ldots, y_s \in K_\beta$；
\item 生成 $K_\gamma$ 的容许序列 $z_1, \ldots, z_t \in K_\gamma$；
\item 元素 $\tilde x_i \in M$，使得 $\beta\gamma\tilde x_i = x_i$；
\item 元素 $\tilde y_i \in M$，使得 $\alpha\gamma\tilde y_i = y_i$；
\item 元素 $\tilde z_i \in M$，使得 $\beta\alpha\tilde z_i = z_i$。
\end{enumerate}
在这些选取下，序列
$y_1, \ldots, y_s, \alpha\tilde z_1, \ldots, \alpha\tilde z_t$
是 $I_\alpha$ 中生成它的容许序列。因此由注
\ref{chow-remark-more-elementary}，行列式
$D = \det_\kappa(M, \alpha, \beta\gamma)$ 是满足下式的唯一元素
$\kappa^*$：
\begin{align*}
[y_1, \ldots, y_s,
\alpha\tilde z_1, \ldots, \alpha\tilde z_s,
\tilde x_1, \ldots, \tilde x_r] \\
= (-1)^{r(s + t)} D
[x_1, \ldots, x_r,
\gamma\tilde y_1, \ldots, \gamma\tilde y_s,
\tilde z_1, \ldots, \tilde z_t]
\end{align*}
由同一注，$D_1 = \det_\kappa(M/K_\beta, \alpha, \gamma)$ 由下式刻画：
$$
[y_1, \ldots, y_s,
\alpha\tilde z_1, \ldots, \alpha\tilde z_t,
\tilde x_1, \ldots, \tilde x_r]
=
(-1)^{rt} D_1
[y_1, \ldots, y_s,
\gamma\tilde x_1, \ldots, \gamma\tilde x_r,
\tilde z_1, \ldots, \tilde z_t]
$$
同样，$D_2 = \det_\kappa(I_\gamma, \alpha, \beta)$ 由下式刻画：
$$
[y_1, \ldots, y_s,
\gamma\tilde x_1, \ldots, \gamma\tilde x_r,
\tilde z_1, \ldots, \tilde z_t]
=
(-1)^{rs} D_2
[x_1, \ldots, x_r,
\gamma\tilde y_1, \ldots, \gamma\tilde y_s,
\tilde z_1, \ldots, \tilde z_t]
$$
合并上述公式即得 $D = D_1 D_2$，正如所需。
\end{proof}


\begin{lemma}
\label{chow-lemma-tricky}
设 $R$ 为剩余域为 $\kappa$ 的局部环。设
$\alpha : (M, \varphi, \psi) \to (M', \varphi', \psi')$
为 $R$ 上 $(2, 1)$-周期复形的态射。假设
\begin{enumerate}
\item $M$、$M'$ 具有有限长度；
\item $(M, \varphi, \psi)$、$(M', \varphi', \psi')$ 正合；
\item 映射 $\varphi$、$\psi$ 在 $K = \Ker(\alpha)$ 上诱导零映射；
\item 映射 $\varphi$、$\psi$ 在 $Q = \Coker(\alpha)$ 上诱导零映射。
\end{enumerate}
记 $N = \alpha(M) \subset M'$。得到 $(2, 1)$-周期复形的两个短正合列
$$
\begin{matrix}
0 \to (N, \varphi', \psi') \to (M', \varphi', \psi') \to (Q, 0, 0) \to 0 \\
0 \to (K, 0, 0) \to (M, \varphi, \psi) \to (N, \varphi', \psi') \to 0
\end{matrix}
$$
它们诱导两个同构 $\alpha_i : Q \to K$，$i = 0, 1$。则
$$
\det\nolimits_\kappa(M, \varphi, \psi)
=
\det\nolimits_\kappa(\alpha_0^{-1} \circ \alpha_1)
\det\nolimits_\kappa(M', \varphi', \psi')
$$
特别地，若 $\alpha_0 = \alpha_1$，则
$\det\nolimits_\kappa(M, \varphi, \psi) =
\det\nolimits_\kappa(M', \varphi', \psi')$。
\end{lemma}

\begin{proof}
证明此引理（至少）有两种方法。一种是利用行列式的性质构造一个庞大的交换图；
另一种是使用以元素的容许序列表述的行列式刻画。我们采用第二种方法。

\medskip\noindent
先精确说明映射 $\alpha_i$。具体而言，$\alpha_0$ 是复合
$$
\alpha_0 : Q = H^0(Q, 0, 0) \to H^1(N, \varphi', \psi') \to H^2(K, 0, 0) = K
$$
而 $\alpha_1$ 是复合
$$
\alpha_1 : Q = H^1(Q, 0, 0) \to H^2(N, \varphi', \psi') \to H^3(K, 0, 0) = K
$$
它们来自引理中所示复形短正合列的边界映射。复形
$(M, \varphi, \psi)$、$(M', \varphi', \psi')$ 的正合性蕴含这些映射
都是同构。

\medskip\noindent
我们采用记号 $I_\varphi = \Im(\varphi)$、
$K_\varphi = \Ker(\varphi)$，其他映射的记号类似。$M$ 与 $M'$ 的正合性
意味着 $K_\varphi = I_\psi$ 及另外三个类似等式。记
$k = \text{length}_R(K)$、$a = \text{length}_R(I_\varphi)$、
$b = \text{length}_R(I_\psi)$。于是
$\text{length}_R(M) = a + b$、
$\text{length}_R(N) = a + b - k$、$\text{length}_R(Q) = k$，并且
$\text{length}_R(M') = a + b$。下面的正合列还将表明
$\text{length}_R(I_{\varphi'}) = a$ 与
$\text{length}_R(I_{\psi'}) = b$。

\medskip\noindent
假设 $K \subset K_\varphi = I_\psi$ 意味着 $\varphi$ 经 $N$ 分解，给出
正合列
$$
0 \to \alpha(I_\psi) \to N \xrightarrow{\varphi\alpha^{-1}} I_\psi \to 0.
$$
这里 $\varphi\alpha^{-1}(x') = y$ 意指 $x' = \alpha(x)$ 且
$y = \varphi(x)$。类似地，有
$$
0 \to \alpha(I_\varphi) \to N \xrightarrow{\psi\alpha^{-1}} I_\varphi \to 0.
$$
$\psi'$ 在 $Q$ 上诱导零映射这一假设意味着
$I_{\psi'} = K_{\varphi'} \subset N$。这意味着商
$\varphi'(N) \subset I_{\varphi'}$ 与 $Q$ 认同。注意
$\varphi'(N) = \alpha(I_\varphi)$。因此得到同构
$$
\varphi' : Q \to I_{\varphi'}/\alpha(I_\varphi)
$$
它可简单描述为
$\varphi'(x' \bmod N) = \varphi'(x') \bmod \alpha(I_\varphi)$。
完全类似地得到
$$
\psi' : Q \to I_{\psi'}/\alpha(I_\psi)
$$
最后，注意 $\alpha_0$ 是复合
$$
\xymatrix{
Q \ar[r]^-{\varphi'} &
I_{\varphi'}/\alpha(I_\varphi)
\ar[rrr]^-{\psi\alpha^{-1}|_{I_{\varphi'}/\alpha(I_\varphi)}} & & &
K
}
$$
而类似地，
$\alpha_1 = \varphi\alpha^{-1}|_{I_{\psi'}/\alpha(I_\psi)} \circ \psi'$。

\medskip\noindent
为缩短下面的公式，以下将写 $\alpha x$ 而不写 $\alpha(x)$。这不会引起混淆，
因为所有映射都用希腊字母表示，而元素用罗马字母表示。选取
\begin{enumerate}
\item 生成 $K$ 的容许序列 $z_1, \ldots, z_k \in K$；
\item 元素 $z'_i \in M$，使得 $\varphi z'_i = z_i$；
\item 元素 $z''_i \in M$，使得 $\psi z''_i = z_i$；
\item 元素 $x_{k + 1}, \ldots, x_a \in I_\varphi$，使得
$z_1, \ldots, z_k, x_{k + 1}, \ldots, x_a$ 是生成 $I_\varphi$ 的
容许序列；
\item 元素 $\tilde x_i \in M$，使得 $\varphi \tilde x_i = x_i$；
\item 元素 $y_{k + 1}, \ldots, y_b \in I_\psi$，使得
$z_1, \ldots, z_k, y_{k + 1}, \ldots, y_b$ 是生成 $I_\psi$ 的
容许序列；
\item 元素 $\tilde y_i \in M$，使得 $\psi \tilde y_i = y_i$；
\item 元素 $w_1, \ldots, w_k \in M'$，使得
$w_1 \bmod N, \ldots, w_k \bmod N$ 是 $Q$ 中生成 $Q$ 的容许序列。
\end{enumerate}
由注 \ref{chow-remark-more-elementary}，元素
$D = \det_\kappa(M, \varphi, \psi) \in \kappa^*$ 由下式刻画：
\begin{eqnarray*}
& &
[z_1, \ldots, z_k,
x_{k + 1}, \ldots, x_a,
z''_1, \ldots, z''_k,
\tilde y_{k + 1}, \ldots, \tilde y_b] \\
& = &
(-1)^{ab} D
[z_1, \ldots, z_k,
y_{k + 1}, \ldots, y_b,
z'_1, \ldots, z'_k,
\tilde x_{k + 1}, \ldots, \tilde x_a]
\end{eqnarray*}
由上面的讨论，
$\alpha x_{k + 1}, \ldots, \alpha x_a, \varphi w_1, \ldots, \varphi w_k$
是生成 $I_{\varphi'}$ 的容许序列，而
$\alpha y_{k + 1}, \ldots, \alpha y_b, \psi w_1, \ldots, \psi w_k$
是生成 $I_{\psi'}$ 的容许序列。因此由注
\ref{chow-remark-more-elementary}，元素
$D' = \det_\kappa(M', \varphi', \psi') \in \kappa^*$ 由下式刻画：
\begin{eqnarray*}
& &
[\alpha x_{k + 1}, \ldots, \alpha x_a,
\varphi' w_1, \ldots, \varphi' w_k,
\alpha \tilde y_{k + 1}, \ldots, \alpha \tilde y_b,
w_1, \ldots, w_k]
\\
& = &
(-1)^{ab} D'
[\alpha y_{k + 1}, \ldots, \alpha y_b,
\psi' w_1, \ldots, \psi' w_k,
\alpha \tilde x_{k + 1}, \ldots, \alpha \tilde x_a,
w_1, \ldots, w_k]
\end{eqnarray*}
注意，在第一个，resp.\ 第二个显示公式中，两边符号的最前，resp.\ 最后
$k$ 个分量相同。因此这些公式实际上等价于下列等式：
\begin{eqnarray*}
& &
[\alpha x_{k + 1}, \ldots, \alpha x_a,
\alpha z''_1, \ldots, \alpha z''_k,
\alpha \tilde y_{k + 1}, \ldots, \alpha \tilde y_b] \\
& = &
(-1)^{ab} D
[\alpha y_{k + 1}, \ldots, \alpha y_b,
\alpha z'_1, \ldots, \alpha z'_k,
\alpha \tilde x_{k + 1}, \ldots, \alpha \tilde x_a]
\end{eqnarray*}
以及
\begin{eqnarray*}
& &
[\alpha x_{k + 1}, \ldots, \alpha x_a,
\varphi' w_1, \ldots, \varphi' w_k,
\alpha \tilde y_{k + 1}, \ldots, \alpha \tilde y_b]
\\
& = &
(-1)^{ab} D'
[\alpha y_{k + 1}, \ldots, \alpha y_b,
\psi' w_1, \ldots, \psi' w_k,
\alpha \tilde x_{k + 1}, \ldots, \alpha \tilde x_a]
\end{eqnarray*}
这些等式在 $\det_\kappa(N)$ 中成立。注意，
$\varphi' w_1, \ldots, \varphi' w_k$ 与
$\alpha z''_1, \ldots, z''_k$ 是生成模
$I_{\varphi'}/\alpha(I_\varphi)$ 的容许序列。写成
$$
[\varphi' w_1, \ldots, \varphi' w_k]
= \lambda_0 [\alpha z''_1, \ldots, \alpha z''_k]
$$
在 $\det_\kappa(I_{\varphi'}/\alpha(I_\varphi))$ 中成立，其中
$\lambda_0 \in \kappa^*$。类似地，写成
$$
[\psi' w_1, \ldots, \psi' w_k]
= \lambda_1 [\alpha z'_1, \ldots, \alpha z'_k]
$$
在 $\det_\kappa(I_{\psi'}/\alpha(I_\psi))$ 中成立，其中
$\lambda_1 \in \kappa^*$。一方面，由上面对 $\alpha_i$ 的描述，显然有
$$
\alpha_i([w_1, \ldots, w_k]) = \lambda_i[z_1, \ldots, z_k]
$$
对 $i = 0, 1$ 成立。这意味着
$$
\det\nolimits_\kappa(\alpha_0^{-1} \circ \alpha_1)
=
\lambda_1/\lambda_0
$$
另一方面，显然有
\begin{eqnarray*}
& &
\lambda_0 [\alpha x_{k + 1}, \ldots, \alpha x_a,
\alpha z''_1, \ldots, \alpha z''_k,
\alpha \tilde y_{k + 1}, \ldots, \alpha \tilde y_b] \\
& = &
[\alpha x_{k + 1}, \ldots, \alpha x_a,
\varphi' w_1, \ldots, \varphi' w_k,
\alpha \tilde y_{k + 1}, \ldots, \alpha \tilde y_b]
\end{eqnarray*}
以及
\begin{eqnarray*}
& &
\lambda_1[\alpha y_{k + 1}, \ldots, \alpha y_b,
\alpha z'_1, \ldots, \alpha z'_k,
\alpha \tilde x_{k + 1}, \ldots, \alpha \tilde x_a] \\
& = &
[\alpha y_{k + 1}, \ldots, \alpha y_b,
\psi' w_1, \ldots, \psi' w_k,
\alpha \tilde x_{k + 1}, \ldots, \alpha \tilde x_a]
\end{eqnarray*}
这些等式蕴含 $\lambda_0 D = \lambda_1 D'$，引理得证。
\end{proof}








\subsection{符号}
\label{chow-subsection-symbols}

\noindent
这一构造适当的一般性也许由下列引理中的情形给出。

\begin{lemma}
\label{chow-lemma-pre-symbol}
设 $A$ 为诺特局部环，$M$ 为维数为 $1$ 的有限 $A$-模。假设
$\varphi, \psi : M \to M$ 是两个单射 $A$-模映射，并且
$\varphi(\psi(M)) = \psi(\varphi(M))$；例如，$\varphi$ 与 $\psi$
交换时即如此。则 $\text{length}_A(M/\varphi\psi M) < \infty$，而且
$(M/\varphi\psi M, \varphi, \psi)$ 是正合 $(2, 1)$-周期复形。
\end{lemma}

\begin{proof}
设 $\mathfrak q$ 为 $M$ 的支集的一个极小素理想。则
$M_{\mathfrak q}$ 是有限长度 $A_{\mathfrak q}$-模，见《Algebra》引理
\ref{algebra-lemma-support-point}。因此 $\varphi$ 与 $\psi$ 都诱导同构
$M_{\mathfrak q} \to M_{\mathfrak q}$。所以 $M/\varphi\psi M$ 的支集为
$\{\mathfrak m_A\}$，从而它具有有限长度（见上述引理）。最后，
$\varphi$ 在 $M/\varphi\psi M$ 上的核显然为
$\psi M/\varphi\psi M$，所以 $\varphi$ 的核就是 $\psi$ 在
$M/\varphi\psi M$ 上的像。反向情形类似，因为根据假设，
$M/\varphi\psi M = M/\psi\varphi M$。
\end{proof}

\begin{lemma}
\label{chow-lemma-symbol-defined}
设 $A$ 为诺特局部环，$a, b \in A$。
\begin{enumerate}
\item 若 $M$ 是维数为 $1$ 的有限 $A$-模，且 $a, b$ 在 $M$ 上均为
非零因子，则 $\text{length}_A(M/abM) < \infty$，而且
$(M/abM, a, b)$ 是正合 $(2, 1)$-周期复形；
\item 若 $a, b$ 是非零因子且 $\dim(A) = 1$，则
$\text{length}_A(A/(ab)) < \infty$，而且
$(A/(ab), a, b)$ 是正合 $(2, 1)$-周期复形。
\end{enumerate}
特别地，在这些情形，
$\det_\kappa(M/abM, a, b) \in \kappa^*$，
resp.\ $\det_\kappa(A/(ab), a, b) \in \kappa^*$ 均有定义。
\end{lemma}

\begin{proof}
由引理 \ref{chow-lemma-pre-symbol} 得出。
\end{proof}

\begin{definition}
\label{chow-definition-symbol-M}
设 $A$ 为剩余域为 $\kappa$ 的诺特局部环，$a, b \in A$。
设 $M$ 为维数为 $1$ 的有限 $A$-模，且 $a, b$ 在 $M$ 上均为非零因子。
把{\it 与 $M, a, b$ 关联的符号}定义为元素
$$
d_M(a, b) =
\det\nolimits_\kappa(M/abM, a, b) \in \kappa^*
$$
\end{definition}

\begin{lemma}
\label{chow-lemma-multiplicativity-symbol}
设 $A$ 为诺特局部环，$a, b, c \in A$。设 $M$ 为满足
$\dim(\text{Supp}(M)) = 1$ 的有限 $A$-模，并假设 $a, b, c$ 在 $M$
上均为非零因子。则
$$
d_M(a, bc) = d_M(a, b) d_M(a, c)
$$
且 $d_M(a, b)d_M(b, a) = 1$。
\end{lemma}

\begin{proof}
把引理 \ref{chow-lemma-multiplicativity-determinant} 应用于 $M/abcM$ 以及
分别由乘以 $a, b, c$ 给出的自同态 $\alpha, \beta, \gamma$，即得第一个
断言。第二个断言来自引理 \ref{chow-lemma-periodic-determinant-shift}。
\end{proof}

\begin{definition}
\label{chow-definition-tame-symbol}
设 $A$ 为剩余域为 $\kappa$ 的 $1$ 维诺特局部整环，$K$ 为 $A$ 的分式域。
把 $A$ 的{\it 驯符号}定义为映射
$$
K^* \times K^* \longrightarrow \kappa^*,
\quad
(x, y) \longmapsto d_A(x, y)
$$
其中利用引理 \ref{chow-lemma-multiplicativity-symbol} 的乘法性，把
$d_A(x, y)$ 延拓到 $K^* \times K^*$。
\end{definition}

\noindent
显然，可以更一般地把 $d_M(-, -)$ 延拓到 $A$ 的某些分式环
（即使 $A$ 不是整环）。

\begin{lemma}
\label{chow-lemma-symbol-when-equal}
设 $A$ 为诺特局部环，$M$ 为维数为 $1$ 的有限 $A$-模。设
$a \in A$ 在 $M$ 上为非零因子。则
$d_M(a, a) = (-1)^{\text{length}_A(M/aM)}$。
\end{lemma}

\begin{proof}
立即由引理 \ref{chow-lemma-periodic-determinant-sign} 得出。
\end{proof}

\begin{lemma}
\label{chow-lemma-symbol-when-one-is-a-unit}
设 $A$ 为诺特局部环，$M$ 为维数为 $1$ 的有限 $A$-模。设
$b \in A$ 在 $M$ 上为非零因子，并设 $u \in A^*$。则
$$
d_M(u, b) = u^{\text{length}_A(M/bM)} \bmod \mathfrak m_A.
$$
特别地，若 $M = A$，则
$d_A(u, b) = u^{\text{ord}_A(b)} \bmod \mathfrak m_A$。
\end{lemma}

\begin{proof}
注意，在此情形 $M/ubM = M/bM$，而乘以 $b$ 在其上为零。因此由引理
\ref{chow-lemma-periodic-determinant-easy-case}，
$d_M(u, b) = \det_\kappa(u|_{M/bM})$。再由引理
\ref{chow-lemma-times-u-determinant} 得出结论。
\end{proof}

\begin{lemma}
\label{chow-lemma-symbol-short-exact-sequence}
设 $A$ 为诺特局部环，$a, b \in A$。设
$$
0 \to M \to M' \to M'' \to 0
$$
为维数为 $1$ 的 $A$-模的短正合列，并且 $a, b$ 在这三个 $A$-模上
均为非零因子。则
$$
d_{M'}(a, b) = d_M(a, b) d_{M''}(a, b)
$$
在 $\kappa^*$ 中成立。
\end{lemma}

\begin{proof}
容易看出，这给出正合 $(2, 1)$-周期复形的短正合列
$$
0 \to
(M/abM, a, b) \to
(M'/abM', a, b) \to
(M''/abM'', a, b) \to 0
$$
因此结论由引理 \ref{chow-lemma-periodic-determinant} 得出。
\end{proof}

\begin{lemma}
\label{chow-lemma-symbol-compare-modules}
设 $A$ 为诺特局部环，$\alpha : M \to M'$ 为维数为 $1$ 的有限
$A$-模之间的同态。设 $a, b \in A$，并假设
\begin{enumerate}
\item $a$、$b$ 在 $M$ 与 $M'$ 上均为非零因子；
\item $\dim(\Ker(\alpha)), \dim(\Coker(\alpha)) \leq 0$。
\end{enumerate}
则 $d_M(a, b) = d_{M'}(a, b)$。
\end{lemma}

\begin{proof}
若 $a \in A^*$，则所需等式由等式
$\text{length}(M/bM) = \text{length}(M'/bM')$ 与引理
\ref{chow-lemma-symbol-when-one-is-a-unit} 得出。类似地，若 $b$ 是单位，
引理也成立（由引理 \ref{chow-lemma-multiplicativity-symbol} 的对称性）。
因此可假设 $a, b \in \mathfrak m_A$。这特别蕴含 $\mathfrak m$ 不是
$M$ 的伴随素理想，因而 $\alpha : M \to M'$ 是单射。于是可把 $M$
视为 $M'$ 的子模。根据假设，$M'/M$ 是支集为 $\{\mathfrak m_A\}$ 的
有限 $A$-模，因此具有有限长度。注意，对任意满足
$M \subset M'' \subset M'$ 的第三个模 $M''$，映射 $M \to M''$ 与
$M'' \to M'$ 也满足引理的假设。因此对 $M'/M$ 的长度归纳，可归结到
$\text{length}_A(M'/M) = 1$ 的情形。最后，在这种情形考虑映射
$$
\overline{\alpha} : M/abM \longrightarrow M'/abM'.
$$
按构造，$\overline{\alpha}$ 的余核 $Q$ 的长度为 $1$。由于
$a, b \in \mathfrak m_A$，它们在 $Q$ 上作用为零。还可推出
$\overline{\alpha}$ 的核 $K$ 长度为 $1$，因而 $a$、$b$ 在 $K$ 上也
作用为零。因此可以应用引理 \ref{chow-lemma-tricky}。所以只需看出两个映射
$\alpha_i : Q \to K$ 相同。事实上，它们都等于映射
$q = x' \bmod \Im(\overline{\alpha}) \mapsto abx' \in K$。
验证从略。
\end{proof}

\begin{lemma}
\label{chow-lemma-compute-symbol-M}
设 $A$ 为诺特局部环，$M$ 为满足 $\dim(\text{Supp}(M)) = 1$ 的有限
$A$-模。设 $a, b \in A$ 在 $M$ 上为非零因子，并设
$\mathfrak q_1, \ldots, \mathfrak q_t$ 为 $M$ 的支集中的极小素理想。
则在 $\kappa^*$ 中有
$$
d_M(a, b)
=
\prod\nolimits_{i = 1, \ldots, t}
d_{A/\mathfrak q_i}(a, b)^{
\text{length}_{A_{\mathfrak q_i}}(M_{\mathfrak q_i})}
$$
\end{lemma}

\begin{proof}
选取由 $A$-子模组成的滤过
$$
0 = M_0 \subset M_1 \subset \ldots \subset M_n = M
$$
使得每个商 $M_j/M_{j - 1}$ 都与 $A/\mathfrak p_j$ 同构，其中
$\mathfrak p_j$ 是 $A$ 的某个素理想。见《Algebra》引理
\ref{algebra-lemma-filter-Noetherian-module}。对每个 $j$，或者对某个 $i$
有 $\mathfrak p_j = \mathfrak q_i$，或者
$\mathfrak p_j = \mathfrak m_A$。此外，对固定的 $i$，满足
$\mathfrak p_j = \mathfrak q_i$ 的 $j$ 的数目等于
$\text{length}_{A_{\mathfrak q_i}}(M_{\mathfrak q_i})$，见《Algebra》引理
\ref{algebra-lemma-filter-minimal-primes-in-support}。因此对每个 $j$，
$d_{M_j}(a, b)$ 均有定义，而且
$$
d_{M_j}(a, b)
=
\left\{
\begin{matrix}
d_{M_{j - 1}}(a, b) d_{A/\mathfrak q_i}(a, b) &
\text{若} & \mathfrak p_j = \mathfrak q_i \\
d_{M_{j - 1}}(a, b) & \text{若} & \mathfrak p_j = \mathfrak m_A
\end{matrix}
\right.
$$
第一种情形使用引理 \ref{chow-lemma-symbol-short-exact-sequence}，第二种情形
使用引理 \ref{chow-lemma-symbol-compare-modules}。引理得证。
\end{proof}

\begin{lemma}
\label{chow-lemma-symbol-is-usual-tame-symbol}
设 $A$ 为分式域为 $K$ 的离散赋值环。对非零 $x, y \in K$，有
$$
d_A(x, y)
=
(-1)^{\text{ord}_A(x)\text{ord}_A(y)}
\frac{x^{\text{ord}_A(y)}}{y^{\text{ord}_A(x)}} \bmod \mathfrak m_A,
$$
换言之，该符号等于通常的驯符号。
\end{lemma}

\begin{proof}
由乘法性，只需证明 $x, y \in A$ 的情形。设 $t \in A$ 为一致化元。
把 $x = t^bu$ 与 $y = t^bv$ 写成这种形式，其中 $a, b \geq 0$ 且
$u, v \in A^*$。令 $l = a + b$。于是
$t^{l - 1}, \ldots, t^b$ 是 $(x)/(xy)$ 中的容许序列，而
$t^{l - 1}, \ldots, t^a$ 是 $(y)/(xy)$ 中的容许序列。因此由注
\ref{chow-remark-more-elementary}，$d_A(x, y)$ 由等式
$$
[t^{l - 1}, \ldots, t^b, v^{-1}t^{b - 1}, \ldots, v^{-1}]
=
(-1)^{ab} d_A(x, y)
[t^{l - 1}, \ldots, t^a, u^{-1}t^{a - 1}, \ldots, u^{-1}].
$$
刻画。于是由符号 $[x_1, \ldots, x_l]$ 的容许关系可见
$$
d_A(x, y) = (-1)^{ab} u^a/v^b \bmod \mathfrak m_A
$$
正如所需。
\end{proof}

\begin{lemma}
\label{chow-lemma-symbol-is-steinberg-prepare}
设 $A$ 为诺特局部环，$a, b \in A$。设 $M$ 为维数为 $1$ 的有限
$A$-模，并且 $a$、$b$、$b - a$ 在该模上均为非零因子。则
$$
d_M(a, b - a)d_M(b, b) = d_M(b, b - a)d_M(a, b)
$$
在 $\kappa^*$ 中成立。
\end{lemma}

\begin{proof}
由引理 \ref{chow-lemma-compute-symbol-M}，只需在
$M = A/\mathfrak q$ 的情形证明该关系，其中
$\mathfrak q \subset A$ 是满足 $\dim(A/\mathfrak q) = 1$ 的某个素理想。

\medskip\noindent
当 $M = A/\mathfrak q$ 时，可以把 $A$ 换成 $A/\mathfrak q$，并把
$a, b$ 换成它们在 $A/\mathfrak q$ 中的像。因此可假设 $A = M$，而
$A$ 是维数为 $1$ 的诺特局部整环。原因是 $A$ 与 $A/\mathfrak q$ 的
剩余域 $\kappa$ 相同，而且对任意 $A/\mathfrak q$-模 $M$，在 $A$ 上
或 $A/\mathfrak q$ 上取得的行列式典范地认同。见引理
\ref{chow-lemma-determinant-quotient-ring}。

\medskip\noindent
只需在 $a, b$ 都属于极大理想时证明该关系。事实上，其中一个或两个
都是单位的情形由引理 \ref{chow-lemma-symbol-when-one-is-a-unit} 与
\ref{chow-lemma-symbol-when-equal} 得出。

\medskip\noindent
按照引理 \ref{chow-lemma-Noetherian-domain-dim-1-two-elements}，选取扩张
$A \subset A'$ 与分解 $a = ta'$、$b = tb'$。注意还有
$b - a = t(b' - a')$，并且
$A' = (a', b') = (a', b' - a') = (b' - a', b')$。
这里以及下文把 $A'$ 视为 $A$-模，并把 $a, b, a', b', t$ 视为 $A'$
的 $A$-模自同态。我们用记号 $d^A_{A'}(a', b')$ 等表示
$$
d^A_{A'}(a', b')
=
\det\nolimits_\kappa(A'/a'b'A', a', b')
$$
它由引理 \ref{chow-lemma-pre-symbol} 定义。上标 ${}^A$ 用于把它与已经定义的
符号 $d_{A'}(a', b')$ 区分开；后者不同，例如其取值在 $A'$ 的剩余域中，
而该剩余域可能与 $\kappa$ 不同。由引理
\ref{chow-lemma-symbol-compare-modules} 可见
$d_A(a, b) = d^A_{A'}(a, b)$，其他组合也类似。利用这一点与乘法性，
只需证明
$$
d^A_{A'}(a', b' - a') d^A_{A'}(b', b')
=
d^A_{A'}(b', b' - a') d^A_{A'}(a', b')
$$
现在，由于 $(a', b') = A'$ 等，有
$$
\begin{matrix}
A'/(a'(b' - a')) & \cong & A'/(a') \oplus A'/(b' - a') \\
A'/(b'(b' - a')) & \cong & A'/(b') \oplus A'/(b' - a') \\
A'/(a'b') & \cong & A'/(a') \oplus A'/(b')
\end{matrix}
$$
此外，注意乘以 $b' - a'$ 在 $A/(a')$ 上等于乘以 $b'$，而乘以
$b' - a'$ 在 $A/(b')$ 上等于乘以 $-a'$。使用引理
\ref{chow-lemma-periodic-determinant-easy-case} 与
\ref{chow-lemma-periodic-determinant}，得到
$$
\begin{matrix}
d^A_{A'}(a', b' - a') & = &
\det\nolimits_\kappa(b'|_{A'/(a')})^{-1}
\det\nolimits_\kappa(a'|_{A'/(b' - a')}) \\
d^A_{A'}(b', b' - a') & = &
\det\nolimits_\kappa(-a'|_{A'/(b')})^{-1}
\det\nolimits_\kappa(b'|_{A'/(b' - a')}) \\
d^A_{A'}(a', b') & = &
\det\nolimits_\kappa(b'|_{A'/(a')})^{-1}
\det\nolimits_\kappa(a'|_{A'/(b')})
\end{matrix}
$$
因此
$$
(-1)^{\text{length}_A(A'/(b'))}
d^A_{A'}(a', b' - a')
=
d^A_{A'}(b', b' - a') d^A_{A'}(a', b')
$$
其中符号来自上述第二个等式中的 $-a'$。另一方面，由引理
\ref{chow-lemma-periodic-determinant-sign}，有
$d^A_{A'}(b', b') = (-1)^{\text{length}_A(A'/(b'))}$，引理得证。
\end{proof}

\noindent
驯符号是 Steinberg 符号。

\begin{lemma}
\label{chow-lemma-symbol-is-steinberg}
设 $A$ 为分式域为 $K$ 的维数为 $1$ 的诺特局部整环。对
$x \in K \setminus \{0, 1\}$，有
$$
d_A(x, 1 -x) = 1
$$
\end{lemma}

\begin{proof}
把 $x = a/b$ 写成这种形式，其中 $a, b \in A$。由于
$1 - x = (b - a)/b$，假设还蕴含 $b - a \not = 0$。因此计算得
$$
d_A(x, 1 - x)
=
d_A(a, b - a)d_A(a, b)^{-1}d_A(b, b - a)^{-1}d_A(b, b)
$$
所以必须证明
$d_A(a, b - a) d_A(b, b) = d_A(b, b - a) d_A(a, b)$。
这正是引理 \ref{chow-lemma-symbol-is-steinberg-prepare}。
\end{proof}










\subsection{长度与行列式}
\label{chow-subsection-length-determinant}

\noindent
本节用行列式比较格。关键引理如下。

\begin{lemma}
\label{chow-lemma-key-lemma}
设 $R$ 为诺特局部环。设 $\mathfrak q \subset R$ 为满足
$\dim(R/\mathfrak q) = 1$ 的素理想。设
$\varphi : M \to N$ 为有限 $R$-模的同态。假设存在
$x_1, \ldots, x_l \in M$ 与 $y_1, \ldots, y_l \in M$，使得
\begin{enumerate}
\item $M = \langle x_1, \ldots, x_l\rangle$,
\item 对 $i = 1, \ldots, l$，有
$\langle x_1, \ldots, x_i\rangle / \langle x_1, \ldots, x_{i - 1}\rangle
\cong R/\mathfrak q$,
\item $N = \langle y_1, \ldots, y_l\rangle$,
\item 对 $i = 1, \ldots, l$，有
$\langle y_1, \ldots, y_i\rangle / \langle y_1, \ldots, y_{i - 1}\rangle
\cong R/\mathfrak q$.
\end{enumerate}
则 $\varphi$ 为单射当且仅当 $\varphi_{\mathfrak q}$ 为同构，而且此时
$$
\text{length}_R(\Coker(\varphi)) = \text{ord}_{R/\mathfrak q}(f)
$$
其中 $f \in \kappa(\mathfrak q)$ 是满足
$$
[\varphi(x_1), \ldots, \varphi(x_l)] = f [y_1, \ldots, y_l]
$$
的元，上式是 $\det_{\kappa(\mathfrak q)}(N_{\mathfrak q})$ 中的等式。
\end{lemma}

\begin{proof}
首先注意，$l = 1$ 时引理成立。事实上，此时 $x_1$ 是 $M$ 在
$R/\mathfrak q$ 上的一组基，$y_1$ 是 $N$ 在 $R/\mathfrak q$ 上的一组基，
并且对某个 $f \in R$ 有 $\varphi(x_1) = fy_1$。因此 $\varphi$ 为单射
当且仅当 $f \not \in \mathfrak q$。再者，
$\Coker(\varphi) = R/(f, \mathfrak q)$，故由 $\text{ord}_{R/q}(f)$ 的定义
（见代数，定义 \ref{algebra-definition-ord}），引理成立。

\medskip\noindent
实际上，更一般地，假设对某些 $f_i \in R$，$f_i \not \in \mathfrak q$，有
$\varphi(x_i) = f_iy_i$。则诱导映射
$$
\langle x_1, \ldots, x_i\rangle / \langle x_1, \ldots, x_{i - 1}\rangle
\longrightarrow
\langle y_1, \ldots, y_i\rangle / \langle y_1, \ldots, y_{i - 1}\rangle
$$
都是单射，且其余核同构于 $R/(f_i, \mathfrak q)$。从而
$$
\text{length}_R(\Coker(\varphi)) = \sum \text{ord}_{R/\mathfrak q}(f_i).
$$
另一方面，由符号的容许关系 (b)，此时显然有
$$
[\varphi(x_1), \ldots, \varphi(x_l)] = f_1 \ldots f_l [y_1, \ldots, y_l]
$$
因此结论在这种情形也成立。

\medskip\noindent
一般情形对 $l$ 用归纳法证明。假设 $l > 1$。取最小的
$i \in \{1, \ldots, l\}$，使得
$\varphi(x_1) \in \langle y_1, \ldots, y_i\rangle$。我们对 $i$ 作归纳。
若 $i = 1$，则得到交换图
$$
\xymatrix{
0 \ar[r] &
\langle x_1 \rangle \ar[r] \ar[d] &
\langle x_1, \ldots, x_l \rangle \ar[r] \ar[d] &
\langle x_1, \ldots, x_l \rangle / \langle x_1 \rangle \ar[r] \ar[d] &
0 \\
0 \ar[r] &
\langle y_1 \rangle \ar[r] &
\langle y_1, \ldots, y_l \rangle \ar[r] &
\langle y_1, \ldots, y_l \rangle / \langle y_1 \rangle \ar[r] &
0
}
$$
此时引理由蛇形引理和对 $l$ 的归纳得出。现假设 $i > 1$。写成
$\varphi(x_1) = a_1 y_1 + \ldots + a_{i - 1} y_{i - 1} + a y_i$，
其中 $a_j, a \in R$ 且 $a \not \in \mathfrak q$（否则 $i$ 不是最小的）。令
$$
x'_j =
\left\{
\begin{matrix}
x_j & \text{若} & j = 1 \\
ax_j & \text{若} & j \geq 2
\end{matrix}
\right.
\quad\text{且}\quad
y'_j =
\left\{
\begin{matrix}
y_j & \text{若} & j < i \\
ay_j & \text{若} & j \geq i
\end{matrix}
\right.
$$
并令 $M' = \langle x'_1, \ldots, x'_l \rangle$、
$N' = \langle y'_1, \ldots, y'_l \rangle$。由构造有
$\varphi(x'_1) = a_1 y'_1 + \ldots + a_{i - 1} y'_{i - 1} + y'_i$，
而对 $j > 1$ 有
$\varphi(x'_j) = a\varphi(x_i) \in \langle y'_1, \ldots, y'_l\rangle$，
故得到如下 $R$-模与映射的交换图：
$$
\xymatrix{
M' \ar[d] \ar[r]_{\varphi'} & N' \ar[d] \\
M \ar[r]^\varphi & N
}
$$
由证明第二段的结果，我们知道
$\text{length}_R(M/M') = (l - 1)\text{ord}_{R/\mathfrak q}(a)$，并且类似地
$\text{length}_R(M/M') = (l - i + 1)\text{ord}_{R/\mathfrak q}(a)$。
追图给出
$$
\text{length}_R(\Coker(\varphi')) =
\text{length}_R(\Coker(\varphi)) + i\ \text{ord}_{R/\mathfrak q}(a).
$$
另一方面，写作
$$
[\varphi(x_1), \ldots, \varphi(x_l)] = f [y_1, \ldots, y_l],
\quad
[\varphi'(x'_1), \ldots, \varphi(x'_l)] = f' [y'_1, \ldots, y'_l]
$$
时，显然 $f' = a^if$。因此只需对
$\varphi(x_1) = a_1y_1 + \ldots a_{i - 1}y_{i - 1} + y_i$ 的情形，
即 $a = 1$ 的情形证明引理。其次，回忆由符号的容许关系有
$$
[y_1, \ldots, y_l] = [y_1, \ldots, y_{i - 1},
a_1y_1 + \ldots a_{i - 1}y_{i - 1} + y_i, y_{i + 1}, \ldots, y_l]
$$
序列
$y_1, \ldots, y_{i - 1},
a_1y_1 + \ldots + a_{i - 1}y_{i - 1} + y_i, y_{i + 1}, \ldots, y_l$
也满足引理的条件 (3)、(4)。因此实际上可以假设
$\varphi(x_1) = y_i$。在这种情形，注意有 $\mathfrak q x_1 = 0$，
从而也有 $\mathfrak q y_i = 0$。又由定义
$\det_{\kappa(\mathfrak q)}(N_{\mathfrak q})$ 的第三个容许关系，有
$$
[y_1, \ldots, y_l] =
- [y_1, \ldots, y_{i - 2}, y_i, y_{i - 1}, y_{i + 1}, \ldots, y_l]
$$
因此可以把 $y_1, \ldots, y_l$ 换成序列
$y'_1, \ldots, y'_l =
y_1, \ldots, y_{i - 2}, y_i, y_{i - 1}, y_{i + 1}, \ldots, y_l$
（它也满足引理的条件 (3) 和 (4)）。这显然使不变量 $i$ 减少 $1$，
从而对 $i$ 的归纳完成证明。
\end{proof}

\noindent
为了应用前一引理，我们证明具有所需性质的元序列往往存在。

\begin{lemma}
\label{chow-lemma-good-sequence-exists}
设 $R$ 为诺特局部环，$\mathfrak q \subset R$ 为素理想。设 $M$ 为有限
$R$-模，且 $\mathfrak q$ 是 $M$ 的支集的一个极小素理想。则存在
$x_1, \ldots, x_l \in M$，使得
\begin{enumerate}
\item $M / \langle x_1, \ldots, x_l\rangle$ 的支集不包含 $\mathfrak q$；
\item 对 $i = 1, \ldots, l$，有
$\langle x_1, \ldots, x_i\rangle / \langle x_1, \ldots, x_{i - 1}\rangle
\cong R/\mathfrak q$。
\end{enumerate}
而且此时 $l = \text{length}_{R_\mathfrak q}(M_\mathfrak q)$。
\end{lemma}

\begin{proof}
$\mathfrak q$ 是 $M$ 的支集中的极小素理想这一条件意味着
$l = \text{length}_{R_\mathfrak q}(M_\mathfrak q)$ 有限（见代数，引理
\ref{algebra-lemma-support-point}）。因此可以找到
$y_1, \ldots, y_l \in M_{\mathfrak q}$，使得对 $i = 1, \ldots, l$ 有
$\langle y_1, \ldots, y_i\rangle / \langle y_1, \ldots, y_{i - 1}\rangle
\cong \kappa(\mathfrak q)$。可以找到 $f_i \in R$、$f_i \not \in \mathfrak q$，
使 $f_i y_i$ 是某个元 $z_i \in M$ 的像。另外，由于 $R$ 是诺特的，
可对某些 $g_j \in R$ 写成 $\mathfrak q = (g_1, \ldots, g_t)$。由假设，
在模 $M_{\mathfrak q}$ 中有
$g_j y_i \in \langle y_1, \ldots, y_{i - 1} \rangle$。由 $z_i$ 的选取，
可以再找到 $f_{ji} \in R$、$f_{ij} \not \in \mathfrak q$，使得
$f_{ij} g_j z_i \in \langle z_1, \ldots, z_{i - 1} \rangle$
（这是模 $M$ 中的等式）。取
$$
x_1 = f_{11}f_{12}\ldots f_{1t}z_1,
\quad
x_2 = f_{11}f_{12}\ldots f_{1t}f_{21}f_{22}\ldots f_{2t}z_2,
$$
并依此类推，即得引理。事实上，所有元 $f_i, f_{ij}$ 在
$R_{\mathfrak q}$ 中都可逆，因而仍有
$R_{\mathfrak q}x_1 + \ldots + R_{\mathfrak q}x_i /
R_{\mathfrak q}x_1 + \ldots + R_{\mathfrak q}x_{i - 1}
\cong \kappa(\mathfrak q)$，其中 $i = 1, \ldots, l$。由构造，
$\mathfrak q x_i \in \langle x_1, \ldots, x_{i - 1}\rangle$。因此
$\langle x_1, \ldots, x_i\rangle / \langle x_1, \ldots, x_{i - 1}\rangle$
是一个由一个元生成且被 $\mathfrak q$ 消去的 $R$-模，将其在
$\mathfrak q$ 处局部化会得到一个在 $\kappa(\mathfrak q)$ 上的 $q$-维向量空间。
故它同构于 $R/\mathfrak q$。
\end{proof}

\noindent
以下是本节的主要结果。下面将看到这一命题的各种不同推论。
建议读者先自行证明较容易的引理
\ref{chow-lemma-application-herbrand-quotient}。

\begin{proposition}
\label{chow-proposition-length-determinant-periodic-complex}
设 $R$ 为剩余域为 $\kappa$ 的诺特局部环。假设
$(M, \varphi, \psi)$ 是 $R$ 上的一个 $(2, 1)$-周期复形。假设
\begin{enumerate}
\item $M$ 是有限 $R$-模；
\item $(M, \varphi, \psi)$ 的上同调模都具有有限长度；
\item $\dim(\text{Supp}(M)) = 1$。
\end{enumerate}
设 $\mathfrak q_i$、$i = 1, \ldots, t$ 是 $M$ 的支集的极小素理想。则\footnote{
显然，可以通过把 $\det_\kappa(M, \varphi, \psi)$ 重新定义为它当前值的逆
而去掉负号；见定义 \ref{chow-definition-periodic-determinant}。}
$$
- e_R(M, \varphi, \psi) =
\sum\nolimits_{i = 1, \ldots, t}
\text{ord}_{R/\mathfrak q_i}\left(
\det\nolimits_{\kappa(\mathfrak q_i)}
(M_{\mathfrak q_i}, \varphi_{\mathfrak q_i}, \psi_{\mathfrak q_i})
\right)
$$
\end{proposition}

\begin{proof}
首先如下化归到 $t = 1$ 的情形。注意
$\text{Supp}(M) = \{\mathfrak m, \mathfrak q_1, \ldots, \mathfrak q_t\}$，
其中 $\mathfrak m \subset R$ 是极大理想。令 $M_i$ 表示
$M \to M_{\mathfrak q_i}$ 的像，从而
$\text{Supp}(M_i) = \{\mathfrak m, \mathfrak q_i\}$。映射 $\varphi$（分别地，\ $\psi$）
诱导出 $R$-模映射 $\varphi_i : M_i \to M_i$（分别地，\ 
$\psi_i : M_i \to M_i$）。因此得到 $(2, 1)$-周期复形的态射
$$
(M, \varphi, \psi) \longrightarrow
\bigoplus\nolimits_{i = 1, \ldots, t} (M_i, \varphi_i, \psi_i).
$$
该映射的核与余核的支集都包含在 $\{\mathfrak m\}$ 中。故由引理
\ref{chow-lemma-compare-periodic-lengths}，
$$
e_R(M, \varphi, \psi) =
\sum\nolimits_{i = 1, \ldots, t}
e_R(M_i, \varphi_i, \psi_i)
$$
另一方面，显然有 $M_{\mathfrak q_i} = M_{i, \mathfrak q_i}$，因此引理公式
右边的各项等于
$$
\text{ord}_{R/\mathfrak q_i}\left(
\det\nolimits_{\kappa(\mathfrak q_i)}
(M_{i, \mathfrak q_i}, \varphi_{i, \mathfrak q_i}, \psi_{i, \mathfrak q_i})
\right)
$$
换言之，若能对每个模 $M_i$ 证明引理，则引理成立。这便化归到
$t = 1$ 的情形。

\medskip\noindent
假设在诺特局部环上有一个 $(2, 1)$-周期复形
$(M, \varphi, \psi)$，其中 $M$ 为有限 $R$-模，
$\text{Supp}(M) = \{\mathfrak m, \mathfrak q\}$，且上同调模具有有限长度。
在这种情形，证明由引理 \ref{chow-lemma-key-lemma} 和仔细的计数得出。记
$K_\varphi = \Ker(\varphi)$、
$I_\varphi = \Im(\varphi)$、
$K_\psi = \Ker(\psi)$ 以及
$I_\psi = \Im(\psi)$。由于 $R$ 是诺特的，它们都是有限 $R$-模。令
$$
a = \text{length}_{R_{\mathfrak q}}(I_{\varphi, \mathfrak q})
= \text{length}_{R_{\mathfrak q}}(K_{\psi, \mathfrak q}),
\quad
b = \text{length}_{R_{\mathfrak q}}(I_{\psi, \mathfrak q})
= \text{length}_{R_{\mathfrak q}}(K_{\varphi, \mathfrak q}).
$$
等式成立是因为复形在 $\mathfrak q$ 处局部化后变得正合。注意
$l = \text{length}_{R_{\mathfrak q}}(M_{\mathfrak q})$ 等于 $l = a + b$。

\medskip\noindent
我们将应用引理 \ref{chow-lemma-good-sequence-exists}，在支集包含于
$\{\mathfrak m, \mathfrak q\}$ 的有限 $R$-模 $N$ 中选取元序列。在这种情形，
$N_{\mathfrak q}$ 有有限长度，记为 $n \in \mathbf{N}$。我们把满足引理
\ref{chow-lemma-good-sequence-exists} 的性质 (1) 和 (2) 的序列
$w_1, \ldots, w_n \in N$ 称为“良序列”。注意，$N$ 对由良序列生成的子模所取的商
$N/\langle w_1, \ldots, w_n \rangle$ 的支集包含于 $\{\mathfrak m\}$，
因而它具有有限长度（代数，引理 \ref{algebra-lemma-support-point}）。再者，符号
$[w_1, \ldots, w_n] \in \det_{\kappa(\mathfrak q)}(N_{\mathfrak q})$
是一个生成元；见引理 \ref{chow-lemma-determinant-dimension-one}。

\medskip\noindent
在此约定下，选取良序列
$$
\begin{matrix}
x_1, \ldots, x_b & \text{于} & K_\varphi, &
t_1, \ldots, t_a & \text{于} & K_\psi, \\
y_1, \ldots, y_a & \text{于} & I_\varphi \cap \langle t_1, \ldots t_a\rangle, &
s_1, \ldots, s_b & \text{于} & I_\psi \cap \langle x_1, \ldots, x_b\rangle.
\end{matrix}
$$
我们将如下对这些选择作少许调整。选取 $y_i \in I_\varphi$ 在 M 中的提升
$\tilde y_i \in M$，以及 $s_i \in I_\psi$ 在 M 中的提升
$\tilde s_i \in M$。未必有
$\mathfrak q \tilde y_1 \subset \langle x_1, \ldots, x_b\rangle$，也未必有
$\mathfrak q \tilde s_1 \subset \langle t_1, \ldots, t_a\rangle$。然而，利用
$\mathfrak q$ 有限生成（如引理 \ref{chow-lemma-good-sequence-exists} 的证明中那样），
可以找到 $d \in R$、$d \not \in \mathfrak q$，使得
$\mathfrak q d\tilde y_1 \subset \langle x_1, \ldots, x_b\rangle$
且
$\mathfrak q d\tilde s_1 \subset \langle t_1, \ldots, t_a\rangle$。
因此，用 $dy_i$ 替换 $y_i$，用 $d\tilde y_i$ 替换 $\tilde y_i$，
用 $ds_i$ 替换 $s_i$，并用 $d\tilde s_i$ 替换 $\tilde s_i$ 后，还可以假设
$x_1, \ldots, x_b, \tilde y_1, \ldots, \tilde y_b$
与 $t_1, \ldots, t_a, \tilde s_1, \ldots, \tilde s_b$
都是 $M$ 中的良序列。

\medskip\noindent
最后，在有限 $R$-模
$$
\langle
x_1, \ldots, x_b, \tilde y_1, \ldots, \tilde y_a
\rangle
\cap
\langle
t_1, \ldots, t_a, \tilde s_1, \ldots, \tilde s_b
\rangle.
$$
中选取一个良序列 $z_1, \ldots, z_l$。注意，这也是 $M$ 中的良序列。

\medskip\noindent
由于 $I_{\varphi, \mathfrak q} = K_{\psi, \mathfrak q}$，存在唯一元
$h \in \kappa(\mathfrak q)$，使得在
$\det_{\kappa(\mathfrak q)}(K_{\psi, \mathfrak q})$ 中有
$[y_1, \ldots, y_a] = h [t_1, \ldots, t_a]$。类似地，由于
$I_{\psi, \mathfrak q} = K_{\varphi, \mathfrak q}$，存在唯一元
$h \in \kappa(\mathfrak q)$，使得在
$\det_{\kappa(\mathfrak q)}(K_{\varphi, \mathfrak q})$ 中有
$[s_1, \ldots, s_b] = g [x_1, \ldots, x_b]$。对 $M$ 中的三个良序列也可以这样做。
总之，得到恒等式
\begin{align*}
[y_1, \ldots, y_a]
& =
h [t_1, \ldots, t_a] \\
[s_1, \ldots, s_b]
& =
g [x_1, \ldots, x_b] \\
[z_1, \ldots, z_l]
& =
f_\varphi [x_1, \ldots, x_b, \tilde y_1, \ldots, \tilde y_a] \\
[z_1, \ldots, z_l]
& =
f_\psi [t_1, \ldots, t_a, \tilde s_1, \ldots, \tilde s_b]
\end{align*}
其中 $g, h, f_\varphi, f_\psi \in \kappa(\mathfrak q)$。

\medskip\noindent
记号已全部建立，现在计算
$\det_{\kappa(\mathfrak q)}(M, \varphi, \psi)$。考虑元
$[z_1, \ldots, z_l]$。在定义 \ref{chow-definition-periodic-determinant} 的映射
$\gamma_\psi \circ \sigma \circ \gamma_\varphi^{-1}$ 下，有
\begin{eqnarray*}
[z_1, \ldots, z_l] & = &
f_\varphi [x_1, \ldots, x_b, \tilde y_1, \ldots, \tilde y_a] \\
& \mapsto & f_\varphi [x_1, \ldots, x_b] \otimes [y_1, \ldots, y_a] \\
& \mapsto &
f_\varphi h/g [t_1, \ldots, t_a] \otimes [s_1, \ldots, s_b] \\
& \mapsto &
f_\varphi h/g [t_1, \ldots, t_a, \tilde s_1, \ldots, \tilde s_b] \\
& = &
f_\varphi h/f_\psi g [z_1, \ldots, z_l]
\end{eqnarray*}
这意味着，除去一个符号后，
$\det_{\kappa(\mathfrak q)}
(M_{\mathfrak q}, \varphi_{\mathfrak q}, \psi_{\mathfrak q})$
等于 $f_\varphi h/f_\psi g$。

\medskip\noindent
对以下量作简记：
\begin{eqnarray*}
k_\varphi & = & \text{length}_R(K_\varphi/\langle x_1, \ldots, x_b\rangle) \\
k_\psi    & = & \text{length}_R(K_\psi/\langle t_1, \ldots, t_a\rangle) \\
i_\varphi & = & \text{length}_R(I_\varphi/\langle y_1, \ldots, y_a\rangle) \\
i_\psi    & = & \text{length}_R(I_\psi/\langle s_1, \ldots, s_a\rangle) \\
m_\varphi & = & \text{length}_R(M/
\langle x_1, \ldots, x_b, \tilde y_1, \ldots, \tilde y_a\rangle) \\
m_\psi    & = & \text{length}_R(M/
\langle t_1, \ldots, t_a, \tilde s_1, \ldots, \tilde s_b\rangle) \\
\delta_\varphi & = & \text{length}_R(
\langle x_1, \ldots, x_b, \tilde y_1, \ldots, \tilde y_a\rangle
\langle z_1, \ldots, z_l\rangle) \\
\delta_\psi & = & \text{length}_R(
\langle t_1, \ldots, t_a, \tilde s_1, \ldots, \tilde s_b\rangle
\langle z_1, \ldots, z_l\rangle)
\end{eqnarray*}
使用正合列 $0 \to K_\varphi \to M \to I_\varphi \to 0$，得
$m_\varphi = k_\varphi + i_\varphi$。类似地，有
$m_\psi = k_\psi + i_\psi$。又有
$\delta_\varphi + m_\varphi = \delta_\psi + m_\psi$，因为两边都等于
$\langle z_1, \ldots, z_l \rangle$ 在 $M$ 中的余长度。最后，首次应用关键引理
\ref{chow-lemma-key-lemma} 得
$$
\delta_\varphi = \text{ord}_{R/\mathfrak q}(f_\varphi),
\quad
\delta_\psi = \text{ord}_{R/\mathfrak q}(f_\psi)
$$

\medskip\noindent
接下来计算周期复形的重数：
\begin{eqnarray*}
e_R(M, \varphi, \psi) & = &
\text{length}_R(K_\varphi/I_\psi) - \text{length}_R(K_\psi/I_\varphi) \\
& = &
\text{length}_R(
\langle x_1, \ldots, x_b\rangle/
\langle s_1, \ldots, s_b\rangle)
+ k_\varphi - i_\psi \\
& & -
\text{length}_R(
\langle t_1, \ldots, t_a\rangle/
\langle y_1, \ldots, y_a\rangle)
- k_\psi + i_\varphi \\
& = &
\text{ord}_{R/\mathfrak q}(g/h) + k_\varphi - i_\psi - k_\psi + i_\varphi \\
& = &
\text{ord}_{R/\mathfrak q}(g/h) + m_\varphi - m_\psi \\
& = &
\text{ord}_{R/\mathfrak q}(g/h) + \delta_\psi - \delta_\varphi \\
& = &
\text{ord}_{R/\mathfrak q}(f_\psi g/f_\varphi h)
\end{eqnarray*}
其中在第三个等式中两次用到了关键引理
\ref{chow-lemma-key-lemma}。根据前面对
$\det_{\kappa(\mathfrak q)}
(M_{\mathfrak q}, \varphi_{\mathfrak q}, \psi_{\mathfrak q})$ 的计算，命题得证。
\end{proof}

\noindent
在多数应用中，以下引理已经足够。

\begin{lemma}
\label{chow-lemma-application-herbrand-quotient}
设 $R$ 为极大理想为 $\mathfrak m$ 的诺特局部环。设 $M$ 为有限
$R$-模，$\psi : M \to M$ 为 $R$-模映射。假设
\begin{enumerate}
\item $\Ker(\psi)$ 与 $\Coker(\psi)$ 具有有限长度；
\item $\dim(\text{Supp}(M)) \leq 1$。
\end{enumerate}
写成
$\text{Supp}(M) = \{\mathfrak m, \mathfrak q_1, \ldots, \mathfrak q_t\}$，
并用 $f_i \in \kappa(\mathfrak q_i)^*$ 表示使得映射
$\det_{\kappa(\mathfrak q_i)}(\psi_{\mathfrak q_i}) :
\det_{\kappa(\mathfrak q_i)}(M_{\mathfrak q_i})
\to \det_{\kappa(\mathfrak q_i)}(M_{\mathfrak q_i})$
等于乘以 $f_i$ 的那个元。则
$$
\text{length}_R(\Coker(\psi))
-
\text{length}_R(\Ker(\psi))
=
\sum\nolimits_{i = 1, \ldots, t}
\text{ord}_{R/\mathfrak q_i}(f_i).
$$
\end{lemma}

\begin{proof}
回忆 $H^0(M, 0, \psi) = \Coker(\psi)$ 且
$H^1(M, 0, \psi) = \Ker(\psi)$；见定义
\ref{chow-definition-periodic-length} 之前的注记。联合命题
\ref{chow-proposition-length-determinant-periodic-complex} 与引理
\ref{chow-lemma-periodic-determinant-easy-case}，即得本引理。

\medskip\noindent
另一证明。如命题 \ref{chow-proposition-length-determinant-periodic-complex} 的证明那样，
化归到 $\text{Supp}(M) = \{\mathfrak m, \mathfrak q\}$ 的情形。然后直接联合引理
\ref{chow-lemma-key-lemma} 与 \ref{chow-lemma-good-sequence-exists}，证明命题
\ref{chow-proposition-length-determinant-periodic-complex} 的这一特殊情形。此时需要的计数少得多，
建议读者自行完成。细节从略。
\end{proof}




\subsection{关键引理的应用}
\label{chow-subsection-application-tame-symbol}

\noindent
本节应用上述结果，利用上面构造的驯符号 $d_A$ 证明关键引理
（引理 \ref{chow-lemma-milnor-gersten-low-degree}）的类似结果。其与 Milnor $K$-理论的关系见注
\ref{chow-remark-gersten-complex-milnor}。

\begin{lemma}[关键引理]
\label{chow-lemma-secondary-ramification}
\begin{reference}
当 $A$ 为优良环时，这是 \cite[Proposition 1]{Kato-Milnor-K}。
\end{reference}
设 $A$ 为分式域为 $K$ 的 $2$-维诺特局部整环。设 $f, g \in K^*$。设
$\mathfrak q_1, \ldots, \mathfrak q_t$ 是 $A$ 的高为 $1$ 的素理想 $\mathfrak q$，
且 $f$ 或 $g$ 中至少有一个不是 $A^*_{\mathfrak q}$ 的元。则
$$
\sum\nolimits_{i = 1, \ldots, t}
\text{ord}_{A/\mathfrak q_i}(d_{A_{\mathfrak q_i}}(f, g))
=
0
$$
也可以写成
$$
\sum\nolimits_{\text{高度}(\mathfrak q) = 1}
\text{ord}_{A/\mathfrak q}(d_{A_{\mathfrak q}}(f, g))
=
0
$$
因为对 $A$ 的任意高为一的素理想 $\mathfrak q$，若
$f, g \in A^*_{\mathfrak q}$，则由引理
\ref{chow-lemma-symbol-when-one-is-a-unit} 有 $d_{A_{\mathfrak q}}(f, g) = 1$。
\end{lemma}

\begin{proof}
驯符号 $d_{A_{\mathfrak q}}(f, g)$ 具有可加性（引理
\ref{chow-lemma-multiplicativity-symbol}），而阶函数
$\text{ord}_{A/\mathfrak q}$ 也具有可加性（代数，引理
\ref{algebra-lemma-ord-additive}），故只需在 $f = a \in A$ 且
$g = b \in A$ 时证明公式。在这种情形，必须证明
$$
\sum\nolimits_{\text{高度}(\mathfrak q) = 1}
\text{ord}_{A/\mathfrak q}(\det\nolimits_\kappa(A_{\mathfrak q}/(ab), a, b))
= 0
$$
由命题 \ref{chow-proposition-length-determinant-periodic-complex}，这等价于证明
$$
e_A(A/(ab), a, b) = 0.
$$
由于复形
$A/(ab) \xrightarrow{a} A/(ab) \xrightarrow{b} A/(ab) \xrightarrow{a} A/(ab)$
是正合的，结论得证。
\end{proof}





\section{附录 B：其他方法}
\label{chow-section-appendix-chow}

\noindent
在本附录中，我们首先简要尝试将正文材料与相干层的 $K$-理论联系起来。
特别地，我们说明与可逆模的 $c_1$ 作杯积如何与张量乘这个可逆模相关；
见引理 \ref{chow-lemma-coherent-sheaf-cap-c1}。这些材料显然非常有趣，值得作更详细、
更充分的阐述。


\subsection{有理等价与 K-群}
\label{chow-subsection-rational-equivalence-K-groups}

\noindent
本节是第 \ref{chow-section-chow-and-K} 节的继续。以下引理的动机来自同调，引理
\ref{homology-lemma-serre-subcategory-K-groups}。

\begin{lemma}
\label{chow-lemma-maps-between-coherent-sheaves}
设 $(S, \delta)$ 如情形 \ref{chow-situation-setup} 中所述。设 $X$ 为 $S$ 上局部有限型的概形，
$\mathcal{F}$ 为 $X$ 上的相干层。设
$$
\xymatrix{
\ldots \ar[r] &
\mathcal{F} \ar[r]^\varphi &
\mathcal{F} \ar[r]^\psi &
\mathcal{F} \ar[r]^\varphi &
\mathcal{F} \ar[r] & \ldots
}
$$
是同调，等式 (\ref{homology-equation-cyclic-complex}) 中的复形。假设
\begin{enumerate}
\item $\dim_\delta(\text{Supp}(\mathcal{F})) \leq k + 1$。
\item 对 $i = 0, 1$，有
$\dim_\delta(\text{Supp}(H^i(\mathcal{F}, \varphi, \psi))) \leq k$。
\end{enumerate}
则作为 $X$ 上的 $k$-维循环，有
$$
[H^0(\mathcal{F}, \varphi, \psi)]_k
\sim_{rat}
[H^1(\mathcal{F}, \varphi, \psi)]_k
$$
\end{lemma}

\begin{proof}
令 $\{W_j\}_{j \in J}$ 为 $\text{Supp}(\mathcal{F})$ 中具有 $\delta$-维数
$k + 1$ 的不可约分支所成的集合。由引理 \ref{chow-lemma-length-finite}，
$\{W_j\}$ 是 $X$ 的闭子集的局部有限族。对每个 $j$，令
$\xi_j \in W_j$ 为一般点，并令
$$
f_j = \det\nolimits_{\kappa(\xi_j)}
(\mathcal{F}_{\xi_j}, \varphi_{\xi_j}, \psi_{\xi_j})
\in
R(W_j)^*.
$$
记号见定义 \ref{chow-definition-periodic-determinant}。我们声称
$$
- [H^0(\mathcal{F}, \varphi, \psi)]_k + [H^1(\mathcal{F}, \varphi, \psi)]_k
=
\sum (W_j \to X)_*\text{div}(f_j)
$$
若证明这一等式，引理即随之得出。

\medskip\noindent
设 $Z \subset X$ 为 $\delta$-维数为 $k$ 的整闭子概形。要证明上述等式，
只需证明 $[Z]$ 在
$
[H^0(\mathcal{F}, \varphi, \psi)]_k - [H^1(\mathcal{F}, \varphi, \psi)]_k
$
中的系数 $n$ 等于 $[Z]$ 在
$
\sum (W_j \to X)_*\text{div}(f_j)
$
中的系数 $m$。令 $\xi \in Z$ 为一般点。考虑局部环
$A = \mathcal{O}_{X, \xi}$。把 $M = \mathcal{F}_\xi$ 视为 $A$-模，并用
$\varphi, \psi : M \to M$ 表示 $\varphi, \psi$ 在茎上的作用。由 $\xi \in Z$ 的选取，
有 $\delta(\xi) = k$，从而 $\dim(\text{Supp}(M)) = 1$。最后，经过 $\xi$ 的整闭子概形
$W_j$ 与 $\text{Supp}(M)$ 的极小素理想 $\mathfrak q_i$ 对应。在每种情形，
元 $f_j \in R(W_j)^*$ 对应于 $\kappa(\mathfrak q_i)^*$ 中的元
$\det_{\kappa(\mathfrak q_i)}(M_{\mathfrak q_i}, \varphi, \psi)$。因此
$$
n = - e_A(M, \varphi, \psi)
$$
且
$$
m =
\sum
\text{ord}_{A/\mathfrak q_i}
(\det\nolimits_{\kappa(\mathfrak q_i)}(M_{\mathfrak q_i}, \varphi, \psi))
$$
于是结论由命题 \ref{chow-proposition-length-determinant-periodic-complex} 得出。
\end{proof}

\begin{lemma}
\label{chow-lemma-cycles-rational-equivalence-K-group}
设 $(S, \delta)$ 如情形 \ref{chow-situation-setup} 中所述，$X$ 为 $S$ 上局部有限型的概形。
引理 \ref{chow-lemma-from-chow-to-K} 中的映射
$$
\CH_k(X) \longrightarrow
K_0(\textit{Coh}_{\leq k + 1}(X)/\textit{Coh}_{\leq k - 1}(X))
$$
诱导出从 $\CH_k(X)$ 到下列映射的像 $B_k(X)$ 上的双射：
$$
K_0(\textit{Coh}_{\leq k}(X)/\textit{Coh}_{\leq k - 1}(X))
\longrightarrow
K_0(\textit{Coh}_{\leq k + 1}(X)/\textit{Coh}_{\leq k - 1}(X)).
$$
\end{lemma}

\begin{proof}
由引理 \ref{chow-lemma-cycles-k-group}，我们有
$Z_k(X) = K_0(\textit{Coh}_{\leq k}(X)/\textit{Coh}_{\leq k - 1}(X))$，
且这一等同与引理 \ref{chow-lemma-from-chow-to-K} 的映射相容。因此，假设
$K_0(\textit{Coh}_{\leq k}(X)/\textit{Coh}_{\leq k - 1}(X))$ 中有一个元
$[A] - [B]$，它在 $B_k(X)$ 中映到零，即在
$K_0(\textit{Coh}_{\leq k + 1}(X)/\textit{Coh}_{\leq k - 1}(X))$ 中映到零。
我们必须证明 $[A] - [B]$ 对应于 $X$ 上与零有理等价的循环。假设
$[A] = [\mathcal{A}]$ 且 $[B] = [\mathcal{B}]$，其中 $\mathcal{A}, \mathcal{B}$ 是 $X$ 上的相干层，
支撑在 $\delta$-维数 $\leq k$ 的部分上。元 $[A] - [B]$ 在群
$K_0(\textit{Coh}_{\leq k + 1}(X)/\textit{Coh}_{\leq k - 1}(X))$ 中映到零这一假设意味着：
存在 $X$ 上支撑在 $\delta$-维数 $\leq k - 1$ 的部分上的相干层
$\mathcal{A}', \mathcal{B}'$，使得
$[\mathcal{A} \oplus \mathcal{A}'] - [\mathcal{B} \oplus \mathcal{B}']$ 在
$K_0(\textit{Coh}_{k + 1}(X))$ 中为零（使用同调，引理
\ref{homology-lemma-serre-subcategory-K-groups} 的 (1)）。由同调，引理
\ref{homology-lemma-serre-subcategory-K-groups} 的 (2)，这意味着在范畴
$\textit{Coh}_{\leq k + 1}(X)$ 中存在一个 $(2, 1)$-周期复形
$(\mathcal{F}, \varphi, \psi)$，使得
$\mathcal{A} \oplus \mathcal{A}' = H^0(\mathcal{F}, \varphi, \psi)$
且 $\mathcal{B} \oplus \mathcal{B}' = H^1(\mathcal{F}, \varphi, \psi)$。
由引理 \ref{chow-lemma-maps-between-coherent-sheaves}，这意味着
$$
[\mathcal{A} \oplus \mathcal{A}']_k
\sim_{rat}
[\mathcal{B} \oplus \mathcal{B}']_k
$$
这证明了，在引理 \ref{chow-lemma-cycles-k-group} 的映射
$$
K_0(\textit{Coh}_{\leq k}(X)/\textit{Coh}_{\leq k - 1}(X))
\longrightarrow Z_k(X)
$$
下，$[A] - [B]$ 映到一个与零有理等价的循环。这正是所需证明的，证明完成。
\end{proof}














\subsection{Cartier 除子与 K-群}
\label{chow-subsection-cartier-coherent}

\noindent
本节说明：与可逆层 $\mathcal{L}$ 的第一 Chern 类相交，如何对应于在
$K$-群中张量乘 $\mathcal{L} - \mathcal{O}$。

\begin{lemma}
\label{chow-lemma-no-embedded-points-modules}
设 $A$ 为诺特局部环，$M$ 为有限 $A$-模，$a, b \in A$。假设
\begin{enumerate}
\item $\dim(A) = 1$,
\item $a$ 和 $b$ 都是 $A$ 中的非零因子；
\item $A$ 没有嵌入素理想；
\item $M$ 没有嵌入伴随素理想；
\item $\text{Supp}(M) = \Spec(A)$。
\end{enumerate}
令 $I = \{x \in A \mid x(a/b) \in A\}$。设
$\mathfrak q_1, \ldots, \mathfrak q_t$ 为 $A$ 的极小素理想。则
$(a/b)IM \subset M$，且
$$
\text{length}_A(M/(a/b)IM)
-
\text{length}_A(M/IM)
=
\sum\nolimits_i
\text{length}_{A_{\mathfrak q_i}}(M_{\mathfrak q_i})
\text{ord}_{A/\mathfrak q_i}(a/b)
$$
\end{lemma}

\begin{proof}
因为 $M$ 没有嵌入伴随素理想，且 $M$ 的支集是 $\Spec(A)$，故
$\text{Ass}(M) = \{\mathfrak q_1, \ldots, \mathfrak q_t\}$。因此 $a$、$b$ 在 $M$ 上都是非零因子。注意
\begin{align*}
& \text{length}_A(M/(a/b)IM) \\
& = \text{length}_A(bM/aIM) \\
& = \text{length}_A(M/aIM)
-
\text{length}_A(M/bM) \\
& = \text{length}_A(M/aM) + \text{length}_A(aM/aIM) - \text{length}_A(M/bM) \\
& = \text{length}_A(M/aM) + \text{length}_A(M/IM) - \text{length}_A(M/bM)
\end{align*}
这是因为单射 $b : M \to bM$ 把 $(a/b)IM$ 映到 $aIM$，而单射
$a : M \to aM$ 把 $IM$ 映到 $aIM$。因此引理等式的左边等于
$$
\text{length}_A(M/aM) - \text{length}_A(M/bM).
$$
分别以 $x = a, b$ 应用引理 \ref{chow-lemma-additivity-divisors-restricted} 的第二个公式，
并使用代数，定义 \ref{algebra-definition-ord} 中 $\text{ord}$-函数的定义，即得结论。
\end{proof}

\begin{lemma}
\label{chow-lemma-no-embedded-points}
设 $(S, \delta)$ 如情形 \ref{chow-situation-setup} 中所述。设 $X$ 为 $S$ 上局部有限型的概形，
$\mathcal{L}$ 为可逆 $\mathcal{O}_X$-模，$\mathcal{F}$ 为相干 $\mathcal{O}_X$-模，而
$s \in \Gamma(X, \mathcal{K}_X(\mathcal{L}))$ 为 $\mathcal{L}$ 的一个亚纯截面。假设
\begin{enumerate}
\item $\dim_\delta(X) \leq k + 1$,
\item $X$ 没有嵌入点；
\item $\mathcal{F}$ 没有嵌入伴随点；
\item $\mathcal{F}$ 的支集是 $X$；
\item 截面 $s$ 是正则亚纯的。
\end{enumerate}
在这种情形，令 $\mathcal{I} \subset \mathcal{O}_X$ 为 $s$ 的分母理想；见除子，定义
\ref{divisors-definition-regular-meromorphic-ideal-denominators}。则：
\begin{enumerate}
\item 存在短正合列
$$
\begin{matrix}
0 &
\to &
\mathcal{I}\mathcal{F} &
\xrightarrow{1} &
\mathcal{F} &
\to &
\mathcal{Q}_1 &
\to &
0 \\
0 &
\to &
\mathcal{I}\mathcal{F} &
\xrightarrow{s} &
\mathcal{F} \otimes_{\mathcal{O}_X} \mathcal{L} &
\to &
\mathcal{Q}_2 &
\to &
0
\end{matrix}
$$
\item 相干层 $\mathcal{Q}_1$、$\mathcal{Q}_2$ 支撑在 $\delta$-维数 $\leq k$ 的部分上；
\item 在 $X$ 的每个 $\delta$-维数为 $k + 1$ 的不可约分支 $X_i$ 上，截面 $s$ 限制为
一个正则亚纯截面 $s_i$；
\item 写成 $[\mathcal{F}]_{k + 1} = \sum m_i[X_i]$，则在 $Z_k(X)$ 中
$$
[\mathcal{Q}_2]_k - [\mathcal{Q}_1]_k
=
\sum m_i(X_i \to X)_*\text{div}_{\mathcal{L}|_{X_i}}(s_i)
$$
特别地，在 $\CH_k(X)$ 中
$$
[\mathcal{Q}_2]_k - [\mathcal{Q}_1]_k
=
c_1(\mathcal{L}) \cap [\mathcal{F}]_{k + 1}
$$
\end{enumerate}
\end{lemma}

\begin{proof}
回忆除子，引理 \ref{divisors-lemma-make-maps-regular-section} 给出了单射
$1 : \mathcal{I}\mathcal{F} \to \mathcal{F}$ 与
$s : \mathcal{I}\mathcal{F} \to \mathcal{F} \otimes_{\mathcal{O}_X}\mathcal{L}$，它们的余核支撑在闭的稀疏子集 $T$ 上。
按引理的记号，用 $\mathcal{Q}_i$ 表示这些余核。从而
$\dim_\delta(\text{Supp}(\mathcal{Q}_i)) \leq k$。由除子，引理
\ref{divisors-lemma-pullback-meromorphic-sections-defined} 与
\ref{divisors-lemma-meromorphic-sections-pullback}，拉回 $s_i$ 有定义，且是
$\mathcal{L}|_{X_i}$ 的正则亚纯截面。(4) 中循环的等式意味着 (4) 中循环类的等式。
因此只剩下证明
$$
[\mathcal{Q}_2]_k - [\mathcal{Q}_1]_k
=
\sum m_i(X_i \to X)_*\text{div}_{\mathcal{L}|_{X_i}}(s_i)
$$
在 $Z_k(X)$ 中成立。为此，设 $Z \subset X$ 为 $\delta$-维数为 $k$ 的整闭子概形，
$\xi \in Z$ 为一般点。令 $A = \mathcal{O}_{X, \xi}$ 且 $M = \mathcal{F}_\xi$。再选取生成元
$s_\xi \in \mathcal{L}_\xi$。则可写成 $s = (a/b) s_\xi$，其中 $a, b \in A$ 是非零因子。
此时 $I = \mathcal{I}_\xi = \{x \in A \mid x(a/b) \in A\}$。左边中 $[Z]$ 的系数为
$$
\text{length}_A(M/(a/b)IM) - \text{length}_A(M/IM)
$$
而右边中 $[Z]$ 的系数为
$$
\sum
\text{length}_{A_{\mathfrak q_i}}(M_{\mathfrak q_i})
\text{ord}_{A/\mathfrak q_i}(a/b)
$$
其中 $\mathfrak q_1, \ldots, \mathfrak q_t$ 是 $1$-维局部环 $A$ 的极小素理想。因此结论由引理
\ref{chow-lemma-no-embedded-points-modules} 得出。
\end{proof}

\begin{lemma}
\label{chow-lemma-coherent-sheaf-cap-c1}
设 $(S, \delta)$ 如情形 \ref{chow-situation-setup} 中所述。设 $X$ 为 $S$ 上局部有限型的概形，
$\mathcal{L}$ 为可逆 $\mathcal{O}_X$-模，$\mathcal{F}$ 为相干 $\mathcal{O}_X$-模。假设
$\dim_\delta(\text{Supp}(\mathcal{F})) \leq k + 1$。则元
$$
[\mathcal{F} \otimes_{\mathcal{O}_X} \mathcal{L}]
-
[\mathcal{F}]
\in
K_0(\textit{Coh}_{\leq k + 1}(X)/\textit{Coh}_{\leq k - 1}(X))
$$
属于引理 \ref{chow-lemma-cycles-rational-equivalence-K-group} 中的子群 $B_k(X)$，并且在映射
$B_k(X) \to \CH_k(X)$ 下映到元 $c_1(\mathcal{L}) \cap [\mathcal{F}]_{k + 1}$。
\end{lemma}

\begin{proof}
设
$$
0 \to \mathcal{K} \to \mathcal{F} \to \mathcal{F}' \to 0
$$
是除子，引理 \ref{divisors-lemma-remove-embedded-points} 中构造的短正合列。特别地，
$\mathcal{F}'$ 没有嵌入伴随点。由于 $\mathcal{K}$ 的支集在 $\mathcal{F}$ 的支集中稀疏，有
$\dim_\delta(\text{Supp}(\mathcal{K})) \leq k$。从 $\mathcal{K}$ 出发，再次应用除子，引理
\ref{divisors-lemma-remove-embedded-points}，得到短正合列
$$
0 \to \mathcal{K}'' \to \mathcal{K} \to \mathcal{K}' \to 0
$$
其中 $\dim_\delta(\text{Supp}(\mathcal{K}'')) < k$，而 $\mathcal{K}'$ 没有嵌入伴随点。
假设能对相干层 $\mathcal{F}'$ 和 $\mathcal{K}'$ 证明引理。使用等式
$$
[\mathcal{F}]_{k + 1}
=
[\mathcal{F}']_{k + 1}
+ [\mathcal{K}']_{k + 1}
+ [\mathcal{K}'']_{k + 1}
$$
（使用引理 \ref{chow-lemma-additivity-sheaf-cycle}），以及
$$
[\mathcal{F} \otimes_{\mathcal{O}_X} \mathcal{L}]
-
[\mathcal{F}]
=
[\mathcal{F}' \otimes_{\mathcal{O}_X} \mathcal{L}]
-
[\mathcal{F}']
+
[\mathcal{K}' \otimes_{\mathcal{O}_X} \mathcal{L}]
-
[\mathcal{K}']
+
[\mathcal{K}'' \otimes_{\mathcal{O}_X} \mathcal{L}]
-
[\mathcal{K}'']
$$
（使用 $\otimes \mathcal{L}$ 的正合性），再利用 $[\mathcal{K}'']_{k + 1}$ 以及
$[\mathcal{K}'' \otimes_{\mathcal{O}_X} \mathcal{L}]
- [\mathcal{K}'']$ 在
$K_0(\textit{Coh}_{\leq k + 1}(X)/\textit{Coh}_{\leq k - 1}(X))$ 中显然为零，可见结论对
$\mathcal{F}$ 也成立。这说明可以假设层 $\mathcal{F}$ 没有嵌入伴随点。

\medskip\noindent
假设 $X$、$\mathcal{F}$ 如引理中所述，并进一步假设 $\mathcal{F}$ 没有嵌入伴随点。考虑除子，
引理 \ref{divisors-lemma-no-embedded-points-endos} 中构造的理想层
$\mathcal{I} \subset \mathcal{O}_X$、相应的闭子概形 $i : Z \to X$ 以及相干 $\mathcal{O}_Z$-模
$\mathcal{G}$。回忆 $Z$ 是一个没有嵌入点的局部诺特概形，$\mathcal{G}$ 是没有嵌入伴随点的相干层，
且 $\text{Supp}(\mathcal{G}) = Z$、$i_*\mathcal{G} = \mathcal{F}$。再令
$\mathcal{N} = \mathcal{L}|_Z$。

\medskip\noindent
由除子，引理 \ref{divisors-lemma-regular-meromorphic-section-exists}，可逆层 $\mathcal{N}$ 在 $Z$ 上有一个
正则亚纯截面 $s$。用 $\mathcal{J} \subset \mathcal{O}_Z$ 表示 $s$ 的分母层。由引理
\ref{chow-lemma-no-embedded-points}，存在短正合列
$$
\begin{matrix}
0 &
\to &
\mathcal{J}\mathcal{G} &
\xrightarrow{1} &
\mathcal{G} &
\to &
\mathcal{Q}_1 &
\to &
0 \\
0 &
\to &
\mathcal{J}\mathcal{G} &
\xrightarrow{s} &
\mathcal{G} \otimes_{\mathcal{O}_Z} \mathcal{N} &
\to &
\mathcal{Q}_2 &
\to &
0
\end{matrix}
$$
使得 $\dim_\delta(\text{Supp}(\mathcal{Q}_i)) \leq k$，且循环
$
[\mathcal{Q}_2]_k - [\mathcal{Q}_1]_k
$
是 $c_1(\mathcal{N}) \cap [\mathcal{G}]_{k + 1}$ 的一个代表。由投影公式（见上同调，引理
\ref{cohomology-lemma-projection-formula}）有
$i_*(\mathcal{G} \otimes \mathcal{N}) = \mathcal{F} \otimes \mathcal{L}$，从而
$$
[\mathcal{F} \otimes_{\mathcal{O}_X} \mathcal{L}]
-
[\mathcal{F}]
=
[i_*\mathcal{Q}_2] - [i_*\mathcal{Q}_1]
$$
在 $K_0(\textit{Coh}_{\leq k + 1}(X)/\textit{Coh}_{\leq k - 1}(X))$ 中成立。这已经说明
$[\mathcal{F} \otimes_{\mathcal{O}_X} \mathcal{L}] - [\mathcal{F}]$ 是 $B_k(X)$ 的元。再者，
\begin{eqnarray*}
[i_*\mathcal{Q}_2]_k - [i_*\mathcal{Q}_1]_k
& = &
i_*\left( [\mathcal{Q}_2]_k - [\mathcal{Q}_1]_k \right) \\
& = &
i_*\left(c_1(\mathcal{N}) \cap [\mathcal{G}]_{k + 1} \right) \\
& = &
c_1(\mathcal{L}) \cap i_*[\mathcal{G}]_{k + 1} \\
& = &
c_1(\mathcal{L}) \cap [\mathcal{F}]_{k + 1}
\end{eqnarray*}
这由上述结果及引理 \ref{chow-lemma-pushforward-cap-c1} 和
\ref{chow-lemma-cycle-push-sheaf} 得出。按定义，它与该元在
$B_k(X) \to \CH_k(X)$ 下的像一致。因此引理得证。
\end{proof}
